\documentclass[11pt]{article}
\usepackage[margin=1.02in]{geometry}
\usepackage{amsmath,amssymb,amsthm,mathtools}
\usepackage{microtype}
\usepackage{enumitem}
\usepackage{booktabs,longtable,array}
\usepackage{xcolor}
\usepackage[hidelinks]{hyperref}
\usepackage{lmodern}

\definecolor{deepblue}{RGB}{31,63,104}
\definecolor{softgray}{RGB}{245,247,250}
\newtheorem{theorem}{Theorem}[section]
\newtheorem{lemma}[theorem]{Lemma}
\newtheorem{proposition}[theorem]{Proposition}
\newtheorem{corollary}[theorem]{Corollary}
\newtheorem{conjecture}[theorem]{Conjecture}
\theoremstyle{remark}
\newtheorem{remark}[theorem]{Remark}
\newcommand{\norm}[1]{\left\lVert #1\right\rVert}
\newcommand{\ip}[2]{\left\langle #1,#2\right\rangle}


\usepackage{mathrsfs}
\hypersetup{pdftitle={Fixed-Space Prime Compatibility in Burnol's Sonine Quotients -- v1.64},pdfauthor={Alexander Eastwood}}
\newtheorem{assumption}[theorem]{Assumption}
\newtheorem{conditionalresult}[theorem]{Conditional result}
\newcommand{\Q}{\mathcal Q}
\newcommand{\Jc}{J^{\mathrm c}}
\newcommand{\Ju}{J^{\mathrm u}}
\newcommand{\Jr}{\widetilde J}
\newcommand{\E}{\mathscr E}
\newcommand{\Z}{\mathcal Z}
\newcommand{\dd}{\mathrm d}
\setlength{\emergencystretch}{2em}
\setcounter{tocdepth}{1}
\title{\textbf{Fixed-Space Prime Compatibility in Burnol's Sonine Quotients}\\[5pt]
\large Complete working manuscript, version 1.64\\
\normalsize Finite sampling, endpoint recovery, and the remaining graph defect}
\author{Alexander Eastwood}
\date{September 23, 2026}
\begin{document}
\maketitle
\begin{abstract}
We study fixed-space prime actions on Burnol's Sonine quotients and
finite Weil forms. Source covariance, compact Mellin division and
Fredholm reduction identify arithmetic boundary and range defects.
For the repaired fixed-order prolate source, tail, endpoint, correlation
and domain estimates establish weak G1. Its polynomial Fourier
projection recovers the physical endpoint and has an exponentially
small ordinary operator residual. These estimates do not identify
the lowest eigenspace or control a growing source block.

The centered Hardy sampler with an endpoint shell preserves ordinary
Gram and Weil-form limits while retaining an explicit graph defect.
Reflected tests and sharp cuts expose endpoint and complex-strip
obstructions. A finite-core metric criterion bypasses the global
arithmetic operator realization. Full-space projection preserves
strip-product bounds and Weil transfer, but its signed commutator
remains uncontrolled.

For the canonical semilocal Weil operator, signed Schur certificates
prove complete positivity and ground ordering at $\lambda=3$.
At $\lambda=4$ the entire tail $|n|>16$ is bounded below by its
positive archimedean energy. A structured inverse estimate retains
this certified tail metric in the low-mode correction, with complete
remote-row and dimension-aware error bounds. Finite-support trials
certify a zero head-error term at this window. Parity signs sharpen the complete inverse
iteration, with finite-prefix bounds retaining all remote rows.
Complete mixed-term estimates rule out the first inverse degree
with the saved head certificate, including its positive update.
Retaining that update gives a monotone family of complete inverse
upper bounds for subsequent degrees. Energy-scaled frozen trials now certify positivity of the entire even
form at this window, with all residual cross-Grams and every omitted
row retained. A complete odd-complement certificate and a repaired
single direction now prove positivity of the entire form at
$\lambda=4$, with exact fixed-window error $\varepsilon_4=0$. A complete shifted
Schur certificate now proves that its ground state is simple and even,
with $0<\mu_0<2.454\cdot10^{-75}$ and $\mu_1>10^{-73}$.
Connes--van Suijlekom's theorem then gives only real zeros for the
entire Fourier transform of this complete ground. This is a fixed-window
statement, with no convergence to Riemann's $\Xi$ asserted. The present scalar far
majorant requires exponentially growing cutoffs as the window grows;
polynomial-rank repairs of that unsigned comparison cannot remove
this cost. Historical computations report failure of the unchanged
tail metric at $\lambda=5$ in both parities; the original witnesses
and verifiers are not archived, so that historical report is not presently
replayable. Version 1.52 supplies new frozen witnesses and fresh two-precision
certificates for this same stipulated comparison, with positive Weil energies. Source-$L^2$ transport through the fixed point-value repair
is nonclosable, so physical concentration blocks require explicit
finite image Grams. Signed infinite block metrics provide a sufficient
alternative. A second historical computation reports that at $\lambda=5$
its first four dyadic blocks above $N=\lambda^2$ violate the criterion
in both parities, including all admissible block metrics and positive
row reweightings. Its original evidence is not archived either. Version 1.52 supplies a new
certificate using only the first three blocks of that same partition;
it does not certify the historical four-block tables.
The reported obstruction concerns the selected block comparison;
other signed estimates and growing-window positivity remain open.
An exact source-weighted jump transform now gives a project-new
candidate for the remaining complement estimate.  In the even sector,
the archimedean and prime terms form a nonnegative jump energy plus a
scalar potential, while the pole is one positive square.  Conjugation
by a positive source converts its orthogonal complement into a
weighted constraint and reduces G2.6 to a nonlocal
capacity inequality.  The reduction is proved; the capacity inequality
and its cofinal decay are not.
The PNT continuum model of the prime correlation has kernel
$e^{|x-y|/2}$. Its positive spectral channel is cancelled exactly by
the existing even pole: the difference is a strictly positive Green
operator with no uniform gap. Thus the even arithmetic obstruction is
isolated in an explicit prime-discrepancy multiplier. Its required
one-sided form estimate is unproved. A source-compressed
time--frequency hierarchy now gives exact sufficient concentration
tests for that multiplier. Its scalar trace and generic thick-set
specializations lose the favorable spectral mass and do not close the
estimate, while the full weighted operator test survives explicit
compactly supported wave-packet falsification attempts reported in v1.31
(the original scripts and outputs are not archived; only a report survives).
An exact source-compressed Hilbert--Schmidt identity sharpens the
negative-level trace bound, but a scale-invariant Dirichlet model proves
that this entire negative-only Schatten hierarchy need not vanish even
when the form is nonnegative and the source is an exact null vector.
A second no-go shows that no nontrivial positive frequency mollifier can
be form-exact on a physical interval: its characteristic function would
have to equal one on the full difference interval.  These results leave
the signed weighted concentration operator, rather than a negative-only
or positive-smoothing surrogate, as the live even-sector mechanism.
An exact Toeplitz--leakage identity now retains every cross-channel and
rank-one source term in products of the concentration operators.  It
also gives a cofinal obstruction: two nested reciprocal-scale bands have
source-compressed commutators bounded away from zero even for the actual
repaired source.  Thus nesting, physical scaling and source compression
do not yield generic approximate simultaneous diagonalization.  The
arithmetic signed hierarchy remains possible, but it still requires a
joint weighted coverage estimate rather than commutator control alone.
A direct Cotlar--Stein regrouping of the joint channels is now excluded
as a new mechanism: atomization leaves the complete channel Gram
unchanged, and disjoint reciprocal-scale channels retain a nonzero
cofinal cross norm after removal of the actual source.  A distinct
signed-primitive transport estimate is proved and optimized exactly.
Its strongest scalar error is $a\Delta_a$, where $\Delta_a$ is the
largest negative interval integral of the exact symbol.  A positive
Dirichlet model with an exact null source has $a\Delta_a=\pi^3/6$;
historical numerical arithmetic stress tests also reported no decay,
but their original script and output are not archived.  The method
is therefore recorded as a rejected scalar candidate, not as the live
G2 mechanism.
The same continuum symbol now represents both complex parity sectors.
A uniformly conditioned lattice of translated Gaussian radical sources
gives a growing physical block of dimension proportional to $\log\lambda$
with full operator residual at most $C\log\lambda\,e^{-\lambda/100}$.
Its odd subblock lies in every complement formed from even image columns.
Thus a uniform or polynomial positive gap on that complement is impossible;
the signs of the crowded near-zero eigenvalues remain undetermined.
The full translate span of the same radical is now shown to be dense
in ordinary $L^2$. The complete zero-extended Weil operators consequently
converge strong-resolvent to zero, a topology which does not control their
lower spectral edges. More usefully, a negative compact test could be
amplified to arbitrarily negative normalized tests. Therefore one uniform
finite lower floor on a cofinal family already forces complete Weil
positivity: the complementary error need only stay bounded, not tend to
zero. This sharper reduction leaves its arithmetic hypothesis unproved.
Uniformly bounded positive comparisons with vanishing error must themselves
disappear strongly on the actual source complement.
The relaxed bounded-error objective does not revive the scalar
primitive route: a new endpoint-vanishing probe proves, for the exact
arithmetic symbol, that its optimal error $a\Delta(\beta_a)$ tends to
infinity. The scalar symbol's infimum also tends to minus infinity.
These unconditional no-go statements leave the physical signed
concentration estimate open; the probe is not a physical squared
Paley--Wiener transform.
A fixed additive shift now gives complete ordinary lower floors
$\mathsf W_\lambda\succeq-8I$ at $\lambda=5,6,8$, in both
parities with the same moderate cutoffs and complete remote enclosures.
The corresponding trial upper bounds are below $10^{-16}$. This finite
experiment is inexpensive, but a bounded shift retains the exponential
cost of the present scalar far comparison. No cofinal floor follows.
Two precise meta-obstructions identify additional information that a
successful cofinal estimate must retain. The optimal full-distribution
mass-cap relaxation, and even the displayed mass/height/variation
relaxation, have divergent negative budgets for the actual symbol.
A separate domain-preserving unitary construction shows that exact
finite head and source-action data do not rescue a bound blind to
spectral orientation. These are scoped impossibility theorems, not a
claim that all recorded closures share one hypothesis.
Strong cancellation persists in the exact
source complement. Every complete eigenfunction has zero physical
endpoint. A cofinal complete-Weil lower error tending to zero would
imply RH directly, but remains unproved.
\end{abstract}
\paragraph{Evidence availability (NS-38, September 22, 2026).}
The original $\lambda=5$ computational evidence (v1.27--v1.28) and
original concentration-era companion evidence (v1.31--v1.34) are
\emph{not archived}, except for one recovered v1.31 report. The v1.29
capacity and v1.30 continuum companion artifacts are also not archived.
Affected numerical assertions and archival descriptions are explicitly
marked below; historical verification accounts are not currently
replayable certificates. Complete analytic proofs printed here remain
available independently of missing companion reviews. The scientific
v1.35 files are available under \path{evidence/v135/} and match their
recorded hashes. See \path{evidence/MISSING.md} for per-group disposition.
This disclosure alone does not change the manuscript version. The new
v1.52 certificates resolve the two specified comparison questions using
new witnesses; they do not recover the original evidence.

\clearpage
\paragraph{Status of this version.}
The latest checkpoint is v1.64 (NS-81), stated below; earlier version
paragraphs record the historical scope and remain part of the audit trail.

Version 1.58 supplies an unconditional adjacent-difference block bound
for the Nyman--Beurling problem, retaining the full projection and every
cross term. The raw operator norm is $O_\varepsilon(N^{-5/3+\varepsilon})$.
A bound $O_\varepsilon(N^{-2+\varepsilon})$ for all exponents is equivalent
to Lindel\"of, and every fixed logarithmic $N^{-2}$ raw norm bound is false.
These statements do not apply automatically to the projected Gram or its
actual residual direction. Two-precision finite checks expose the loss in
the old trace bound and the cost of transforming the numerator. The
asymptotic lower correlation input is still missing; RH and G2 remain open.
See Section~\ref{sec:v158-nb-difference}.

Version 1.57 completes the bounded translated-radical boundary test.
The full exterior pairing is defined on the complete form domain, with
both poles and every prime row retained. Its vanishing for every translate
is exactly equivalent to the original nullvector equation. Finite
prime-weight changes give an explicit forcing profile, necessarily nonzero
for nonzero tests; compactness supplies no uniform inverse estimate on the
whole fixed-window space. No independent arithmetic sign or uniqueness
estimate emerged. NS-58 stops at its stated condition; API, the uniform
floor, G2 and RH remain open. Section~\ref{sec:ns58-boundary} gives the
precise identities and scope.

Version 1.56 adds complete-operator controls for proposed uniqueness
arguments. Every nonzero finite change to prime-shift coefficients gives
negative compact tests in both parities and a finite first-zero window,
while preserving translation, reflection, support consistency and an
appropriately modified affine-potential equation. The original radical
identity is not preserved. For a fixed perturbation pattern of size
$\epsilon$, the first-zero parameter has an analytic upper bound
$\lambda_*=O(\sqrt{\log(1/\epsilon)})$ with pattern-dependent constants;
no numerical threshold is supplied. A fixed bounded perturbation preserves
the uniform-floor objective. Exact scalar compensation and an RH-conditional
perturbed spectral limit follow from the existing dense-radical and floor
lemmas. These controls do not give a negative original Weil direction or
an arithmetic sign estimate. API, G2 and RH remain open.

Version 1.55 tests the complete arithmetic kernel equation. Local open-set
unique continuation fails because a prime shift can cancel the integral
response on an empty patch; the construction works in both parities and
survives finitely many source constraints. The full semigroup also fails
to preserve ordinary pointwise positivity. These are not whole-window
nullvectors or negative diagonal Weil directions. First-zero support
saturation and an exact inward-dilation identity leave a signed arithmetic
overlap estimate missing. API, the uniform floor, G2 and RH remain open.
Section~\ref{sec:v155-local-uniqueness} states the precise scopes.

Version 1.54 supplies a complete-domain affine screw-potential equation
for nullvectors and a rank-one countermodel to a generic $H_0^1$ bootstrap.
Paired heat-zero transport removes internal inverse-gap singularities but
retains the external-zero interaction; a same-test heat comparison has an
explicit unmatched-zero obstruction. Exact Nyman--Beurling block gains admit
an unconditional trace lower bound. Arithmetic injectivity, full transport
and asymptotic correlation estimates remain open. These are scoped identities
and implications, not a uniform floor, G2 or RH proof.
Section~\ref{sec:v154-new-shapes} records the hypotheses and proofs.

Version 1.53 proves a larger complete near-zero block, jointly with
Astra-2 (NS-46). Fixed-range, increasingly close translates have dimension
$2\lfloor\lambda^2/(40000\log\lambda)\rfloor+1$, full residual at most
$C\exp(-\lambda^2/2000)$, and explicit even/odd dimensions. The finite
Gram lower bound deteriorates with the spacing and is retained throughout.
This strengthens the necessary removal rank for a uniform or polynomial
positive-gap strategy
to order $\lambda^2/\log\lambda$. Neither the energy signs nor a lower
bound on the remaining complete complement are supplied. The threshold
and constant are analytic existence statements, not certified finite-window
numbers. Section~\ref{sec:v153-fine-block} contains the full proof.
No new numerical window, tail metric, G2 or RH claim.

Version 1.52 adds fresh certificates for failure of the existing
$\lambda=5$ stipulated $10^{-8}D$ tail comparison and the existing
$N=25$ dyadic block-norm comparison, in both parities
(Propositions~\ref{prop:v152-tail-comparison} and
\ref{prop:v152-block-comparison}). The witnesses are new; the historical
v1.27--v1.28 artifacts remain unarchived. Their Weil energies are positive,
so no negative Weil direction is produced. The analytic results extend
the Gaussian prime-channel loss obstruction, prove completeness of the
unrestricted exact adaptive hierarchy as an equivalent floor target,
and give precise complete CCM selection criteria and scoped abstract
nonselection obstructions. A particular restricted bump-discrepancy
estimate already carries the full off-line-zero obstruction; that
estimate is not proved. Section~\ref{sec:v152-route-selection} states
all hypotheses and limitations. No new window, tail metric, cofinal
signed floor, G2 or RH result is asserted.

Version 1.51 compares the three open inputs in one proposition,
\ref{prop:v151-input-comparison}. The positive-weight potential is an
inverse multiplier of the complete corrected zero field, with the finite
window pole mass and exterior contribution retained. ZLD implies the
quantization-growth clause of QG on the same prescribed family, with
lower bound $\kappa\theta^2D_a/(2K)$; the separate packet input is not
supplied. CAE is exactly the uniform complement-floor target in these
coordinates, not a new estimate. General measure countermodels do not
prove arithmetic independence. Universal-head QG is false by cutoff
dilution, whether or not CAE holds; no converse for a specified natural
rule is established. All three arithmetic inputs remain unproved.
NS-38 separately makes the missing historical evidence explicit, without
claiming recovery. No new window, tail metric, G2 or RH claim is made.

Version 1.50 tests the explicit positive Gaussian arithmetic radical in
the physical prime-shift capacity identity. Its exact transformed form
retains both poles and exterior edges. A scoped obstruction proves that
no fixed fraction of its positive transformed energy, or its prime-2
channel alone, can be discarded with a uniform ordinary-norm error.
The full critical energy comparison remains an equivalent named open
input for the bounded-floor target, not a new sign theorem. No new
window, tail metric or numerical experiment. G2 and RH remain open.


Version 1.49 records the review of the unmerged v1.48 candidate and
replaces its peak-density interpretation by an exact quantization input.
The level density has an integrable unbounded singularity at any negative
critical value in the head range; a histogram peak is not its supremum.
The weighted-cost implication survives without a bounded-density assumption,
but one-level cost growth alone does not control higher fixed level counts.
Exact quantization growth and an admissible bounded-energy weighted packet
remain named open inputs. Diagnostic feasibility and projection limitations
are explicit. The route remains blocked; G2 and RH remain open.

Version 1.47 separates the arithmetic level-band measure from the
Fourier coverage and original energy of an admissible packet. Cumulative
sublevel estimates do not imply density on every interval of levels.
The exact zero formula gives the named open input ZLD, not an
unconditional value-distribution theorem. Local depth-to-width bounds and
a broad-packet energy calculation show what the existing arguments do
supply and where uniformity is missing. Both the arithmetic distribution
and bounded-energy packet coverage remain open; the minorant route stays
blocked, not closed. G2 and RH remain open.

Version 1.46 records the exact explicit-formula identity for the actual
discrepancy symbol, including the strict prime-power endpoint and the
Perron remainder. The lattice comb has a varying envelope, not exact
equal crests or quarter-point zeros. NS-26 supplies an identity and a
numerical illustration. NS-27 proves local packet bounds and a conditional
level-count inequality, but the required arithmetic level distribution and
bounded packet energy are not supplied by trough-depth growth. The proposed
bounded-complexity closure is therefore not proved; the route remains
blocked, and the weighted criterion remains valid. G2 and RH remain open.

Version 1.45 makes only the v1.43 audit's MINOR M1/M2 wording corrections:
the first-zero proposition now distinguishes new $1024/1280$-bit gates
from inherited $320/768/896$-bit ingredients, and the limitation of the
direct bound is paired with the separate transfer enclosure's frozen-trial
floor, approximately $6.65\cdot10^{-38}$. No bound is sharpened. G2 and RH
remain open.

Version 1.44 records the exact pencil-direction demand of the weighted
concentration criterion (Proposition~\ref{prop:v144-pencil-concentration}).
The admissible directional loss is at most
$\nu_kq_a^+[v_k]+\eta_a\norm{v_k}^2$; mixtures also require the full
restricted matrix inequality. The diagnostic uncompressed head is not
silently identified with the actual source complement. These are
fixed-window statements; the live sufficient criterion still requires
complete-space cofinal control, with bounded errors already sufficient
under the earlier hypotheses. No new numerical bound or refinement of
v1.43 is made. G2 and RH remain open.

Version 1.43 closes NS-19 with a complete first-positive-zero enclosure
$|z_1(4)-\gamma_1|<8.752082\cdot10^{-33}<9\cdot10^{-33}$, verified
at $1024/1280$ bits with the archived tail metric. The evaluator bound,
all earlier-zero exclusions, and its resolution accounting are stated
in Propositions~\ref{prop:v143-first-zero-bound}--\ref{prop:v143-fixed-window-meaning}.
A lower-energy transfer at resolution $10^{-71}$ would need a Rayleigh
gap of about $3.2032\cdot10^{-153}$; that input is absent and would
still require a suitably centered trial. The benchmark's finite
$10^{-71}$ discrepancy is not assigned to the complete ground.
No signed discrepancy, cofinal convergence, G2 or RH claim is made.

Version 1.42 integrates the five NS-17 positioning/goals/bibliography
inserts at their specified anchors, with the continuation claim qualified
and the Zhu bibliography deduplicated. It also proves a unique simple local zero of the complete
$\lambda=4$ ground transform within $1/10$ of the first zeta ordinate
(Proposition~\ref{prop:v142-ground-zero}). The archived complete even
Schur data certify the stronger even-sector separator $10^{-67}$;
a spectral-energy projection estimate transfers endpoint and derivative
signs at $1024/1280$ bits. The full old tail bounds are inherited,
not reassembled. No tiny-discrepancy, zero-ordering or cofinal claim
is made. G2 and RH remain open.

Version 1.41 completes the CCM numerical semantic lock at the existing
$\lambda=3$, $N=120$ compression. Proposition~\ref{prop:v141-ccm-lock}
reproduces eight published discrepancies using interval root gates
replayed at $768$ and $1024$ bits. Lemma~\ref{lem:v141-centering}
identifies the missing centering phase in the prior NS-14 diagnostic;
its sinc-lattice conclusion is withdrawn. The comparison of published
windows is made explicit, and the v1.39 coverage table is linked.
The original later evidence groups remain OPEN. G2 and RH remain open.

Version 1.40 completes NS-5: Proposition~\ref{prop:v140-ground4}
certifies the simple even ground of the complete $\mathsf W_4$,
with ordinary spectral gap greater than $9.7546\cdot10^{-74}$.
The exact frozen trials and archived complete Schur bounds are retained;
a resolvent comparison accounts for the entire change under the shift
$10^{-73}$. All new gates pass at $1024$ and $1280$ bits.
Corollary~\ref{cor:v140-real-zeros} checks the precise distribution
and operator-core hypotheses of Connes--van Suijlekom's Theorem~6.1,
so the complete ground transform has only real zeros. Its zeros are
not identified with Riemann zeros, and no cofinal G2 gap is closed.

Version 1.39 proves the two scoped meta-obstructions in
Theorems~\ref{thm:v139-probe-relaxation} and
\ref{thm:v139-protected-orbit}. The former rules out the strongest
full-distribution mass-cap certificate and the stated variation
refinement. The latter preserves arbitrary finite protected operator
actions, including the actual source residual, while varying spectral
orientation. Its conjugates need not be multipliers. A positive
location-independent counterexample shows why the unrestricted slogan
is false. The evidence report audits the ten historical closures and
does not claim that all are subsumed. No G2 sign gap is closed.

Version 1.38 tests the finite-floor reduction rather than requiring
another positivity certificate. A common shift $8$, head cutoff $256$
and remote cutoff $4096$ certify both complete parity forms at
$\lambda=5,6,8$, with identical moment and verification rules. The
same exact dyadic witnesses pass at $160$ and $256$ bits. Separate
generalized-margin lower bounds are reported; whitening margins and
resolution margins are not substituted for them. The small positive
trial quotients provide upper bounds, not a sign conclusion. A proved
shifted cutoff estimate shows that any bounded shift still has the
old asymptotic scalar-comparison barrier. No G2 gap is closed.

Version 1.37 closes the bounded scalar-primitive shortcut.
Corollary~\ref{cor:v137-scalar-no-go} proves unconditionally that
$a\Delta(\beta_a)\to\infty$ and $\inf\beta_a\to-\infty$ for the actual
arithmetic symbol. The proof first derives a negative cutoff side lobe
under RH with an endpoint-vanishing probe; any bounded cofinal scalar
certificate would itself imply RH by v1.36 and then contradict that
calculation. The conditional linear rate is not asserted unconditionally.
This is not a negative physical Weil test. The complete joint signed
concentration lower floor remains unproved in both parities, and no G2
sign gap is closed. Current delivery is LaTeX only.

Version 1.36 sharpens the remaining sufficient target. By
Proposition~\ref{prop:v136-bounded-floor} and
Corollary~\ref{cor:v136-complement-floor}, one finite ordinary lower bound
on the complete small-residual complement, uniform along a cofinal family,
already suffices for Weil positivity. The error need not be proved to
vanish. The proof amplifies any hypothetical negative compact test by
subtracting ordinary-close radical approximations. No such uniform
arithmetic lower bound is established here. The same dense radical family
proves ordinary strong-resolvent collapse to zero and rules out nonzero
persistent bounded positive comparisons with vanishing error. These are
reductions and obstructions, not a new positivity mechanism or an RH proof.
Both parity sectors and all physical blocks remain required. No G2 sign
gap is closed.

Version 1.35 proves a quantitative growing-block obstruction, rather
than another fixed-window positivity result.  Proposition~\ref{prop:v135-growing-radical}
constructs $2\lfloor a/(2D)\rfloor+1$ independent physical columns with
uniform ordinary Gram lower bound and full operator residual
$Ca e^{-\lambda/100}$, where $a=\log\lambda$ and $D$ is fixed by the
limiting source.  Corollary~\ref{cor:v135-coercivity-obstruction} bounds
any positive complementary gap by that error after fewer columns are
removed, and applies to (G2.6) with arbitrarily many even image columns
through its untouched odd block.  The same number of complete eigenvalues
lies near zero, without a conclusion about their signs.
Proposition~\ref{prop:v135-both-parities} also extends the exact continuum
symbol to both parities and gives their combined gap-free error budget.
The translated radical family is a quantified continuation of a previously
logged fixed-rank obstruction, not a novel positivity mechanism.  The
signed cofinal complement estimate remains unproved. No G2 sign gap is closed.

Version 1.34 proves two further limits on proposed simplifications of
the signed concentration route.  First, disjoint reciprocal-scale
frequency channels have an actual-source-compressed cross norm with
cofinal lower bound $73/(375\pi)>0.0619$; fixed-ratio separation does
not make the Cotlar cross terms small.  More fundamentally, channel
atomization or regrouping leaves $Z_a^*Z_a$ unchanged, so an ordinary
Cotlar row sum is not a new signed mechanism unless its complete
constant meets the original coverage budget.  Second,
Proposition~\ref{prop:v134-primitive} proves an exact scalar
signed-primitive estimate.  Its optimal ordinary error is
$a\Delta_a$, with $\Delta_a$ the maximum negative primitive drawdown.
The scale-invariant Dirichlet model gives the exact obstruction
$a\Delta_a=\pi^3/6$ despite complete positivity and an exact null
source.  Arithmetic diagnostics at $\lambda=3,4,5,6$ likewise increase
from about $2.54$ to $5.51$; their original script and output are not
archived, and they are not proof evidence.  The scalar
candidate is rejected, while the full signed weighted concentration
inequality remains the live conditional even-sector mechanism.  No G2
sign gap is closed.

Version 1.33 derives the exact Toeplitz--Hankel product defect for the
source-compressed concentration hierarchy, including the two rank-one
source terms required by the first-slot-linear convention.  The proved
ordinary source limit then yields a cofinal stress test.  For the nested
bands $[-1/(10a),1/(10a)]\subset[-1/(5a),1/(5a)]$, dilation reduces the
uncompressed operators to two fixed concentration operators on
$(-1,1)$, while source compression vanishes in operator norm on these
compact channels.  Exact Taylor bounds prove that their commutator norm
exceeds $9\cdot10^{-6}$.  Hence generic reciprocal-band approximate
commutation is false for the manuscript's actual source.  This does not
identify the arithmetic superlevel sets with those bands and does not
refute Proposition~\ref{prop:v131-concentration}; it rules out a proposed
automatic simplification.  The exact surviving obligation is the joint
weighted coverage bound in \eqref{eq:v133-joint-channel}.  It is
algebraically the existing signed concentration target, so no renamed
candidate is registered.  No G2 sign gap is closed.

Version 1.32 tests two tempting simplifications of the source-compressed
concentration route.  Proposition~\ref{prop:v132-schatten} proves an
exact Hilbert--Schmidt formula after removal of the actual source and a
strictly sharper negative-level sufficient bound.  The same proposition
gives a scale-invariant nonnegative Dirichlet model with an exact null
source for which that sufficient coefficient stays bounded away from
zero.  Thus Schatten refinement alone cannot supply the cofinal decay.
Proposition~\ref{prop:v132-flat-top} proves that a nontrivial positive
frequency smoothing cannot preserve every form on the physical
difference interval.  A signed flat-top extension remains possible in
principle, but positivity of such an extension is an additional sign
problem, not a consequence of smoothing. Historical numerical tests
at $\lambda=3,4,5$ reported agreement; their original scripts and outputs
are not archived, and those reports are only falsification diagnostics.  The full signed weighted-superlevel
criterion of Proposition~\ref{prop:v131-concentration} remains the live
conditional even-sector mechanism.  Its cofinal arithmetic inequality
and the odd sector remain open.

Version 1.31 develops the exact discrepancy symbol into a
source-compressed time--frequency concentration hierarchy.
Proposition~\ref{prop:v131-concentration} proves both a rigorous
negative-level layer-cake bound and a sharper weighted superlevel
operator criterion. The scalar trace specialization is finite and
correct, but historical numerical diagnostics at $\lambda=3,4,5$
reported that it discards most favorable symbol mass; the original
scripts and outputs are not archived (only the v1.31 report survives). Generic Logvinenko--Sereda
constants likewise do not meet the necessary signed budget. Explicit
even packets centred at the deepest located wells nevertheless retain
positive complete signed energy in that historical report, including
a reported independent physical-space calculation whose script and
output are not archived. Those numerical reports do not close a route; the
weighted concentration operator remains a live conditional mechanism.
Its required cofinal arithmetic concentration inequality and the odd
sector remain open.

Version 1.30 develops the capacity candidate through the exact PNT
continuum kernel. Proposition~\ref{prop:v130-continuum-cancellation}
proves that the even pole cancels the sole positive channel of
$e^{|x-y|/2}$ and identifies the positive remainder as a Green
operator. Proposition~\ref{prop:v130-remainder} isolates the complete
arithmetic discrepancy as a scalar Fourier multiplier and states the
ordinary-error estimate sufficient for the gap-free route. The Green
operator is compact, so this produces no hidden spectral gap. A
stronger pointwise multiplier criterion was reported numerically false
at already certified positive fixed windows; the original v1.30
diagnostic script and output are not archived.
The arithmetic remainder estimate and the odd sector remain open.

Version 1.29 isolates a new candidate mechanism for G2.6 after screening
the project history and primary literature. Proposition~\ref{prop:v129-capacity}
rewrites the inversion-even Weil form by a source-weighted nonlocal
ground-state transform and gives an exact capacity-absorption condition
for asymptotic nonnegativity on the source complement.  The mechanism
uses the simultaneous signs of the physical jump kernel and is not a
new Schur solve, block metric, finite-rank repair or scalar prime-norm
bound. A historical finite $\lambda=3$, $N=64$ diagnostic reports
falsification of the stronger pointwise-supersolution shortcut; its
recomputation script and output are not archived; it neither proves nor disproves the
complete capacity condition.  The odd pole has the opposite sign and
is explicitly outside this candidate's present scope.

Version 1.28 tested the signed block-metric criterion at $\lambda=5$
with the explicitly chosen polynomial cutoff $N=\lambda^2=25$.
The historical computation reports individually positive first four
dyadic blocks and coupling lower bounds contradicting the criterion in
both parities. Its original witnesses, verifiers and outputs are not
archived. Conditional on those reported gates, the obstruction survives
every admissible smaller block metric and positive scalar row reweighting.
It concerns that sufficient comparison, not a negative Weil direction
or other cutoffs and partitions.
Complete positivity at $\lambda=4$, weak G1, all 72 historical
dispositions, the physical endpoint and the sampler's explicit graph
defect are retained. The unchanged $10^{-8}D$ metric at $\lambda=5$
is historically reported to fail; its original evidence is not archived; no complete $\lambda=5$ positivity margin or
successful growing-window tuning rule is asserted. No G2 sign gap
is closed, and RH is not proved.

\paragraph{Reading conventions.}
The Hermitian product is linear in the first variable:
$\ip{x}{cy}=\overline c\,\ip{x}{y}$.
In the Sonine development the multiplicative cutoff is $R>1$.
In the finite Weil development the multiplicative cutoff is $\lambda>1$
and its logarithmic length is $L=2\log\lambda$.
The source graph functional $\eta_{\mathrm B}$, the equal-coefficient
boundary functional $\eta_\partial$, and physical evaluation
$\delta=L^{-1/2}\eta_\partial$ are different objects.
The smoothed Cauchy value is $\mathscr E(F)=\sigma F(\sigma)$.
An off-line root space is used only under a contradiction hypothesis;
no off-line zero is asserted to exist. Independence from a zero list is
a property of a construction, not a substitute for proving its estimates.

\clearpage
\paragraph{Version 1.59 checkpoint (NS-60--66).}
Explicit Euler-grid corrections retain a scoped raw-norm obstruction.
One fixed Mellin integration preserves the RH criterion and gives an
elementary $3\kappa/(8N^2)$ difference-block trace budget. No cofinal positive lower estimate is proved here for the actual
optimized correlations. The canonical sharp M\"obius interpolant still fails along the
full sequence after any number of integrations; selected subsequences
and optimized coefficients are outside that exclusion. Increasing the
smoothing order can give vanishing relative error for one fixed atom
while its absolute error diverges. All original Weil targets, G2 and RH
remain open. A quantitative dual separator excludes the specified positive
coefficient cone. Fixed ordinary Ces\`aro averages also fail for the
canonical rule, but its weaker fixed-order norm-growth exponent is exactly
$\beta_*-1/2$. The required subpolynomial bound remains open; this is a
reformulation, not a new zero-free estimate.

\paragraph{Version 1.60 checkpoint (NS-67--71).}
Complete canonical cell energies and a polynomial cutoff retain the open
subpolynomial-growth criterion. Centering the two complete Abel vectors
shows that cancellation between those whole vectors cannot change the
power exponent; their internal arithmetic remains unestimated.
A positive renewal equation transfers, but does not improve, the forcing
exponent. Its critical weighted total energy is unconditionally infinite.
For optimized residuals, the exact endpoint remainder exceeds the leading
model on two certified blocks. Correcting an elementary Gram cusp restores
the adjacent power, but uniform all-coefficient comparison still fails.
The repaired kernel can instead select a correction direction measured
with the actual complete energy. Its finite gains do not supply the open
cofinal numerator and selected-energy estimates. These are scoped
identities, implications and finite checks, not an RH or G2 proof.
All cofinal arithmetic inputs remain open.

\paragraph{Version 1.61 checkpoint (NS-72--74).}
The repaired model has a complete positive local differential energy and
an explicit discrete-curvature numerator comparison with constants
independent of size. The resulting local rule retains more than $97.63\%$
of the full gain at $N=256$, with complete arithmetic energies.
Three prescribed M\"obius templates, even in their optimized span,
perform less well on all five checked blocks; no asymptotic exclusion
is inferred. An explicit unitary altered family preserves every Gram
entry and admits arbitrarily close finite statistics, yet has a positive
error floor. This controls the interpretation of finite geometric
success, not the original zeta family. The cofinal arithmetic contraction,
uniform Weil floor, G2 and RH remain open.

\paragraph{Version 1.62 checkpoint (NS-75--77).}
Exact divisor feedback cancels the next derivative jumps when old
coefficients are held fixed, but its unit step increases the complete
error on every checked block. At $N=256$ its optimized three-direction
span retains only $56.02\%$ of full gain versus curvature's $97.63\%$.
The continuous relaxation has an exact Hardy inner-factor distance;
its closure is proved strictly larger than the integer closure.
A complete zero-tail bound, using published finite-height information,
gives continuous squared distance below $2.08\cdot10^{-32}$.
This uses external zero verification and is neither an upper bound
for integer error nor a proof that either distance is zero.
The cofinal arithmetic contraction, original Weil floor, G2 and RH
remain open. See Sections~\ref{sec:ns75-divisor-feedback}--\ref{sec:ns77-continuous-tail}.

\paragraph{Version 1.63 checkpoint (NS-78--80).}
The complete squared residual norm is uniformly comparable to a weighted
sum of all reciprocal-integer samples, with exact divisor inversion and
complete physical and sampling tails. Independent cell integration
confirms the finite Gram norms. Tail-balanced divisor feedback still
underperforms curvature at $N=256$; its endpoint-augmented span captures
$71.18\%$ of full gain. The full arithmetic sampling Gram gives a
size-independent preconditioner comparison, yielding at least $21/121$
of available block gain in one ideal step and more than $97.79\%$ in
twenty. No fast inverse is supplied, and this is not a lower bound on
available gain relative to current error. The cofinal signed arithmetic
estimate, original Weil floor, G2 and RH remain open.

\paragraph{Version 1.64 checkpoint (NS-81).}
Finite reciprocal samples bound coefficient size and give ordinary Gram
condition number $O(N^4\log^2(2N))$. A complete discarded-sample bound
makes the arithmetic comparator explicitly finite: cutoff
$M=2^{16}(N+1)^6$ suffices for relative tail below one and projected
condition ratio at most $42$. This cutoff is enormous, and the large
comparator is not assembled. Certified small-prefix and scalar replays
support the constants; no practical solver, cofinal error decay, signed
Weil floor, G2 or RH result follows.

\tableofcontents
\clearpage
\section{Objective and correction}
Fix a prime number $p$ and write
\[
  \lambda_\rho=p^{1/2-\rho}.
\]
The target is a fixed positive Hilbert-space action that retains this multiplier on the arithmetic modes attached to the nontrivial zeros, with positivity or contractivity proved independently of their locations.

A convention correction is essential. Burnol's $Z_w$ is defined by a complex bilinear pairing, whereas the project uses a positive Hermitian product linear in the first variable. If $k_w$ denotes the Hermitian Riesz vector, then
\[
  k_w=Z_{\bar w}.
\]
The project mode is therefore
\[
  v_\rho^R=k_{\bar\rho}^R=Z_\rho^R,
\]
not $Z_{\bar\rho}^R$. With this convention the original compression has the correct critical-line phase. The historical reciprocal-phase rejection resulted from a second, erroneous conjugation.

\section{Burnol framework}
Work in
\[
 K=L^2((0,\infty),dx),
\]
with Hermitian inner product linear in the first slot. Let
\[
 If(x)=x^{-1}f(1/x),\qquad G=I\mathcal F_+I,
\]
where
\[
 (\mathcal F_+f)(x)=2\int_0^\infty \cos(2\pi xt)f(t)\,dt
\]
is the unitary cosine transform, initially on integrable square-integrable
functions and then by its $L^2$ extension. The Mellin conventions are
\[
 M_0f(s)=\int_0^\infty f(x)x^{s-1}\,dx,\qquad
 c_0(s)=\pi^{-s/2}\Gamma(s/2),\qquad M f(s)=c_0(s)M_0f(s).
\]
These formulas are first used where the integrals converge and elsewhere
by the Sonine continuation. In particular $M(E\psi)=c_0\zeta M_0\psi$.
The evaluation and division properties of the completed Sonine space are
the framework inputs from \cite{Burnol2001,Burnol2002Sonine,BurnolCoPoisson2004}.
They do not include the additional quantitative assumptions below. For $R>1$ define
\[
 H_R=\{f:\operatorname{supp}f,\operatorname{supp}Gf\subset(0,R]\},
\]
\[
 V_R=C_c^\infty((1/R,R)),\qquad
 \mathcal D=x\frac{d^2}{dx^2}x,
\]
and
\[
 E(\psi)(x)=\sum_{n\ge1}\psi(nx)-\frac{1}{x}\int_0^\infty\psi(t)\,dt.
\]
The arithmetic co-Poisson space and quotient are
\[
 W_R=\overline{\{E(\mathcal D\phi):\phi\in V_R\}},\qquad
 Q_R=H_R\ominus W_R.
\]
Let $S_R$ and $P_R$ be the orthogonal projections onto $H_R$ and $Q_R$ respectively.

For a nontrivial zero $\rho$, Burnol's construction gives a nonzero vector
\[
 v_\rho^R=Z_\rho^R=k_{\bar\rho}^R\in Q_R.
\]
Put
\[
 v=v_\rho^R,\qquad v_+=v_\rho^{pR},\qquad b=v_+-v.
\]
Because the ambient Sonine spaces are nested and the two vectors represent the same evaluation on $H_R$,
\begin{equation}\label{eq:nested}
 S_Rv_+=v,\qquad b\perp H_R,\qquad
 \norm b^2=\norm{v_+}^2-\norm v^2.
\end{equation}

The unitary dilation is
\[
 (D_pf)(x)=p^{-1/2}f(x/p),
\]
and the completed Mellin transform obeys
\[
 M(D_pf)(s)=p^{s-1/2}M(f)(s).
\]

\subsection*{Quantitative evaluator inputs and compact division}
\begin{assumption}[Evaluator estimates used in the asymptotic sections]\label{ass:evaluators}
For each fixed critical-line parameter $w$, the chosen completed-Mellin
Riesz vectors obey $\|v_w^R\|^2=C_w\log R+O_w(1)$, $C_w>0$.
For fixed distinct critical parameters their cross inner product is
$O_{w,z}(1)$. For a fixed off-line parameter with
$d=|\Re w-1/2|$, $\|v_w^R\|\asymp_w R^d$.
For each fixed jet order $m$, the corresponding ambient jet Gram matrix
$V_{w,R}$ satisfies $\det V_{w,R}\ll_{w,m}R^{2md}$.
Where raw representatives are used below, their local expansion and
Sonine projection error are assumed to have the stated uniform bounds,
including the $O_w(R^{-1})$ projection correction.
\end{assumption}
These quantitative statements were attributed to earlier calculations in
the previous draft without a complete derivation in the supplied files.
They are not consequences of completeness or minimality alone. Sections
using their rates are conditional on this assumption. Exact compression,
covariance, and projection identities do not require it.
Proposition~\ref{prop:v110-normalized-shell} below proves the off-line
\emph{upper} norm bound directly and uses only nested nonzero evaluation
for a lower bound. It does not prove the matching asymptotic, the
critical-line estimates, or the projection-error assertions above.

\begin{lemma}[Compact Mellin division]\label{lem:compact-division}
Let $\theta\in C_c^\infty((1/R,R))$, let $\Theta=M_0\theta$, and let
$w\notin\{0,1\}$. There exists a source $\eta$ in the same compact
support interval with $M_0\eta(s)=\Theta(s)/(s-w)^r$ if and only if
$\Theta$ vanishes to order at least $r$ at $w$. If
$\Theta(0)=\Theta(1)=0$, the divided source has the same two moments.
Moreover $\mathcal D V_R$ is exactly the compact smooth source class with
these two moments zero.
\end{lemma}
\begin{proof}
For one division choose $b$ below the support and set
\[
 \eta(x)=-x^{-w}\int_b^x t^{w-1}\theta(t)\,dt.
\]
The function vanishes below the support. It vanishes above it precisely
when $\Theta(w)=0$, and then is smooth and compactly supported in the
same interval. It solves $-x\eta'-w\eta=\theta$. Integration by parts
gives $(s-w)M_0\eta=\Theta$. Repeat the argument for higher order; the
successive removability conditions are exactly the $r$ vanishing jets.
Necessity follows because a compact source has an entire Mellin transform.
Evaluation at 0 and 1 proves moment preservation when $w\notin\{0,1\}$.
Finally $M_0(\mathcal D\phi)=s(s-1)M_0\phi$. Conversely divide first by
$s$ and then by $s-1$ using the two zero moments and the same compact
construction. The resulting $\phi$ satisfies $\mathcal D\phi=\theta$.
\end{proof}

\begin{lemma}[Finite-jet realization and passage to the closed kernel]\label{lem:jet-density}
Fix $w\notin\{0,1\}$ and a finite order $m$. Admissible compact sources
realize arbitrary jets $(\Psi(w),\ldots,\Psi^{(m-1)}(w))$.
If their associated functionals extend continuously to $W_R$, the
intersection of their kernels in the source core is dense in the
intersection of their kernels in $W_R$.
\end{lemma}
\begin{proof}
For $\psi=\mathcal D\phi$, multiplication by $s(s-1)$ is an invertible
triangular map on the jet at $w$. The distributions
$x^{w-1}(\log x)^j$, $0\le j<m$, are linearly independent on every open
source interval. Hence the jet map from compact smooth $\phi$ is onto:
otherwise a nonzero linear combination of those distributions would
annihilate all test functions. Choose core vectors $u_1,\ldots,u_m$
whose joint jets are the coordinate vectors. If $f_n$ are core
approximants to a vector $f$ with zero jet, replace them by
$f_n-\sum_j\ell_j(f_n)u_j$. Continuity makes the correction tend to zero
and each corrected approximant has zero jet.
\end{proof}


\section{The corrected fixed compression}
Define
\[
 T_{R,p}=P_RD_p\big|_{Q_R}.
\]
It is a contraction independently of RH. The corrected exact boundary identity is
\begin{equation}\label{eq:boundary}
 (T_{R,p}-\lambda_\rho I)v_\rho^R
 =-P_RD_p(v_\rho^{pR}-v_\rho^R).
\end{equation}
Indeed the first-slot Riesz convention gives
\[
 S_RD_pv_\rho^{pR}=\lambda_\rho v_\rho^R,
\]
and \eqref{eq:boundary} follows from \eqref{eq:nested}.

Under Assumption~\ref{ass:evaluators}, the critical-line evaluator norm law is
\[
 \norm{v_\rho^R}^2=C_\rho\log R+O_\rho(1),\qquad C_\rho>0.
\]
Hence \eqref{eq:nested} gives $\norm b=O_{\rho,p}(1)$ and
\begin{equation}\label{eq:criticalrate}
 \frac{\norm{(T_{R,p}-\lambda_\rho I)v_\rho^R}}{\norm{v_\rho^R}}
 =O_{\rho,p}((\log R)^{-1/2}).
\end{equation}
Thus every fixed critical-line mode is asymptotically retained with the correct multiplier. The all-zero version of this estimate remains unproved and is RH-strength.

\section{New source covariance at adjacent scales}
The next lemma is a project deduction from the exact co-Poisson identities. It does not use a zero location.

For a source function define
\[
 (\tau_p\phi)(x)=\phi(x/p).
\]
A direct calculation gives
\begin{equation}\label{eq:covariance}
 \mathcal D\tau_p=\tau_p\mathcal D,
 \qquad
 D_pE(\psi)=p^{-1/2}E(\tau_p\psi).
\end{equation}
Consequently
\begin{equation}\label{eq:DWinclusion}
 D_pW_R\subset W_{pR}.
\end{equation}
Also $V_R\subset V_{pR}$, so $W_R\subset W_{pR}$. Burnol's identity $GE=EI$, together with commutation of $I$ and $\mathcal D$ on the compact source domain, gives
\[
 GW_{pR}=W_{pR}.
\]

\begin{theorem}[Three-piece arithmetic source decomposition]\label{thm:threepiece}
If $R^2>p$, then
\begin{equation}\label{eq:Wdecomp}
 W_{pR}=\overline{W_R+D_pW_R+GD_pW_R}.
\end{equation}
Since $GW_R=W_R$ and $GD_p=D_{1/p}G$, the last summand may equivalently be written $D_{1/p}W_R$.
\end{theorem}

\begin{proof}
The inclusion from right to left follows from \eqref{eq:DWinclusion}, $W_R\subset W_{pR}$, and $G$-invariance of $W_{pR}$.

For the reverse inclusion, the three source intervals
\[
 [1/(pR),R/p],\qquad [1/R,R],\qquad [p/R,pR]
\]
cover $[1/(pR),pR]$ when $R^2>p$. A smooth partition of unity therefore splits every $\phi\in V_{pR}$ as
\[
 \phi=\phi_-+\phi_0+\phi_+,
\]
where $\phi_0\in V_R$, $\phi_+=\tau_p\psi_+$ for some $\psi_+\in V_R$, and $\phi_-=I\tau_p\psi_-$ for some $\psi_-\in V_R$. Applying $E\mathcal D$, using \eqref{eq:covariance} and $GE=EI$, places the three pieces in $W_R$, $D_pW_R$, and $GD_pW_R$ respectively. Taking closures proves \eqref{eq:Wdecomp}.
\end{proof}

\begin{corollary}[Arithmetic projection disappears in two adjoint channels]\label{cor:projectiondrop}
Assume $R^2>p$ and let $y\in Q_{pR}$. Then
\begin{equation}\label{eq:drop}
 P_Ry=S_Ry,\qquad
 P_RD_py=S_RD_py,\qquad
 P_RD_{1/p}y=S_RD_{1/p}y.
\end{equation}
\end{corollary}

\begin{proof}
Since $Q_{pR}\perp W_{pR}$, Theorem~\ref{thm:threepiece} gives
$y\perp W_R$, $y\perp D_pW_R$, and $y\perp D_{1/p}W_R$. These imply respectively that $S_Ry$, $S_RD_{1/p}y$, and $S_RD_py$ are orthogonal to $W_R$. Each therefore lies in $Q_R$, yielding \eqref{eq:drop}.
\end{proof}

\section{A quotient-native prime operator pencil}
Define two contractions from the old quotient into the enlarged quotient,
\[
 J_{R,p}=P_{pR}\big|_{Q_R}:Q_R\to Q_{pR},
\]
\[
 B_{R,p}=P_{pR}D_{1/p}\big|_{Q_R}:Q_R\to Q_{pR}.
\]
Their adjoints are
\[
 J_{R,p}^*=P_R\big|_{Q_{pR}},\qquad
 B_{R,p}^*=P_RD_p\big|_{Q_{pR}}.
\]
For the zero vector $v_+=v_\rho^{pR}$, nested evaluation and the corrected boundary calculation give the exact identities
\begin{equation}\label{eq:pencildual}
 J_{R,p}^*v_+=v,
 \qquad
 B_{R,p}^*v_+=\lambda_\rho v.
\end{equation}
Equivalently,
\begin{equation}\label{eq:pencil}
 (B_{R,p}^*-\lambda_\rho J_{R,p}^*)v_\rho^{pR}=0.
\end{equation}
Thus the arithmetic prime multiplier already appears as an exact generalized eigenvalue of a fixed, zero-independent operator pencil. The unresolved issue is converting this generalized eigenrelation into a single positive fixed-space action without inserting the zero location by hand.

\section{Exact defect kernel and a sufficient fixed-space theorem}
Define on $Q_{pR}$
\begin{equation}\label{eq:Delta}
 \Delta_{R,p}=J_{R,p}J_{R,p}^*-B_{R,p}B_{R,p}^*.
\end{equation}
This is self-adjoint. From \eqref{eq:pencildual}, for every nontrivial zero,
\begin{equation}\label{eq:Deltadiag}
 \ip{\Delta_{R,p}v_\rho^{pR}}{v_\rho^{pR}}
 =(1-|\lambda_\rho|^2)\norm{v_\rho^R}^2.
\end{equation}
More generally, for zeros $\rho,\sigma$,
\begin{equation}\label{eq:Pick}
 \ip{\Delta_{R,p}v_\sigma^{pR}}{v_\rho^{pR}}
 =\bigl(1-\lambda_\sigma\overline{\lambda_\rho}\bigr)
   \ip{v_\sigma^R}{v_\rho^R}.
\end{equation}
This is a Pick-type difference-of-Gram kernel on the zero-evaluation system.

\begin{proposition}[Finite-cutoff span positivity would imply RH]\label{prop:positivity}
If, for one fixed $p>1$ and admissible cutoffs, $\Delta_{R,p}$ were nonnegative on the span of the zero-evaluation vectors $v_\rho^{pR}$, then every nontrivial zero would satisfy $|\lambda_\rho|\le1$. Functional-equation symmetry would then imply RH.
\end{proposition}

\begin{proof}
Apply nonnegativity to a single vector in \eqref{eq:Deltadiag}. Since $v_\rho^R\ne0$,
$1-|\lambda_\rho|^2\ge0$. Hence $p^{1/2-\Re\rho}\le1$, or $\Re\rho\ge1/2$. The functional equation supplies the reflected zero $1-\rho$, forcing equality.
\end{proof}

The span-positivity hypothesis is substantially stronger than RH and should not be treated as the main target. Indeed, assume RH for the moment. Then every diagonal quantity in \eqref{eq:Deltadiag} is zero. If $\Delta_{R,p}$ were positive semidefinite on the whole zero span, positivity would force each $v_\rho^{pR}$ to lie in the nullspace of the restricted form. Hence for all critical zeros $\rho,\sigma$,
\begin{equation}\label{eq:crossvanish}
 \bigl(1-\lambda_\sigma\overline{\lambda_\rho}\bigr)
 \ip{v_\sigma^R}{v_\rho^R}=0.
\end{equation}
Thus distinct prime phases would require the corresponding evaluation vectors to be orthogonal. Burnol's zero systems are studied as complete/minimal systems, not supplied here as an orthogonal basis. Therefore RH by itself does not furnish the extra cross-term vanishing in \eqref{eq:crossvanish}. The safe target is the \emph{diagonal} sign on each zero evaluator, not positivity of the full Pick matrix.

For the same reason, the stronger global operator inequality
\begin{equation}\label{eq:globalpositive}
 B_{R,p}B_{R,p}^*\le J_{R,p}J_{R,p}^*
\end{equation}
should be regarded only as an over-strong sufficient condition. Were it true, the Douglas factorization lemma would produce a contraction $C_{R,p}:Q_R\to Q_R$ with $B_{R,p}=J_{R,p}C_{R,p}$ and hence $C_{R,p}^*v_\rho^R=\lambda_\rho v_\rho^R$. Nothing in the present work proves the cross-term conditions required for such a global factorization, even under RH.

\section{The defect becomes a pure Sonine comparison}
Theorem~\ref{thm:threepiece} removes the arithmetic projection from the defect form. For $y\in Q_{pR}$ and $R^2>p$, Corollary~\ref{cor:projectiondrop} gives
\[
 J_{R,p}^*y=S_Ry,
 \qquad
 B_{R,p}^*y=S_RD_py.
\]
Therefore
\begin{equation}\label{eq:puresonine}
 \ip{\Delta_{R,p}y}{y}
 =\norm{S_Ry}^2-\norm{S_RD_py}^2.
\end{equation}
The diagonal sign problem has thus been reduced to the following concrete question:
\begin{quote}
For each zero evaluator $y=v_\rho^{pR}$, can an independent arithmetic estimate prove $\|S_RD_py\|\le\|S_Ry\|$?
\end{quote}
A proof of
\[
 \norm{S_RD_py}\le\norm{S_Ry}
\]
on each zero evaluator would prove RH by \eqref{eq:Deltadiag}. Positivity on the whole span is a stronger sufficient condition, as explained in Proposition~\ref{prop:positivity}. A proof on all of $Q_{pR}$ would give the stronger Douglas factorization above.

This is the current fixed-space compatibility target. It is more explicit than the earlier boundary-distance statement because the co-Poisson projection has disappeared from the quadratic form exactly. No sign for \eqref{eq:puresonine} has yet been proved.


\section{Exact Fredholm reduction to a boundary-shell imbalance}
We now make \eqref{eq:puresonine} explicit using Burnol's finite-interval Fredholm projection. Put
\[
 a=R^{-1},\qquad q=Iy\in K_{a/p},
\]
where $I f(x)=x^{-1}f(1/x)$ and $K_{a/p}$ is the cosine-Sonine space with gap $a/p$. Let $\mathcal F_+$ be the unitary cosine transform, let $P_a$ denote restriction/projection to $(0,a)$, and set
\[
 F_a=P_a\mathcal F_+P_a,\qquad D_a=F_a^2.
\]
Under inversion, $S_R$ becomes the orthogonal projection $\pi_a$ onto $K_a$, while $ID_p=D_{1/p}I$. Hence
\begin{equation}\label{eq:fredholm-start}
 \norm{S_Ry}^2-\norm{S_RD_py}^2
 =\norm{\pi_a q}^2-\norm{\pi_aD_{1/p}q}^2.
\end{equation}

For a general $f\in L^2(0,\infty)$ define
\[
 u=P_af,\qquad v=P_a\mathcal F_+f.
\]
Burnol's projection theorem (equivalently the $2\times2$ Fredholm inverse in Lemma 2 of \cite{Burnol2002Sonine}) gives the exact lost energy
\begin{align}
 \mathcal E_a(f)&:=\norm{f-\pi_af}^2\\
 &=\frac12\ip{u+v}{(I+F_a)^{-1}(u+v)}
 +\frac12\ip{u-v}{(I-F_a)^{-1}(u-v)}.\label{eq:FredholmEnergy}
\end{align}
In particular $\mathcal E_a(f)\ge0$ and
\[
 \norm{\pi_af}^2=\norm f^2-\mathcal E_a(f).
\]
Since dilation is unitary, \eqref{eq:fredholm-start} becomes
\begin{equation}\label{eq:fredholm-difference}
 \norm{S_Ry}^2-\norm{S_RD_py}^2
 =\mathcal E_a(D_{1/p}q)-\mathcal E_a(q).
\end{equation}

The Sonine condition $q\in K_{a/p}$ now produces a one-channel collapse after dilation. Indeed $\mathcal F_+q$ vanishes on $(0,a/p)$ and
\[
 \mathcal F_+D_{1/p}=D_p\mathcal F_+,
\]
so
\[
 P_a\mathcal F_+D_{1/p}q=0.
\]
Writing
\[
 u=P_aq,\qquad v=P_a\mathcal F_+q,
 \qquad u'=P_aD_{1/p}q,
\]
we obtain exactly
\begin{equation}\label{eq:dilated-energy}
 \mathcal E_a(D_{1/p}q)=\ip{u'}{(I-D_a)^{-1}u'}.
\end{equation}
Equations \eqref{eq:FredholmEnergy}--\eqref{eq:dilated-energy} are the desired finite-$R$ Fredholm formula for the sign problem.

There is also a transparent shell approximation. The Hilbert--Schmidt estimate gives $\norm{F_a}\le2a$. Consequently, for $a<1/2$,
\begin{align}
 \left|\mathcal E_a(q)-(\norm u^2+\norm v^2)\right|
 &\le \frac{2a}{1-2a}(\norm u^2+\norm v^2),\label{eq:Eapprox1}\\
 \left|\mathcal E_a(D_{1/p}q)-\norm{u'}^2\right|
 &\le \frac{4a^2}{1-4a^2}\norm{u'}^2.\label{eq:Eapprox2}
\end{align}
Because $q$ and $\mathcal F_+q$ both vanish below $a/p$,
\[
 \norm{u'}^2-\norm u^2
 =\int_a^{pa}|q(t)|^2\,dt,
 \qquad
 \norm v^2=\int_{a/p}^{a}|\mathcal F_+q(t)|^2\,dt.
\]
Thus
\begin{equation}\label{eq:shell-balance}
 \boxed{
 \norm{S_Ry}^2-\norm{S_RD_py}^2
 =\int_a^{pa}|q(t)|^2dt
 -\int_{a/p}^{a}|\mathcal F_+q(t)|^2dt
 +R_{a,p}(q),}
\end{equation}
where, explicitly,
\begin{align*}
 |R_{a,p}(q)|
 &\le \frac{4a^2}{1-4a^2}\int_{a/p}^{pa}|q(t)|^2dt\\
 &\quad+\frac{2a}{1-2a}
 \left(\int_{a/p}^{a}|q(t)|^2dt+
 \int_{a/p}^{a}|\mathcal F_+q(t)|^2dt\right).
\end{align*}
The sign problem has therefore been reduced, up to a quantitatively small Fredholm correction, to a signed imbalance between two adjacent multiplicative shells: a physical outer shell and a Fourier inner shell.

The arithmetic decomposition from Theorem~\ref{thm:threepiece} implies more than mere Sonine membership. For $y\in Q_{pR}$,
\[
 S_Ry,\quad S_RD_py,\quad S_RD_{1/p}y\in Q_R.
\]
Equivalently, the shell pair in \eqref{eq:shell-balance} comes from a vector orthogonal to the old, prime-expanded, and inverse-prime arithmetic source directions simultaneously. This is the precise additional arithmetic structure available for a future sign theorem.

\paragraph{What this does not prove.}
There is no generic Sonine inequality forcing the first shell integral in \eqref{eq:shell-balance} to dominate the second. On a zero evaluator $y=v_\rho^{pR}$ the exact dual identity already gives
\[
 \norm{S_Ry}^2-\norm{S_RD_py}^2
 =(1-|\lambda_\rho|^2)\norm{v_\rho^R}^2.
\]
Hence positivity on all zero evaluators is itself equivalent to the required one-sided zero exclusion. The Fredholm reduction is useful because it identifies the missing mechanism as a sign-sensitive arithmetic shell theorem; it does not manufacture that theorem from generic uncertainty geometry.


\subsection{Signed moments and an unconditional relative error budget}
The $O(a)$ correction in \eqref{eq:shell-balance} contains a signed
cross term. It can be retained explicitly, and even resummed with all
constant-kernel contributions, before estimating the error.

\begin{lemma}[Signed Fredholm formula and rank-one resummation]\label{lem:v110-signed}
Let $0<a<1/2$ and use $u,v,u'$ as above. Put
$A_a=(I-F_a^2)^{-1}$ and
\[
 H_a=\norm{u'}^2+\norm u^2+\norm v^2,\qquad
 \mathcal B_{a,p}(q)=\int_a^{pa}|q|^2-\int_{a/p}^a|\mathcal F_+q|^2.
\]
For $\mathfrak d_a(q)=\mathcal E_a(D_{1/p}q)-\mathcal E_a(q)$,
\begin{align}
 \mathfrak d_a(q)
 &=\ip{u'}{A_au'}-\ip{u}{A_au}-\ip{v}{A_av}
       +2\Re\ip{u}{F_aA_av}.\label{eq:v110-exact-signed}
\end{align}
Define $m=\int_0^a u$, $n=\int_0^a v$, and $m'=\int_0^a u'$.
Then
\begin{align}
 \mathfrak d_a(q)&=\mathcal B_{a,p}(q)+\mathcal C_a(q)
                      +\mathcal R_a(q),\label{eq:v110-rankone}\\
 \mathcal C_a(q)&=
 \frac{4\Re(m\overline n)+4a(|m'|^2-|m|^2-|n|^2)}{1-4a^2},
 \qquad
 |\mathcal R_a(q)|\le
 \frac{4\pi^2a^5}{5(1-2a)^2}H_a.\label{eq:v110-rankone-error}
\end{align}
In particular, $m'=p^{-1/2}\int_{a/p}^{pa}q(t)\,dt$.
These formulas require Sonine membership, not arithmetic orthogonality.
\end{lemma}
\begin{proof}
The identities $(I\pm F_a)^{-1}=A_a(I\mp F_a)$ in
\eqref{eq:FredholmEnergy} give \eqref{eq:v110-exact-signed}.
Self-adjointness and commutation of $F_a,A_a$ make the two mixed
terms conjugate; they enter the \emph{difference} with a plus sign.

Let $T_af=2\int_0^a f(t)\,dt$, the constant-kernel operator on $(0,a)$.
The actual cosine kernel gives
\[
 |2\cos(2\pi xt)-2|\le4\pi^2x^2t^2,\qquad
 \norm{F_a-T_a}\le\norm{F_a-T_a}_{\mathrm{HS}}
 \le\frac{4\pi^2}{5}a^5.
\]
Both $\norm{F_a}$ and $\norm{T_a}$ are at most $2a$.
On $L^2(0,a)\oplus L^2(0,a)$ set
\[
 K_F=\begin{pmatrix}I&F_a\\F_a&I\end{pmatrix},\qquad
 K_T=\begin{pmatrix}I&T_a\\T_a&I\end{pmatrix}.
\]
The resolvent identity yields
\[
 \norm{K_F^{-1}-K_T^{-1}}
 \le\frac{4\pi^2a^5}{5(1-2a)^2}.
\]
Apply this bound separately to $(u,v)$ and $(u',0)$ in the lost-energy
quadratic forms. With $e=1_{(0,a)}$ and
$e\otimes e:f\mapsto\ip{f}{e}e$, one has
\[
 (I-T_a^2)^{-1}=I+\frac{4a}{1-4a^2}e\otimes e,
 \qquad T_a(I-T_a^2)^{-1}=\frac{2}{1-4a^2}e\otimes e.
\]
Substitution proves the stated moment formula and error. The final
identity for $m'$ follows from $D_{1/p}q(x)=\sqrt p\,q(px)$ and
the gap below $a/p$.
\end{proof}

\begin{proposition}[Normalized shell reduction without evaluator asymptotics]\label{prop:v110-normalized-shell}
Fix a nontrivial off-line zero $\rho$, put
$\beta=\Re\rho$, $d=|\beta-1/2|\in(0,1/2)$, and fix a prime $p$.
For $y_R=v_\rho^{pR}$, $q_R=Iy_R$, $a=R^{-1}$ and
$N_R=\norm{v_\rho^R}^2$, the following estimates do not use
Assumption~\ref{ass:evaluators}:
\begin{align}
 \frac{|\mathcal R_a(q_R)|}{N_R}&=O_{\rho,p}(R^{2d-5}),
 &\frac{|\mathcal C_a(q_R)|}{N_R}&=O_{\rho,p}(R^{2d-1}),
 \label{eq:v110-normalized-errors}\\
 \frac{\mathcal B_{1/R,p}(q_R)}{N_R}
 &=1-p^{1-2\beta}+O_{\rho,p}(R^{2d-1}).
 \label{eq:v110-shell-limit}
\end{align}
Constants may also depend on a fixed initial cutoff with nonzero
evaluation. In particular all displayed error terms tend to zero.
\end{proposition}
\begin{proof}
For $\Re s=b>1/2$ and $f\in H_R$, support in $(0,R]$ and
Cauchy--Schwarz give
\[
 |M f(s)|\le |c_0(s)|\left(\frac{R^{2b-1}}{2b-1}\right)^{1/2}\norm f.
\]
The Mellin cosine multiplier, with the left-Mellin convention of
Section 2, gives
\[
 M_0(Gf)(s)=\frac{c_0(1-s)}{c_0(s)}M_0f(1-s),
 \qquad M(Gf)(s)=M f(1-s).
\]
These identities hold first on the Mellin line and extend by the
Sonine continuation; see \cite[\S1, (7)]{Burnol2004}, converting its
right-Mellin convention to the present one. Since $G$ is a unitary
involution preserving $H_R$, the same bound on the reflected parameter
applies when $b<1/2$. The norm of the evaluator at $\bar\rho$ is thus
\[
 \norm{v_\rho^R}^2\le C_\rho R^{2d},\qquad
 C_\rho=
 \begin{cases}
 |c_0(\rho)|^2/(2\beta-1),&\beta>1/2,\\
 |c_0(1-\rho)|^2/(1-2\beta),&\beta<1/2.
 \end{cases}
\]
Nested evaluation gives $N_R\ge N_{R_0}>0$ for $R\ge R_0$;
this uses only the nonzero evaluator and orthogonal projection, not
a lower growth asymptotic. Also
$H_a\le3\norm{q_R}^2=3\norm{v_\rho^{pR}}^2\ll_{\rho,p}R^{2d}$.
The remainder estimate follows. By Cauchy--Schwarz,
$|m|^2\le a\norm u^2$, and likewise for $n,m'$, whence
\[
 |\mathcal C_a(q_R)|\le
 \frac{2a(\norm u^2+\norm v^2)+4a^2H_a}{1-4a^2}.
\]
This proves the second error estimate. Finally use the exact arithmetic
diagonal identity \eqref{eq:Deltadiag} in \eqref{eq:v110-rankone}.
\end{proof}

\paragraph{The remaining arithmetic obligation, now with no Fredholm ambiguity.}
It suffices to fix $p=2$ and prove independently, for each hypothetical
zero with $\Re\rho<1/2$, that
\begin{equation}\label{eq:v110-sign-target}
 \liminf_{R\to\infty}
 \frac{\displaystyle\int_{1/R}^{2/R}|Iv_\rho^{2R}|^2
       -\displaystyle\int_{1/(2R)}^{1/R}|\mathcal F_+Iv_\rho^{2R}|^2}
      {\norm{v_\rho^R}^2}\ge0.
\end{equation}
Indeed \eqref{eq:v110-shell-limit} would instead give a strictly negative
limit, and reflection then excludes right-hand zeros too. Multiplicity
does not alter this criterion, since the value evaluator already exists
at every zero. No uniformity in $\rho$, no full-span positivity, and no
matching evaluator lower asymptotic is required. The rates are not
uniform as the parameter approaches the critical line or strip edges.
The three arithmetic source orthogonalities remain available, but neither
they nor the signed moment formula prove \eqref{eq:v110-sign-target}.
In particular the explicit cross term is asymptotically negligible on
each fixed off-line evaluator; controlling it alone cannot settle the
leading shell sign. This closes a normalization prerequisite of this
route, not the signed G2 estimate or RH.

\section{Projection-angle form of the same obstruction}
There is a second exact reformulation that removes the Fredholm coordinates altogether. For any closed subspace $M\subset K$ and unitary $U$, the projection onto $U^{-1}M$ is $U^{-1}P_MU$. Hence
\[
 D_{1/p}S_RD_p
\]
is the orthogonal projection onto $D_{1/p}H_R$. If $y\in Q_{pR}$, Theorem~\ref{thm:threepiece} gives $y\perp W_R$ and $y\perp D_{1/p}W_R$. Therefore
\[
 S_Ry\in Q_R,
 \qquad
 D_{1/p}S_RD_py\in D_{1/p}Q_R.
\]
The two vectors are in fact the orthogonal projections of $y$ onto the indicated arithmetic quotient spaces. Thus
\begin{equation}\label{eq:angleform}
 \boxed{
 \norm{S_Ry}^2-\norm{S_RD_py}^2
 =\norm{P_{Q_R}y}^2-\norm{P_{D_{1/p}Q_R}y}^2.}
\end{equation}
Here the second $P$ denotes the ordinary Hilbert projection onto the closed subspace $D_{1/p}Q_R$.

Equation~\eqref{eq:angleform} is useful conceptually. The two comparison spaces $Q_R$ and $D_{1/p}Q_R$ are unitarily congruent; neither is generically contained in the other. The missing sign is therefore not a consequence of dimension, positivity, or ordinary projection monotonicity. It asks for a genuinely prime-oriented arithmetic preference inside $Q_{pR}$.

For a zero evaluator $y=v_\rho^{pR}$ the two projections are explicit:
\[
 P_{Q_R}y=v_\rho^R,
 \qquad
 P_{D_{1/p}Q_R}y
 =\lambda_\rho D_{1/p}v_\rho^R,
\]
so \eqref{eq:angleform} reduces again to
\[
 (1-|\lambda_\rho|^2)\norm{v_\rho^R}^2.
\]
Thus any proof based only on the abstract fact that the two spaces are unitary images of one another cannot decide the sign. A successful argument must exploit how the co-Poisson arithmetic range sits relative to the dilation direction, not merely the ambient Sonine geometry.

\subsection*{The next concrete arithmetic object}
Theorem~\ref{thm:threepiece} identifies the enlarged arithmetic space with the closure of the range of the synthesis map
\[
 \mathcal A_{R,p}:W_R\oplus W_R\oplus W_R\longrightarrow H_{pR},
 \qquad
 \mathcal A_{R,p}(w_0,w_+,w_-)
 =w_0+D_pw_++D_{1/p}w_-.
\]
Thus $W_{pR}=\overline{\operatorname{Ran}\mathcal A_{R,p}}$; closed range is not asserted. Its Gram operator is the $3\times3$ Toeplitz-type block matrix
\[
 \mathcal A_{R,p}^*\mathcal A_{R,p}
 =\begin{pmatrix}
 I & C_p & C_{1/p}\\
 C_p^* & I & C_{1/p^2}\\
 C_{1/p}^* & C_{1/p^2}^* & I
 \end{pmatrix},
\]
where, for example,
\[
 C_p=P_{W_R}D_p\big|_{W_R}.
\]
The exact placement of adjoints depends only on the chosen first-slot convention. Positivity of this Gram matrix is automatic because it is a Gram operator; what is not automatic is the Schur-complement inequality needed to orient \eqref{eq:angleform}. This identifies a new, genuinely arithmetic target: obtain sign-sensitive control of the cross-scale correlation operators $C_p$ and $C_{p^2}$ from the co-Poisson source formula. Generic Hilbert-space positivity alone cannot provide it.


\subsection{A convergent arithmetic adjoint for the shell constraints}
To use the source ranges in a signed estimate, their infinite sums must
first be paired legitimately with arbitrary $L^2$ vectors. The formal
sum $\sum_{n\ge1}y(t/n)/n$ need not converge. The following cutoff
supplies an explicit renormalization in the conventions of Section 2.
The integer $M$ here is a co-Poisson summation cutoff, not a Fourier cut
or the Hardy smoothing order used in later sections.

\begin{lemma}[Integral-corrected source cutoff]\label{lem:v111-cutoff}
Let $h\in C_c^\infty((\alpha,b))$, $0<\alpha<b$, put
$V(h)=\int|h'|$, and define
\[
 E_Mh(x)=\sum_{n=1}^M h(nx)-\frac1x\int_0^{Mx}h(t)\,dt.
\]
Then $E_Mh=Eh$ for $x\ge b/M$, and
\begin{equation}\label{eq:v111-source-error}
 \norm{E_Mh-Eh}_2\le V(h)\sqrt{b/M}.
\end{equation}
If $h\ne0$ and $\int h=0$, put $H(t)=\int_0^t h$.
The uncorrected sum $U_Mh=\sum_{n\le M}h(nx)$ instead satisfies
\begin{equation}\label{eq:v111-raw-divergence}
 \norm{U_Mh-Eh}_2
 =\sqrt M\left(\int_\alpha^b\frac{|H(t)|^2}{t^2}\,dt\right)^{1/2}
       +O_h(M^{-1/2}).
\end{equation}
The leading constant is strictly positive.
\end{lemma}
\begin{proof}
Subtract the defining co-Poisson counterterm to obtain
\[
 E_Mh-Eh=-\sum_{n>M}h(nx)+\frac1x\int_{Mx}^\infty h(t)\,dt.
\]
On each interval $((n-1)x,nx)$, the error between the right endpoint
value and the integral average is at most the variation of $h$ on that
interval. Thus the displayed difference is bounded by
$\int_{Mx}^\infty|h'|\le V(h)$ and vanishes above $b/M$.
This proves \eqref{eq:v111-source-error}. The same quadrature bound
shows $Eh$ bounded near zero; above $b$ it is $-x^{-1}\int h$, so
$Eh\in L^2$. If $\int h=0$, then $H$ is supported in $[\alpha,b]$ and
$U_Mh-Eh=(E_Mh-Eh)+H(Mx)/x$.
Changing variables gives
$\norm{H(Mx)/x}_2=\sqrt M\norm{H(t)/t}_2$.
The triangle inequality proves \eqref{eq:v111-raw-divergence}; its
constant cannot vanish unless $h=H'=0$.
\end{proof}

\begin{proposition}[Regularized adjoint and arithmetic source equation]\label{prop:v111-adjoint}
For $y\in L^2(0,\infty)$ set
\begin{equation}\label{eq:v111-adjoint-finite}
 (\mathscr A_My)(t)=\sum_{n=1}^M\frac{y(t/n)}n
                       -\int_{t/M}^\infty\frac{y(x)}x\,dx.
\end{equation}
There is a distribution $\mathscr Ay\in H^{-1}_{\rm loc}(0,\infty)$
such that $\mathscr A_My\to\mathscr Ay$ locally in $H^{-1}$, and
\begin{align}
 \ip{E_Mh}{y}&=\int h(t)\overline{\mathscr A_My(t)}\,dt,
 &\ip{Eh}{y}&=\ip{h}{\mathscr Ay},\label{eq:v111-duality}\\
 |\ip{h}{(\mathscr A_M-\mathscr A)y}|
 &\le V(h)\sqrt{b/M}\norm y_2.\label{eq:v111-adjoint-error}
\end{align}
The distribution pairing in the second slot extends the first-slot-linear
Hermitian convention. For $y\in H_R$, one has the exact characterization
\begin{equation}\label{eq:v111-arithmetic-equation}
 y\in Q_R\quad\Longleftrightarrow\quad
 \mathscr Ay(t)=A_y+B_y/t\quad\hbox{on }(1/R,R)
\end{equation}
as distributions, for some constants $A_y,B_y$.
For every dilation factor $c>0$, $\mathscr A D_c=D_c\mathscr A$.
\end{proposition}
\begin{proof}
For finite $M$ all sums may be integrated term by term. In the integral
term, Fubini is justified by
\[
 \int_{\alpha/M}^\infty\frac{|y(x)|}{x}\,dx
 \le\sqrt{M/\alpha}\norm y_2.
\]
Changing $t=nx$ in each summand and reversing the counterterm integral
gives \eqref{eq:v111-duality}. Lemma~\ref{lem:v111-cutoff} gives
\eqref{eq:v111-adjoint-error}. Since
$V(h)\le\sqrt{b-\alpha}\norm{h'}_2$, the differences are Cauchy in
the dual of $H_0^1(\alpha,b)$, using its derivative norm. These local
limits agree on overlaps and define $\mathscr Ay$.

By the definition of $W_R$, arithmetic orthogonality is exactly
$\ip{\mathcal D\phi}{\mathscr Ay}=0$ for all
$\phi\in C_c^\infty((1/R,R))$. The operator
$\mathcal D\phi=(t^2\phi')'$ is formally self-adjoint. Hence this is
$(t^2(\mathscr Ay)')'=0$ in distributions. Integrating twice on the
connected interval gives precisely $A_y+B_y/t$. The converse follows
by integration by parts and passage to the closure defining $W_R$.
Finally a change of variables in \eqref{eq:v111-adjoint-finite} gives
$\mathscr A_MD_c=D_c\mathscr A_M$ exactly. Distributional passage
to the limit proves the covariance assertion.
\end{proof}

For the three-scale problem this becomes
$\mathscr Ay=A_y+B_y/t$ on $(1/(pR),pR)$, in addition to
$y\in H_{pR}$. This makes all three source constraints explicit;
it does not add three independent scalar conditions to that one local
equation. In particular the two coefficients need not vanish.
No pointwise or local $L^2$ convergence of $\mathscr A_My$ is claimed
for arbitrary $y$.
There is also an explicit growing-interval budget. With
$I_R=(1/(pR),pR)$ and the derivative norm on $H_0^1(I_R)$,
\eqref{eq:v111-adjoint-error} gives
\[
 \norm{(\mathscr A_M-\mathscr A)y}_{(H_0^1(I_R))'}
 \le\sqrt{pR(pR-1/(pR))/M}\norm y_2.
\]
For the fixed off-line evaluator $y=v_\rho^{pR}$,
Proposition~\ref{prop:v110-normalized-shell} therefore bounds this error
divided by $\norm{v_\rho^R}$ by $O_{\rho,p}(R^{1+d}/\sqrt M)$.
Choosing $M=\lceil R^4\rceil$ makes it tend to zero. This controls
source testing in the stated dual norm, not shell energies or a source
Gram inverse; it does not change any later Fourier cutoff.

\begin{corollary}[Controlled three-channel source Gram limits]\label{cor:v111-gram}
Fix finitely many smooth sources $h_i$ supported in $(\alpha,b)$, and
write $w_i=Eh_i$, $w_i^{(M)}=E_Mh_i$ and
$\epsilon_i=V(h_i)\sqrt{b/M}$. For
$c,d\in\{1,p,1/p\}$,
\begin{align*}
 &\left|\ip{D_cw_i^{(M)}}{D_dw_j^{(M)}}-
                  \ip{D_cw_i}{D_dw_j}\right|\\
 &\hspace{15mm}\le
 \epsilon_i\norm{w_j}+\epsilon_j\norm{w_i}+\epsilon_i\epsilon_j.
\end{align*}
Thus every fixed finite source-Gram matrix has a convergent arithmetic
quadrature representation. No convergence of the associated inverse or
uniform lower singular-value bound is implied.
\end{corollary}
\begin{proof}
Expand the difference into its two single-error terms and its double-error
term, use unitarity of both dilations and apply Cauchy--Schwarz.
\end{proof}

\paragraph{A check against a circular sign argument.}
For $1/2<\Re w<1$ and $y(x)=x^{w-1}1_{(0,B)}(x)$, which is in $L^2$,
the finite formula on $0<t<B$ gives
\[
 \mathscr A_My(t)=t^{w-1}
 \left(\sum_{n\le M}n^{-w}-\frac{M^{1-w}}{1-w}\right)
                  +\frac{B^{w-1}}{1-w}.
\]
The elementary Euler-summation continuation of $\zeta$ therefore gives
\begin{equation}\label{eq:v111-power-response}
 \mathscr Ay(t)=\zeta(w)t^{w-1}+\frac{B^{w-1}}{1-w}.
\end{equation}
Indeed, comparison of $n^{-w}$ with its integral over $(n,n+1)$
gives a locally uniform tail $O_w(M^{-\Re w})$ on $\Re w>0$,
$w\ne1$; the limiting analytic function agrees with $\zeta$ on
$\Re w>1$ and hence throughout that domain.
This is consistent with the Mellin co-Poisson identities in
\cite{Burnol2004,BurnolCoPoisson2004}. At a hypothetical zero the
nonconstant power disappears, but this only re-encodes the zero condition.
The raw power is not asserted to satisfy the second Sonine support
condition. Thus neither this scalar calculation nor positivity of the
source Gram matrix proves the signed shell target
\eqref{eq:v110-sign-target}. The new adjoint supplies a justified
arithmetic constraint for a future quadratic estimate, not that estimate.

\section{A local-prime sieve identity inside the three-scale range}
The three-piece decomposition contains more arithmetic information than support covariance alone. Let $h=\mathcal D\phi$ with $\phi\in V_R$. Then $\int h=0$, so the co-Poisson counterterm vanishes and
\[
 E(h)(x)=\sum_{n\ge1}h(nx).
\]
Since
\[
 p^{-1/2}D_{1/p}f(x)=f(px),
\]
we obtain the exact identity
\begin{equation}\label{eq:psieve}
 \boxed{
 \left(I-p^{-1/2}D_{1/p}\right)E(h)(x)
 =\sum_{\substack{n\ge1\\ p\nmid n}}h(nx).}
\end{equation}
Thus the difference between the old arithmetic direction and its inverse-prime dilation is precisely the co-Poisson sum with the $p$-divisible lattice points removed.

On Mellin transforms, \eqref{eq:psieve} is the familiar local Euler-factor cancellation
\begin{equation}\label{eq:eulerremove}
 \widehat{E(h)}(s)\longmapsto
 (1-p^{-s})\widehat{E(h)}(s)
 =\zeta(s)(1-p^{-s})\widehat h(s),
\end{equation}
within the strip where the co-Poisson/Mellin identity is interpreted. The factor $\zeta(s)(1-p^{-s})$ is the zeta Euler product with the local factor at $p$ removed.

Applying the functional-equation involution $G$ gives the companion direction
\[
 G\left(I-p^{-1/2}D_{1/p}\right)G
 =I-p^{-1/2}D_p,
\]
whose Mellin factor is $1-p^{s-1}$. On the critical line these two local factors are complex conjugates and have equal modulus; off the line their radial behavior differs. This is a genuine prime-local structure absent from a generic pair of Sonine subspaces.

The identity does not by itself prove the shell sign: $y\in Q_{pR}$ is already orthogonal to both $W_R$ and $D_{1/p}W_R$, hence also to their $p$-sieved difference. Its value is that it identifies a concrete arithmetic correlation operator to study. If
\[
 C_c=P_{W_R}D_c\big|_{W_R},
\]
then norms of the $p$-sieved vectors determine the real part of $C_{1/p}$ through
\[
 \left\|w-p^{-1/2}D_{1/p}w\right\|^2
 =\left(1+p^{-1}\right)\norm w^2
 -2p^{-1/2}\Re\ip{w}{D_{1/p}w}.
\]
Consequently a quantitative lower bound for the $p$-sieved co-Poisson norm, uniform over the relevant source class and strong enough to survive the three-space Schur complement, would give new control of the cross-scale Gram operator appearing above. This is a narrower arithmetic target than an unrestricted positivity claim.


\section{Why a finite Euler-product sieve fails in the Hilbert norm}
The one-prime identity \eqref{eq:psieve} suggests iterating over all primes up to a cutoff $z$. Put
\[
 \mathfrak P(z)=\prod_{p\le z}p,
 \qquad
 \mathcal S_z=\prod_{p\le z}\left(I-p^{-1/2}D_{1/p}\right).
\]
For an admissible compact source $h=\mathcal D\phi$ one obtains exactly
\begin{equation}\label{eq:finite-sieve}
 \mathcal S_zE(h)(x)
 =\sum_{(n,\mathfrak P(z))=1}h(nx).
\end{equation}
For each fixed $x>0$ the right side eventually equals $h(x)$, because only finitely many integers occur and every fixed $n>1$ eventually has a prime divisor at most $z$. Thus the finite Euler-product sieve converges pointwise to the bare source. The Hilbert norm behaves very differently.

Let
\[
 \Psi(s)=\int_0^\infty h(t)t^{s-1}\,dt.
\]
Admissibility gives $\Psi(1)=0$ (and also $\Psi(0)=0$). If $h\ne0$, analyticity implies that there is a least integer $m\ge1$ such that $\Psi^{(m)}(1)\ne0$. Define
\[
 R_z(x)=\mathcal S_zE(h)(x)-h(x).
\]
Choose a compact interval $U\Subset(0,\inf\operatorname{supp}h)$. For $x=u/z$, $u\in U$, and $z$ sufficiently large, every integer contributing to $R_z(u/z)$ lies between $z$ and $C_Uz<z^2$. Hence an integer coprime to $\mathfrak P(z)$ in this range is necessarily prime. The prime number theorem with a smooth weight then gives, uniformly for $u\in U$,
\begin{equation}\label{eq:prime-shell-asymp}
 R_z(u/z)
 =\frac{(-1)^m\Psi^{(m)}(1)}{u}\,
   \frac{z}{(\log z)^{m+1}}
 +O_{h,U}\!\left(\frac{z}{(\log z)^{m+2}}\right).
\end{equation}
Indeed, the main term is
\[
 \frac zu\int \frac{h(t)}{\log(zt/u)}\,dt,
\]
and expansion of the denominator in powers of $1/\log z$ kills precisely the first $m$ logarithmic moments. Consequently
\begin{equation}\label{eq:sieve-R2-divergence}
 \int_{U/z}|R_z(x)|^2\,dx
 \sim
 |\Psi^{(m)}(1)|^2
 \left(\int_U\frac{du}{u^2}\right)
 \frac{z}{(\log z)^{2m+2}},
\end{equation}
which tends to $+\infty$.

Thus the sharp finite Euler-product inversion has the paradoxical behavior
\[
 \mathcal S_zE(h)(x)\to h(x)\quad\text{for every fixed }x>0,
 \qquad
 \norm{\mathcal S_zE(h)-h}_2\not\to0;
\]
in fact a fixed nonzero admissible source has a diverging prime-shell contribution. This is a scoped no-go theorem for the raw finite-prime sieve, not for all smoothed or nonlocal arithmetic inverses. It also explains why imposing any \emph{fixed} finite number of additional logarithmic moments cannot repair the mechanism: powers of $\log z$ cannot beat the factor $z^{1/2}$ in the $L^2$ norm. A successful Euler-factor construction would need a genuinely nonlocal tail or a number of matched conditions growing with the scale.

\section{The arithmetic-reflection defect as a dual zero channel}
Let $w$ be a hypothetical off-line zero and put $w^\sharp=1-\overline w$. The ambient de Branges reflection multiplier
\[
 B_w(s)=\frac{s-w^\sharp}{s-w}
\]
has modulus one on the critical line. The compact-division argument below shows that, at an off-line zero of multiplicity $m$, the finite-cutoff leakage
\[
 K_{w,R}=P_{Q_R}B_w\big|_{W_R}
\]
has rank one and factors as
\[
 K_{w,R}f=\ell_{w,R}(f)q_{w,R},
 \qquad
 \ell_{w,R}(E(\psi))=\Psi(w).
\]
At a simple zero the vector $q_{w,R}$ is characterized by vanishing completed Mellin values at all other zero constraints and a fixed nonzero value at $w$. This is the finite-cutoff analogue of the biorthogonal zero direction appearing in Burnol's complete/minimal systems; at the critical cutoff $a=1$ the dual system is represented, up to triangular changes within a multiple zero, by the functions $\zeta(s)/(s-\rho)^k$ \cite{Burnol2004}.

The rank-one channel has an exact prime transport law. Dilation covariance of the source gives
\[
 \ell_{w,pR}(D_pf)=p^{w-1/2}\ell_{w,R}(f),
 \qquad
 \ell_{w,pR}(D_{1/p}f)=p^{1/2-w}\ell_{w,R}(f).
\]
Since $B_w$ commutes with dilation and $D_{\pm1/p}W_R\subset W_{pR}$, comparison of the two rank-one factorizations yields
\begin{equation}\label{eq:dual-prime-transport}
 \boxed{P_{Q_{pR}}D_pq_{w,R}=p^{w-1/2}q_{w,pR}},
 \qquad
 \boxed{P_{Q_{pR}}D_{1/p}q_{w,R}=p^{1/2-w}q_{w,pR}}.
\end{equation}
Writing $d=|\Re w-1/2|$, contractivity gives
\begin{equation}\label{eq:q-upper}
 p^d\norm{q_{w,pR}}\le\norm{q_{w,R}}.
\end{equation}
For a simple zero, $M(q_{w,R})(w)$ is the fixed nonzero number
\[
 (w-w^\sharp)c_0(w)\zeta'(w),
\]
when the normalizing source has $\Psi(w)=1$. Cauchy--Schwarz against the zero-evaluation Riesz vector, together with the off-line evaluator growth in Assumption~\ref{ass:evaluators} $\norm{v_w^R}\asymp_w R^d$, therefore gives
\begin{equation}\label{eq:q-lower}
 \norm{q_{w,R}}\gg_w R^{-d}.
\end{equation}
Iterating \eqref{eq:q-upper} supplies the matching upper estimate along geometric cutoffs $R=p^nR_0$. Hence
\begin{equation}\label{eq:q-asymp}
 \norm{q_{w,p^nR_0}}\asymp_{w,p,R_0}(p^nR_0)^{-d}.
\end{equation}

Let $r_{w,R}\in W_R$ be the Riesz vector of the source-moment functional $\ell_{w,R}$. The same two dilation inclusions imply
\[
 \norm{r_{w,pR}}\ge p^d\norm{r_{w,R}}.
\]
Since $B_w$ is unitary on its ambient domain,
\[
 \norm{K_{w,R}}=\norm{r_{w,R}}\norm{q_{w,R}}\le1.
\]
Combining these estimates with \eqref{eq:q-lower} yields, for a hypothetical simple off-line zero,
\begin{equation}\label{eq:rankone-persistent}
 0<c_{w,p,R_0}
 \le \norm{K_{w,p^nR_0}}\le1
 \qquad(n\ge0).
\end{equation}
Thus the arithmetic-reflection defect does not disappear in operator norm even though every fixed source can be repaired asymptotically in the enlarging arithmetic spaces. The source-moment Riesz vector grows like $R^d$, the dual zero channel shrinks like $R^{-d}$, and the product retains order-one strength. This gives a precise moving-boundary explanation for why fixed-input convergence is insufficient.

\subsection{Full multiplicity channel: exact jet transport and persistent leakage volume}
The rank-one calculation above uses one division by $s-w$ and therefore exposes only the first source moment, even when the zero $w$ has higher multiplicity.  To treat a zero of multiplicity $m$ without discarding its jet structure, use the $m$-fold ambient reflection
\[
 \mathcal B_w^{(m)}(s)=\left(\frac{s-w^\sharp}{s-w}\right)^m,
 \qquad w^\sharp=1-\overline w,
\]
and define
\[
 \mathcal K_{w,R}^{(m)}=P_{Q_R}\mathcal B_w^{(m)}\big|_{W_R}.
\]
For $f=E(\psi)$ write $\Psi=M_0\psi$ and introduce the factorially normalized source jet
\[
 \mathcal J_{w,R}f=
 \left(\Psi(w),\Psi'(w),\ldots,\frac{\Psi^{(m-1)}(w)}{(m-1)!}\right)^T\in\mathbb C^m.
\]
Because $\zeta$ has order exactly $m$ at $w$, these functionals extend continuously from the source core to $W_R$: they are fixed triangular linear combinations of the completed-Mellin output derivatives of orders $m,\ldots,2m-1$.

\begin{theorem}[Exact multiplicity-$m$ leakage factorization]\label{thm:mleak}
Conditionally suppose that $w$ is a nontrivial zero of multiplicity $m$ with $\Re w\ne1/2$.  Then there is an injective map
\[
 \mathcal Q_{w,R}:\mathbb C^m\longrightarrow Q_R
\]
such that
\begin{equation}\label{eq:mfactor}
 \mathcal K_{w,R}^{(m)}=\mathcal Q_{w,R}\mathcal J_{w,R},
 \qquad
 \operatorname{rank}\mathcal K_{w,R}^{(m)}=m.
\end{equation}
If $h=\log p$ and $N$ is the nilpotent shift on factorial jet coordinates, set
\[
 M_p=p^{w-1/2}e^{hN}.
\]
Then the leakage basis obeys the exact prime transport laws
\begin{equation}\label{eq:mtransport}
 P_{Q_{pR}}D_p\mathcal Q_{w,R}=\mathcal Q_{w,pR}M_p,
 \qquad
 P_{Q_{pR}}D_{1/p}\mathcal Q_{w,R}=\mathcal Q_{w,pR}M_p^{-1}.
\end{equation}
\end{theorem}

\begin{proof}
Write locally
\[
 \zeta(s)=(s-w)^m a_w(s),\qquad a_w(w)\ne0.
\]
For a core vector $f=E(\psi)$,
\[
 M(\mathcal B_w^{(m)}f)(s)
 =c_0(s)a_w(s)(s-w^\sharp)^m\Psi(s).
\]
The first $m$ completed-Mellin derivatives at $w$ are therefore an invertible lower-triangular transformation of the source jet $\mathcal J_{w,R}f$. Every element of $W_R$ has those first $m$ derivatives equal to zero.
This alone does not prove factorization. If the source jet vanishes,
Lemma~\ref{lem:compact-division} shows that
$\Psi(s)/(s-w)^m$ is again the transform of an admissible compact source.
Multiplication by $(s-w^\sharp)^m$ preserves that class, so
$\mathcal B_w^{(m)}f\in W_R$. The kernel inclusion extends to $W_R$ by
Lemma~\ref{lem:jet-density}, using continuity of the jet and of the ambient
unitary multiplier. Therefore the leakage factors through the source jet.
The jet map is onto by Lemma~\ref{lem:jet-density}. The displayed triangular
output map is invertible because $w\ne w^\sharp$, so its factor into $Q_R$
is injective and the rank is exactly $m$. On the critical line the
multiplier is the identity and the leakage is zero; that case is excluded.

Dilation gives
\[
 \mathcal J_{w,pR}(D_pf)=p^{w-1/2}e^{hN}\mathcal J_{w,R}f=M_p\mathcal J_{w,R}f.
\]
The multiplier $\mathcal B_w^{(m)}$ commutes with dilations, and $D_pW_R,D_{1/p}W_R\subset W_{pR}$. Projecting the identity $D_{\pm p}\mathcal B_w^{(m)}f=\mathcal B_w^{(m)}D_{\pm p}f$ to $Q_{pR}$ gives \eqref{eq:mtransport}.  Notice that $\det e^{hN}=1$, hence
\[
 |\det M_p|=p^{m(\Re w-1/2)}.
\]
\end{proof}

The theorem has a basis-independent consequence.  Let
\[
 G_R=\mathcal Q_{w,R}^*\mathcal Q_{w,R},
 \qquad
 H_R=\mathcal J_{w,R}\mathcal J_{w,R}^*.
\]
Thus $G_R$ is the Gram matrix of the dual leakage directions and $H_R$ is the Gram matrix of the Riesz representers of the source-jet functionals.  Contractivity of the two projected dilation maps in \eqref{eq:mtransport} gives
\begin{align}
 M_p^*G_{pR}M_p&\preceq G_R,
 &M_p^{-*}G_{pR}M_p^{-1}&\preceq G_R,\label{eq:Gineq}\\
 H_{pR}&\succeq M_pH_RM_p^*,
 &H_{pR}&\succeq M_p^{-1}H_RM_p^{-*}.\label{eq:Hineq}
\end{align}
Writing
\[
 d=\left|\Re w-\frac12\right|,
\]
taking determinants yields
\begin{equation}\label{eq:det-one-sided}
 \det G_{pR}\le p^{-2md}\det G_R,
 \qquad
 \det H_{pR}\ge p^{2md}\det H_R.
\end{equation}

The remaining quantitative leakage conclusions in this section assume Assumption~\ref{ass:evaluators}. Only a one-sided ambient jet bound is needed.  Let $V_{w,R}$ be the Gram matrix of the first $m$ completed-Mellin evaluation jets at $w$ on the ambient Sonine space $H_R$.  For the raw representatives specified in Assumption~\ref{ass:evaluators}, the fixed finite raw jet family has the exact cutoff scaling
\[
 R^{1/2-w}e^{(\log R)N}
 \bigl(1_{t\ge1}t^{w-1},\,1_{t\ge1}t^{w-1}\log t,\ldots,\,
 1_{t\ge1}t^{w-1}(\log t)^{m-1}\bigr)
\]
for $\Re w<1/2$; for $\Re w>1/2$ the Fourier involution gives the reflected model.  The nilpotent factor has determinant one, and the fixed raw Gram matrix is finite.  Orthogonal Sonine projection can only decrease the Gram matrix in Loewner order.  Therefore
\begin{equation}\label{eq:jetgram-upper}
 \boxed{\det V_{w,R}\ll_{w,m}R^{2md}.}
\end{equation}
No relative $o(1)$ estimate for the whole projected jet family is used here.

Let $A_w$ be the fixed invertible triangular matrix taking the source jet to the first $m$ output jets in the proof of Theorem~\ref{thm:mleak}.  The Gram determinant inequality applied to
\[
 \Gamma_{w,R}\mathcal Q_{w,R}=A_w
\]
gives
\[
 |\det A_w|^2\le \det V_{w,R}\det G_R.
\]
Combining this with \eqref{eq:jetgram-upper} and the upper bound obtained by iterating \eqref{eq:det-one-sided}, along every geometric sequence $R=p^nR_0$,
\begin{equation}\label{eq:Gdet-asymp}
 \boxed{\det G_R\asymp_{w,m,p,R_0}R^{-2md}.}
\end{equation}

The nonzero singular values of $\mathcal K_{w,R}^{(m)}=\mathcal Q_{w,R}\mathcal J_{w,R}$ satisfy
\begin{equation}\label{eq:singprod}
 \prod_{j=1}^m\sigma_j(\mathcal K_{w,R}^{(m)})^2
 =\det G_R\det H_R.
\end{equation}
Since $\mathcal B_w^{(m)}$ is unitary on the ambient critical-line space and $P_{Q_R}$ is contractive, every $\sigma_j\le1$.  Equations \eqref{eq:Hineq} and \eqref{eq:Gdet-asymp} imply
\[
 \det H_R\gg R^{2md},
\]
whereas \eqref{eq:singprod}, $\sigma_j\le1$, and the lower bound for $\det G_R$ give
\[
 \det H_R\ll R^{2md}.
\]
Thus
\begin{equation}\label{eq:Hdet-asymp}
 \boxed{\det H_R\asymp_{w,m,p,R_0}R^{2md}}
\end{equation}
and, most importantly,
\begin{equation}\label{eq:singular-persistent}
 \boxed{
 0<c_{w,m,p,R_0}
 \le \prod_{j=1}^m\sigma_j(\mathcal K_{w,R}^{(m)})
 \le1
 }
 \qquad(R=p^nR_0).
\end{equation}
Because each singular value is at most one, \eqref{eq:singular-persistent} also bounds the smallest active singular value away from zero.  Hence the entire multiplicity-$m$ arithmetic-reflection leakage channel remains uniformly nondegenerate along geometric cutoffs.

\begin{remark}
This is the multiplicity analogue of the simple-zero compensation law.  The dual leakage volume decays like $R^{-md}$, the source-jet volume grows like $R^{md}$, and the nilpotent logarithmic factors cancel from the determinant.  Higher multiplicity therefore does not make the off-line defect disappear; it converts the scalar renormalization into a Jordan-coupled $m$-channel renormalization.
\end{remark}

\begin{remark}[What this still does not prove]
The persistent leakage theorem is conditional on the existence of an off-line zero and describes what the arithmetic reflection channel would then do.  It is not itself a contradiction: a unitary ambient reflection can have a uniformly nontrivial finite-dimensional leakage into $Q_R$.  To deduce RH one still needs an independent arithmetic or geometric principle forcing this leakage channel to vanish, become asymptotically singular, or conflict with a positive fixed-space prime action.  No such principle is proved here.
\end{remark}


\section{Bridge to the prime-compression defect}
The persistent reflection leakage is not merely an auxiliary finite-dimensional construction.  It is itself the defect of a canonical compression of the unitary zero-reflection multiplier, and it intertwines exactly with the cross-scale prime compression.

Let $R_R$ denote orthogonal projection onto $W_R$, and for a fixed off-line parameter $w$ define
\[
 A_{w,R}=R_RB_w\big|_{W_R},
 \qquad
 K_{w,R}=P_RB_w\big|_{W_R}.
\]
At a zeta zero $w$, every $f\in W_R$ has completed Mellin value zero at $w$, so Burnol's division theorem places $B_wf$ back in $H_R=W_R\oplus Q_R$.  Since $B_w$ is unitary on the ambient critical-line $L^2$ space,
\begin{equation}\label{eq:reflection-defect}
 \boxed{I_{W_R}-A_{w,R}^*A_{w,R}=K_{w,R}^*K_{w,R}.}
\end{equation}
Thus the rank-one leakage for one reflection, and the multiplicity-$m$ leakage for $B_w^m$, are canonical positive defect operators of the arithmetic reflection compression.

Now define the cross-scale prime compression
\[
 C_{R,p}=P_{pR}D_p\big|_{Q_R}:Q_R\longrightarrow Q_{pR}.
\]
Because $D_pW_R\subset W_{pR}$ and $B_wD_p=D_pB_w$, for every $f\in W_R$ one has
\begin{equation}\label{eq:leak-intertwine}
 \boxed{K_{w,pR}D_pf=C_{R,p}K_{w,R}f.}
\end{equation}
The same identity holds for the multiplicity-$m$ leakage operator $\mathcal K_{w,R}^{(m)}$.  In the jet basis of Theorem~\ref{thm:mleak},
\[
 C_{R,p}\mathcal Q_{w,R}=\mathcal Q_{w,pR}M_p,
 \qquad M_p=p^{w-1/2}e^{(\log p)N}.
\]
Consequently the cross-scale defect restricted to the leakage channel has the exact Gram form
\begin{equation}\label{eq:cross-defect-gram}
 \boxed{
 \mathcal Q_{w,R}^*(I-C_{R,p}^*C_{R,p})\mathcal Q_{w,R}
 =G_R-M_p^*G_{pR}M_p\succeq0.
 }
\end{equation}
This is the first exact positive operator identity in the project in which the persistent arithmetic-reflection channel and the prime compression appear in the same formula.

The fixed-space compression is nevertheless different.  Since $D_pQ_R\subset H_{pR}$, write
\[
 F_{R,p}=R_{pR}D_p\big|_{Q_R}.
\]
Then
\begin{equation}\label{eq:cross-defect-factor}
 \boxed{I-C_{R,p}^*C_{R,p}=F_{R,p}^*F_{R,p}.}
\end{equation}
Let
\[
 \mathsf R_{R,p}=P_R\big|_{Q_{pR}}:Q_{pR}\to Q_R,
 \qquad
 E_{R,p}=P_RR_{pR}D_p\big|_{Q_R}:Q_R\to Q_R.
\]
Decomposing $D_pq$ into its $Q_{pR}$ and $W_{pR}$ components and then projecting to $Q_R$ gives the exact return formula
\begin{equation}\label{eq:return-alias}
 \boxed{T_{R,p}=\mathsf R_{R,p}C_{R,p}+E_{R,p}.}
\end{equation}
Moreover $E_{R,p}=P_RF_{R,p}$, hence
\begin{equation}\label{eq:alias-bound}
 \boxed{E_{R,p}^*E_{R,p}\preceq I-C_{R,p}^*C_{R,p}.}
\end{equation}
The term $E_{R,p}$ is the arithmetic \emph{aliasing} created by projecting the newly available relations in $W_{pR}$ back into the older quotient $Q_R$.  It vanishes if the quotients are nested in the required direction, but no such nesting is available here.

Equations \eqref{eq:cross-defect-gram}--\eqref{eq:alias-bound} isolate the precise missing bridge.  Reflection leakage controls the cross-scale prime defect, while the RH-sensitive fixed-space operator also depends on the angle map $\mathsf R_{R,p}$ and the aliasing $E_{R,p}$.  Hilbert-space positivity alone does not compare the two defects.

\begin{proposition}[No generic domination between cross-scale and fixed-space defects]\label{prop:no-domination}
Even with nested ambient spaces $H_R\subset H_{pR}$, nested arithmetic spaces $W_R\subset W_{pR}$, and a unitary $U$ satisfying $UW_R\subset W_{pR}$, there is no universal operator inequality in either direction between the defects of
\[
 P_{Q_{pR}}U\big|_{Q_R}
 \quad\text{and}\quad
 P_{Q_R}U\big|_{Q_R}.
\]
\end{proposition}
\begin{proof}
Take an orthonormal basis $e_1,e_2,e_3$, set
\[
 H_R=\operatorname{span}\{e_1,e_2\},\quad W_R=\operatorname{span}\{e_2\},\quad Q_R=\operatorname{span}\{e_1\},
\]
and
\[
 H_{pR}=\mathbb C^3,\qquad
 W_{pR}=\operatorname{span}\{e_2,\cos\theta\,e_1+\sin\theta\,e_3\}.
\]
Then
\[
 Q_{pR}=\operatorname{span}\{-\sin\theta\,e_1+\cos\theta\,e_3\}.
\]
For $U=I$, one has $UW_R\subset W_{pR}$, the fixed-space compression on $Q_R$ is the identity, while the cross-scale compression has norm $|\sin\theta|$.  Thus the fixed-space defect is zero while the cross-scale defect is $\cos^2\theta$.

Conversely choose a unitary $U$ fixing $e_2$ and sending $e_1$ to the unit generator of $Q_{pR}$.  Again $UW_R\subset W_{pR}$, but now the cross-scale compression is isometric while the fixed-space compression has norm $|\sin\theta|$.  The two defects therefore reverse their sizes.  This rules out a comparison based only on positivity, nesting, and unitary transport.
\end{proof}

The proposition does not use the additional arithmetic structure of Burnol's spaces, so it is not a no-go theorem for the RH program.  Its role is diagnostic: any successful passage from \eqref{eq:cross-defect-gram} to the fixed-space defect must exploit a genuinely arithmetic estimate on $\mathsf R_{R,p}$ or $E_{R,p}$, not a generic projection inequality.


\section{Infinitesimal dilation and the exact source-range defect}
The finite-prime picture has an infinitesimal counterpart which unifies the source-moment obstruction, the reflection leakage, and the Jordan transport law.  Let
\[
 (\mathsf A f)(x)=-\frac12 f(x)-x f'(x)
\]
be the skew-adjoint generator of the ambient unitary dilation group $D_{e^t}=e^{t\mathsf A}$.  On the uncompleted Mellin transform,
\begin{equation}\label{eq:generator-mellin}
 M_0(\mathsf A f)(s)=(s-\tfrac12)M_0f(s).
\end{equation}
Let $\mathcal A_R=\mathcal D V_R$ be the admissible compact source class and let
\[
 \mathfrak C_R=E(\mathcal A_R),
\]
which is dense in $W_R$ by definition.  Dilation covariance implies
\[
 \mathsf A E(\psi)=E(\mathsf A\psi),
 \qquad
 \mathsf A\mathcal D=\mathcal D\mathsf A.
\]
Moreover, if $\Psi(0)=\Psi(1)=0$, then $(s-\tfrac12)\Psi(s)$ has the same two zero moments.  Hence
\begin{equation}\label{eq:generator-core-invariance}
 \boxed{\mathsf A\mathfrak C_R\subset\mathfrak C_R.}
\end{equation}
The restriction $\mathsf A_{W,R}=\mathsf A|_{\mathfrak C_R}$ is therefore a densely defined skew-symmetric operator in $W_R$.

Fix $w$ in the open strip, put $\lambda_w=w-\tfrac12$, and let $r\ge1$.  For $f=E(\theta)$ with $\Theta=M_0\theta$, equation \eqref{eq:generator-mellin} gives
\[
 (\mathsf A_{W,R}-\lambda_w)^r E(\eta)=E(\theta)
 \quad\Longleftrightarrow\quad
 (s-w)^r H(s)=\Theta(s),
\]
where $H=M_0\eta$.  Lemma~\ref{lem:compact-division} shows that $H=\Theta/(s-w)^r$ is again the Mellin transform of an admissible source in the same cutoff exactly when
\[
 \Theta(w)=\Theta'(w)=\cdots=\Theta^{(r-1)}(w)=0.
\]
Thus the source-core range is exactly
\begin{equation}\label{eq:generator-range-core}
 \operatorname{Ran}(\mathsf A_{W,R}-\lambda_w)^r
 =\left\{E(\theta):\theta\in\mathcal A_R,\ \Theta^{(j)}(w)=0\ (0\le j<r)\right\}.
\end{equation}

At a zeta zero $w$ of multiplicity $m$, the source-jet functionals
\[
 \ell_j(E(\theta))=\frac{\Theta^{(j)}(w)}{j!},\qquad 0\le j<m,
\]
extend continuously to $W_R$.  Indeed they are triangular linear combinations of completed-Mellin output derivatives of orders $m,\ldots,2m-1$, with nonzero leading coefficient determined by $\zeta^{(m)}(w)$.  They are linearly independent: choose one interior source $\eta$ with $H(w)\ne0$ and several sufficiently small dilations $D_{a_k}\eta$; the jet matrix is a nonzero triangular factor times a Vandermonde matrix in $\log a_k$.  Therefore
\begin{theorem}[Exact generator range defect]\label{thm:generator-defect}
Conditionally suppose that $w$ is a zeta zero of multiplicity $m$.  Then
\begin{equation}\label{eq:generator-codim}
 \boxed{
 \overline{\operatorname{Ran}(\mathsf A_{W,R}-\lambda_w)^m}
 =\bigcap_{j=0}^{m-1}\ker\ell_j,
 \qquad
 \operatorname{codim}_{W_R}=m.
 }
\end{equation}
The codimension is exact at every finite cutoff for which the source class is nontrivial.
\end{theorem}
\begin{proof}
The range on the left is the range of the algebraic power on the invariant
source core, not an asserted power of a closed operator. Compact division
gives \eqref{eq:generator-range-core}. The continuous independent source
jets have a surjective finite-dimensional joint map. Lemma~\ref{lem:jet-density}
identifies the closure of its core kernel with its full kernel, which has
codimension $m$.
\end{proof}

This theorem identifies the multiplicity-$m$ leakage channel as a higher-order range defect rather than an ad hoc correction.  The basic deficiency index is one; multiplicity $m$ appears through a length-$m$ generalized deficiency chain, as made precise below.  Indeed on the Mellin line
\begin{equation}\label{eq:cayley-reflection}
 B_w(s)=\frac{s-w^\sharp}{s-w}
 =1+\frac{w-w^\sharp}{s-w},
 \qquad w^\sharp=1-\overline w,
\end{equation}
so $B_w$ is the Cayley--Blaschke transform of the dilation generator at the spectral parameter $\lambda_w$.  The condition $\ell_0=0$ is precisely the condition under which the resolvent term in \eqref{eq:cayley-reflection} remains inside the same compact arithmetic source space.  At a zeta zero, the output factor $\zeta(w)=0$ guarantees that $B_wW_R\subset H_R$, and the unresolved source coordinate is therefore forced to appear internally as the leakage map $K_{w,R}:W_R\to Q_R$.  Repeating the division $m$ times produces the full jet channel of Theorem~\ref{thm:mleak}.

The prime transport matrix from that theorem is exactly the exponential of the jet generator:
\begin{equation}\label{eq:jet-exponential}
 M_p
 =p^{w-1/2}e^{(\log p)N}
 =\exp\!\left((\log p)\bigl((w-\tfrac12)I+N\bigr)\right).
\end{equation}
Thus the persistent finite-cutoff leakage, its Jordan factors, and its prime scaling are the finite-time form of the same infinitesimal range defect.

This also rules out a tempting shortcut.  The ambient generator $\mathsf A$ is skew-adjoint, but its arithmetic restriction on $W_R$ is not automatically maximal skew-symmetric: equation \eqref{eq:generator-codim} exhibits a finite-codimensional resolvent obstruction.  Merely invoking the ambient unitary dilation group therefore cannot yield a Hilbert--Polya proof.  One would have to show that the relevant deficiency channel is eliminated by an additional arithmetic boundary condition while retaining every zero mode.  That is precisely the compatibility theorem still missing from the fixed-space prime construction.


\section{Boundary-triplet audit and an exact zero-incidence criterion}
The infinitesimal formulation permits a sharper operator-theoretic audit.  Let
\[
 \mathsf S_R=\overline{\mathsf A_{W,R}}
\]
be the closure in $W_R$ of the skew-symmetric arithmetic dilation generator on the compact source core.  The first conclusion is that the basic deficiency dimension is one, not the multiplicity of a zeta zero.

For $w$ with $0<\Re w<1$ and $\zeta(w)\ne0$, define the source-evaluation functional on $W_R$ by
\[
 \ell_w(f)=\frac{\mathcal M f(w)}{c_0(w)\zeta(w)}.
\]
It is continuous because completed Mellin evaluation is continuous on $H_R$.  At a zeta zero the same functional extends by the derivative formula already used in the source-moment audit.  On the source core,
\[
 \operatorname{Ran}(\mathsf A_{W,R}-(w-\tfrac12))
 =\ker\ell_w\cap\mathfrak C_R.
\]
Since one may correct a core approximant by a fixed core vector with source moment one, closure gives
\begin{equation}\label{eq:first-range-kernel}
 \overline{\operatorname{Ran}(\mathsf S_R-(w-\tfrac12))}
 =\ker\ell_w.
\end{equation}
Choose $0<a<1/2$ and the real points $w_\pm=1/2\pm a$, where $\zeta(w_\pm)\ne0$.  Because $\mathsf S_R$ is skew-symmetric,
\[
 \norm{(\mathsf S_R\mp a)f}^2=\norm{\mathsf S_Rf}^2+a^2\norm f^2,
\]
so these ranges are closed.  Equation \eqref{eq:first-range-kernel} therefore implies
\[
 \dim\ker(\mathsf S_R^*-a)=\dim\ker(\mathsf S_R^*+a)=1.
\]
Thus $\mathsf S_R$ has deficiency indices $(1,1)$.  If $r_w\in W_R$ denotes the Riesz vector of $\ell_w$, then
\begin{equation}\label{eq:deficiency-riesz}
 r_w\in\ker\!\left(\mathsf S_R^*-(\overline w-\tfrac12)\right).
\end{equation}
The involution $G$ preserves $W_R$ and anti-commutes with the dilation generator.  On source moments it sends $\Psi(s)$ to $\Psi(1-s)$, hence
\begin{equation}\label{eq:G-deficiency}
 Gr_w=r_{1-w}.
\end{equation}
The deficiency vectors are total: if $f\in W_R$ is orthogonal to $r_w$ on any open set of parameters avoiding zeta zeros, then the entire completed Mellin transform of $f$ vanishes on an open set and therefore vanishes identically.  Hence $\mathsf S_R$ is simple in the standard symmetric-operator sense.

Standard Liv\v{s}ic/von Neumann theory therefore supplies a one-parameter family of maximal skew-adjoint extensions of $\mathsf S_R$.  If $e_\pm$ are normalized generators of $\ker(\mathsf S_R^*\mp a)$ and
\[
 f=f_0+c_+e_++c_-e_-,\qquad f_0\in\mathcal D(\mathsf S_R),
\]
then the boundary form is
\begin{equation}\label{eq:boundary-form}
 \ip{\mathsf S_R^*f}{g}+\ip{f}{\mathsf S_R^*g}
 =2a\left(c_+\overline{d_+}-c_-\overline{d_-}\right).
\end{equation}
A maximal skew-adjoint extension is the restriction of $-\mathsf S_R^*$ (not $\mathsf S_R^*$) to the domain specified by one fixed unitary boundary relation $c_-=e^{i\vartheta}c_+$.  This statement is ordinary extension theory; it does not identify any such extension with the zeta zeros.

The zeta condition enters through an additional ambient incidence statement.  Fix an off-critical parameter $w$ and choose an admissible source $\eta_w\in\mathcal A_R$ with Mellin transform $H_w$ satisfying $H_w(w)\ne0$; the bump construction from the source-moment repair gives such a source.  Put
\[
 f_w=E(\eta_w)\in W_R,
 \qquad
 B_w(s)=\frac{s-(1-\overline w)}{s-w}.
\]
Then
\begin{theorem}[Exact zero-incidence criterion]\label{thm:zero-incidence}
For $0<\Re w<1$, $\Re w\ne1/2$,
\begin{equation}\label{eq:zero-incidence}
 \boxed{\zeta(w)=0\quad\Longleftrightarrow\quad B_wf_w\in H_R.}
\end{equation}
Equivalently, if $S_R^{(H)}$ denotes orthogonal projection onto $H_R$,
\begin{equation}\label{eq:zero-incidence-distance}
 \boxed{\zeta(w)=0\quad\Longleftrightarrow\quad
 \norm{(I-S_R^{(H)})B_wf_w}=0.}
\end{equation}
At a zero the non-arithmetic component therefore lies entirely in $Q_R$; away from a zero it has a nonzero $H_R^\perp$ component.
\end{theorem}
\begin{proof}
The completed Mellin transform of $f_w$ is
\[
 \mathcal M f_w(s)=c_0(s)\zeta(s)H_w(s).
\]
If $\zeta(w)=0$, then $\mathcal M f_w(w)=0$, and Burnol's ambient division theorem places $B_wf_w$ in $H_R$.  Conversely, if $B_wf_w\in H_R$, its completed Mellin transform is entire.  On the critical line it agrees with
\[
 \frac{s-(1-\overline w)}{s-w}\,c_0(s)\zeta(s)H_w(s),
\]
and hence with its meromorphic continuation.  Since $w\ne1-\overline w$ and $c_0(w)H_w(w)\ne0$, removability of the apparent pole at $s=w$ forces $\zeta(w)=0$.
\end{proof}

The theorem clarifies both the promise and the limitation of the boundary-triplet route.  The source generator $\mathsf S_R$ itself has the familiar deficiency pattern $(1,1)$ and therefore admits fixed maximal skew-adjoint boundary conditions.  But the zeta zeros are not presently characterized as eigenvalues of one of those extensions.  They are characterized by the external incidence condition \eqref{eq:zero-incidence}: the Cayley--Blaschke defect generated from $W_R$ must land in the larger Sonine space $H_R$, and at a zero its residual part is precisely the $Q_R$ leakage channel studied above.  The zero-evaluation vectors themselves live in $Q_R$, not in $W_R$.

Consequently a genuine Hilbert--P\'olya conclusion from this picture still requires one new theorem: identify the incidence condition \eqref{eq:zero-incidence}, or an equivalent condition on the $Q_R$ leakage, with a \emph{single zero-independent} maximal skew-adjoint boundary condition or with the characteristic function of one fixed conservative coupling.  No such identification is proved here.  In particular, choosing a boundary condition separately as a function of $w$ would simply repackage the zero problem and would be circular.

The correction above also explains the role of multiplicity.  The operator $\mathsf S_R$ has one-dimensional basic deficiency spaces.  A zero of multiplicity $m$ produces the codimension-$m$ closure of the algebraic source-core range of $(\mathsf A_{W,R}-(w-1/2))^m$ and the length-$m$ Jordan/source-jet chain of Theorem~\ref{thm:mleak}; it does not change the basic deficiency indices from $(1,1)$.


\section{A nested rank-one defect ladder and defect transfer}
The zero-incidence theorem admits a sharper quantitative form.  There are two nested compression defects associated with the same ambient Blaschke multiplier $B_w$: one for the Sonine space $H_R$, and one for its arithmetic subspace $W_R$.

Let $S_R^{(H)}$ and $R_R$ denote the orthogonal projections onto $H_R$ and $W_R$, respectively, and define
\[
 \Delta^{H}_{w,R}=(I-S_R^{(H)})B_w\big|_{H_R},
 \qquad
 \Delta^{W}_{w,R}=(I-R_R)B_w\big|_{W_R}.
\]
Burnol's division property implies that $\Delta^H_{w,R}$ vanishes on the codimension-one evaluation kernel
\[
 \{f\in H_R:\mathcal M f(w)=0\}.
\]
It is not the zero operator: if $k_w^R$ is the Riesz vector for completed Mellin evaluation, then $B_wk_w^R$ has an uncancelled pole at $w$ and therefore cannot belong to $H_R$.  Hence $\Delta^H_{w,R}$ has rank exactly one.  Put
\[
 d^H_{w,R}=\frac{\Delta^H_{w,R}k_w^R}{\norm{k_w^R}^2}.
\]
Then
\begin{equation}\label{eq:H-defect-factor}
 \boxed{\Delta^H_{w,R}f=\mathcal M f(w)\,d^H_{w,R}\qquad(f\in H_R).}
\end{equation}

Likewise the exact source-division identity shows that $\Delta^W_{w,R}$ vanishes on $\ker\ell_w$.  Thus it also has rank one.  If $r_w^R$ is the Riesz vector of the source-evaluation functional $\ell_w$ on $W_R$, set
\[
 d^W_{w,R}=\frac{\Delta^W_{w,R}r_w^R}{\norm{r_w^R}^2}.
\]
Then
\begin{equation}\label{eq:W-defect-factor}
 \boxed{\Delta^W_{w,R}f=\ell_w(f)\,d^W_{w,R}\qquad(f\in W_R).}
\end{equation}
For $f\in W_R$ the arithmetic Mellin factorization gives
\[
 \mathcal M f(w)=c_0(w)\zeta(w)\ell_w(f).
\]
Projecting \eqref{eq:W-defect-factor} further onto $H_R^\perp$ and comparing with \eqref{eq:H-defect-factor} therefore gives the exact vector identity
\begin{equation}\label{eq:defect-transfer-vector}
 \boxed{
 d^W_{w,R}=q_{w,R}+c_0(w)\zeta(w)d^H_{w,R},
 \qquad q_{w,R}=P_{Q_R}d^W_{w,R}\in Q_R.
 }
\end{equation}
The two summands are orthogonal because
\[
 W_R^\perp=Q_R\oplus H_R^\perp.
\]
Consequently
\begin{equation}\label{eq:defect-transfer-pythagoras}
 \boxed{
 \norm{d^W_{w,R}}^2
 =\norm{q_{w,R}}^2
 +|c_0(w)\zeta(w)|^2\norm{d^H_{w,R}}^2.
 }
\end{equation}

Equation \eqref{eq:defect-transfer-pythagoras} is a defect-conservation law.  The source generator always has one nontrivial Blaschke/Cayley defect direction off the critical line.  At a zeta zero its $H_R^\perp$ component vanishes and the entire defect lies internally in the arithmetic quotient $Q_R$; away from a zero a nonzero portion necessarily escapes the Sonine space.  Thus a zero does not annihilate the defect.  It transfers the defect from the outer Sonine channel to the inner arithmetic channel.

Theorem~\ref{thm:zero-incidence} can therefore be restated as the angle condition
\begin{equation}\label{eq:defect-angle-zero}
 \zeta(w)=0
 \quad\Longleftrightarrow\quad
 d^W_{w,R}\in Q_R
 \quad\Longleftrightarrow\quad
 \frac{|c_0(w)\zeta(w)|^2\norm{d^H_{w,R}}^2}{\norm{d^W_{w,R}}^2}=0.
\end{equation}
This resembles a scalar characteristic or transmission coefficient for the nested pair $W_R\subset H_R$, but no identification with the Liv\v{s}ic characteristic function of $\mathsf S_R$ is asserted here.  Establishing such an identification, with a zero-independent conservative colligation, would be a genuinely new theorem rather than a change of notation.


\section{Exact same-label returns are overstrong; the asymptotic regime is calibrated}
The complete/minimal zero systems suggest a formally natural same-label return between two scale fibers.  A necessary caution is that \emph{exact contractivity} of such a return is stronger than RH, whereas the asymptotic compression used earlier is correctly calibrated.

\begin{proposition}[Exact diagonal contractions force orthogonality]\label{prop:exact-return-overstrong}
Let $\mathcal H$ be a Hilbert space and let $(x_j)_{j\in J}$ have dense linear span.  Suppose a contraction $A$ satisfies
\[
 Ax_j=\lambda_jx_j,\qquad |\lambda_j|=1
\]
for every $j$.  Then $A$ is unitary.  In particular,
\[
 \lambda_i\ne\lambda_j\quad\Longrightarrow\quad \ip{x_i}{x_j}=0.
\]
\end{proposition}
\begin{proof}
For every $j$, equality of norms gives
\[
 0=\norm{x_j}^2-\norm{Ax_j}^2
   =\ip{(I-A^*A)x_j}{x_j}.
\]
Since $I-A^*A\succeq0$, this implies $(I-A^*A)x_j=0$.  Density and continuity give $A^*A=I$, so $A$ is an isometry.  Every $x_j$ belongs to $\operatorname{Ran}A$ because $x_j=\lambda_j^{-1}Ax_j$; hence the range is dense.  The range of an isometry is closed, so $A$ is onto and therefore unitary.  Distinct eigenspaces of a unitary operator are orthogonal.
\end{proof}

Under RH the desired prime multipliers $\lambda_\rho=p^{1/2-\rho}$ all have modulus one.  Thus an \emph{exact} contractive fixed-space operator diagonal on a complete zero system would force exact orthogonality between all zero vectors carrying distinct prime phases.  That orthogonality is an additional requirement not contained in RH.  Consequently exact same-label contractivity is not the correct target for the Burnol zero systems.

This is consistent with Burnol's later study of the biorthogonal family $\zeta(s)/(s-\rho)^k$: the system and its dual are complete and minimal, but bounded unconditional reconstruction is not supplied by completeness/minimality alone.  The hereditary-completeness problem is genuinely subtler; in the admissible mixed systems studied there the possible defect is controlled only up to codimension one.  Thus a formal biorthogonal return must not be treated as a bounded, let alone contractive, operator without a separate theorem.

Under Assumption~\ref{ass:evaluators}, the moving-cutoff construction avoids the preceding overconstraint because fixed critical-line zero kernels become asymptotically orthogonal. The following calculation records the consequence of those quantitative inputs, not a new proof of them.  For two fixed distinct critical-line parameters
\[
 w=\tfrac12+i\tau,\qquad z=\tfrac12+i\sigma,\qquad \tau\ne\sigma,
\]
let $Z_w^R,Z_z^R$ be Burnol's completed-Mellin evaluation vectors.  The raw Sonine representatives satisfy
\[
 C_1(v,w)=\gamma_+(w)v^{-w}+O_w(1),
 \qquad
 C_1(v,z)=\gamma_+(z)v^{-z}+O_z(1)
\]
near zero.  Hence their cross term has leading integral
\[
 \gamma_+(w)\overline{\gamma_+(z)}
 \int_{R^{-2}}^1 v^{-1-i(\tau-\sigma)}\,dv=O_{w,z}(1),
\]
while all remaining local and tail terms are bounded.  The Sonine projection changes each raw evaluator by $O_w(R^{-1})$, so
\begin{equation}\label{eq:critical-cross-bounded}
 \ip{Z_w^R}{Z_z^R}=O_{w,z}(1).
\end{equation}
On the other hand
\[
 \norm{Z_w^R}^2=C_w\log R+O_w(1),
 \qquad
 \norm{Z_z^R}^2=C_z\log R+O_z(1),
 \qquad C_w,C_z>0.
\]
Therefore
\begin{equation}\label{eq:critical-asymptotic-orthogonality}
 \boxed{
 \frac{\ip{Z_w^R}{Z_z^R}}
 {\norm{Z_w^R}\norm{Z_z^R}}
 =O_{w,z}\!\left(\frac1{\log R}\right).}
\end{equation}
For any fixed finite set of distinct critical-line parameters, the normalized Gram matrix therefore tends to the identity.

\begin{corollary}[Finite-mode asymptotic calibration]\label{cor:finite-mode-calibration}
Assume Assumption~\ref{ass:evaluators}. Fix a finite set $F$ of distinct critical-line parameters and unit-modulus numbers $\lambda_w$.  On
\[
 \mathcal V_R(F)=\operatorname{span}\{Z_w^R:w\in F\}
\]
define the diagonal operator $A_R(F)Z_w^R=\lambda_w Z_w^R$.  Then
\[
 \norm{A_R(F)}=1+O_F\!\left(\frac1{\log R}\right).
\]
In contrast, if a finite set contains a parameter with $|\lambda_w|>1$, then every operator having that vector as an exact eigenvector satisfies $\norm A\ge |\lambda_w|>1$.
\end{corollary}
\begin{proof}
After normalizing the basis vectors, equation \eqref{eq:critical-asymptotic-orthogonality} gives a Gram matrix $G_R=I+O_F(1/\log R)$.  The squared norm of the diagonal phase operator is the largest generalized eigenvalue of the pencil $(D^*G_RD,G_R)$.  Here $D=\operatorname{diag}(\lambda_w)$, so this eigenvalue equals $1+O_F(1/\log R)$.  The second assertion is immediate from the norm of an eigenvector.
\end{proof}

The correct fixed-space target is therefore \emph{asymptotic} rather than exact contractivity.  A useful non-circular formulation would construct, from the arithmetic source geometry and not from the zero list itself, operators $R_{R,p,T}$ on a fixed finite spectral window such that
\[
 \limsup_{R\to\infty}\norm{R_{R,p,T}}\le1
\]
for every fixed $T$.  Any hypothetical off-line zero inside the window would force a fixed eigenvalue modulus $p^{1/2-\Re\rho}>1$ and contradict this bound, while under RH the finite-mode Gram calibration above is compatible with norm $1+o(1)$.  The remaining difficulty is to define such a return by a zero-independent co-Poisson/prime formula.  Defining it spectrally from the zero vectors would merely encode the desired conclusion.


\section{Off-line resolvent localization and why the determinant shortcut is tautological}
The fixed Hilbert-space resolvent construction admits a useful extension away from the critical line.  This gives a zero-independent spectral window, but it also shows that determinant or spectral-radius inequalities on that window are endpoints equivalent to zero exclusion rather than independent mechanisms.

\begin{assumption}[Fixed closed arithmetic realization]\label{ass:closed-operator}
A fixed Hilbert space $\mathscr H_{\mathrm B}$ of functions holomorphic
in the nontrivial strip is specified, with continuous off-line evaluations
and an isometric critical-line realization. A closed densely defined
operator $A$ has graph formula
\[
 AF(s)=sF(s)-\eta_{\mathrm B}(F)\zeta(s)
\]
on its stated domain. The finite span of
$F_{\rho,j}=\zeta(s)/(s-\rho)^j$ is a graph core, with
$AF_{\rho,j}=\rho F_{\rho,j}+F_{\rho,j-1}$ for $j>1$ and
$AF_{\rho,1}=\rho F_{\rho,1}$.
The bounded divided-difference resolvents below preserve the realization,
and the resulting off-line meromorphic resolvent has no additional
spectral subspaces beyond the indicated finite root blocks.
\end{assumption}
This is an explicit realization hypothesis, not a completed operator
construction in this manuscript. The inherited critical-line norm, the
domain, closure, and graph-core theorem must be supplied together before
using the spectral claims that follow. Completeness and minimality of the
root system do not imply the graph-core assertion. Later finite root-space
sampling identities can be defined algebraically without this assumption;
calling their source space a Riesz window requires it.


Let $\mathscr H_{\mathrm B}$ denote the fixed right-Mellin Hilbert realization in Assumption~\ref{ass:closed-operator}, let $A$ be the closed operator defined there, and let
\[
 F_{\rho,j}(s)=\frac{\zeta(s)}{(s-\rho)^j},\qquad 1\le j\le m_\rho,
\]
be its complete minimal zero/jet system.  For a point $z$ satisfying
\[
 \Re z\ne\tfrac12,\qquad \zeta(z)\ne0,
\]
define
\begin{equation}\label{eq:complex-resolvent}
 (R_zF)(s)=\frac{F(s)-F(z)\zeta(s)/\zeta(z)}{s-z}.
\end{equation}
Point evaluation at $z$ is continuous on $\mathscr H_{\mathrm B}$.  Moreover, if
$\delta_z=|\Re z-1/2|$, then multiplication by $(s-z)^{-1}$ on the critical line has norm at most $\delta_z^{-1}$.  Since $\zeta(s)/s$ belongs to the ambient Mellin $L^2$ space and $s/(s-z)$ is bounded on the critical line, one also has
\[
 q_z(s):=\frac{\zeta(s)}{s-z}\in L^2.
\]
Consequently $R_z$ is bounded in the ambient norm, with the explicit coarse estimate
\begin{equation}\label{eq:complex-resolvent-bound}
 \norm{R_zF}
 \le\left(\delta_z^{-1}+\frac{\norm{k_z}\norm{q_z}}{|\zeta(z)|}\right)\norm F,
\end{equation}
where $k_z$ is the Riesz vector for evaluation at $z$.

On the finite zero/jet core the numerator in \eqref{eq:complex-resolvent} factors algebraically, so $R_z$ maps the core to itself and is the triangular Jordan resolvent.  Boundedness therefore extends $R_z$ to $\mathscr H_{\mathrm B}$.  The separately assumed graph-core property of the zero/jet system then upgrades the algebraic identities to
\[
 (A-zI)R_z=I,\qquad R_z(A-zI)=I
\]
on their natural domains.  Thus:

\begin{conditionalresult}[Off-line spectrum under a closed-operator realization]\label{thm:off-line-spectrum}
Assume Assumption~\ref{ass:closed-operator}. In the nontrivial strip away from the critical line,
\[
 \sigma(A)\cap\{0<\Re z<1:\Re z\ne\tfrac12\}
 =\{\rho:\zeta(\rho)=0,\ \Re\rho\ne\tfrac12\},
\]
with algebraic multiplicities matching the zeta multiplicities.
\end{conditionalresult}
\begin{proof}
If $\zeta(z)\ne0$, equations \eqref{eq:complex-resolvent}--\eqref{eq:complex-resolvent-bound} and the graph-core argument above show that $z\in\rho(A)$.  If $z=\rho$ is a nontrivial zeta zero of order $m_\rho$, $F_{\rho,1}$ is an eigenvector of $A$ and $F_{\rho,1},\ldots,F_{\rho,m_\rho}$ form the displayed Jordan chain.  The chain cannot be extended: if $G\in D(A)$ satisfied $(A-\rho I)G=F_{\rho,m_\rho}$, then the graph formula for $A$ would force
\[
 G(s)=\frac{\zeta(s)}{(s-\rho)^{m_\rho+1}}
      +c\,\frac{\zeta(s)}{s-\rho}
\]
for some scalar $c$.  The first term has a nonremovable pole at the nontrivial zero $\rho$, whereas elements of $\mathscr H_{\mathrm B}$ have no such pole.  Thus the root space has exactly the zeta multiplicity.  No assertion about spectrum exactly on the critical line is needed for this theorem.
\end{proof}

Under the same realization assumption, a compact contour $\Gamma=\partial\Omega$ lying strictly to one side of the critical line and avoiding zeta zeros defines a genuinely zero-independent Riesz projection
\begin{equation}\label{eq:riesz-window}
 \Pi_\Omega=\frac{1}{2\pi i}\oint_\Gamma (zI-A)^{-1}\,dz.
\end{equation}
Its range is precisely the finite-dimensional sum of the zero/jet blocks inside $\Omega$, without listing those zeros in the definition.

For a prime $p$ define on this window
\begin{equation}\label{eq:window-prime-calculus}
 U_{p,\Omega}
 =\frac{1}{2\pi i}\oint_\Gamma
 p^{1/2-z}(zI-A)^{-1}\,dz.
\end{equation}
On a zero block at $\rho$ this has the exact action
\[
 p^{1/2-\rho}\exp[-(\log p)N].
\]
Thus \eqref{eq:window-prime-calculus} is an exact fixed-space, zero-independent finite-window prime action.

It is tempting to seek $r(U_{p,\Omega})\le1$ or $|\det U_{p,\Omega}|\le1$ for a window strictly left of the critical line.  The determinant calculation shows why this is not by itself a new mechanism.  Holomorphic functional calculus gives
\begin{equation}\label{eq:window-det}
 \det\!\left(U_{p,\Omega}|_{\operatorname{Ran}\Pi_\Omega}\right)
 =\prod_{\rho\in\Omega}p^{m_\rho(1/2-\rho)},
\end{equation}
so
\begin{equation}\label{eq:window-det-modulus}
 \log\left|\det U_{p,\Omega}\right|
 =(\log p)\sum_{\rho\in\Omega}m_\rho(1/2-\Re\rho).
\end{equation}
Equivalently, when the contour avoids the pole and trivial zeros,
\begin{equation}\label{eq:window-det-contour}
 \log\det U_{p,\Omega}
 =\frac{\log p}{2\pi i}
 \oint_\Gamma\left(\frac12-z\right)\frac{\zeta'(z)}{\zeta(z)}\,dz,
\end{equation}
with a consistent branch on the finite-dimensional determinant.  For a left window every summand in \eqref{eq:window-det-modulus} is strictly positive.  Therefore the inequality $|\det U_{p,\Omega}|\le1$ is equivalent to the emptiness of that window.  The spectral-radius inequality has the same issue.  These quantities are useful diagnostics and exact spectral selectors, but a proof still requires an independent arithmetic estimate before taking the determinant.

\subsection*{A global hyponormality shortcut is impossible}
Let $T_{R,p}=P_RD_p|_{Q_R}$ and use Burnol's involution $G$, which preserves $Q_R$ and satisfies $GT_{R,p}G=T_{R,p}^*$.  Its self-commutator
\[
 C_{R,p}=T_{R,p}^*T_{R,p}-T_{R,p}T_{R,p}^*
\]
therefore obeys the exact anti-symmetry
\begin{equation}\label{eq:commutator-antisym}
 \boxed{GC_{R,p}G=-C_{R,p}.}
\end{equation}
Hence $C_{R,p}\succeq0$ or $C_{R,p}\preceq0$ on all of $Q_R$ would force $C_{R,p}=0$, i.e. normality of the compression.  Thus global hyponormality/cohyponormality cannot supply the missing one-sided zero exclusion unless one first proves the much stronger normality statement.  We do not claim a side-dependent asymptotic value for the diagonal of $C_{R,p}$: the existing inward-mode audit determines one of the two norm limits in each off-line half-plane, but not the complementary outward norm required for such a formula.

\section{Comparison with the finite-prime Weil-form program}
A recent and independently developed spectral program of Connes--Consani--Moscovici constructs, from the Weil quadratic form on $[\lambda^{-1},\lambda]$, self-adjoint finite-prime operators whose regularized determinants have only critical-line zeros.  Their finite form on this interval involves only prime-power terms $n\le\lambda^2$, the same multiplicative scale ratio that appears naturally in the Burnol source window $[1/R,R]$.  Their theorem produces a self-adjoint rank-one perturbation of the logarithmic scaling operator in a positive inner product obtained from the truncated Weil form, and hence an entire determinant with real zeros.  The stated missing step is convergence of those finite-prime determinants/spectra toward the Riemann $\Xi$ function; a rigorous proof of that convergence would establish RH \cite{ConnesConsaniMoscovici2025}.

\subsection*{A conditional obstruction to norm-resolvent convergence}
A unitary norm-resolvent comparison has the following necessary condition in the inherited Burnol norm.  Set
\[
 B=A-\frac12 I.
\]
On a critical-line zero block its eigenvalues are $i\gamma$.  If $D_n$ are self-adjoint finite-prime operators and $U_n$ are unitary identifications into one fixed Hilbert space, then $iU_nD_nU_n^*$ are skew-adjoint.  A norm-resolvent limit of skew-adjoint operators is again skew-adjoint.  But the exact boundary form from the resolvent audit gives
\[
 \ip{BF_\rho}{F_\sigma}+\ip{F_\rho}{BF_\sigma}
 =i(\gamma_\rho-\gamma_\sigma)\ip{F_\rho}{F_\sigma}
\]
for critical-line zero vectors.  If a pair of distinct critical ordinates satisfies
$\ip{F_\rho}{F_\sigma}\ne0$, this identity rules out skew-symmetry of $B$
and hence such a norm-resolvent limit to $B$ (equivalently a shifted limit
to $A$). The earlier draft referred to a certified two-zero calculation
but supplied no reproducible certificate. This revision does not assert
that numerical premise. The obstruction is conditional on verifying it
in the specified Hilbert realization. Finite-dimensional approximations
also require a precise common-space realization before norm-resolvent
convergence is a well-typed statement.

This does not rule out a comparison through a changing positive metric.  A nonnormal operator may be similar to a skew-adjoint one in an equivalent Hilbert norm, and the finite-prime Weil construction naturally provides such positive forms at each cutoff.

\subsection*{A finite-window Lyapunov metric criterion}
The preceding observation suggests a weaker and correctly typed bridge.  Let $\Omega$ be a compact Riesz window contained strictly in one half of the critical strip, let $\Pi_\Omega$ be \eqref{eq:riesz-window}, and keep $B=A-\frac12I$.

\begin{theorem}[Uniform metric defect excludes an off-line window]\label{thm:metric-window}
Assume the Riesz realization in Assumption~\ref{ass:closed-operator}. Suppose there are bounded positive operators $G_n$ on $\mathscr H_{\mathrm B}$ and numbers $c_\Omega>0$, $\varepsilon_n\to0$ such that for all sufficiently large $n$,
\begin{equation}\label{eq:metric-coercive}
 \ip{G_nf}{f}\ge c_\Omega\norm{f}^2,
 \qquad f\in\operatorname{Ran}\Pi_\Omega,
\end{equation}
and, on this finite-dimensional spectral subspace,
\begin{equation}\label{eq:metric-lyap}
 \left|\ip{G_nBf}{f}+\ip{G_nf}{Bf}\right|
 \le \varepsilon_n\norm{f}^2.
\end{equation}
Then $\operatorname{Ran}\Pi_\Omega=\{0\}$.
\end{theorem}
\begin{proof}
If $\rho\in\Omega$ is a zero and $F_\rho$ is a value eigenvector, then $BF_\rho=(\rho-\frac12)F_\rho$.  Substituting $f=F_\rho$ in \eqref{eq:metric-lyap} gives
\[
 2\left|\Re\rho-\frac12\right|\ip{G_nF_\rho}{F_\rho}
 \le\varepsilon_n\norm{F_\rho}^2.
\]
By \eqref{eq:metric-coercive},
\[
 2c_\Omega\left|\Re\rho-\frac12\right|\le\varepsilon_n,
\]
which is impossible for large $n$ because $\Omega$ stays a positive distance from the critical line.  Thus no zero block occurs in $\Omega$, and Theorem~\ref{thm:off-line-spectrum} gives $\Pi_\Omega=0$.
\end{proof}

The theorem is deliberately finite-window.  A single globally bounded positive metric making $B$ skew-adjoint would be substantially stronger and would impose global basis geometry on Burnol's complete/minimal zero system.  For RH it is enough to obtain \eqref{eq:metric-coercive}--\eqref{eq:metric-lyap} separately on each compact off-line window.

The finite-prime Weil forms make this target concrete.  On the Connes--Consani--Moscovici finite-dimensional space, under their stated hypotheses, including the required ground-state and boundary properties, subtracting the lowest form eigenvalue and quotienting its null direction makes the corrected logarithmic operator self-adjoint in the resulting positive form \cite{ConnesConsaniMoscovici2025}.  To turn that finite-prime positivity into Theorem~\ref{thm:metric-window}, one would need a zero-independent transport of those forms into the Burnol fixed space for which two properties survive simultaneously: a cutoff-uniform coercive lower bound on $\operatorname{Ran}\Pi_\Omega$, and a Lyapunov defect tending to zero.  This is weaker than direct norm-resolvent convergence and avoids the exact-orthogonality obstruction, but no such transport estimate is proved here.

This formulation also identifies what would count as genuine progress from the finite-prime side.  Mere convergence of eigenvalue lists, or convergence of the determinant after zeros are already matched, is not enough.  The required new input is an operator/form comparison that remains positive before any zero locations are inserted.

\subsection*{A finite-core criterion without a closed-operator assumption}
For zero exclusion the Riesz realization can be bypassed. Let $\mathcal V$
be the algebraic span of finitely many genuine zero chains
$F_{\rho,j}=\zeta(s)/(s-\rho)^j$, $1\le j\le m_\rho$. Distinct partial
fractions are linearly independent, and the formula
\[
 B F_{\rho,1}=(\rho-\tfrac12)F_{\rho,1},\qquad
 B F_{\rho,j}=(\rho-\tfrac12)F_{\rho,j}+F_{\rho,j-1}\quad(j>1)
\]
defines an operator on this finite-dimensional space without a global
domain or graph-core theorem. Equip $\mathcal V$ with any fixed positive
Hilbert norm. The source functional $\eta_{\mathrm B}$ is one on each
$F_{\rho,1}$ and zero on the higher members of each canonical chain.

\begin{theorem}[Finite-core metric criterion and exact transport budget]
\label{thm:v15-finite-metric}
Let $J_n:\mathcal V\to E_n$, let $D_n=D_n^*$, and let $S_n=S_n^*$
be finite matrices. Define
\[
 Y_n=iD_nJ_n-J_nB,\qquad G_n=J_n^*S_nJ_n,
 \qquad \mathcal L_n=G_nB+B^*G_n.
\]
If $G_n\succeq c_n I$ for some $c_n>0$ and
\begin{equation}\label{eq:v15-relative-criterion}
 \frac{\|\mathcal L_n\|}{c_n}\longrightarrow0,
\end{equation}
then every eigenvalue of $B|_{\mathcal V}$ has real part zero.
In particular a space consisting of off-line chains must be zero.
No global closed-operator or quantitative evaluator assumption is used.
The exact finite identity is
\begin{equation}\label{eq:v15-budget}
 \boxed{\mathcal L_n
 =-iJ_n^*[D_n,S_n]J_n-J_n^*S_nY_n-Y_n^*S_nJ_n.}
\end{equation}
Thus a sufficient condition for \eqref{eq:v15-relative-criterion} is
\[
 \|J_n^*[D_n,S_n]J_n\|+2\|J_n^*S_nY_n\|=o(c_n).
\]
Estimating the combined signed expression in \eqref{eq:v15-budget}
can be weaker than estimating its terms separately.
\end{theorem}
\begin{proof}
For a genuine eigenvector $Bf=\mu f$, first-slot linearity gives
\[
 \ip{\mathcal L_n f}{f}=2\Re\mu\,\ip{G_nf}{f}.
\]
Consequently $2|\Re\mu|c_n\le\|\mathcal L_n\|$. Every zero chain
contains its first, genuine eigenvector, so multiplicity creates no gap
in this argument. To obtain the budget, substitute
$J_nB=iD_nJ_n-Y_n$ into both terms of $G_nB+B^*G_n$ and use
$D_n=D_n^*$ and $S_n=S_n^*$. This proves the displayed identity with
the indicated sign and adjoints. The norm bound is the triangle inequality.
\end{proof}

The criterion is invariant under multiplication of $S_n$ by a positive
scalar. A vanishing unnormalized defect obtained by collapsing the metric
does not suffice. On the corrected sampler of
\eqref{v1:eq:graphdefect}, its total source is exactly
\begin{equation}\label{eq:v15-total-source}
 Y_n=z^{\mathrm c}_{L,2K}\eta_{\mathrm B}+\mathcal R_{L,K}.
\end{equation}
Both terms must enter \eqref{eq:v15-budget}. Weak G1 concerns particular
Weil residual pairings and does not bound this budget for an arbitrary
positive $S_n$. Applying the criterion to every hypothetical off-line
zero would prove RH, but the required small signed budget remains an
independent arithmetic obligation. The old Riesz theorem is still useful
for spectral projections; its realization assumption has not been proved.

\subsection*{A positive metric that annihilates the source and shell}
There is a concrete zero-list-independent benchmark for this criterion.
Use the physical scaling $t_n=2\pi n/L$, $D=D_L$, the centered output
$H_F=\kappa F$ of \eqref{v1:eq:hardy}, and the corrected sampler
$J=\Jr_{L,K}$ on $E_{2K}$. Fix $T>0$, write $P=P_{T,L}$ for the
orthogonal projection to $|t_n|\le T$, and choose
$K\ge\lfloor TL/(2\pi)\rfloor$ so the endpoint shell lies outside
this band. Put
\[
 z=z^{\mathrm c}_{L,2K},\qquad z_T=Pz,\qquad
 w=\frac{z_T}{\|z_T\|},\qquad
 S_{T,L}=P-|w\rangle\langle w|.
\]
For all sufficiently large $L$, $z_T\ne0$. Here the rank-one map is
$x\mapsto w\ip{x}{w}$, in the fixed first-slot-linear convention.
The construction uses values of $\kappa\zeta$ on a physical band, not a
list of nontrivial zeros. It is an analytic benchmark, not a newly proved
finite-prime Weil metric or positivity transfer.

\begin{proposition}[Source-annihilating metric with fixed-core coercivity]
\label{prop:v15-projector}
The matrix $S_{T,L}$ is an orthogonal projection and satisfies exactly
\begin{equation}\label{eq:v15-annihilation}
 S_{T,L}z=0,\qquad S_{T,L}\mathcal R_{L,K}=0,
 \qquad S_{T,L}J=S_{T,L}\Jc_{L,2K}.
\end{equation}
For every fixed finite root space $\mathcal V$ as above, there is
$c_{T,\mathcal V}>0$ such that, for sufficiently large $L$,
\begin{equation}\label{eq:v15-projector-coercivity}
 \ip{S_{T,L}JF}{JF}\ge c_{T,\mathcal V}\|F\|^2
 \qquad(F\in\mathcal V).
\end{equation}
These assertions are uniform in the ambient cutoff once its shell lies
outside the fixed band. The exact remaining Lyapunov defect is
\begin{equation}\label{eq:v15-commutator-budget}
 G_LB+B^*G_L=-iJ^*[D,S_{T,L}]J,
 \qquad G_L=J^*S_{T,L}J.
\end{equation}
\end{proposition}
\begin{proof}
Since $w\in\operatorname{Ran}P$ is a unit vector, $S_{T,L}$ is the
orthogonal projection onto $\operatorname{Ran}P\cap w^\perp$.
It annihilates $z$. Both $c_{L,K}$ and $D_Lc_{L,K}$ are supported on
$K<|n|\le2K$, so it also annihilates the two summands of
$\mathcal R_{L,K}F$ and the endpoint correction itself. This proves
\eqref{eq:v15-annihilation}, without any estimate on their sizes.

For $U,V\in\mathcal V+\mathbb C\zeta$, set
\[
 [U,V]_T=\frac1{2\pi}\int_{-T}^T
 H_U(\tfrac12+it)\overline{H_V(\tfrac12+it)}\,dt.
\]
The centered factors $(-1)^n$ cancel in ordinary Gram pairings.
Riemann sums with coefficient $L^{-1/2}$ therefore give, uniformly on
the fixed finite-dimensional unit sphere,
\begin{equation}\label{eq:v15-schur-gram}
 \ip{S_{T,L}JF}{JG}\longrightarrow
 [F,G]_T-\frac{[F,\zeta]_T[\zeta,G]_T}{[\zeta,\zeta]_T}.
\end{equation}
The denominator is positive because $H_\zeta$ is not identically zero
on an interval. Equality in the Cauchy--Schwarz inequality on the right
would give $H_F=aH_\zeta$ there. Since $F=\zeta R_F$ for a proper
rational function $R_F$, analytic continuation away from isolated zeros
forces $R_F\equiv a$. Its decay at infinity gives $a=0$ and $F=0$.
Thus the limiting Gram form is positive definite. Finite-dimensional
compactness proves \eqref{eq:v15-projector-coercivity}. Finally use
\eqref{eq:v15-budget} and \eqref{eq:v15-annihilation}.
\end{proof}

This disposes of the original source and shell terms for this particular
metric, not of the whole compatibility problem. The remaining
commutator can be computed exactly.

\begin{proposition}[Variance of the remaining commutator]
\label{prop:v15-variance}
Put $m_{T,L}=\ip{Dw}{w}\in\mathbb R$, $h=(D-m_{T,L})w$, and
$\varsigma_{T,L}=\|h\|$. Then
\begin{equation}\label{eq:v15-variance}
 [D,S_{T,L}]=-|h\rangle\langle w|+|w\rangle\langle h|,
 \qquad \|[D,S_{T,L}]\|=\varsigma_{T,L}.
\end{equation}
At each fixed $T>0$,
\begin{equation}\label{eq:v15-variance-limit}
 \varsigma_{T,L}^2\longrightarrow
 \frac{\int_{-T}^T(t-m_T)^2|H_\zeta(\tfrac12+it)|^2\,dt}
      {\int_{-T}^T|H_\zeta(\tfrac12+it)|^2\,dt}>0,
 \quad
 m_T=\frac{\int_{-T}^T t|H_\zeta(\tfrac12+it)|^2\,dt}
              {\int_{-T}^T |H_\zeta(\tfrac12+it)|^2\,dt}.
\end{equation}
Moreover the self-adjoint compression $\widehat D=S_{T,L}DS_{T,L}$
on $\operatorname{Ran}S_{T,L}$, with $\widehat J=S_{T,L}J$, has
the exact new graph source
\begin{equation}\label{eq:v15-projected-graph}
 i\widehat D\widehat JF-\widehat JBF
       =-i h\,\ip{JF}{w}.
\end{equation}
\end{proposition}
\begin{proof}
$D$ commutes with $P$ and $h\perp w$. Expanding the commutator gives
\eqref{eq:v15-variance}; on the span of $w$ and $h$ its two singular
values are $\|h\|$, and it vanishes on the orthogonal complement.
The Riemann-sum limits for the first two moments give
\eqref{eq:v15-variance-limit}. The variance is strictly positive since
the nonzero continuous density is not supported at a single frequency.
For \eqref{eq:v15-projected-graph}, first use $S_{T,L}Y=0$, then
\[
 i\widehat D\widehat J-\widehat JB
 =-i S_{T,L}D(I-S_{T,L})J.
\]
The part outside $P$ is killed by $S_{T,L}D$, and the part along $w$
gives $S_{T,L}Dw=h$, proving the formula.
\end{proof}

Thus the elementary operator-norm bound for this commutator does not
tend to zero. This does not by itself decide a restricted signed pairing.
For a hypothetical eigenvector $BF=\mu F$, however, the exact restriction is
\[
 -i\ip{[D,S_{T,L}]JF}{JF}
       =2\Re\mu\,\ip{S_{T,L}JF}{JF}.
\]
It has the required nonzero size under an off-line hypothesis; it is not
an independent upper estimate contradicting that hypothesis. Positivity
and source annihilation alone have simply moved the defect into the
commutator, or equivalently into \eqref{eq:v15-projected-graph}.

\subsection*{Why removing successive source derivatives loses coercivity}
Projecting away $Dz,D^2z,\ldots$ is a natural proposed repair of the last
commutator. On a fixed physical band it encounters an explicit obstruction.

\begin{proposition}[Polynomial source annihilation suppresses zero vectors]
\label{prop:v15-krylov}
Fix $T>0$ and a hypothetical off-line zero
$\rho=1/2+\mu$, where $\mu=d+i\gamma$ and $d\ne0$.
Let $Q_{m,L}$ be contractions supported on $\operatorname{Ran}P_{T,L}$
such that
\[
 Q_{m,L}D^jz_T=0\qquad(0\le j\le2m+1).
\]
Set $M^2=d^2+(T+|\gamma|)^2$ and
$r=1-d^2/M^2\in(0,1)$. For the genuine eigenvector
$F_{\rho,1}=\zeta/(s-\rho)$ one has
\begin{equation}\label{eq:v15-krylov-bound}
 \|Q_{m,L}JF_{\rho,1}\|
 \le \frac{r^{m+1}}{|d|}\,\|z_T\|.
\end{equation}
Consequently bounded projection metrics obtained by removing an increasing
number of these source derivatives cannot retain a uniform positive lower
bound on this vector when $m\to\infty$ at fixed $T$.
\end{proposition}
\begin{proof}
On $[-T,T]$ define $\xi(t)=1-(d^2+(t-\gamma)^2)/M^2\in[0,r]$.
The polynomial
\[
 p_m(t)=\frac{-d-i(t-\gamma)}{M^2}
                   \sum_{j=0}^m\xi(t)^j
\]
has degree $2m+1$. A finite geometric sum gives the exact remainder
\[
 \frac1{it-\mu}-p_m(t)
       =\frac{\xi(t)^{m+1}}{it-\mu},\qquad
 \sup_{|t|\le T}\left|\frac1{it-\mu}-p_m(t)\right|
       \le |d|^{-1}r^{m+1}.
\]
The centered source and root samples have the same half-period phase, so
$P_{T,L}JF_{\rho,1}=(iD-\mu)^{-1}z_T$ exactly; the high shell
does not meet the band. Applying $Q_{m,L}$ removes $p_m(D)z_T$.
The diagonal remainder estimate proves \eqref{eq:v15-krylov-bound}.
Finally $\|z_T\|$ stays bounded by its Riemann-sum limit.
\end{proof}

This is a limited obstruction to bounded polynomial-source projection
metrics, not to all positive metrics, changing bands, or signed arithmetic
estimates. Rescaling a collapsing metric must be assessed by the ratio in
\eqref{eq:v15-relative-criterion}; it cannot be judged by an unnormalized
commutator alone. Section 19 therefore now has a domain-free finite-core
criterion and a coercive positive benchmark with exact source removal,
but no independent small Lyapunov defect. G2 and RH remain open.


\subsection*{Why a sharp band projection does not inherit strip transfer}
The ordinary Gram convergence in Proposition~\ref{prop:v15-projector}
does not authorize applying the holomorphic Weil-transfer theorem after
the physical band projection. The following calculation identifies the
new cut terms on exactly that finite object.

\begin{lemma}[Transform of a sharp physical band]\label{lem:v16-band-transform}
Fix $T>0$, set $\Delta=2\pi/L$, $N=\lfloor T/\Delta\rfloor$ and
$T_L=N\Delta$. Let $a_L$ be functions with uniformly bounded $C^2$
norm on $[-T,T]$. In centered physical coordinates define
\[
 f_L(y)=\frac1L\sum_{n=-N}^N a_L(t_n)e^{it_ny},\quad |y|\le L/2,
 \qquad t_n=n\Delta,
\]
and extend by zero. For fixed $v\in\mathbb C$ with $\Re v\ne0$, put
$\mathcal H_L(v)=\int f_L(y)e^{-vy}\,dy$. Then exactly
\begin{equation}\label{eq:v16-exact-cut}
 \mathcal H_L(v)=-\frac{2\sinh(vL/2)}L
       \sum_{n=-N}^N\frac{(-1)^n a_L(t_n)}{it_n-v}.
\end{equation}
Moreover, uniformly when $v$ ranges over a compact subset of
$\{\Re v\ne0\}$,
\begin{equation}\label{eq:v16-cut-leading}
 \mathcal H_L(v)=-\frac{(-1)^N\sinh(vL/2)}L
 \left\{\frac{a_L(T_L)}{iT_L-v}
       +\frac{a_L(-T_L)}{-iT_L-v}+O(L^{-1})\right\}.
\end{equation}
If $a_L=a+O(L^{-1})$ in $C^2$, define
\[
 C_a(v)=\frac{a(T)}{iT-v}+\frac{a(-T)}{-iT-v}.
\]
Whenever $C_a(v)\ne0$,
\begin{equation}\label{eq:v16-cut-growth}
 |\mathcal H_L(v)|\sim\frac{|C_a(v)|}{2L}
                  e^{|\Re v|L/2}.
\end{equation}
This is compatible with a bounded ordinary norm, since
$\|f_L\|_2^2=L^{-1}\sum_{|n|\le N}|a_L(t_n)|^2=O(1)$.
\end{lemma}
\begin{proof}
Integrate each exponential over $[-L/2,L/2]$; both endpoint phases
are $(-1)^n$, which proves \eqref{eq:v16-exact-cut} with its sign
and factor two. Set $q(t)=a_L(t)/(it-v)$ and
$q_j=q(-T_L+j\Delta)$. An exact telescoping identity is
\[
 2\sum_{j=0}^{2N}(-1)^j q_j-q_0-q_{2N}
 =\sum_{k=0}^{N-1}(q_{2k}-2q_{2k+1}+q_{2k+2}).
\]
The right side has absolute value at most
$N\Delta^2\|q''\|_\infty=T_L\Delta\|q''\|_\infty$.
The denominators are uniformly separated from zero on the stated compact
sets of $v$, giving \eqref{eq:v16-cut-leading}. Now
$T_L=T+O(L^{-1})$ and the assumed $C^2$ convergence replace the
braces by $C_a(v)+O(L^{-1})$. Finally
$|\sinh(vL/2)|\sim\tfrac12 e^{|\Re v|L/2}$. Parseval on the
finite interval proves the ordinary norm assertion.
\end{proof}

\begin{proposition}[Source removal does not remove the sharp-cut growth]
\label{prop:v16-projected-cut}
Use the corrected sampler and the projection $S_{T,L}$ of
Proposition~\ref{prop:v15-projector}, with its endpoint shell outside
the fixed band. Suppose $\rho=\tfrac12+\mu$ is a hypothetical off-line
zero, take its genuine first eigenvector $F=\zeta/(s-\rho)$, and choose
$T>0$ such that $H_\zeta(\tfrac12\pm iT)\ne0$. Such bands exist
because the excluded endpoint zeros are isolated. Put
\[
 Z(t)=H_\zeta(\tfrac12+it),\qquad R(t)=\frac1{it-\mu},\qquad
 \alpha_T=\frac{\int_{-T}^T |Z(t)|^2R(t)\,dt}
                 {\int_{-T}^T |Z(t)|^2\,dt}.
\]
The centered physical representative of $S_{T,L}JF$ has the coefficients
of Lemma~\ref{lem:v16-band-transform} with
\[
 a_L(t)=Z(t)\{R(t)-\alpha_L\},\qquad
 \alpha_L=\frac{\sum_{|n|\le N}|Z(t_n)|^2R(t_n)}
                   {\sum_{|n|\le N}|Z(t_n)|^2}
 =\alpha_T+O(L^{-1}).
\]
Let $a(t)=Z(t)\{R(t)-\alpha_T\}$. For every fixed real $b>0$,
at least one of $C_a(b)$ and $C_a(-b)$ is nonzero. Thus at least one
of the two transforms at $s=\tfrac12\pm b$ grows as a nonzero
constant times $e^{bL/2}/L$. In particular these source-projected
vectors have no bounded full-strip transform limit, even on a compact
substrip containing those two points when $0<b<1/2$.
\end{proposition}
\begin{proof}
The projection coefficient is the first-slot-linear ratio
$\ip{PJF}{z_T}/\|z_T\|^2$, giving the displayed $\alpha_L$; the
half-period factors cancel in this ratio. The shell is absent from the
band. Since $\Re\mu\ne0$, $R$ is smooth there. Ordinary Riemann-sum
estimates give the $O(L^{-1})$ rate, including the positive limiting
denominator. Consequently $a_L=a+O(L^{-1})$ in $C^2$.

Write $A=a(T)$ and $B=a(-T)$. The two equations
$C_a(b)=C_a(-b)=0$ form a nonsingular two-by-two system in $A,B$:
its determinant is
\[
 -\frac{4ibT}{(T^2+b^2)^2}\ne0.
\]
They would force $A=B=0$. Since $Z(\pm T)\ne0$, this would imply
$R(T)=\alpha_T=R(-T)$, which is impossible for $T>0$.
Apply \eqref{eq:v16-cut-growth} to a nonzero coefficient.
\end{proof}

The functional-equation partner does not, by itself, cancel this edge
mechanism. The next statement concerns one reflected pair of summands in
the zero-side formula, not the complete Weil form.

\begin{proposition}[Reflected-pair edge contribution]
\label{prop:v17-pair-edge}
In Lemma~\ref{lem:v16-band-transform}, assume
$a_L=a+O(L^{-1})$ in $C^2[-T,T]$. Fix
$v=d+i\gamma$ with $d>0$, put $v^\sharp=-\overline v$, and suppose
the zeros $\tfrac12+v$ and $\tfrac12+v^\sharp$ have common
multiplicity $m$. The contribution of this reflected pair to the
value-level zero formula is
\[
 \mathcal P_{v,L}=m\left\{
 \mathcal H_L(v)\overline{\mathcal H_L(v^\sharp)}
 +\mathcal H_L(v^\sharp)\overline{\mathcal H_L(v)}\right\}.
\]
Writing
\[
 A_{a,T}(v)=C_a(v)\overline{C_a(v^\sharp)},
\]
one has
\begin{equation}\label{eq:v17-pair-edge}
 \boxed{
 \mathcal P_{v,L}
 =-\frac{m e^{dL}}{2L^2}
   \Re\!\left(e^{i\gamma L}A_{a,T}(v)\right)
   +O\!\left(\frac{e^{dL}}{L^3}\right).}
\end{equation}
Thus reflection converts the two growing transforms into a real
oscillatory term; it does not cancel them identically. If
$A_{a,T}(v)\ne0$ and $\gamma\ne0$, this pair contribution has
subsequences of both signs and magnitude comparable to $e^{dL}/L^2$.
For the source-projected root vector of
Proposition~\ref{prop:v16-projected-cut}, the coefficient is the explicit
product formed from $a(t)=Z(t)\{(it-\mu)^{-1}-\alpha_T\}$.
No assertion is made when one of these two edge coefficients vanishes.
\end{proposition}
\begin{proof}
Lemma~\ref{lem:v16-band-transform} gives, with
$\epsilon_L=(-1)^N$,
\[
 \mathcal H_L(v)=-\frac{\epsilon_L\sinh(vL/2)}L
                   \{C_a(v)+O(L^{-1})\}.
\]
Since $v^\sharp=-\overline v$,
\[
 \overline{\mathcal H_L(v^\sharp)}
 =\frac{\epsilon_L\sinh(vL/2)}L
       \{\overline{C_a(v^\sharp)}+O(L^{-1})\}.
\]
Their product therefore equals
\[
 -\frac{\sinh^2(vL/2)}{L^2}
  \{A_{a,T}(v)+O(L^{-1})\}.
\]
The other reflected summand is its complex conjugate. Finally
$\sinh^2(vL/2)=\tfrac14e^{vL}\{1+O(e^{-dL})\}$, which proves
\eqref{eq:v17-pair-edge}. Choosing the phase of $\gamma L$ near
$-\arg A_{a,T}(v)$ and near $\pi-\arg A_{a,T}(v)$ gives the two
subsequences.
\end{proof}

This does not prove divergence or negativity of the complete Weil form:
other zero summands may cancel the displayed contribution. The combined
prime, pole and archimedean expression is an equivalent representation
of that same form, not an additional contribution. Nor does this rule out a new
signed product estimate just on the zero set. It proves that the existing
bounded holomorphic-strip transfer cannot simply be inherited by these
sharply projected vectors. There is no conflict with weak G1, which
estimates a residual against band tests without asserting that those tests
retain the original Hardy output's strip bounds.

For a proposed correction $M=S_{T,L}(W+a_0I)S_{T,L}$, source and shell
annihilation still hold, but all three terms in
\[
 [D,M]=[D,S_{T,L}](W+a_0I)S_{T,L}
       +S_{T,L}[D,W]S_{T,L}
       +S_{T,L}(W+a_0I)[D,S_{T,L}]
\]
remain in the metric budget. Neither positivity of this compression nor
a small signed budget is proved. A valid next step must control its
actual finite prime, pole and archimedean terms together, or introduce
a different filter with independently justified product estimates.

\subsection*{Full-space projection with a proved strip-product bound}
The preceding obstruction concerns a fixed physical band. It can be
avoided without discarding the graph defect. Keep the original ambient
space $E_{2K}$, the physical $D$, the source $z=\Jc_{L,2K}\zeta$,
the shell $c=c_{L,K}$ and the corrected $J=\Jr_{L,K}$. Define
\[
 e_0=c/\|c\|,\qquad e_1=Dc/\|Dc\|,\qquad
 P_{\rm sh}=|e_0\rangle\langle e_0|+|e_1\rangle\langle e_1|,
 \qquad Q_0=I-P_{\rm sh}.
\]
Parity makes $e_0,e_1$ orthonormal. Put
\begin{equation}\label{eq:v18-full-projection}
 u=Q_0z,\qquad
 S_{L,K}=Q_0-\frac{|u\rangle\langle u|}{\|u\|^2},\qquad
 \alpha_{L,K}(F)=\frac{\ip{Q_0\Jc_{L,2K}F}{u}}{\|u\|^2}.
\end{equation}
The nonzero denominator for large $L$ is proved below. This construction
uses analytic samples of $\kappa\zeta$ and the two known shell vectors;
it does not use a zero list. It changes the metric, not $J$, its physical
endpoint, or the finite Weil matrix.

\begin{proposition}[Full-space source removal and Weil transfer]
\label{prop:v18-full-projection}
Let $\mathcal V$ be a fixed finite span of genuine root chains,
on either side of or on the critical line. Choose the fixed Hardy
smoothing so that, on $\mathcal U=\mathcal V+\mathbb C\zeta$,
\[
 |H_U(\tfrac12+it)|\le C\|U\|(1+|t|)^{-R},\quad R\ge3,
 \qquad |g_U^{(j)}(y)|\le C\|U\|e^{-p|y|},\quad j\le2,
\]
with $p>1/2$, as in \eqref{v1:eq:tails}. Let
$K\ge\lceil e^{2L}L^4\rceil$. Then $S=S_{L,K}$ is an orthogonal
projection and, exactly,
\begin{equation}\label{eq:v18-full-annihilation}
 Sz=Sc=SDc=0,\qquad S\mathcal R=0,\qquad
 SJF=Q_0\Jc_{L,2K}(F-\alpha_{L,K}(F)\zeta).
\end{equation}
If $\mathcal H_{L,K,F}$ is the transform of $SJF$ translated back
to $[-L/2,L/2]$ and zero extended, then throughout $0\le\Re s\le1$,
\begin{equation}\label{eq:v18-full-strip}
 \left|\mathcal H_{L,K,F}(s)-H_F(s)
             +\alpha_{L,K}(F)H_\zeta(s)\right|
 \le \frac{C\varepsilon_{L,K}\|F\|}{1+|\Im s|},
\end{equation}
where, uniformly over these cutoffs,
\[
 \varepsilon_{L,K}
 =e^{-(p-1/2)L/2}
       +Le^{L/4}(1+K/L)^{2-R}\longrightarrow0.
\]
Write $[U,V]_\infty=(2\pi)^{-1}\int_{\mathbb R}
H_U(1/2+it)\overline{H_V(1/2+it)}\,dt$. Uniformly on bounded
subsets of $\mathcal V$,
\begin{align}
 \ip{SJF}{SJG}&\longrightarrow
 [F,G]_\infty-
 \frac{[F,\zeta]_\infty[\zeta,G]_\infty}{[\zeta,\zeta]_\infty},
       \label{eq:v18-full-gram}\\
 \ip{WSJF}{SJG}&\longrightarrow\Q(g_F,g_G).
       \label{eq:v18-full-weil}
\end{align}
The first limiting form is positive definite on $\mathcal V$.
The second limit has error $O(\varepsilon_{L,K}+\varepsilon_{L,K}^2)$.
\end{proposition}
\begin{proof}
Let $X=1+K/L$. A centered sample vector $a=\Jc_{L,2K}U$ obeys
\[
 \|1_{K<|n|\le2K}a\|_2\le C\|U\|X^{1/2-R},\qquad
 \sum_n(1+|t_n|)|(P_{\rm sh}a)_n|
       \le C\sqrt L\|U\|X^{2-R}.
\]
For the second inequality use Cauchy--Schwarz on the $2K$ shell
coordinates, $\|P_{\rm sh}a\|\le\|1_{\rm sh}a\|$ and
$1+|t_n|\le CX$ there. The omitted coefficients of the infinite
centered Fourier series satisfy the same weighted sum bound, since
$R>2$. Riemann sums on the line and these tail bounds give
\[
 \|u\|^2\longrightarrow[\zeta,\zeta]_\infty>0,\qquad
 \alpha_{L,K}(F)\longrightarrow
       [F,\zeta]_\infty/[\zeta,\zeta]_\infty.
\]
In particular $\alpha_{L,K}$ is uniformly bounded on the fixed space.
The projection and annihilation identities now follow directly from
\eqref{eq:v18-full-projection}; $\mathcal R F$ lies in
$\operatorname{span}\{c,Dc\}$ exactly. Since $Sc=0$, the endpoint
correction in $J$ is annihilated, which proves the last identity in
\eqref{eq:v18-full-annihilation}.

Here is the strip estimate, including the physical endpoint jumps.
Set $U=F-\alpha_{L,K}(F)\zeta$. Compare the centered representative
of $Q_0\Jc_{L,2K}U$ with $1_{[-L/2,L/2]}P_Lg_U$.
Their difference $b$ consists of the omitted Fourier tail and the
negative shell projection. The weighted coefficient bounds above,
with the physical basis factor $L^{-1/2}$, give on this interval
\[
 \|b\|_\infty+\|b'\|_\infty
       \le C\|F\|X^{2-R}.
\]
For $|d|\le1/2$, zero extension therefore has
\[
 \|e^{-d\,\cdot}b\|_1+
 \operatorname{Var}(e^{-d\,\cdot}b)
       \le CL e^{L/4}\|F\|X^{2-R}.
\]
The variation includes both cut atoms. The $L^1$ estimate at bounded
ordinates and one integration by parts against this finite derivative
measure at large ordinates give the second term of
\eqref{eq:v18-full-strip}. Lemma~\ref{v1:lem:strip} supplies the
first term. At $R=3$ the second term is at most
$C e^{-7L/4}L^{-2}$ for the stated cutoffs. No fixed-band edge
asymptotic is used.

Ordinary Gram convergence and the projection formula prove
\eqref{eq:v18-full-gram}. Equality in its Cauchy--Schwarz bound
would give $H_F=aH_\zeta$ almost everywhere on the line.
Analyticity then forces the proper rational function $F/\zeta$
to be constant, hence zero. This proves positive definiteness and
thus a uniform lower bound on the fixed finite-dimensional unit sphere.

For the Weil limit, $H_\zeta(\rho)=0$ at every nontrivial zero;
the pole at $1$ has already been canceled by $\kappa$.
Thus subtracting $\alpha_{L,K}(F)H_\zeta$ changes none of the
global zero-form values. Equation~\eqref{eq:v18-full-strip} and
the bounded holomorphic transforms bound each product error by
\[
 C\|F\|\|G\|
 (\varepsilon_{L,K}+\varepsilon_{L,K}^2)
                       (1+|\Im\rho|)^{-2}.
\]
This is summable over all zeros, without RH. Apply
\eqref{v1:eq:zeros} and translate back to $[0,L]$; reflected
translation factors cancel in each product. This proves
\eqref{eq:v18-full-weil} on the actual matrix $W_{L,2K}$.
\end{proof}

The projection above resolves a transfer issue for a different metric;
it does not repair the fixed-band vectors of
Propositions~\ref{prop:v16-projected-cut}--\ref{prop:v17-pair-edge}.
Nor does it make the original graph defect small: it annihilates that
defect inside a particular metric. Its remaining signed budget is exactly
\begin{equation}\label{eq:v18-full-budget}
 G_LB+B^*G_L=-iJ^*[D,S_{L,K}]J,\qquad G_L=J^*S_{L,K}J.
\end{equation}
In particular $Q_0$ must not be commuted through $D$; removing the
two shell directions does not make their orthogonal complement invariant.

\begin{corollary}[A fixed scalar shift does not close the repaired metric]
\label{cor:v18-shift}
Let $\mathcal V$ be a nonempty hypothetical one-sided off-line root
space and fix $a_0>0$. With $S=S_{L,K}$ as above, put
\[
 M_L=S(W+a_0I)S,\quad A_L=J^*SWSJ,\quad
 G_L=J^*SJ,\quad H_L=J^*M_LJ.
\]
Then $A_L\to0$, $G_L\to G_\infty\succ0$ and
$H_L\to a_0G_\infty\succ0$, but
\begin{equation}\label{eq:v18-shift-budget}
 H_LB+B^*H_L\longrightarrow
                  a_0(G_\infty B+B^*G_\infty)\ne0.
\end{equation}
Consequently this candidate does not satisfy the relative defect
criterion of Theorem~\ref{thm:v15-finite-metric}.
\end{corollary}
\begin{proof}
The global Weil form vanishes on a one-sided root space, so
\eqref{eq:v18-full-weil} gives $A_L\to0$ in its finite matrix norm.
Since $B$ is fixed, $A_LB+B^*A_L\to0$ as well. The identity
$H_L=A_L+a_0G_L$ proves the limits. On a first root vector $f$,
$Bf=\mu f$ with $\Re\mu\ne0$, the limiting signed pairing is
$2a_0\Re\mu\,\ip{G_\infty f}{f}\ne0$. Its ratio to the
limiting metric pairing is $2\Re\mu$. The pullback metrics are
bounded above as well as below, so their positive lower bounds cannot
turn this nonzero limit into a vanishing relative defect.
\end{proof}

This is restricted coercivity, not positivity of $M_L$ on the entire
ambient space. The exact budget for $M_L$, whose source terms also
vanish, still contains all three terms
\[
 [D,M_L]=[D,S](W+a_0I)S+S[D,W]S+S(W+a_0I)[D,S].
\]
The new transfer theorem controls their combined pulled-back value via
\eqref{eq:v18-shift-budget}; it supplies no independent small signed
estimate. A successful G2 metric must do more than preserve the
one-sided zero Weil limit and add a fixed ordinary positive shift.

\subsection*{Norm-level Weil action and the squared-metric candidate}
Squaring the finite Weil matrix supplies positivity without a scalar
shift. Its behavior cannot be inferred from form convergence alone.
We first prove the needed norm estimate. The classical unconditional
zero count gives, with multiplicities,
\begin{equation}\label{eq:v19-local-count}
 \#\{\rho:j\le\Im\rho<j+1\}\le C\log(2+|j|),\qquad j\in\mathbb Z.
\end{equation}
For example, this follows by subtracting the estimates for $N(T)$ at
the two endpoints in Trudgian~\cite[Corollary 1]{Trudgian2012}; finitely
many small ordinates are absorbed into $C$. Neither simplicity nor RH
is used.

\begin{lemma}[Local synthesis from decaying zero coefficients]
\label{lem:v19-synthesis}
Suppose $|a_\rho|\le\epsilon/(1+|\Im\rho|)$ on the multiset of
nontrivial zeros. Then, for $A\ge1$, the series
\[
 r(y)=\sum_{\rho\in\Z}a_\rho e^{(\rho-1/2)y}
\]
converges in $L^2(-A,A)$, and
\begin{equation}\label{eq:v19-synthesis}
 \|r\|_{L^2(-A,A)}\le C\epsilon\sqrt{A+1}\,e^{A/2}.
\end{equation}
The constant is independent of $A$, $\epsilon$ and the coefficients.
\end{lemma}
\begin{proof}
Choose a nonnegative smooth function $\chi_A$ equal to one on
$[-A,A]$, zero outside $[-A-1,A+1]$, with its first two derivatives
bounded independently of $A$. Write $\rho=1/2+d_\rho+i\gamma_\rho$.
Since $|d_\rho|<1/2$, two integrations by parts, together with the
direct integral bound, give
\[
 \left|\int\chi_A(y)e^{(d_\rho+d_\tau)y}
                  e^{i(\gamma_\rho-\gamma_\tau)y}\,dy\right|
 \le\frac{C(A+1)e^A}{1+|\gamma_\rho-\gamma_\tau|^2}.
\]
For a finite partial sum put
$b_j=\sum_{j\le\gamma_\rho<j+1}|a_\rho|$. Grouping both
indices by these unit intervals bounds its squared norm by
\[
 C(A+1)e^A\sum_{j,k}\frac{b_jb_k}{1+|j-k|^2}
 \le C(A+1)e^A\sum_j b_j^2.
\]
The last inequality is the $\ell^1*\ell^2$ convolution bound.
Equation~\eqref{eq:v19-local-count} gives
$b_j\le C\epsilon\log(2+|j|)/(1+|j|)$, whose square is summable.
The same bound on omitted indices makes the partial sums Cauchy in
$L^2(-A,A)$ and proves \eqref{eq:v19-synthesis}. It also permits
arbitrary finite exhaustions of the zero multiset.
\end{proof}

Let $\mathcal T_L$ denote the unitary translation from the physical
interval $[0,L]$ to $[-L/2,L/2]$; on coefficient vectors it represents
$x$ as $L^{-1/2}\sum x_n(-1)^ne^{it_ny}$. Let $\Pi_{L,2K}$
be the ordinary orthogonal projection onto the corresponding centered
Fourier space. These operations do not change the finite matrix $W$.

\begin{proposition}[Weil-action transfer on a fixed root family]
\label{prop:v19-action}
Use the objects and cutoffs of Proposition~\ref{prop:v18-full-projection},
but choose the fixed smoothing with tail rate $p>1$. Define the finite sum
\[
 q_F(y)=\sum_{\rho\in\Z}H_F(\rho)e^{(\rho-1/2)y},
\]
where zeros are counted with multiplicity; only the finitely many
selected root locations can contribute. Then uniformly on the fixed
root family,
\begin{equation}\label{eq:v19-action}
 \|\mathcal T_LWSJF-\Pi_{L,2K}q_F\|_2
 \le C\Delta_{L,K}\|F\|,\qquad
 \Delta_{L,K}=\sqrt{L+1}\,e^{L/4}\varepsilon_{L,K}\longrightarrow0.
\end{equation}
In particular the first part of $\Delta_{L,K}$ is
$\sqrt{L+1}\,e^{-(p-1)L/2}$. The original condition $p>1/2$
alone does not make this bound tend to zero.
\end{proposition}
\begin{proof}
Let $x=SJF$ and let $\mathcal H_x$ be the transform of
$\mathcal T_Lx$. At a zero, Proposition~\ref{prop:v18-full-projection}
and $H_\zeta(\rho)=0$ give
\[
 |\mathcal H_x(\rho)-H_F(\rho)|
       \le C\varepsilon_{L,K}\|F\|/(1+|\Im\rho|).
\]
Lemma~\ref{lem:v19-synthesis} synthesizes these errors into a function
$r_F$ of norm at most the right side of \eqref{eq:v19-action}.
For each centered finite Fourier test $v$, the first-slot-linear
zero formula gives exactly
\[
 \Q(\mathcal T_Lx,v)=\sum_\rho\mathcal H_x(\rho)
          \overline{H_v(\rho^\sharp)}
       =\ip{q_F+r_F}{v}.
\]
Indeed $\overline{\rho^\sharp}-1/2=-(\rho-1/2)$, so the
exponential in the first slot is $e^{(\rho-1/2)y}$ with a positive
sign. The product series is absolutely convergent for these tests;
the synthesized series also converges in $L^2$. Translation invariance
of the paired form identifies the left side with
$\ip{\mathcal T_LWx}{v}$. Thus
$\mathcal T_LWx=\Pi_{L,2K}(q_F+r_F)$, and projection contractivity
proves the bound. At $R=3$ the cutoff contribution to $\Delta_{L,K}$
is $O(e^{-3L/2}L^{-3/2})$. The other contribution tends to zero
because $p>1$.
\end{proof}

The stronger tail margin is available, for example with fixed
$\sigma=2$, $a=4$ and sufficient $M$, taking $1<p\le3/2$ in
\eqref{v1:eq:tails}. Critical-line root chains are allowed here and
in Proposition~\ref{prop:v18-full-projection}: the zeros remove their
apparent transform poles, and the proper-rational independence proof
works away from the isolated zero locations. This extension assumes
genuine root functions, not inserted meromorphic poles.

\begin{corollary}[Mass of the squared-Weil metric]
\label{cor:v19-square}
On the same finite space set $M_L^{\square}=S W^2 S\succeq0$.
For a zero $\rho=1/2+\mu$ of multiplicity $m$, write
$F_{\rho,j}=\zeta/(s-\rho)^j$, $1\le j\le m$, and
\[
 C_\rho=m\kappa(\rho)\frac{\zeta^{(m)}(\rho)}{m!}\ne0.
\]
If $d=\Re\mu\ne0$, then
\begin{equation}\label{eq:v19-square-top}
 \ip{M_L^{\square}JF_{\rho,m}}{JF_{\rho,m}}
 =\|WSJF_{\rho,m}\|^2
 \sim |C_\rho|^2\frac{\sinh(|d|L)}{|d|}.
\end{equation}
For a critical-line zero the corresponding asymptotic is
$|C_\rho|^2L$. For every lower jet $j<m$,
\begin{equation}\label{eq:v19-square-lower}
 \ip{M_L^{\square}JF_{\rho,j}}{JF_{\rho,j}}
       =O(\Delta_{L,K}^2)\longrightarrow0.
\end{equation}
These assertions concern the displayed normalization; a rescaled metric
must still satisfy the relative criterion, not merely an absolute bound.
\end{corollary}
\begin{proof}
Self-adjointness of $W$ gives positivity and the squared norm identity.
For the top root, $q_F=C_\rho e^{\mu y}$; for lower jets $q_F=0$,
because all their zero values vanish. Moreover
\[
 \|e^{\mu y}\|_{L^2(-L/2,L/2)}^2
 =\begin{cases}\sinh(|d|L)/|d|,&d\ne0,\\L,&d=0.\end{cases}
\]
Its centered Fourier coefficient at index $n$ is
$2(-1)^n\sinh(\mu L/2)/(\sqrt L(\mu-it_n))$, with the removable
value used when $\mu=it_n$. For $|n|>2K$ this gives
\[
 \|(I-\Pi_{L,2K})e^{\mu\,\cdot}\|_2
       \le C_\mu e^{|d|L/2}\sqrt{L/K}\longrightarrow0.
\]
The cutoff dominates this exponential since $|d|<1/2$.
Apply Proposition~\ref{prop:v19-action} to obtain both asymptotics
and the lower-jet bound.
\end{proof}

This candidate is globally positive and annihilates the exact total
source, since $S(z\eta_{\mathrm B}+\mathcal R)=0$. Nevertheless
its full commutator is
\[
 [D,M_L^{\square}]=[D,S]W^2S
     +S\bigl([D,W]W+W[D,W]\bigr)S+SW^2[D,S].
\]
All terms belong to the signed G2 budget. At a hypothetical simple
off-line root $Bf=\mu f$, putting $G_L^{\square}=J^*M_L^{\square}J$
gives the exact quotient
\[
 \frac{\ip{(G_L^{\square}B+B^*G_L^{\square})f}{f}}
      {\ip{G_L^{\square}f}{f}}=2\Re\mu
\]
whenever its denominator is nonzero. At a multiple zero, the lower
jets additionally lose an unscaled coercive lower bound by
\eqref{eq:v19-square-lower}. Squaring has not supplied a small
relative defect. The norm-action theorem is a stronger transfer result,
not a proof of the independent signed arithmetic estimate or RH.

\section{Finite Weil residuals: compact-source radicality and G1 transfer}\label{sec:g1}
The finite-prime comparison can be narrowed further.  The following subsection isolates the part of the Connes--Consani--Moscovici comparison that is controlled by classical prolate asymptotics and the prime number theorem, separately from the still-open spectral-ordering problem.

\subsection*{A rank-two symmetry defect for an arbitrary even candidate}
Let $W_{\lambda,N}$ denote the finite Weil matrix, let $D_N$ be the Fourier-index matrix, and identify the equal-coefficient boundary covector with its Riesz vector $\eta_N$, writing $\eta_N(x)=\langle x,\eta_N\rangle$ in the first-slot-linear convention.  The exact CCM commutator has the form
\begin{equation}\label{eq:ccm-commutator}
 D_NW_{\lambda,N}-W_{\lambda,N}D_N
 =|\beta_N\rangle\langle\eta_N|-|\eta_N\rangle\langle\beta_N|,
\end{equation}
where $\beta_N$ is odd and $\eta_N$ is even.  Let $k_N$ be any real even vector satisfying $\eta_N(k_N)=1$, and define
\[
 \mathcal D_{k,N}=D_N-|D_Nk_N\rangle\langle\eta_N|.
\]
Then parity gives $\langle k_N,\beta_N\rangle=0$, and applying \eqref{eq:ccm-commutator} to $k_N$ gives
\[
 W_{\lambda,N}D_Nk_N=D_NW_{\lambda,N}k_N-\beta_N.
\]
Substitution yields the exact identity
\begin{equation}\label{eq:ranktwo-weil-defect}
 \boxed{
 W_{\lambda,N}\mathcal D_{k,N}
 -\mathcal D_{k,N}^{*}W_{\lambda,N}
 =-|D_NW_{\lambda,N}k_N\rangle\langle\eta_N|
  +|\eta_N\rangle\langle D_NW_{\lambda,N}k_N|.}
\end{equation}
Thus no ground-state vector is needed to identify the metric-symmetry defect: for an explicit even candidate it is rank at most two and is controlled by one linear residual.

It is useful to state the scale invariant version with the physical logarithmic generator $D_{\log}=(2\pi/L)D_N$, $L=2\log\lambda$, and the actual boundary-evaluation covector $\delta_N=L^{-1/2}\eta_N$.  On the fixed physical test band $|D_{\log}|\le T$, the restriction of the boundary-evaluation covector $\delta_N$ has norm bounded by a constant depending only on $T$, independently of how large the ambient section $E_N$ is.  Thus one may later choose a much larger diagonal section $N=N(\lambda)$ to resolve the prolate boundary layer without changing the fixed-window estimate.  Hence the relevant projective quantity is
\begin{equation}\label{eq:physical-projective-residual}
 \frac{\|P_TD_{\log}W_{\lambda,N}k_N\|}
 {|\delta_N(k_N)|},
\end{equation}
where $P_T$ is projection to the physical-frequency window $[-T,T]$.  This formulation avoids the factor-of-$\sqrt L$ ambiguity that results from simultaneously rescaling the covector and the candidate.

\subsection*{The compact-source radical passage}
The classical radical theorem applies to the even Schwartz class with
$h(0)=\widehat h(0)=0$~\cite[Section 3]{ConnesConsani2023}.
The following extension supplies the domain passage needed for the compact
prolate source. It is a statement about explicit-formula pairings and the
closed \emph{semilocal} form, not a claim of a positive global closed form.

\begin{theorem}[Radicality with a nonzero compact-source endpoint]
\label{thm:v14-radical}
Fix $b>0$. Let $h$ be even and smooth on $[-b,b]$, extended by zero
outside that interval, and suppose
\[
 h(0)=0,\qquad \int_0^b h(u)\,du=0.
\]
No vanishing of the interior trace $h(b^-)$ is required. In logarithmic
coordinates put
\[
 k(x)=\mathcal E(h)(e^x)=e^{x/2}\sum_{n\ge1}h(ne^x).
\]
For every $v$ obtained by zero-extending an $H^1$ function on a compact
interval $J$, all terms of the polarized explicit formula $QW(k,v)$ are
finite and $QW(k,v)=0$. In particular, for a compact interval $I$
containing the support of $v$, set $k_{\rm in}=1_I k$ and
$r=k-k_{\rm in}$. Then $k_{\rm in}$ and $v$ belong to the closed
semilocal form domain and
\[
 QW_I(k_{\rm in},v)=-QW(r,v).
\]
The assertion holds at each fixed $b$; its approximation constants need
not be uniform as $b$ varies.
\end{theorem}
\begin{proof}
\emph{A primitive gains one power.}
Write $\beta(t)=\tfrac12-\{t\}$ away from integers. Euler summation,
with $N=\lfloor b/u\rfloor$, gives for almost every sufficiently small $u$
\[
 \sum_{n=1}^N h(nu)
 =u^{-1}\int_0^{Nu}h(t)\,dt+
       \frac{h(Nu)-h(0)}2+O_h(u)
 =h(b^-)\beta(b/u)+O_h(u).
\]
For example, the composite trapezoid remainder is bounded by
$Cbu\|h''\|_\infty$ in the first equality; Taylor expansion of the
missing interval $[Nu,b]$ gives the second. Values at the arithmetic
cuts affect neither an integral nor the correlations below. Consequently
\begin{equation}\label{eq:v14-euler-tail}
 \mathcal E(h)(u)
 =h(b^-)u^{1/2}\beta(b/u)+O_h(u^{3/2}).
\end{equation}
In particular $k\in L^2(\mathbb R)$, it vanishes for $x>\log b$, and
it is of bounded variation on every compact logarithmic interval.
Let $K(x)=\int_{-\infty}^x k(t)\,dt$. Since $\beta$ has mean zero
and a bounded primitive,
\[
 \int_0^u v^{-1/2}\beta(b/v)\,dv
 =\sqrt b\int_{b/u}^{\infty}\beta(w)w^{-3/2}\,dw
 =O_b(u^{3/2}).
\]
The last bound follows by one integration by parts against that bounded
primitive. Integrating the remainder in \eqref{eq:v14-euler-tail} yields
\begin{equation}\label{eq:v14-primitive}
 K(x)=O_h(e^{3x/2})\qquad(x\longrightarrow-\infty).
\end{equation}
Both moments matter: without $h(0)=0$ an additional term
$-\tfrac12h(0)u^{1/2}$ generally prevents this gain.

\emph{Absolute convergence of the geometric side.}
The zero extension $v_0$ of $v$ has bounded variation, including its two
endpoint jumps. In the first-slot-linear convention its correlation is
\[
 C_{k,v}(y)=\int k(x)\overline{v_0(x-y)}\,dx
          =-\int K(x)\,d\overline{v_0(x-y)}.
\]
Thus \eqref{eq:v14-primitive} gives
$C_{k,v}(y)=O_{h,v}(e^{3y/2})$ as $y\to-\infty$; it is zero for
all sufficiently large positive $y$. Its reflected correlation has the
opposite support orientation. The even combinations therefore decay as
$O_{h,v}(e^{-3y/2})$ for $y\to+\infty$. Each growing pole integral
is absolutely convergent, and so is the prime series, since its tail is
bounded by a constant times $\sum_{n\ge2}(\log n)n^{-2}$.

At the archimedean singularity use, for $0<t\le1$,
\begin{equation}\label{eq:v14-translation}
 \|v_0(\cdot-t)-v_0\|_2
 \le C\left[t\|v'\|_2+
       \sqrt t\{|v(\inf J)|+|v(\sup J)|\}\right].
\end{equation}
This follows from the fundamental theorem of calculus on the overlap
and the two endpoint strips. Cauchy--Schwarz makes the subtracted
archimedean numerator $O(t^{1/2})$, which is integrable against its
$O(t^{-1})$ kernel. Its large-$t$ part is integrable by $L^2$ bounds
and exponential kernel decay.

\emph{Moment-preserving smoothing.}
Choose $\phi_\epsilon\ge0$ in $C_c^\infty((\epsilon,2\epsilon))$
with integral one, and define
\begin{equation}\label{eq:v14-smoothing}
 h_\epsilon(u)=\int\phi_\epsilon(t)e^{t/2}h(e^t u)\,dt,
 \qquad
 k_\epsilon(x)=\int\phi_\epsilon(t)k(x+t)\,dt.
\end{equation}
Multiplicative convolution smooths the endpoint; near zero it preserves
the smooth even germ. Hence $h_\epsilon$ is even and compactly supported
smooth, with support in $[-be^{-\epsilon},be^{-\epsilon}]$.
Both moments remain exactly zero: dilation multiplies the integral by
$e^{-t/2}$ and the central value is identically zero. Direct substitution
in the finite source sum proves
$\mathcal E(h_\epsilon)(e^x)=k_\epsilon(x)$, with all scale factors
as in \eqref{eq:v14-smoothing}. This is the scaling covariance used in
\cite[Lemma 6.2]{ConnesConsani2023}. Translation continuity gives
$k_\epsilon\to k$ in $L^2(\mathbb R)$, while
\[
 K_\epsilon(x)=\int\phi_\epsilon(t)K(x+t)\,dt
              =O_h(e^{3x/2})
\]
uniformly for sufficiently small $\epsilon$. In particular none of the
estimates needed for the passage uses derivatives of $h_\epsilon$ that
might blow up near its moving endpoint.

For a smooth compact logarithmic test $w$, the classical Schwartz
radical theorem gives $QW(k_\epsilon,w)=0$. To allow the stated $v$,
convolve its zero extension with a nonnegative smooth additive
approximate identity. These tests converge in $L^2$, have uniformly
bounded variation and support, and satisfy the same translation bound
\eqref{eq:v14-translation} by convolution contraction. The preceding
correlation bounds therefore dominate separately the prime series, both
pole integrals and the subtracted archimedean integral. Dominated
convergence gives $QW(k_\epsilon,v)=0$. The same argument, now using
the uniform primitive bound for $K_\epsilon$ and $L^2$ convergence,
allows $\epsilon\downarrow0$ and proves $QW(k,v)=0$.

\emph{The semilocal domain and truncation.}
On any fixed compact interval $k$ is a finite sum of smooth pieces with
finitely many jumps, so $k_{\rm in}$ is compactly supported BV. Such a
function has Fourier transform $O((1+|\xi|)^{-1})$, and hence
\[
 \int_{\mathbb R}(1+\log(2+|\xi|))
               |\widehat{k_{\rm in}}(\xi)|^2\,d\xi<\infty.
\]
The same holds for $v_0$. On a fixed compact support the archimedean
multiplier has logarithmic growth; the finitely many prime translations
and pole terms are bounded forms. This is the closed semilocal form
domain of the explicit formula~\cite{ConnesConsani2023}.
The smooth approximations converge there as well: on a fixed enlarged
interval they have uniform BV bounds and $L^2$ convergence, so a Fourier
tail split with $\log(2+|\xi|)/(1+|\xi|)^2$ proves convergence in the
logarithmic form norm. Compact truncation leaves the far lower primitive
unchanged, so $QW(r,v)$ is also finite. Linearity now gives the asserted
truncation identity, with the semilocal term equal to its closed-form
realization.
\end{proof}

The endpoint of the approximating sources need not equal that of $h$.
The theorem is proved at fixed $b$ \emph{before} normalizing by the
original source's separately evaluated one-sided endpoint. No endpoint
trace is inferred from form-norm convergence. This corrects the stronger,
unnecessary endpoint-preserving approximation requirement in version 1.3.

\subsection*{The focused PSWF cross-tail estimates}
The fixed-order radial estimates, the two-sided exterior-energy argument,
and Fuchs's concentration asymptotic below establish the source scale.
Together with Theorem~\ref{thm:v14-radical}, they close the fixed-band
and Sobolev-pairing G1 statements of this section for the explicit repaired
source. They do not give a uniform operator-norm bound on a growing deep
prolate block, spectral ordering, or the G2 graph/form comparison.


CCM's proposed source is the linear combination of the two fixed-order prolate modes $h_{0,\lambda},h_{4,\lambda}$ with vanishing integral.  The order-zero radial normalization can be matched exactly to the endpoint amplitude.  In Dunster's notation, for $m=0$ his matching constant satisfies
\[
 c_n^0(\gamma)=K_n^0(\gamma),
\]
where $K_n^0(\gamma)=Ps_n^0(1,\gamma^2)$.  His Bessel approximation is uniform for the full radial interval $1<x<\infty$.

\begin{proposition}[Endpoint-normalized radial bound]\label{prop:endpoint-radial}
Fix $n\in\{0,4\}$.  For all sufficiently large $\gamma$ there is a constant $C_n$, independent of $\gamma$, such that
\[
 |Ps_n^0(x,\gamma^2)|\le C_n\frac{|Ps_n^0(1,\gamma^2)|}{x},
 \qquad x\ge1.
\]
\end{proposition}
\begin{proof}
Dunster's formula (3.5) gives, uniformly for $1<x<\infty$,
\[
 Ps_n^0(x,\gamma^2)
 =c_n^0(\gamma)
 \left\{\frac{\eta}{(x^2-1)(x^2-\sigma^2)}\right\}^{1/4}
 \left[J_0(\gamma\eta^{1/2})
 +O(\gamma^{-1})\operatorname{env}J_0(\gamma\eta^{1/2})\right],
\]
with $\eta=\xi^2$.  His local relation
\[
 \eta=2(1-\sigma^2)(x-1)+O((x-1)^2)
\]
shows that the geometric prefactor tends to one as $x\downarrow1$; since $J_0(0)=1$ and $c_n^0(\gamma)=Ps_n^0(1,\gamma^2)$, the normalization is exactly the endpoint scale up to the uniform $O(\gamma^{-1})$ error.  For fixed $n$, Dunster also gives $\sigma^2=O_n(\gamma^{-1})$.

Use the standard envelope bound
\[
 \operatorname{env}J_0(y)\ll \min(1,y^{-1/2}).
\]
If $\gamma\xi\le1$, then $x-1=O_n(\gamma^{-2})$, so the geometric prefactor is uniformly bounded and $x^{-1}\asymp1$.  If $\gamma\xi\ge1$, the Bessel-envelope bound gives
\[
 \left\{\frac{\xi^2}{(x^2-1)(x^2-\sigma^2)}\right\}^{1/4}
 \operatorname{env}J_0(\gamma\xi)
 \ll_n
 \frac{\gamma^{-1/2}}{\{(x^2-1)(x^2-\sigma^2)\}^{1/4}}.
\]
On the boundary layer where $x-1=O(1)$ this is $O_n(1)$, while for $x$ bounded away from $1$ it is $O_n(x^{-1})$.  Enlarging the constant across the compact transition interval proves the claim.
\end{proof}

The preceding endpoint estimate can be sharpened away from the radial turning point, and the resulting gain is the key to summing the co-Poisson lattice.

\begin{lemma}[Away-boundary and exterior-energy bounds]\label{lem:pswf-away-energy}
Fix $n\in\{0,4\}$ and write
\[
 \Psi_{n,c}(x)=Ps_n^0(x,c^2),\qquad B_{n,c}=|\Psi_{n,c}(1)|.
\]
For all sufficiently large $c$,
\begin{align}
 |\Psi_{n,c}(x)|&\ll_n \frac{B_{n,c}}{\sqrt c\,x},\qquad x\ge2,\label{eq:pswf-away}\\
 \int_1^\infty |\Psi_{n,c}(x)|^2\,dx&\ll_n \frac{B_{n,c}^2}{c}.\label{eq:pswf-exterior-energy}
\end{align}
\end{lemma}
\begin{proof}
In Dunster's uniform formula (3.5), for $x\ge2$ one has $\xi(x)\asymp_n x$ and
$\{(x^2-1)(x^2-\sigma^2)\}^{1/4}\asymp_n x$.  The Bessel-envelope estimate therefore gives
\[
 \left\{\frac{\xi^2}{(x^2-1)(x^2-\sigma^2)}\right\}^{1/4}
 \operatorname{env}J_0(c\xi)
 \ll_n \frac1{\sqrt c\,x},
\]
which proves \eqref{eq:pswf-away}.

For $1<x<2$, use $d\xi/dx=((x^2-\sigma^2)/(x^2-1))^{1/2}$.  Squaring Dunster's formula and changing variables from $x$ to $\xi$ reduces the main term, up to bounded fixed-$n$ factors, to
\[
 B_{n,c}^2\int_0^{\xi(2)}
 \xi\,\operatorname{env}J_0(c\xi)^2\,d\xi.
\]
The bound $\operatorname{env}J_0(t)^2\ll\min(1,t^{-1})$ makes this $O_n(B_{n,c}^2/c)$; the uniform $O(c^{-1})$ remainder in Dunster's formula is smaller.  On $[2,\infty)$, \eqref{eq:pswf-away} is square integrable and contributes another $O_n(B_{n,c}^2/c)$.  This proves \eqref{eq:pswf-exterior-energy}.
\end{proof}

\subsection*{Normalized source coefficients and endpoint noncancellation}
We now settle the coefficient and repair-size parts of the former source
assumption. These estimates concern the actual two-mode source, not a
choice of independently rescaled endpoint values. The form-domain issue
is treated separately from the elementary correlation consequences below.

Let $\psi_{n,c}$ be the real angular PSWF with
$\|\psi_{n,c}\|_{L^2(-1,1)}=1$, extended to the real line by its finite
Fourier eigenrelation. For $n=0,4$ choose its sign so that
$\psi_{n,c}(0)>0$. Write $\vartheta_n(c)=\chi_n(c)^2$ for its
concentration eigenvalue, and put $b_n(c)=|\psi_{n,c}(1)|$.
These $\psi_{n,c}$ are scalar multiples of the preceding Dunster
normalization; its homogeneous radial estimates are unchanged.

\begin{lemma}[Endpoint size from exterior energy]
\label{lem:v12-endpoint-energy}
For each $n\in\{0,4\}$, as $c\to\infty$,
\begin{equation}\label{eq:v12-endpoint-energy}
 b_n(c)^2\asymp_n c\{1-\vartheta_n(c)\},\qquad
 b_n(c)\asymp_n c^{n/2+3/4}e^{-c}.
\end{equation}
In particular $b_4(c)/b_0(c)\asymp c^2$.
\end{lemma}
\begin{proof}
For clarity, the eigenrelation on $[-1,1]$ is
\[
 \int_{-1}^{1}e^{-icxy}\psi_{n,c}(y)\,dy
       =\mu_n(c)\psi_{n,c}(x),\qquad
 |\mu_n(c)|^2=\frac{2\pi}{c}\vartheta_n(c).
\]
Using the same formula to extend the function to all real $x$, Plancherel
gives the exact identity
\begin{equation}\label{eq:v12-exterior-identity}
 \int_1^\infty |\psi_{n,c}(x)|^2\,dx
     =\frac{1-\vartheta_n(c)}{2\vartheta_n(c)}.
\end{equation}
The upper bound for this integral by $C_nb_n(c)^2/c$ was proved in
Lemma~\ref{lem:pswf-away-energy}. We also need a lower bound; an upper
bound alone would not justify an endpoint comparison.

On the fixed interval $2\le x\le3$, Dunster's radial formula
\cite[eqs.~(3.5)--(3.7)]{Dunster2017} and the large-argument expansion
of $J_0$ give, with a real choice of the endpoint sign,
\[
 \psi_{n,c}(x)=\psi_{n,c}(1)\sqrt{\frac{2}{\pi c}}
 w_c(x)\left\{\cos(c\xi_c(x)-\pi/4)+O_n(c^{-1})\right\},
\]
where $w_c(x)=\{(x^2-1)(x^2-\sigma_n(c)^2)\}^{-1/4}$ and
$\xi_c'(x)=\{(x^2-\sigma_n(c)^2)/(x^2-1)\}^{1/2}$.
Here $\sigma_n(c)^2=O_n(c^{-1})$. On this interval $w_c$ and
$\xi_c'$ have uniform positive lower bounds and bounded derivatives.
One integration by parts therefore yields
\[
 \int_2^3 w_c(x)^2\cos^2(c\xi_c(x)-\pi/4)\,dx
      =\frac12\int_2^3w_c(x)^2\,dx+O_n(c^{-1})\gg_n1.
\]
The integrated cross and squared remainder terms are smaller by an
additional factor $c^{-1}$. Hence the exterior energy is also at least
$c_n b_n(c)^2/c$ for some $c_n>0$. Equation
\eqref{eq:v12-exterior-identity} and $\vartheta_n(c)\to1$ prove the
first comparison in \eqref{eq:v12-endpoint-energy}.

Fuchs's fixed-index asymptotic \cite[Theorem~1]{Fuchs1964}, in this
concentration normalization, is
\begin{equation}\label{eq:v12-fuchs}
 1-\vartheta_n(c)\sim
 \frac{4\sqrt\pi\,8^n}{n!}\,c^{n+1/2}e^{-2c}.
\end{equation}
Combining it with the two-sided energy comparison proves the second
comparison and the ratio assertion. No differentiation of an asymptotic
remainder is used.
\end{proof}

\begin{proposition}[An explicit normalized two-mode source]
\label{prop:v12-source}
Set $c=2\pi\lambda^2$ and
$h_{n,\lambda}(u)=\lambda^{-1/2}\psi_{n,c}(u/\lambda)$ on
$[-\lambda,\lambda]$, zero-extended for Fourier transformation.
Endpoint values denote interior traces. Define
\begin{align}
 a_{4,\lambda}&=\chi_0(c)h_{0,\lambda}(0),&
 a_{0,\lambda}&=-\chi_4(c)h_{4,\lambda}(0),\label{eq:v12-coefficients}\\
 h_\lambda^\circ&=a_{0,\lambda}h_{0,\lambda}
                 +a_{4,\lambda}h_{4,\lambda},&
 B_\lambda^+&=\sqrt\lambda\,h_\lambda^\circ(\lambda).
 \notag
\end{align}
The source has exactly zero integral. Both coefficients have finite
nonzero limits. For all sufficiently large $\lambda$,
\begin{align}
 |B_\lambda^+|&\asymp c^{11/4}e^{-c},\label{eq:v12-combined-endpoint}\\
 \frac{\sum_{n=0,4}|a_{n,\lambda}|\sqrt\lambda
                 |h_{n,\lambda}(\lambda)|}{|B_\lambda^+|}
     &=1+O(c^{-2}),\label{eq:v12-noncancellation}\\
 \frac{|h_\lambda^\circ(0)|}{|B_\lambda^+|}
     &=O(c^{7/4}e^{-c})=O_A(\lambda^{-A})
       \quad\text{for every fixed }A.\label{eq:v12-zero-value}
\end{align}
The last two conclusions are invariant under a common nonzero scalar
rescaling of the source.
\end{proposition}
\begin{proof}
The compressed physical Fourier eigenvalue is $\chi_n(c)>0$ for
$n=0,4$, so evaluation at zero gives exactly
\[
 \int_{-\lambda}^{\lambda}h_{n,\lambda}(u)\,du
       =\chi_n(c)h_{n,\lambda}(0).
\]
Thus \eqref{eq:v12-coefficients} cancels the integral. The fixed-order
Hermite limit, with the unit-norm scaling just specified, gives
\[
 h_{0,\lambda}(0)\longrightarrow2^{1/4},\qquad
 h_{4,\lambda}(0)\longrightarrow
       \frac{\sqrt3}{2\,2^{1/4}}.
\]
This central-value normalization follows, for example, from the
fixed-order limit used in \cite[Lemma~7.2]{ConnesConsaniMoscovici2025};
the uniform error there also makes the physical $L^2$ norm tend to one,
so unit-norm renormalization preserves these limits. Since $\chi_n\to1$,
both coefficients converge to the stated nonzero central constants,
with a minus sign for $a_{0,\lambda}$.

By Lemma~\ref{lem:v12-endpoint-energy}, the magnitude of the order-zero
endpoint contribution divided by that of order four is $O(c^{-2})$.
The triangle and reverse triangle inequalities prove
\eqref{eq:v12-combined-endpoint} and \eqref{eq:v12-noncancellation},
without needing a claim about their relative signs. Finally,
\[
 h_\lambda^\circ(0)
   =\{\chi_0(c)-\chi_4(c)\}
       h_{0,\lambda}(0)h_{4,\lambda}(0).
\]
Equation~\eqref{eq:v12-fuchs} and
$1-\chi_n=(1-\vartheta_n)/(1+\chi_n)$ give
$|h_\lambda^\circ(0)|=O(c^{9/2}e^{-2c})$.
Division by \eqref{eq:v12-combined-endpoint} gives
\eqref{eq:v12-zero-value}.
\end{proof}

\begin{corollary}[Size of the exact value-and-integral repair]
\label{cor:v12-repair}
Fix an even $\psi\in C_c^\infty((-1,1))$ with $\psi(0)=1$ and
$\int\psi=0$. For the source in Proposition~\ref{prop:v12-source},
\[
 \widetilde h_\lambda=h_\lambda^\circ-h_\lambda^\circ(0)\psi
\]
has zero value at zero, zero integral and the same endpoint as
$h_\lambda^\circ$. In lower-tail coordinates $u=\lambda^{-1}e^{-s}$,
for every fixed $A$ and every $S\ge0$ its co-Poisson correction obeys
\begin{equation}\label{eq:v12-repair-tail}
 \left\|1_{s\ge S}\,h_\lambda^\circ(0)
            \mathcal E(\psi)(\lambda^{-1}e^{-s})\right\|_2
     \le C_A |B_\lambda^+|\lambda^{-A}e^{-S/2}.
\end{equation}
This proves the repair's lower-tail size. Its repaired compact source
also meets Theorem~\ref{thm:v14-radical}, which supplies the radical and
semilocal form-domain passage used below.
\end{corollary}
\begin{proof}
The moment and endpoint statements follow directly from the fixed bump.
Poisson summation, retaining its zero-value term, gives
\[
 \mathcal E(\psi)(u)=-\frac12u^{1/2}
           +u^{-1/2}\sum_{m\ge1}\widehat\psi(m/u)=O(u^{1/2})
       \qquad(0<u\le1).
\]
The last bound follows from Schwartz decay, taking any fixed power
greater than one in the sum. The $L^2$ tail of $u^{1/2}$ in the
$s$ coordinate equals $\lambda^{-1/2}e^{-S/2}$. Apply
\eqref{eq:v12-zero-value}, increasing its arbitrary power if needed.
\end{proof}

The coefficient noncancellation and repair-size bounds are therefore no
longer assumptions for this explicitly normalized source. The tail-to-
correlation consequences and translated endpoint terms are proved below;
the global radical/form-domain passage is supplied by
Theorem~\ref{thm:v14-radical}.


The possible lattice resonance can be isolated entirely in a range where absolute summation is cheap.  Beyond that range the radial differential equation itself supplies a genuinely linear phase with a summable remainder.

\begin{lemma}[Uniform far field with controlled $c$-dependence]\label{lem:pswf-far-volterra}
For fixed $n\in\{0,4\}$ there is a scalar $A_{n,c}$ such that, uniformly for $x\ge c^2$,
\begin{equation}\label{eq:pswf-far-controlled}
 \Psi_{n,c}(x)
 =A_{n,c}\frac{\sin(cx-n\pi/2)}{x}
 +O_n\!\left(\frac{B_{n,c}\sqrt c}{x^2}\right),
 \qquad
 |A_{n,c}|\ll_n\frac{B_{n,c}}{\sqrt c}.
\end{equation}
\end{lemma}
\begin{proof}
Dunster's radial equation, with $m=0$, becomes after the standard substitution
$w(x)=(x^2-1)^{1/2}\Psi_{n,c}(x)$,
\[
 w''+c^2w=V_{n,c}(x)w,
 \qquad
 V_{n,c}(x)=\frac{\lambda_n^0(c^2)}{x^2-1}-\frac1{(x^2-1)^2}.
\]
His fixed-order relation $\sigma^2=2(n+\tfrac12)c^{-1}+O_n(c^{-2})$, together with
$\sigma^2=1+c^{-2}\lambda_n^0(c^2)$, gives
\[
 \lambda_n^0(c^2)=-c^2+(2n+1)c+O_n(1).
\]
Hence $|V_{n,c}(x)|\ll_n c^2x^{-2}$ for $x\ge2$ and
\[
 \int_{c^2}^\infty \frac{|V_{n,c}(t)|}{c}\,dt\ll_n c^{-1}.
\]
By Lemma~\ref{lem:pswf-away-energy}, $|w(x)|\ll_n B_{n,c}/\sqrt c$ for $x\ge2$.  On an interval of length $c^{-1}$ about $x_0=c^2$, the differential equation then gives
$|w''|\ll_n B_{n,c}c^{3/2}$; Taylor's formula yields
$|w'(x_0)|\ll_n B_{n,c}\sqrt c$.  Variation of constants for the perturbed oscillator and Gronwall therefore give
\[
 |w(x)|+c^{-1}|w'(x)|\ll_n B_{n,c}/\sqrt c,
 \qquad x\ge c^2.
\]
The usual incoming/outgoing amplitudes have derivatives bounded by
$|V_{n,c}(x)w(x)|/c$, which is integrable.  They therefore converge at infinity, and integration from $x$ to infinity gives an error
\[
 \ll_n \frac1c\int_x^\infty |V_{n,c}(t)w(t)|\,dt
 \ll_n \frac{B_{n,c}\sqrt c}{x}.
\]
After division by $(x^2-1)^{1/2}\asymp x$, this is the remainder in
\eqref{eq:pswf-far-controlled}.  Finally Dunster's exact radial phase at infinity, equation (1.24), identifies the limiting real combination with
$\sin(cx-n\pi/2)$, proving the stated form and amplitude bound.
\end{proof}

\begin{theorem}[Endpoint-normalized radial lattice bound]\label{thm:lattice-tail}
Put $c=2\pi\lambda^2$.  For fixed $n\in\{0,4\}$ and all $s\ge0$,
\begin{equation}\label{eq:lattice-log-target}
 \boxed{
 \left|\sum_{m\ge2}\Psi_{n,c}(me^s)\right|
 \ll_n
 \frac{B_{n,c}}{\sqrt c\,e^s}(1+\log c).}
\end{equation}
The estimate is uniform in $s$ and hence contains no unproved nonresonance hypothesis.
\end{theorem}
\begin{proof}
Set $x=e^s$ and
\[
 M=\max\!\left(2,\left\lceil\frac{c^2}{x}\right\rceil\right).
\]
For $2\le m<M$, Lemma~\ref{lem:pswf-away-energy} gives, by absolute summation,
\[
 \sum_{2\le m<M}|\Psi_{n,c}(mx)|
 \ll_n\frac{B_{n,c}}{\sqrt c\,x}\sum_{m<M}\frac1m
 \ll_n\frac{B_{n,c}(1+\log c)}{\sqrt c\,x}.
\]
For $m\ge M$, Lemma~\ref{lem:pswf-far-volterra} applies.  Since $n=0$ or $4$, the fixed phase $n\pi/2$ is an integral multiple of $2\pi$.  The elementary harmonic-sine bound
\[
 \sup_{\theta\in\mathbb R}\sup_{M\ge1}
 \left|\sum_{m=M}^{\infty}\frac{\sin(m\theta)}m\right|<\infty
\]
then gives
\[
 \left|\frac{A_{n,c}}x\sum_{m\ge M}\frac{\sin(cmx-n\pi/2)}m\right|
 \ll_n\frac{B_{n,c}}{\sqrt c\,x}.
\]
(The displayed uniform bound follows, for example, by splitting at $m\asymp |\theta|^{-1}$ and applying Dirichlet summation to the remaining tail.)
The remainder in \eqref{eq:pswf-far-controlled} contributes
\[
 \ll_n\frac{B_{n,c}\sqrt c}{x^2}\sum_{m\ge M}m^{-2}.
\]
If $M\asymp c^2/x$ this is $O_n(B_{n,c}/(c^{3/2}x))$; if $M=2$, then $x\gg c^2$ and it is even smaller than $B_{n,c}/(\sqrt c\,x)$.  Combining the two ranges proves \eqref{eq:lattice-log-target}.
\end{proof}

The $m=1$ Poisson sample is a boundary-layer term rather than a lattice-summation problem.  The finite-Fourier normalization introduces no additional power of $\lambda$: in the standard PSWF convention
\[
 \mathcal F_c\Psi_{n,c}=\mu_n(c)\Psi_{n,c},\qquad
 \frac{c}{2\pi}|\mu_n(c)|^2=\chi_n(c)^2,
\]
where $\chi_n^2$ is the concentration eigenvalue.  Since $c=2\pi\lambda^2$,
$|\mu_n(c)|=\chi_n(c)/\lambda$.  After scaling from $[-1,1]$ to
$[-\lambda,\lambda]$, if $u=\lambda^{-1}e^{-s}$ then the Fourier-lattice part of the lower tail is, up to the harmless fixed even-mode Fourier phase,
\begin{equation}\label{eq:scaled-poisson-radial}
 r_{n,\lambda}^{\rm lat}(s)
 =\chi_n(c)e^{s/2}\sum_{m\ge1}\Psi_{n,c}(me^s).
\end{equation}
For an even source with $\int h=0$ but $h(0)\ne0$, Poisson summation has in addition the elementary term
\begin{equation}\label{eq:poisson-zero-defect}
 \mathcal E(h)(u)
 =-\frac12u^{1/2}h(0)
   +u^{-1/2}\sum_{m\ge1}\widehat h(m/u).
\end{equation}
Thus the raw CCM two-mode source has a small zero-value defect in addition to the radial lattice term.  It will be estimated explicitly below rather than silently omitted.  In particular, at $s=0$ the $m=1$ lattice term is $\chi_n(c)$ times the upper-window boundary amplitude.

\begin{corollary}[Fixed-order co-Poisson tail and endpoint comparison]\label{cor:pswf-cross-tail}
Let $h_\lambda^{\circ}$ be the explicitly normalized two-mode source of
Proposition~\ref{prop:v12-source}, and let
$B_\lambda^{+}=\sqrt\lambda\,h_\lambda^{\circ}(\lambda)$ be its upper-window co-Poisson boundary amplitude.  Then the lower logarithmic tail $r_\lambda^{\circ}$ satisfies, for every $S\ge0$,
\begin{equation}\label{eq:lower-tail-l2}
 \left\|1_{\{s\ge S\}}r_\lambda^{\circ}\right\|_2
 \ll \frac{|B_\lambda^{+}|}{\lambda}(1+\log\lambda)e^{-S/2}.
\end{equation}
The same estimate holds for the exact value-and-integral repair. If
$B_\lambda$ is the common one-sided endpoint value of the inversion-even
finite-window proxy, then
\begin{equation}\label{eq:v13-endpoint-comparison}
 \frac{B_\lambda}{B_\lambda^+}
 =1+O\!\left(\frac{1+\log c}{\sqrt c}\right),
 \qquad c=2\pi\lambda^2.
\end{equation}
In particular $B_\lambda\asymp B_\lambda^+$ and it is nonzero for all
sufficiently large $\lambda$.
\end{corollary}
\begin{proof}
For the $m=1$ term in \eqref{eq:scaled-poisson-radial}, Lemma~\ref{lem:pswf-away-energy} gives the exact log-coordinate estimate
\[
 \int_0^\infty
 \left|\chi_n(c)e^{s/2}\Psi_{n,c}(e^s)\right|^2ds
 =\chi_n(c)^2\int_1^\infty|\Psi_{n,c}(x)|^2dx
 \ll_n\frac{B_{n,c}^2}{c}.
\]
For $s\ge\log2$, \eqref{eq:pswf-away} also gives the pointwise bound
$O_n(B_{n,c}c^{-1/2}e^{-s/2})$.  Thus its $L^2$ tail beyond depth $S$ is
$O_n(B_{n,c}c^{-1/2}e^{-S/2})$ for all $S\ge0$, after enlarging the constant on $0\le S\le\log2$.

For $m\ge2$, Theorem~\ref{thm:lattice-tail} and \eqref{eq:scaled-poisson-radial} give directly
\[
 |r_{n,\lambda}^{(\ge2)}(s)|
 \ll_n \frac{B_{n,c}(1+\log c)}{\sqrt c}e^{-s/2}.
\]
Since $\sqrt c\asymp\lambda$, this proves \eqref{eq:lower-tail-l2} mode by mode. The endpoint noncancellation bound in Proposition~\ref{prop:v12-source} controls the sum of the individual endpoint magnitudes by $|B_\lambda^+|$, so the modewise estimate passes to the actual normalized combination. This part no longer uses an unproved coefficient estimate.

For completeness, the exact radical source condition can be imposed without changing this scale.  Choose a fixed even $\psi\in C_c^\infty((-1,1))$ with $\psi(0)=1$ and $\int\psi=0$, and put
\[
 \widetilde h_\lambda=h_\lambda^{\circ}-h_\lambda^{\circ}(0)\psi.
\]
Then $\widetilde h_\lambda(0)=0$ and $\int\widetilde h_\lambda=0$, while its endpoint at $\pm\lambda$ is unchanged. Proposition~\ref{prop:v12-source} proves
$|h_\lambda^\circ(0)|/|B_\lambda^+|=O_A(\lambda^{-A})$ for every fixed $A$,
and Corollary~\ref{cor:v12-repair} supplies the lower-tail repair bound.
The repaired compact source satisfies Theorem~\ref{thm:v14-radical},
so its global radical and semilocal form-domain use is justified.
For the fixed bump one must retain its zero-value term: since $\widehat\psi(0)=\int\psi=0$ and $\psi(0)=1$, Poisson summation gives
\[
 \mathcal E(\psi)(u)
 =-\frac12u^{1/2}
   +u^{-1/2}\sum_{m\ge1}\widehat\psi(m/u).
\]
The second term is rapidly decreasing as $u\downarrow0$, while the first is $O(u^{1/2})$.  Because $|h_\lambda^{\circ}(0)|/|B_\lambda^{+}|=O_A(\lambda^{-A})$ for every fixed $A$, the whole correction is still $o(|B_\lambda^{+}|\lambda^{-1}e^{-s/2})$ at $u=\lambda^{-1}e^{-s}$ and is negligible in \eqref{eq:lower-tail-l2}.  The same estimate also absorbs the raw zero-defect term in \eqref{eq:poisson-zero-defect}.

It remains to identify the endpoint after inversion. Put
$f_\lambda=\mathcal E(\widetilde h_\lambda)$. At the upper endpoint only
the $m=1$ source summand remains, so
$f_\lambda(\lambda^-)=B_\lambda^+$. At the lower endpoint, Poisson
summation and \eqref{eq:scaled-poisson-radial} give
\[
 f_\lambda((\lambda^{-1})^+)
 =\sum_{n=0,4}a_{n,\lambda}\chi_n(c)
       \sqrt\lambda h_{n,\lambda}(\lambda)
 +O\!\left(\frac{1+\log c}{\sqrt c}
       \sum_{n=0,4}|a_{n,\lambda}|\sqrt\lambda
                    |h_{n,\lambda}(\lambda)|\right)
 +o(|B_\lambda^+|).
\]
The last term includes the fixed-bump repair. If $\lambda^2$ is an
integer, taking the indicated interior trace can also remove the single
source summand at the cut; its size is $O(|B_\lambda^+|/\lambda)$ and is
absorbed above. Since $\chi_n(c)=1+o(1)$ and
\eqref{eq:v12-noncancellation} controls the sum of the individual endpoint
magnitudes by $|B_\lambda^+|$, the displayed expression equals
$B_\lambda^+\{1+O((1+\log c)/\sqrt c)\}$. Finally
$p_\lambda=(f_\lambda+Jf_\lambda)/2$, so its two common one-sided endpoint
values are the average of the upper and lower traces. This proves
\eqref{eq:v13-endpoint-comparison}.
\end{proof}

\begin{lemma}[Uniform zero-extension correlations and the local term]
\label{lem:v13-correlations}
Let $L=2\log\lambda$, let $r_\lambda$ be either oriented logarithmic tail
of the repaired source, and zero-extend $v\in H^1([-L,0])$. For
\[
 C_{r,v}(y)=\int_{\mathbb R}r_\lambda(t)
                  \overline{v(t-y)}\,dt
\]
and for each reflected correlation occurring in the polarized even Weil
kernel, the interior-test derivative contribution satisfies
\begin{equation}\label{eq:pswf-cross-tail}
 \boxed{
 |C_{r,v,\mathrm{bulk}}'(y)|
 \le C\frac{|B_\lambda|}{\lambda}(1+\log\lambda)
       \|v'\|_2
 \begin{cases}
  1,&0\le y\le L,\\[2mm]
  e^{-(y-L)/2},&y>L.
 \end{cases}}
\end{equation}
Here $C_{r,v,\mathrm{bulk}}'(y)=-\int r_\lambda(t)
\overline{v'(t-y)}1_{(-L,0)}(t-y)\,dt$.
The full distributional derivative has, in addition, precisely the two terms
\begin{equation}\label{eq:v13-correlation-jumps}
 -\overline{v(-L)}\,r_\lambda(y-L)
 +\overline{v(0)}\,r_\lambda(y),
\end{equation}
up to the harmless sign and reflection dictated by the chosen orientation;
a term with a negative tail argument is zero. These translated-tail terms
are themselves locally integrable densities, not singular atoms in the
derivative of the correlation. They must not be omitted from its
absolutely continuous part. Moreover the local
archimedean functional of the even polarized sum $q_\lambda$ obeys
\begin{equation}\label{eq:v13-archimedean}
 |\mathcal A_\infty(q_\lambda)|
 \le C\frac{|B_\lambda|}{\lambda}(1+\log\lambda)
 \bigl(\|v\|_2+\|v'\|_2+|v(-L)|+|v(0)|\bigr).
\end{equation}
These estimates are uniform in $v$ and require no form-domain assumption.
\end{lemma}
\begin{proof}
Inside the support interval, differentiation in $y$ puts one derivative
on $v$ and never on the prolate tail. The overlap uses tail depths
$t\ge(y-L)_+$. Cauchy--Schwarz and \eqref{eq:lower-tail-l2}, followed by
\eqref{eq:v13-endpoint-comparison}, prove
\eqref{eq:pswf-cross-tail}. In distributions,
\[
 D(1_{[-L,0]}v)=1_{(-L,0)}v'
       +v(-L)\delta_{-L}-v(0)\delta_0.
\]
Substitution into the correlation, with the minus sign from differentiating
$v(t-y)$, gives \eqref{eq:v13-correlation-jumps}. Reflection gives the
other orientation and changes none of the bounds.

For the singularity at the origin use the elementary zero-extension
translation estimate, valid for $0<x\le1$,
\begin{equation}\label{eq:v13-translation}
 \|v(\cdot-x)-v\|_{L^2(\mathbb R)}
 \le C\left[x\|v'\|_2
   +\sqrt{x}\{ |v(-L)|+|v(0)|\}\right].
\end{equation}
It follows by applying the fundamental theorem of calculus on the overlap
and estimating the two boundary strips separately. Thus, with
$R_\lambda=|B_\lambda|(1+\log\lambda)/\lambda$,
\[
 |q_\lambda(x)+q_\lambda(-x)-2e^{-x/2}q_\lambda(0)|
 \le C R_\lambda\left[
  \sqrt{x}\{|v(-L)|+|v(0)|\}
  +x\{\|v'\|_2+\|v\|_2\}\right].
\]
The archimedean kernel is $O(x^{-1})$ on $(0,1]$, so this is integrable.
For $x\ge1$ the kernel is $O(e^{-x/2})$, while Cauchy--Schwarz and
\eqref{eq:lower-tail-l2} bound every correlation by
$CR_\lambda\|v\|_2$. Integration proves
\eqref{eq:v13-archimedean}.
\end{proof}

The truncated-Mellin identity remains a useful independent check on the powers and normalizations.  For an exact two-moment source,
\begin{equation}\label{eq:truncated-mellin-tail}
 \mathcal M^-_\lambda(\tau)
 =\int_0^{1/\lambda}\mathcal E(h_\lambda)(u)u^{-i\tau}\,\frac{du}{u},
\end{equation}
and Poisson summation gives
\begin{equation}\label{eq:truncated-mellin-poisson}
 \mathcal M^-_\lambda(\tau)
 =\sum_{m\ge1}m^{-1/2-i\tau}
   \int_{m\lambda}^{\infty}\widehat h_\lambda(x)
   x^{-1/2+i\tau}\,dx.
\end{equation}
The preceding argument unconditionally supplies the physical/logarithmic
tail and correlation estimates needed below, given the cited classical
fixed-order PSWF inputs. Theorem~\ref{thm:v14-radical} justifies passing
the global radical identity to the closed semilocal form at each fixed
$\lambda$; the quantitative estimates, not its approximation constants,
provide the required uniformity as $\lambda\to\infty$.


\begin{lemma}[Quantitative endpoint-jump bound]\label{lem:endpoint-jumps}
Let $E_{\lambda,T}$ be the periodic trigonometric space on $[-a,a]$, $a=\log\lambda$, with physical logarithmic frequencies in $[-T,T]$.  Extend $v\in E_{\lambda,T}$ by zero outside $[-a,a]$, and let $q_\lambda$ be its logarithmic correlation with the lower tail $r_\lambda$.  In the distributional derivative of $q_\lambda$, the two endpoint jumps contribute translates of $r_\lambda$ multiplied by the common periodic endpoint value $v(a)=v(-a)$.  There is a constant $c_1>0$, independent of $\lambda$, such that uniformly for $\|v\|_2=1$ their total contribution to the centered pole--prime Stieltjes pairing is
\begin{equation}\label{eq:endpoint-jump-rate}
 \ll_T |B_\lambda|\left((1+L)e^{-L/2}
 +(1+L)^{3/2}e^{-c_1\sqrt L}\right),
 \qquad L=2\log\lambda.
\end{equation}
In particular the endpoint-jump contribution is $o_T(|B_\lambda|)$.
\end{lemma}
\begin{proof}
The endpoint evaluation functional on $E_{\lambda,T}$ has norm $O_T(1)$ because the normalized Fourier basis has $O_T(L)$ modes and each basis vector has endpoint magnitude $L^{-1/2}$.  For the jump meeting the lower tail at depth $s=y$, split $0\le s\le\log2$ from $s\ge\log2$.  On the first interval, $V(s)=O(s)$ and Cauchy--Schwarz with \eqref{eq:lower-tail-l2} gives
\[
 O_T\!\left(|B_\lambda|\frac{1+\log\lambda}{\lambda}\right)
 =O_T\!\left(|B_\lambda|(1+L)e^{-L/2}\right).
\]
On the second, Theorem~\ref{thm:lattice-tail} together with the $m=1$ away-bound gives
\[
 |r_\lambda(s)|\ll |B_\lambda|\frac{1+\log\lambda}{\lambda}e^{-s/2},
 \qquad s\ge\log2,
\]
and \eqref{eq:centered-prime-cumulative} below supplies an integrable factor $e^{-c\sqrt s}$.  This contributes the same $O_T(|B_\lambda|(1+L)e^{-L/2})$ scale.

The second endpoint enters only after a shift $y=L+s$.  There
\[
 |V(L+s)|\ll \lambda e^{s/2}e^{-c\sqrt{L+s}}.
\]
For $0\le s\le\log2$, Cauchy--Schwarz and \eqref{eq:lower-tail-l2} cancel the factor $\lambda$ and give
$O_T(|B_\lambda|(1+L)e^{-c\sqrt L})$.  For $s\ge\log2$, the preceding pointwise tail estimate cancels $e^{s/2}$ and leaves
\[
 O_T\!\left(|B_\lambda|(1+L)
 \int_0^\infty e^{-c\sqrt{L+s}}\,ds\right).
\]
The integral is explicit:
\[
 \int_0^\infty e^{-c\sqrt{L+s}}\,ds
 =2e^{-c\sqrt L}\left(\frac{\sqrt L}{c}+\frac1{c^2}\right).
\]
After decreasing the exponential constant if necessary, these estimates give \eqref{eq:endpoint-jump-rate}.
\end{proof}

\subsection*{Pole--prime reduction and closure of the weak G1 estimate}
Let $k_\lambda$ be the exact-radical co-Poisson vector obtained from the corrected source, and write
\[
 k_\lambda=k_{\lambda,\mathrm{in}}+r_\lambda,
 \qquad k_{\lambda,\mathrm{in}}=1_{I_\lambda}k_\lambda.
\]
For every zero-extended $H^1$ test vector $v$ supported in $I_\lambda$,
Theorem~\ref{thm:v14-radical} gives
\begin{equation}\label{eq:radical-truncation}
 QW(k_{\lambda,\mathrm{in}},v)=-QW(r_\lambda,v).
\end{equation}

The dominant point-mass part and the prime part of the explicit formula pair the logarithmic correlation $q_\lambda$ with
\[
 e^{y/2}\,dy
 \quad\text{and}\quad
 \sum_{n\ge2}\frac{\Lambda(n)}{\sqrt n}\,\delta_{\log n},
\]
respectively; the additional $e^{-y/2}$ part is absolutely integrable.  Put
\[
 V(Y)=2(e^{Y/2}-1)-\sum_{2\le n\le e^Y}\frac{\Lambda(n)}{\sqrt n}.
\]
Partial summation together with the classical de la Vall\'ee-Poussin prime-number-theorem remainder gives constants $c,C>0$ such that
\begin{equation}\label{eq:centered-prime-cumulative}
 |V(Y)|\le C e^{Y/2}e^{-c\sqrt Y}\qquad(Y\ge1).
\end{equation}
Stieltjes integration by parts can be made quantitative.  Put
\[
 A_{\lambda,T}=C_T|B_\lambda|(1+L)e^{-L/2},
\]
so that \eqref{eq:pswf-cross-tail} bounds the interior-test derivative part by
$A_{\lambda,T}$ on $[0,L]$ and by
$A_{\lambda,T}e^{-(y-L)/2}$ on $[L,\infty)$.  On $1\le y\le L$ split at
$L/4$.  The first part is
$O_T(|B_\lambda|(1+L)L e^{-3L/8})$, while on $[L/4,L]$ the PNT factor gives
$e^{-c\sqrt y}\le e^{-c\sqrt L/2}$ and the $e^{y/2}$ integral cancels the
factor $e^{-L/2}$.  Beyond $L$ one obtains
\[
 A_{\lambda,T}e^{L/2}\int_L^\infty e^{-c\sqrt y}\,dy
 =2A_{\lambda,T}e^{L/2-c\sqrt L}
 \left(\frac{\sqrt L}{c}+\frac1{c^2}\right).
\]
Together with Lemma~\ref{lem:endpoint-jumps}, this yields constants
$c_1,C_T>0$ such that
\begin{equation}\label{eq:pole-prime-small}
 \left|\int_0^\infty q_\lambda(y)e^{y/2}\,dy
 -\sum_{n\ge2}\frac{\Lambda(n)}{\sqrt n}q_\lambda(\log n)\right|
 \le C_T|B_\lambda|\left((1+L)^{3/2}e^{-c_1\sqrt L}
 +(1+L)e^{-L/2}\right).
\end{equation}
The constants use only the classical unconditional zero-free-region/PNT remainder; no
RH assumption or information about individual nontrivial zeros is inserted.

For completeness, the archimedean local term may be estimated directly rather than hidden inside a multiplier statement.  In logarithmic coordinates the standard Weil formula contributes, after polarization,
\[
 \mathcal A_\infty(q)=-(\log 4\pi+\gamma_0)q(0)
 -\int_0^\infty
 \{q(x)+q(-x)-2e^{-x/2}q(0)\}
 \frac{e^{x/2}}{e^x-e^{-x}}\,dx .
\]
The kernel is $O(e^{-x/2})$ for $x\ge1$ and is $O(x^{-1})$ at the origin.
For zero-extended $H^1$ tests the numerator need only be $O(x^{1/2})$
because of the endpoint strips; this is still integrable. Lemma~\ref{lem:v13-correlations}
gives precisely the required uniform estimate:
\[
 |q(0)|\ll |B_\lambda|e^{-L/2}(1+\log\lambda),\qquad
 |\mathcal A_\infty(q_\lambda)|\ll_T |B_\lambda|e^{-L/2}(1+\log\lambda)
 =O_T\!\left(|B_\lambda|\frac{1+\log\lambda}{\lambda}\right).
\]
The remaining $e^{-y/2}$ pole contribution is absolutely integrable and obeys the same bound.

Define the explicit zero-independent envelope
\begin{equation}\label{eq:G1-envelope}
 \mathfrak r(\lambda)
 :=(1+L)^{3/2}e^{-c_1\sqrt L}+(1+L)e^{-L/2},
 \qquad L=2\log\lambda.
\end{equation}
Combining the pole--prime, archimedean, and decaying-pole pieces with global radicality gives the quantitative continuum G1 estimate
\begin{equation}\label{eq:G1-continuum}
 \boxed{
 \mathcal E_T(\lambda):=
 \sup_{\substack{\|v\|_2=1\\ v\in E_{\lambda,T}}}
 \frac{|QW(k_{\lambda,\mathrm{in}},v)|}{|B_\lambda|}
 \le C_T\mathfrak r(\lambda)\longrightarrow0.}
\end{equation}
The quantitative estimates use only fixed-order PSWF asymptotics, the
radial differential equation, Poisson summation, and the classical
unconditional prime-number-theorem remainder. The global radical and
semilocal domain passage is now Theorem~\ref{thm:v14-radical}; hence
\eqref{eq:G1-continuum} is proved for the explicitly repaired source,
with no remaining source assumption. The construction and all displayed
cutoff rules are independent of a nontrivial-zero list. This is the weak
G1 conclusion for a fixed physical band, not a uniform bound on growing
prolate blocks or an RH implication by itself.

\begin{lemma}[Inversion symmetrization preserves the G1 envelope]\label{lem:g1-inversion}
Let
\[
 p_\lambda=\tfrac12(I+J)k_{\lambda,\mathrm{in}},
 \qquad Jf(u)=f(u^{-1}).
\]
Then
\begin{equation}\label{eq:G1-sym-envelope}
 \sup_{\substack{\|v\|_2=1\\v\in E_{\lambda,T}}}
 \frac{|QW(p_\lambda,v)|}{|B_\lambda|}
 \le \mathcal E_T(\lambda)
 \le C_T\mathfrak r(\lambda).
\end{equation}
\end{lemma}
\begin{proof}
The Weil form is invariant under simultaneous inversion,
$QW(Jf,Jg)=QW(f,g)$, and $J$ is a unitary involution preserving the physical
frequency window $E_{\lambda,T}$; this is the semilocal form of the inversion
symmetry in Connes--Consani--Moscovici~\cite{ConnesConsaniMoscovici2025}.
Hence
\[
 QW(Jk_{\lambda,\mathrm{in}},v)
 =QW(k_{\lambda,\mathrm{in}},Jv).
\]
Taking the average with $QW(k_{\lambda,\mathrm{in}},v)$ and then the supremum
proves \eqref{eq:G1-sym-envelope}.
\end{proof}


\begin{proposition}[Sobolev form of the continuum G1 estimate]\label{prop:g1-sobolev}
Use the explicit repaired source above. Let $u$ be a periodic $H^1$ function on the logarithmic window $[-a,a]$,
$a=\log\lambda$, and interpret it by zero extension when it enters the
semilocal Weil form.  Put
\[
 \mathcal S_1(u):=\|u\|_2+\|D_{\log}u\|_2
 +|u(-a)|+|u(a)|.
\]
Then the same zero-independent envelope as in \eqref{eq:G1-continuum} obeys
\begin{equation}\label{eq:G1-sobolev}
 \boxed{
 \frac{|QW_\lambda(p_\lambda,u)|}{|B_\lambda|}
 \le C\,\mathfrak r(\lambda)\,\mathcal S_1(u).}
\end{equation}
The constant is independent of $\lambda$ and $u$.  In particular the
fixed-physical-band hypothesis in \eqref{eq:G1-sym-envelope} is a convenient
way to bound the Sobolev and trace factors, not an essential spectral
restriction on the continuum G1 pairing.
\end{proposition}
\begin{proof}
Repeat the proof of the cross-tail and pole--prime estimates without applying
a Bernstein inequality.  On the open logarithmic interval, differentiation
of the correlation is placed on $u$, so Cauchy--Schwarz gives the factor
$\|D_{\log}u\|_2$ in place of the earlier band constant $T$.  The two jumps
created by zero extension have coefficients exactly $u(-a)$ and $u(a)$, so
Lemma~\ref{lem:endpoint-jumps} carries over with the endpoint-evaluation norm
replaced by $|u(-a)|+|u(a)|$.  Finally
$|q(0)|\le\|r_\lambda\|_2\|u\|_2$, and the archimedean numerator is
controlled near the origin by the same $H^1$ correlation estimate.  The
Stieltjes integration-by-parts calculation leading to
\eqref{eq:pole-prime-small} is otherwise unchanged and gives
\[
 |QW_\lambda(k_{\lambda,\mathrm{in}},u)|
 \le C|B_\lambda|\mathfrak r(\lambda)\mathcal S_1(u).
\]
Simultaneous inversion preserves the $L^2$ norm, the logarithmic derivative
norm, and the two endpoint traces.  Averaging with the inverted vector as in
Lemma~\ref{lem:g1-inversion} proves \eqref{eq:G1-sobolev}.
\end{proof}

The ordinary Fourier finite-section issue requires more care.  The relevant regularity is supplied directly by the arithmetic finite-sum formula, rather than by a global continuity assumption.

\begin{lemma}[Arithmetic BV structure and absence of a periodic boundary jump]\label{lem:g1-bv-structure}
Let $\widetilde h_\lambda$ be the compactly supported exact-moment prolate source used above, and let
\[
 \mathcal E(\widetilde h_\lambda)(u)
 =u^{1/2}\sum_{n\ge1}\widetilde h_\lambda(nu)
\]
be its co-Poisson/Riemann-sum vector.  On the window
$I_\lambda=[\lambda^{-1},\lambda]$, the restriction of
$\mathcal E(\widetilde h_\lambda)$ is of bounded variation and has only finitely many possible interior jump points.  Consequently
\[
 p_\lambda=\tfrac12(I+J)k_{\lambda,\mathrm{in}},
 \qquad Jf(u)=f(u^{-1}),
\]
is also of bounded variation.  Moreover the inversion identity
$p_\lambda(u)=p_\lambda(u^{-1})$ implies
\begin{equation}\label{eq:g1-endpoint-limits}
 p_\lambda(\lambda^-)=p_\lambda((\lambda^{-1})^+)=:B_\lambda,
\end{equation}
so the periodic representative in logarithmic coordinates has no jump at the identified boundary.  No assertion of global continuity is needed.
\end{lemma}
\begin{proof}
The scale-invariant Riemann-sum formula
$\mathcal E(f)(u)=u^{1/2}\sum_{n>0}f(nu)$ is the standard co-Poisson map used in the prolate construction~\cite{ConnesConsani2023}.  Because
$\operatorname{supp}\widetilde h_\lambda\subset[-\lambda,\lambda]$, for
$u\in[\lambda^{-1},\lambda]$ only indices
$n\le\lambda/u\le\lambda^2$ can occur.  Thus the sum has only finitely many summands throughout the window.  Between the finitely many threshold points at which an argument $nu$ meets an endpoint or another piece boundary of the compact source, it is a finite sum of smooth functions; its one-sided limits are finite.  Hence its total variation on the compact window is finite.  The fixed smooth exact-moment repair used above does not alter this conclusion.

Inversion preserves bounded variation, so the symmetrized proxy $p_\lambda$ is BV.  Finally, by construction $p_\lambda=Jp_\lambda$.  Taking one-sided limits at the two ends of $I_\lambda$ gives \eqref{eq:g1-endpoint-limits}.  After the logarithmic endpoints are identified, this is exactly the absence of a periodic boundary jump.  Interior arithmetic jumps remain allowed.
\end{proof}

Thus for each \emph{fixed} cutoff ordinary Fourier projection converges both in the closed semilocal Weil form norm and at the endpoint.  This gives the following fixed-cutoff existence statement, but no quantitative diagonal compatible with the growing-band estimate is supplied.

\begin{proposition}[Fixed-cutoff boundary-preserving Galerkin approximation]\label{prop:boundary-galerkin}
Assume $p_\lambda$ is of bounded variation on $[-a,a]$ and that its one-sided endpoint limits agree, with common nonzero value
\[
 B_\lambda=p_\lambda(a^-)=p_\lambda((-a)^+).
\]
Let $P_N$ be ordinary Fourier projection onto the trigonometric space $E_N$.  For each fixed $\lambda$,
\[
 P_Np_\lambda\to p_\lambda
\]
in the form norm of the closed semilocal Weil form and
\[
 (P_Np_\lambda)(a)\to B_\lambda.
\]
Consequently, for every fixed $\lambda$ and every $\varepsilon>0$ there is an $N=N(\lambda,\varepsilon)$ such that
\[
 \|P_Np_\lambda-p_\lambda\|_{QW_\lambda}
 +|(P_Np_\lambda)(a)-B_\lambda|
 \le \varepsilon |B_\lambda|.
\]
No rate relating $N$ to $\lambda$ is asserted.
\end{proposition}
\begin{proof}
A bounded-variation function on the log interval has Fourier coefficients $\widehat p_\lambda(n)=O_\lambda(|n|^{-1})$.  The archimedean part of the semilocal Weil form has Fourier weight comparable to $1+\log(1+n^2)$, while at fixed $\lambda$ the pole term and the finite prime-power sum are bounded perturbations.  Hence
\[
 \sum_n(1+\log(1+n^2))|\widehat p_\lambda(n)|^2<\infty,
\]
and the Fourier tails tend to zero in the form norm.  Since the two one-sided endpoint limits agree, the periodic representative has no boundary jump; the Dirichlet--Jordan theorem therefore gives convergence at the identified endpoint.  Interior jumps do not affect this conclusion.
\end{proof}

The lack of a rate is not cosmetic.  The physical-band estimate proved above remains uniform only while the physical frequency grows sufficiently slowly compared with the quantitative PNT saving.  A generic Fourier approximation of a function with an exponentially small endpoint value may require a much larger cutoff.  Therefore Proposition~\ref{prop:boundary-galerkin} by itself does \emph{not} transfer \eqref{eq:G1-continuum} to a diagonal $N(\lambda)$.

There is nevertheless a zero-independent finite-dimensional transfer of the
\emph{fixed-band residual}, and the quantitative envelope above makes its exact
finite-section dependence transparent.  For $N$ large enough to contain the
physical band, set
\[
 k_{\lambda,N}=P_Np_\lambda,
\]
and define the relative boundary and band-form defects
\begin{align}
 \beta_{\lambda,N}
 &:=\frac{|\delta_N(k_{\lambda,N})-B_\lambda|}{|B_\lambda|},
 \label{eq:g1-beta}\\
 \gamma_{\lambda,N,T}
 &:=\frac1{|B_\lambda|}
 \sup_{\substack{v\in E_{\lambda,T}\\\|v\|_2=1}}
 |QW_\lambda(k_{\lambda,N}-p_\lambda,D_{\log}v)|.
 \label{eq:g1-gamma}
\end{align}
Both quantities are defined entirely from the prolate/co-Poisson proxy and the
finite Weil form; no zeta zero is used.

\begin{proposition}[Quantitative and cofinal finite-section transfer of G1]\label{prop:finite-g1-transfer}
Fix $T>0$.  Whenever $\beta_{\lambda,N}<1$,
\begin{equation}\label{eq:g1-master-bound}
 \boxed{
 \frac{\|P_TD_{\log}W_{\lambda,N}k_{\lambda,N}\|_2}
 {|\delta_N(k_{\lambda,N})|}
 \le
 \frac{T\mathcal E_T(\lambda)+\gamma_{\lambda,N,T}}
 {1-\beta_{\lambda,N}}.}
\end{equation}
For each fixed $\lambda$,
\begin{equation}\label{eq:g1-defects-fixed-lambda}
 \beta_{\lambda,N}\longrightarrow0,
 \qquad
 \gamma_{\lambda,N,T}\longrightarrow0
 \qquad(N\to\infty).
\end{equation}
Consequently, for all sufficiently large $\lambda$ there is a zero-independent
threshold $N_{G1}(\lambda,T)$ such that for \emph{every}
$N\ge N_{G1}(\lambda,T)$,
\begin{align}
 \beta_{\lambda,N}&\le\mathfrak r(\lambda),
 \label{eq:g1-threshold-beta}\\
 \gamma_{\lambda,N,T}&\le\mathfrak r(\lambda),
 \label{eq:g1-threshold-gamma}
\end{align}
and hence
\begin{equation}\label{eq:finite-g1-transfer}
 \boxed{
 \sup_{N\ge N_{G1}(\lambda,T)}
 \frac{\|P_TD_{\log}W_{\lambda,N}k_{\lambda,N}\|_2}
 {|\delta_N(k_{\lambda,N})|}
 \le C'_T\mathfrak r(\lambda)\longrightarrow0.}
\end{equation}
In particular one may choose any diagonal $N(\lambda)\ge N_{G1}(\lambda,T)$;
G1 does not require an economical or finely tuned Fourier cutoff.
\end{proposition}
\begin{proof}
Hermiticity of the finite Weil matrix and self-adjointness of $D_{\log}$ on the
periodic Fourier section give
\[
 \|P_TD_{\log}W_{\lambda,N}k_{\lambda,N}\|_2
 =\sup_{\substack{v\in E_{\lambda,T}\\\|v\|_2=1}}
 |QW_\lambda(k_{\lambda,N},D_{\log}v)|.
\]
Split $k_{\lambda,N}=p_\lambda+(k_{\lambda,N}-p_\lambda)$.  Since
$D_{\log}v\in E_{\lambda,T}$ and
$\|D_{\log}v\|_2\le T$, Lemma~\ref{lem:g1-inversion} bounds the first part by
$T|B_\lambda|\mathcal E_T(\lambda)$, while the second is exactly bounded by
$|B_\lambda|\gamma_{\lambda,N,T}$.  Also
\[
 |\delta_N(k_{\lambda,N})|
 \ge |B_\lambda|(1-\beta_{\lambda,N}).
\]
This proves \eqref{eq:g1-master-bound}.

For fixed $\lambda$, Proposition~\ref{prop:boundary-galerkin} gives the boundary
convergence and form-norm convergence of $P_Np_\lambda-p_\lambda$, proving the
first limit in \eqref{eq:g1-defects-fixed-lambda}.  To obtain the second without
an unspecified pointwise continuity constant, choose $m_\lambda$ so that
\[
 QW_\lambda(u,u)+m_\lambda\|u\|_2^2
\]
is a positive closed form.  Its Cauchy--Schwarz inequality, together with the
finite-dimensional boundedness of
$\{D_{\log}v:v\in E_{\lambda,T},\ \|v\|_2=1\}$ in the corresponding form norm,
shows that the supremum in \eqref{eq:g1-gamma} tends to zero with the form norm
of $P_Np_\lambda-p_\lambda$.

The convergence in \eqref{eq:g1-defects-fixed-lambda} is convergence of the
full sequence, so for each fixed $\lambda$ and the positive tolerance
$\mathfrak r(\lambda)$ there is a threshold after which
\eqref{eq:g1-threshold-beta}--\eqref{eq:g1-threshold-gamma} hold for all larger
$N$.  Since $\mathfrak r(\lambda)\to0$, it is $<1/2$ for large $\lambda$.
Insert these inequalities and
$\mathcal E_T(\lambda)\le C_T\mathfrak r(\lambda)$ into
\eqref{eq:g1-master-bound} to obtain \eqref{eq:finite-g1-transfer}.
\end{proof}


\subsection*{Quantitative high--low decay and an exact-boundary G1 repair}
Endpoint recovery of the ordinary Fourier partial sums in
Proposition~\ref{prop:finite-g1-transfer} was initially nonquantitative.
Proposition~\ref{prop:v114-polynomial-endpoint} below now supplies an
asymptotic polynomial cutoff for the fixed source. For the G1
residual itself one can also bypass endpoint selection.  The key observation is that the
explicit CCM matrix has an additional off-diagonal factor, so only bounded
variation of the co-Poisson proxy is needed.

Write $U_n=L^{-1/2}e^{2\pi inx/L}$ for the normalized logarithmic Fourier basis and $V_n=\kappa(U_n)$ for its multiplicative image, and set
\[
 \tau^{(\lambda)}_{n,m}:=QW_\lambda(V_n,V_m).
\]

\begin{lemma}[Zero-independent high--low decay of the Weil matrix]\label{lem:g1-high-low}
There is an absolute constant $C$ such that, for $\lambda\ge2$,
$L=2\log\lambda$, and $n\ne m$,
\begin{equation}\label{eq:g1-high-low}
 \boxed{
 |\tau^{(\lambda)}_{n,m}|
 \le C\frac{\lambda}{|n-m|}.}
\end{equation}
The constant uses no information about the nontrivial zeros of $\zeta$.
\end{lemma}
\begin{proof}
For $n\ne m$, the explicit convolution kernel of
Connes--Consani--Moscovici~\cite{ConnesConsaniMoscovici2025} is
\[
 q(U_n,U_m)(y)
 =\frac{\sin(2\pi my/L)-\sin(2\pi ny/L)}{\pi(n-m)},
 \qquad 0\le y\le L,
\]
so $q(U_n,U_m)(0)=0$ and
$\|q(U_n,U_m)\|_\infty\le 2/(\pi|n-m|)$.  The pole contribution is therefore
bounded by
\[
 \frac{C}{|n-m|}\int_0^L(e^{y/2}+e^{-y/2})\,dy
 \le C\frac{\lambda}{|n-m|}.
\]
For the nonarchimedean contribution, partial summation together with the
classical Chebyshev bound $\psi(x)\ll x$ gives
\[
 \sum_{1<k\le e^L}\frac{\Lambda(k)}{\sqrt{k}}
 \ll e^{L/2}=\lambda,
\]
and hence the same bound.

For the archimedean term set
\[
 \rho(y)=\frac{e^{y/2}}{e^y-e^{-y}},\qquad
 J_L(r)=\int_0^L\sin(2\pi r y/L)\rho(y)\,dy.
\]
Because the off-diagonal kernel vanishes at the origin, the CCM formula reduces
to
\[
 W_{\mathbb R}(V_n,V_m)
 =\frac{J_L(m)-J_L(n)}{\pi(n-m)}.
\]
Uniformly for $L\ge2\log2$ and integers $r$, $J_L(r)=O(1)$: on $0<y\le1$,
write $\rho(y)=(2y)^{-1}+O(1)$ and use the uniform boundedness of the sine
integral, while on $y\ge1$ use $\rho(y)\ll e^{-y/2}$.  Thus the archimedean
piece is $O(|n-m|^{-1})$.  Combining the three contributions proves
\eqref{eq:g1-high-low}.
\end{proof}

\begin{corollary}[Complete Fourier-matrix bound]\label{cor:g1-complete-matrix}
For all integers $n,m$ and $\lambda\ge2$,
\begin{equation}\label{eq:g1-complete-matrix}
 \boxed{
 |\tau^{(\lambda)}_{n,m}|
 \le
 \begin{cases}
 C\lambda/|n-m|,&n\ne m,\\[2mm]
 C\bigl(\lambda+1+\log(2+|n|/L)\bigr),&n=m.
 \end{cases}}
\end{equation}
The constant is absolute and uses no nontrivial-zero information.
\end{corollary}
\begin{proof}
The off-diagonal case is Lemma~\ref{lem:g1-high-low}.  The diagonal is a
separate CCM formula, not the limit of the off-diagonal divided difference:
for $0\le y\le L$,
\[
 q(U_n,U_n)(y)=2\left(1-\frac yL\right)
 \cos\frac{2\pi ny}{L},
 \qquad q(U_n,U_n)(0)=2.
\]
Thus the pole term is $O(\lambda)$ and, by the classical Chebyshev estimate,
the prime-power term is also $O(\lambda)$.

It remains to control the archimedean term.  Put $t_n=2\pi n/L$.  On
$0<x\le1$ the numerator in the logarithmic Weil kernel is, up to an absolute
constant,
\[
 \left(1-\frac xL\right)\cos(t_nx)-e^{-x/2}.
\]
Since $\rho(x)=e^{x/2}/(e^x-e^{-x})\ll x^{-1}$ there,
\[
 \left|\left(1-\frac xL\right)\cos(t_nx)-e^{-x/2}\right|
 \le C\{\min(1,t_n^2x^2)+x\}.
\]
Consequently
\[
 \int_0^1
 \left|\left(1-\frac xL\right)\cos(t_nx)-e^{-x/2}\right|
 \rho(x)\,dx
 \le C\bigl(1+\log(2+|t_n|)\bigr).
\]
For $x\ge1$ the kernel is $O(e^{-x/2})$ and the diagonal convolution is
bounded, so the remaining integral is $O(1)$.  Therefore the archimedean
contribution is
$O(1+\log(2+|n|/L))$, proving \eqref{eq:g1-complete-matrix}.
\end{proof}

\subsection*{A closed Weil operator and its omitted Fourier complement}
The entrywise bound above can be strengthened to an operator statement.
Here the closed operator is the canonical semilocal Weil operator; it is
not the arithmetic generator whose realization is conditional in Section~18.
Fix $\lambda>1$, put $L=2\log\lambda$, $t_n=2\pi n/L$, and write
\[
 P_\lambda=\sum_{1<m\le\lambda^2}\frac{\Lambda(m)}{\sqrt m},
 \qquad \rho(y)=\frac{e^{y/2}}{2\sinh y},
 \qquad h_L=\sinh^2(L/4).
\]
For later constants use the complete archimedean diagonal
\begin{align}
 A_n={}&\int_0^L\rho(y)
 \{2(1-y/L)\cos(t_ny)-2e^{-y/2}\}\,dy \notag\\
 &+\log(4\pi)+\gamma_{\rm E}+\log\tanh(L/2).
 \label{eq:v114-arch-diagonal}
\end{align}
The cancellation at zero makes this an ordinary convergent integral.

\begin{proposition}[Bounded commutator and the operator core]
\label{prop:v114-weil-core}
In the normalized logarithmic Fourier basis the canonical semilocal
Weil operator $\mathsf W_\lambda$ is
\begin{equation}\label{eq:v114-weil-decomposition}
 \mathsf W_\lambda=D_d+[M_b,H],\qquad
 H_{nm}=\begin{cases}(n-m)^{-1},&n\ne m,\\0,&n=m,\end{cases}
\end{equation}
where $D_d$ is real diagonal, $M_b$ is multiplication by the real odd
sequence
\begin{align}
 b_n={}&\frac{32Lh_Ln}{L^2+16\pi^2n^2}
 +\frac1\pi\left\{\int_0^L\rho(y)\sin(t_ny)\,dy
 +\sum_{1<m\le\lambda^2}\frac{\Lambda(m)}{\sqrt m}
                         \sin(t_n\log m)\right\},\notag\\
 d_n={}&\frac{32Lh_L(L^2-16\pi^2n^2)}{(L^2+16\pi^2n^2)^2}
 -2\sum_{1<m\le\lambda^2}\frac{\Lambda(m)}{\sqrt m}
        (1-\log m/L)\cos(t_n\log m)-A_n.
 \label{eq:v114-b-d}
\end{align}
In particular the diagonal is not filled in by a divided-difference limit.
One has
\begin{equation}\label{eq:v114-b-bound}
 \|b\|_\infty\le B_\lambda^{\rm mat}
 :=\frac{4h_L+P_\lambda+1+\pi/4}{\pi},\qquad
 \|[M_b,H]\|\le2\pi B_\lambda^{\rm mat}.
\end{equation}
At fixed $\lambda$, $d_n=\log|n|+O_\lambda(1)$ as $|n|\to\infty$.
Consequently
\begin{align}
 \mathcal D(\mathsf W_\lambda)
 &=\left\{x:\sum_n\log^2(2+|n|)|x_n|^2<\infty\right\},\notag\\
 \mathcal D[\mathsf W_\lambda]
 &=\left\{x:\sum_n\log(2+|n|)|x_n|^2<\infty\right\}.
 \label{eq:v114-weil-domains}
\end{align}
Finite Fourier sums form an operator core. The operator is lower bounded
and has compact resolvent.
\end{proposition}
\begin{proof}
The exact CCM entries give $\tau_{nm}=(b_n-b_m)/(n-m)$ for $n\ne m$
and $\tau_{nn}=d_n$; see also \eqref{eq:beta-bcoeff-direct}.
For $n>0$, expansion $\rho(y)=\sum_{k\ge0}e^{-(2k+1/2)y}$ and
$t_nL=2\pi n$ give
\[
 S_n:=\int_0^L\rho(y)\sin(t_ny)\,dy
 =\sum_{k\ge0}\frac{(1-e^{-(2k+1/2)L})t_n}
                         {(2k+1/2)^2+t_n^2}.
\]
Interchange is justified by integrability of
$\rho(y)|\sin(t_ny)|$. The summands are nonnegative. The $k=0$
term is at most $1$. For $k\ge1$ discard the factor
$1-e^{-(2k+1/2)L}\le1$ and compare the decreasing majorant
$t_n/\{(2k+1/2)^2+t_n^2\}$ with an integral; this gives at most
$\pi/4$. Thus
$0\le S_n\le1+\pi/4$. Oddness treats negative $n$, and $S_0=0$.
The pole part of $b_n$ has absolute value at most $4h_L/\pi$;
the prime sine sum is bounded by $P_\lambda$. This proves the first
bound in \eqref{eq:v114-b-bound}.

On $\ell^2(\mathbb Z)$, $H$ is the discrete Hilbert transform.
The Fourier series $\sum_{k\ne0}e^{ik\theta}/k=i(\pi-\theta)$
for $0<\theta<2\pi$ proves $H^*=-H$ and $\|H\|=\pi$.
Hence $[M_b,H]$ is bounded self-adjoint and has the asserted norm.
Moreover
\[
 A_n=A_0-\int_0^L w_L(y)(1-\cos(t_ny))\,dy,
 \qquad w_L(y)=2\rho(y)(1-y/L).
\]
Since $w_L(y)-1/y$ is integrable on $(0,L)$, the last integral is
$\log|t_n|+O_L(1)$. For example, split
$\int_0^{|t_n|L}(1-\cos v)\,dv/v$ at $1$ and integrate its
oscillatory tail by parts. The other terms in $d_n$ are bounded at
fixed $\lambda$, proving its logarithmic asymptotic.

The diagonal operator is self-adjoint on the first domain in
\eqref{eq:v114-weil-domains}, has finite sequences as an operator
core, and has compact resolvent. Bounded self-adjoint perturbation
preserves these statements and gives the displayed form domain.
This identifies the actual physical Weil realization: the finite sums
are a form core for its lower-bounded closed form by
\cite[Propositions~3.2--3.4]{ConnesConsaniMoscovici2025}, and both
forms have exactly the same entries there. Uniqueness of the form
closure proves equality, without imposing a new boundary condition.
\end{proof}

Every periodic BV proxy, including the repaired $p_\lambda$, therefore
belongs to $\mathcal D(\mathsf W_\lambda)$ at fixed $\lambda$:
its Fourier coefficients are $O_\lambda((1+|n|)^{-1})$.
Furthermore $P_Np_\lambda\to p_\lambda$ in this operator graph norm,
by the diagonal description and boundedness of the commutator. This
fixed-window statement supplies no endpoint-normalized uniform rate,
and does not remove the separate sampler graph defect.

\begin{proposition}[An explicit positive high-frequency tail]
\label{prop:v114-weil-tail}
Set $a=\min(1,L)$ and
\begin{align}
 C_\lambda^{\rm tail}
 :={}&|A_0|+4+\pi/2+a+2a/L+2P_\lambda
                   +2\sinh(L/2)-L,\notag\\
 \Gamma_{\lambda,N}
 :={}&\log\frac{2\pi(N+1)a}{L}-C_\lambda^{\rm tail}.
 \label{eq:v114-tail-constant}
\end{align}
For every vector $f$ in the form domain supported on Fourier indices
$|n|>N$,
\begin{equation}\label{eq:v114-tail-positive}
 QW_\lambda(f,f)\ge\Gamma_{\lambda,N}\|f\|_2^2.
\end{equation}
In particular this omitted complement is strictly positive whenever
$\Gamma_{\lambda,N}>0$. The constants use no nontrivial-zero data.
\end{proposition}
\begin{proof}
On $0<y<a$, $2\rho(y)\ge e^{-y/2}/y\ge1/y-1/2$, whence
\[
 w_L(y)\ge1/y-(1/2+1/L).
\]
Since $w_L\ge0$ on the full interval and
$F(T):=\int_0^T(1-\cos v)\,dv/v\ge\log T-2$ for $T>0$,
\[
 -A_n\ge\log(|t_n|a)-|A_0|-2-a-2a/L\qquad(n\ne0).
\]
The bound for $F$ is immediate for $T\le1$; for $T\ge1$ the
integral of $\cos v/v$ on $[1,T]$ has absolute value at most $2$
by integration by parts. The off-diagonal archimedean operator has
norm at most $2+\pi/2$, using $b_n^{\rm arch}=S_n/\pi$ in the
preceding commutator proof.

The complete prime operator is a sum of negative weighted truncated
translations and their adjoints, so its lower bound is $-2P_\lambda I$.
In centered logarithmic coordinates $x\in[-L/2,L/2]$, the pole form is
\[
 2|\ip f{\cosh(x/2)}|^2-2|\ip f{\sinh(x/2)}|^2
 \ge-\{2\sinh(L/2)-L\}\|f\|_2^2.
\]
Combining these three estimates gives
\eqref{eq:v114-tail-positive} first for finite Fourier sums and then
on the tail form domain by closure.
\end{proof}

\subsection*{A sharper positive Fourier complement at the certified window}
The previous bound loses constants by bounding each place on the full
space. On the omitted Fourier tail one can retain the limiting
archimedean commutator, the physical incidence of prime shifts, and the
decay of the negative pole vector. These are bounds for the same Weil
operator and its original Fourier projection.

For $L=2\log\lambda$ put
\begin{align}
 R_L&=\frac{e^{-L/2}}{1-e^{-2L}},&
 C_L&=\max\{1,1/4+R_L\},\notag\\
 E_L(t)&=\frac1{15t}+\frac7{2t^2}
       +\frac\pi{2Lt}+\frac{2+2R_L}{Lt^2},\qquad t>0.
 \label{eq:v115-arch-errors}
\end{align}
Define the weighted prime degree
\begin{equation}\label{eq:v115-prime-degree}
 M_\lambda=\mathop{\rm ess\,sup}_{0<x<L}
 \left\{\sum_{\substack{1<m<\lambda^2\\\log m\le x}}
                    \frac{\Lambda(m)}{\sqrt m}
       +\sum_{\substack{1<m<\lambda^2\\\log m\le L-x}}
                    \frac{\Lambda(m)}{\sqrt m}\right\}.
\end{equation}
The strict upper cut omits only a possible full-length translation,
which is zero on $L^2([0,L])$.

\begin{proposition}[A practical positive omitted complement]
\label{prop:v115-sharp-tail}
For $N\ge1$ let $t_*=2\pi(N+1)/L$, and set
\begin{align}
 \Gamma_{\lambda,N}^{\rm even}
   &=\log\frac{N+1}{L}-M_\lambda
                -\frac{2C_L}{t_*}-E_L(t_*),\notag\\
 \Gamma_{\lambda,N}^{\rm all}
   &=\Gamma_{\lambda,N}^{\rm even}-\frac\pi2
                   -\frac{4\sinh^2(L/4)L}{\pi^2N}.
 \label{eq:v115-sharp-tail}
\end{align}
For every form-domain vector $f$ whose Fourier coefficients vanish on
$|n|\le N$,
\[
 QW_\lambda(f,f)\ge\Gamma_{\lambda,N}^{\rm all}\|f\|_2^2.
\]
For an inversion-even $f$, the stronger lower bound with
$\Gamma_{\lambda,N}^{\rm even}$ holds. In particular, at the actual
window $\lambda=3$,
\begin{equation}\label{eq:v115-lambda3-tail}
 \Gamma_{3,256}^{\rm all}>0.0148,\qquad
 \Gamma_{3,256}^{\rm even}>1.5867,\qquad
 \Gamma_{3,512}^{\rm all}>0.7085.
\end{equation}
These inequalities control the entire omitted complement, not merely a
finite further block.
\end{proposition}
\begin{proof}
Write $a_k=2k+1/2$ and use the exact Fourier-grid series for $S_n$
in the proof of Proposition~\ref{prop:v114-weil-core}.
For $t>0$, the function $x\mapsto t/(x^2+t^2)$ decreases on $x>0$.
Comparing its values on $1/2+2\mathbb N$ with its integral gives
\[
 \frac\pi4-\frac1{4t}
 \le\sum_{k\ge0}\frac{t}{a_k^2+t^2}
 \le\frac\pi4+\frac1t.
\]
The exponentially weighted part of $S_n$ is nonnegative and at most
$R_L/|t_n|$. Oddness therefore proves
\begin{equation}\label{eq:v115-S-sharp}
 \left|S_n-\frac\pi4\operatorname{sgn}n\right|
       \le\frac{C_L}{|t_n|}\qquad(n\ne0).
\end{equation}
On the tail, the archimedean off-diagonal multiplier is consequently
$b_n^{\rm arch}=\operatorname{sgn}(n)/4+e_n$ with
$\|e\|_\infty\le C_L/(\pi t_*)$.
Since the compressed discrete Hilbert transform has norm at most $\pi$,
its commutator has norm at most $\pi/2+2C_L/t_*$.
For even coefficients the limiting commutator is instead positive:
under the normalized identification with positive indices its matrix is
\[
 K_{nm}=\frac1{2(n+m)},\qquad n,m>N.
\]
Positivity follows from
$1/(n+m)=\int_0^1x^{n+m-1}\,dx$. Thus only $2C_L/t_*$ is lost in
the even sector. This uses the unshifted Fourier basis of
\eqref{eq:v114-weil-decomposition}; translation to centered physical
coordinates conjugates by $(-1)^n$ and preserves both parity and the
positive quadratic form.

Next retain the complete archimedean diagonal. Its exact series, already
used in the independent matrix certificate, gives for
$z=1/4+it/2$ and $t=|t_n|>0$,
\begin{equation}\label{eq:v115-arch-exact}
 -A_n=\Re\psi_0(z)-\log\pi+\frac{\Re\psi_1(z)}{2L}
 -\frac2L\sum_{k\ge0}e^{-a_kL}
                   \frac{a_k^2-t^2}{(a_k^2+t^2)^2}.
\end{equation}
Here $\psi_0$ and $\psi_1$ are digamma and trigamma.
Binet's exact integral formula~\cite[eq.~5.9.15]{NISTDigamma} at
$w=z+1=5/4+it/2$ has remainder of modulus at most
\[
 2\int_0^\infty\frac{u\,du}
 {|u^2+w^2|(e^{2\pi u}-1)}
 \le\frac8{5t}\int_0^\infty\frac{u\,du}{e^{2\pi u}-1}
 =\frac1{15t},
\]
since $|u^2+w^2|\ge5t/4$. The recurrence
$\psi_0(z)=\psi_0(z+1)-1/z$, together with
$\Re\log w\ge\log(t/2)$,
$\Re(1/(2w))\le5/(2t^2)$, and $\Re(1/z)\le1/t^2$, yields
\[
 \Re\psi_0(z)\ge\log(t/2)-\frac1{15t}-\frac7{2t^2}.
\]
The convergent trigamma series and a decreasing-integrand comparison give
\[
 |\psi_1(z)|\le\sum_{k\ge0}\frac1{(k+1/4)^2+t^2/4}
                 \le\frac4{t^2}+\frac\pi t.
\]
The last term in \eqref{eq:v115-arch-exact} has modulus at most
$2R_L/(Lt^2)$. Hence
\begin{equation}\label{eq:v115-diagonal-sharp}
 -A_n\ge\log\frac{|n|}{L}-E_L(|t_n|).
\end{equation}
Its right side increases with $|n|$, so $|n|=N+1$ supplies the uniform
tail bound. No logarithmic constant or finite-$L$ diagonal term was
discarded from the exact identity.

For the prime part, set $w_m=\Lambda(m)/\sqrt m$ and $a_m=\log m$.
The elementary inequality $2\Re(z\overline w)\le|z|^2+|w|^2$ gives
\[
 -2\sum_m w_m\Re\int_0^{L-a_m}
      f(x+a_m)\overline{f(x)}\,dx
 \ge-\int_0^L d_\lambda(x)|f(x)|^2dx
 \ge-M_\lambda\|f\|_2^2,
\]
where $d_\lambda$ is the function inside braces in
\eqref{eq:v115-prime-degree}. This applies equally to complex vectors.

In centered physical coordinates the negative pole term is
$-2|\ip f{\sinh(x/2)}|^2$. Direct Fourier integration gives
\[
 |\widehat{\sinh(x/2)}(n)|^2
 =\frac{4\sinh^2(L/4)t_n^2}{L(t_n^2+1/4)^2}.
\]
Consequently
\[
 2\|Q_N\sinh(x/2)\|_2^2
 \le\frac{4\sinh^2(L/4)L}{\pi^2}
                         \sum_{n>N}n^{-2}
 \le\frac{4\sinh^2(L/4)L}{\pi^2N}.
\]
It vanishes identically on the even sector; the positive pole term can
be retained or omitted in a lower bound. Combining these estimates
proves \eqref{eq:v115-sharp-tail} first on finite Fourier sums and then
on the tail form domain by Proposition~\ref{prop:v114-weil-core}.

For the numerical constants at $\lambda=3$, symmetry reduces the prime
degree to $0<x<\log3$. Its breakpoints, expressed in $e^x$, are
$9/8,9/7,9/5,2,9/4$. With
$P=w_2+w_3+w_4+w_5+w_7+w_8$, the six interval values are
\[
 P,\quad P-w_8,\quad P-w_8-w_7,\quad
 w_2+w_3+w_4,\quad2w_2+w_3+w_4,\quad2w_2+w_3.
\]
All are at most $P$, using $w_5>w_2$ for the fifth value, and the first
occurs on an interval of positive length. Thus $M_3=P$ exactly.
Also $\sinh^2(L/4)=1/3$, $R_L=27/80$ and $C_L=1$.
Outward interval evaluation of the displayed elementary expressions
gives \eqref{eq:v115-lambda3-tail}; the reproducibility bundle contains
the 256-bit scalar certificate. These are not new finite-matrix
eigenvalue computations.
\end{proof}

\subsection*{A positive weight for the prime shifts and a diagonal inverse bound}
The flat degree bound can be sharpened without changing the Weil operator.
Let $\mathsf T_{\rm pr}$ denote the sum of the truncated translations and
their adjoints with weights $w_m=\Lambda(m)/\sqrt m$. Its coefficients
are nonnegative, and the prime contribution to the Weil form is
$-\mathsf T_{\rm pr}$. For every bounded function $\phi$ bounded away
from zero and every complex $f$,
\begin{equation}\label{eq:v116-weighted-prime-general}
 |\ip{\mathsf T_{\rm pr}f}{f}|
 \le\int_0^L\frac{(\mathsf T_{\rm pr}\phi)(x)}{\phi(x)}|f(x)|^2\,dx
 \le M_\phi\norm f_2^2,
 \qquad
 M_\phi:=\mathop{\rm ess\,sup}_{0<x<L}
       \frac{(\mathsf T_{\rm pr}\phi)(x)}{\phi(x)}.
\end{equation}
Here $\phi$ is required to be strictly positive. This follows edge by
edge from
\[
 2|f(x+s)f(x)|\le
 \frac{\phi(x)}{\phi(x+s)}|f(x+s)|^2
 +\frac{\phi(x+s)}{\phi(x)}|f(x)|^2.
\]
Thus the argument is a weighted Schur bound on the translation sum;
it assumes no sign for the complete Weil form.

\begin{proposition}[A certified weighted prime bound at $\lambda=3$]
\label{prop:v116-weighted-prime}
At $L=2\log3$, take $\phi(x)=\cosh(x-\log3)$. With
$w_m=\Lambda(m)/\sqrt m$, put
\begin{equation}\label{eq:v116-prime-constant}
 \mathfrak m_3=
 \frac{85}{116}w_2+\frac{65}{87}w_3+
 \frac{193}{232}w_4+\frac{137}{145}w_5+
 \frac{35}{29}w_7.
\end{equation}
Then $M_\phi=\mathfrak m_3$ and
\[
 2.6890557719319011796<\mathfrak m_3
       <2.6890557719319011797.
\]
In particular $\norm{\mathsf T_{\rm pr}}\le\mathfrak m_3$.
Replacing $M_3$ by $\mathfrak m_3$ in
Proposition~\ref{prop:v115-sharp-tail} gives
\begin{equation}\label{eq:v116-weighted-tail-values}
 \widetilde\Gamma_{3,256}^{\rm all}>0.4970,
 \qquad \widetilde\Gamma_{3,256}^{\rm even}>2.0690,
 \qquad \widetilde\Gamma_{3,512}^{\rm all}>1.1907.
\end{equation}
These are omitted-complement bounds, with no assertion yet about the
sign of the head/complement Schur matrix.
\end{proposition}
\begin{proof}
Write $u=e^x\in[1,9]$. The ratio for a forward shift by $\log m$ is
\[
 \frac{m u^2+9/m}{u^2+9},
\]
and for its backward shift it is $(u^2/m+9m)/(u^2+9)$.
On an interval with fixed active shifts, their weighted sum has the
form $A+B\tanh(x-\log3)$ and is monotone or constant. The endpoints
are $u=1,9,m,9/m$, for $m=2,3,4,5,7,8$.
By reflection it suffices to examine the six intervals in $1<u<3$,
whose successive upper endpoints are
\[
 9/8,\quad9/7,\quad9/5,\quad2,\quad9/4,\quad3.
\]
On each such interval the ratio is increasing. Its largest one-sided
value is approached from $u<9/7$ and is
\[
 \sum_{m=2,3,4,5,7}
       w_m\frac{9m^2+49}{58m}=\mathfrak m_3.
\]
The finite comparisons with the other endpoint values use rational
coefficients and are certified at 256-bit Arb precision by
\texttt{certify\_weighted\_prime\_tail.py}; its exact coefficient lists
and interval enclosures are included in the reproducibility bundle.
The $m=8$ term is retained on every interval where that shift is active.
The $m=9$ translation has full interval length and is zero in $L^2$.
This proves the asserted essential supremum and norm bound by
\eqref{eq:v116-weighted-prime-general}. Substitution in the explicit
formulas of Proposition~\ref{prop:v115-sharp-tail}, with
$h_L=1/3$, $R_L=27/80$ and $C_L=1$, gives
\eqref{eq:v116-weighted-tail-values}; the same scalar interval
certificate verifies the strict inequalities.
\end{proof}

\begin{proposition}[Diagonal domination of the tail inverse]
\label{prop:v116-diagonal-inverse}
Continue at $\lambda=3$ and retain $E_L$ and $C_L$ from
\eqref{eq:v115-arch-errors}. Fix $N\ge1$ and set
$t_*=2\pi(N+1)/L$. For the even sector define
\[
 \kappa_{\rm even}=\mathfrak m_3+2C_L/t_*;
\]
for the whole space, and hence also the odd sector, define
\[
 \kappa_{\rm all}=\kappa_{\rm even}+\pi/2
                         +4h_LL/(\pi^2N).
\]
In the selected sector put
\begin{equation}\label{eq:v116-diagonal-weight}
 g_n=\log\frac{|n|}{L}-E_L(2\pi|n|/L)-\kappa,
            \qquad |n|>N.
\end{equation}
If $g_{N+1}>0$, the canonical omitted-tail operator $T$ satisfies
the form inequality $T\succeq D_g$, where $D_g$ is this diagonal
operator, and therefore
\begin{equation}\label{eq:v116-inverse-order}
 T^{-1}\preceq D_g^{-1}.
\end{equation}
The same weights apply in the normalized positive-index parity basis.
\end{proposition}
\begin{proof}
The proof of Proposition~\ref{prop:v115-sharp-tail} bounds the
archimedean diagonal at each index by
$\log(|n|/L)-E_L(2\pi|n|/L)$ before any infimum is taken.
The remaining archimedean error has norm at most $2C_L/t_*$.
Its limiting commutator is nonnegative in the even sector and is
bounded below by $-\pi/2$ on the whole tail. Apply
Proposition~\ref{prop:v116-weighted-prime} to the prime term and the
existing negative-pole tail estimate to the whole-space case.
These bounds prove $T\succeq D_g$ on finite tail sums and hence on
the common logarithmic form domain. The weights $g_n$ increase with
$|n|$, since $E_L(t)$ is decreasing, so $D_g\succeq g_{N+1}I>0$.
The inverse inequality follows, for example, from the variational
identity
\[
 \ip{T^{-1}v}{v}
 =\sup_w\{2\Re\ip vw-\mathfrak t[w]\}
 \le\sup_w\left\{2\Re\ip vw-\sum_n g_n|w_n|^2\right\}
 =\sum_n\frac{|v_n|^2}{g_n}.
\]
The quadratic quantities are real in the fixed first-slot-linear
convention. Restriction to a parity sector is isometric and leaves
the diagonal weights unchanged.
\end{proof}


The preceding tail estimates alone do not establish full Weil positivity.
The complete signed certificate below also controls the effective head.
Its absolute arithmetic bound still gives a sufficient cutoff
exponential in $O(\lambda)$; it does not reach the natural prolate
scale. For the signed correction use $Q_N=I-P_N$ and let
$\Gamma_{\lambda,N}>0$ denote any of the proved tail bounds. Write
\[
 \mathsf W_\lambda=
 \begin{pmatrix}F&B^*\\B&T\end{pmatrix},\qquad
 F=P_N\mathsf W_\lambda P_N,\quad
 B=Q_N\mathsf W_\lambda P_N,\quad T\succeq\Gamma_{\lambda,N}I.
\]
The cross operator $B$ is bounded by \eqref{eq:v114-weil-decomposition}.
The remaining exact sign problem is the finite matrix
$K=F-B^*T^{-1}B$. There is a useful verified-solve alternative to
bounding $B$ alone. For any $Z:E_N\to\mathcal D(T)$ put
$R=B-TZ$ and $K_Z=F-B^*Z-Z^*B+Z^*TZ$. Then
\begin{equation}\label{eq:v114-schur-residual}
 K=K_Z-R^*T^{-1}R
 \succeq K_Z-\Gamma_{\lambda,N}^{-1}R^*R
 \succeq K_Z-\frac{\|R\|^2}{\Gamma_{\lambda,N}}I.
\end{equation}
Expansion proves the identity and $T^{-1}\preceq\Gamma_{\lambda,N}^{-1}I$
proves both bounds. Retaining the residual Gram matrix $R^*R$ preserves
directional cancellation lost by its scalar-norm majorant.
If the final finite lower bound is
$-\varepsilon I$, $\varepsilon\ge0$, completing the form square
gives $\mathsf W_\lambda\succeq-\varepsilon I$.
Proving such a signed bound along growing windows is still open;
positivity of $F$ alone is insufficient.

\subsection*{A directional certificate for the entire Fourier complement}

The residual Gram in \eqref{eq:v114-schur-residual} can be enclosed
without replacing its small directions by a scalar norm.  We retain
the actual matrix \eqref{eq:v114-b-d}, including its separate diagonal.
All adjoints below are the usual coordinate adjoints in the fixed
orthonormal Fourier basis; quadratic-form order is unchanged by our
first-slot-linear convention.

\begin{proposition}[Moment enclosure of the remote residual]
\label{prop:v116-remote-moment}
Fix $\lambda$, integers $1\le M<J$, and $r\ge1$.  Let
$G:\mathbb C^d\to E_M$ be any coefficient matrix, with row $G_m$ at
Fourier index $m$, and choose $B_*\ge\sup_n|b_n|$, for example
$B_*=B_\lambda^{\rm mat}$ from \eqref{eq:v114-b-bound}.
Define the $r$-by-$d$ moment matrices by
\begin{equation}\label{eq:v116-moments}
 S_j=\sum_{m=-M}^M(m/M)^jG_m,\qquad
 T_j=\sum_{m=-M}^M(m/M)^j b_mG_m,
 \qquad 0\le j<r,
\end{equation}
using $0^0=1$.  Set
\begin{align}
 H_{jk}&=M^{j+k}\{1+(-1)^{j+k}\}
                  \zeta(j+k+2,J+1),\notag\\
 c_{M,r}&=\sum_{m=-M}^M(|m|/M)^{2r},\notag\\
 \delta_{M,J,r}^{2}
 &=\frac{8B_*^2c_{M,r}M^{2r}\zeta(2r+2,J+1)}
                    {(1-M/(J+1))^2},
 \label{eq:v116-moment-constants}
\end{align}
where $\zeta(s,a)=\sum_{k\ge0}(k+a)^{-s}$ is the Hurwitz zeta
function.  For every $\tau>0$, the explicitly finite matrix
\begin{equation}\label{eq:v116-remote-bound}
 U_J=(1+\tau)\{2B_*^2S^*HS+2T^*HT\}
       +(1+\tau^{-1})\delta_{M,J,r}^{2}G^*G
\end{equation}
satisfies
\[
 (Q_J\mathsf W_\lambda G)^*(Q_J\mathsf W_\lambda G)\preceq U_J.
\]
The assertion also holds on either normalized parity sector, provided
the moments are those of the corresponding full Fourier coefficients.
\end{proposition}
\begin{proof}
Since $|n|>J>M\ge|m|$, no diagonal entry occurs.  The finite geometric
identity for $(n-m)^{-1}$ gives exactly
\begin{align*}
 (\mathsf W_\lambda G)_n&=A_n+E_n,\\
 A_n&=\sum_{j=0}^{r-1}\frac{M^j}{n^{j+1}}(b_nS_j-T_j),\\
 E_n&=\sum_{m=-M}^M\left(\frac mn\right)^r
                         \frac{b_n-b_m}{n-m}G_m.
\end{align*}
The matrix $H$ is the Gram matrix of the sequences
$M^j n^{-j-1}$ on $|n|>J$; pairing the two signs of $n$ proves the
formula in \eqref{eq:v116-moment-constants} and $H\succeq0$.
For a column $x$, apply
$|b_na_n-t_n|^2\le2B_*^2|a_n|^2+2|t_n|^2$ and sum over the remote
indices.  This bounds the Gram of $A$ by
$2B_*^2S^*HS+2T^*HT$.

For the remainder, Cauchy--Schwarz in $m$ yields
\[
 |E_nx|^2\le
 \frac{4B_*^2(M/|n|)^{2r}c_{M,r}}{(|n|-M)^2}\|Gx\|^2.
\]
Use $|n|-M\ge|n|(1-M/(J+1))$ and sum over both tails.
The resulting bound is
$\sum_{|n|>J}|E_nx|^2\le\delta_{M,J,r}^{2}\|Gx\|^2$;
this accounts for the factor $8$ in
\eqref{eq:v116-moment-constants}.  Finally apply
$|a+e|^2\le(1+\tau)|a|^2+(1+\tau^{-1})|e|^2$ before summing.
Every estimate is a quadratic-form inequality for arbitrary $x$, so
it proves the stated matrix order, not merely entrywise bounds.

For completeness, an even normalized parity vector has full
coordinates $G_0=v_0$, $G_m=G_{-m}=v_m/\sqrt2$ for $m>0$;
an odd vector has $G_0=0$, $G_m=-G_{-m}=v_m/\sqrt2$.
As $b$ is odd, $S_j=0$ for odd $j$ and $T_j=0$ for even $j$ in the
even sector, with the opposite zeros in the odd sector.
These zeros are exact.  The two tails are already counted in $H$;
no further factor of two is inserted in the parity calculation.
\end{proof}

\begin{proposition}[Diagonal-weighted Schur certificate]
\label{prop:v116-directional-schur}
Write $P=P_N$, $Q=I-P$ and
$F=P\mathsf W_\lambda P$, $B=Q\mathsf W_\lambda P$,
$T=Q\mathsf W_\lambda Q$.  Suppose an independent form estimate gives
\[
 T\succeq D_g:=\operatorname{diag}_{|n|>N}(g_{|n|})
       \succeq\gamma I>0,
\]
where $g_k$ is nondecreasing for $k>N$; the estimate of
Proposition~\ref{prop:v116-diagonal-inverse} is one source of such a
sequence.  Let $N<M<J$, let $Z:E_N\to E_M\cap Q\mathcal H$ be any
finite matrix, and put
\[
 G=I_{E_N\hookrightarrow E_M}-Z,\qquad
 K_Z=F-B^*Z-Z^*B+Z^*TZ.
\]
For $N<|n|\le J$ compute the rows
$R_n=(\mathsf W_\lambda G)_n$ using the complete actual matrix,
and form
\begin{equation}\label{eq:v116-directional-lower}
 \mathcal L_Z=K_Z-
     \sum_{N<|n|\le J}\frac{R_n^*R_n}{g_{|n|}}
                    -\frac{U_J}{g_{J+1}},
\end{equation}
with $U_J$ from Proposition~\ref{prop:v116-remote-moment}.
Then the exact Schur complement satisfies
$F-B^*T^{-1}B\succeq\mathcal L_Z$.
In particular, a certified inequality
$\mathcal L_Z\succeq-\varepsilon I$, $\varepsilon\ge0$, implies
$\mathsf W_\lambda\succeq-\varepsilon I$ on the entire form domain.
The same implication holds after restriction to either parity sector.
\end{proposition}
\begin{proof}
The cross operator $B$ is bounded by
Proposition~\ref{prop:v114-weil-core}; finite $Z$ has range in the
operator domain.  Thus $R=B-TZ=Q\mathsf W_\lambda G$ is well-defined.
Inverse order gives $T^{-1}\preceq D_g^{-1}$.  For unbounded
positive operators this follows, for example, by comparing the
variational formulas for their inverse quadratic forms; the positive
lower bound $\gamma$ makes both inverses bounded.  The exact
verified-solve identity therefore gives
\[
 F-B^*T^{-1}B=K_Z-R^*T^{-1}R
                  \succeq K_Z-R^*D_g^{-1}R.
\]
Split the last quadratic form into $N<|n|\le J$ and $|n|>J$.
Monotonicity gives $g_{|n|}^{-1}\le g_{J+1}^{-1}$ on the remote part,
where \eqref{eq:v116-remote-bound} applies.  This proves
\eqref{eq:v116-directional-lower}.  Writing $\mathfrak w_\lambda$ for the closed form of
$\mathsf W_\lambda$, completion of the square gives, for head $x$
and form-domain tail $y$,
\[
 \mathfrak w_\lambda[x+y]
 =\|T^{1/2}(y+T^{-1}Bx)\|^2
       +\langle(F-B^*T^{-1}B)x,x\rangle.
\]
Interpreted as the closed form on its form domain, this is at least
$-\varepsilon\|x\|^2\ge-\varepsilon\|x+y\|^2$.
For a calculation confined to one parity sector, all ingredients
are taken in that sector and the same proof applies.
\end{proof}

The use of $g_{|n|}$ on the intervening rows, and of the full moment
Grams on the remote rows, is substantive: replacing either by a
single operator norm can destroy the small head directions that the
certificate must preserve.  The finite rows in
\eqref{eq:v116-directional-lower} include the true diagonal $d_n$;
only the disjoint remote rows use the divided-difference expansion.
These propositions supply a finite sufficient test, not a
uniform-in-$\lambda$ positivity theorem.  A passed test at one window
does not close G2 or justify any transfer without its existing graph
defect.

\subsection*{A certificate for the complete fixed-window form}
The preceding estimates now permit a signed certificate that includes
the entire Fourier complement. This is a computer-assisted result for
one physical window; its finite arithmetic and its analytic tail bounds
are both part of the proof.

\begin{proposition}[Complete positivity at $\lambda=3$]
\label{prop:v116-window-positive}
For the canonical closed Weil form of
Proposition~\ref{prop:v114-weil-core}, at $\lambda=3$ there is a
constant $\epsilon_3>0$ such that
\[
 QW_3(f,f)\ge\epsilon_3\|f\|_2^2
\]
for every vector in its form domain. The result includes both parity
sectors and all complex vectors. No numerical value of $\epsilon_3$
is asserted here. In particular, the LDL pivots below are not lower
bounds for the smallest eigenvalue.
\end{proposition}
\begin{proof}[Analytic reduction and interval certificate]
Use the actual unshifted orthonormal Fourier basis on $[0,L]$,
$L=2\log3$, and its orthonormal parity combinations
$e_0$, $(e_m+e_{-m})/\sqrt2$ and
$(e_m-e_{-m})/\sqrt2$. The real matrix commutes with reflection, so
positivity in both sectors implies positivity for the whole complex
space. For each sector split the head at $N$, use an intermediate
inverse-action support through $M$, and compute residual rows through
$J$, with the following exact parameters:
\begin{center}
\begin{tabular}{lrrrrr}
\toprule
Sector & Head dimension & $N$ & $M$ & $J$ & Moment order $r$\\
\midrule
Even &257&256&512&4096&80\\
Odd  &512&512&1024&4096&100\\
\bottomrule
\end{tabular}
\end{center}
The tail lower functions $g_n$ are those of
Proposition~\ref{prop:v116-diagonal-inverse}, with the common base cut $N$
retained in every loss term. In particular their smallest values exceed
$2.0690$ in the even case and $1.1907$ in the odd case.

Let $T_M$ be the compression to $N<|n|\le M$, and let $B_M$ be the
corresponding head-to-tail block. A 768-bit Arb interval solve encloses
$T_M^{-1}B_M$; freeze its midpoint as the exact dyadic matrix $Z$.
The certificate then uses the complete expression
\[
 K_Z=F-B_M^*Z-Z^*B_M+Z^*T_MZ,
\]
and computes every residual row $N<|n|\le J$ from
$R_n=(\mathsf W_3G)_n$, where $G=I_{\rm head}-Z$.
This retains the nonzero intermediate residual from freezing $Z$.
The correct separate logarithmic diagonal is used whenever $n=m$.
For the remote rows, use
Proposition~\ref{prop:v116-remote-moment} with
\[
 B_* =\frac{4/3+P_3+1+\pi/4}{\pi},\qquad
 P_3=\sum_{1<m\le9}\frac{\Lambda(m)}{\sqrt m},\qquad
 \tau=10^{-6}.
\]
The optional full-length term in this conservative $B_*$ does not
enter the actual matrix. It only enlarges the absolute sequence bound.
The squared geometric remainder factors satisfy
\[
 \delta^2_{\rm even}<1.656\cdot10^{-148},\qquad
 \delta^2_{\rm odd}<3.288\cdot10^{-124}.
\]
The full matrix majorant $U_{\rm remote}$, including its leading
moment terms, is retained; these small remainder factors alone are not
estimates for the whole remote residual.

Proposition~\ref{prop:v116-directional-schur} gives the finite lower form
\begin{equation}\label{eq:v116-certified-lower}
 \mathcal L=K_Z-
 \sum_{N<|n|\le J}\frac{R_n^*R_n}{g_n}
 -\frac{U_{\rm remote}}{g_{J+1}}
 \preceq F-B^*T^{-1}B.
\end{equation}
In a parity basis the sum uses one normalized row per positive index;
the remote formula is evaluated in the isometric full-index embedding.
There is no additional factor two or omitted opposite-sign tail.
Interval LDL decomposition certifies every pivot of $\mathcal L$ to
be strictly positive. The smallest pivot enclosures are contained in
\[
 (9.17498,9.17499)\cdot10^{-9}\quad\hbox{(257 even pivots)},
 \qquad
 (9.68026,9.68027)\cdot10^{-8}\quad\hbox{(512 odd pivots)}.
\]
The reproducibility bundle records every pivot, the exact dyadic
witnesses, their hashes, the matrix assembly, and all analytic error
constants. Replaying both identical witnesses at 896-bit precision
again certifies all pivots positive.

For completeness, the coefficient evaluation does not truncate an
uncontrolled archimedean series. With $a_k=2k+1/2$, the exact
digamma and trigamma expressions from
\eqref{eq:v115-arch-exact} leave exponentially weighted sums.
After $k=0,\ldots,127$, put $a_{128}=256+1/2$ and
$e_{128}=3^{-513}$. The radii added to the sine sequence and
complete archimedean diagonal respectively are
\[
 \frac{e_{128}}{2a_{128}(1-1/81)},\qquad
 \frac{2e_{128}}{La_{128}^2(1-1/81)}.
\]
They follow from $t/(a^2+t^2)\le1/(2a)$ and
$|a^2-t^2|/(a^2+t^2)^2\le1/a^2$, uniformly in frequency.
The full-length $m=9$ prime translation is exactly zero; every
shorter prime-power shift is present. All remaining function evaluations
and matrix operations use outward ball enclosures.

Thus the exact Schur complement is positive definite in each finite
head and $T\succeq g_{N+1}I>0$ on its entire complement. Completing
the form square proves coercivity, since the triangular map
$(x,y)\mapsto(x,y+T^{-1}Bx)$ is boundedly invertible and the finite
positive Schur matrix has a positive smallest eigenvalue. The established
Fourier form core extends the result to the complete closed form domain.
No zero-location or Weil-positivity assumption is used.
\end{proof}

\subsection*{Fixed-window spectral ordering and ordinary ground overlap}

The fixed-window computer-assisted positivity certificate can be shifted to count its low
eigenvalues.  The thresholds in the two reflection sectors need not
coincide.  In particular, an even trial vector only sees the even spectral
gap when its overlap with the ground state is estimated.

\begin{proposition}[A fixed-window simple even ground state]
\label{prop:v117-ground-order}
Let $\mu_0\le\mu_1\le\cdots$ be the eigenvalues, counted with
multiplicity, of the complete canonical operator $\mathsf W_3$.
Then its ground state is simple and even, and
\begin{equation}\label{eq:v117-ground-order}
 0<\mu_0<3.644\cdot10^{-38},\qquad \mu_1>10^{-36}.
\end{equation}
Within the even sector alone, every eigenvalue other than $\mu_0$
exceeds $10^{-34}$.  These statements concern the complete operator,
including its infinite Fourier complement, at the single window
$\lambda=3$.
\end{proposition}
\begin{proof}[Shifted inertia and a matching trial direction]
For a threshold $a>0$ replace every diagonal $d_n$ by $d_n-a$ and
every tail lower weight $g_n$ by $g_n-a$.  Keep the same exact dyadic
witness $Z$ and the same head, intermediate and remote cutoffs as in
Proposition~\ref{prop:v116-window-positive}.  If $T_a=T-aI$, then
\[
 K_{Z,a}=K_{Z,0}-a(I+Z^*Z),\qquad
 R_a=B-T_aZ=R_0+aZ.
\]
Thus all intermediate residual rows change and are recomputed.  The
remote moment bound is unchanged, because those rows have no diagonal
entry and $G=\binom{I_{\rm head}}{-Z}$ is unchanged.  The diagonal inverse comparison gives
\[
 S_a:=F-aI-B^*T_a^{-1}B\succeq
 L_a:=K_{Z,a}-\sum_{N<n\le J}\frac{R_{a,n}^*R_{a,n}}{g_n-a}
                      -\frac{U_J}{g_{J+1}-a}.
\]
The sum uses normalized parity rows.  All denominators are strictly
positive by the earlier entire-tail bounds.

The signed interval LDL certificate for the even lower matrix at
$a_{\rm e}=10^{-34}$ has $257$ nonzero pivots, exactly one of which
is negative (zero-based index $24$).  For the odd lower matrix at
$a_{\rm o}=10^{-36}$ all $512$ pivots are strictly positive.
The complete signed elimination retains negative pivots with their
signs; it does not stop at the first negative pivot.  The exact dyadic
witnesses and every pivot enclosure are retained in the bundle;
\texttt{certify\_shifted\_inertia.py} recomputes all shifted residual rows.
The exact parameters are those of Proposition~\ref{prop:v116-window-positive}:
$(N,M,J,r)=(256,512,4096,80)$ for the even sector and
$(512,1024,4096,100)$ for the odd sector. Both certificates pass
at 768 bits and again at 896 bits using identical dyadic witnesses.
By finite-dimensional min--max (see also \cite[Section~4.3]{Teschl2009}), $S_{a_{\rm e}}$ has at most one
nonpositive eigenvalue, because the second eigenvalue of its lower
matrix is strictly positive.  The odd Schur matrix is positive definite.

A matching negative even direction is supplied by the existing
finite-supported candidate of
Proposition~\ref{prop:v112-finite-certificate}. More precisely, in that
$N=64$ certificate let $u$ be its normalized even vector, let $Q$ project
onto its even orthogonal complement. Set
$W_{64}=P_{64}\mathsf W_3P_{64}$, $\alpha=\ip{W_{64}u}{u}$,
$C=QW_{64}Q$ on that complement, and $r=QW_{64}u$. Write $A=C-\alpha I$, $z=A^{-1}r$, $q=\norm z$ and
$E=\ip zr$.  The vector
\[
 v=\frac{u-z}{\sqrt{1+q^2}}
\]
has norm one and belongs to the finite Fourier core.  Direct expansion
gives its exact complete-form Rayleigh value
\[
 \beta=QW_3(v,v)=\alpha-\frac{E}{1+q^2}.
\]
The previously certified enclosures imply
\[
 \alpha<5.312382\cdot10^{-38},\quad
 E>1.668982\cdot10^{-38},\quad q<0.000216,
 \qquad \beta<3.644\cdot10^{-38}.
\]
The last strict scalar comparison is separately checked by
\texttt{certify\_ground\_overlap\_constants.py}.  Because $v$ has
finite support, its finite-compression Rayleigh value equals its
complete-form value exactly.  No approximation of the omitted output
coefficients is used in this identity.  In particular $\beta<a_{\rm e}$.

On the closed form domain, square completion reads
\[
 q_a[x+y]=\ip{S_ax}{x}
       +\norm{T_a^{1/2}(y+T_a^{-1}Bx)}^2.
\]
The map $(x,y)\mapsto(x,y+T_a^{-1}Bx)$ is boundedly invertible and
preserves that domain.  It transfers the negative index to $S_a$;
also $\ker(\mathsf W_3-aI)$ is isomorphic to $\ker S_a$.
The negative even trial direction therefore makes the first even
Schur eigenvalue strictly negative, while its second is strictly
positive.  There is no even eigenvalue exactly at $a_{\rm e}$, and
exactly one even eigenvalue lies below it, counted with multiplicity.
Every odd eigenvalue exceeds $a_{\rm o}$.  The established compact
resolvent makes these full spectral statements.  Finally
Proposition~\ref{prop:v116-window-positive} supplies $\mu_0>0$,
and min--max applied to $v$ supplies $\mu_0\le\beta$.
Since $\beta<a_{\rm o}<a_{\rm e}$, the unique even low eigenvalue is
the simple ground state of the whole operator.  This proves
\eqref{eq:v117-ground-order} and the stronger even-sector threshold.
\end{proof}

\begin{corollary}[An ordinary ground-overlap bound at one window]
\label{cor:v117-ground-overlap}
For the normalized rational candidate $u$ of
Proposition~\ref{prop:v112-finite-certificate} and the complete unit
ground vector $\xi_0$ at $\lambda=3$,
\[
 \sin\angle(u,\xi_0)<0.01931.
\]
The same strict upper bound holds for the normalized exact $N=64$
source projection after the already certified normalization error is
included.  This is an ordinary Hilbert-space bound; it asserts no
endpoint comparison with $\xi_0$.
\end{corollary}
\begin{proof}
The improved trial $v$ in the proof above is even.  In its spectral
expansion, every non-ground component has energy greater than
$a_{\rm e}=10^{-34}$, and the ground energy is positive.  Hence
\[
 \sin\angle(v,\xi_0)
 <\sqrt{3.644\cdot10^{-38}/10^{-34}}.
\]
Also $z\perp u$, so the norm of the difference of the two rank-one
orthogonal projections onto $u$ and $v$ equals
$q/\sqrt{1+q^2}<0.000216$.  The triangle inequality for these
projectors proves
\[
 \sin\angle(u,\xi_0)
 <0.000216+\sqrt{3.644\cdot10^{-4}}
 <0.019305264<0.01931.
\]
For the unnormalized rational candidate $c$,
Proposition~\ref{prop:v113-source-certificate} gives
$\|P_{64}p_3-c\|<\varepsilon:=2.60\cdot10^{-36}$, while the
exact rational norm calculation gives $\|c\|>0.65172555$.
Writing $\rho=\varepsilon/\|c\|$, elementary ball geometry yields
\[
 \left\|\frac{P_{64}p_3}{\|P_{64}p_3\|}-\frac c{\|c\|}\right\|^2
 \le\frac{2\rho^2}{1+\sqrt{1-\rho^2}}<(4\cdot10^{-36})^2.
\]
Indeed the real inner product of these two unit vectors is positive
and at least $\sqrt{1-\rho^2}$, by minimizing the squared distance
from $c$ along the positive ray through $P_{64}p_3$.
This normalization error fits within the final margin.  The ordinary
norm controls projector distance but not the physical endpoint
functional; the earlier endpoint obstruction is therefore unchanged.
\end{proof}

This proves parity, simplicity and quantitative ordinary overlap at one
fixed window. It does not identify the unprojected source with the
ground vector. There is still no uniform family of spectral thresholds
or source overlap estimates along unbounded windows. The physical
endpoint is resolved below by a separate regularity argument; it does
not follow from this ordinary overlap bound. The sampler graph defect
and full-strip requirements remain separate; G2 and RH are not closed
by the fixed-window spectral result.

\subsection*{The physical endpoint of the complete Weil eigenfunctions}

The endpoint of an actual eigenfunction is not determined by ordinary
norm convergence of finite eigenvectors.  For the complete operator it
can instead be determined from the singular archimedean kernel.
Throughout this subsection $\Omega=(0,L)$, physical functions are
extended by zero outside $\Omega$, and Fourier frequency is normalized
as in $U_n(x)=L^{-1/2}e^{2\pi inx/L}$.

\begin{proposition}[Restricted logarithmic Laplacian decomposition]
\label{prop:v118-log-dirichlet}
Let $\mathcal A=\tfrac12L_\Delta^{\rm Dir}$ be the restricted,
exterior-Dirichlet logarithmic Laplacian on $\Omega$, whose full-line
symbol is $\log|t|$.  This is not the spectral logarithm of the
Dirichlet Laplacian.  The actual closed Weil operator satisfies
\begin{equation}\label{eq:v118-log-decomposition}
 \mathsf W_\lambda=\mathcal A+\mathcal K_\lambda,
 \qquad \mathcal D(\mathsf W_\lambda)=\mathcal D(\mathcal A),
\end{equation}
where, writing $k(y)=\rho(y)-(2y)^{-1}$, $a_m=\log m$, and
$w_m=\Lambda(m)/\sqrt m$,
\begin{align}
 (\mathcal K_\lambda f)(x)={}&-\log(2\pi)f(x)
       -\int_0^L k(|x-y|)f(y)\,dy\notag\\
 &-\sum_{1<m<\lambda^2}w_m
       \{\mathbf1_{x+a_m<L}f(x+a_m)
                   +\mathbf1_{x>a_m}f(x-a_m)\}\notag\\
 &+2\int_0^L\cosh((x-y)/2)f(y)\,dy.
 \label{eq:v118-bounded-physical-part}
\end{align}
The operator $\mathcal K_\lambda$ is real and self-adjoint on $L^2$.
For every $1\le p\le\infty$ it is bounded on $L^p(\Omega)$, with
\begin{equation}\label{eq:v118-K-bound}
 \|\mathcal K_\lambda\|_{p\to p}\le C_\lambda
 :=\log(2\pi)+2\int_0^L|k(y)|\,dy
                  +2P_\lambda+4\sinh(L/2).
\end{equation}
\end{proposition}
\begin{proof}
In dimension one the integral normalization of the logarithmic
Laplacian~\cite[Theorem~1.1]{ChenWethLog} gives
\[
 \mathcal A f(x)=\int_0^\infty
 \frac{2f(x)\mathbf1_{y<1}-f(x+y)-f(x-y)}{2y}\,dy
                   -\gamma_{\rm E}f(x)
\]
on compactly supported smooth tests.  The actual negative-archimedean
part of the Weil operator is
\[
 \int_0^\infty\rho(y)
 \{2e^{-y/2}f(x)-f(x+y)-f(x-y)\}\,dy
             -\{\log(4\pi)+\gamma_{\rm E}\}f(x).
\]
Its Fourier diagonal is precisely the negative of
\eqref{eq:v114-arch-diagonal}; the diagonal integral over $y>L$
equals $-\log\tanh(L/2)$.  Since
\[
 \int_0^\infty\left\{\frac1{\sinh y}
                -\frac{\mathbf1_{y<1}}y\right\}dy=\log2,
\]
subtracting $\mathcal A$ leaves $-\log(2\pi)I$ and the kernel $-k$.
The integral follows by evaluating $\log\tanh(y/2)$ with a positive
lower cutoff and then taking its limit.  In particular the constant
part of the logarithmic diagonal has been retained.
The prime shifts and the pole kernel $2\cosh((x-y)/2)$ give the
remaining terms of \eqref{eq:v118-bounded-physical-part}.
The latter kernel is the difference of the centered rank-one
$2\cosh\otimes\cosh$ and $2\sinh\otimes\sinh$ terms, so no pole
sign has been removed.

The function $k$ extends continuously to zero, with $k(0)=1/4$.
The integral-kernel row and column bounds, and the contraction of
each truncated translation, prove \eqref{eq:v118-K-bound}.
For the pole kernel the largest row integral is
$4\sinh(L/2)$.  Symmetry proves self-adjointness.

It remains to identify the closed forms, without imposing an
unproved endpoint condition.  Smooth compactly supported functions
are a form core for the restricted logarithmic
Laplacian~\cite[Theorem~3.1]{ChenWethLog}.  They are also a form core
for the operator in Proposition~\ref{prop:v114-weil-core}.
Indeed, approximate each finite periodic Fourier sum $f$ by smooth
cutoffs that remove endpoint strips of width $\epsilon$.
The discarded periodic function has $L^2$ norm $O(\epsilon^{1/2})$
and $H^1$ norm $O(\epsilon^{-1/2})$, hence $H^s$ norm
$O(\epsilon^{1/2-s})$ for any fixed $0<s<1/2$.
This tends to zero in the logarithmically weighted form norm.
The two closed lower-bounded forms therefore agree on a common form
core.  Bounded perturbation by $\mathcal K_\lambda$ proves
\eqref{eq:v118-log-decomposition}, including its operator-domain assertion.
\end{proof}

\begin{theorem}[Zero boundary trace of every complete Weil eigenfunction]
\label{thm:v118-eigen-endpoint-zero}
Fix $\lambda>1$.  Every $L^2$ eigenfunction
$\mathsf W_\lambda u=\mu u$ has a bounded representative whose zero
extension is continuous on $\mathbb R$.  With
\[
 \ell(r)=\frac1{|\log(\min\{r,0.1\})|},\qquad
 d(x)=\min\{x,L-x\},
\]
there is a finite constant $C$ such that
\begin{equation}\label{eq:v118-eigen-boundary-decay}
 |u(x)|\le C\sqrt{\ell(d(x))}\quad(0<x<L).
\end{equation}
In particular its actual physical endpoints satisfy $u(0)=u(L)=0$.
No uniform bound for $C$ as $\lambda$ grows is asserted.
\end{theorem}
\begin{proof}
First we prove boundedness, without assuming a positive Sobolev gain.
For $0<\epsilon<\min\{1,L\}$ let $\mathcal A_\epsilon$ be the
nonnegative killed small-jump generator
\[
 \mathcal A_\epsilon f(x)=\int_0^\epsilon
       \frac{2f(x)-f(x+y)-f(x-y)}{2y}\,dy
\]
in its closed Dirichlet-form realization, and put
\[
 J_\epsilon f(x)=\int_{\substack{y\in\Omega\\|x-y|\ge\epsilon}}
                    \frac{f(y)}{2|x-y|}\,dy.
\]
There is an exact decomposition
\[
 \mathcal A=\mathcal A_\epsilon+
       \{\log(1/\epsilon)-\gamma_{\rm E}\}I-J_\epsilon.
\]
The small-jump form is Markovian, including its nonnegative exterior
killing term.  Consequently its resolvent is consistent on $L^2$
and $L^\infty$, and
$\|(\mathcal A_\epsilon+aI)^{-1}\|_{p\to p}\le a^{-1}$ for
$a>0$, $p=2,\infty$.  One may obtain this directly by first cutting
off jumps smaller than a positive number, using the contraction
resolvent of that finite-activity killed jump generator, and then
passing to the increasing closed forms.  Also Cauchy--Schwarz gives
\[
 \|J_\epsilon\|_{2\to\infty}\le(2\epsilon)^{-1/2}.
\]
Choose $\epsilon$ so that
$a=\log(1/\epsilon)-\gamma_{\rm E}-\mu>C_\lambda$.
The eigen equation becomes
\[
 (\mathcal A_\epsilon+aI+\mathcal K_\lambda)u=J_\epsilon u.
\]
The inverse on the left exists on both $L^2$ and $L^\infty$ by the
Neumann series for
$[I+(\mathcal A_\epsilon+aI)^{-1}\mathcal K_\lambda]^{-1}
 (\mathcal A_\epsilon+aI)^{-1}$, with norm at most
$(a-C_\lambda)^{-1}$.  Since $J_\epsilon u\in L^2\cap L^\infty$,
the two series agree.  Uniqueness in $L^2$ identifies the bounded
solution with $u$, and proves
\[
 \|u\|_\infty\le
   \frac{\|u\|_2}{(a-C_\lambda)\sqrt{2\epsilon}}.
\]
Thus boundedness has not been inferred from the logarithmic operator
domain alone.

By Proposition~\ref{prop:v118-log-dirichlet}, the zero extension lies
in the logarithmic Dirichlet form space and satisfies
\[
 L_\Delta u=2(\mu u-\mathcal K_\lambda u)\in L^\infty(\Omega),
 \qquad u=0\quad\hbox{outside }\Omega.
\]
The interval satisfies an exterior uniform sphere condition.
The boundary regularity theorem of Hern\'andez-Santamar\'ia,
L\'opez R\'ios and Salda\~na~\cite[Theorem~1.1]{HLSLogBoundary}
therefore gives continuity of the zero extension and
\eqref{eq:v118-eigen-boundary-decay}.  For complex $u$, apply the
real theorem to its real and imaginary parts.  This proves the endpoint
assertion for every eigenfunction, independently of its spectral order.
\end{proof}

\begin{corollary}[Filtered endpoint decay and the nonzero-source obstruction]
\label{cor:v118-ground-fejer}
For an eigenfunction as above and $N\ge4$, define its Fej\'er endpoint
in the actual normalized Fourier coordinates by
\[
 \sigma_{N-1}u(0)=L^{-1/2}
       \sum_{|n|<N}(1-|n|/N)\widehat u(n).
\]
Then
\begin{equation}\label{eq:v118-fejer-endpoint}
 |\sigma_{N-1}u(0)|
 \le C\sqrt{\ell(L/\sqrt N)}+\frac{\|u\|_\infty}{2\sqrt N}
 =O_{\lambda,\mu,u}((\log N)^{-1/2}).
\end{equation}
For the exact repaired source with $B_\lambda=p_\lambda(0)\ne0$,
any complete-operator eigenfunction $u$ and scalar $c$ satisfy
\begin{equation}\label{eq:v118-ground-source-mismatch}
 \frac{|p_\lambda(0)-cu(0)|}{|B_\lambda|}=1.
\end{equation}
\end{corollary}
\begin{proof}
The positive Fej\'er kernel in physical coordinates is
$L^{-1}F_N(2\pi x/L)$, where
$F_N(\theta)=N^{-1}\{\sin(N\theta/2)/\sin(\theta/2)\}^2$;
it has integral one.  Put $h=L/\sqrt N\le L/2$.
On $d(x)<h$, \eqref{eq:v118-eigen-boundary-decay} bounds the
convolution by $C\sqrt{\ell(h)}$.
On $d(x)\ge h$, $\sin(\pi x/L)\ge2d(x)/L$ gives
\[
 L^{-1}F_N(2\pi x/L)\le\frac{L}{4N d(x)^2},\qquad
 \int_{d(x)\ge h}L^{-1}F_N(2\pi x/L)\,dx
       \le\frac{L}{2Nh}=\frac1{2\sqrt N}.
\]
This proves \eqref{eq:v118-fejer-endpoint}.
The mismatch identity follows from $u(0)=0$.
\end{proof}

Theorem~\ref{thm:v118-eigen-endpoint-zero} applies in particular to
the simple even ground at $\lambda=3$.  It disproves a nonzero physical
endpoint comparison between the repaired source and the complete
operator's ground, while leaving their ordinary-norm overlap intact.
It does not identify the endpoint of a finite-compression eigenvector
on a joint growing-window and growing-cutoff diagonal.
Nor does continuity alone prove absolute Fourier convergence or
convergence of sharp Fourier endpoint sums: the quantitative statement
\eqref{eq:v118-fejer-endpoint} concerns the stated positive filter.
The finite endpoint and its graph-defect terms therefore remain
separate from the complete eigenfunction's trace.

\begin{proposition}[The Weil graph norm does not control the physical trace]
\label{prop:v118-graph-trace}
At every fixed $\lambda>1$, physical endpoint evaluation on finite
Fourier polynomials has no continuous extension to
$\mathcal D(\mathsf W_\lambda)$ with its operator graph norm.
Indeed there are real even finite polynomials $h_N$ with
\[
 \delta(h_N)=1,\qquad
 \|h_N\|_2+\|\mathsf W_\lambda h_N\|_2
       =O_\lambda\!\left(\frac{\log(2+N)}{\sqrt N}\right)\longrightarrow0.
\]
The same failure holds with $\mathsf W_\lambda-\mu I$ for any fixed
$\mu\in\mathbb R$ and with the weaker closed-form norm.
\end{proposition}
\begin{proof}
In the physical logarithmic basis $U_n=L^{-1/2}e^{2\pi inx/L}$ set
\[
 h_N=\frac{\sqrt L}{2N}\sum_{N<|n|\le2N}U_n.
\]
There are exactly $2N$ terms. At either physical endpoint all their
phases are one, so $\delta(h_N)=1$ and
$\|h_N\|_2=\sqrt{L/(2N)}$. The exact decomposition in
Proposition~\ref{prop:v114-weil-core} gives
\[
 \|\mathsf W_\lambda h_N\|_2
 \le\left(\max_{N<|n|\le2N}|d_n|
              +2\pi B_\lambda^{\rm mat}\right)\sqrt{\frac L{2N}}.
\]
Since $d_n=\log|n|+O_\lambda(1)$, this proves the estimate.
A fixed scalar shift changes only the bounded constant. Graph convergence
implies form convergence for a lower-bounded self-adjoint operator.
In fact the trace is not closable on the polynomial core in the graph
norm: this sequence tends to zero while its trace tends to one.
\end{proof}

Let
\[
 \mathcal V_\lambda:=\operatorname{Var}(p_\lambda)
\]
be the total variation of the periodic logarithmic representative of the
inversion-even proxy.  Let
\[
 M_T(\lambda):=\max\{m\ge0:2\pi m/L\le T\}.
\]

\begin{proposition}[Explicit BV rate for the fixed-band Galerkin defect]\label{prop:g1-bv-rate}
Fix $T>0$.  If $N\ge2M_T(\lambda)+1$, then
\begin{equation}\label{eq:g1-gamma-rate}
 \boxed{
 \gamma_{\lambda,N,T}
 \le C_T\frac{\lambda L\,\mathcal V_\lambda}
 {|B_\lambda|\,N}.}
\end{equation}
In particular the form defect in G1 has a fully quantitative,
zero-independent rate; no $C^2$ or higher regularity of $p_\lambda$ is used.
\end{proposition}
\begin{proof}
Use the translated logarithmic coordinate $[0,L]$.  Since $p_\lambda$ has
bounded variation and matching periodic endpoint values, Stieltjes integration
by parts gives, for $n\ne0$,
\[
 |\widehat p_\lambda(n)|
 \le\frac{\sqrt L\,\mathcal V_\lambda}{2\pi|n|}.
\]
Let $g=D_{\log}v$ with $v\in E_{\lambda,T}$ and $\|v\|_2=1$.  Then
$\|g\|_2\le T$ and $g_m=0$ for $|m|>M_T(\lambda)$.  If $|n|>N\ge2M_T+1$,
Lemma~\ref{lem:g1-high-low} and Cauchy--Schwarz imply
\[
 \left|\sum_{|m|\le M_T}\tau^{(\lambda)}_{n,m}g_m\right|
 \le C_T\frac{\lambda\sqrt L}{|n|}.
\]
The resulting tail is absolutely summable, and hence
\[
 |QW_\lambda((P_N-I)p_\lambda,g)|
 \le C_T\lambda L\mathcal V_\lambda
       \sum_{|n|>N}\frac1{n^2}
 \le C_T\frac{\lambda L\mathcal V_\lambda}{N}.
\]
Divide by $|B_\lambda|$ and take the supremum over $v$.
\end{proof}

\subsection*{A compatibility-friendlier filtered finite section}
The exact-boundary repair below is useful for G1 alone, but it adds a new
high-frequency vector.  For the later same-object G1/G2 comparison it is
preferable to stay inside a canonical smoothing of the literal
prolate/co-Poisson proxy.  A de la Vall\'ee--Poussin filter does this while
preserving all sufficiently low Fourier modes exactly.

Work in the translated logarithmic coordinate $[0,L]$, with the endpoints
identified.  Let $\widehat p_\lambda(n)$ be the normalized Fourier
coefficients and define the multiplier
\begin{equation}\label{eq:g1-vp-multiplier}
 \nu_N(n):=
 \begin{cases}
  1,& |n|\le N,\\
  2-|n|/N,&N<|n|<2N,\\
  0,&|n|\ge2N.
 \end{cases}
\end{equation}
The de la Vall\'ee--Poussin mean is
\begin{equation}\label{eq:g1-vp-mean}
 \mathsf V_Np_\lambda
 :=\sum_{n\in\mathbb Z}\nu_N(n)\widehat p_\lambda(n)U_n\in E_{2N}.
\end{equation}
Equivalently its convolution kernel is the difference of two Fej\'er kernels,
$2F_{2N}-F_N$.  In particular its $L^1$ norm is bounded by an absolute
constant and, for $0<\delta\le L/2$, the kernel mass outside the boundary arc
$|x|<\delta$ is $O(L/(N\delta))$.

Define the endpoint modulus
\begin{equation}\label{eq:g1-edge-modulus}
 \omega_\lambda(\delta)
 :=\sup_{0\le t\le\delta}
 \max\{|p_\lambda(t)-B_\lambda|,
       |p_\lambda(L-t)-B_\lambda|\}.
\end{equation}
The no-boundary-jump property gives $\omega_\lambda(\delta)\to0$ as
$\delta\downarrow0$ for every fixed $\lambda$.

\begin{lemma}[Quantitative endpoint recovery by de la Vall\'ee--Poussin filtering]\label{lem:g1-vp-endpoint}
For every $N\ge1$ and $0<\delta\le L/2$,
\begin{equation}\label{eq:g1-vp-endpoint}
 \boxed{
 |\delta_{2N}(\mathsf V_Np_\lambda)-B_\lambda|
 \le C\left(
  \omega_\lambda(\delta)
  +\frac{\mathcal V_\lambda L}{N\delta}
 \right).}
\end{equation}
The constant is absolute.  Hence the endpoint recovery is quantitative in
terms of the explicit boundary modulus of the chosen prolate/co-Poisson
proxy and uses no zero information.
\end{lemma}
\begin{proof}
Let $K_N^{\mathrm{VP}}$ be the period-$L$ convolution kernel of
$\mathsf V_N$.  Since
$|K_N^{\mathrm{VP}}|\le2F_{2N}+F_N$, its $L^1$ norm is $O(1)$ and the
standard Fej\'er estimate gives
\[
 \int_{|x|\ge\delta}|K_N^{\mathrm{VP}}(x)|\,dx
 \le C\frac{L}{N\delta}.
\]
On the boundary arc use
$|p_\lambda(x)-B_\lambda|\le\omega_\lambda(\delta)$.  On its complement,
bounded variation and the absence of a periodic boundary jump give
$|p_\lambda(x)-B_\lambda|\le\mathcal V_\lambda$.  Splitting the convolution
integral proves \eqref{eq:g1-vp-endpoint}.
\end{proof}

For later bookkeeping it is useful to note that the endpoint modulus is itself
zero-independently computable from the source.  The compact support of the
prolate source makes the co-Poisson sum finite on the window; away from the
finitely many arithmetic thresholds it is a finite sum of smooth functions.
Thus there are explicit numbers $\Delta_\lambda>0$ and $H_\lambda<\infty$,
determined only by the chosen source and $\lambda$, for which
\begin{equation}\label{eq:g1-edge-lipschitz}
 \omega_\lambda(\delta)\le H_\lambda\delta,
 \qquad0<\delta\le\Delta_\lambda.
\end{equation}
For example one may take $\Delta_\lambda$ to be half the logarithmic distance
from the identified endpoint to the nearest interior arithmetic threshold;
$H_\lambda$ is then the supremum of the ordinary logarithmic derivative on
those two endpoint arcs.  No uniform-in-$\lambda$ estimate for $H_\lambda$ is
being asserted here.

\begin{proposition}[Filtered BV transfer with explicit boundary control]\label{prop:g1-vp-transfer}
Fix $T>0$ and suppose $N\ge2M_T(\lambda)+1$.  Set
\[
 k^{\mathrm{VP}}_{\lambda,N}:=\mathsf V_Np_\lambda\in E_{2N},
 \qquad
 \beta^{\mathrm{VP}}_{\lambda,N}
 :=\frac{|\delta_{2N}(k^{\mathrm{VP}}_{\lambda,N})-B_\lambda|}
 {|B_\lambda|}.
\]
Then
\begin{equation}\label{eq:g1-vp-form-defect}
 \boxed{
 \frac1{|B_\lambda|}
 \sup_{\substack{v\in E_{\lambda,T}\\\|v\|_2=1}}
 |QW_\lambda(k^{\mathrm{VP}}_{\lambda,N}-p_\lambda,
             D_{\log}v)|
 \le C_T\frac{\lambda L\mathcal V_\lambda}{|B_\lambda|N}.}
\end{equation}
Whenever $\beta^{\mathrm{VP}}_{\lambda,N}<1$,
\begin{equation}\label{eq:g1-vp-master}
 \boxed{
 \frac{\|P_TD_{\log}W_{\lambda,2N}
 k^{\mathrm{VP}}_{\lambda,N}\|_2}
 {|\delta_{2N}(k^{\mathrm{VP}}_{\lambda,N})|}
 \le
 \frac{T\mathcal E_T(\lambda)
 +C_T\lambda L\mathcal V_\lambda/(|B_\lambda|N)}
 {1-\beta^{\mathrm{VP}}_{\lambda,N}}.}
\end{equation}
Moreover, if $H_\lambda>0$ and
$N\ge C\mathcal V_\lambda L/(H_\lambda\Delta_\lambda^2)$, then
\begin{equation}\label{eq:g1-vp-beta-rate}
 \boxed{
 \beta^{\mathrm{VP}}_{\lambda,N}
 \le
 \frac{C}{|B_\lambda|}
 \sqrt{\frac{H_\lambda\mathcal V_\lambda L}{N}}.}
\end{equation}
Consequently one obtains a deterministic, zero-independent one-scale G1
diagonal by taking $N=N_\lambda$ large enough that
\begin{equation}\label{eq:g1-vp-cutoff}
 N_\lambda\ge C_T\max\left\{
 2M_T(\lambda)+1,
 \frac{\lambda L\mathcal V_\lambda}
      {|B_\lambda|\mathfrak r(\lambda)},
 \frac{H_\lambda\mathcal V_\lambda L}
      {|B_\lambda|^2\mathfrak r(\lambda)^2},
 \frac{\mathcal V_\lambda L}{H_\lambda\Delta_\lambda^2}
 \right\},
\end{equation}
with the last term omitted when it is vacuous.  Along such a diagonal,
\begin{equation}\label{eq:g1-vp-vanishing}
 \boxed{
 \frac{\|P_TD_{\log}W_{\lambda,2N_\lambda}
 k^{\mathrm{VP}}_{\lambda,N_\lambda}\|_2}
 {|\delta_{2N_\lambda}(k^{\mathrm{VP}}_{\lambda,N_\lambda})|}
 \le C'_T\mathfrak r(\lambda)\longrightarrow0.}
\end{equation}
\end{proposition}
\begin{proof}
The multiplier \eqref{eq:g1-vp-multiplier} is exactly one for $|n|\le N$.
Therefore the Fourier coefficients of
$k^{\mathrm{VP}}_{\lambda,N}-p_\lambda$ vanish in that range and, for
$|n|>N$, have modulus at most $|\widehat p_\lambda(n)|$.  The Stieltjes
integration-by-parts bound from Proposition~\ref{prop:g1-bv-rate} and the
high--low estimate of Lemma~\ref{lem:g1-high-low} therefore give
\eqref{eq:g1-vp-form-defect} with the same $O(N^{-1})$ tail as the ordinary
Galerkin projection.

Equation~\eqref{eq:g1-vp-master} is the same Hermitian fixed-band argument as
\eqref{eq:g1-master-bound}.  Lemma~\ref{lem:g1-vp-endpoint}, followed by
\eqref{eq:g1-edge-lipschitz} and the choice
$\delta\asymp(\mathcal V_\lambda L/(H_\lambda N))^{1/2}$, gives
\eqref{eq:g1-vp-beta-rate} once that $\delta$ lies inside the endpoint arc.
The cutoff \eqref{eq:g1-vp-cutoff} makes both the form defect and boundary
defect $O_T(\mathfrak r(\lambda))$; inserting the continuum envelope
\eqref{eq:G1-continuum} proves \eqref{eq:g1-vp-vanishing}.
\end{proof}

The filtered vector has another advantage over an external boundary repair:
it changes no Fourier coefficient below $N$.  Its derivative can still grow,
because the literal co-Poisson proxy is only BV across its arithmetic interior
thresholds, but the growth is intrinsic rather than injected by a separate
corrector.  The coefficient estimate used above gives the explicit bound
\begin{equation}\label{eq:g1-vp-derivative}
 \|D_{\log}k^{\mathrm{VP}}_{\lambda,N}\|_2
 \le C\mathcal V_\lambda\sqrt{\frac{N}{L}}.
\end{equation}

There is in fact a sharp obstruction to improving the derivative price by
pushing an exact boundary repair farther into the Fourier tail.

\begin{lemma}[High-frequency boundary repair has unavoidable square-root derivative cost]\label{lem:g1-repair-no-go}
Let
\[
 c(x)=\sum_{j=N+1}^{K}w_j\cos(2\pi jx/L),
 \qquad c(0)=\sum_{j=N+1}^{K}w_j=1.
\]
Then
\begin{equation}\label{eq:g1-repair-no-go}
 \boxed{
 \|D_{\log}c\|_2^2
 \ge
 \frac{2\pi^2}{L\displaystyle\sum_{j=N+1}^{K}j^{-2}}
 \ge\frac{2\pi^2N}{L}.}
\end{equation}
The first lower bound is sharp: equality is attained by weights
$w_j\propto j^{-2}$.
\end{lemma}
\begin{proof}
Orthogonality of the sine modes gives
\[
 \|D_{\log}c\|_2^2
 =\frac{2\pi^2}{L}\sum_{j=N+1}^{K}j^2|w_j|^2.
\]
Weighted Cauchy--Schwarz and $\sum w_j=1$ give
\[
 1\le
 \left(\sum j^2|w_j|^2\right)
 \left(\sum j^{-2}\right).
\]
This proves the first bound, with equality precisely for $w_j\propto j^{-2}$.
Finally
$\sum_{j>N}j^{-2}\le\int_N^\infty x^{-2}\,dx=N^{-1}$.
\end{proof}

Lemma~\ref{lem:g1-repair-no-go} shows that the square-root derivative growth
seen in the shell construction is not merely an artifact of equal weights: it
is unavoidable for every even exact-boundary corrector supported entirely
above frequency $N$.  For a norm-level same-object estimate, the filtered
literal proxy of Proposition~\ref{prop:g1-vp-transfer} is therefore still the
preferred G1 candidate.  This derivative lower bound is not, however, a
pairing-level no-go: the shell estimates in Section~\ref{sec:matching} below proves that the
shell's explicit G1 and scalar-shift pairings against the higher-order Hardy
sampler can be made arbitrarily small by moving the shell outward.

The endpoint defect can also be removed exactly without asking for a
quantitative Dirichlet--Jordan theorem.  Put
\[
 q_{\lambda,N}=P_Np_\lambda,
 \qquad
 d_{\lambda,N}=B_\lambda-q_{\lambda,N}(0),
\]
and, for an even integer $K>2N$, define the even high-frequency shell
corrector
\begin{equation}\label{eq:g1-shell-corrector}
 c^{\mathrm{sh}}_{\lambda,K}(x)
 :=\frac{2}{K}\sum_{j=K/2+1}^{K}\cos(2\pi jx/L)
 \quad\text{in logarithmic coordinates}.
\end{equation}
Every summand has boundary value one, so
$c^{\mathrm{sh}}_{\lambda,K}(0)=1$.  Moreover all of its Fourier modes lie
above $N$.  Then set
\begin{equation}\label{eq:g1-bc-candidate}
 k^{\mathrm{bc}}_{\lambda;N,K}
 :=q_{\lambda,N}+d_{\lambda,N}c^{\mathrm{sh}}_{\lambda,K}\in E_K.
\end{equation}
By construction it is real, inversion-even, and
\begin{equation}\label{eq:g1-bc-boundary}
 \delta_K(k^{\mathrm{bc}}_{\lambda;N,K})=B_\lambda
\end{equation}
exactly.

\begin{theorem}[Constructive finite fixed-band G1 closure]\label{thm:g1-boundary-corrected}
Use the explicit repaired source above and fix $T>0$. For $N\ge2M_T(\lambda)+1$ and an even integer
$K>2N$,
\begin{equation}\label{eq:g1-bc-master}
 \boxed{
 \begin{aligned}
 \frac{\|P_TD_{\log}W_{\lambda,K}
 k^{\mathrm{bc}}_{\lambda;N,K}\|_2}{|B_\lambda|}
 \le{}&T\mathcal E_T(\lambda)
 +C_T\frac{\lambda L}{N}
       \frac{\mathcal V_\lambda}{|B_\lambda|}\\
 &+C_T\frac{\lambda L}{K}
 \left(
 1+\sqrt{\frac{2N+1}{L}}
       \frac{\|p_\lambda\|_2}{|B_\lambda|}
 \right).
 \end{aligned}}
\end{equation}
Consequently there are completely zero-independent, deterministic cutoffs
$N_\lambda<K_\lambda$ for which
\begin{equation}\label{eq:g1-bc-vanishing}
 \boxed{
 \frac{\|P_TD_{\log}W_{\lambda,K_\lambda}
 k^{\mathrm{bc}}_{\lambda;N_\lambda,K_\lambda}\|_2}
 {|\delta_{K_\lambda}(k^{\mathrm{bc}}_{\lambda;N_\lambda,K_\lambda})|}
 \le C'_T\mathfrak r(\lambda)\longrightarrow0.}
\end{equation}
One admissible prescription is to choose, after enlarging $C_T$ if necessary,
\begin{align}
 N_\lambda&\ge
 \max\left\{2M_T(\lambda)+1,
 C_T\frac{\lambda L}{\mathfrak r(\lambda)}
       \frac{\mathcal V_\lambda}{|B_\lambda|}\right\},
 \label{eq:g1-N-explicit}\\
 K_\lambda&\text{ even},\qquad K_\lambda>2N_\lambda,\qquad
 K_\lambda\ge
 C_T\frac{\lambda L}{\mathfrak r(\lambda)}
 \left(1+\sqrt{\frac{2N_\lambda+1}{L}}
       \frac{\|p_\lambda\|_2}{|B_\lambda|}\right).
 \label{eq:g1-K-explicit}
\end{align}
All quantities on the right are explicit functions of the chosen
prolate/co-Poisson proxy and $\lambda$; no nontrivial zeta zero enters the
choice.
\end{theorem}
\begin{proof}
For $v\in E_{\lambda,T}$ with $\|v\|_2=1$, Hermiticity gives
\[
 |\langle P_TD_{\log}W_{\lambda,K}
 k^{\mathrm{bc}}_{\lambda;N,K},v\rangle|
 =|QW_\lambda(k^{\mathrm{bc}}_{\lambda;N,K},D_{\log}v)|.
\]
Split the candidate as
\[
 k^{\mathrm{bc}}_{\lambda;N,K}
 =p_\lambda+(P_N-I)p_\lambda+d_{\lambda,N}c^{\mathrm{sh}}_{\lambda,K}.
\]
The first term is at most $T|B_\lambda|\mathcal E_T(\lambda)$ by
Lemma~\ref{lem:g1-inversion}; Proposition~\ref{prop:g1-bv-rate} controls the
second.  For each mode in the shell, Lemma~\ref{lem:g1-high-low} and
Cauchy--Schwarz on the fixed physical band give a bound
$C_T\lambda L/j$; since $K/2<j\le K$ and the weights in
\eqref{eq:g1-shell-corrector} sum to one, averaging yields
\[
 |QW_\lambda(c^{\mathrm{sh}}_{\lambda,K},D_{\log}v)|
 \le C_T\frac{\lambda L}{K}.
\]
Finally the endpoint-evaluation norm on $E_N$ is
$\sqrt{(2N+1)/L}$, so
\[
 \frac{|d_{\lambda,N}|}{|B_\lambda|}
 \le1+\sqrt{\frac{2N+1}{L}}
       \frac{\|p_\lambda\|_2}{|B_\lambda|}.
\]
These estimates prove \eqref{eq:g1-bc-master}.  The choices
\eqref{eq:g1-N-explicit}--\eqref{eq:g1-K-explicit}, together with
$\mathcal E_T(\lambda)\le C_T\mathfrak r(\lambda)$, make each of the last two
terms $O_T(\mathfrak r(\lambda))$.  Equation~\eqref{eq:g1-bc-boundary} removes
the denominator defect exactly, proving \eqref{eq:g1-bc-vanishing}.
\end{proof}

Theorem~\ref{thm:g1-boundary-corrected} transfers the proved continuum G1
bound to a finite vector with exact endpoint normalization. It is a weak
fixed-band estimate. Increasing the cutoff does not make the full
derivative of the endpoint repair small: its squared norm is
\[
 |d_{\lambda,N}|^2\frac{8\pi^2}{K^2L}
 \sum_{j=K/2+1}^{K}j^2\le
 4\pi^2|d_{\lambda,N}|^2K/L.
\]
The growing derivative is relevant to any later graph argument.
The scalar-shift identity remains exact,
\[
 D_{\log}(W-\alpha I)k=D_{\log}Wk-\alpha D_{\log}k,
\]
but a weak bound on the first term supplies no norm bound on the second.

For filtered endpoint recovery, if the local Lipschitz constant is zero,
the near-endpoint error is identically zero; the cutoff is chosen from
the far-tail error alone. A formula that divides by that Lipschitz
constant is used only in its positive case. All other displayed cutoff
thresholds depend on the specified source, its variation, and the desired
tolerance, without using a list of zeta zeros.

The next section restores an explicit weighted-sampler estimate from the
earlier manuscript. The subsequent sampler construction then separates
ordinary norm, physical endpoint, Weil form, and graph behavior. These
four notions of compatibility cannot be interchanged.


\subsection*{Ordinary source mass and a full operator residual}
The Sobolev G1 estimate and the polynomial Fourier tail also give an
absolute ordinary-norm operator residual for this fixed source. This supplies the
Rayleigh and residual prerequisites in the spectral-overlap criterion;
it does not supply the overlap itself.

\begin{lemma}[A nonzero ordinary-norm limit of the repaired proxy]
\label{lem:v115-source-mass}
Put
\[
 h_\infty(u)=\frac{4\pi}{\sqrt3}u^2(2\pi u^2-3)e^{-\pi u^2},
 \qquad f_\infty=\mathcal E(h_\infty).
\]
Extend $p_\lambda$ by zero outside $I_\lambda$. Then
\begin{equation}\label{eq:v115-source-limit}
 \norm{p_\lambda-f_\infty}_{L^2(du/u)}=O(\lambda^{-1/2}),
 \qquad Jf_\infty=f_\infty.
\end{equation}
In particular $0<\norm{f_\infty}<\infty$. For
$N=N_\lambda=\lceil\lambda^8(1+L)\rceil$ and
$k_\lambda=P_Np_\lambda$, both $\norm{p_\lambda}$ and
$\norm{k_\lambda}$ tend to $\norm{f_\infty}$ and are bounded above
and away from zero. A valid eventual common lower bound is
\begin{equation}\label{eq:v115-mass-lower}
 m_*:=\frac14\left(\int_1^2h_\infty(u)^2\,du\right)^{1/2}>0.
\end{equation}
No numerical threshold in $\lambda$ is asserted.
\end{lemma}
\begin{proof}
The uniform fixed-mode Hermite approximation in
\cite[Lemma~7.2]{ConnesConsaniMoscovici2025} gives suitably normalized
$q_{j,\lambda}$ with
$\sup_{|u|\le\lambda}|q_{j,\lambda}(u)-H_j(u)|=O(\lambda^{-2})$,
$j=0,4$, where $H_j$ has unit $L^2(\mathbb R)$ norm. This rate survives
the manuscript's exact unit normalization: expanding the squared norm,
the cross term is $O(\lambda^{-2})\norm{H_j}_1$, the squared-error term
is $O(\lambda^{-3})$, and the omitted Hermite tail is Gaussian.
Thus $\norm{q_{j,\lambda}}=1+O(\lambda^{-2})$.

The limits of the coefficients in \eqref{eq:v12-coefficients} are
$-H_4(0)=-\sqrt3/(2\,2^{1/4})$ and $H_0(0)=2^{1/4}$, with error
$O(\lambda^{-2})$. Substitution of the normalized Hermite polynomials
gives exactly the displayed $h_\infty$, including its scalar factor.
The fixed repair has exponentially small amplitude by
\eqref{eq:v12-zero-value}. Consequently
\[
 \sup_{|u|\le\lambda}|\widetilde h_\lambda(u)-h_\infty(u)|
       \le C\lambda^{-2}.
\]
The limiting source is a Fourier-invariant combination of orders $0,4$,
with zero value at zero and zero integral. Poisson summation therefore
gives $Jf_\infty=f_\infty$. Its Gaussian decay at infinity, followed by
inversion at zero, gives square integrability.

On the physical window there are at most $\lambda/u$ compact-source
summands. Their accumulated approximation error is at most
$C\lambda^{-1}u^{-1/2}$. Its squared Haar norm is bounded by
\[
 C\lambda^{-2}\int_{1/\lambda}^\lambda u^{-2}\,du
       \le C\lambda^{-1}.
\]
The omitted Gaussian summands $nu>\lambda$ contribute an exponentially
smaller norm: bound their decreasing polynomial-Gaussian tail by its
first term plus $u^{-1}$ times its tail integral. The same estimate and
inversion control $f_\infty$ outside $I_\lambda$. Inversion averaging
is an $L^2$ contraction, proving \eqref{eq:v115-source-limit}.
For $u\ge1$, all summands $h_\infty(nu)$ are positive. Thus
$\norm{f_\infty}^2\ge\int_1^2h_\infty(u)^2\,du>0$.
Finally \eqref{eq:v114-high-tail} gives
\[
 \norm{(I-P_N)p_\lambda}
       \le C|B_\lambda|\lambda\sqrt{L/N}\longrightarrow0,
\]
so the projected vectors have the same nonzero norm limit. Choosing a
common sufficiently large threshold proves \eqref{eq:v115-mass-lower}.
\end{proof}

\begin{proposition}[A full residual after ordinary norm normalization]
\label{prop:v115-absolute-residual}
Use the same $N_\lambda,k_\lambda$ and the canonical closed operator
$\mathsf W_\lambda$ of Proposition~\ref{prop:v114-weil-core}.
For all sufficiently large $\lambda$, with $c=2\pi\lambda^2$,
\begin{equation}\label{eq:v115-absolute-residual}
 \norm{\mathsf W_\lambda p_\lambda}
 +\norm{\mathsf W_\lambda k_\lambda}
 +\norm{W_{\lambda,N_\lambda}k_\lambda}
       \le C\lambda^6 e^{-c/3}.
\end{equation}
For either source normalized to ordinary norm one, its Rayleigh value
and its full centered residual tend to zero at this rate. The finite
matrix assertion uses that same normalization and the original physical
Fourier cutoff.
\end{proposition}
\begin{proof}
All source vectors belong to $\mathcal D(\mathsf W_\lambda)$ by
Proposition~\ref{prop:v114-weil-core}. The weighted Chebyshev estimate
$P_\lambda=O(\lambda)$, together with $h_L=O(\lambda)$, gives
$\norm{[M_b,H]}=O(\lambda)$ in \eqref{eq:v114-b-bound}.
Corollary~\ref{cor:g1-complete-matrix}, including its separately proved
diagonal, gives uniformly
\[
 |d_n|\le C\{\lambda+1+\log(2+|n|/L)\}.
\]
Let $N=N_\lambda$ and $e=(I-P_N)p_\lambda$.
Using \eqref{eq:v114-high-tail} and comparing
$\sum_{n>N}\log^2(2+n/L)/n^2$ with its integral yields
\begin{align}
 \norm{\mathsf W_\lambda e}
 &\le C|B_\lambda|\lambda\sqrt{L/N}
        \{\lambda+1+\log(2+N/L)\}\notag\\
 &\le C|B_\lambda|\lambda^{-2}.
 \label{eq:v115-uniform-graph-tail}
\end{align}
Both the diagonal part and the bounded commutator are retained.

Next split the \emph{output} at an auxiliary integer $M>2N$.
This does not change the candidate or its finite matrix. For any
$v=P_Mv$ with $\norm v=1$,
\[
 \norm{D_{\log}v}\le2\pi M/L,\qquad
 |v(-a)|+|v(a)|\le2\sqrt{(2M+1)/L}.
\]
The full Sobolev estimate \eqref{eq:G1-sobolev}, whose constant is
independent of its test function, therefore implies
\begin{equation}\label{eq:v115-low-output}
 \norm{P_M\mathsf W_\lambda p_\lambda}
 \le C|B_\lambda|\mathfrak r(\lambda)
       \{1+M/L+\sqrt{M/L}\}.
\end{equation}
This includes the zero Fourier mode. It uses the already established
compact-source radicality and centered pole--prime cancellation, with
both source moments and inversion averaging intact. No fixed-band
constant is extrapolated to a growing band.

The diagonal has no cross block between $E_N$ and $E_M^\perp$.
For the complete off-diagonal matrix, Lemma~\ref{lem:g1-high-low}
gives the Hilbert--Schmidt bound
\[
 \norm{(I-P_M)\mathsf W_\lambda P_N}^2
 \le C\lambda^2\sum_{|m|\le N}\sum_{|n|>M}\frac1{(n-m)^2}
 \le C\lambda^2N/M.
\]
Lemma~\ref{lem:v115-source-mass} and
\eqref{eq:v115-uniform-graph-tail} now give
\begin{equation}\label{eq:v115-high-output}
 \norm{(I-P_M)\mathsf W_\lambda p_\lambda}
 \le C\lambda^5\sqrt{1+L}\,M^{-1/2}
         +C|B_\lambda|\lambda^{-2}.
\end{equation}
Thus all prime, pole and archimedean pieces remain in the actual matrix;
no divergent exterior pole or prime norm sum is estimated separately.

Take $M=\lceil e^{2c/3}\rceil$, which exceeds $2N$ eventually.
By \eqref{eq:v12-combined-endpoint} and
\eqref{eq:v13-endpoint-comparison},
$|B_\lambda|\le C\lambda^{11/2}e^{-c}$, while
$\mathfrak r(\lambda)$ is bounded. Equations
\eqref{eq:v115-low-output}--\eqref{eq:v115-high-output} are therefore
at most $C\lambda^6e^{-c/3}$ in total. This proves the first term of
\eqref{eq:v115-absolute-residual};
\eqref{eq:v115-uniform-graph-tail} proves the second, and compression
contractivity proves the third.
Finally ordinary normalization is legitimate by
\eqref{eq:v115-mass-lower}. For a unit vector $u$, letting
$\alpha=\ip{\mathsf W_\lambda u}u$, one has
$|\alpha|\le\norm{\mathsf W_\lambda u}$ and
$\norm{(\mathsf W_\lambda-\alpha)u}\le2\norm{\mathsf W_\lambda u}$.
The same argument applies to the compressed matrix.
\end{proof}

The preceding proposition supplies no endpoint-normalized operator residual:
dividing its bound by $|B_\lambda|$ leaves an exponentially growing
upper bound. Nor does it identify the spectral point near zero as the
lowest one. The elementary family $\operatorname{diag}(-1,\epsilon)$,
with candidate $(0,1)$, has Rayleigh value $\epsilon$ and centered
residual zero while its lowest level stays at $-1$. An independent
quantitative ground-space overlap or signed complementary lower bound
is still necessary. No estimate on the corrected sampler's graph defect,
or on a growing family of prolate modes, is inferred here.

\subsection*{G2 calibration from the screw-function realization}
Recent work of Suzuki~\cite{Suzuki2026} gives a useful, but more limited, structural calibration of G2.  Let $A_a$ denote the lower-bounded self-adjoint operator associated with the semilocal Weil form on $(-a,a)$, and write
\[
 \lambda_a=\inf\sigma(A_a).
\]
Suzuki proves that $A_a$ has discrete spectrum bounded from below and that $a\mapsto\lambda_a$ is continuous.  He also proves positivity, simplicity, and evenness of the ground state for \emph{sufficiently small} $a$ by a separate small-support asymptotic argument.  Those simplicity/parity conclusions must not be extrapolated to arbitrary growing support; the large-$a$ spectral ordering needed here remains open.

There is nevertheless an elementary global monotonicity which is important for claim discipline.  If $a_2>a_1$, extension by zero identifies the smaller compactly supported form domain with a subspace of the larger one, while the global Weil form is unchanged on such vectors.  The min--max principle therefore gives
\begin{equation*}
 \boxed{\lambda_{a_2}\le \lambda_{a_1}.}
 \tag{G2.1}\label{eq:g2-monotone}
\end{equation*}
In particular, once $\lambda_a$ becomes negative it remains bounded above by that same negative value at every larger support.  This is compatible with Suzuki's global continuity theorem and his observation that failure of RH is equivalent to the existence of a support with $\lambda_a<0$.

Weak fixed-band G1 alone does not test the form on the changing
proxy itself. The stronger operator estimate in
Proposition~\ref{prop:v115-absolute-residual} now supplies that missing
Rayleigh prerequisite for the actual normalized source and its literal
polynomial projection: $q_a(u_a,u_a)\to0$. Hence monotonicity gives
\[
 \lambda_a\to0\quad\Longleftrightarrow\quad
 \lambda_a\ge0\text{ for every finite }a.
\]
The forward implication uses only monotonicity; the reverse combines
positivity with the proved Rayleigh upper bound. The right side is the
compact-support Weil positivity criterion. The new source estimate
proves the upper bound, not positivity or the asserted limit of the
bottom level.

A quantitative spectral-overlap formulation isolates one possible extra ingredient without assuming simplicity.  Let $P_a^{(0)}$ be the orthogonal projection onto the (finite-dimensional) lowest eigenspace of $A_a$, let $p_a\in\mathcal D(A_a)$ be a normalized explicit proxy, and put
\[
 \alpha_a=\langle A_ap_a,p_a\rangle,
 \qquad
 r_a=(A_a-\alpha_a)p_a,
 \qquad
 \kappa_a=\|P_a^{(0)}p_a\|.
\]
Spectral expansion gives the exact bound
\[
 \boxed{
 |\lambda_a-\alpha_a|\,\kappa_a\le \|r_a\|.
 }
 \tag{G2.3}
\]
Proposition~\ref{prop:v115-absolute-residual} now proves both
$\alpha_a\to0$ and $\|r_a\|\to0$ for the normalized actual source,
with $\lambda=e^a$ and an upper envelope
\[
 \varepsilon_\lambda=C\lambda^6e^{-2\pi\lambda^2/3}.
\]
Therefore an independently proved overlap bound
$\kappa_a\ge\kappa_0>0$ would imply $\lambda_a\to0$ and close G2.
More generally it suffices that
$\varepsilon_\lambda/\kappa_a\to0$ along an unbounded sequence;
for example, $\kappa_a\ge e^{-\theta c}$ with fixed
$\theta<1/3$ and $c=2\pi\lambda^2$ would suffice.
No such lower bound is proved. If instead
$\lambda_a\le-\delta<0$, spectral projection of
$A_au_a$ gives $\delta\kappa_a\le\|A_au_a\|$.
Thus a negative ground state can coexist with the new residual estimate
by having exponentially small overlap with this source.

At small support Suzuki's positivity-improving analysis gives a simple
positive ground state, but it does not furnish a uniform large-support
overlap bound. Even positivity of both vectors would imply only a
nonzero overlap, without its required quantitative size. The model
$\operatorname{diag}(-1,\epsilon)$ with source $(0,1)$ shows directly
why a vanishing source residual does not determine the bottom sign.

This calibration is deliberately restrictive: continuity and monotonicity of the bottom level are structural, while large-support simplicity/parity and quantitative spectral overlap remain part of the open RH-strength content.  A successful continuation must therefore find new large-support structure, not extrapolate the small-support Dirichlet-form theorem.

\subsection*{A Schur-complement formulation of the large-support obstruction}
The preceding overlap criterion singles out one vector.  A more robust formulation does not require identification of the ground state and should be stated at the level of the closed semilocal Weil form $q_a$, rather than on the operator domain of $A_a$.  Let $P_a$ be any finite-rank orthogonal projection whose range is contained in the form domain, and decompose
\[
 L^2(-a,a)=\mathcal R_a\oplus\mathcal R_a^\perp,
 \qquad \mathcal R_a=\operatorname{Ran}P_a.
\]
The following elementary estimate isolates exactly the diagonal, cross, and complement quantities that a prolate near-radical construction would have to control.

\begin{proposition}[Form-Schur lower bound]\label{prop:g2-schur}
Suppose that for all $r\in\mathcal R_a$ and all $s\in\mathcal R_a^\perp$ in the form domain,
\[
 q_a(r,r)\ge-\alpha_a\norm r^2,
 \qquad
 |q_a(r,s)|\le\beta_a\norm r\norm s,
 \qquad
 q_a(s,s)\ge c_a\norm s^2
\]
with $c_a>0$, $\alpha_a\ge0$, and $\beta_a\ge0$. Then
\[
 \boxed{
 \lambda_a\ge -\alpha_a-\frac{\beta_a^2}{c_a}.
 }
 \tag{G2.4}
\]
In particular, $\alpha_a\to0$ and $\beta_a^2/c_a\to0$ imply $\liminf_{a\to\infty}\lambda_a\ge0$.
\end{proposition}
\begin{proof}
For $x=r+s$ the hypotheses give
\[
 q_a(x,x)
 \ge -\alpha_a\norm r^2-2\beta_a\norm r\norm s+c_a\norm s^2.
\]
Completing the square in $\norm s$ yields
\[
 q_a(x,x)
 \ge -\left(\alpha_a+\frac{\beta_a^2}{c_a}\right)\norm r^2
 \ge -\left(\alpha_a+\frac{\beta_a^2}{c_a}\right)\norm x^2.
\]
Taking the infimum of the Rayleigh quotient proves the claim.
\end{proof}

\subsection*{A complement bound without a positive spectral gap}
The ordinary operator residual permits a different reduction from the
coercive Schur estimate.  It does not require a positive gap on the
complement and therefore introduces no division by a collapsing level.

\begin{proposition}[A gap-free near-radical block bound]
\label{prop:v118-gap-free}
Let $A$ be a lower-bounded self-adjoint operator with closed form $q$,
and let $P$ be a finite-rank orthogonal projection with
$\operatorname{Ran}P\subset\mathcal D(A)$.  Put $Q=I-P$ and assume
$\norm{AP}\le\epsilon$.  Then $Q$ preserves the form domain and
\begin{equation}\label{eq:v118-gap-free-difference}
 \left|q(f,f)-q(Qf,Qf)\right|
 \le\frac{2}{\sqrt3}\epsilon\norm f^2
 \qquad(f\in\mathcal D(q)).
\end{equation}
Consequently, if
\[
 q(g,g)\ge-\eta\norm g^2\quad
 (g\in\operatorname{Ran}Q\cap\mathcal D(q)),\qquad \eta\ge0,
\]
then
\begin{equation}\label{eq:v118-gap-free-lower}
 \inf\sigma(A)\ge-\eta-\frac{2}{\sqrt3}\epsilon.
\end{equation}
The constant $2/\sqrt3$ in both statements is sharp.
\end{proposition}
\begin{proof}
Write $h=Pf$ and $g=Qf$.  The vector $h$ lies in the operator domain,
so $g$ belongs to the form domain and, with the fixed first-slot-linear
convention,
\[
 q(f,f)-q(g,g)
 =\ip{Ah}{h}+2\Re\ip{Ah}{g}
 =\Re\ip{Ah}{h+2g}.
\]
Since $h\perp g$, the absolute value is at most
\[
 \epsilon\norm h\sqrt{\norm h^2+4\norm g^2}
 \le\frac{2}{\sqrt3}\epsilon
           (\norm h^2+\norm g^2).
\]
For the last inequality, the maximum of
$t\sqrt{4-3t^2}$ on $0\le t\le1$ is $2/\sqrt3$.
This proves \eqref{eq:v118-gap-free-difference}; the complementary
lower bound and $\norm{Qf}\le\norm f$ prove
\eqref{eq:v118-gap-free-lower}.  Sharpness already holds on
$\mathbb C^2$: take $P$ onto the first coordinate and
\[
 A=-\frac{\epsilon}{\sqrt3}
   \begin{pmatrix}1&\sqrt2\\\sqrt2&0\end{pmatrix}.
\]
Then $\norm{AP}=\epsilon$, the complementary form is zero, and the
lowest eigenvalue is $-2\epsilon/\sqrt3$.
\end{proof}

For the rank-one projection onto either normalized source of
Proposition~\ref{prop:v115-absolute-residual},
$\norm{\mathsf W_\lambda P}=\norm{\mathsf W_\lambda u_\lambda}$
already tends to zero with the proved ordinary residual bound.
It therefore suffices to prove asymptotic nonnegativity
$q_\lambda|_{u_\lambda^\perp}\succeq-\eta_\lambda I$ with
$\eta_\lambda\to0$ along an unbounded window sequence.  Strict
complementary coercivity is unnecessary for this version of the
reduction.  For an arbitrary growing source block the required uniform
$\norm{\mathsf W_\lambda P_\lambda}\to0$ must still be established.
Proposition~\ref{prop:v135-growing-radical} and
Corollary~\ref{ns46-a2-cor-large-block} prove it for their specified
translate families. Neither complementary sign estimate is proved here.

The missing sign cannot be concealed by removing the near-radical
block.  Indeed, suppose a fixed unit compact test $f$ has
$q(f,f)=-\delta<0$, and extend it by zero to every larger window.
Then \eqref{eq:v118-gap-free-difference} gives
\[
 q_\lambda((I-P_\lambda)f,(I-P_\lambda)f)
 \le-\delta+\frac{2}{\sqrt3}\epsilon_\lambda.
\]
Whenever $\epsilon_\lambda<\sqrt3\delta/2$, the projected vector
is nonzero and negative; its normalized Rayleigh value is bounded
above by the same negative quantity, since its norm is at most one.
Thus a negative direction persists in the complement of any uniformly
near-radical block, regardless of its rank.  By the existing cofinal
Weil-positivity criterion, the desired asymptotic complementary sign
for the proved rank-one source remains RH-strength.  The reduction
removes a gap denominator; it does not supply the arithmetic sign.

\subsection*{Strong cancellation persists in the exact source complement}

The proved ordinary source residual does not mean that significant
prime/non-prime cancellation is confined to one direction. The following
fixed-window certificate tests that possible shortcut on the literal
source, rather than on sampled basis or random vectors.

\begin{proposition}[A certified cancelling direction off the source]
\label{prop:v121-complement-cancellation}
At $\lambda=3$ there exists an inversion-even unit Fourier polynomial
$v\in E_{64}$ exactly orthogonal to the repaired source $p_3$ such that
\begin{equation}\label{eq:v121-complement-cancellation}
 \begin{gathered}
 \ip{v}{p_3}=0,\qquad
 0<QW_3(v,v)<3\cdot10^{-31},\\
 q_{\rm np}(v)>0.4,\qquad q_{\rm pr}(v)<-0.4.
 \end{gathered}
\end{equation}
Here $q_{\rm np}=q_{\rm ar}+q_{\rm po}$ includes both the pure
archimedean term and the pole term, and $q_{\rm pr}$ is the signed
prime contribution, so $QW_3=q_{\rm np}+q_{\rm pr}$.
\end{proposition}
\begin{proof}
Let $V$ have the two real dyadic columns supplied in
Appendix~\ref{app:v121-audit}, in the normalized even basis of $E_{64}$.
Let $c$ be the exact rational source candidate from
Proposition~\ref{prop:v113-source-certificate}, converted to the same
parity coordinates. Set $a=V^*c$ and $x=(-a_2,a_1)^t$. The exact
algebra gives $\ip{Vx}{c}=0$. The certified Gram matrix $V^*V$ is
positive definite and $\norm{Vx}>0$, so put $w=Vx/\norm{Vx}$.
All vectors are real here; the stated Hermitian convention is unchanged.

Outward interval evaluation of the separate finite matrices gives
\[
 \begin{aligned}
 2.5844134\cdot10^{-31}&<QW_3(w,w)<2.5844135\cdot10^{-31},\\
 q_{\rm np}(w)&>0.4173237590,\\
 q_{\rm pr}(w)&<-0.4173237590.
 \end{aligned}
\]
These are inequalities for exact dyadic/rational constructions;
approximate eigenvectors were used only to propose the dyadic columns.

Write $u=P_{64}p_3/\norm{P_{64}p_3}$ and $\widehat c=c/\norm c$.
The existing exact-source certificate, matched to $c$ by its file hash,
gives $\norm{P_u-P_{\widehat c}}<\delta=4\cdot10^{-36}$.
Since $w\perp c$, define
\[
 v=\frac{(I-P_u)w}{\norm{(I-P_u)w}}.
\]
Then $\norm{v-w}<2\delta$. Also $v\in E_{64}$, hence $v\perp u$
implies $v\perp p_3$, not just orthogonality to an approximate source.
For any finite Hermitian matrix $H$,
\[
 |\ip{Hv}{v}-\ip{Hw}{w}|
 \le2\norm H\norm{v-w}<4\delta\norm H.
\]
The certified all-row bounds for the compressed matrices are
\[
 \norm{W_{64}}<8.926,\qquad
 \norm{H_{\rm np,64}}<4.176,\qquad
 \norm{H_{\rm pr,64}}<5.490.
\]
Applying these bounds to the full interval enclosures, before rounding,
gives
\[
 2.5829853\cdot10^{-31}<QW_3(v,v)<2.5858415\cdot10^{-31}
\]
and preserves both strict $0.4$ bounds. A finite-support polynomial's
complete quadratic form equals its finite matrix quadratic form.
No invariant-subspace or full operator-residual assertion is required.
\end{proof}

Thus $QW_3(v,v)/q_{\rm np}(v)<7.5\cdot10^{-31}$ inside the exact
source complement. This rules out an explanation in which substantial
archimedean-plus-pole dominance eliminates the need for cancellation
everywhere off that source. It does not rule out positivity on the
complement, and exhibits neither a negative Weil direction nor an
additional exact null vector. Positive basis-vector or typical
random-vector energies cannot control these coherent linear combinations.
The pole term must be included explicitly: calling $q_{\rm np}$ the
pure archimedean contribution would change the asserted comparison.


\subsection*{Why a finite source block cannot reduce the unsigned prime norm}

The arithmetic bound in the preceding tail estimates is not merely a
poor choice of positive weight.  There is an exact recurrence obstruction
to replacing it by a small norm on the source complement.

\begin{proposition}[Prime norm survives every finite-dimensional removal]
\label{prop:v119-prime-essential}
Fix $\lambda>1$ and let $\mathsf T_{\rm pr}$ be the positive-coefficient
sum of truncated prime translations in
\eqref{eq:v116-weighted-prime-general}.  For every finite-rank
orthogonal projection $P$ on $L^2(0,L)$, $Q=I-P$, one has
\begin{equation}\label{eq:v119-prime-compression}
 \|Q\mathsf T_{\rm pr}Q\|=\|\mathsf T_{\rm pr}\|.
\end{equation}
For every compact operator $C$ on this space,
\begin{equation}\label{eq:v119-prime-compact}
 \|Q(\mathsf T_{\rm pr}-C)Q\|\ge\|\mathsf T_{\rm pr}\|.
\end{equation}
In particular the essential norm of $\mathsf T_{\rm pr}$ equals its
ordinary norm.  The projection need not be a Fourier projection and
may remove any prescribed finite source block.
\end{proposition}
\begin{proof}
For each positive integer $k$, let
$M_k f(x)=e^{2\pi ikx/L}f(x)$.  Simultaneous recurrence of the finitely
many phases gives integers $k_j\to\infty$ such that
\[
 e^{2\pi ik_j\log m/L}\longrightarrow1
 \quad\hbox{for every active }m.
\]
For completeness, partition the unit cube in the finitely many phase
coordinates into $H^d$ cubes and compare $H^d+1$ successive multiples.
Two lie in one cube, giving a positive return with coordinate errors
at most $1/H$.  If such returns have an unbounded subsequence, use it;
otherwise a fixed positive return has all coordinate errors zero,
and its increasing multiples give the required sequence.  No rational
independence of prime logarithms is assumed.

Conjugating a truncated translation by $M_k$ only multiplies it by
its corresponding phase.  Since each translation has norm at most one,
\begin{equation}\label{eq:v119-prime-recurrence}
 \|M_{k_j}^*\mathsf T_{\rm pr}M_{k_j}-\mathsf T_{\rm pr}\|
 \le2\sum_m w_m|e^{2\pi ik_j\log m/L}-1|\longrightarrow0.
\end{equation}
For every fixed $f\in L^2(0,L)$, $M_{k_j}f\rightharpoonup0$:
each scalar product is a Fourier coefficient of an $L^1$ function.
Thus $PM_{k_j}f\to0$, $PM_{k_j}\mathsf T_{\rm pr}f\to0$, and
$CM_{k_j}f\to0$ strongly.  Combining these statements with
\eqref{eq:v119-prime-recurrence} gives
\[
 Q(\mathsf T_{\rm pr}-C)Q M_{k_j}f
       -M_{k_j}\mathsf T_{\rm pr}f\longrightarrow0.
\]
For unit $f$, the input has norm at most one, so taking norms and then
the supremum over $f$ proves \eqref{eq:v119-prime-compact}.
Putting $C=0$ and using compression contractivity proves
\eqref{eq:v119-prime-compression}.  Taking the infimum over compact
$C$, with $P=0$, proves the essential-norm assertion.
\end{proof}

\begin{proposition}[The sharp growing-window scale of the unsigned prime norm]
\label{prop:v119-prime-scale}
With the actual weights $w_m=\Lambda(m)/\sqrt m$ and $L=2\log\lambda$,
\begin{equation}\label{eq:v119-prime-scale}
 \|\mathsf T_{\rm pr}\|=(1+o(1))\lambda
       \qquad(\lambda\longrightarrow\infty).
\end{equation}
Consequently \eqref{eq:v119-prime-scale} also holds for
$\|Q_\lambda\mathsf T_{\rm pr}Q_\lambda\|$ for any family of
finite-rank orthogonal projections $I-Q_\lambda$, with no restriction on their
finite ranks or Fourier cutoffs.
\end{proposition}
\begin{proof}
Write
\[
 a(x)=e^{-x/2},\quad b(x)=e^{-(L-x)/2},\quad\phi(x)=a(x)+b(x),
\]
and
\[
 \Psi(X)=\sum_{1<m<X}\Lambda(m),\quad
 S(X)=\sum_{1<m<X}\frac{\Lambda(m)}m.
\]
The strict cutoff agrees with the truncated translations almost
everywhere.  Direct evaluation gives
\begin{equation}\label{eq:v119-prime-weight-ratio}
 \frac{(\mathsf T_{\rm pr}\phi)(x)}{\phi(x)}
 =\frac{a(x)\{S(e^{L-x})+\Psi(e^x)\}
       +b(x)\{\Psi(e^{L-x})+S(e^x)\}}{a(x)+b(x)}.
\end{equation}
The exact identity
\[
 a(x)e^x+b(x)e^{L-x}=\lambda\{a(x)+b(x)\}
\]
identifies the leading constant.  The classical prime number
theorem~\cite{NISTPrimeAsymptotics}, equivalently $\Psi(X)\sim X$,
implies that for each $\delta>0$ there is $C_\delta<\infty$ with
\[
 |\Psi(X)-X|\le\delta X+C_\delta\qquad(X\ge1).
\]
The Chebyshev bound and partial summation give
$0\le S(X)\le C(1+\log X)$, with an absolute $C$.
Therefore, uniformly for $0<x<L$,
\[
 (1-\delta)\lambda-C_\delta
 \le\frac{\mathsf T_{\rm pr}\phi}{\phi}
 \le(1+\delta)\lambda+C_\delta+C(1+L).
\]
At each fixed window $\phi$ is bounded and bounded away from zero.
The upper estimate and the weighted Schur inequality
\eqref{eq:v116-weighted-prime-general} bound the operator norm from
above.  Integrating the lower estimate against $\phi^2$ gives the
same lower bound for its Rayleigh quotient and hence for the norm.
Divide by $\lambda$, let $\lambda\to\infty$, and then let
$\delta\downarrow0$. This proves \eqref{eq:v119-prime-scale}.
The projection assertion follows from
Proposition~\ref{prop:v119-prime-essential} at each window.
\end{proof}

These statements concern the unsigned prime operator, not the sign of
the full Weil form.  They make precise a limitation of the particular
place-by-place bound \eqref{eq:v115-sharp-tail}. Even replacing its
prime constant by the exact tail norm cannot certify positivity unless
\begin{equation}\label{eq:v119-unsigned-cutoff}
 \log\frac{N+1}{L}>\|\mathsf T_{\rm pr}\|
       =(1+o(1))\lambda.
\end{equation}
Here we refer to the whole-space bound, whose remaining subtracted
terms are nonnegative. Thus this method still requires an exponential
Fourier cutoff in $\lambda$, rather than the polynomial source cutoff.
This is a limitation of that sufficient scalar bound, not a necessary
cutoff for actual Weil positivity or for a frequency-weighted argument.

The physical pole operator has rank at most two.  Equation
\eqref{eq:v119-prime-compact} shows that combining it with the prime
operator does not make their ordinary operator norm smaller than
$\|\mathsf T_{\rm pr}\|$, even after a finite source removal.
The recurrence frequencies in the proof may be arbitrarily large;
there the archimedean logarithmic energy also grows.  No negative
direction for the complete Weil operator is inferred.  A scalable G2
estimate must control the arithmetic term together with that energy
or retain more precise signed correlations, rather than assert a small
unweighted norm on the complement.  The required signed estimate
remains open.

\subsection*{Retaining the archimedean energy in a one-sided estimate}

Choose a literal Fourier tail on which $a_n:=-A_n>0$ and write
$D=\operatorname{diag}(a_n)$ there. By
\eqref{eq:v114-weil-decomposition}, the compressed closed form has
$\mathsf W=D+R$, with $R$ bounded self-adjoint. It contains the
actual signed prime and pole terms and the archimedean off-diagonal;
no arithmetic diagonal is absorbed into $D$. Put
\begin{equation}\label{eq:v120-weighted-remainder}
 \mathcal B=D^{-1/2}R D^{-1/2}.
\end{equation}
The operator $\mathcal B$ is compact self-adjoint, since
$a_n\sim\log|n|$ at the fixed window. The form identity is
\[
 q_{\mathsf W}[f]
 =\ip{(I+\mathcal B)D^{1/2}f}{D^{1/2}f}.
\]
Thus the exact positivity target is $\mathcal B\succeq-I$;
the stronger two-sided condition $\|\mathcal B\|\le1$ can fail
because of a harmless positive direction.

\begin{proposition}[The scale of an approximate weighted sign]
\label{prop:v120-weighted-error}
Suppose $D\succeq d_0I>0$, $R\succeq-CI$, $C\ge0$, and
$D^{-1/2}RD^{-1/2}\succeq-(1+\eta)I$, $\eta\ge0$.
Then, on the entire form domain,
\begin{equation}\label{eq:v120-weighted-error}
 D+R\succeq-\frac{\eta C}{1+\eta}I.
\end{equation}
The constant is sharp. In particular, a dimensionless error
$\eta_\lambda\to0$ is not by itself an ordinary lower error
$o(1)$ when $C_\lambda$ grows.
\end{proposition}
\begin{proof}
The hypotheses give $D+R\succeq-\eta D$ and
$D+R\succeq D-CI$. Add these inequalities with weights
$1/(1+\eta)$ and $\eta/(1+\eta)$ respectively. For $C>0$,
the scalar example $D=C/(1+\eta)$, $R=-C$ attains equality.
Taking $C_\lambda=\lambda$ and $\eta_\lambda=1/\lambda$
shows why the missing scale factor cannot be omitted. A sufficient
conversion requires $\eta_\lambda C_\lambda\to0$, or a more
precise directional energy estimate.
\end{proof}

There is also a limitation on how the compact weighted remainder can
be estimated. A Hilbert--Schmidt budget for its whole infinite matrix
is unavailable, despite the valid finite-column residual Grams above.

\begin{proposition}[Weighted compactness without a square-summable matrix]
\label{prop:v120-not-schatten}
At any fixed $\lambda>\sqrt2$, let $D$ be a positive diagonal on a
fixed Fourier tail with $a_n\sim\log|n|$. The compact operator
$D^{-1/2}\mathsf T_{\rm pr}D^{-1/2}$ belongs to no Schatten class
$\mathcal S_p$, $0<p<\infty$. For the exact archimedean diagonal
$a_n=-A_n$, the same assertion holds for
$\mathcal B$ in \eqref{eq:v120-weighted-remainder}.
\end{proposition}
\begin{proof}
The prime diagonal is the nonzero finite trigonometric sequence
\[
 \tau_n=2\sum_{1<m<\lambda^2}
 \frac{\Lambda(m)}{\sqrt m}
 \left(1-\frac{\log m}{L}\right)
 \cos\frac{2\pi n\log m}{L}.
\]
Collect equal unit-circle frequencies to write
$\tau_n=\sum_{z\in F}c_z z^n$, with all $c_z>0$.
Finite geometric sums give
\[
 \frac1X\sum_{n=1}^X|\tau_n|^2\longrightarrow
 \sum_{z\in F}c_z^2>0.
\]
Boundedness then implies that $|\tau_n|$ exceeds a positive
constant on a set of positive lower density. Consequently, for each
$p\ge1$, $\sum_n|\tau_n/a_n|^p=\infty$: the partial sum on
that set is bounded below by a constant times $X/(\log X)^p$.
For self-adjoint $K\in\mathcal S_p$, Jensen's inequality in its
spectral measure gives
$\sum_n|\ip{KU_n}{U_n}|^p\le\operatorname{Tr}|K|^p$.
This is a contradiction. The inclusions
$\mathcal S_p\subset\mathcal S_1$ for $0<p<1$ settle the
remaining exponents. For the complete remainder its diagonal is
$(-\tau_n+O_\lambda(n^{-2}))/a_n$, since the exact pole diagonal
has this decay and the archimedean remainder has zero diagonal.
The same argument applies. None of this invalidates the separated
finite-column moment bounds, whose off-diagonal factor
$(n-m)^{-1}$ is square summable in the remote row index.
\end{proof}

\begin{proposition}[Sharp fixed-window decay of the weighted prime tail]
\label{prop:v120-weighted-tail-asymptotic}
Under the diagonal assumptions above, let $Q_J$ remove $|n|\le J$.
At each fixed window,
\begin{equation}\label{eq:v120-weighted-tail-asymptotic}
 \lim_{J\to\infty}(\log J)
 \|Q_JD^{-1/2}\mathsf T_{\rm pr}D^{-1/2}Q_J\|
       =\|\mathsf T_{\rm pr}\|.
\end{equation}
Subtracting any fixed compact operator from $\mathsf T_{\rm pr}$
before weighting leaves this limit unchanged. No uniformity in
$\lambda$ is asserted.
\end{proposition}
\begin{proof}
The upper bound is
$\|\mathsf T_{\rm pr}\|/\inf_{|n|>J}a_n$.
For the lower bound, fix $\epsilon>0$ and choose a unit Fourier
polynomial $f$, supported on $|m|\le M$, with
$|\ip{\mathsf T_{\rm pr}f}{f}|>\|\mathsf T_{\rm pr}\|-\epsilon$.
Returns of the finite vector of shift phases to any fixed neighborhood
of the identity are relatively dense. Indeed its cyclic closure is a
compact group; finitely many integer translates of that neighborhood
cover it. The largest difference between the corresponding translation
indices bounds the gap between return times. This also covers rational
phase relations.

Choose the neighborhood so that the conjugation error in
\eqref{eq:v119-prime-recurrence} is below $\epsilon$.
For every large $J$ there is therefore a return
$k\in[J+M+1,J+M+1+K]$, with $K$ independent of $J$.
Set $g=M_kf$ and $x=D^{1/2}g$. Both are supported outside $J$,
$\|x\|^2=(1+o(1))\log J$, and the weighted Rayleigh numerator
at $x$ equals $\ip{\mathsf T_{\rm pr}g}{g}$, whose absolute value
is greater than $\|\mathsf T_{\rm pr}\|-2\epsilon$.
Let $J\to\infty$ and then $\epsilon\downarrow0$.
For compact $C$ the change in the left side is at most
\[
 \frac{\log J}{\inf_{|n|>J}a_n}\|Q_JCQ_J\|\longrightarrow0.
\]
The return-length bound $K$ depends on the window and the tolerance.
Thus one may not combine \eqref{eq:v120-weighted-tail-asymptotic}
with $\|\mathsf T_{\rm pr}\|\sim\lambda$ to infer a bound on
an arbitrary joint scale $J=J(\lambda)$.
\end{proof}

\begin{proposition}[A signed energy bound at the larger window $\lambda=4$]
\label{prop:v120-lambda4-tail}
Let $L=2\log4$, retain the exact diagonal $a_n=-A_n$, and let
$Q_{16}$ be the literal projection onto $|n|>16$. For every complex
vector in its complete closed form domain,
\begin{equation}\label{eq:v120-lambda4-tail}
 QW_4(f,f)\ge10^{-8}\sum_{|n|>16}a_n|f_n|^2,
 \qquad f=Q_{16}f.
\end{equation}
Every weight in this sum is positive. This is an infinite-tail
computer-assisted bound, not a claim about the remaining low modes
or about an unbounded family of windows.
\end{proposition}
\begin{proof}
Apply the directional verified-solve proof to
$\mathsf W_4-cD$, $c=10^{-8}$, on $Q_{16}$.
Only the archimedean diagonal is changed, to $(1-c)a_n$;
the exact off-diagonal coefficients $b_n$, both pole signs, and every
prime power $2,3,4,5,7,8,9,11,13$ remain. The full-length $m=16$
translation is zero. The positive physical weight
$\phi(x)=e^{-x/2}+e^{-(L-x)/2}$ gives the certified Schur bound
\begin{equation}\label{eq:v120-lambda4-prime-bound}
 \|\mathsf T_{\rm pr}\|
 \le M_\phi<4.624042316766478.
\end{equation}
To verify it, set $u=e^x$. On each interval between
$1,16,m,16/m$ the ratio $(\mathsf T_{\rm pr}\phi)/\phi$
is a linear-fractional function of $u$, hence monotone or constant.
Every one-sided endpoint value is evaluated with outward rounding;
the chosen upper endpoint is checked against every candidate.

Use $E_L,C_L,h_L$ from \eqref{eq:v115-arch-errors}, with
$C_L=1$, $h_L=9/16$ and $R_L=64/255$. For $n>N$ the proved
remote inverse bound applies with
\[
 g_n=(1-c)\{\log(n/L)-E_L(2\pi n/L)\}-\kappa,
\]
where $\kappa=M_\phi+2C_L/t_*$ in the even sector and
$\kappa=M_\phi+2C_L/t_*+\pi/2+4h_LL/(\pi^2N)$ in the odd
sector, $t_*=2\pi(N+1)/L$. These weights are positive and
increasing for the choices below. Also the same diagonal estimate
at $n=17$ proves positivity of $a_n$ on the whole claimed tail.

The finite head now consists of $17\le n\le N$ in the normalized
parity basis; it does not include the omitted physical low modes.
The solve witness has support $N<n\le M$. With
$G=I-Z$, form the exact signed $K_Z$ and every residual row up to
$J$, and use the moment majorant $U_J$ beyond $J$. The lower matrix is
\begin{equation}\label{eq:v120-lambda4-lower}
 \mathcal L=K_Z-
  \sum_{N<n\le J}\frac{R_n^*R_n}{g_n}-\frac{U_J}{g_{J+1}}.
\end{equation}
The same full-index parity conversion and both infinite tails as in
Proposition~\ref{prop:v116-remote-moment} are used. Bounding its
geometric-remainder constant by including the absent indices
$1,\ldots,16$ is conservative. The certificate parameters are
\begin{center}
\begin{tabular}{lrrrrr}
\toprule
sector&head size&$N$&$M$&$J$&moment order\\
\midrule
even&496&512&1024&4096&48\\
odd&1520&1536&2048&4096&64\\
\bottomrule
\end{tabular}
\end{center}
Both $Z$ and an inverse-Cholesky congruence witness $C$ are frozen
to exact dyadic entries. Their numerical method of construction has
no proof role: outward interval evaluation verifies every row of
$C\mathcal L C^*$ has positive diagonal minus the sum of the
absolute off-diagonal entries. Thus $C\mathcal L C^*\succ0$
and $\mathcal L\succ0$. Square completion with the positive
remote operator proves \eqref{eq:v120-lambda4-tail} separately
in each parity sector and hence for arbitrary complex vectors.

The saved witnesses, analytic series radii, all finite residual rows
and the infinite moment bound are reproduced by
\texttt{assembly\_general.py} and
\texttt{certify\_weighted\_tail.py}. Replaying the same witnesses
at higher precision verifies the same strict inequalities.
These are internal computer-assisted certificates, not independent
external verification. No zero locations or RH assumptions enter them.
\end{proof}

\begin{proposition}[The two-sided weighted contraction fails at this cutoff]
\label{prop:v120-norm-counterexample}
For the same exact $D$ and $Q_{16}$ at $\lambda=4$,
\[
 \|D^{-1/2}Q_{16}(\mathsf W_4-D)Q_{16}D^{-1/2}\|>1.
\]
This does not contradict Proposition~\ref{prop:v120-lambda4-tail}.
\end{proposition}
\begin{proof}
The archived exact dyadic even vector $v$, supported on
$17\le|n|\le256$, has the rigorously enclosed quotient
\[
 2.05254409540345<\frac{QW_4(v,v)}{\ip{Dv}{v}}
                         <2.05254409540346.
\]
Its positive numerator excess over $2\ip{Dv}{v}$ exceeds
$0.0525$ in the saved normalization. The same analytic matrix
assembly, independent of the floating eigenvector that suggested
$v$, proves these interval inequalities. Testing the weighted
remainder on $D^{1/2}v$ gives a Rayleigh value greater than one.
This is a positive direction, not a negative Weil direction.
\end{proof}

The larger-window tail sign reduces full positivity at $\lambda=4$
to the signed effective matrix on
$E_{16}$:
\[
 F-B^*T^{-1}B,\qquad
 F=P_{16}\mathsf W_4P_{16},\quad B=Q_{16}\mathsf W_4P_{16},
 \quad T=Q_{16}\mathsf W_4Q_{16}.
\]
It has $17$ even and $16$ odd coordinates.
Proposition~\ref{prop:v125-even-complete} below settles the complete
even sector; Proposition~\ref{prop:v126-full-window} completes the
odd sector. It is not the raw finite
Weil matrix: the inverse-tail correction must remain.
The next growing-window estimate must also satisfy the ordinary-error
scale in Proposition~\ref{prop:v120-weighted-error}.
A related recent preprint~\cite{Zhu2026} studies certified
compact-window positivity and a pointwise prime-comb envelope barrier.
Our conclusion here concerns the actual Fourier-tail form and does
not use that preprint's computational certificates. Neither fixed-window
result supplies a uniform signed arithmetic estimate.

\subsection*{Reusing the signed tail certificate in the low-mode correction}
The scalar conclusion $T\succeq10^{-8}D$ discards most of the
information in Proposition~\ref{prop:v120-lambda4-tail}. The following
inverse bound retains its finite signed factorization. All adjoints
below use the fixed first-slot-linear convention.

\begin{proposition}[A structured inverse bound for a certified tail]
\label{prop:v120-structured-inverse}
Let a closed positive tail operator have a finite-head decomposition
\[
 T=\begin{pmatrix}F&B^*\\ B&U\end{pmatrix},
 \qquad U\succeq D_g\succeq\gamma I>0,
\]
where $B$ is bounded and $D_g$ is a positive diagonal operator.
Let $Z$ map the finite head into $\mathcal D(U)$, and set
\[
 R=B-UZ,\qquad K_Z=F-B^*Z-Z^*B+Z^*UZ.
\]
Suppose a finite Hermitian matrix $\underline L$ and invertible $C$
satisfy
\begin{equation}\label{eq:v120-inner-inverse-certificate}
 0\prec\underline L\preceq K_Z-R^*D_g^{-1}R,
 \qquad C\underline L C^*\succeq\mu I,\quad\mu>0.
\end{equation}
Then $T$ is bounded below by a strictly positive constant and, for
any tail right-hand side $r=(h,k)$,
\begin{equation}\label{eq:v120-structured-inverse}
 \ip{T^{-1}r}{r}
 \le \norm{D_g^{-1/2}k}^2+
 \mu^{-1}\norm{C(h-Z^*k-R^*D_g^{-1}k)}^2.
\end{equation}
The same statement holds as a Gram inequality for finitely many
right-hand sides simultaneously.
\end{proposition}
\begin{proof}
Write the original variable as $(x,y-Zx)$. Its exact form is
\[
 \ip{K_Zx}{x}+2\Re\ip{Rx}{y}+q_U[y],
\]
which is bounded below by
\[
 \ip{\underline Lx}{x}
       +\norm{D_g^{1/2}(y+D_g^{-1}Rx)}^2.
\]
The first change preserves the original form domain because
$Zx\in\mathcal D(U)$. The displayed lower form is defined on the
possibly larger domain consisting of the finite head and
$\mathcal D(D_g^{1/2})$. Its triangular change is bounded and
invertible there because $D_g^{-1}Rx\in\mathcal D(D_g)$ and the
head is finite dimensional. Thus that lower form is closed and
strictly positive; the original form dominates it on its own domain.
In its inverse variational formula the head part of the linear
functional is $h-Z^*k-R^*D_g^{-1}k$. Consequently its inverse
quadratic form equals the first term of
\eqref{eq:v120-structured-inverse} plus the quadratic form of
$\underline L^{-1}$ on that head vector. Inverse order and
\eqref{eq:v120-inner-inverse-certificate} give
$\underline L^{-1}\preceq\mu^{-1}C^*C$. This proves the bound.
Applying it to every linear combination of the right-hand sides
proves the Gram version. The argument concerns closed forms and
does not replace the logarithmic operator by a bounded matrix.
\end{proof}

If the saved factors instead certify $T-cD$, with $cD\succeq0$,
they still bound $T^{-1}$ because
$T^{-1}\preceq(T-cD)^{-1}$. The factors and residual in the
right side must then be those of $T-cD$ throughout; a shifted
residual is not an unshifted residual. In particular, the existing
$\lambda=4$ certificates are usable without losing their directional
information to the coarse upper bound $10^8D^{-1}$.

Here is a finite-row version that also retains every omitted mode.
When $Z$ has finite support, take $J$ past that support and split the inner far-tail rows
at $J$. For a matrix of right-hand sides define
\[
 H_J=h-Z^*k-\sum_{n\le J}R_n^*k_n/g_n,
 \qquad V_J=\sum_{n\le J}k_n^*k_n/g_n.
\]
The sums run over the inner far tail in normalized parity coordinates;
both Fourier signs are included by that decomposition. Suppose certified positive
matrices satisfy
\[
 \sum_{n>J}R_n^*R_n/g_n\preceq U_R,
 \qquad \sum_{n>J}k_n^*k_n/g_n\preceq U_k,
 \qquad \rho\ge\norm{CU_RC^*}.
\]
Then, for every $\tau>0$,
\begin{equation}\label{eq:v120-structured-finite-rows}
 r^*T^{-1}r\preceq V_J+U_k+
 \frac{(1+\tau)H_J^*C^*CH_J+
             (1+\tau^{-1})\rho U_k}{\mu}.
\end{equation}
Indeed, if $E=\sum_{n>J}R_n^*k_n/g_n$, then
$E^*C^*CE\preceq\rho U_k$ by the operator Cauchy--Schwarz
inequality. Apply the matrix Young inequality to $C(H_J-E)$.
The trace of $CU_RC^*$ is a permissible choice of $\rho$;
its possible growth with the inner head dimension must then be
retained. These are finite-column remote Grams, as in
Proposition~\ref{prop:v116-remote-moment}, not an invalid
Hilbert--Schmidt estimate for the whole weighted Weil operator.
A joint remote Gram can be sharper than separating this mixed term.

\begin{proposition}[A convergent arithmetic refinement of the far inverse]
\label{prop:v120-inverse-polynomial}
If $D_g\preceq U\preceq mD_g$, $D_g\succeq\gamma I>0$,
$m>1$, let $A=D_g^{-1/2}UD_g^{-1/2}$ denote the bounded
operator represented by the preconditioned form, and set
\[
 S=I-A/m,\qquad
 P_k(A)=\frac1m\sum_{j=0}^{k-1}S^j+S^k,\quad k\ge0.
\]
The empty sum gives $P_0=I$. Then
\begin{equation}\label{eq:v120-inverse-polynomial}
 A^{-1}\preceq P_{k+1}(A)\preceq P_k(A)\preceq I,
 \qquad
 0\preceq P_k(A)-A^{-1}\preceq(1-1/m)^k I.
\end{equation}
Thus $D_g^{-1/2}P_k(A)D_g^{-1/2}$ are decreasing upper
bounds for $U^{-1}$, with error at most
$(1-1/m)^kD_g^{-1}$ in form order. In particular,
\begin{equation}\label{eq:v120-first-inverse-refinement}
 U^{-1}\preceq D_g^{-1}
       -m^{-1}D_g^{-1}(U-D_g)D_g^{-1}.
\end{equation}
Products involving $U$ in this formula are understood through the
bounded preconditioned form.

For the actual far operators used in
Proposition~\ref{prop:v120-lambda4-tail}, including their shift
$\mathsf W_4-10^{-8}D$, one can take
\[
 m=67\quad\hbox{on the even tail }n>512,\qquad
 m=335\quad\hbox{on the odd tail }n>1536.
\]
These bounds cover all omitted Fourier modes at this fixed window.
\end{proposition}
\begin{proof}
The spectral theorem gives $I\preceq A\preceq mI$ and
$0\preceq S\preceq(1-1/m)I$. The geometric identity is
\[
 P_k(A)-A^{-1}=S^k(I-A^{-1}).
\]
Its factors commute and are positive. This proves the error bound;
$P_k-P_{k+1}=m^{-1}S^k(A-I)\succeq0$ proves monotonicity.
Conjugate by $D_g^{-1/2}$, using the inverse identity for the
positive closed form $U$. The case $k=1$ is
\eqref{eq:v120-first-inverse-refinement}.

To verify the two concrete constants, first strengthen
\eqref{eq:v115-diagonal-sharp} to
\begin{equation}\label{eq:v120-two-sided-arch}
 |a_n-\log(|n|/L)|\le E_L(|t_n|),\qquad n\ne0.
\end{equation}
For the upper bound use the same Binet formula at
$w=5/4+it/2$. It gives
\[
 \Re\log w-\log(t/2)
 =\tfrac12\log(1+25/(4t^2))\le25/(8t^2),
\]
while $-\Re(1/(2w))-\Re(1/z)\le0$, $z=1/4+it/2$.
The Binet remainder is at most $1/(15t)$ in absolute value.
Because $25/8<7/2$, the existing absolute trigamma and exponential
series bounds give the same $E_L$ for the upper error.

Write $\mathsf W_4=D+R_0$. From
\eqref{eq:v114-b-d}--\eqref{eq:v114-b-bound}, a valid bound is
\[
 \norm{R_0}\le C_R:=2\pi B_*+32h_L/L+2P_4.
\]
Here $P_4$ may omit the zero full-length translation $m=16$.
The pole diagonal bound is $32h_L/L$; its factor $L$ is essential.
The commutator bound is taken before the parity compression, so it
also bounds either normalized parity sector. With $c=10^{-8}$,
$g_n=(1-c)\{\log(n/L)-E_L(t_n)\}-\kappa$, and
$t_*=2\pi(N+1)/L$, the already proved lower bound and
\eqref{eq:v120-two-sided-arch} imply
\[
 D_g\preceq U\preceq
 D_g+\{\kappa+2(1-c)E_L(t_*)+C_R\}I.
\]
Since $g_n$ is positive and increasing, it suffices that
\begin{equation}\label{eq:v120-inverse-m-bound}
 m\ge1+\frac{\kappa+2(1-c)E_L(t_*)+C_R}{g_{N+1}}.
\end{equation}
Outward interval evaluation gives $C_R<34.450462$. The right side
of \eqref{eq:v120-inverse-m-bound} is below $66.76640$ for the
even $N=512$ choice, and below $334.71708$ for the odd $N=1536$
choice. The same prime bound, diagonal errors and inverse weights
as the tail certificate are used. Strict interval checks at two
precisions verify the stated integers.
\end{proof}

For a matrix of residuals $Y$, the error of the $k$th inverse bound
is at most
\[
 (1-1/m)^k Y^*D_g^{-1}Y.
\]
Along growing windows this part of the error is therefore controlled
if $\norm{D_{g,\lambda}^{-1/2}Y_\lambda}^2
 e^{-k_\lambda/m_\lambda}\to0$, with the actual values of
$m_\lambda$ retained. The explicit integers above apply only at
$\lambda=4$. Proposition~\ref{prop:v122-parity-inverse} below improves these
constants to $20$ and $90$ using the actual parity signs. Evaluating
the signed powers and their complete remote contributions, and
proving the remaining low-head sign, remain separate obligations. The inverse iteration proves neither of those signs.


\subsection*{Parity signs and certified evaluation of the inverse refinement}
The global remainder norm in the preceding inverse estimate loses useful
signs. Retaining them improves the complete far-operator sandwich;
it does not determine the sign of the low Schur matrix.

\begin{proposition}[Sharper inverse constants at the fixed window]
\label{prop:v122-parity-inverse}
For the actual shifted far operators and the unchanged weights $D_g$
in Proposition~\ref{prop:v120-lambda4-tail}, one has
\[
 D_g\preceq U\preceq20D_g\quad\text{in the even sector }n>512,
 \qquad
 D_g\preceq U\preceq90D_g\quad\text{in the odd sector }n>1536.
\]
Thus the constants in Proposition~\ref{prop:v120-inverse-polynomial}
can be replaced by $20$ and $90$, respectively, with error factors
$(19/20)^k$ and $(89/90)^k$. These are complete-operator bounds
at $\lambda=4$, with no assertion of growing-window uniformity.
\end{proposition}
\begin{proof}
Put $h=9/16$, $c=10^{-8}$, $L=2\log4$ and
$t_N=2\pi(N+1)/L$. Let $M_\phi$ be the certified prime bound
in \eqref{eq:v120-lambda4-prime-bound}. The limiting archimedean
off-diagonal in the normalized parity basis is the matrix
$1/[2(n+m)]$ in the even sector and its negative in the odd sector.
It is positive and has norm at most $\pi/2$; the remaining
archimedean commutator has norm at most $2/t_N$, by the proof of
Proposition~\ref{prop:v115-sharp-tail}.

In centered coordinates the pole form is
$2|\ip f{\cosh(x/2)}|^2-2|\ip f{\sinh(x/2)}|^2$.
Only the positive term acts on even vectors, and only the negative
term on odd vectors. Direct integration gives
\[
 |\widehat{\cosh(x/2)}(n)|^2
 =\frac{h}{L(t_n^2+1/4)^2},\qquad
 2\|Q_N\cosh(x/2)\|^2
 \le\frac{hL^3}{4\pi^4}\sum_{n>N}n^{-4}
 \le\frac{hL^3}{12\pi^4N^3}.
\]
The factor two for the form and both Fourier signs are included.
Consequently the remainder above the shifted archimedean diagonal
has the following upper bounds:
\begin{equation}\label{eq:v122-parity-upper}
 \begin{split}
 C_{\rm e}(N)&=M_\phi+2/t_N+\pi/2+hL^3/(12\pi^4N^3),\\
 C_{\rm o}(N)&=M_\phi+2/t_N.
 \end{split}
\end{equation}
The odd limiting Hilbert matrix and odd pole term are nonpositive
and may be omitted in an upper bound. The shift $cD$ changes only
the exact archimedean diagonal; none of these remainder terms is
multiplied by $1-c$.

For the same $\kappa$ and increasing $g_n$ as before, the two-sided
diagonal estimate \eqref{eq:v120-two-sided-arch} now yields
\[
 U\preceq D_g+
 \{\kappa+2(1-c)E_L(t_N)+C_{\rm e/o}(N)\}I.
\]
It suffices to bound one plus the braced quantity divided by
$g_{N+1}$. Outward interval evaluation gives values below
$19.215607$ for even $N=512$, and below $89.840960$ for odd
$N=1536$. The integer gates $20$ and $90$ pass at both 320 and
384 bits. The lower bounds and closed form domains are unchanged.
\end{proof}

\begin{proposition}[Complete residual bounds from a finite prefix]
\label{prop:v122-prefix-refinement}
Assume $D_g\preceq U\preceq mD_g$, with
$D_g\succeq\gamma I>0$. Put
$A=D_g^{-1/2}UD_g^{-1/2}$, $H=A-I$, and let $P$ be an
orthogonal prefix projection, $Q=I-P$. Suppose
$QHQ\preceq\nu Q$ with $\nu\ge0$. For a residual $k$ write
\[
 y=D_g^{-1/2}k=p+q,\quad p=Py,\quad q=Qy,\qquad
 v=\norm p^2,\quad a=\ip{Hp}{p},\quad \norm q^2\le E.
\]
With $[s]_+=\max(s,0)$, the following scalar bound includes
every omitted residual row:
\begin{equation}\label{eq:v122-prefix-scalar}
 \ip{U^{-1}k}{k}
 \le v+E-\frac1m[\sqrt a-\sqrt{\nu E}]_+^2.
\end{equation}
For a matrix of residual columns $Y=D_g^{-1/2}\mathcal R$, put
$Y_P=PY$, $Y_Q=QY$, $V=Y_P^*Y_P$ and
$A_P=Y_P^*HY_P$. If $Y_Q^*Y_Q\preceq\mathcal E$, then, for
every fixed $0<\theta\le1$,
\begin{equation}\label{eq:v122-prefix-matrix}
 \mathcal R^*U^{-1}\mathcal R\preceq
 V+\mathcal E-\frac1m
 \{(1-\theta)A_P-(\theta^{-1}-1)\nu\mathcal E\}.
\end{equation}
If a certified $0\le\rho<1$ instead satisfies
$\nu\mathcal E\preceq\rho^2A_P$, the braced lower bound can
be replaced by $(1-\rho)^2A_P$.
\end{proposition}
\begin{proof}
The bounded operator $H$ is positive, so the reverse triangle
inequality for $H^{1/2}p+H^{1/2}q$ gives
\[
 \ip{Hy}{y}\ge
 [\sqrt a-\|H^{1/2}q\|]_+^2
 \ge[\sqrt a-\sqrt{\nu E}]_+^2.
\]
Combine this with the first inverse polynomial
$A^{-1}\preceq I-H/m$ and
$\norm y^2=\norm p^2+\norm q^2\le v+E$.
For the matrix version, apply the lower Young inequality to
$H^{1/2}Y_Px+H^{1/2}Y_Qx$ for every head vector $x$.
Since $1-\theta^{-1}\le0$, the remote upper bound supplies the
stated lower form. The relative variant follows from the scalar
reverse triangle inequality in every head direction. No factor
equal to the number of columns is inserted. The scalar positive
part in \eqref{eq:v122-prefix-scalar} is not a matrix positive-part
operation; the latter does not follow from these arguments.
\end{proof}

In the actual parity tail based at $N$, the prefix can end at any
$J>N$. The same proof as \eqref{eq:v122-parity-upper} gives the
complete remote compression bound
\begin{equation}\label{eq:v122-remote-H}
 Q_JHQ_J\preceq\nu_J Q_J,\qquad
 \nu_J=\frac{\kappa_N+2(1-c)E_L(t_J)+C_{\rm e/o}(J)}{g_{J+1}}.
\end{equation}
Here $t_J=2\pi(J+1)/L$, whereas the constant $\kappa_N$ in
the original $D_g$ is kept fixed. The upper numerator decreases
and the positive denominator increases with the index. The moment
bound of Proposition~\ref{prop:v116-remote-moment}, divided by
$g_{J+1}$, supplies $\mathcal E$. Thus these formulas evaluate a
signed inverse improvement without setting the infinite residual
to zero. All occurrences of $U$ are interpreted through its
bounded preconditioned form.

\begin{proposition}[One complete far-residual comparison]
\label{prop:v122-directional-far}
At $\lambda=4$, take the archived exact dyadic even solve $X$
on modes $17,\ldots,256$, and the archived exact dyadic head
vector $v$ of length $17$ specified in
Appendix~\ref{app:v122-inverse}. Put
\[
 f=\binom{v}{-Xv},\qquad K=QW_4(f,f),\qquad
 k=Q_{512}\mathsf W_4 f.
\]
For the complete shifted even far operator
$U=Q_{512}(\mathsf W_4-10^{-8}D)Q_{512}$ one has
\begin{equation}\label{eq:v122-directional-far}
 \ip{U^{-1}k}{k}<5.758\cdot10^{-20}
 <6.469\cdot10^{-20}<K.
\end{equation}
This is one far-energy comparison. It does not include the mixed
term of \eqref{eq:v122-refined-structured}, and it is not a lower
bound for the full low Schur matrix.
\end{proposition}
\begin{proof}[Outward interval evaluation with a complete remote bound]
Use $m=20$ and $J=65536$ in
Proposition~\ref{prop:v122-prefix-refinement}. The frozen $f$ has
support through $256$. Thus its retained Weil energy and residual
rows are evaluated from the actual finite matrix coefficients,
without truncating $f$. The outer residual is unshifted; the far
inverse and its $D_g$ are those of the shifted tail certificate.
The following outward enclosures suffice:
\begin{center}
\begin{tabular}{@{}ll@{}}
\toprule Quantity & Certified enclosure\\\midrule
$K$ & $(6.4692843941,6.4692843942)\,10^{-20}$\\
$v_J=\|P_JD_g^{-1/2}k\|^2$ & $(6.9332416283,6.9332416284)\,10^{-20}$\\
$a_J=\ip{HP_JD_g^{-1/2}k}{P_JD_g^{-1/2}k}$
 & $(4.0338113722,4.0338113723)\,10^{-19}$\\
$\|Q_JD_g^{-1/2}k\|^2$ & $<3.4797080873\,10^{-21}$\\
$\|Q_JHQ_J\|$ & $<1.987323730$\\\bottomrule
\end{tabular}
\end{center}
The remote norm uses the complete order-$64$ moment majorant,
divided by $g_{J+1}$; the last row uses
\eqref{eq:v122-remote-H}. Substitution into
\eqref{eq:v122-prefix-scalar} gives the stricter computed upper
bound $5.757887789\cdot10^{-20}$, proving the displayed claim.
All gates replay with identical exact input witnesses at $768$ and
$896$ bits. The prefix matrix action uses exact Arb polynomial
products for its divided-difference and Hankel parts, as detailed
in the reproduction appendix. No floating-point sign decision,
zero-location assumption or omitted-residual cutoff enters the gate.
\end{proof}

\begin{proposition}[Refining both terms of the structured inverse]
\label{prop:v122-refined-structured}
Keep the factors $Z,R,C,\underline L,\mu$ of
Proposition~\ref{prop:v120-structured-inverse}, and assume also
$U\preceq mD_g$. Let $Q_*$ be a bounded positive operator such
that
\[
 U^{-1}\preceq Q_*\preceq D_g^{-1}.
\]
Then, for every $r=(h,k)$,
\begin{equation}\label{eq:v122-refined-structured}
 \ip{T^{-1}r}{r}
 \le\ip{Q_*k}{k}
       +\mu^{-1}\norm{C(h-Z^*k-R^*Q_*k)}^2.
\end{equation}
The same inequality holds for the simultaneous residual Gram.
In particular all inverse polynomials of
Proposition~\ref{prop:v120-inverse-polynomial} can be used while
retaining the original certified $C$ and $\mu$.
\end{proposition}
\begin{proof}
The assumed sandwich implies
$m^{-1}D_g^{-1}\preceq Q_*\preceq D_g^{-1}$.
Consequently $Q_*$ is injective and the closed form of $Q_*^{-1}$
has the same domain as $D_g^{1/2}$, with equivalent form norms.
Inverse order gives $U\succeq Q_*^{-1}$ as forms. Moreover
\[
 K_Z-R^*Q_*R\succeq K_Z-R^*D_g^{-1}R\succeq\underline L.
\]
Repeat the lower-form square completion in
Proposition~\ref{prop:v120-structured-inverse}, replacing the far
form by $Q_*^{-1}$ and its inverse by $Q_*$. Its exact inverse
contains $Q_*$ in both displayed terms, proving the assertion.
The change preserves the lower form domain since $Q_*Rx$ belongs
to the operator domain of $Q_*^{-1}$ for every finite head vector.
Applying the result to all linear combinations proves the Gram
version. If the saved factors certify $T-cD$, use inverse order
as before and retain those shifted factors throughout.
\end{proof}

Improving only the first term while leaving the old mixed term
unchanged is not justified. Positivity of the complete low Schur
matrix requires the resulting full correction, on the entire head,
to fit below $K_X$. The new estimates do not assert that comparison
or convergence along growing windows.

\begin{remark}[The simultaneous head acceptance test]
\label{rem:v122-head-ldl}
For a positive definite trial matrix $K_X$, put
$C_X=r^*T^{-1}r$. Then
\[
 \lambda_{\max}(K_X^{-1/2}C_XK_X^{-1/2})\le1
 \quad\Longleftrightarrow\quad K_X-C_X\succeq0.
\]
Thus an eigenvalue or inverse-square-root estimate is unnecessary.
It suffices to construct a Hermitian \emph{upper enclosure in form
order} $\overline C_X\succeq C_X$ for the complete residual Gram
and certify $K_X-\overline C_X\succeq0$, for example by interval
$LDL^*$ with strictly positive pivots. A fixed invertible dyadic
congruence may improve conditioning without changing the sign.
Strictly positive pivots prove the stronger strict comparison;
an exact zero pivot requires the corresponding semidefinite
factorization conditions, not an interval containing zero.
An entrywise upper bound is not by itself an upper bound in form
order. For $T_M=P_MT P_M$ and $r_M=P_Mr$, the variational identity
gives $r_M^*T_M^{-1}r_M\preceq C_X$, the opposite enclosure
direction. All mixed and remote contributions must be enclosed
before the $LDL^*$ test proves the complete Schur sign.
\end{remark}

\subsection*{Complete mixed correlations and a monotone head refinement}

The successful far-energy comparison in
Proposition~\ref{prop:v122-directional-far} leaves a mixed term.
The next estimate includes both the omitted input and the output
outside the retained Fourier rows. These are different contributions.

\begin{proposition}[A directional enclosure of the complete mixed term]
\label{prop:v123-mixed-tail}
Use the factors of Proposition~\ref{prop:v122-refined-structured},
and put
\[
 A=D_g^{-1/2}UD_g^{-1/2},\quad H=A-I,\quad
 Y=D_g^{-1/2}RC^*,\quad y=D_g^{-1/2}k,\quad d=C(h-Z^*k).
\]
Here $A$ is represented by its bounded closed form and
$0\preceq H\preceq(m-1)I$. Let $P$ be an orthogonal Fourier
prefix projection on the inner far tail, $Q=I-P$, $p=Py$ and
$q=Qy$. Suppose
\[
 \norm q^2\le E,\qquad 0\preceq QHQ\preceq\nu Q,
 \qquad 0\le\nu<m.
\]
Set $v=\norm p^2$, $a=\ip{Hp}{p}$ and
\[
 b=d-Y^*(I-H/m)y,\qquad
 b_P=d-(PY)^*(I-PHP/m)p.
\]
For any finite head vector $z$, let
$\eta_z=\ip{H PYz}{PYz}$ and suppose
$\norm{QYz}^2\le\rho_z$. Then
\begin{equation}\label{eq:v123-mixed-tail}
 |z^*(b-b_P)|\le
 \frac{\sqrt{\nu\eta_z E}+\sqrt{\nu a\rho_z}}{m}
       +\sqrt{\rho_z E}=:\epsilon_z.
\end{equation}
In addition the complete first-polynomial far energy satisfies
\begin{equation}\label{eq:v123-far-lower}
 \ip{(I-H/m)y}{y}\ge v-\frac{a}{m-\nu}.
\end{equation}
In particular, for $z\ne0$ the complete scalar-$\mu$ majorant
in \eqref{eq:v122-refined-structured} is at least
\begin{equation}\label{eq:v123-majorant-lower}
 v-\frac{a}{m-\nu}
 +\frac{[|z^*b_P|-\epsilon_z]_+^2}{\mu\norm z^2}.
\end{equation}
\end{proposition}
\begin{proof}
Write $f=Yz=f_P+f_Q$. The difference of the mixed pairings is
the sum of $m^{-1}f_P^*Hq$, $m^{-1}f_Q^*Hp$, and
$-f_Q^*(I-QHQ/m)q$. Positive-form Cauchy--Schwarz bounds the
first two by the corresponding square roots in
\eqref{eq:v123-mixed-tail}. Since $0\preceq QHQ\preceq\nu Q$
and $\nu<m$, the last factor is a contraction on the remote
subspace. This proves the mixed estimate without setting either
remote contribution to zero.

For $x=\norm q$, positivity of $H$ and its remote compression give
\[
 \ip{(I-H/m)y}{y}
 \ge v-a/m-2\sqrt{a\nu}\,x/m+(1-\nu/m)x^2.
\]
The minimum of this scalar quadratic over $x\ge0$ is
$v-a/(m-\nu)$. The directional lower bound on $\norm b$
then proves \eqref{eq:v123-majorant-lower}. All pairings are
bounded-operator pairings after preconditioning. No application
of the unbounded $U$ to an arbitrary Hilbert vector is used.
\end{proof}

\begin{proposition}[Failure of the first scalar-head majorant]
\label{prop:v123-first-majorant-fails}
On the same exact even head direction and frozen support-$256$
trial as Proposition~\ref{prop:v122-directional-far}, let
$\mathcal B_1$ denote the complete right side of
\eqref{eq:v122-refined-structured} with
$Q_*=D_g^{-1/2}(I-H/20)D_g^{-1/2}$ and the saved
$\mu=0.9999999999$. Then
\begin{equation}\label{eq:v123-first-majorant-fails}
 \mathcal B_1>6.7609\cdot10^{-20}
 >6.4693\cdot10^{-20}>K_X(v,v).
\end{equation}
Thus this particular complete majorant cannot certify that trial
direction. This is not a negative Weil vector and does not exclude
a better inverse polynomial or a stronger initial head certificate.
\end{proposition}
\begin{proof}[Complete interval certificate]
Retain $J=65536$, $m=20$ and the complete remote bounds of
Proposition~\ref{prop:v122-directional-far}. Choose the exact
$160$-bit dyadic head vector $z$ specified in the reproduction
appendix, obtained by rounding the direction of $b_P$. Its
normalization is enclosed, not assumed exact. The following
outward enclosures suffice in addition to the previous values of
$v,a,E,\nu$:
\begin{center}
\begin{tabular}{@{}ll@{}}
\toprule Quantity & Certified enclosure\\\midrule
$\eta_z$ & $(0.4276838592,0.4276838593)$\\
$\rho_z$ & $<0.0672057785$\\
$|z^*b_P|$ & $(1.7339321076,1.7339321077)\,10^{-10}$\\
$\norm z$ & $(1-10^{-45},1+10^{-45})$\\
$\epsilon_z$ & $<2.961708367\,10^{-11}$\\
complete far energy & $>4.693812399\,10^{-20}$\\
$\mu^{-1}\norm b^2$ & $>2.067157472\,10^{-20}$\\\bottomrule
\end{tabular}
\end{center}
The bound on $\rho_z$ uses the complete order-$64$ moment
majorant on the finite vector $[I;-Z]C^*z$, whose support ends
at $1024$, divided by $g_{J+1}$. It includes every row beyond
$J$. The scalar $\eta_z$ and the retained pairings use the exact
Arb polynomial action from the preceding appendix, with the
unchanged archimedean-series remainder and parity normalization.
Equation~\eqref{eq:v123-majorant-lower} gives
$\mathcal B_1>6.760969872\cdot10^{-20}$, whereas the already
certified $K_X(v,v)<6.469284395\cdot10^{-20}$.
The same frozen direction and solve replay at $768$ and $896$
bits. The outer residual is that of the unshifted Weil form;
the inner inverse factors consistently belong to the certified
shifted tail. Their inverse upper bound remains valid by order.
\end{proof}

\begin{proposition}[Retaining the strengthened head metric]
\label{prop:v123-monotone-head}
In the same setup, let a positive definite finite matrix $M_0$
satisfy
\[
 C K_Z C^*-Y^*Y\succeq M_0\succ0.
\]
The existing certificate permits $M_0=\mu I$. Let $P_j(A)$
be the inverse upper polynomials of
Proposition~\ref{prop:v120-inverse-polynomial}, and set
\begin{align*}
 M_j&=M_0+Y^*(I-P_j(A))Y,\\
 b_j&=d-Y^*P_j(A)y,\\
 \mathcal E_j(r)&=\ip{P_j(A)y}{y}+b_j^*M_j^{-1}b_j.
\end{align*}
Then these are complete simultaneous upper bounds with
\begin{equation}\label{eq:v123-monotone-head}
 \ip{T^{-1}r}{r}\le\mathcal E_{j+1}(r)
          \le\mathcal E_j(r).
\end{equation}
They converge as $j\to\infty$. Their limit need not equal
$\ip{T^{-1}r}{r}$: any loss in replacing
$CK_ZC^*-Y^*Y$ by $M_0$ remains.
For the first polynomial, the head improvement is exactly
$Y^*HY/m$. It must be retained together with the changed
mixed vector if monotonicity of the complete bound is required.
\end{proposition}
\begin{proof}
After the finite triangular substitution and head congruence,
the original form has blocks $CK_ZC^*$, $RC^*$ and $U$,
and the dual vector is $(d,k)$. Put
$Q_j=D_g^{-1/2}P_j(A)D_g^{-1/2}$ and consider the lower form
\[
 \mathcal L_j=
 \begin{pmatrix}
 M_0+Y^*Y & CR^*\\
 RC^* & Q_j^{-1}
 \end{pmatrix}.
\]
The assumed head comparison and $Q_j^{-1}\preceq U$ imply
that the transformed original form dominates $\mathcal L_j$.
Its Schur complement is $M_j\succeq M_0\succ0$.
Because $m^{-1}D_g^{-1}\preceq Q_j\preceq D_g^{-1}$,
all far forms have the same domain $\mathcal D(D_g^{1/2})$.
The finite head and bounded coupling give closed, coercive
lower forms by the same square completion as before.

Their inverse quadratic forms at $(d,k)$ are precisely
$\mathcal E_j(r)$. The polynomials decrease in operator order,
so $Q_j^{-1}$ increases. The head of $\mathcal L_j$ is fixed;
therefore $\mathcal L_{j+1}\succeq\mathcal L_j$, proving the
claimed inverse order. Norm convergence of $P_j(A)$ to $A^{-1}$,
boundedness of $Y$, and the common positive floor of $M_j$
give convergence of the displayed scalar and Gram expressions.
The possible head loss explains the stated limitation of the
limit. A certificate for the shifted tail is again transferred
only by the already stated inverse order.
\end{proof}

\begin{corollary}[The first head update does not rescue this trial]
\label{cor:v123-updated-first-fails}
For the same saved data, choose $M_0=\mu I$ in
Proposition~\ref{prop:v123-monotone-head}. Even its exact first
head update satisfies
\[
 \mathcal E_1(r)>6.6588\cdot10^{-20}>K_X(v,v).
\]
This excludes the first inverse degree with that initial certificate,
including its full positive head update. It does not exclude a
stronger initial $M_0$ or a higher inverse degree.
\end{corollary}
\begin{proof}
Use the same frozen $z$ and complete tail bounds. Positive-form
Cauchy gives
\[
 z^*M_1z\le\mu\norm z^2+
 \frac{(\sqrt{\eta_z}+\sqrt{\nu\rho_z})^2}{20}
       <1.051962238.
\]
Metric Cauchy--Schwarz therefore gives
\[
 b_1^*M_1^{-1}b_1
 \ge\frac{[|z^*b_P|-\epsilon_z]_+^2}{z^*M_1z}
       >1.965049122\cdot10^{-20}.
\]
Adding the complete far lower bound gives
$\mathcal E_1(r)>6.658861521\cdot10^{-20}$.
Both precision replays include this gate with the same exact
dyadic witnesses and all remote correlations.
\end{proof}

\begin{proposition}[Finite probes for the positive head improvement]
\label{prop:v123-head-probes}
Keep $H,Y,P,Q,\nu$ as above. Let $V$ have finitely many
finite-prefix columns, suppose $G=V^*HV\succ0$, and set
$B_P=V^*HPY$. If $Y^*QY\preceq E_Y$, then for every
$0<\theta<1$,
\begin{equation}\label{eq:v123-head-probes}
 Y^*HY\succeq
 (1-\theta)B_P^*G^{-1}B_P
       -(\theta^{-1}-1)\nu E_Y.
\end{equation}
Consequently the right side, divided by $m$ and added to $M_0$,
is a valid finite lower enclosure of the first improved head.
It may be inverted only after its positivity is certified.
\end{proposition}
\begin{proof}
The Gram matrix of $H^{1/2}V$ and $H^{1/2}Y$ gives
$Y^*HY\succeq B^*G^{-1}B$, where $B=V^*HY$.
Write $B=B_P+E_B$ with $E_B=V^*HQY$.
The projection onto the range of $H^{1/2}V$ shows
\[
 E_B^*G^{-1}E_B\preceq Y^*QHQY\preceq\nu E_Y.
\]
Young's inequality in the $G^{-1}$ metric proves
\eqref{eq:v123-head-probes}. This keeps the omitted probe
correlation explicitly. No unjustified matrix positive-part
operation is applied.
\end{proof}

These first-degree bounds do not settle the complete low-mode Schur
sign. The later simultaneous and directional-complement certificates
prove both parity signs at $\lambda=4$.
The first-degree obstruction now has a proof even with the
induced head update and the saved initial certificate. Increasing
the inverse degree, or strengthening the initial certificate,
requires a different complete enclosure.
Neither finite-prefix improvements nor convergence at this fixed
window establishes uniformity as $\lambda\to\infty$.

\subsection*{A complete Schur-direction certificate from a finite trial}

The first-degree obstruction above concerns a particular inverse
majorant and saved trial. A finite iterative calculation can instead
propose a new trial, after which the complete residual is verified
afresh. This uses the same verified-solve identity; it does not
identify a power of a finite compression with the full power.

\begin{proposition}[A posteriori verification without a full-power claim]
\label{prop:v124-posterior-identity}
Let the complete outer operator have finite head and coercive
tail $T$, with Schur form $S_\infty=F-B^*T^{-1}B$.
For a finite trial $g=(v,w)$ in its operator domain, set
$e=Bv+Tw$. Then
\begin{equation}\label{eq:v124-posterior-identity}
 S_\infty(v,v)=QW(g,g)-\ip{T^{-1}e}{e}.
\end{equation}
Thus any certified complete upper bound $\mathcal U(e)$ for the
last term gives
\[
 QW(g,g)-\mathcal U(e)\le S_\infty(v,v)\le QW(g,g).
\]
The algorithm that proposed the frozen finite coefficients of $w$
need not be an approximation to any particular infinite polynomial.
\end{proposition}
\begin{proof}
Expand $\ip{T^{-1}(Bv+Tw)}{Bv+Tw}$ and cancel the tail and
cross terms in $QW(g,g)$. Coercivity justifies the inverse.
Equivalently, for any residual $r$ and finite correction $z$,
\[
 \ip{T^{-1}r}{r}
 =2\Re\ip rz-\ip{Tz}{z}
       +\ip{T^{-1}(r-Tz)}{r-Tz}.
\]
If the inner certificate belongs to $T-cD$, use
$T^{-1}\preceq(T-cD)^{-1}$ only in bounding the new residual
term. The energy and $e=P_{\rm tail}Wg$ remain unshifted.
This avoids an erroneous omission of a $cDz$ residual or a shift
energy when the two operators are mixed.
\end{proof}

For completeness, let the new residual split as $e=(h,k)$ in
the already certified inner decomposition. Past the support of
the saved $Z$, retain rows through $J$ and put
\[
 V_J=\sum_{n\le J}|k_n|^2/g_n,\qquad
 b_J=C\left(h-Z^*k-\sum_{n\le J}R_n^*k_n/g_n\right).
\]
The sums run over the inner far rows. If complete remote bounds
give $\sum_{n>J}|k_n|^2/g_n\le E$ and
$\|D_g^{-1/2}Q_JRC^*\|^2\le\rho$, then the old structured
inverse certificate supplies
\begin{equation}\label{eq:v124-posterior-upper}
 \mathcal U(e)=V_J+E+
       \mu^{-1}\bigl(\norm{b_J}+\sqrt{\rho E}\bigr)^2.
\end{equation}
Every row beyond $J$ is included. For a matrix of trial columns,
the simultaneous version retains their Gram matrices and uses
\eqref{eq:v120-structured-finite-rows}; separate positive scalar
tests are insufficient for positivity on their span.

\begin{proposition}[One complete positive Schur direction at $\lambda=4$]
\label{prop:v124-positive-direction}
Let $v$ be the exact archived even head vector used in
Propositions~\ref{prop:v122-directional-far} and
\ref{prop:v123-first-majorant-fails}. For the \emph{complete}
even Schur form with head $0\le n\le16$,
\begin{equation}\label{eq:v124-positive-direction}
 3.0930\cdot10^{-20}<S_\infty(v,v)<3.2356\cdot10^{-20}.
\end{equation}
This certifies one fixed direction. It does not assert that the
whole even head, the odd head, or the full $\lambda=4$ form is
positive, and it provides no growing-window uniformity.
\end{proposition}
\begin{proof}[Frozen finite trial and complete residual certificate]
Start with the previous support-$256$ vector $f=[v;-Xv]$.
Subtract the exact dyadic tail correction $z$, supported on
$17\le n\le4096$, specified in the reproduction appendix.
The new vector $g=f-(0,z)$ has exactly the same head $v$.
A finite preconditioned iteration proposed $z$; neither its
convergence nor its recursive residual is used in this proof.

Evaluate $QW_4(g,g)$ and the unshifted residual
$e=Q_{16}W_4g$ through $J=65536$. Retain the original shifted
inner factors, $\mu=0.9999999999$, and the previously certified
$\rho<0.165102333$. Direct interval evaluation gives
\begin{center}
\begin{tabular}{@{}ll@{}}
\toprule Quantity & Certified enclosure\\\midrule
$QW_4(g,g)$ & $(3.2355843481,3.2355843482)\,10^{-20}$\\
$V_J$ & $<3.742924324\,10^{-22}$\\
$\norm{b_J}^2$ & $<1.617782100\,10^{-23}$\\
$E$ & $<8.080848756\,10^{-22}$\\
$\mathcal U(e)$ & $<1.424889\,10^{-21}$\\\bottomrule
\end{tabular}
\end{center}
The bound $E$ is the complete order-$64$ moment majorant applied
to the finite vector $g$, with its actual support through $4096$,
divided by $g_{J+1}$. The saved remote $\rho$ belongs to the same
inner witness and cutoff. Substitution in
\eqref{eq:v124-posterior-upper} and
\eqref{eq:v124-posterior-identity} proves the lower bound.
The upper bound follows from $T^{-1}\succeq0$.

The frozen correction uses exact dyadic coefficients produced at
$256$ bits; the proof evaluates them anew at $768$ and $896$
bits. Both checks retain the zero-mode $\sqrt2$ factor, normalized
parity, logarithmic diagonal and archimedean-series remainder.
Selected rows, including zero and the support endpoint, are also
checked by dense coefficient assembly. The independently expanded
energy identity agrees with the direct energy. No finite power is
used as a substitute for an infinite operator power.
\end{proof}

For the \emph{old} trial in this same direction the complete inverse
correction is consequently bounded by
\[
 3.2337\cdot10^{-20}<\ip{T^{-1}r}{r}
       <3.3762\cdot10^{-20}.
\]
In particular $S_\infty(v,v)>0.4781K_X(v,v)$ on this one
direction. This explains why the earlier majorant could fail
without identifying an intrinsically difficult Schur direction.
The next target is the simultaneous matrix comparison, including
the remaining directions of small energy. This directional
certificate cannot be extrapolated to that matrix or to RH.

\subsection*{A simultaneous complete Schur certificate}

A finite trial may be chosen in coordinates adapted to its head energy;
the final certificate must still use its actual head and every residual
cross term.  This distinction is essential when some ordinary head
energies are extremely small.

\begin{proposition}[Simultaneous trial and ordinary-error certificate]
\label{prop:v125-simultaneous}
Let the complete closed Weil form have finite head and coercive tail
$T$, and let $G$ be a matrix of finite Fourier trial columns.  Write
$V=P_{16}G$, $E=Q_{16}\mathsf W_4G$ and $K=G^*\mathsf W_4G$.
Then
\begin{equation}\label{eq:v125-simultaneous}
 V^*S_\infty V=K-E^*T^{-1}E.
\end{equation}
Suppose $V$ is square and invertible, and a complete matrix upper bound
$\mathcal U\succeq E^*T^{-1}E$ is certified.  If
$K-\mathcal U\succ0$, the entire parity form is strictly positive.
More generally, if $\varepsilon\ge0$ and
\[
 K-\mathcal U+\varepsilon V^*V\succeq0,
\]
then $S_\infty\succeq-\varepsilon I$ and the entire parity form is
bounded below by $-\varepsilon I$ in ordinary norm.
\end{proposition}
\begin{proof}
Apply Proposition~\ref{prop:v124-posterior-identity} to every linear
combination of the columns.  Invertibility of the actual head $V$
transfers the resulting congruent inequality to $S_\infty$.
For a head vector $x$ and tail vector $y$, the closed form identity is
\[
 q_{\mathsf W_4}[x,y]
 =\ip{S_\infty x}{x}
   +\norm{T^{1/2}(y+T^{-1}Bx)}^2.
\]
The head-to-tail map $B$ is bounded because the head is finite
and its Fourier vectors belong to the operator domain.  Thus the
completion is valid on the form domain.  The tail square is
nonnegative, and $\norm{x}\le\norm{(x,y)}$.  Strict head and tail
positivity also imply fixed-window coercivity by the bounded
invertible triangular change in this identity.  No quantitative
ordinary spectral gap is identified with an $LDL^*$ pivot.
\end{proof}

\begin{proposition}[Retaining a joint remote Gram]
\label{prop:v125-joint-remote}
In the finite-row structured inverse bound, put
$Y=D_g^{-1/2}Q_JRC^*$ and $y=D_g^{-1/2}Q_Jk$, and let
$H=C(h-Z^*k-R_J^*D_g^{-1}k_J)$ be the retained mixed matrix.
Suppose one joint enclosure and a lower bound satisfy
\[
 \Gamma=\begin{pmatrix}\Gamma_{YY}&\Gamma_{Yk}\\
 \Gamma_{Yk}^*&\Gamma_{kk}\end{pmatrix}
 \succeq [Y,y]^*[Y,y],\qquad 0\preceq L_Y\preceq Y^*Y,
 \qquad A_0=\mu I+L_Y-\Gamma_{YY}\succ0.
\]
Then the complete inverse correction is bounded by
\begin{equation}\label{eq:v125-joint-remote}
 E^*T^{-1}E\preceq V_J+\Gamma_{kk}
 +(H-\Gamma_{Yk})^*A_0^{-1}(H-\Gamma_{Yk}).
\end{equation}
Here $V_J=k_J^*D_g^{-1}k_J$.  One may take $L_Y=0$.
If $\Gamma_{YY}\preceq\rho I$ with $\rho<\mu$, replacing
$A_0^{-1}$ by $(\mu-\rho)^{-1}I$ gives a simpler valid bound.
\end{proposition}
\begin{proof}
For an outer coefficient vector $x$, the remote portion of the old
structured majorant is
\[
 \norm{yx}^2+\mu^{-1}\norm{Hx-Y^*yx}^2
 =\sup_z\{\norm{yx-Yz}^2-\norm{Yz}^2
       +2\Re\ip{Hx}{z}-\mu\norm z^2\}.
\]
Apply the joint upper Gram to $(-z,x)$ and the lower bound $L_Y$
to $\norm{Yz}^2$, then complete the finite $A_0$ square.
This proves the claim for every $x$.  The bounded finite-column maps
justify all remote pairings.  The cross block must come from the
\emph{same joint} Gram; unrelated diagonal upper bounds do not
justify an invented mixed term.
\end{proof}

\begin{proposition}[The complete even form at $\lambda=4$]
\label{prop:v125-even-complete}
The canonical complete Weil form $\mathsf W_4$ restricted to its
even parity space is strictly positive on every nonzero vector
in its form domain, and is bounded below by a positive constant
at this fixed window.
\end{proposition}
\begin{proof}[Frozen matrix certificate]
Keep the outer head $0\le n\le16$ and the complete tail certified
in Proposition~\ref{prop:v120-lambda4-tail}.  The archived matrix
trial has $17$ columns and support through $M=4096$.  Its exact
dyadic base and correction coefficients define $G$; the actual
head $V=P_{16}G$ is certified invertible.  The proposal first scales
the old finite trial by its head-energy $LDL^*$ factors, rounds
that base to exact dyadics, and uses $16$ preconditioned finite CG
steps at $512$ bits for each column.  Neither CG convergence nor
its recursive residual enters the proof.

Recompute the unshifted $K=G^*\mathsf W_4G$ and
$E=Q_{16}\mathsf W_4G$ at $768$ and $896$ bits from the same
frozen inputs.  Retain all cross-Grams through $J=65536$ and apply
the order-$64$ moment bound to the entire trial for every omitted
row.  The saved shifted inner factors have $N=512$, support $1024$,
$\mu=0.9999999999$ and complete remote bound
$\rho<0.165102333$.  Inverse order is used only for the residual
inverse, so the trial energy and residual stay unshifted.
For $\tau=1/10$, equation~\eqref{eq:v120-structured-finite-rows}
gives the complete matrix upper bound $\mathcal U_\tau$.
All $17$ interval $LDL^*$ pivots of $K-\mathcal U_\tau$ are positive;
the smallest exceeds $0.7645944101$ in these saved head coordinates.
Both precision replays give this strict gate.  This number certifies
sign and is not an ordinary-norm eigenvalue bound.

The verifier checks the actual head determinant, exact dyadic
reconstruction, inner-certificate provenance, positivity of the
far diagonal, and nine independently assembled physical rows for
every column.  The complete remote moment includes the zero-mode
normalization and both Fourier signs through the parity basis.
Proposition~\ref{prop:v125-simultaneous} therefore proves positivity
of the entire even form, not only a Fourier compression.
\end{proof}

A common remote Gram for the padded inner source $[I;-Z]C^*$ and
this trial also supplies the cross centre in
Proposition~\ref{prop:v125-joint-remote}.  Its scalar-denominator
certificate has all $17$ pivots positive, the smallest greater than
$0.7892519172$; using the full finite denominator improves that
smallest pivot to more than $0.7892519657$.  This is an additional
fixed-witness check, not needed for the preceding even sign proof.
The whole-operator conclusion comes from a simultaneous matrix
inequality and complete tail control; separately positive directions
would not suffice.

The comparison also has a certified dimensionless margin:
\begin{equation}\label{eq:v125-relative-margin}
 \mathcal U_{1/10}\preceq
 \left(1-\frac{62629}{100000}\right)K.
\end{equation}
Interval $LDL^*$ verifies the equivalent positive matrix at both
precisions.  Thus the generalized relative margin exceeds $0.62629$;
unlike a smallest ordinary eigenvalue or a coordinate pivot, this
ratio is invariant under an invertible head congruence.  On the
original physical head direction from
Proposition~\ref{prop:v124-positive-direction}, this new matrix trial
has certified headroom greater than $12.1284$.  Its trial is different
from the earlier single-direction trial, whose headroom was about
$22.71$; changing the trial changes this diagnostic while leaving
$S_\infty(v,v)$ unchanged.

The analogous initial odd-head bounds are inconclusive. The next
subsection completes the odd sector by combining a complete
complement certificate with a larger-support directional trial.
No finite list of such fixed-window statements implies the cofinal
estimate needed for RH.

\subsection*{A complete complement bound and the full sign at $\lambda=4$}

The $\lambda=4$ odd computation can be reduced to one direction
without inferring a matrix sign from separately positive vectors.
The reduction uses positivity of the \emph{exact inverse correction},
which is already known from the coercive complete tail.

% Proposed proof insertion only; no numerical gate is asserted here.
\begin{proposition}[One repaired subspace and its complete complement]
\label{prop:v126-direction-complement}
Let $K$ be a positive definite Hermitian form on a finite-dimensional
complex space, let $\mathcal C\succeq0$, and put $S=K-\mathcal C$.
Let $E$ be a subspace and write $F=E^{\perp_K}$.
Suppose that $s\leq1$, $\sigma\geq0$, and, simultaneously on the
indicated subspaces,
\[
 S|_E\succeq sK|_E,
 \qquad \mathcal C|_F\preceq\sigma K|_F.
\]
Then
\[
 S\succeq(s-\sigma)K.
\]
In particular, when $E=\operatorname{span}\{v\}$ and $K(v,v)=1$,
the first premise is the single inequality $S(v,v)\geq s$.
\end{proposition}
\begin{proof}
For $x=e+f$ with $e\in E$ and $f\in F$, Cauchy--Schwarz for the
positive semidefinite form $\mathcal C$ gives
\begin{align*}
 \mathcal C(x,x)
 &\leq\bigl(\sqrt{\mathcal C(e,e)}+
                 \sqrt{\mathcal C(f,f)}\bigr)^2\\
 &\leq\bigl(\sqrt{1-s}\,\|e\|_K+
                 \sqrt{\sigma}\,\|f\|_K\bigr)^2\\
 &\leq(1-s+\sigma)(\|e\|_K^2+\|f\|_K^2)
  =(1-s+\sigma)K(x,x).
\end{align*}
Subtracting from $K(x,x)$ proves the claim.  Degeneracy of
$\mathcal C$ causes no problem, since Cauchy--Schwarz holds for
its induced seminorm.  The constant cannot be improved under
these premises alone: take $K=I$ on $\mathbb C^2$,
$E=\operatorname{span}\{(1,0)\}$, and
$\mathcal C=uu^*$ with $u=(\sqrt{1-s},\sqrt{\sigma})^T$.
\end{proof}

\begin{corollary}[Ordinary error localized to the unresolved subspace]
\label{cor:v126-local-error}
Under the hypotheses of Proposition~\ref{prop:v126-direction-complement},
assume $0\le\sigma<1$, let $P_E$ be the $K$-orthogonal projection
onto $E$, and set
\[
 \eta=\frac{(\sigma-s)_+}{1-\sigma}.
\]
Then
\begin{equation}\label{eq:v126-local-error}
 S\succeq-\eta P_E^*KP_E.
\end{equation}
If $S$ is the complete Schur form of a trial with invertible
physical head $V$ and coercive complete tail, the full form has
ordinary lower error at most
\begin{equation}\label{eq:v126-local-physical-error}
 \varepsilon^{\rm loc}
 =\eta\,\norm{K^{1/2}P_EV^{-1}}^2.
\end{equation}
The congruence cost is thus restricted to the unresolved subspace;
it is not discarded.
\end{corollary}
\begin{proof}
Write $x=e+f$, $e\in E$, $f\in E^{\perp_K}$, and
$a=\norm e_K$, $b=\norm f_K$. Positivity of $\mathcal C$ gives
\[
 S(x,x)\ge sa^2+(1-\sigma)b^2
                 -2\sqrt{(1-s)\sigma}\,ab.
\]
Completing the square in $b$ bounds this below by
$(s-\sigma)a^2/(1-\sigma)$, which proves
\eqref{eq:v126-local-error}. For physical head $y=Vx$, use
$a^2=\norm{K^{1/2}P_EV^{-1}y}^2$. The positive complete-tail
square and $\norm y\le\norm{\text{full vector}}$ give
\eqref{eq:v126-local-physical-error} on the closed form domain.
\end{proof}

Along a cofinal window family, establishing all these hypotheses
and $\varepsilon^{\rm loc}_\lambda\to0$ would satisfy
Proposition~\ref{prop:v121-cofinal-rh}, using its fixed-support
identification with the same Weil form. This is a sufficient target
which can tolerate a negative restricted estimate; no such
convergence is proved here. Its factors must be controlled for
the complete operator at every window of that family.

\paragraph{Use with two finite trials.}
For a coercive complete tail $T$ and any finite trial $G$ with
invertible physical head $V$, let
$K=G^*WG$, $r=QWG$, and $\mathcal C=r^*T^{-1}r$.
The exact Schur matrix in these head coordinates is $S=K-\mathcal C$.
A second finite trial with the same physical head gives the same $S$.
Consequently a directional estimate from the second trial may be combined
with a complete complement estimate from the first, provided that both
are normalized by the \emph{same first-trial metric} $K$.
An available complete upper bound $\mathcal C\preceq U$ permits the
complement premise to be verified as
$Z^*(\sigma K-U)Z\succeq0$, where $Z$ spans the exact
$K$-orthogonal complement.  No numerical eigenvector assumption is needed.

\begin{proposition}[Complete odd complement]
\label{prop:v126-odd-complement}
At $\lambda=4$, let $G_0$ be the archived odd $16$-column trial
with support through $4096$, physical head $V=P_{16}G_0$,
$K=G_0^*\mathsf W_4G_0$, and complete correction
$\mathcal C=(Q_{16}\mathsf W_4G_0)^*T^{-1}
                (Q_{16}\mathsf W_4G_0)$.
There is a specified nonzero exact dyadic vector $v\in\mathbb R^{16}$
for which
\begin{equation}\label{eq:v126-complement}
 K\succ0,\qquad
 \mathcal C|_{v^{\perp_K}}\preceq10^{-3}K|_{v^{\perp_K}}.
\end{equation}
This is a complete inverse-correction estimate on the indicated
subspace, with every omitted Fourier row retained.
\end{proposition}
\begin{proof}[Outward certificates and the complete-form passage]
Use the same odd inner factors as
Proposition~\ref{prop:v120-lambda4-tail}, with $N=1536$,
inner solve support $2048$, verification cutoff $J=65536$ and
remote moment order $64$.  The frozen trial has $16$ proposal
iterations per column; its construction is not a proof input.
Reconstructing it at $768$ bits gives the complete matrix upper
$\mathcal U_\tau$ of Proposition~\ref{prop:v125-simultaneous}.
The exact frozen dyadics defining $v$ are included in the bundle.

Put $a=v^*Kv>0$ and $P=I-v(v^*K)/a$.  Delete column $12$ of
$P$ to obtain $Z$.  The corresponding coordinate of $v$ is nonzero,
so $Z$ has rank $15$ and range $v^{\perp_K}$: the kernel of $P$
is $\operatorname{span}\{v\}$, which has no nonzero vector
with coordinate $12$ zero.  The projection uses the exact $K$;
its interval enclosure is not frozen at a midpoint.  For
$\tau=1/20$, interval $LDL^*$ proves
\[
 Z^*(10^{-3}K-\mathcal U_{1/20})Z\succ0,
\]
with all $15$ pivots positive and the smallest exceeding
$0.0004086079$.  Since $\mathcal C\preceq\mathcal U_{1/20}$,
this proves \eqref{eq:v126-complement}.  The scalar direction
was suggested numerically, but its proposed eigenvector property
is neither assumed nor needed.

\end{proof}

\begin{proposition}[Complete positivity at $\lambda=4$]
\label{prop:v126-full-window}
The canonical complete Weil form $\mathsf W_4$ is strictly positive
on every nonzero vector in its form domain and is bounded below
by a positive constant at this fixed window. In particular its
nonnegative-error formulation has $\varepsilon_4=0$ exactly.
\end{proposition}
\begin{proof}[A larger-support directional certificate]
Keep $G_0,K,V,v$ and the complete odd tail $T$ from
Proposition~\ref{prop:v126-odd-complement}; put $a=v^*Kv$.
The archived exact dyadic odd trial $f$ has support through $8192$
and physical head exactly $Vv$. Thus, with $r_f=Q_{16}\mathsf W_4f$,
\[
 S(v,v)=QW_4(f,f)-\ip{T^{-1}r_f}{r_f},\qquad S=K-\mathcal C.
\]
This identity uses the same complete Schur form and the old
normalization $a$, although the finite trial has changed.

The same physical diagonal, inner certificate, $J=65536$ and
order-$64$ remote moment give the following conservative outward
intervals, reproduced at $1024$ and $1280$ bits on the same frozen
witness:
\begin{center}
\begin{tabular}{ll}
\toprule
quantity & enclosing interval\\
\midrule
$a$ & $(1,\ 1.000000000000002)$\\
$QW_4(f,f)$ & $(0.9016760830,\ 0.9016760831)$\\
finite far energy upper $V_f$ & $(0.3125027254,\ 0.3125027255)$\\
remote source majorant $E_f$ & $(0.0608540372,\ 0.0608540373)$\\
finite mixed squared norm $H_f$ & $(0.0701478147,\ 0.0701478148)$\\
complete correction upper $U_f$ & $(0.4725352935,\ 0.4725352936)$\\
$(QW_4(f,f)-U_f)/a$ & $(0.4291407894,\ 0.4291407895)$\\
\bottomrule
\end{tabular}
\end{center}
Here $\mu=0.9999999999$ and
$0.0412001510<\rho<0.0412001511$ are the certified odd inner
constants. Proposition~\ref{prop:v125-simultaneous}, or its scalar
Cauchy--Schwarz form, gives
\[
 \ip{T^{-1}r_f}{r_f}\le U_f
 :=V_f+E_f+\mu^{-1}
          \bigl(\sqrt{H_f}+\sqrt{\rho E_f}\bigr)^2.
\]
The bound includes every row past $J$, not only the finite residual.
The verifier checks exact equality of all physical head coefficients,
provenance of the frozen inner factors, and eight independently
assembled rows of the unshifted physical operator. The physical
head equality is also checked with independent exact rational
arithmetic; its largest dyadic mantissa has $912$ bits. The finite
proposal method supplies no hypothesis about its convergence.

The outward directional inequality proves
$S(v,v)\ge(429/1000)a$. Combining it with
\eqref{eq:v126-complement} in
Proposition~\ref{prop:v126-direction-complement} yields the
simultaneous odd-head estimate
\begin{equation}\label{eq:v126-odd-relative-margin}
 S\succeq\frac{428}{1000}K\succ0.
\end{equation}
Since $V$ is invertible and the complete tail is coercive,
square completion proves positivity of the entire odd form.
The triangular change is bounded with bounded inverse, so the
finite positive head and coercive tail also give a positive
ordinary lower constant at this fixed window. The orthogonal
parity decomposition and Proposition~\ref{prop:v125-even-complete}
prove the full claim. The relative constant $428/1000$ is not
an ordinary spectral-gap estimate, and no congruence factor is
needed to pass the exact sign $\varepsilon_4=0$.
\end{proof}

The reduced directional argument closes this fixed-window obligation.
It gives no bound for the heads, inverse factors or residuals at
unbounded $\lambda$. In particular, the constants above are tied
to $L=2\log4$ and are not transferred to other windows.

% NS-17: window_comparison_insert.tex
\subsection*{Position of the certified window}
\label{sec:v141-window-position}
The window parameter $\lambda$ is not the parameter used in the
published compact-window literature, which is indexed by the
logarithmic support half-width. A test supported multiplicatively in
$[\lambda^{-1},\lambda]$ is supported in $[-a,a]$ in the logarithmic
coordinate, with
\[
 a=\log\lambda=L/2 .
\]
Zhu~\cite{Zhu2026} writes this half-width $L$; it is one half of the
logarithmic length $L=2\log\lambda$ of the reading conventions. The
classical positivity range of Yoshida and Connes--Consani is stated
as $\operatorname{supp}f\subset[-(\log2)/2,(\log2)/2]$, as recalled in
\cite{Zhu2026}. On one axis:
\begin{center}
\begin{tabular}{lrr}
\toprule
result & half-width $a$ & $\lambda=e^{a}$\\
\midrule
classical Yoshida / Connes--Consani & $0.3466$ & $1.414$\\
Zhu~\cite{Zhu2026}, certified $L=0.8$ & $0.8000$ & $2.226$\\
Proposition~\ref{prop:v116-window-positive}, $\lambda=3$ & $1.0986$ & $3$\\
Proposition~\ref{prop:v126-full-window}, $\lambda=4$ & $1.3863$ & $4$\\
\bottomrule
\end{tabular}
\end{center}
The $\lambda=4$ certificate is therefore at $1.73$ times the
half-width of the published Zhu window and $4$ times the classical
range. The comparison is legitimate because the semilocal form and the
global Weil form agree on every test supported inside the window;
Proposition~\ref{prop:v121-cofinal-rh} relies on exactly this
extension-by-zero identity, so the two sides certify the same object
on the same functions.

The two results are not of the same kind. Zhu proves a certified
two-sided enclosure
$8.9\cdot10^{-18}\le\lambda_{\min}(0.8)\le2.27\cdot10^{-17}$ of the
normalized window infimum at $a=0.8$, by a one-stroke reduction to a
single finite positive semidefinite matrix.
Proposition~\ref{prop:v126-full-window} proves nonnegativity
$\mathsf W_4\succeq0$ at $a=\log4$, and
Proposition~\ref{prop:v140-ground4} encloses its complete ground,
$0<\mu_0<2.454\cdot10^{-75}$ and $\mu_1>10^{-73}$, by a structured
inverse, a simultaneous seventeen-column certificate and
direction--complement gluing. Zhu's result is stronger in kind, a
quantitative two-sided gap on a smaller window; the present one is
stronger in reach. Neither subsumes the other. The $LDL^*$ pivots and
the margin $1-\lambda_{\max}(\mathcal U,K)$ recorded in the
$\lambda=4$ certificates are not ordinary-norm eigenvalue bounds and are not comparable with
Zhu's enclosure: the pivots depend on coordinates, whereas the relative
margin compares two forms.

In the notation of Connes--Consani--Moscovici
\cite[Corollaries~3.7--3.8]{ConnesConsaniMoscovici2025}, write
$\mu^{\mathrm{CCM}}_\lambda$ for the infimum of the spectrum of the
self-adjoint operator of the form $QW_\lambda$ certified here. They
prove that $\lambda\mapsto\mu^{\mathrm{CCM}}_\lambda$ is nonincreasing,
write that they ``cannot assert that'' it is nonnegative, and prove
that $\lim_{\lambda\to\infty}\mu^{\mathrm{CCM}}_\lambda=0$ implies RH.
By that monotonicity,
\[
 \mathsf W_4\succeq0
 \quad\Longleftrightarrow\quad
 \mu^{\mathrm{CCM}}_\lambda\ge0\ \text{ for every }1<\lambda\le4,
\]
so Proposition~\ref{prop:v126-full-window} settles the sign of
$\mu^{\mathrm{CCM}}_\lambda$ on $1<\lambda\le4$, a range $1.73$ times
Zhu's certified window and beyond the windows of their own numerical
section. The floors of Proposition~\ref{prop:v138-three-floors} are the
lower bounds $\mu^{\mathrm{CCM}}_\lambda\ge-8$ for $1<\lambda\le8$.
This is a target statement, not a shortcut: cofinal nonnegativity of
$\mu^{\mathrm{CCM}}_\lambda$ is RH by
Proposition~\ref{prop:v121-cofinal-rh}, and no cofinal statement is
made here. No G2 gap or RH claim follows.



\subsection*{Reported failure of the unchanged construction at $\lambda=5$}
\label{sec:v127-transfer}
\paragraph{Availability (NS-38).} The original frozen protocols, dyadic
witnesses, coefficient enclosures and $320/448$-bit verification outputs
for this subsection are \emph{not archived}. The following numerical
assertions and verification account are historical reports, not currently
replayable certified computations. See \path{evidence/MISSING.md}.
The transfer protocol was frozen before testing, including the exact
energy-scaling rule, CG depths, preconditioner, support choices, Young
parameters and moment orders. The successful $\lambda=4$ code specifies
literal sizes; it supplies no general window-dependent support rule.
Importing the separate finite-resolution probe's $\lambda^2$ head rule
would therefore be a modification. The surviving protocol description is recorded in
Appendix~\ref{app:v127-transfer}; the original protocol file is not archived.

\begin{proposition}[Reported counterexample; original evidence not archived]
\label{prop:v127-metric-fails}
The historical computation reports the following assertion. Let $D=\operatorname{diag}(a_n)$ be the actual archimedean diagonal
at $\lambda=5$. In each parity the complete inequality
\[
 Q_{16}\mathsf W_5Q_{16}\succeq10^{-8}Q_{16}DQ_{16}
\]
is false. Exact frozen dyadic vectors $v_+$ and $v_-$, supported on
parity indices $17,\ldots,128$, satisfy respectively
\begin{align*}
 2.0806767455\cdot10^{-18}
 &<\frac{QW_5(v_+,v_+)}{\ip{Dv_+}{v_+}}
 <2.0806767456\cdot10^{-18},\\
 6.7808574118\cdot10^{-17}
 &<\frac{QW_5(v_-,v_-)}{\ip{Dv_-}{v_-}}
 <6.7808574120\cdot10^{-17}.
\end{align*}
Both denominators and both Weil numerators are strictly positive.
\end{proposition}
\begin{proof}[Historical verification account; original evidence not archived]
The historical account uses the exact coefficient formula and analytic
coefficient enclosures of Proposition~\ref{prop:v114-weil-core}.
It reports that the outward verifier forms
\[
 W_{nm}^{\pm}=W_{nm}\pm W_{n,-m},\qquad
 W_{nn}=d_n,\quad
 W_{nm}=\frac{b_n-b_m}{n-m}\ (n\ne m),\quad b_{-n}=-b_n.
\]
The metric diagonal is $a_n$, not the full arithmetic diagonal $d_n$.
At $320$ and $448$ bits the frozen vectors reportedly give the stated strict
intervals and negative values of $QW_5(v,v)-10^{-8}\ip{Dv}{v}$.
A separate signed-index quadratic assembly expands each parity vector
to coefficients $v_n/\sqrt2$ at $+n$ and $\pm v_n/\sqrt2$ at $-n$,
and reportedly reproduces the intervals at $448$ bits. The odd expansion differs
from real-sine coordinates only by a common unit phase. Witness hashes,
coefficient enclosures and the independent verifier are described in
Appendix~\ref{app:v127-transfer}; only the descriptions and recorded
hashes survive here. The underlying files are not archived. Each vector has literal finite support,
so its complete quadratic form is exactly this finite sum: no remote
row contributes to the pairing. If the reported ratio enclosures hold,
they are below $10^{-8}$ and imply the asserted complete-tail failure.
\end{proof}

The first remote-weight gate already fails at the frozen even
$N=512$ and odd $N=1536$: its historically reported certified lower weights are approximately
$-1.0764504091$ and $-1.5489108675$. The same historical evaluation
of the weight formula reports first positive cutoffs $1502$ and $7224$.
These changes alone cannot repair Proposition~\ref{prop:v127-metric-fails}:
the reported witnesses, if recovered and verified, disprove the required
metric itself, independently of the verification cutoff. More CG steps, greater moment order, or
larger remote support cannot make that fixed inequality true.

The necessary change is to the outer head or the tail metric. Conditional on the reported witness bounds, keeping
$Q_{16}$ and a common scalar $cD$ forces
$c<2.0806767456\cdot10^{-18}$, a necessary ceiling only, not a certified
admissible coefficient. No modified construction was run after this
obstruction, and no successful tuning rule has been established.
In particular, no positive $\lambda=5$ tail metric or structured
inverse factors are asserted. The complete certificate margin
$1-\lambda_{\max}(\mathcal U,K)$ is unavailable; a finite-prefix value
or a resolution margin cannot replace it. These reported witnesses
would be counterexamples to the frozen metric, not negative Weil vectors.

\subsection*{What a non-scalar metric must change}
A change of coordinates alone cannot turn a negative lower comparison
into a positive one. The next obstruction also rules out repairing
that particular comparison in only polynomially many directions.

\begin{proposition}[Rank cost of repairing the unsigned comparison]
\label{prop:v127-rank-barrier}
In one parity let $H_N=\operatorname{diag}(g_n)_{n>N}$ satisfy
\[
 g_n\le(1-c)\log(n/L)-\kappa_\lambda,\qquad 0\le c<1,
 \quad \kappa_\lambda\ge\|\mathsf T_{{\rm pr},\lambda}\|.
\]
Set $\Theta=L\exp(\kappa_\lambda/(1-c))$ and
$m=\max\{0,\lceil\Theta\rceil-N-1\}$. For every bounded self-adjoint
correction $R$ with rank $r<m$, $H_N+R$ has a strictly negative
form direction. An invertible congruence, on the pulled-back form
domain, cannot remove this negative direction. Consequently, at a
polynomial cutoff $N(\lambda)$, a polynomial-rank correction cannot
make this unsigned comparison nonnegative along unbounded windows.
\end{proposition}
\begin{proof}
The $m$ coordinate vectors with $N<n<\Theta$ span a strictly negative
subspace for $H_N$. The kernel of $R$ has codimension at most $r$.
If $r<m$, their intersection contains a nonzero vector $v$, and
$\ip{(H_N+R)v}{v}=\ip{H_Nv}{v}<0$. Pulling $v$ back through an
invertible congruence preserves its negative form value.
Proposition~\ref{prop:v119-prime-scale} gives
$\kappa_\lambda\ge(1-o(1))\lambda$, so $m$ is exponential when
$N$ is polynomial. This proves the stated rank limitation.
\end{proof}

This result concerns the already weakened, unsigned diagonal
comparison. It imposes no cutoff lower bound on the actual Weil form
and does not exclude an infinite-rank signed metric. In particular,
using the exact prime and archimedean correlations before taking a
norm changes the comparison and is not covered by this obstruction.
Choosing a metric equal to the unknown positive operator would be
circular; a useful metric needs its own independent certificate.

\begin{proposition}[An infinite signed block-metric criterion]
\label{prop:v127-block-metric}
Let a symmetric form $q_T$ be given on finite-block sums,
with Hermitian block entries $T_{jk}$. Suppose every finite diagonal
block satisfies $T_{jj}\succeq M_j\succ0$. Suppose further that
\[
 \beta_{jk}=\|M_j^{-1/2}T_{jk}M_k^{-1/2}\|\quad(j\ne k),\qquad
 \sup_j\sum_{k\ne j}\beta_{jk}\le\rho<1.
\]
Then the form on finite-block sums is closable, its closure has the
diagonal-block form domain, and
\[
 T\succeq(1-\rho)\bigoplus_jT_{jj}
   \succeq(1-\rho)\bigoplus_jM_j.
\]
If the actual closed operator agrees on this form core, this is a
bound for that operator. A positive ordinary lower constant additionally
follows if $\inf_j\lambda_{\min}(M_j)>0$ at the fixed window.
\end{proposition}
\begin{proof}
For a finite-block vector put $a_j=\|M_j^{1/2}f_j\|$.
Self-adjointness gives $\beta_{jk}=\beta_{kj}$. Hence
\[
 |q_{\rm off}[f]|\le2\sum_{j<k}\beta_{jk}a_ja_k
 \le\sum_j a_j^2\sum_{k\ne j}\beta_{jk}
 \le\rho\sum_j a_j^2.
\]
The latter sum is at most $q_{\rm diag}[f]$ with diagonal blocks
$T_{jj}$. Thus $q_T\ge(1-\rho)q_{\rm diag}$, and also
$q_T\le(1+\rho)q_{\rm diag}$. Polarization makes the off-diagonal
form continuous in the diagonal form norm. Completing finite-block
sums in either equivalent form norm gives the claimed closed form
and lower bound. Agreement on a common form core identifies it with
the actual closed form.
\end{proof}

A concrete candidate takes signed frequency blocks
$2^jN<|n|\le2^{j+1}N$ separately in each parity, with
$N=C\lambda^p$. It retains the actual signed operator inside every
block and bounds the normalized couplings between blocks. To prove
a polynomial-cutoff result one must verify every within-block positive
matrix and the entire infinite row-sum bound, with explicit window
dependence. A finite block list does not do this. The whole weighted
prime operator is not Hilbert--Schmidt
(Proposition~\ref{prop:v120-not-schatten}), so a global squared-entry
budget cannot be inserted in place of the missing estimate.
No such uniform block certificate is proved here. The next subsection
records the historical test of its natural $N=\lambda^2$ instance at
$\lambda=5$, whose original evidence is not archived. The abstract implication remains
valid, but does not supply a solution for the Weil operator.


\subsection*{Testing the signed block criterion at a polynomial cutoff}
\label{sec:v128-test}
\paragraph{Availability (NS-38).} The original proposals, dyadic
witnesses, congruence factors, $320/448$-bit outputs and independent
verifier are \emph{not archived}. The numerical failure below is a
historical report, not a currently replayable certificate. The intervening
comparison lemma is analytic and is proved in full here.
The historical test fixes $\lambda=5$, $N=\lambda^2=25$ and the literal
parity blocks
\[
 I_0=\{26,\ldots,50\},\quad I_1=\{51,\ldots,100\},\quad
 I_2=\{101,\ldots,200\},\quad I_3=\{201,\ldots,400\}.
\]
They begin the proposed infinite dyadic partition. This is a new
signed-block construction, not a continuation of the unchanged
$Q_{16}$ metric reported to fail in Proposition~\ref{prop:v127-metric-fails}.
The rule uses the actual signed diagonal blocks as metrics,
$M_j=T_{jj}$. The next lemma shows that this is the most favorable
choice among the metrics allowed by Proposition~\ref{prop:v127-block-metric}.

\begin{lemma}[Optimal block energies and a finite obstruction to all row weights]
\label{lem:v128-comparison-obstruction}
Suppose $T_{jj}\succ0$ on finite blocks. For any metrics
$0\prec M_j\preceq T_{jj}$ put
\[
 \beta_{jk}(M)=\|M_j^{-1/2}T_{jk}M_k^{-1/2}\|,\qquad j\ne k.
\]
Then $\beta_{jk}(M)\ge\beta_{jk}(T)$. Let a finite symmetric
entrywise nonnegative matrix $H$, with zero diagonal, satisfy
$H_{jk}\le\beta_{jk}(T)$. If $\lambda_{\max}(H)>1$, no choice of
these metrics and positive scalar weights $w_j$ can satisfy
\begin{equation}\label{eq:v128-weighted-row}
 \sup_j\sum_{k\ne j}\beta_{jk}(M)\frac{w_k}{w_j}<1
\end{equation}
on any infinite extension of the given blocks.
\end{lemma}
\begin{proof}
The norm has the variational expression
\[
 \beta_{jk}(M)=\sup_{x,y\ne0}
 \frac{|\ip{T_{jk}y}{x}|}
 {\sqrt{\ip{M_jx}{x}\,\ip{M_ky}{y}}}.
\]
Increasing each metric to $T_{jj}$ increases the denominator and
cannot increase the supremum. If \eqref{eq:v128-weighted-row} held
with upper bound $b<1$, restrict its positive weight vector to the
finite index set of $H$. Entrywise nonnegativity gives $Hw\le bw$.
Consequently $\operatorname{diag}(w)^{-1}H\operatorname{diag}(w)$
has maximum absolute row sum at most $b$. Similarity preserves its
eigenvalues, contradicting $\lambda_{\max}(H)>1$. Adding further
nonnegative norm terms cannot remove this obstruction.
\end{proof}

\begin{proposition}[Reported dyadic failure; original evidence not archived]
\label{prop:v128-dyadic-fails}
The historical computation reports the following assertion. For the four blocks above, each actual signed $T_{jj}$ is strictly
positive in each parity. Nevertheless
Proposition~\ref{prop:v127-block-metric} fails for this partition,
for all metrics $0\prec M_j\preceq T_{jj}$. It still fails after
any positive scalar row reweighting.
\end{proposition}
\begin{proof}[Historical verification account; original evidence not archived]
The historical account reconstructs exact coefficients at $320$ and
$448$ bits, with the
same $96$-term analytic archimedean enclosure used above. Frozen
dyadic inverse-Cholesky proposals $C_j$ reportedly satisfy the outward
Gershgorin certificates $C_jT_{jj}C_j^*\succ(99/100)I$ in every
block. In particular each $C_j$ is invertible and $T_{jj}\succ0$.
Numerical singular vectors are used only to propose exact dyadic
pairs $(x_{jk},y_{jk})$. The missing interval evaluation reportedly proves positive
block energies and lower bounds on
\[
 r_{jk}=\frac{|\ip{T_{jk}y_{jk}}{x_{jk}}|}
 {\sqrt{\ip{T_{jj}x_{jk}}{x_{jk}}
              \ip{T_{kk}y_{jk}}{y_{jk}}}}
 \le\beta_{jk}(T).
\]
The following are exact integer numerators of the reported strict lower bounds
$H_{jk}$ with denominator $10^6$, in the order indicated:
\begin{center}
\begin{tabular}{lrrrrrr}
\toprule
Sector & $01$ & $02$ & $03$ & $12$ & $13$ & $23$\\
\midrule
Even &999998&570610&337575&483821&196957&490303\\
Odd  &999998&486261&324229&614752&289700&487903\\
\bottomrule
\end{tabular}
\end{center}
Use symmetry and zero diagonal to define the two $4\times4$ matrices
$H_+$ and $H_-$. The vector ${\bf1}=(1,1,1,1)^t$ gives the exact
Rayleigh lower bounds
\begin{align*}
 \lambda_{\max}(H_+)&\ge
 \frac{{\bf1}^tH_+{\bf1}}4
 =\frac{6158528}{4000000}=1.539632>1,\\
 \lambda_{\max}(H_-)&\ge
 \frac{{\bf1}^tH_-{\bf1}}4
 =\frac{6405686}{4000000}=1.6014215>1.
\end{align*}
Conditional on the asserted block comparisons,
Lemma~\ref{lem:v128-comparison-obstruction} proves the assertion.
Both precision replays reportedly use identical frozen witnesses and give
identical rational lower matrices. Each pairing is a complete
finite-support pairing in the actual operator; no remote row is
missing from it. Every uncomputed block would add nonnegative terms
to the norm comparison, so verified versions of these reported finite
obstructions would settle failure of the infinite criterion on this partition. The witness files,
congruence factors and independent check are described in
Appendix~\ref{app:v128-test}, but are not archived. The integer
Rayleigh arithmetic displayed here is exact; the missing step is
replaying the asserted lower comparison with the actual blocks.
\end{proof}

These reported numbers concern lower bounds for a norm-comparison eigenvalue,
not values of $1-\lambda_{\max}(\mathcal U,K)$ and not resolution
margins. No complete $\lambda=5$ positivity margin is obtained.
The ordinary block eigenvalues and proposed singular values are not
used as certified eigenvalue estimates. The failure concerns this
fixed cutoff and partition; it does not exclude another polynomial
cutoff or a family that starts at a larger window.

There is no paradox in positive diagonal blocks failing this test.
Even the positive matrix
\[
 T=\tfrac14 I_3+\tfrac34{\bf1}{\bf1}^t
\]
has diagonal blocks $1$ and eigenvalues $1/4,1/4,5/2$, whereas its
pair-norm comparison has off-diagonal entries $3/4$ and largest
eigenvalue $3/2$. The norm comparison discards joint signs and
directions. This exact illustrative example is not an asserted
factorization of the actual Weil blocks. A subsequent signed
criterion must preserve such information, or change the partition;
merely shrinking the admissible block metrics, rotating coordinates
inside them, or reweighting scalar rows could not fix a test with the
reported certified lower comparisons.

\subsection*{A complete simple even ground and its real-zero transform at $\lambda=4$}

The fixed-window sign can be sharpened to spectral ordering without
replacing the complete operator by a compression. We retain the actual
physical window, normalized parity bases and ordinary Hilbert norm of
Proposition~\ref{prop:v114-weil-core}. The following comparison accounts
for the entire change in the tail inverse under a spectral shift.

\begin{lemma}[A shifted Schur lower bound from a complete unshifted certificate]
\label{lem:v140-shifted-schur}
Let $W=\left(\begin{smallmatrix}F&B^*\\B&T\end{smallmatrix}\right)$
be a lower-bounded self-adjoint operator with finite head,
$T\succeq\delta I$, $\delta>0$, and bounded head-to-tail coupling $B$.
Let $G=(V,Z)^t$ be a finite trial map, with $V$ invertible and
$\operatorname{Ran}Z\subset\mathcal D(T)$. Put
\[
 K=G^*WG,\quad H_h=V^*V,\quad H_t=Z^*Z,\quad
 S=F-B^*T^{-1}B.
\]
Suppose $K\succ0$ and $V^*SV\succeq\kappa K$, $\kappa>0$.
For $0<\sigma<\delta$, define
\begin{equation}\label{eq:v140-shifted-lower}
 L_\sigma=\kappa K-\sigma\left[H_h+
       \frac{2}{1-\sigma/\delta}(H_t+\delta^{-1}K)\right].
\end{equation}
Then
\[
 V^*\{F-\sigma I-B^*(T-\sigma I)^{-1}B\}V
 \succeq L_\sigma.
\]
\end{lemma}
\begin{proof}
Set $r=BV+TZ$ and $\mathcal C=r^*T^{-1}r$. Square completion gives
$V^*SV=K-\mathcal C$, hence $0\preceq\mathcal C\preceq K$.
With $Y=T^{-1}BV=-Z+T^{-1}r$, the Hilbert-space square inequality
and $T^{-2}\preceq\delta^{-1}T^{-1}$ give
\[
 Y^*Y\preceq2H_t+2r^*T^{-2}r
             \preceq2H_t+2\delta^{-1}K.
\]
The resolvent identity yields the exact shifted Schur identity
\[
 V^*S_\sigma V=V^*SV-\sigma H_h
       -\sigma Y^*(I-\sigma T^{-1})^{-1}Y.
\]
Spectral calculus gives
$(I-\sigma T^{-1})^{-1}\preceq(1-\sigma/\delta)^{-1}I$;
substitution proves the claim. All inverse actions are bounded,
and $T^{-1}B$ maps the finite head into $\mathcal D(T)$.
No omitted residual row or inverse-shift term has been discarded.
\end{proof}

\begin{proposition}[Complete ground ordering at $\lambda=4$]
\label{prop:v140-ground4}
Let $\mu_0\le\mu_1\le\cdots$ be the eigenvalues, counted with
multiplicity, of the canonical complete $\mathsf W_4$ in its ordinary
Hilbert norm. Its ground state is simple and even, and
\begin{equation}\label{eq:v140-ground4-bounds}
 0<\mu_0<2.454\cdot10^{-75},\qquad \mu_1>10^{-73}.
\end{equation}
In particular its spectral gap is greater than $9.7546\cdot10^{-74}$.
\end{proposition}
\begin{proof}[Complete certificates and shifted inertia]
Split at $E_{16}$ separately in the two parities.
Proposition~\ref{prop:v120-lambda4-tail} gives the complete tail floor
$\delta=1.7940\cdot10^{-8}$. Use exactly the frozen $4096$-support
trial matrices of Propositions~\ref{prop:v125-even-complete} and
\ref{prop:v126-odd-complement}. Their head dimensions are $17$ and
$16$. The certified complete Schur inequalities in these same
coordinates are
\[
 V_{\rm ev}^*S_{\rm ev}V_{\rm ev}\succeq
       \frac{62629}{100000}K_{\rm ev},\qquad
 V_{\rm od}^*S_{\rm od}V_{\rm od}\succeq
       \frac{428}{1000}K_{\rm od}.
\]
The even constant is rechecked against the complete residual Gram
upper bound. The odd constant follows from the complete
$K$-orthogonal complement gate with $\sigma_{\rm old}=1/1000$
and the different, exactly matching-head $8192$-support trial with
directional bound $429/1000$, using
Proposition~\ref{prop:v126-direction-complement}. These bounds include
every remote Fourier row, not only the rows through $J=65536$.

The new verifier reconstructs the exact dyadic trial coefficients
and computes $H_h,H_t$ without a midpoint solve or a normalized-coordinate
replacement of the ordinary Gram. Applying
Lemma~\ref{lem:v140-shifted-schur} with $\sigma=10^{-73}$ gives:
\begin{center}
\begin{tabular}{lrrrr}
\toprule
sector & precision & dimension & negative pivots & positive pivots\\
\midrule
even & $1024$ & $17$ & $1$ & $16$\\
even & $1280$ & $17$ & $1$ & $16$\\
odd & $1024$ & $16$ & $0$ & $16$\\
odd & $1280$ & $16$ & $0$ & $16$\\
\bottomrule
\end{tabular}
\end{center}
Every outward interval pivot excludes zero. The exact even trial
$f=G_{\rm ev}e_{16}$ (zero-based column index) has
\[
 2.45361754930\cdot10^{-75}
 <\frac{K_{{\rm ev},16,16}}{(H_h+H_t)_{{\rm ev},16,16}}
 <2.45361754931\cdot10^{-75}.
\]
This is its complete Rayleigh quotient because $f$ is finite
supported; no omitted input coefficient is present.

The finite lower even Schur matrix has exactly one negative eigenvalue
and all others positive, so the exact shifted even Schur form has
at most one nonpositive eigenvalue. The even trial below $\sigma$
forces one strictly negative eigenvalue by square completion.
The odd shifted Schur form is positive definite.
Since $T-\sigma I\succ0$, its triangular square-completion map is
boundedly invertible and preserves the form domain. It identifies
negative indices and identifies the kernels of the complete shifted
operators with the finite Schur kernels. Compact resolvent, proved
in Proposition~\ref{prop:v114-weil-core}, therefore gives exactly one
complete eigenvalue below $\sigma$ and none at $\sigma$. It is even
and simple. Previous complete positivity gives $\mu_0>0$; the
Rayleigh bound gives the asserted upper bound.

The new gates and exact Grams are evaluated at $1024$ and $1280$
bits. They reuse the saved complete even residual enclosures at
$768/896$ bits and the saved odd complement enclosure at $768$ bits;
the matching odd directional enclosures are the $1024/1280$-bit
ones. The small prior final gates and exact shared-head equality
are replayed. This is not represented as a fresh assembly of the
old infinite-tail ingredients. Full input hashes and all pivots are
recorded under \texttt{evidence/v140/}.
\end{proof}

\begin{corollary}[Real zeros of the complete ground transform at $\lambda=4$]
\label{cor:v140-real-zeros}
Let $\xi_4$ be a normalized ground eigenfunction of the complete
$\mathsf W_4$ in centered logarithmic coordinates, extended by zero
outside $[-\log4,\log4]$. The nonzero entire function
\[
 \widehat\xi_4(z)=\frac1{\sqrt{2\pi}}
       \int_{-\log4}^{\log4}\xi_4(x)e^{-izx}\,dx
\]
has only real zeros.
\end{corollary}
\begin{proof}
We check the precise distribution hypothesis of
\cite[Theorem~6.1 and Remark~4.3]{ConnesVanSuijlekom2025}.
Put $L=2\log4$, $w_m=\Lambda(m)/\sqrt m$,
$\rho(t)=e^{t/2}/(2\sinh t)$ and
$c_L=\log(4\pi)+\gamma_{\rm E}+\log\tanh(L/2)$.
For a smooth function $h$ on $[0,L]$ define the real distribution
\begin{align*}
 \mathcal D_L(h)={}&\int_0^L2\cosh(t/2)h(t)\,dt
       -\sum_{1<m\le16}w_m h(\log m)\\
 &-\int_0^L\rho(t)\{h(t)-e^{-t/2}h(0)\}\,dt
       -\tfrac12c_Lh(0).
\end{align*}
The canceled numerator is $O(t)$, so the origin integral converges
and this distribution has order at most one. For a trigonometric
polynomial $f$ on $[0,L]$, extend by zero and let
$R_f(t)=\int f(x+t)\overline{f(x)}\,dx$. Its symmetrization
$h_f(t)=R_f(t)+R_f(-t)$ is smooth from the right on $[0,L]$.
The physical Weil formula is exactly
$QW_4(f,f)=\mathcal D_L(h_f)$. In particular
$h_f(0)=2\norm f_2^2$ retains the logarithmic diagonal and
$h_f(L)=0$, so the endpoint prime-power convention introduces no
extra term. Polarization preserves the first-slot-linear convention.

Proposition~\ref{prop:v114-weil-core} supplies lower boundedness
and essential self-adjointness on trigonometric polynomials, exactly
the operator-core hypothesis of the cited theorem.
Proposition~\ref{prop:v140-ground4} supplies the simple isolated
global minimum with reflection-even eigenfunction. Theorem~6.1
therefore applies. Translation from the centered interval to $[0,L]$
and removal of our unitary normalization multiply the transform by
$\sqrt{2\pi}e^{-izL/2}$, which has no zeros. The asserted zero set
is unchanged.
\end{proof}

The corollary concerns the actual complete ground transform, not a
finite Ritz vector, a repaired source proxy or Riemann's $\Xi$.
It asserts neither simplicity of its zeros nor convergence of these
functions to $\Xi$. The small-window spectral question NS-5 is
closed at $\lambda=4$; the cofinal G2 estimate, endpoint-transfer
qualifications and explicit sampler graph defect are unchanged.
No G2 gap or RH claim follows.


\subsection*{An external numerical semantic lock}

\begin{lemma}[The centering phase in the Fourier dictionary]
\label{lem:v141-centering}
Put $a=L/2$, $\omega_n=2\pi n/L$, and let $v=(v_0,\ldots,v_N)$
be coefficients in the actual unshifted even basis on $[0,L]$:
\[
 f_v(y)=\frac{v_0}{\sqrt L}
       +\sqrt{\frac2L}\sum_{n=1}^Nv_n\cos(\omega_n y).
\]
Its centered representative $g_v(x)=f_v(x+a)$ on $[-a,a]$ has
cosine coefficients $(-1)^nv_n$ for $n\ge1$. For the nonunitary
Fourier transform used by CCM,
\begin{equation}\label{eq:v141-transform}
 \widehat g_v(z)=\frac{2\sin(az)}{\sqrt L}
 \left\{\frac{v_0}{z}+\sqrt2\sum_{n=1}^N
          \frac{v_nz}{z^2-\omega_n^2}\right\}.
\end{equation}
All apparent poles are removable. In particular, for $1\le k\le N$,
\begin{equation}\label{eq:v141-lattice-value}
 \widehat g_v(\omega_k)=\sqrt{L/2}\,(-1)^kv_k.
\end{equation}
Hence the sine factor alone does not impose lattice zeros in the
retained band.
\end{lemma}
\begin{proof}
The phase is the exact identity
$\cos(\omega_n(x+a))=(-1)^n\cos(\omega_nx)$.
Integrating each cosine against $e^{-izx}$ cancels this phase
against the factor $(-1)^n$ in the sine numerators, giving
\eqref{eq:v141-transform}. Orthogonality of exponentials, or the
removable-pole limit, gives \eqref{eq:v141-lattice-value}.
This is CCM's finite partial-fraction formula in orthonormal parity
coordinates. Unitary Fourier normalization multiplies it by
$(2\pi)^{-1/2}$; translating the whole function multiplies its
transform by a zero-free exponential. Neither changes its zeros.
\end{proof}

\begin{proposition}[A certified finite-compression semantic lock]
\label{prop:v141-ccm-lock}
At the already studied window $\lambda=3$, take $N=120$ as in
\cite[Section 6, Figure 1]{ConnesConsaniMoscovici2025}. Assemble
both parity compressions with the existing exact coefficient formulas.
Their common global lowest eigenvalue is simple and even. Normalize
its even coefficient vector by $v_0=1$. For each of the first eight
positive zeta ordinates $\gamma_k$, the transform
\eqref{eq:v141-transform} has a unique root in the archived local
interval $I_k$ of radius at most $1.000001\cdot10^{-110}$.
Its positive discrepancy from $\gamma_k$ rounds to the corresponding
published entry:
\begin{center}
\begin{tabular}{rrr}
\toprule
$k$ & reproduced discrepancy (rounded) & CCM Figure 1\\
\midrule
1 & $1.582329697193127\cdot10^{-34}$ & $1.6\cdot10^{-34}$\\
2 & $2.060550289594801\cdot10^{-31}$ & $2.1\cdot10^{-31}$\\
3 & $1.460904856368295\cdot10^{-29}$ & $1.5\cdot10^{-29}$\\
4 & $8.293686860378121\cdot10^{-27}$ & $8.3\cdot10^{-27}$\\
5 & $1.279455965166440\cdot10^{-25}$ & $1.3\cdot10^{-25}$\\
6 & $1.174389453606719\cdot10^{-23}$ & $1.2\cdot10^{-23}$\\
7 & $7.528486586592898\cdot10^{-22}$ & $7.5\cdot10^{-22}$\\
8 & $6.565874264666424\cdot10^{-21}$ & $6.6\cdot10^{-21}$\\
\bottomrule
\end{tabular}
\end{center}
This is a local finite-compression root certificate, not an exhaustive
ordering of all transform zeros or a transfer to the complete ground.

\paragraph{Normalization scope of the numerical test.}
For any $c>0$ and real $s$, replacing the assembled finite matrix
$A$ by $cA+sI$ preserves its ground eigenspace and hence these
transform zeros. Agreement with the published zero discrepancies
cannot by itself test an overall positive scalar or an identity shift.
The identification of the full form normalization rests separately on
the defining coefficient formulas, not on this numerical benchmark.
\end{proposition}
\begin{proof}[Interval verification]
The verifier in \texttt{evidence/v141/} uses the original archived
coefficient assembler, \texttt{assembly\_general.py}, without modification.
It reassembles $121$ even and $120$ odd coordinates using the
$128$-term archimedean formula and its explicit geometric remainder.
The logarithmic diagonal, normalized zero mode, pole term and all
prime powers below $9$ are retained. The endpoint contribution at
$9$ vanishes on the autocorrelations, so the strict prime cutoff
agrees with CCM's inclusive convention.

Arb's verified eigenpair routine isolates all eigenvalues and proves
that the chosen even value lies strictly below every other even and
odd value. Verified eigenvector ratios enclose the real coefficients.
Midpoint root finding only proposes decimal centers. Acceptance uses
opposite interval signs of the rational factor at the endpoints,
a derivative interval excluding zero on the entire bracket, and a
sine interval excluding zero there. These give existence and local
uniqueness by the intermediate value theorem and monotonicity.
Arb enclosures of the fixed low zeta ordinates then give the saved
discrepancy intervals, each strictly inside its published rounding
cell. Every gate is replayed at $768$ and $1024$ bits, at the same
$\lambda,N$ and series order. Direct termwise transforms provide an
additional algebraic consistency check. No midpoint solve is used.
\end{proof}

\paragraph{What had to change.}
The earlier NS-14 diagnostic applied unshifted eigenvector coefficients
to centered cosines without the required $(-1)^n$ conversion. It
therefore evaluated a different window function. In addition, a small
coefficient is not zero: \eqref{eq:v141-lattice-value} precludes the
inference from concentrated low-mode mass to exact retained-band
lattice zeros. The new certificate disproves that diagnostic inference
at $\lambda=3$ and reproduces CCM's fixed-$N$ benchmark. It does not
establish complete-ground zero locations at $\lambda=3$ or $4$, nor
convergence to $\Xi$. Proposition~\ref{prop:v142-ground-zero} below
adds one coarse complete-ground local enclosure at $\lambda=4$;
high-accuracy discrepancy transfer and cofinal convergence remain open.
G2 and RH remain open. No new window is computed.


\subsection*{One local zero of the complete ground transform at $\lambda=4$}

The finite-compression semantic lock does not itself locate complete-ground
zeros. The following transfer uses the complete even-sector spectral
separation and the energy of an exact archived trial. No new window or
finite eigenvector is computed.

\begin{lemma}[Energy control of a ground projection]
\label{lem:v142-energy-projection}
Let $W\succeq0$ be self-adjoint on a real Hilbert space of functions in
$L^2([-a,a])$, with a simple normalized ground $u$. Suppose an invariant
closed subspace contains $u$ and its remaining spectrum lies in
$[b,\infty)$, where $b>0$. For a real unit trial $f$ in that subspace and
in the form domain, put $\rho=q_W[f]<b$, $h=P_u f$, and
$\eta=\sqrt{\rho/b}$. Then $h\ne0$ and
\[
 \norm{f-h}_2\le\eta.
\]
For the nonunitary Fourier transform and every real $t$,
\begin{equation}\label{eq:v142-transform-budget}
 |\widehat f(t)-\widehat h(t)|\le\sqrt{2a}\,\eta,
 \qquad
 |\widehat f'(t)-\widehat h'(t)|\le\sqrt{2a^3/3}\,\eta.
\end{equation}
If both functions are even, opposite endpoint signs for $\widehat h$
and a strictly positive derivative throughout an interval prove a
unique simple zero there, also a zero of $\widehat u$.
\end{lemma}
\begin{proof}
The spectral theorem and nonnegativity give
$q_W[f]\ge b\norm{(I-P_u)f}_2^2$. Since $\eta<1$, the orthogonal
projection $h$ is nonzero. Cauchy--Schwarz applied to $e^{-itx}$ and
$xe^{-itx}$ gives \eqref{eq:v142-transform-budget}. Compact support
justifies differentiation and makes the transforms entire. Reality
and evenness make them real on the real axis. The intermediate value
theorem and strict monotonicity give the asserted unique zero, whose
nonzero derivative makes it simple. Finally $h$ is a nonzero scalar
multiple of $u$. This comparison uses $h=P_u f$, not a separately
renormalized approximation to $u$; no normalization error is omitted.
\end{proof}

\begin{proposition}[A complete-ground local zero near the first zeta ordinate]
\label{prop:v142-ground-zero}
Let $\gamma_1$ be the first positive ordinate of a Riemann zeta zero,
and let $\xi_4$ be the complete ground of
Corollary~\ref{cor:v140-real-zeros}. Then $\widehat\xi_4$ has exactly
one zero in the real interval
\[
 I=(\gamma_1-1/10,\gamma_1+1/10),
 \qquad \gamma_1=14.13472514173469\ldots,
\]
and that zero is simple. This is a local statement: it does not order
all earlier zeros, determine the sign of the discrepancy, identify
the zero with $\gamma_1$, or establish convergence to Riemann's $\Xi$.
\end{proposition}
\begin{proof}[Complete separation and two-precision interval gates]
Use the complete even-sector data of
Proposition~\ref{prop:v140-ground4} without changing the physical window,
the head cut $16$, the exact trial support $4096$, or the archived
complete residual bounds. Apply Lemma~\ref{lem:v140-shifted-schur} with
\[
 \kappa=62629/100000,\quad \delta=1.7940\cdot10^{-8},
 \quad b=10^{-67}.
\]
At both $1024$ and $1280$ bits, the $17$-dimensional lower matrix
$L_b$ has signed LDL pivots
\[
 \underbrace{+\,,\ldots,\,+}_{15},\ -\,,\ +.
\]
Its positive subspace has dimension $16$. The exact even trial below
$b$ supplies a strictly negative direction of the complete shifted
form, while $T-bI\succ0$. Square completion and compact resolvent,
as in Proposition~\ref{prop:v140-ground4}, show that precisely one
complete even eigenvalue is below $b$ and none equals $b$.
Thus every even vector orthogonal to the ground has spectral support
strictly above $10^{-67}$. This is an even-sector separator, not an
improved bound for the first excited eigenvalue across both parities.

Take the same real trial column $G_{\rm ev}e_{16}$, divided by its
ordinary norm, and center it on $[-\log4,\log4]$. Its complete
Rayleigh quotient is
\[
 2.45361754930\cdot10^{-75}<\rho<2.45361754931\cdot10^{-75}.
\]
With $L=2\log4$, Lemma~\ref{lem:v142-energy-projection} gives the
following outward bounds for $h=P_{\xi_4}f$:
\begin{align*}
 \norm{f-h}_2&<1.566403\cdot10^{-4},\\
 E_0=\sqrt L\sqrt{\rho/b}&<0.000260824,\\
 E_1=\sqrt{L^3/12}\sqrt{\rho/b}&<0.000208757.
\end{align*}
All $4097$ exact trial coefficients are retained in
\eqref{eq:v141-transform}, with its correct centering convention.
Interval evaluation at Arb enclosures of $\gamma_1\pm1/10$ gives
\[
 \widehat f(\gamma_1-1/10)+E_0<-0.00006731,
 \qquad
 \widehat f(\gamma_1+1/10)-E_0>0.00003349.
\]
The differentiated formula, evaluated on eight covering subintervals
of length $1/40$, gives
\[
 \widehat f'(t)-E_1>0.0016\quad(t\in\overline I).
\]
These three strict inequalities transfer to $\widehat h$ and prove
the result. Multiplication by the unitary Fourier factor or by the
nonzero ground-projection scalar preserves the zero and its simplicity.

Both precision runs reconstruct the exact dyadic trial and ordinary
Grams. The inherited even residual enclosures at $768/896$ bits
include every row beyond $J=65536$; their small complete Schur gate
is replayed, but their full assemblies are not recomputed.
A second implementation checks the shifted negative index with
principal determinants and the transform with termwise centered
cosine integrals, including derivative signs on $32$ subintervals.
Saved decimal intervals are reparsed to check that serialization
preserves all gates. Scripts, input hashes and reports are in
\texttt{evidence/v142/}. This implementation cross-check is not an
external audit.
\end{proof}

\paragraph{What had to change, and what remains open.}
Using the old global separator $10^{-73}$ in the same energy comparison
would give $E_0\simeq0.260823$, too large to certify either endpoint
sign. Restricting the comparison to the invariant even sector permits
$10^{-67}$ and reduces this error bound by a factor of $1000$.
The obstruction was the coarse separation bound, not a missing zero
of the complete ground transform. The resulting $0.1$ enclosure
still does not transfer the minute CCM finite-compression discrepancies.
A sharper energy or residual-to-transform estimate would be needed
for that task. Cofinal convergence, G2 and RH remain open.


\subsection*{A complete first-zero enclosure and its archived resolution (NS-19)}

NS-19 records the strongest enclosure extracted here from the archived
comparison, and quantifies its resolution. The former signed, factor-ten
target is no longer pursued. The complete discrepancy is not assigned the
scale of a finite-compression benchmark.

\begin{lemma}[Evaluation bound from the archived Schur comparison]
\label{lem:v143-evaluator-dual}
Use the block operator and trial map of
Lemma~\ref{lem:v140-shifted-schur}. Suppose $T\succeq\delta I$,
$K\succ0$, and $V^*SV\succeq\kappa K$ with $0<\kappa<1$.
For a bounded real functional $\ell$, write $p=\ell G$ and let
$\ell_t$ be its restriction to the tail in ordinary coordinates. Put
\[
 A^2=pK^{-1}p^*,\qquad B^2=\delta^{-1}\norm{\ell_t}^2,
 \qquad
 C_\ell=\frac{A^2+B^2+2\sqrt{1-\kappa}\,AB}{\kappa}.
\]
For every even $v$ in the complete form domain,
\[
 |\ell(v)|^2\le C_\ell q_W[v].
\]
\end{lemma}
\begin{proof}
Write $v=Gx+(0,y)^t$, possible since $V$ is invertible, and put
$r=BV+TZ$, $X^2=x^*Kx$, $Y^2=q_T[y]$. The complete Schur identity
$r^*T^{-1}r=K-V^*SV$ gives
\[
 q_W[v]\ge X^2+Y^2-2\sqrt{1-\kappa}\,XY,
 \qquad |\ell(v)|\le AX+BY.
\]
The first inequality retains the cross term through its complete
inverse-weighted bound. The positive two-by-two matrix with diagonal
$1$ and off-diagonal $-\sqrt{1-\kappa}$ has determinant $\kappa$.
Applying Cauchy--Schwarz in that matrix metric gives the displayed
constant. This uses the already certified tail floor, not a new tail metric.
\end{proof}

\begin{proposition}[Complete first positive zero: an absolute bound]
\label{prop:v143-first-zero-bound}
Assume the complete nonnegative form $\mathsf W_4$ has its certified simple
real even ground, that the remainder of its even spectrum lies in
$[b,\infty)$ with $b=10^{-67}$, and that the archived Schur hypotheses
of Lemma~\ref{lem:v143-evaluator-dual} hold with the constants below.
Let $f$ be the normalized exact column $16$ of the archived $4097$-row
trial, $\rho=q_W[f]$, and $h=P_{\xi_4}f$. Assume the archived local
endpoint gates and $\widehat h'>d=0.0016$ throughout
$[\gamma_1-0.1,\gamma_1+0.1]$. The new gates, including the earlier-zero
exclusion below, are verified at $1024$ and $1280$ bits. The tail floor
$\delta$ is inherited from the $320$-bit v1.26 certificate; $K$ and the
Schur margin $\kappa$ use the v1.25/v1.26 ingredients at $768/896$ bits,
with only their small final gates replayed. These inherited full
assemblies are not recomputed at $1024/1280$ bits.
Then the first positive zero $z_1(4)$ of the complete ground transform
is simple and
\begin{equation}\label{eq:v143-first-zero-absolute}
 -8.752082\cdot10^{-33}<z_1(4)-\gamma_1<8.752082\cdot10^{-33}.
\end{equation}
The interval includes zero; neither the sign of the discrepancy nor
a nonzero lower bound for its magnitude is certified.
\end{proposition}
\begin{proof}[Interval evaluation at $1024$ and $1280$ bits]
Set $L=2\log4$ and use the real functional
$\ell(v)=\int_{-L/2}^{L/2}v(x)\cos(\gamma_1x)\,dx$.
Use the same $G,K,\kappa=62629/100000$ and
$\delta=1.7940\cdot10^{-8}$ as above. The vector $p$ is evaluated
using every trial coefficient, and the solve with $K$ is certified.
For the normalized unshifted cosine basis, the head coordinates of
$\ell$ are
\[
 c_0=\frac{2\sin(\gamma_1L/2)}{\sqrt L\,\gamma_1},\qquad
 c_n=\frac{2\sqrt2\sin(\gamma_1L/2)}{\sqrt L}
       \frac{\gamma_1}{\gamma_1^2-(2\pi n/L)^2}.
\]
Parseval supplies the entire omitted evaluator tail, without truncation:
\[
 \norm{\ell_t}^2=\frac L2+\frac{\sin(\gamma_1L)}{2\gamma_1}
                    -\sum_{n=0}^{16}c_n^2.
\]
The interval calculations give
\[
 A^2\in(0.41411,0.41412),\qquad
 B^2\in(49877,49878),\qquad C_\ell\in(79920.06156,79920.06158)<79921.
\]
For the same unit trial $f$ and $h=P_{\xi_4}f$ used in
Proposition~\ref{prop:v142-ground-zero}, nonnegativity gives
$q_W[h]\le q_W[f]=\rho$. Hence
\[
 |\widehat h(\gamma_1)|<1.400334\cdot10^{-35}.
\]
The already verified $\widehat h'>0.0016$ on
$[\gamma_1-0.1,\gamma_1+0.1]$, the mean value theorem and its unique
root give the displacement budget
\[
 |z-\gamma_1|<8.752082\cdot10^{-33}<9\cdot10^{-33}.
\]

To identify this local root as the first positive zero, the verifier
also covers $[0,\gamma_1-0.1]$ with $304$ dyadic subintervals.
On each, direct sinc integrals for trial modes $0$ through $64$ give
an upper bound strictly below zero after adding both the old
complete-ground projection error and
$\sqrt L(\sum_{n=65}^{4096}|f_n|^2)^{1/2}$. This last term accounts
for every omitted coefficient of the exact finite trial; it is not
an omitted infinite eigenfunction tail. The projection estimate
already covers that entire tail. These signs exclude all earlier
nonnegative transform values. The previous local uniqueness and
simplicity gates finish the claim.

Both precisions use the same trial and inherited complete residual
bounds. A second implementation checks the two-by-two dual solve,
recomputes all earlier-zero signs, verifies exact dyadic coverage
with rational arithmetic, and checks the serialized enclosures.
The files \texttt{discrepancy\_budget\_b1024.json} and
\texttt{discrepancy\_budget\_b1280.json} under \texttt{evidence/v143/}
record the result and explicitly mark the signed target unresolved.
\end{proof}

\begin{proposition}[Resolution of the archived bounds and the missing energy direction]
\label{prop:v143-resolution}
Retain the hypotheses, normalization and functional $\ell$ of
Proposition~\ref{prop:v143-first-zero-bound}. Write
\[
 \rho=2.45361754930714\ldots\cdot10^{-75},\qquad
 C_\ell=79920.06156776\ldots,\qquad d=0.0016.
\]
The two decimal ellipses here denote the enclosures in the archived
$1024/1280$-bit reports, not exact midpoint inputs.
The direct bound and its scalar resolution budget are
\begin{equation}\label{eq:v143-direct-budget}
 |z_1(4)-\gamma_1|\le R_0:=\frac{\sqrt{C_\ell\rho}}d,
 \qquad R_0=8.75208157392882\ldots\cdot10^{-33}.
\end{equation}
In this formula the limiting energy input is the \emph{upper} budget
$q_W[h]\le\rho$, multiplied by the complete evaluator constant
$C_\ell$. It is not the ordinary projection error
$\sqrt{\rho/b}=1.5664027417\ldots\cdot10^{-4}$.
A lower bound on the ground energy alone does not reduce the direct
formula \eqref{eq:v143-direct-budget}, but can improve the separate
transfer enclosure toward its frozen-trial floor
$|\widehat f(\gamma_1)|/d\simeq6.65\cdot10^{-38}$, as quantified below.

If an additional certified lower bound $0\le\underline\mu\le\mu_0$
is supplied, set $\Delta=\rho-\underline\mu$ and $e=f-h$. Then
\begin{align}
 \norm e^2&\le\frac{\Delta}{b-\underline\mu},\label{eq:v143-lower-norm}\\
 q_W[e]&\le\frac{b\Delta}{b-\underline\mu},\label{eq:v143-lower-energy}\\
 |\ell(f)-\ell(h)|&\le
 \sqrt{\frac{C_\ell b\Delta}{b-\underline\mu}}.
 \label{eq:v143-lower-transfer}
\end{align}
Thus allocating a root-resolution budget $\varepsilon=10^{-71}$
entirely to this transfer term requires the scalar condition
\begin{equation}\label{eq:v143-required-lower}
 \Delta\le
 \frac{(d\varepsilon)^2(b-\rho)}{C_\ell b-(d\varepsilon)^2}
 =3.20320065696758\ldots\cdot10^{-153}.
\end{equation}
The rounded condition $\rho-\underline\mu\le3.2031\cdot10^{-153}$
is sufficient for that transfer budget. This is a relative gap of
about $1.3055\cdot10^{-78}$ in the Rayleigh energy, not a lower
bound of size $10^{-153}$. No lower estimate this close to $\rho$
is present in the archive. With only the old ordinary-norm projection
inequality, the corresponding gap requirement would instead be
$9.2332480\ldots\cdot10^{-216}$.

Even the condition \eqref{eq:v143-required-lower} would not alone
put the complete zero within $10^{-71}$ of $\gamma_1$: the same
frozen trial has the certified value
\[
 1.0641\cdot10^{-40}<\widehat f(\gamma_1)<1.0642\cdot10^{-40}.
\]
For a $\gamma_1$-centered enclosure one must also control this center
value, since the resulting bound is
\[
 |z_1(4)-\gamma_1|\le
 \frac{|\widehat f(\gamma_1)|+
       \sqrt{C_\ell b\Delta/(b-\underline\mu)}}d.
\]
Even if its transfer-error term vanished, this last estimate would
retain the frozen-trial floor
\[
 \frac{|\widehat f(\gamma_1)|}{d}
 =\frac{1.064123340460\ldots\cdot10^{-40}}{0.0016}
 =6.650770877876\ldots\cdot10^{-38}.
\]
This is a floor of that enclosure formula, not a lower bound on the
actual zero discrepancy, and does not assert that the floor is attainable.
A Temple--Lehmann-type lower-energy estimate is therefore a missing
\emph{transfer} direction, not by itself a solution of the original
$\gamma_1$-centered discrepancy problem.
\end{proposition}
\begin{proof}
The direct bound is the preceding proposition. To make its right-hand
side at most $10^{-71}$ without changing the evaluator constant would
need an upper ground-energy budget of at most
$(d\varepsilon)^2/C_\ell=3.2032007355\ldots\cdot10^{-153}$;
a lower energy estimate does not provide this upper-energy improvement,
but can improve the transfer enclosure, whose frozen-trial floor is
$|\widehat f(\gamma_1)|/d\simeq6.65\cdot10^{-38}$.
For the transfer alternative, spectral orthogonality gives
$\rho=\mu_0\norm h^2+q_W[e]$ and $q_W[e]\ge b\norm e^2$.
As $\norm h^2=1-\norm e^2$,
\[
 \rho\ge\underline\mu+(b-\underline\mu)\norm e^2,
 \qquad
 q_W[e]\le\Delta+\underline\mu\,q_W[e]/b.
\]
These prove the first two estimates; the evaluator lemma applied to $e$
proves the third. Squaring its value-to-zero budget and using
$b-\underline\mu=b-\rho+\Delta$ gives
\eqref{eq:v143-required-lower}. Replacing the evaluator estimate by
$|\ell(e)|\le\sqrt L\norm e$ gives the ordinary-norm threshold
$(d\varepsilon)^2(b-\rho)/(L-(d\varepsilon)^2)$.
The code \texttt{certify\_resolution.py} replays both thresholds,
the outward rounding and the frozen-trial value at both precisions.
The mean value theorem gives the last displayed enclosure. The
claimed extra lower-energy hypothesis has not been verified.
\end{proof}

\begin{proposition}[Finite-window meaning and the scope of the resolution floor]
\label{prop:v143-fixed-window-meaning}
Under Proposition~\ref{prop:v143-first-zero-bound}'s complete-operator,
tail, endpoint and derivative hypotheses, its enclosure is one
finite-window proximity statement of the type sought in CCM's
step~(b). Its scalar radius $R_0$ is the resolution floor of the
\emph{displayed archived comparison with these inputs unchanged}.
The $\lambda=4$, $N=120$ finite-compression benchmark
\[
 2.9250\cdot10^{-71}<z_{1,120}(4)-\gamma_1
                         <2.9251\cdot10^{-71}
\]
is below that resolution by more than $10^{38}$, and also below
that of the archived ordinary projection estimates. Even formally
omitting the nonnegative tail contribution from the same dual
comparison leaves the radius
\[
 \frac{\sqrt{\rho A^2/\kappa}}d
 =2.5174230434\ldots\cdot10^{-35}>10^{35}(3\cdot10^{-71}).
\]
The finite benchmark is not an enclosure of the complete discrepancy.
Consequently these statements neither assign that discrepancy a sign
or a $10^{-71}$ magnitude, nor prove that no different argument using
the same archived data could resolve it. This is a failure of the
stated bounds to resolve the benchmark scale, not a failure of the
complete-ground zero agreement.
\end{proposition}
\begin{proof}
The first enclosure has already been proved. The separate finite
benchmark was certified before the revised NS-19 task, with the
unchanged coefficient assembler and verified eigenpair/root gates;
its $1024/1280$-bit reports are retained as finite-compression evidence
under \texttt{evidence/v143/}. Comparing its outward interval to
$R_0$ and to the formally reduced head-only budget proves the
stated resolution gaps. Neither comparison is a lower bound on the
actual complete discrepancy. No optimality theorem for every
possible use of the v1.25/v1.26 data is asserted.
\end{proof}

CCM's second missing step additionally needs control of how
$k_\lambda$ approximates a scalar multiple of $\xi_\lambda$
\cite[Section 8]{ConnesConsaniMoscovici2025}. The fixed-window
enclosure supplies no cofinal rate, no convergence of further zeros,
and no locally uniform convergence of normalized transforms to
Riemann's $\Xi$. No new window or tail metric is introduced.
NS-19 closes the enclosure-and-resolution accounting; the former
signed target is not claimed. G2 and RH remain open.

\subsection*{What the present diagonal majorant costs as the window grows}

\begin{proposition}[Exponential cutoff cost of the scalar far bound]
\label{prop:v125-cutoff-cost}
Fix $0\le c<1$, put $L=2\log\lambda$, and use the present far
weights
\[
 g_{N+1}=(1-c)\{\log((N+1)/L)-E_L(2\pi(N+1)/L)\}
               -\kappa_{\lambda,N},
 \qquad \kappa_{\lambda,N}\ge M_{\phi,\lambda}
                    \ge\|\mathsf T_{{\rm pr},\lambda}\|.
\]
If $g_{N+1}>0$, then
\begin{equation}\label{eq:v125-cutoff-cost}
 N+1>L\exp\!\left(\frac{M_{\phi,\lambda}}{1-c}\right).
\end{equation}
Consequently, along any unbounded family using these positive weights,
\[
 \liminf_{\lambda\to\infty}
 \frac{\log((N(\lambda)+1)/L)}{\lambda}\ge\frac1{1-c}.
\]
This is a cost of this particular scalar diagonal majorant, not a
lower bound on the cutoff required by every possible argument.
\end{proposition}
\begin{proof}
The error $E_L$ is positive, so positivity of $g_{N+1}$ gives
$(1-c)\log((N+1)/L)>\kappa_{\lambda,N}\ge M_{\phi,\lambda}$.
Use Proposition~\ref{prop:v119-prime-scale}, which proves
$\|\mathsf T_{{\rm pr},\lambda}\|/\lambda\to1$, for the limit.
\end{proof}

% NS-17: cutoff_cost_cite_insert.tex
The doubly exponential shape of \eqref{eq:v125-cutoff-cost}, namely
$N\sim\exp(e^{a})$ in the half-width $a=\log\lambda$, is that of a
published family: Zhu~\cite{Zhu2026} shows that any one-stroke
certificate must resolve frequencies up to $2\pi e^{A_L}$ with
$A_L\sim4e^{L}$ in the half-width $L$, drawing the same conclusion that
the positivity route cannot reach RH unassisted, and
Groskin~\cite{Groskin2026} gives the corresponding brute-cutoff cost
for the Connes--van Suijlekom truncations (an archimedean cutoff of
order $10^{63}$ to resolve a spectral scale $10^{-59}$ at prime cutoff
$c=100$), resolved there instead by a cutoff-free interval $LDL^{T}$
factorization, the factorization technique used here.
The objects differ: Zhu bounds the frequency resolution of a
one-stroke certificate, Groskin the archimedean cutoff of a Galerkin
truncation, and the proposition above the remote cutoff of one specific
scalar majorant; the proposition is placed in that family and is not
presented as new.

For $c=10^{-8}$, freshly evaluating the physical weighted prime bound
and the parity constants at each window gives the following smallest
integer $N$ for which the stated $g_{N+1}$ is positive.  These are
certificates for the remote block only; they do not certify the
intervening inner head or its inverse factors.
\begin{center}
\begin{tabular}{rrrr}
\toprule
$\lambda$ & $M_{\phi,\lambda}$ (rounded) & even $N$ & odd $N$\\
\midrule
3&2.980079&44&209\\
4&4.624042&283&1360\\
5&6.145139&1502&7224\\
6&7.644648&7488&36021\\
8&10.306338&124442&598623\\
\bottomrule
\end{tabular}
\end{center}
Each threshold has positive interval margin, while its preceding
integer has negative margin; monotonicity in $N$ proves minimality
for this formula.  The displayed prime values are rounded summaries,
not the outward endpoints used by the verifier.

The appropriate whole-head diagnostic is
$1-\lambda_{\max}(\mathcal U,K)$, with $K\succ0$ and a
\emph{complete} inverse upper bound $\mathcal U$.  The separately
useful ratios $\lambda_{\min}(K-\mathcal U)/\lambda_{\min}(K)$ and
$\|\mathcal U\|/\lambda_{\min}(K)$ depend on the chosen coordinates;
they should be recorded with the head dimension and condition of
its congruence.  A strong single direction cannot distinguish all
possible causes of a weak whole-head margin.  For example,
\[
 K=\operatorname{diag}(1,\epsilon),\qquad
 U=\operatorname{diag}(\delta,\epsilon(1-\eta)),
 \quad 0<\epsilon,\delta,\eta<1,
\]
has generalized margin $\min(1-\delta,\eta)$, while the headroom
in its first direction is $1/\delta$ and the condition of a
normalizing congruence is $\epsilon^{-1/2}$.  These parameters
may vary independently.  The actual error terms
$\alpha_\lambda$, $e_\lambda^2$ and $t_\lambda^2/\mu_\lambda$
still require complete window-specific witnesses.  The same
numerical cutoff or moment order may be used at different windows
only after its inequalities are verified afresh.


\begin{proposition}[An ordinary-error target for growing windows]
\label{prop:v120-structured-uniform}
For each window, let the complete form have an outer finite head
and positive complement $T_\lambda$, with bounded coupling
$B_\lambda$. For a finite-support solve $X_\lambda$ put
\[
 \begin{split}
 K_{X,\lambda}&=F_\lambda-B_\lambda^*X_\lambda
       -X_\lambda^*B_\lambda+X_\lambda^*T_\lambda X_\lambda,\\
 r_\lambda&=B_\lambda-T_\lambda X_\lambda.
 \end{split}
\]
Assume the inner tail admits the factors of
Proposition~\ref{prop:v120-structured-inverse}, and write
$r_\lambda=(h_\lambda,k_\lambda)$ in that inner decomposition.
If, in operator norm on the entire outer head,
\[
 \begin{gathered}
 K_{X,\lambda}\succeq-\alpha_\lambda I,\qquad
 \norm{D_{g,\lambda}^{-1/2}k_\lambda}\le e_\lambda,\\
 \norm{C_\lambda(h_\lambda-Z_\lambda^*k_\lambda
       -R_\lambda^*D_{g,\lambda}^{-1}k_\lambda)}\le t_\lambda,
 \qquad \alpha_\lambda\ge0,
 \end{gathered}
\]
then the whole closed form satisfies
\begin{equation}\label{eq:v120-structured-uniform}
 \mathsf W_\lambda\succeq
 -\left(\alpha_\lambda+e_\lambda^2+
               t_\lambda^2/\mu_\lambda\right)I.
\end{equation}
Thus convergence of the displayed error to zero along a cofinal
window family is sufficient for the large-support sign criterion.
No such convergence is asserted here.
\end{proposition}
\begin{proof}
The exact outer Schur matrix is
\[
 S_\lambda=K_{X,\lambda}-r_\lambda^*T_\lambda^{-1}r_\lambda.
\]
Apply \eqref{eq:v120-structured-inverse} to each outer head vector.
Completing the outer form square leaves $S_\lambda$ on the head
and a nonnegative complementary square. Since the head projection
is orthogonal, its norm is at most the full ordinary norm, proving
\eqref{eq:v120-structured-uniform}.
\end{proof}

Two scales cannot be suppressed in applying this criterion.
If only individual bounds
$\norm{D^{-1/2}r e_j}\le\delta$ are known for a head of
dimension $d$, the valid conclusion is
$\norm{D^{-1/2}r}^2\le d\delta^2$. Equal residual columns
attain this factor. In particular, a single source estimate or
columnwise decay does not prove the required growing-head bound.
Likewise, $C\mathcal L C^*\succeq-\eta I$ implies only
\[
 \mathcal L\succeq-\eta C^{-1}C^{-*}
             \succeq-\eta\norm{C^{-1}}^2I.
\]
Congruence preserves positivity, but it does not preserve the size
of an ordinary negative error. Using merely $T\succeq cD$
yields the coarser correction
$c^{-1}\norm{D^{-1/2}r}^2$, and its division by $c$ is sharp
already for a scalar tail $T=cD$.

For clarity, the remote truncation itself has an explicit uniform
budget. If $G$ has Fourier support $|m|\le M$, $J>M$, and
$|b_n|\le B_*$, the actual off-diagonal divided difference gives
\begin{equation}\label{eq:v120-separated-uniform-budget}
 (Q_J\mathsf W_\lambda G)^*(Q_J\mathsf W_\lambda G)
 \preceq\frac{8B_*^2(2M+1)}{J-M}\,G^*G.
\end{equation}
Its proof bounds the squared Hilbert--Schmidt norm of the separated
row/column block by
$4B_*^2(2M+1)\sum_{|n|>J}(|n|-M)^{-2}$ and uses the
integral bound $\sum_{n>J}(n-M)^{-2}\le1/(J-M)$.
The diagonal is absent. If the diagonal lower weights $g_n$ are increasing,
dividing by $g_{J+1}$ gives the weighted Gram bound. For
$G=[I;-X]$ its exact final factor is $I+X^*X$, with no extra
number-of-columns factor. The moment bound is usually much sharper,
but a growing moment order controls only its geometric remainder;
it does not force the leading moment Gram to vanish.

These estimates control certification of the low-mode correction.
The simultaneous certificate above settles the complete even form at
$\lambda=4$. The complete odd sign is established above, while the corresponding
growing-window head/inverse estimates remain unproved; positivity of a remote tail
alone does not supply them.



\subsection*{Physical concentration blocks and the repair topology}
The Slepian concentration operator acts on a source space, whereas
$A_a$ acts on the logarithmic physical window. Fix $\lambda=e^a$,
the exact source parameter $c=2\pi\lambda^2$, and thresholds
$\delta_\lambda=\lambda^{-K}$, $K>0$. The source eigenvalues
$\theta_j(c)$ in the plunge interval
$\delta_\lambda<\theta_j(c)<1-\delta_\lambda$ obey
\[
 \#I_{\rm plunge}
 =O\!\left(\log c\,\log(1/\delta_\lambda)\right)=O_K(a^2).
 \tag{G2.5}
\]
This is a source index count \cite{SlepianTransition}, not a physical
orthogonal decomposition. Connes--Consani define physical projections
through finite corrected image spans \cite{ConnesConsani2023}; no
unitary transport of the full source decomposition is supplied by that
construction. The following obstruction makes the distinction necessary
for the fixed repair used in this manuscript.

\begin{proposition}[The fixed point-value repair is not closable in source $L^2$]
\label{prop:v127-nonclosable}
Fix an even $\psi\in C_c^\infty((-1,1))$ with $\psi(0)=1$,
$\int\psi=0$, and actual positive support radius $b<1$.
For $\lambda>1/b$, on smooth even mean-zero sources supported in
$(-\lambda,\lambda)$ put
\[
 Rh=h-h(0)\psi,\qquad
 (\mathcal Eh)(u)=\sqrt u\sum_{n\ge1}h(nu),\qquad
 B_\lambda h(x)=\tfrac12\{\mathcal Eh(e^x)+\mathcal Eh(e^{-x})\}.
\]
The operator $B_\lambda R$ from the mean-zero source $L^2$ space
to $L^2(-\log\lambda,\log\lambda)$ is not closable.
\end{proposition}
\begin{proof}
Let $h_\epsilon(t)=\psi(t/\epsilon)$ with $\epsilon b<1/\lambda$.
It is even and mean zero, $h_\epsilon(0)=1$, and
$\|h_\epsilon\|_2=\sqrt\epsilon\|\psi\|_2\to0$.
For $u\in[1/\lambda,\lambda]$, all $nu$ are outside its support.
Thus $B_\lambda h_\epsilon=0$ and
$B_\lambda R h_\epsilon=-B_\lambda\psi$ for all such $\epsilon$.
This fixed image is nonzero: choose $u$ arbitrarily close to $b$
from below with $u>\max\{b/2,1/\lambda\}$ and $\psi(u)\ne0$.
Only $n=1$ contributes to $\mathcal E\psi(u)$, and
$\mathcal E\psi(1/u)=0$ since $1/u>b$. Hence
$B_\lambda\psi(\log u)=\sqrt u\,\psi(u)/2\ne0$.
Continuity gives a nonzero $L^2$ vector. A graph limit above the zero
source is therefore nonzero, contradicting closability.
\end{proof}

The uncorrected map is bounded at every fixed window: only
$n\le\lambda^2$ contribute, and
\[
 \|\sqrt u\,h(nu)\|_{L^2([1/\lambda,\lambda],du/u)}
 \le n^{-1/2}\|h\|_{L^2(0,\lambda)}.
\]
The obstruction comes from point evaluation in $R$. It does not
invalidate a fixed smooth source, weak G1, or any finite repaired
span. A stronger source topology may control evaluation but requires
its own conversion to the concentration norm. In particular, this
result closes the proposed bounded source-$L^2$ transport route, not
G2 itself.

Here is a literal physical definition that remains valid. Choose the
finite source lists $I_{\rm deep}=\{j:\theta_j\ge1-\delta_\lambda\}$
and $I_{\rm plunge}$, and explicitly specify both moment repairs.
For example, choose an even $\chi\in C_c^\infty((-1,1))$ with
$\int\chi=1$ and $\chi(0)=0$; then
$h-(\int h)\chi-h(0)\psi$ has both required zero moments.
Form the corrected, windowed physical columns $G_{\rm deep}$ and
$G_{\rm plunge}$ using $B_\lambda$ above. Establish their form or
operator domain membership whenever those stronger norms are used.
If $\Gamma_{\rm deep}=G_{\rm deep}^*G_{\rm deep}\succ0$, define
\[
 P_{\rm deep}=G_{\rm deep}\Gamma_{\rm deep}^{-1}G_{\rm deep}^*.
\]
Set $H_{\rm plunge}=(I-P_{\rm deep})G_{\rm plunge}$ and define
$P_{\rm plunge}$ from its positive Gram in the same way. Then
\[
 \mathcal H_a=\operatorname{Ran}P_{\rm deep}\oplus
 \operatorname{Ran}P_{\rm plunge}\oplus\mathcal C_a^{\rm low},
 \quad
 \mathcal C_a^{\rm low}:=\operatorname{Ran}(I-P_{\rm deep}-P_{\rm plunge}).
 \tag{G2.6}
\]
If a Gram is singular, its actual range still defines an orthogonal
projection, but the dimension is smaller. A pseudoinverse defines that
range without proving source-to-image rank or uniform stability.
Consequently the deep source count of order $\lambda^2$ does not yet
prove that physical dimension; (G2.5) only bounds the number of supplied
plunge columns. In the inversion-even construction, the orthogonal
complement also includes any untouched parity sector; this definition
does not silently certify its sign.

For full-rank physical columns $G\subset\mathcal D(A_a)$ and
$\Gamma=G^*G$, the exact growing-block residual is
\begin{equation}\label{eq:v127-gram-residual}
 \|A_aP_G\|^2
 =\|\Gamma^{-1/2}(A_aG)^*(A_aG)\Gamma^{-1/2}\|.
\end{equation}
Indeed $G\Gamma^{-1/2}$ is an isometry onto the range, and its
adjoint has norm one and is onto the coefficient space. The
fixed-source estimates do not bound this growing normalized matrix.
The sufficient coercivity target remains
\[
 P_{\mathcal C_a^{\rm low}}A_aP_{\mathcal C_a^{\rm low}}\succeq c_a I
 \quad\hbox{on }\mathcal C_a^{\rm low},
 \tag{G2.7}
\]
with the other blocks and their couplings controlled. Alternatively,
Proposition~\ref{prop:v118-gap-free} needs only a small ordinary
residual on the removed block and asymptotic nonnegativity of its
whole complement, not a positive gap. The proved rank-one source
residual is available for that rank-one choice; the whole deep block
still requires \eqref{eq:v127-gram-residual}.

\begin{proposition}[The signed weighted target on a physical complement]
\label{prop:v127-weighted-complement}
Choose a positive logarithmic reference diagonal $D\succeq d_0I>0$
so $A_a=D+R_a$ with $R_a$ bounded at the fixed window. For a full-rank
finite physical image matrix $G$ put
\[
 F=D^{-1/2}G,\qquad
 \Pi=I-F(F^*F)^{-1}F^*,\qquad
 K=D^{-1/2}A_aD^{-1/2}.
\]
The last expression is defined by the closed form. For $\eta\ge0$,
\begin{equation}\label{eq:v127-weighted-complement}
 q_{A_a}[f]\ge-\eta\|f\|^2\quad(G^*f=0)
 \quad\Longleftrightarrow\quad
 \Pi(K+\eta D^{-1})\Pi\succeq0.
\end{equation}
The condition on the left ranges over the full form domain.
\end{proposition}
\begin{proof}
Put $z=D^{1/2}f$. This bijects the reference form domain with the
Hilbert space. Then $G^*f=F^*z$,
$q_{A_a}[f]=\ip{Kz}{z}$ and $\|f\|^2=\ip{D^{-1}z}{z}$.
Restriction to $\operatorname{Ran}\Pi$ gives the equivalence.
\end{proof}

The target keeps the signed prime, pole and archimedean terms together.
At fixed window $K-I$ is compact, but compactness gives neither a useful
uniform cutoff nor the required sign. Replacing $\Pi$ by projection off
$G$ imposes the wrong weighted constraint. Replacing $\eta D^{-1}$ by
$\eta I$ yields a negative error relative to the unbounded $D$-energy,
not the required ordinary norm. No cofinal decay or complementary sign
is established here. Proposition~\ref{prop:v127-block-metric} is one
sufficient signed mechanism, but Proposition~\ref{prop:v128-dyadic-fails}
reports failure of its tested $\lambda=5$, $N=25$ dyadic instance,
including admissible metric changes and scalar row reweighting; the
original numerical evidence is not archived. Its
uniform hypotheses remain unproved.

\subsection*{A source-weighted nonlocal capacity target}

This subsection records a different candidate for G2.6. It is not a
Schur-complement estimate, a matrix tail metric or a scalar bound on the
prime norm. Instead it preserves the physical jump geometry and makes
a ground-state change of variables before estimating the signed
potential. Ground-state representations for nonlocal forms are
classical; see Frank--Seiringer~\cite{FrankSeiringer2008}. Nonlocal
Poincare and Lyapunov criteria such as those of Chen--Wang
\cite{ChenWang2013} explain the type of estimate sought below, but do
not supply its arithmetic constant.

Put $I=(-a,a)$, $\lambda=e^a$, and extend functions by zero outside
$I$. Set
\[
 \rho(t)=\frac{e^{t/2}}{2\sinh t},\quad
 \ell_m=\log m,\quad w_m=\frac{\Lambda(m)}{\sqrt m},\quad
 P_\lambda=\sum_{1<m<\lambda^2}w_m,
\]
and
\[
 c_{\rm ar}=\psi(1/4)-\log\pi
 =-\gamma-\log(8\pi)-\frac\pi2.
\]
On the inversion-even sector the exact physical form can be written
\begin{equation}\label{eq:v129-jump-decomposition}
 q_{W_\lambda}[f]=\mathcal E_\lambda[f]
 +(c_{\rm ar}-2P_\lambda)\norm f^2
 +2\left|\ip f{\cosh(\cdot/2)}\right|^2,
\end{equation}
where, for the zero extension $\widetilde f$,
\begin{align}
 \mathcal E_\lambda[f]
 ={}&\frac12\iint_{\mathbb R^2}\rho(|x-y|)
       |\widetilde f(x)-\widetilde f(y)|^2\,dx\,dy \notag\\
 &+\sum_{1<m<\lambda^2}w_m\int_{\mathbb R}
       |\widetilde f(x+\ell_m)-\widetilde f(x)|^2\,dx .
 \label{eq:v129-killed-energy}
\end{align}
The constant $-2P_\lambda$ is accompanied by the exterior killing
terms in \eqref{eq:v129-killed-energy}. Retaining that constant after
discarding those terms would give an incorrect diagonal.

Let $H_\lambda$ be associated with
$\mathcal E_\lambda+(c_{\rm ar}-2P_\lambda)\norm{\cdot}^2$, so that
$W_\lambda=H_\lambda+2|c\rangle\langle c|$ on the even sector, where
$c(x)=\cosh(x/2)$. Fix a real even $h\in C^1(\overline I)$ with
$h\ge m_h>0$ and $\norm h=1$. Define
\[
 d\nu_h=h^2dx,\qquad r_h=\frac{H_\lambda h}{h},
\]
and, for $f=hg$,
\begin{align}
 \mathcal E_{\lambda,h}[g]
 ={}&\frac12\iint_{I^2}\rho(|x-y|)h(x)h(y)
                |g(x)-g(y)|^2\,dx\,dy \notag\\
 &+\sum_{1<m<\lambda^2}w_m
   \int_{\substack{x\in I\\x+\ell_m\in I}}
       h(x)h(x+\ell_m)|g(x+\ell_m)-g(x)|^2\,dx .
 \label{eq:v129-transformed-energy}
\end{align}

\begin{proposition}[Exact nonlocal Picone identity]
\label{prop:v129-picone}
For every even $hg$ in the form domain,
\begin{equation}\label{eq:v129-picone}
 q_{W_\lambda}[hg]=\mathcal E_{\lambda,h}[g]
 +\int_I r_h|g|^2\,d\nu_h
 +2\left|\int_I g(x)h(x)c(x)\,dx\right|^2.
\end{equation}
\end{proposition}
\begin{proof}
On each symmetric archimedean edge and each oriented positive prime
edge use
\[
 |h_xg_x-h_yg_y|^2-h_xh_y|g_x-g_y|^2
 =(h_x-h_y)(h_x|g_x|^2-h_y|g_y|^2).
\]
After integration the exterior terms enter $(H_\lambda h)/h$. This
proves the formula on smooth core functions. Since $h$ and $h^{-1}$
are bounded and $C^1$, multiplication by either preserves the interior
logarithmic form domain. Moreover $r_h$ has only positive logarithmic
endpoint growth from exterior killing. Closure gives
\eqref{eq:v129-picone} on the stated domain.
\end{proof}

The candidate mechanism is the following source-weighted capacity
estimate. It is written with the actual removed vector $u_\lambda$;
there is no assumption that $h$ is that vector or that $u_\lambda$ is
positive.

\begin{proposition}[Capacity absorption sufficient for the gap-free route]
\label{prop:v129-capacity}
Let $u_\lambda$ be a normalized even vector in the operator domain.
Suppose that, for some $0\le\delta_\lambda\le1$ and
$\eta_\lambda\ge0$, every even $g$ with $hg$ in the form domain and
\begin{equation}\label{eq:v129-constraint}
 \int_I g(x)h(x)\overline{u_\lambda(x)}\,dx=0
\end{equation}
satisfies
\begin{align}
 \int_I(r_h)_-|g|^2\,d\nu_h
 \le{}&(1-\delta_\lambda)\mathcal E_{\lambda,h}[g]
 +2\left|\int_I ghc\,dx\right|^2
 +\eta_\lambda\int_I|g|^2\,d\nu_h .
 \label{eq:v129-capacity}
\end{align}
Then
\[
 q_{W_\lambda}[f]\ge-\eta_\lambda\norm f^2
 \quad(f\perp u_\lambda,\ f\text{ even}).
\]
If also $\norm{W_\lambda u_\lambda}\le\epsilon_\lambda$, then the
even-sector conclusion of Proposition~\ref{prop:v118-gap-free} is
\begin{equation}\label{eq:v129-gapfree-conclusion}
 W_\lambda\succeq-\left(\eta_\lambda+\frac{2}{\sqrt3}
 \epsilon_\lambda\right)I.
\end{equation}
Consequently, cofinal $\eta_{\lambda_j},\epsilon_{\lambda_j}\to0$
would give asymptotic nonnegativity in this sector without a positive
spectral gap.
\end{proposition}
\begin{proof}
The constraint is exactly $hg\perp u_\lambda$. Substitute
\eqref{eq:v129-capacity} into \eqref{eq:v129-picone}; the retained
terms $\delta_\lambda\mathcal E_{\lambda,h}[g]$ and
$\int(r_h)_+|g|^2d\nu_h$ are nonnegative. The last assertion is
Proposition~\ref{prop:v118-gap-free}.
\end{proof}

This controls the whole even rank-one complement, hence any physical
low-concentration subspace proved to lie in that complement. We do
\emph{not} identify $\mathcal C_a^{\rm low}$ with it: if the physical
decomposition is not subordinate to $u_\lambda^\perp$, the remaining
blocks and couplings still require estimates. The odd sector is also
untouched; there the pole is
$-2|\ip f{\sinh(\cdot/2)}|^2$ and needs a separate argument.

The new project-level ingredient is the exact reduction to the
weighted jump capacity \eqref{eq:v129-capacity}. It retains the joint
archimedean and prime edge signs and converts source orthogonality into
\eqref{eq:v129-constraint}. The ground-state transform itself is
literature-known, so no worldwide novelty is claimed. Formula
comparison distinguishes this target from the earlier signed
Schur/structured-inverse routes, scalar prime-norm bounds, finite-rank
metric repairs, source-only residual arguments, dyadic pair-norm
metrics and shared-channel common-Gram factorization.

Two adversarial checks delimit the claim. First, $W_\lambda h\ge0$
pointwise is not enough: replacing $H_\lambda h/h$ by
$W_\lambda h/h$ creates a negative weighted variance term. Second,
the historically reported calculation from the saved $\lambda=3$, $N=64$
compression gives a projected ratio with
minimum about $-0.414$ and negative values on about $44.95\%$ of a
dense grid (about $32.33\%$ in $h^2dx$ mass). This falsifies the
stronger finite pointwise-supersolution surrogate, but does not
evaluate the complete physical ratio, which has positive logarithmic
endpoint growth. The transformed finite absorption ratio is one to
binary64 resolution, so it supplies no margin. These are numerical
diagnostics, not certificates. The v1.29 recomputation script, output
and companion review are not archived; the underlying source compression
is available, but does not reproduce these derived values by itself.

Thus the exact remaining lemma is a uniform proof of
\eqref{eq:v129-capacity}, for an explicit zero-free $h_\lambda$ and the
actual source constraint, with $\eta_{\lambda_j}\to0$ along a cofinal
family. A generic Poincare constant, an ordinary residual estimate or
another positive fixed window does not imply this lemma. No G2 sign
gap is closed here.

\subsection*{Exact continuum cancellation and the arithmetic remainder}

The capacity target has a sharper translation-invariant normal form.
It identifies what the PNT continuum model settles exactly and what
remains genuinely arithmetic. Define on $I=(-a,a)$, using zero
extension,
\[
 (K_a^\pm f)(x)=\int_I e^{\pm|x-y|/2}f(y)\,dy,\qquad
 c(x)=\cosh(x/2),\quad s(x)=\sinh(x/2).
\]

\begin{proposition}[The continuum prime channel is cancelled by the even pole]
\label{prop:v130-continuum-cancellation}
On $L^2(I)$,
\begin{equation}\label{eq:v130-kernel-factorization}
 2|c\rangle\langle c|-K_a^+=K_a^-+2|s\rangle\langle s|\succeq0.
\end{equation}
In particular, for even $f$,
\[
 2|\ip f c|^2-\ip{K_a^+f}f=\ip{K_a^-f}f\ge0.
\]
The inequality is strict for nonzero $f$, but has no positive
ordinary-norm gap.
\end{proposition}
\begin{proof}
The kernel identity is
\[
 2\cosh(x/2)\cosh(y/2)-e^{|x-y|/2}
 =e^{-|x-y|/2}+2\sinh(x/2)\sinh(y/2).
\]
The Fourier transform of $e^{-|x|/2}$ is
$(\xi^2+1/4)^{-1}>0$, so its compression $K_a^-$ is strictly
positive. The last rank-one term is positive and vanishes on the even
sector. The compressed convolution is compact, and its eigenvalues
tend to zero; therefore strict positivity supplies no uniform gap.
\end{proof}

Let the exact prime-correlation operator be
\[
 (C_af)(x)=\sum_{1<m<e^{2a}}\frac{\Lambda(m)}{\sqrt m}
 \{\widetilde f(x+\log m)+\widetilde f(x-\log m)\},
\]
and put $R_a=C_a-K_a^+$. For the unitary Fourier transform of the zero
extension, define
\begin{align}
 \mathfrak r_a(\xi)
 :={}&2\sum_{1<m<e^{2a}}\frac{\Lambda(m)}{\sqrt m}
                 \cos(\xi\log m)
       -2\int_0^{2a}e^{t/2}\cos(\xi t)\,dt, \label{eq:v130-r-symbol}\\
 \beta_a(\xi)
 :={}&\Re\psi\!\left(\frac54+\frac{i\xi}{2}\right)-\log\pi
       -\mathfrak r_a(\xi). \label{eq:v130-beta}
\end{align}

\begin{proposition}[Exact even-sector discrepancy symbol]
\label{prop:v130-remainder}
For every even $f$ in the physical form domain,
\begin{equation}\label{eq:v130-symbol-form}
 q_{W_\lambda}[f]
   =\int_{\mathbb R}\beta_a(\xi)
          |\widehat{\widetilde f}(\xi)|^2\,d\xi .
\end{equation}
Consequently either of the following is sufficient for the even
conclusion of Proposition~\ref{prop:v118-gap-free}:
\begin{align}
 \beta_a(\xi)&\ge-\eta_a\quad(\xi\in\mathbb R),\label{eq:v130-pointwise}\\
 \int_{\mathbb R}\beta_a|\widehat{\widetilde f}|^2
   &\ge-\eta_a\norm f^2
   \quad(f\perp u_\lambda,\ f\ {\rm even}).\label{eq:v130-constrained}
\end{align}
If \eqref{eq:v130-constrained} holds along a cofinal family with
$\eta_a\to0$, and the normalized source has
$\norm{W_\lambda u_\lambda}\le\epsilon_a\to0$, then
\[
 W_\lambda\succeq-
 \left(\eta_a+\frac{2}{\sqrt3}\epsilon_a\right)I
\]
on the even sector.
\end{proposition}
\begin{proof}
The prime part of \eqref{eq:v129-jump-decomposition} is
$-\ip{C_af}f$. Replacing the Stieltjes prime measure by $dx$ gives
$K_a^+$ exactly, since $x^{-1/2}dx=e^{t/2}dt$ under $x=e^t$.
Proposition~\ref{prop:v130-continuum-cancellation} replaces the pole
minus $K_a^+$ by $K_a^-$ on even functions.

The archimedean multiplier is
\[
 2\int_0^\infty\rho(t)(1-\cos(\xi t))\,dt+c_{\rm ar}
 =\Re\psi\!\left(\frac14+\frac{i\xi}{2}\right)-\log\pi.
\]
The multiplier of $K_a^-$ is $(\xi^2+1/4)^{-1}$. The recurrence
$\psi(z+1)=\psi(z)+z^{-1}$ combines these two terms into the first
two terms of \eqref{eq:v130-beta}. The discrepancy multiplier is
\eqref{eq:v130-r-symbol}, proving \eqref{eq:v130-symbol-form}.
Plancherel gives \eqref{eq:v130-pointwise}; the constrained version
and the last assertion use Proposition~\ref{prop:v118-gap-free}.
\end{proof}

The tempting next step would be to declare $R_a$ a vanishing
perturbation. That is false even after imposing the actual rank-one
constraint.

\begin{proposition}[A fixed local obstruction to perturbative PNT replacement]
\label{prop:v130-local-no-go}
Let $u_a$ be any normalized family in $L^2(-a,a)$. For all sufficiently
large $a$ there is a normalized real even
$g_a\in C_c^\infty(-a,a)\cap u_a^\perp$ such that
\begin{equation}\label{eq:v130-local-no-go}
 \ip{R_ag_a}{g_a}>\frac{\log2}{\sqrt2}.
\end{equation}
The $g_a$ may be chosen in a fixed finite-dimensional family, so every
fixed archimedean reference-form norm of $g_a$ is bounded uniformly.
Therefore no estimate
\[
 R_a\preceq\eta_a I
 \quad\hbox{or}\quad
 R_a\preceq\alpha_a A_{\rm ref}+\eta_a I
\]
can hold on $u_a^\perp$ with respectively $\eta_a\to0$, or with both
$\alpha_a,\eta_a\to0$, where $A_{\rm ref}$ is any fixed positive
archimedean reference form on the displayed finite-dimensional family.
The first conclusion remains true after removal of any
finite-dimensional subspace at each window.
\end{proposition}
\begin{proof}
Put $\ell=\log2$, $d=1/10$, and choose
$0<\varepsilon<1/100$. For an even
$\phi\in C_c^\infty(-\varepsilon,\varepsilon)$ let
\[
 c_{-1}=c_1=1,\qquad c_0=-2\cosh(d/2),
\]
\[
 g_\phi(x)=\sum_{j=-1}^1c_j
 \{\phi(x-\ell/2-jd)+\phi(x+\ell/2-jd)\}.
\]
The six supports are disjoint and their total diameter is less than
$\log3$. Hence among all prime-power shifts only $\log2$ couples two
supports, and direct translation gives
\[
 \ip{C_ag_\phi}{g_\phi}
   =\frac{\log2}{\sqrt2}\norm{g_\phi}^2.
\]
For either sign,
\[
 \int e^{\pm x/2}g_\phi(x)\,dx=0,
\]
because
$\sum_{j=-1}^1c_je^{\pm jd/2}
=2\cosh(d/2)-2\cosh(d/2)=0$.
Thus both the $c$ and $s$ moments vanish. The kernel identity in
Proposition~\ref{prop:v130-continuum-cancellation} now gives
$\ip{K_a^+g_\phi}{g_\phi}=-\ip{K_a^-g_\phi}{g_\phi}<0$.
This proves the strict inequality in \eqref{eq:v130-local-no-go}
before normalization.

Choose two linearly independent such bumps. Their images form a fixed
two-dimensional space, so for every $u_a$ a nonzero linear combination
is orthogonal to $u_a$. Normalization preserves the strict bound, and
all fixed reference-form norms are uniformly bounded on this
finite-dimensional space. For the last statement, allow $\phi$ to
range over the infinite-dimensional space of even smooth bumps in
$(-\varepsilon,\varepsilon)$. The displayed construction is injective
because its six supports are disjoint. Its intersection with the
orthogonal complement of any finite-dimensional removed space remains
nonzero, and every such vector obeys the same ordinary lower bound.
\end{proof}

Proposition~\ref{prop:v130-local-no-go} does not disprove the full
capacity mechanism: the complete archimedean term may absorb this
fixed local excess with a coefficient that does not tend to zero.
It does prove that ordinary residual smallness, an $o(1)$ relative
perturbation or the continuum cancellation alone cannot close G2.
There is no hidden positive baseline:
\[
 \Re\psi(5/4)-\log\pi\approx-1.372<0.
\]
The exact live lemma is now \eqref{eq:v130-constrained}: a
source-constrained, support-sensitive lower estimate for the full
symbol $\beta_a$, not a norm estimate on $\mathfrak r_a$. The
pointwise strengthening \eqref{eq:v130-pointwise} is numerically false
at the already positive fixed windows, because $\beta_a$ has narrow
negative wells. Thus a proof must use the physical support constraint,
for example through a quantitative uncertainty or concentration
estimate. The relevant classical tools include the sharp
time--frequency concentration theory of Landau--Pollak
\cite{LandauPollak1961}, Nazarov's finite-measure uncertainty
inequality \cite{Nazarov1994}, and Logvinenko--Sereda estimates for
unions of frequency intervals \cite{Kovrijkine2000}. None applies
until the depths, widths, number and spacing of the negative wells of
$\beta_a$ are controlled with their full $a$ dependence; importing a
generic constant would not prove \eqref{eq:v130-constrained}.
Numerical wells are diagnostic only. The odd sector remains separate.
No G2 sign gap is closed.

A historical rational computation reports failure of a global
conditional-positive-definiteness shortcut at a certified positive physical
window. Its original verifier and output are not archived; the analytic
negative-index implication below is conditional on the reported enclosure.

\begin{proposition}[Negative-index implication of a reported enclosure]
\label{prop:v132-global-index}
The unarchived rational computation reports that at $a=\log4$ the
exact symbol in \eqref{eq:v130-beta} satisfies
\[
 -0.64784909<\beta_a(1)<-0.6321938<-\frac35.
\]
If this reported enclosure holds, the full-line multiplier by $\beta_a$ has infinite
negative index, even on the even sector and after imposing any fixed
finite number of linear constraints. This does not contradict the
proved positivity of the physical $\lambda=4$ form.
\end{proposition}
\begin{proof}
The historical account attributes the displayed enclosure to exact
rational interval arithmetic (the code and output are not archived): Machin's formula encloses $\pi$, the atanh series encloses
$\log2$ and all prime weights below the strict cutoff $16$, and the
convergent digamma series at $5/4+i/2$ has its positive tail bounded by
the corresponding integral. All remainders are outward rational
enclosures.

Continuity makes $\beta_a$ negative on a nonempty interval about one.
Choose arbitrarily many smooth even Fourier tests with pairwise
disjoint symmetric supports in that interval. Their inverse transforms
span a negative subspace of the same dimension; finitely many linear
constraints reduce that dimension only finitely. These inverse
transforms are Schwartz but are not supported in $(-a,a)$. Thus the
argument rules out a full-line finite-negative-index extension, not the
physical Paley--Wiener form already certified positive at this window.
\end{proof}

\subsection*{Source-compressed concentration: a rigorous hierarchy and its scalar loss}

The support constraint in \eqref{eq:v130-constrained} can be expressed
without replacing the signed symbol by its pointwise minimum.  Let
$\mathcal H_a^{\rm ev}=L^2_{\rm ev}(-a,a)$, let $\mathcal F_a$ denote
zero extension followed by the unitary Fourier transform, and let
$u_a\in\mathcal H_a^{\rm ev}$ be the normalized repaired source.  On
$\mathcal H_a^{\rm ev}$ put
\[
 P_a=I-|u_a\rangle\langle u_a|,
 \qquad
 \mathcal C_{a,u}(E)
 =P_a\mathcal F_a^*1_E\mathcal F_aP_a
\]
for a symmetric measurable set $E\subset\mathbb R$.  This is the
source-compressed time--frequency concentration operator.  It is a
positive contraction; no Fourier truncation is made in its definition.

\begin{proposition}[Negative-level and weighted concentration criteria]
\label{prop:v131-concentration}
Let $\beta_a$ be \eqref{eq:v130-beta}, put
\[
 D_a=\norm{(\beta_a)_-}_{L^\infty},\qquad
 E_a(t)=\{\xi:\beta_a(\xi)<-t\}\quad(0<t<D_a).
\]
Then $D_a<\infty$, every $E_a(t)$ is bounded, and on
$P_a\mathcal H_a^{\rm ev}$ one has
\begin{equation}\label{eq:v131-layercake}
 q_{W_\lambda}[f]\ge-\eta_a^{\rm op}\norm f^2,
 \qquad
 \eta_a^{\rm op}:=
 \int_0^{D_a}\norm{\mathcal C_{a,u}(E_a(t))}\,dt.
\end{equation}
If $E_a(t)$ has finite measure, then
\begin{align}
 \operatorname {tr}\mathcal C_{a,u}(E_a(t))
  =\int_{E_a(t)}\left\{
    \frac{a+\sin(2a\xi)/(2\xi)}{2\pi}
    -|\widehat{\widetilde u_a}(\xi)|^2\right\}d\xi,
 \label{eq:v131-trace}\\
 \norm{\mathcal C_{a,u}(E_a(t))}
  \le \min\{1,\operatorname {tr}\mathcal C_{a,u}(E_a(t))\}.
 \label{eq:v131-trace-bound}
\end{align}
The quotient at $\xi=0$ is interpreted continuously.

More generally, suppose that symmetric sets $G_{a,1},\ldots,G_{a,J}$
and positive weights $w_{a,j}$ obey the pointwise lower minorant
\begin{equation}\label{eq:v131-step-minorant}
 -D_a+\sum_{j=1}^Jw_{a,j}1_{G_{a,j}}(\xi)
 \le\beta_a(\xi).
\end{equation}
If
\begin{equation}\label{eq:v131-weighted-concentration}
 -D_aP_a+\sum_{j=1}^Jw_{a,j}\mathcal C_{a,u}(G_{a,j})
 \succeq-\eta_aP_a,
\end{equation}
then \eqref{eq:v130-constrained} holds with the same $\eta_a$.
Consequently a cofinal family satisfying
\eqref{eq:v131-step-minorant}--\eqref{eq:v131-weighted-concentration}
with $\eta_a\to0$ proves the even conclusion of
Proposition~\ref{prop:v118-gap-free}, without a positive complement
gap.
\end{proposition}
\begin{proof}
For fixed $a$ the prime sum in \eqref{eq:v130-r-symbol} is finite, the
continuum term is $O_a((1+|\xi|)^{-1})$, and
$\Re\psi(5/4+i\xi/2)=\log|\xi|+O(1)$ as $|\xi|\to\infty$.
Thus $\beta_a$ is continuous, bounded below and tends to $+\infty$;
the first assertions follow.

For $f=P_af$, Tonelli's theorem and the layer-cake identity give
\[
 \int(\beta_a)_-|\widehat{\widetilde f}|^2
 =\int_0^{D_a}\ip{\mathcal C_{a,u}(E_a(t))f}f\,dt.
\]
Dropping the positive part of $\beta_a$ proves
\eqref{eq:v131-layercake}.  The even evaluation vector at frequency
$\xi$ is $(2\pi)^{-1/2}\cos(\xi x)$ on $(-a,a)$; its squared norm is
$\{a+\sin(2a\xi)/(2\xi)\}/(2\pi)$.  Compression off the normalized
source subtracts $|\widehat{\widetilde u_a}(\xi)|^2$.  Integrating the
diagonal gives \eqref{eq:v131-trace}; positivity gives
\eqref{eq:v131-trace-bound}.

Finally, Plancherel and \eqref{eq:v131-step-minorant} yield
\[
 q_{W_\lambda}[f]
 \ge -D_a\norm f^2+
       \sum_{j=1}^Jw_{a,j}\ip{\mathcal C_{a,u}(G_{a,j})f}f.
\]
Equation \eqref{eq:v131-weighted-concentration} proves the claim.
\end{proof}

The scalar layer-cake estimate is rigorous. Historical first tests
reported the following diagnostic losses; their original scripts and
outputs are not archived. The surviving v1.31 report alone is available
at \path{evidence/v131/g2_weighted_concentration_report.md}.  On grids truncated at $|\xi|\le5000$, its already accumulated
coefficients are about $2.1363$, $3.2440$ and $4.5622$ for
$\lambda=3,4,5$ respectively; these are lower diagnostics for the
complete scalar upper bounds, not certificates.  They are large
because \eqref{eq:v131-layercake} discards every favorable value of
$\beta_a$.

A two-level specialization makes the signed budget explicit.  If
$G_a(\theta)=\{\beta_a\ge\theta D_a\}$ and
$\mathcal C_{a,u}(G_a(\theta))\succeq p_aP_a$, then
\begin{equation}\label{eq:v131-two-level}
 q_{W_\lambda}[f]
 \ge D_a\{(1+\theta)p_a-1\}\norm f^2.
\end{equation}
Thus mere thickness of $\{\beta_a\ge0\}$ is insufficient: one needs
$p_a\ge(1-o(1)/D_a)/(1+\theta)$.  A reciprocal-scale thickness
length $C/a$ keeps the exponent in a Logvinenko--Sereda inequality
bounded, but the generic constants from
\cite{Kovrijkine2000} do not supply this near-unit concentration
budget.  This is a failure of that scalar sufficient estimate, not a
no-go for \eqref{eq:v131-weighted-concentration}.

A historical adversarial wave-packet test also reported no falsification
of the weighted mechanism; its original script and output are not archived.  For the even packets
$\cos(\pi x/(2a))\cos(\omega x)$ centred at the deepest located wells,
the numerical complete Rayleigh quotients were approximately
$0.951081$, $2.757308$ and $3.683514$ at $\lambda=3,4,5$, although the
point values of $\beta_a$ there were approximately $-2.37816$,
$-3.41944$ and $-5.01788$.  Direct physical archimedean/prime/pole
quadrature and Fourier integration agreed to about $3\cdot10^{-10}$.
These are non-rigorous falsification diagnostics.  They show exactly
why the favorable superlevels must be retained; they prove neither a
fixed-window sign nor the cofinal estimate.

The live arithmetic lemma is therefore
\eqref{eq:v131-weighted-concentration} for an explicit finite or
convergent hierarchy of superlevel sets, with all $a$ dependence and
$\eta_a\to0$.  It acts on the complete even source complement, the
larger space required by Proposition~\ref{prop:v118-gap-free}; it does
not identify that complement with $C_a^{\rm low}$.  The construction
is a development of candidate G2.6-CAP-01, not a renamed Schur or tail
metric: its operators are physical Paley--Wiener concentration
operators for the exact arithmetic symbol, with the actual source
removed before comparison.  Concentration operators and
Logvinenko--Sereda estimates are classical; only this arithmetic
specialization is new within the project, and no worldwide novelty is
claimed.  The odd sector remains a separate obligation.

\subsection*{The weighted criterion on a level pencil: finite tests and cofinal quantifiers}

This is an exact translation of the existing sufficient criterion, not a
new concentration estimate or a numerical certificate. The pencil in
\texttt{evidence/diag\_true\_symbol/}, Section~7, motivates the test class;
its quadrature and midpoint generalized eigenvalues remain diagnostic only.

\begin{proposition}[Directional loss budget and the scope of the pencil test]
\label{prop:v144-pencil-concentration}
Fix $a=\log4$ and let $H_{48}\subset\mathcal D(q_a)$ be the even
head spanned by cosine modes $0,\ldots,48$ (dimension $49$).
Use zero extension and the unitary Fourier normalization of
\eqref{eq:v135-full-symbol}. Define the exact level forms
\[
 q_a^\pm[f]=\int_{\mathbb R}(\beta_a)_\pm
                    |\widehat{\widetilde f}|^2,
 \qquad q_a=q_a^+-q_a^-.
\]
Let $W,W^+$ be their respective exact head matrices for $q_a,q_a^+$,
and assume $W^+\succ0$. Choose a complete generalized eigenbasis
\[
 Wv_k=\nu_k W^+v_k,\qquad
 V^*W^+V=I,\qquad V^*WV=\operatorname{diag}(\nu_k),
 \quad 0\le k\le48,
\]
identifying each column $v_k$ with its physical head function.
Thus $\nu_k=q_a[v_k]/q_a^+[v_k]$ and
$q_a^-[v_k]/q_a^+[v_k]=1-\nu_k$.
Let $u_a\in\mathcal D(q_a)$ be the actual normalized even source,
$P_a=I-|u_a\rangle\langle u_a|$, and put
$E=H_{48}\cap u_a^\perp$.

Suppose symmetric measurable sets $G_{a,j}$ and positive weights $w_{a,j}$
give the pointwise minorant in \eqref{eq:v131-step-minorant}. Write
\[
 m_a=-D_a+\sum_jw_{a,j}1_{G_{a,j}}\le\beta_a,\qquad
 r_a=\beta_a-m_a\ge0,\qquad
 \mathcal L_a[f]=\int r_a|\widehat{\widetilde f}|^2.
\]
Here $\mathcal L_a$ is the nonnegative energy lost by this minorant;
it is not a new tail metric. For $\eta_a\ge0$ and each pencil
vector satisfying $P_av_k=v_k$, the directional restriction of
\eqref{eq:v131-weighted-concentration} is exactly
\begin{equation}\label{eq:v144-direction-budget}
 \frac{\mathcal L_a[v_k]}{q_a^+[v_k]}
 \le\nu_k+\eta_a\frac{\norm{v_k}^2}{q_a^+[v_k]}.
\end{equation}
In particular, a zero-error test ($\eta_a=0$) must lose no more
than the relative surplus $\nu_k$ in that direction. With nonzero
allowed error, the second term must not be omitted.

This directional demand is necessary, but is not by itself sufficient
on the span. Define matrices $G,R$ by
\[
 c^*Gc=\norm{Vc}^2,\qquad c^*Rc=\mathcal L_a[Vc],
 \qquad b_k=\ip{v_k}{u_a}.
\]
Then the restriction of \eqref{eq:v131-weighted-concentration} to
\emph{all} of $E$ is equivalent to
\begin{equation}\label{eq:v144-pencil-matrix}
 c^*\{\operatorname{diag}(\nu_k)-R+\eta_a G\}c\ge0
 \quad\hbox{for every }c\hbox{ with }\sum_k b_kc_k=0.
\end{equation}
Thus the minorant-loss and ordinary-norm cross terms must also be
controlled. The uncompressed diagnostic does not establish $b_k=0$.
For a nonadmissible $v_k$, substituting $P_av_k$ changes both level
energies; its ratio is not certified to remain $\nu_k$.

Conditionally, if the first twelve \emph{exact} values satisfy
$\nu_k<10^{-8}$ for $0\le k\le11$, their span intersects $u_a^\perp$
in dimension at least eleven, and every nonzero vector in that
intersection has $q_a[f]/q_a^+[f]<10^{-8}$. This conditional statement
uses no numerical threshold as a proved hypothesis.

All statements so far concern one fixed window and a finite head.
The full live criterion instead demands the operator inequality on
the complete source complement along a cofinal family $a\to\infty$,
with uniform control of its permitted errors. The original
Proposition~\ref{prop:v131-concentration} asks for $\eta_a\to0$;
the weaker sufficient objective in
Corollary~\ref{cor:v136-complement-floor} allows
$\sup_a\eta_a<\infty$, with both parity blocks and the stated source
residual/coupling hypotheses. Neither cofinal quantifier follows
from \eqref{eq:v144-direction-budget} or
\eqref{eq:v144-pencil-matrix} at $a=\log4$.
\end{proposition}
\begin{proof}
For $f=P_af$, Plancherel gives the exact identity
\[
 -D_a\norm f^2+\sum_jw_{a,j}
       \ip{\mathcal C_{a,u}(G_{a,j})f}f
 =\int m_a|\widehat{\widetilde f}|^2
 =q_a[f]-\mathcal L_a[f].
\]
Testing the left-hand side against $-\eta_a\norm f^2$ and substituting
$q_a[v_k]=\nu_kq_a^+[v_k]$ proves
\eqref{eq:v144-direction-budget}. For $f=Vc$ the source constraint is
$\sum_kb_kc_k=0$, and the same identity proves
\eqref{eq:v144-pencil-matrix}. Diagonal tests cannot replace that
matrix inequality: even with $q=\nu I$ and the positive loss matrix
$R=\nu\left(\begin{smallmatrix}1&1\\1&1\end{smallmatrix}\right)$,
$\nu>0$, the two zero-error diagonal tests pass while $q-R$ has
negative determinant $-\nu^2$. This is an algebraic illustration,
not a counterexample for the arithmetic symbol.

The twelve-vector span loses at most one dimension under the source
constraint. In it,
\[
 \frac{q_a[Vc]}{q_a^+[Vc]}
   =\frac{\sum_{k=0}^{11}\nu_k|c_k|^2}
          {\sum_{k=0}^{11}|c_k|^2}<10^{-8},
\]
under the stated exact hypotheses. The last assertions separate the
quantifiers in the cited concentration and bounded-floor criteria;
no extension from a finite head to its complement is implicit.
\end{proof}

\paragraph{Diagnostic scope.}
Section~7 of \path{evidence/diag_true_symbol/results.md} reports
at $\lambda=4$, $N=48$ twelve pencil values below $10^{-8}$, with
$\nu_0\approx2.9\cdot10^{-65}$. These are diagnostic values:
$W^-$ uses floating quadrature and the generalized eigensolve uses
midpoints. They do not certify exact thresholds, a source-compressed
pencil, or any proposed minorant. They illustrate the relative
loss budgets that a zero-error minorant would face if the exact
values and admissibility hypotheses were established.
For the weaker bounded-error goal the allowance is instead
$\nu_k+\eta_a\norm{v_k}^2/q_a^+[v_k]$; tiny $\nu_k$ alone does not
force relative accuracy $\nu_k$, or refute that goal.

The finding is a finite-window diagnostic test class, not a replacement
for the cofinal analytic input. A concentration proof still needs an
explicit minorant and control on every admissible mixture and on the
complete complement, uniformly over the cofinal family (including the
other parity). No such family or new bound is supplied here. The
translation does not show that the weighted criterion fails, and
certainly does not produce a negative direction of $W_4$. The task
stops at this quantifier distinction. No new window, tail metric or
zero-enclosure refinement is attempted; G2 and RH remain open.

\subsection*{The lattice comb is the sharply truncated explicit formula}

\begin{proposition}[Exact zero expansion of the discrepancy symbol]
\label{prop:v146-beta-explicit}
Let $a>0$, $L=2a$, $x=e^L>1$, and let $\beta_a$ be exactly
\eqref{eq:v130-beta}, with the strict prime cutoff $n<x$. Write
$s=\tfrac12-i\xi$, $\xi\in\mathbb R$, and put
\[
 A(\xi)=\Re\psi(\tfrac54+\tfrac{i\xi}2)-\log\pi,
 \quad e_x=\begin{cases}\Lambda(x)/\sqrt{x},&x\in\mathbb N,\\0,&x\notin\mathbb N,\end{cases}
 \quad T_x(s)=\sum_{n=1}^\infty\frac{x^{-2n-s}}{2n+s}.
\]
Here $\Lambda(x)=0$ at integers which are not prime powers. Sum all
nontrivial zeros $\rho$, with multiplicity, by symmetric height limits;
no restriction $\Re\rho=1/2$ is imposed. Away from $\zeta(s)=0$,
\begin{align}
 \beta_a(\xi)={}&A(\xi)+2\Re\frac{\zeta'}\zeta(s)
       -2\Re\frac1{1-s}
       +2\Re\sum_\rho\frac{x^{\rho-s}}{\rho-s}
       -2\Re T_x(s)+e_x\cos(\xi L) \notag\\
 ={}&2\Re\sum_\rho\frac{x^{\rho-s}}{\rho-s}
       -2\Re T_x(s)+e_x\cos(\xi L).
 \label{eq:v146-beta-explicit}
\end{align}
The full real expression extends continuously at the excluded values.
At finite zero height $T$, add the Perron remainder $-2\Re R_{x,T}(s)$
to the right side; it is not permissible to discard this term merely
because a finite zero list has been smoothed.

More explicitly, for $c>1$, $B>0$ not an even integer, and $T>|\xi|$
with no boundary zero, set
\[
 F_s(z)=-\frac{\zeta'}\zeta(z)\frac{x^{z-s}}{z-s},\qquad
 E_{c,T}=\lim_{H\to\infty}\frac1{2\pi i}
 \left(\int_{c-iH}^{c-iT}+\int_{c+iT}^{c+iH}\right)F_s(z)\,dz.
\]
Let $\mathcal C_{B,T}$ run from $c+iT$ to $-B+iT$, down to
$-B-iT$, and back to $c-iT$. Then an exact remainder is
\begin{equation}\label{eq:v146-perron-remainder}
 R_{x,T}(s)=E_{c,T}-\frac1{2\pi i}\int_{\mathcal C_{B,T}}F_s(z)\,dz
       -\sum_{2n>B}\frac{x^{-2n-s}}{2n+s}.
\end{equation}
For the infinite expansion use the usual Perron symmetric-height
limit (or its Cesaro limit), in which this remainder tends to zero.
\end{proposition}
\begin{proof}
Put $S_x(s)=\sum_{n<x}\Lambda(n)n^{-s}$ and
$S_x^0(s)=S_x(s)+\tfrac12\Lambda(x)x^{-s}$, with the extra term zero
unless $x$ is an integer. Perron inversion applied to $-\zeta'/\zeta$
gives $S_x^0$. The residues of $F_s$ at $1,s,\rho,-2n$ are,
respectively,
\[
 \frac{x^{1-s}}{1-s},\quad-\frac{\zeta'}\zeta(s),\quad
 -\frac{x^{\rho-s}}{\rho-s},\quad\frac{x^{-2n-s}}{2n+s}.
\]
The rectangle with its right edge upwards gives exactly
\eqref{eq:v146-perron-remainder}. Letting the contours tend to infinity
in the standard Perron order gives the displayed zero expansion;
this is the classical explicit-formula argument, with the shifted
Perron denominator retained (compare \cite{HarperZetaNotes}).
Subtracting $(x^{1-s}-1)/(1-s)$ and taking $-2\Re$ gives the first
line of \eqref{eq:v146-beta-explicit}. In particular the strict
endpoint contributes $+e_x\cos(\xi L)$, not zero and not twice this.

The functional equation, not RH, gives on $\Re s=1/2$
\[
 2\Re\frac{\zeta'}\zeta(s)=\log\pi-\Re\psi(s/2).
\]
This follows by logarithmically differentiating the completed zeta
functional equation and using conjugation; the real parts of
$1/s+1/(s-1)$ cancel. The digamma recurrence now gives
$A+2\Re(\zeta'/\zeta)(s)=2\Re(1/(1-s))$, proving the second line
(\cite{DLMFZetaReflection}). Continuity follows from the original
finite prime sum. At a critical-line zero, the real part of the
individual singular quotient is interpreted by this limit; one does
not evaluate an infinite imaginary pole numerically.
\end{proof}

\paragraph{Exact lattice identity; no exact constant-amplitude comb.}
For $\xi=2\pi u/L$ away from the excluded ordinates, define
\[
 H_x(\xi)=\sum_\rho\frac{x^{\rho-1/2}}{\rho-1/2+i\xi}
       -\sum_{n=1}^\infty\frac{x^{-2n-1/2}}{2n+1/2-i\xi},\quad
 C_x(\xi)=e_x+2\Re H_x(\xi),\quad S_x(\xi)=-2\Im H_x(\xi).
\]
Then the exact identity is
\begin{equation}\label{eq:v146-lattice}
 \beta_a(2\pi u/L)=C_x(2\pi u/L)\cos(2\pi u)
                         +S_x(2\pi u/L)\sin(2\pi u).
\end{equation}
Thus integers sample $C_x$, half-integers sample $-C_x$, and odd
quarter-integers sample $\pm S_x$. The phase supplies alternation;
equal crests and exactly opposite troughs would additionally require
$C_x$ to be constant at those sample points, and quarter-point zeros
would require $S_x=0$ there. Neither is an identity. Conjugation makes
$C_x$ even and $S_x$ odd, but supplies neither missing condition.
The zero-frequency amplitude has the exact arithmetic expression
\begin{equation}\label{eq:v146-amplitude}
 C_x(0)=\beta_a(0)=\psi(5/4)-\log\pi
        -2\sum_{n<x}\frac{\Lambda(n)}{\sqrt n}+4(\sqrt x-1).
\end{equation}
It is the weighted remainder plus the archimedean constant, not
$\psi(x)-x$, and it is not a single amplitude for every lattice cell.
The archimedean constant is about $-1.372183419$, correcting the
$-2.11$ printed in the earlier diagnostic discussion.

For any conjugate pair actually on the critical line, its contribution
to the zero term is exactly
\[
 2\left\{\frac{\sin((\gamma+\xi)L)}{\gamma+\xi}
        +\frac{\sin((\gamma-\xi)L)}{\gamma-\xi}\right\}.
\]
Replacing all zeros by such pairs would assume RH. In the unconditional
formula the factors satisfy $|x^{\rho-1/2}|=x^{\Re\rho-1/2}$; setting
all of them equal to one is precisely where RH would enter. Even for
critical-line pairs the denominators depend on $\xi$. The expression
$4\sum_{\gamma>0}\sin(\gamma L)/\gamma$ is only their zero-frequency
amplitude contribution, before endpoint and trivial-zero terms; a
constant-amplitude crest-minus-trough would be twice the amplitude.

\paragraph{Numerical illustration, not a certificate (NS-26).}
The archived windows have $x=9,16,36,64$, all integers. Their endpoint
amplitudes are $\log3/3,\log2/4,0,\log2/8$, respectively. The
noninteger case is covered above by $e_x=0$, but cannot be substituted
for these archived inputs. At $\lambda=3$ the decomposition of
$\beta_a(2\pi/L)-\beta_a(\pi/L)$ is, numerically,
\[
 \underbrace{-0.159765789563}_{\text{full zero contribution}}
 \;\underbrace{-0.003943883015}_{\text{trivial zeros}}
 \;+\underbrace{0.732408192445}_{\text{strict endpoint}}
 =0.568698519867.
\]
With the first $400$ positive zeros and height-Cesaro weight
$1-\gamma/680$, the first entry is $-0.157163432709$ and its remaining
correction is $-0.002602356854$. At zero frequency the full zero
contribution is $-0.077669995516$; index-Cesaro summation of $400$
zeros gives $-0.077720711653$. This accounts for the reported
``$-0.08$ versus $+0.57$'': endpoint convention and amplitude-versus-
difference normalization are substantial, while shifted denominators,
trivial zeros and finite-sum remainders account for the residual.
The archimedean block cancels exactly; it is not an unspecified fitted
correction.

The ten $u$ values in \path{evidence/diag_true_symbol/selfsim_output.txt}
were checked at each of its four windows, using Gaussian-regularized
Perron residues with the exact desmoothing correction back to $n<x$.
The derivation and replay are in
\path{evidence/diag_beta_lattice_identity/}. The saved file has only
three decimals: all $40$ roundings agree, and the recomputed legacy
formula and zero-side evaluation agree to better than $10^{-6}$
(maximum observed discrepancy below $4\cdot10^{-21}$). The numerical
illustration is not an interval certificate or a new window run.

This is the explicit formula restated for the actual symbol. It
supplies no new sign, uniform minorant or ground-energy bound. The
complete energy still involves the state's full Fourier distribution,
not a scalar amplitude or just lattice samples. The concentration
criterion, G2 and RH remain open.

\paragraph{Deep troughs and sufficient zero counts (NS-26 follow-up).}
The sharp sum in Proposition~\ref{prop:v146-beta-explicit} is the
symmetric-height Perron sum over all zeros with multiplicity, or its
Cesaro limit; a finite prefix keeps \eqref{eq:v146-perron-remainder}.
For the numerical check the convention is instead Gaussian Perron
weight $\exp(\varepsilon^2(\rho-s)^2/2)$ with
$\varepsilon=0.025$, together with the pole, trivial residues, left
contour and exact desmoothing back to $n<x$. This full zero series
is absolutely convergent without RH. A raw $400$-zero prefix is not
assigned the same accuracy.

At $192/256$-bit zero input precision and $45/65$ decimal digits, the
ten low-$u$ samples at each of $\lambda=3,4,6,8$ respectively need
no more than $83,83,80,80$ positive zeros and their conjugates in this
numerical convention. The five recorded deep troughs and two binding
centroids at each of $\lambda=4,8$ need no more than $530,609$.
These are empirical sufficient counts for the listed samples, not
minimal counts or interval tail certificates. The selection rule
bounds the absolute remaining paired contributions through a
$1200$-zero reference prefix by $2\cdot10^{-7}$ and then checks the
answer against the recomputed original symbol. The reference prefix
was compared with $1000$ zeros; its finite-region zero count was
also checked. No hypothesis about unseen zeros is used as RH.

\begin{center}
\small
\begin{tabular}{rrrrrr}
$\lambda$&$\xi$&$u$&$\beta_a(\xi)$&$N_+$&absolute residual\\\hline
4&$526.500000$&232.329&$-3.419437295$&431&$1.46\cdot10^{-7}$\\
4&$39.247000$&17.319&$-3.365250756$&99&$6.96\cdot10^{-9}$\\
4&$222.840000$&98.333&$-3.300770464$&215&$4.58\cdot10^{-8}$\\
4&$81.020000$&35.752&$-3.206301313$&124&$7.27\cdot10^{-9}$\\
4&$652.130000$&287.766&$-3.136706249$&530&$1.84\cdot10^{-8}$\\
4&$64.554323^{\dagger}$&28.486&$+4.608588674$&115&$4.64\cdot10^{-8}$\\
4&$46.512756^{\dagger}$&20.525&$-0.481323287$&103&$3.21\cdot10^{-9}$\\
8&$362.255000$&239.779&$-6.220830020$&311&$9.61\cdot10^{-9}$\\
8&$290.745000$&192.446&$-5.658707866$&261&$9.61\cdot10^{-8}$\\
8&$177.390000$&117.416&$-5.543126380$&182&$4.07\cdot10^{-8}$\\
8&$634.630000$&420.066&$-5.515281321$&517&$7.12\cdot10^{-9}$\\
8&$751.895000$&497.684&$-5.513353521$&609&$6.42\cdot10^{-8}$\\
8&$220.024046^{\dagger}$&145.635&$-2.259398834$&214&$6.36\cdot10^{-8}$\\
8&$213.278015^{\dagger}$&141.170&$+7.424038966$&207&$2.49\cdot10^{-8}$\\
\end{tabular}
\end{center}
Here $N_+$ counts positive zeros, with conjugates also included;
residuals use that row's prefix. Daggers mark binding-direction
centroids, not trough locations: two lie on positive symbol values.
The deep-point reference residuals are below $6\cdot10^{-27}$ and
$4\cdot10^{-43}$ for $\lambda=4,8$, respectively, at both precisions.
The archived grid depths agree within $7\cdot10^{-13}$.
All scripts, components, count rules and residuals are recorded in
\path{evidence/diag_beta_lattice_identity/verify_deep.py} and its
\path{deep_b192.json}, \path{deep_b256.json} outputs.

At $\lambda=8$, $\xi=751.895$ the sharp $400$-zero approximation,
including endpoint and trivial terms but omitting its Perron remainder,
still errs by about $5.567$; at $\lambda=4$, $\xi=652.130$ the error
is about $0.233$. Thus the deep troughs also obey the full identity,
but cannot be inferred from an uncorrected short prefix. Their density
is a level-set question for the whole phased zero sum, not an established
consequence of nearest-neighbour zero spacing. The other zeros, shifted
denominators, cutoff phase, and without RH the factors
$x^{\Re\rho-1/2}$ all remain. These finitely many checks prove no
cofinal density, width, or minorant estimate.

\subsection*{Step-minorant packets: an exact local bound and the missing cofinal input}

The loss in \eqref{eq:v131-weighted-concentration} is an operator
lower-floor loss, not the norm of the discarded multiplier
$\beta_a-m_a$. These quantities cannot be interchanged. The following
statements are proved implications with the hypotheses displayed;
the NS-22 floating head experiment motivates them but proves none of
the arithmetic hypotheses below.

\begin{proposition}[Fixed-window packet bounds for a step minorant]
\label{prop:v146-packet-loss}
Fix $a>0$ and the even source complement
$P_a\mathcal H_a^{\rm ev}$ of Proposition~\ref{prop:v131-concentration}.
Let $m=-D_a+\sum_{j=1}^Jw_j1_{G_j}\le\beta_a$ be a finite step
minorant with positive finite weights and symmetric measurable sets.
Define
\[
 T_m=P_a\mathcal F_a^*m\mathcal F_aP_a\big|_{P_a\mathcal H_a^{\rm ev}},
 \quad \eta(m)=\max\{0,-\inf\sigma(T_m)\},\quad
 c_{a,u}(E)=\norm{\mathcal C_{a,u}(E)}.
\]
Any admissible $\eta_a$ in \eqref{eq:v131-weighted-concentration}
must be at least $\eta(m)$. For a symmetric set $E$:
\begin{enumerate}
\item If $m\le-b$ on $E$ and $m\le B$ everywhere, with $b,B\ge0$, then
\begin{equation}\label{eq:v146-spill-bound}
 \eta(m)\ge\bigl[(b+B)c_{a,u}(E)-B\bigr]_+.
\end{equation}
In particular $\eta(m)\ge b c_{a,u}(E)$ if $m\le0$ globally.
If the local level is $-D_a+r$, then $b=D_a-r$: the proposed
$c_{a,u}(E)(D_a-r)$ bound needs this nonpositivity assumption, or
separate control of the positive contribution outside $E$. Here $r$
is the uplift from $-D_a$, not a trough's absolute depth.
\item Put $\ell=\beta_a-m\ge0$. If $\ell\ge\delta$ on $E$ and a
unit $g\in P_a\mathcal H_a^{\rm ev}\cap\mathcal D(q_a)$ satisfies
$\int_E|\mathcal F_ag|^2\ge c$ and $q_a[g]\le Q$, then
\begin{equation}\label{eq:v146-gap-packet}
 \eta(m)\ge[c\delta-Q]_+.
\end{equation}
Thus a quantization gap supplies a bound on discarded energy; the
original packet energy $Q$ must also be controlled to bound $\eta(m)$.
\end{enumerate}
\end{proposition}
\begin{proof}
For a unit source-admissible vector of $E$-mass $p$,
$\ip{T_mg}g\le-bp+B(1-p)$. Take the supremum of $p$ to obtain
\eqref{eq:v146-spill-bound}. For the second assertion the exact identity
is $-\ip{T_mg}g=\int\ell|\mathcal F_ag|^2-q_a[g]$.
Nonnegativity of $\ell$ on the complement of $E$ gives
\eqref{eq:v146-gap-packet}. No uncontrolled cross term is dropped.
\end{proof}

\begin{lemma}[Explicit cell concentration and the rank-one cost]
\label{lem:v146-packet-concentration}
Put $W=a\omega>h>0$ and
\[
 E_{\omega,h}=[\omega-h/a,\omega+h/a]
                  \cup[-\omega-h/a,-\omega+h/a],\qquad
 f_\omega(t)=\frac{\cos(\omega t)}{\sqrt{a(1+\operatorname{sinc}(2W))}}
                  \quad(|t|<a),
\]
where $\operatorname{sinc}t=\sin(t)/t$. Set
\[
 p_h=\frac1\pi\int_{-h}^h\operatorname{sinc}^2t\,dt,\qquad
 c_0(W,h)=\frac{\left[\sqrt{p_h}
              -\sqrt{2h/\pi}/(2W-h)\right]_+^2}{1+1/(2W)}.
\]
Then $\norm{f_\omega}=1$ and
$\int_{E_{\omega,h}}|\mathcal F_af_\omega|^2\ge c_0(W,h)$.
If $r=|\ip{f_\omega}{u_a}|<1$, the normalized projection
$g=P_af_\omega/\sqrt{1-r^2}$ gives the explicit lower bound
\begin{equation}\label{eq:v146-source-concentration}
 c_{a,u}(E_{\omega,h})\ge
   \frac{[\sqrt{c_0(W,h)}-r]_+^2}{1-r^2}.
\end{equation}
For $W\ge10\pi$, the uncompressed lower bounds are $7/10$ when
$h=\pi/2$ and $11/25$ when $h=\pi/4$. If also $r\le1/20$, the
source-compressed lower bounds are respectively $3/5$ and $3/8$.
The same uncompressed bounds hold for the normalized odd sine packet.
The repaired source is even, so the odd packet needs no source projection.
\end{lemma}
\begin{proof}
Direct integration gives the norm and, after setting
$t=a(\xi-\omega)$ in the positive band, its paired-band mass is
\[
 \frac1{\pi(1+\operatorname{sinc}(2W))}
 \int_{-h}^h\{\operatorname{sinc}t+\operatorname{sinc}(t+2W)\}^2dt.
\]
Use the reverse triangle inequality in $L^2(-h,h)$,
$|\operatorname{sinc}(t+2W)|\le(2W-h)^{-1}$, and
$|\operatorname{sinc}(2W)|\le(2W)^{-1}$. Projection costs at most
$r\norm{1_E\mathcal F_au_a}\le r$ in the Fourier norm on $E$,
which proves \eqref{eq:v146-source-concentration}.
For $h\le\pi/2$, $\sin t/t\ge1-t^2/6\ge0$ gives
\[
 p_h\ge\frac{2h}{\pi}\left(1-\frac{h^2}{9}+\frac{h^4}{180}\right).
\]
At $h=\pi/2,\pi/4$ this is respectively greater than $3/4,7/15$.
Substitution with $W\ge10\pi$ proves the stated rational lower bounds
(using $3.14<\pi<22/7$); substituting $r\le1/20$ gives the compressed
bounds. For sine packets the second sinc has a minus sign and the norm
uses $1-\operatorname{sinc}(2W)$; the same inequalities apply.
\end{proof}

The lattice spacing is $\pi/a$. Thus $h=\pi/2$ describes a full
lattice-cell width $\pi/a$, whereas an actual half-cell has
$h=\pi/4$ and width $\pi/(2a)$. This distinction matters. The even
trace formula, before source subtraction, gives
\begin{equation}\label{eq:v146-halfcell-upper}
 c_{a,u}(E_{\omega,\pi/4})
 \le\operatorname{tr}\mathcal F_a^*1_{E_{\omega,\pi/4}}\mathcal F_a
       \big|_{\mathcal H_a^{\rm ev}}
 \le\frac12+\frac1{39\pi}<0.509\qquad(W\ge10\pi).
\end{equation}
Indeed the general trace upper bound is $2h/\pi+h/(\pi(W-h))$.
The diagnostic values $0.6$--$0.8$ in Section~8 of
\path{evidence/diag_true_symbol/results.md} refer to the entire
negative set, not one isolated half-cell. They cannot serve as that
half-cell's lower bound.

The rank-one cost is not automatically zero at a fixed window. Along
a cofinal family, the actual source satisfies $J_au_a\to\phi$ in
$L^2(\mathbb R)$ with $\phi\in L^1$, as used in
Proposition~\ref{prop:v136-positive-comparison}. Hence, uniformly in
$W\ge10\pi$,
\[
 r\le\norm{J_au_a-\phi}_2
       +\frac{\norm\phi_1}{\sqrt{a(1-1/(20\pi))}}\longrightarrow0.
\]
If a bound on $q_a[f_\omega]$ is available, it can also be projected
without losing its cross term: with
$\epsilon_a=\norm{\mathsf W_{e^a}u_a}$,
\[
 q_a[g]\le\frac{q_a[f_\omega]+(2r+r^2)\epsilon_a}{1-r^2}.
\]
A concentration estimate alone does not supply the numerator's packet
energy bound.

\begin{proposition}[What a level-count obstruction would require]
\label{prop:v146-level-complexity}
Fix a window with $D=D_a>0$ and a source-admissible unit vector
$g\in\mathcal D(q_a)$, and suppose $q_a[g]\le Q$, $Q\ge0$.
Define its arithmetic level distribution by
\[
 \mu_g(B)=\int_{\{\xi:\beta_a(\xi)\in B\}}
                      |\mathcal F_ag(\xi)|^2\,d\xi.
\]
Suppose, as an additional hypothesis, that for every Borel
$B\subset[-D,0]$,
\begin{equation}\label{eq:v146-level-density}
 \mu_g(B)\ge\frac{c}{D}|B|,\qquad c>0.
\end{equation}
Then every step minorant in \eqref{eq:v131-step-minorant} having at
most $K$ distinct values obeys
\begin{equation}\label{eq:v146-level-lower}
 \eta(m)\ge\left[\frac{cD}{2K}-Q\right]_+.
\end{equation}
Consequently, along a cofinal family for which these hypotheses hold
with $c\ge c_*>0$ and $Q\le Q_*<\infty$, bounded $K$ forces
$\eta(m)\to\infty$ by Corollary~\ref{cor:v137-scalar-no-go}.
A bound $\eta(m)\le C$ forces
$K\ge c_*D_a/(2(C+Q_*))$ when $C+Q_*>0$.
These conclusions are conditional on the displayed level-distribution
and packet-energy inputs, not on RH.
\end{proposition}
\begin{proof}
The minorant is at least $-D$, and its finitely many values must include
$-D$ since $\inf\beta_a=-D$. Insert its negative levels and the endpoint
$0$ into $[-D,0]$. The resulting $n\le K$ gaps have lengths
$\Delta_j$ summing to $D$. On a gap $(b,b+\Delta_j)$ every admissible
minorant is at most $b$, irrespective of the geometry of its sets.
Thus \eqref{eq:v146-level-density} yields
\[
 \int(\beta_a-m)|\mathcal F_ag|^2
 \ge\frac{c}{D}\sum_{j=1}^n\int_0^{\Delta_j}t\,dt
 =\frac{c}{2D}\sum_j\Delta_j^2\ge\frac{cD}{2K}.
\]
Subtract $q_a[g]\le Q$ as in \eqref{eq:v146-gap-packet}. This proves
both the local inequality and its conditional cofinal consequences.
\end{proof}

Hypothesis~\eqref{eq:v146-level-density} has concrete geometric
content. For example, it follows with $c=b/M$ if, on one frequency
interval, $\beta_a$ is monotone through the full range $[-D,0]$,
$0<|\beta_a'|\le MaD$ almost everywhere, and
$|\mathcal F_ag|^2\ge ba$ there. This is the change-of-variables
formula, not an assertion about the arithmetic troughs. A count of
near-$D$ troughs does not supply this distribution over all levels,
and it does not supply $q_a[g]\le Q_*$. Also $K$ counts distinct values:
$J$ arbitrary weighted sets can encode as many as $2^J$ values. A
linear lower bound on $K$ is not a linear lower bound on $J$.

For the particular nonpositive ``$J$ deepest troughs'' construction,
a simpler exact bound is available. Let the true negative components
have depths $d_1\ge d_2\ge\cdots$, and put $m=-d_j$ on the first
$J$ components and $m=-d_{J+1}$ everywhere else. Then
\begin{equation}\label{eq:v146-lobe-floor}
 m\le-d_{J+1}\quad\Longrightarrow\quad\eta(m)\ge d_{J+1}.
\end{equation}
This holds on every source complement without a packet argument.
In particular $\eta(m)\le C$ requires patching every component deeper
than $C$. To infer divergence for fixed $J$ one still needs
$d_{J+1}(a)\to\infty$; to infer a linear component count one needs a
corresponding lower bound for $\#\{j:d_j(a)>C\}$. The unconditional
growth of $d_1=D_a$ alone proves neither.

\paragraph{A countermodel to an inference from depth and concentration alone.}
Fix a bounded symmetric band set $E$ of positive measure and let
$0<c=\norm{\mathcal C_{a,u}(E)}<1$. The strict upper bound follows
from compactness and Fourier uniqueness: a compactly supported
nonzero physical vector cannot also have Fourier support in $E$.
For any $d>0$, put $B=dc/(1-c)$ and
\[
 m_d=B-(B+d)1_E=-d+(B+d)1_{E^c}.
\]
This is a two-value, one-set minorant of itself, with depth $d$, yet
\[
 T_{m_d}=B I-(B+d)\mathcal C_{a,u}(E)\succeq0,
 \qquad \eta(m_d)=0.
\]
One may take a continuous even symbol above $m_d$, equal to $-d$ on
an interior subband, rising to $B$ before the boundary of $E$, and
tending to infinity outside $E$. Its negative depth is still $d$ and
its minorant still has zero loss. This is not the arithmetic symbol;
it proves that depth growth and a positive concentration fraction
alone cannot imply the proposed unrestricted level-count conclusion.
The positive spill in \eqref{eq:v146-spill-bound} is essential.

\paragraph{Finding (NS-27): the scoped proof does not close the route.}
The unconditional inputs are the local packet inequalities, the explicit
concentration and projection bounds, the lobe-floor inequality, and
$D_a\to\infty$ from v1.37. The missing arithmetic inputs are a
uniformly populated level distribution seen by a source-admissible
bounded-energy packet (or an alternative joint signed estimate), and,
for trough-by-trough patching, the stated multiplicity/depth growth.
v1.37 gives an infimum tending to minus infinity, not uniformly wide
near-deepest troughs, a density of distinct such troughs, or these
packet-energy estimates. Its conditional RH probe is explicitly not
a squared transform of a physical window vector. It cannot fill this
gap. No RH assumption enters the new local results or the conditional
level-count implication.

The even argument applies to the odd symbol using
\eqref{eq:v135-full-symbol} and the sine packet; the same missing
arithmetic hypotheses remain, with no even-source constraint in that
sector. The NS-22 finite-head observations (including their finite
frequency scan, approximate tail charge, and optimized rather than
exhaustive level choices) remain diagnostics. They do not prove a
lower bound for every minorant. Accordingly the bounded-complexity
route stays \emph{blocked}, not \emph{closed}, on the map.
Proposition~\ref{prop:v131-concentration} remains a valid sufficient
criterion. This is a gap in the proposed obstruction proof, not a
failure of the physical form. No new window, tail metric or zero-bound
refinement is made. G2 and RH remain open.

\subsection*{Level distribution is a value-distribution and packet input (NS-29)}

The depth is $D_a=\| (\beta_a)_-\|_\infty=-\inf\beta_a$ when
$\inf\beta_a<0$, not the signed infimum itself. Fix $a>0$, an integer
head size $N\ge1$, and $X=2\pi N/L$, where $L=2a$. The head cutoff
specifies the observation interval $[0,X]$; it does not truncate the
complete symbol, the Fourier transform, or the form. In any cofinal
assertion below $N=N(a)$ must be specified as part of the family.
In particular the global depth $D_a$ must not be replaced silently by
$D_{a,X}:=\max(0,-\min_{[0,X]}\beta_a)\le D_a$.

\begin{proposition}[Separation of level measure and packet coverage]
\label{prop:v147-level-separation}
Fix $D=D_a>0$ and $0<\theta\le1$. Put
\[
 F_{a,X}(t)=|\{\xi\in[0,X]:\beta_a(\xi)<-t\}|,
 \qquad E(t_1,t_2)=\{\xi\in[0,X]:-t_2<\beta_a(\xi)<-t_1\}.
\]
The following is an explicit sufficient level-measure hypothesis:
\begin{equation}\label{eq:v147-level-increments}
 F_{a,X}(t_1)-F_{a,X}(t_2)
 =|E(t_1,t_2)|\ge\kappa X\frac{t_2-t_1}{D}
 \quad(0\le t_1<t_2\le\theta D),\qquad\kappa>0.
\end{equation}
Suppose also that a source-admissible even unit vector
$g\in\mathcal D(q_a)$, a measurable $S\subset[0,X]$, and constants
$b,\alpha>0$, $Q\ge0$ satisfy
\begin{align}
 |\mathcal F_ag(\xi)|^2&\ge b/X\quad\hbox{a.e. on }S,
 &q_a[g]&\le Q, \notag\\
 |S\cap E(t_1,t_2)|&\ge\alpha|E(t_1,t_2)|
 &&(0\le t_1<t_2\le\theta D).
 \label{eq:v147-packet-coverage}
\end{align}
Then, for every Borel $B\subset[-\theta D,0]$,
\begin{equation}\label{eq:v147-weighted-density}
 \mu_g(B)\ge\frac{2b\alpha\kappa}{D}|B|.
\end{equation}
For $\theta=1$ this is precisely the hypothesis of
Proposition~\ref{prop:v146-level-complexity}, with $c=2b\alpha\kappa$.
For $0<\theta<1$ it is a weaker, sufficient replacement: every minorant
with at most $K$ distinct values satisfies
\begin{equation}\label{eq:v147-partial-level-loss}
 \eta(m)\ge\left[\frac{b\alpha\kappa\theta^2D}{K}-Q\right]_+.
\end{equation}
Thus uniform positive $b,\alpha,\kappa,\theta$ and bounded $Q$ on a
cofinal family imply the same linear-in-$D_a$ level-count obstruction.
These are additional hypotheses, not conclusions about the arithmetic
symbol. Taking $S=[0,X]$ permits $\alpha=1$. Merely requiring
$|S|\ge\alpha X$, or an upper bound on Fourier density, does not
replace \eqref{eq:v147-packet-coverage}.
\end{proposition}
\begin{proof}
The exact $\beta_a$ is nonconstant and real analytic on the real axis,
and tends to $+\infty$ at infinity. Its level sets in $[0,X]$ have
Lebesgue measure zero, so the equality in
\eqref{eq:v147-level-increments} is valid, including at the endpoints.
On each band, integration of the Fourier lower bound gives at least
$b\alpha\kappa(t_2-t_1)/D$. Evenness doubles this contribution on the
reflected interval; all omitted frequency mass is nonnegative.
Domination on intervals extends to Borel sets (equivalently, apply the
Radon--Nikodym derivative of the level pushforward). This proves
\eqref{eq:v147-weighted-density}.

For the last bound partition $[-\theta D,0]$ at the minorant's values
in its interior. There are at most $K$ gaps: at least one value is
$-D\le-\theta D$, since the minorant is at least $-D$, has finitely
many values, and is below a symbol with infimum $-D$. On each gap the
minorant is at most its lower endpoint; if that endpoint is the
inserted $-\theta D$, this remains an upper bound. For gap lengths
$\Delta_j$ the loss is therefore at least
\[
 \frac{b\alpha\kappa}{D}\sum_j\Delta_j^2
 \ge\frac{b\alpha\kappa\theta^2D}{K}.
\]
Subtract $q_a[g]\le Q$ using \eqref{eq:v146-gap-packet}.
Unconditional $D_a\to\infty$ proves the conditional cofinal assertion.
\end{proof}

A cumulative estimate such as
$F_{a,X}(t)\ge\kappa X(1-t/D)$ is not enough. For example, the
nonarithmetic affine symbol $\beta(\xi)=-D/2-D\xi/(2X)$ has
$F(t)=X$ for $0<t\le D/2$ and $F(t)=2X(1-t/D)$ for $D/2<t<D$.
It satisfies that cumulative estimate with $\kappa=1$, but has no
mass in any level interval contained in $(-D/2,0)$. This is a
counterexample to the measure-theoretic inference, not to the
arithmetic symbol. Subtracting two lower bounds is not a lower bound
for their difference. Likewise two sets each having positive measure
need not intersect in every level band.

There is also an exact local description. At almost every regular level
$y$, the unweighted and weighted densities on $[0,X]$ are
\begin{equation}\label{eq:v147-coarea}
 h_{a,X}(y)=\sum_{\substack{0<\xi<X\\\beta_a(\xi)=y}}
          \frac1{|\beta_a'(\xi)|},\qquad
 h_{g,X}(y)=\sum_{\substack{0<\xi<X\\\beta_a(\xi)=y}}
          \frac{|\mathcal F_ag(\xi)|^2}{|\beta_a'(\xi)|}.
\end{equation}
This follows by splitting at the finitely many critical points and
changing variables on the monotone pieces. The increment condition is
equivalent to $h_{a,X}(y)\ge\kappa X/D$ almost everywhere on the
specified level interval. It asks for enough crossings with controlled
slopes, not just zeros of $\zeta$ or one deepest trough. Necessarily
$D_{a,X}\ge\theta D_a$; the finite head need not see the global depth.

This separation is sufficient, not necessary for the original weighted
hypothesis. The single-interval geometry following
Proposition~\ref{prop:v146-level-complexity} could supply the weighted
level density on a frequency interval of length $O(1/a)$ with packet
density of order $a$, without a positive fraction of the whole $[0,X]$.
There is also no assertion uniform over arbitrary choices of $N$:
for fixed $a$ the negative set of $\beta_a$ is bounded, so
$F_{a,X}(0)/X\to0$ as $X\to\infty$. Choosing excessively large
heads can therefore defeat \eqref{eq:v147-level-increments} without
affecting the complete weighted condition. Prescribing the cofinal
range is part of ZLD, not a harmless choice made after the estimate.

\paragraph{Exact reduction to a zero statement.}
Use all nontrivial zeros with multiplicity and the symmetric-height
Perron convention of Proposition~\ref{prop:v146-beta-explicit}. Define
\begin{align}
 Z_a(\xi)&=2\Re\lim_{T\to\infty}
       \sum_{|\Im\rho|<T}\frac{x^{\rho-1/2+i\xi}}{\rho-1/2+i\xi},
 &x&=e^{2a}, \notag\\
 C_a(\xi)&=e_x\cos(2a\xi)
       -2\Re\sum_{n\ge1}\frac{x^{-2n-1/2+i\xi}}{2n+1/2-i\xi},
 &\mathcal Z_a(\xi)&=Z_a(\xi)+C_a(\xi)=\beta_a(\xi).
 \label{eq:v147-corrected-zero-field}
\end{align}
Real continuous limits are used at critical-line zeros. No zero-location
hypothesis enters this identity. In particular the endpoint is retained
for prime-power $x$. Uniformly in real $\xi$,
\begin{equation}\label{eq:v147-correction-size}
 |C_a(\xi)|\le\varepsilon_a:=
 e_x+\frac{4x^{-5/2}}{5(1-x^{-2})}\longrightarrow0.
\end{equation}
Indeed $|2n+1/2-i\xi|\ge2n+1/2\ge5/2$, and
$e_x\le(\log x)/\sqrt x$. Consequently
\[
 \{Z_a<-t-\varepsilon_a\}\subset\{\beta_a<-t\}
 \subset\{Z_a<-t+\varepsilon_a\}.
\]
For a band one must instead shrink both ends:
\[
 \{-t_2+\varepsilon_a<Z_a<-t_1-\varepsilon_a\}
 \subset E(t_1,t_2)
 \quad\hbox{within }[0,X].
\]
The inner band can be empty when $t_2-t_1\le2\varepsilon_a$.
Thus even a small correction cannot simply be discarded in a demand
about \emph{every} level interval. A finite zero prefix also needs the
Perron remainder $-2\Re R_{x,T}$; \eqref{eq:v147-correction-size} is
not an estimate for that remainder.

\begin{assumption}[Corrected zero level distribution, ZLD]
\label{ass:v147-zld}
Fix a cofinal set of windows, a prescribed head-size function $N(a)$,
and $0<\theta\le1$. There is a constant $\kappa>0$ independent of
$a$ such that, eventually, for every $0\le r_1<r_2\le\theta$,
\begin{equation}\label{eq:v147-zld}
 \frac1{X_a}\left|\left\{0\le\xi\le X_a:
  -r_2D_a<\mathcal Z_a(\xi)<-r_1D_a\right\}\right|
 \ge\kappa(r_2-r_1),\qquad X_a=2\pi N(a)/L.
\end{equation}
This is a named \emph{open input}, not an established conjecture from
the literature. The independent packet input is
\eqref{eq:v147-packet-coverage} with $b,\alpha$ bounded below and
$Q$ bounded above on that same family. Neither input is asserted here.
\end{assumption}
Substituting \eqref{eq:v147-corrected-zero-field} proves that ZLD is
exactly \eqref{eq:v147-level-increments}. No information is gained by
changing from the prime expression to the zero expression. Only ZLD
with $\theta=1$ addresses the full original interval $[-D_a,0]$;
fixed $\theta<1$ suffices for the modified bound
\eqref{eq:v147-partial-level-loss}.

\paragraph{What the zero literature supplies, and the missing reduction.}
The following assessment concerns the stated theorems, not a proof
that no future consequence of zero theory can yield ZLD.

Riemann--von Mangoldt gives
$N(T)=\frac{T}{2\pi}\log\frac{T}{2\pi e}+O(\log T)$; explicit
versions are given by Hasanalizade--Shen--Wong
\cite{HasanalizadeShenWong2022}. It counts ordinates, without a signed
value-distribution conclusion for \eqref{eq:v147-corrected-zero-field}.
At large height $T$, an interval of width $O(1/a)$ has main count
$O((\log T)/a)$, while differencing the displayed error leaves
$O(\log T)$. In particular this count estimate alone does not isolate
or control cancellation at every window-scale side lobe.

Zero-density theorems bound the number of zeros off a specified
vertical line. For example Ingham's bound is
$N(\sigma,T)\ll T^{3(1-\sigma)/(2-\sigma)}\log^5T$; see
\cite{Ingham1940,ChourasiyaSimonic2025}. This is an upper count,
not a lower signed band measure. In the present sum an off-line zero
also has the factor $x^{\Re\rho-1/2}$, which may be large. Replacing
all those factors by one uses RH. Removing this issue by assuming RH
still leaves phase cancellation and the normalization by the extreme
value $D_a$; it does not supply a proof of ZLD.

Landau--Gonek estimates control the linear sums
$\sum_{0<\Im\rho\le T}x^\rho$, including an arithmetic main term
at prime powers; see \cite[Introduction, (1.1)]{Aryan2019}.
They do not state a signed distribution theorem for a moving resolvent.
Even on RH, with $A_x(v)=\sum_{0<\gamma\le v}x^{i\gamma}$ and
$\xi<U<V$, partial summation reads
\[
 \sum_{U<\gamma\le V}\frac{x^{i\gamma}}{\gamma-\xi}
 =\frac{A_x(V)}{V-\xi}-\frac{A_x(U)}{U-\xi}
       +\int_U^V\frac{A_x(v)}{(v-\xi)^2}\,dv.
\]
An upper error bound here loses control as the interval approaches
$\xi$ and gives no lower signed quantile. The real, paired expression
has removable singularities; the problem is the loss of its cancellation
in such absolute estimates, not a pole in $\beta_a$. Controlling a
first or second moment likewise does not give a lower density on
\emph{every} interval of values, particularly relative to a supremum.

Montgomery's theorem and pair-correlation conjecture concern averages
of differences of ordinates \cite{Montgomery1973}; the unconditional
complex-zero extension also remains a correlation average
\cite[Theorem 1]{BaluyotEtAl2023}. ZLD instead concerns values of a
phase-sensitive complete sum, with $x=e^{2a}$ varying, on a prescribed
frequency interval and normalized by its global negative depth. No
reduction from these correlation statements to ZLD, or conversely, is
proved here or supplied by the cited results. The same applies to the
classical density hypothesis
$N(\sigma,T)\ll_\epsilon T^{2(1-\sigma)+\epsilon}$ (and to
Lindel\"of's hypothesis): these control counts or upper growth, not the
required signed level distribution. ZLD is therefore a separate input
with \emph{unclassified logical strength}; it is not justified to call
it weaker than RH, equivalent to RH or pair correlation, or logically
independent of them. Its conjunction with the packet input obstructs
a minorant method; it is not a positivity assertion.

\begin{proposition}[What depth growth actually says about measure]
\label{prop:v147-depth-measure}
Unconditionally, at a fixed window define
\[
 M_a=\tfrac12\psi'(5/4)
       +2\sum_{1<n<x}\frac{\Lambda(n)\log n}{\sqrt n}
       +2\int_0^L t e^{t/2}\,dt.
\]
Then $|\beta_a'|\le M_a$. If $d=D_{a,X}>t\ge0$, then
\begin{equation}\label{eq:v147-depth-width}
 F_{a,X}(t)\ge\min\{X,(d-t)/M_a\}.
\end{equation}
If $\beta_a([0,X])$ contains $[-\theta D_a,0]$, then
\eqref{eq:v147-level-increments} holds at this window with
$\kappa=D_a/(M_aX)$. No uniform positive lower bound for this
$\kappa$ follows from $D_a\to\infty$.

Under the additional RH hypothesis of
Proposition~\ref{prop:v137-conditional-mass}, let
$w_a(\xi)=w(2a(\xi-\gamma_0))$. Eventually, for $0\le t<D_a$,
\begin{equation}\label{eq:v147-probe-measure}
 |\{\beta_a<-t\}\cap\operatorname{supp}w_a|
 \ge\frac{[1/(6\pi)-\pi t/a]_+}{D_a-t}.
\end{equation}
This is a conditional, possibly zero, cumulative measure bound. It
is not ZLD or a physical packet estimate.
\end{proposition}
\begin{proof}
Differentiate the finite prime sum and the continuum integral in
\eqref{eq:v130-r-symbol}. The trigamma series implies
$|\psi'(5/4+i\xi/2)|\le\psi'(5/4)$, proving the derivative bound.
From a minimizer on $[0,X]$, the Lipschitz estimate supplies a
one-sided interval of length at least the right side of
\eqref{eq:v147-depth-width}; strict endpoints have zero measure.
If the range contains the indicated levels, almost every such level
has a regular preimage. Equation~\eqref{eq:v147-coarea} then gives
$h_{a,X}\ge1/M_a$. This proves the fixed-window increment estimate.

For the RH statement, $0\le w_a\le1$, $\int w_a=\pi/a$ and,
eventually, $J_a=\int\beta_aw_a\le-1/(6\pi)$. Splitting the
integral at $\beta_a=-t$ gives
\[
 J_a\ge-\frac{\pi t}{a}
        -(D_a-t)\int_{\{\beta_a<-t\}}w_a.
\]
Rearranging and using $w_a\le1$ proves the claim.
\end{proof}
At $t=\theta D_a$, the numerator in
\eqref{eq:v147-probe-measure} need not be positive: the argument has
no upper control on $D_a/a$. Even a positive value is not a bound on
all band differences normalized by $X_a$. The unconditional v1.37
argument is a contradiction: a bounded cofinal scalar floor would
imply RH and then contradict the conditional probe. It proves global
depth divergence, not a uniform sublevel measure, a head-range depth,
or the Fourier mass of a realizable packet. Continuity supplies local
width only with its window-dependent slope cost, as made explicit above.

\begin{proposition}[Broad Fourier coverage does not bound the energy]
\label{prop:v147-broad-packet}
Fix a nonnegative, even $h\in C_c^\infty(-1/2,1/2)$ of $L^2$ norm
one. For $X\to\infty$ and $a>1/(2X)$ put
$f_X(t)=\sqrt X h(Xt)$, a unit even physical window vector. For a
constant $b_h>0$ depending only on $h$,
\begin{equation}\label{eq:v147-broad-packet}
 |\mathcal F_af_X(\xi)|^2\ge b_h/X\quad(0\le\xi\le X),
 \qquad q_a[f_X]=\log X+C_h+o(1).
\end{equation}
The energy asymptotic is uniform in these windows once
$1/X<\log2$. It gives no uniform $Q$ for this construction.
These vectors are not asserted source-admissible.
\end{proposition}
\begin{proof}
The unitary transform is $X^{-1/2}\widehat h(\xi/X)$. On $|v|\le1$,
$\widehat h(v)\ge(2\pi)^{-1/2}\cos(1/2)\int h>0$, proving the
lower bound. The autocorrelation of $f_X$ is supported in
$[-1/X,1/X]$. Every prime lag $\log n\ge\log2$ therefore contributes
zero exactly; the continuum term contributes $O(1/X)$, since the
autocorrelation has absolute value at most one. Finally
\[
 \int_{\mathbb R}A(Xv)|\widehat h(v)|^2\,dv
 =\log X-\log(2\pi)
       +\int_{\mathbb R}\log|v|\,|\widehat h(v)|^2\,dv+o(1).
\]
This follows from the digamma asymptotic; split $|v|<1/X$, whose
logarithmic contribution is $O((\log X)/X)$, and use the rapid decay
of $\widehat h$ on the remainder. The displayed integral defining
$C_h$ is finite. No prime cancellation or RH is used.
\end{proof}

Source projection has a pointwise cost as well as an energy cost. For
any unit even $f\in\mathcal D(q_a)$, with the actual source
$u_a\in\mathcal D(\mathsf W_{e^a})$, write $r=|\langle f,u_a\rangle|<1$ and
$g=P_af/\sqrt{1-r^2}$. Then
\begin{equation}\label{eq:v147-projected-density}
 |\mathcal F_ag(\xi)|\ge
 \frac{[|\mathcal F_af(\xi)|-r|\mathcal F_au_a(\xi)|]_+}
      {\sqrt{1-r^2}}.
\end{equation}
In particular a raw $b/X$ lower bound persists with $b/(4X)$ if
$r\|\mathcal F_au_a\|_\infty\le\sqrt{b/X}/2$. Small ordinary
source overlap alone is not this scale-dependent inequality.
The exact form expansion also gives
\[
 \left|(1-r^2)q_a[g]-q_a[f]\right|
 \le(2r+r^2)\|\mathsf W_{e^a}u_a\|.
\]
Thus for the broad pulses, if $r\to0$ and the source residual is
bounded, projection still leaves $q_a[g]/\log X\to1$.
It cannot repair their energy cost. The packet of a single trough in
Lemma~\ref{lem:v146-packet-concentration}, conversely, supplies local
Fourier mass but not coverage of every level across $[0,X]$ or a
uniform original energy. At one fixed window every nonzero physical
packet has an analytic transform and hence can be bounded below off
small neighborhoods of its finitely many zeros in $[0,X]$; the bound
and the energy may depend on the window, and those neighborhoods
can meet precisely the needed level bands. This does not give
\eqref{eq:v147-packet-coverage}.

\paragraph{Finding and scope.}
The outcome of NS-29 is a reduction to the named open input ZLD
\emph{together with} bounded-energy packet coverage. The existing
unconditional zero theorems assessed above do not complete that reduction
into a proved bound. Nor can one identify packet spread alone as the
remaining obstacle: the uniform arithmetic band measure (a) is itself
unproved, and even granting it leaves the joint coverage and energy
requirement (b). The modification from cumulative sublevel size to
\emph{increments}, and from an arbitrary positive fraction to
\emph{coverage of each level band}, is essential.

The new fixed-window implications require no RH assumption. RH is used
explicitly in \eqref{eq:v147-probe-measure} and when interpreting the
critical-line exponential sums in the literature discussion. The two cofinal
inputs remain open. The argument and factors of two apply to odd
packets as well because their Fourier modulus is even; the even-source
constraint is then absent. NS-22 and the concurrent NS-28 measurements
are diagnostics, not substitutes for these cofinal hypotheses. No new
window, head computation, tail metric or zero enclosure is made. This
is a gap in a proposed minorant obstruction, not failure of the true
weighted concentration criterion or of the physical object. The route
remains blocked, not closed. G2 and RH remain open.

\subsection*{The exact weighted cost, not a histogram peak (NS-32/33)}

The review of the v1.48 candidate (PR~14, revised head \texttt{f3611ef})
retains its weighted-cost implication and concentration-function approach,
but rejects its two displayed rearrangements for $K$ and its identification
of aligned-bin maxima with the full concentration function. The saved
$\lambda=4$ optimizer also fails its own bin constraint. The candidate
is not merged by this review; the audit is
\path{audits/2026-09-21-v1.48-pr14-v3.html}.
This subsection gives a self-contained corrected formulation and answers
the zero-statement question. None of the diagnostic optimizations is rerun.

Fix $a>0$, $D=D_a>0$, a prescribed integer head size $N\ge1$, and $X=2\pi N/L$,
$L=2a$, $x=e^L$. Let $\mathcal Z_a$ be the \emph{complete corrected zero field}
of \eqref{eq:v147-corrected-zero-field}; thus $\mathcal Z_a=\beta_a$,
with all nontrivial zeros, the strict-cutoff endpoint and the trivial
terms retained. Its symmetric-height Perron convention and continuous
real limits are unchanged. Define finite positive Borel measures on
$[-D,0]$ by
\begin{equation}\label{eq:v149-measures}
 \widetilde m_a(B)=\frac1X\int_0^X
               1_{\{\mathcal Z_a(\xi)\in B\}}\,d\xi,
 \qquad
 \mu_g(B)=\int_{\{\beta_a\in B\}}|\mathcal F_ag(\xi)|^2\,d\xi.
\end{equation}
For a unit vector both measures have total mass at most one. The negative
fraction and the exact density at almost every regular negative level are
\begin{align}
 \varphi_-(a)&=\frac1X\int_0^X1_{\{\mathcal Z_a(\xi)<0\}}\,d\xi,
 \notag\\
 \widetilde\rho_a(t)&=\frac1X
  \int_{\{\xi\in(0,X):\mathcal Z_a(\xi)=t\}}
       \frac{d\mathcal H^0(\xi)}{|\mathcal Z_a'(\xi)|},
 &R_a&=\mathop{\rm ess\,sup}_{-D<t<0}\widetilde\rho_a(t).
 \label{eq:v149-zero-density}
\end{align}
Here $\mathcal H^0$ counts preimages, so the integral is a sum, and
$R_a$ is allowed to be infinite. The measures have no atoms at levels
$-D$ or $0$. Equation~\eqref{eq:v147-coarea} proves absolute continuity,
not an $L^\infty$ bound for the density.

\begin{proposition}[Critical levels make the peak-density bound vacuous]
\label{prop:v149-critical-density}
Suppose $\xi_0\in(0,X)$ is a critical point with
$t_0=\beta_a(\xi_0)<0$. Then $R_a=\infty$. More precisely, if
\[
 \beta_a(\xi)=t_0+A(\xi-\xi_0)^r+O((\xi-\xi_0)^{r+1}),
 \qquad A\ne0,\quad r\ge2,
\]
one local inverse branch contributes
\begin{equation}\label{eq:v149-critical-asymptotic}
 \frac{1+o(1)}{Xr|A|^{1/r}}
               |t-t_0|^{1/r-1}
\end{equation}
to $\widetilde\rho_a(t)$ on the side reached by that branch.
In particular every complete negative component contained inside
$(0,X)$ supplies an interior minimum and makes $R_a$ infinite.
The same conclusion holds for a one-sided boundary critical point if
its interior branch has negative levels.
\end{proposition}
\begin{proof}
The actual symbol is nonconstant real analytic. Thus the first nonzero
Taylor derivative at a critical point has finite order $r\ge2$.
Inverting on either available monotone branch gives
$|\xi-\xi_0|=(|t-t_0|/|A|)^{1/r}(1+o(1))$ and hence
\eqref{eq:v149-critical-asymptotic}. The exponent is strictly between
$-1$ and $0$: the singularity is integrable but unbounded on sets of
positive level measure. Other preimages add nonnegative terms, so
there is no cancellation in the density. A complete negative component
has zero boundary values and a negative interior minimum, proving the
last assertion, with the boundary case treated one-sidedly.
\end{proof}

Thus there need not be a finite ``most common level'' near zero. The
largest \emph{bin average} in the NS-28 histogram is a different,
resolution-dependent statistic. Near the critical value above its
average on an adjacent bin of width $h$ grows like $h^{1/r-1}$ as
$h\downarrow0$. A finite bin peak cannot upper-bound $R_a$. The theorem
is conditional only on the stated geometric critical point, not on RH;
it does not certify the location of any diagnostic trough.

A finite density upper bound would itself be an all-band statement:
\[
 \widetilde m_a(B)\le R_a|B|\quad\hbox{for every Borel }B\subset[-D,0].
\]
To make $\varphi_-^2/R_a\to\infty$ using a finite upper estimate $r_a$
would suffice to prove, eventually,
\begin{equation}\label{eq:v149-slope-demand}
 \frac1{X_a}\sum_{\mathcal Z_a(\xi)=t}\frac1{|\mathcal Z_a'(\xi)|}
 \le r_a=o(\varphi_-(a)^2)
 \quad\hbox{for almost every }-D_a<t<0.
\end{equation}
A bound on the number of crossings by $n_a$, together with
$|\mathcal Z_a'|\ge s_a>0$ at all their negative regular preimages,
would give $r_a=n_a/(X_as_a)$. A lower slope bound at one level or
one preimage is insufficient; multiplicity of crossings matters too.
Near a negative critical point the displayed uniform slope mechanism
fails. The derivative means the derivative of the \emph{whole} corrected
field, not an unjustified termwise derivative of a conditionally
convergent zero prefix. It can equivalently be evaluated from the exact
finite-prime formula:
\[
 \beta_a'(\xi)=-\tfrac12\Im\psi'(5/4+i\xi/2)
   +2\sum_{1<n<x}\frac{\Lambda(n)\log n}{\sqrt n}\sin(\xi\log n)
   -2\int_0^L t e^{t/2}\sin(\xi t)\,dt.
\]
No replacement $|x^{\rho-1/2}|=1$ is used; that replacement for every
zero would require RH.

\paragraph{The arithmetic of the scalar and its relation to ZLD.}
For $0<\varphi_-\le1$ and $0<R_a<\infty$,
\begin{equation}\label{eq:v149-scalar-arithmetic}
 R_a\ge\frac{\varphi_-}{D_a},\qquad
 0\le\frac{\varphi_-^2}{R_a}\le\varphi_-D_a\le D_a,
 \qquad
 \frac{\varphi_-^2}{R_a}\to\infty
 \ \Longleftrightarrow\ \frac{R_a}{\varphi_-^2}\to0.
\end{equation}
In particular divergence requires $R_a\to0$. A positive lower bound
on $R_a$ prevents it; an upper bound on $R_a$ does \emph{not} prevent
it. Uniform level density has $R_a=\varphi_-/D_a$, and gives the
scalar $\varphi_-D_a$, which can diverge even with $\varphi_-\le1$.
With $R_a=\infty$ the packing lower bound is simply zero. Thus bounded
negative fraction alone is not the reason for rejecting the scalar;
its actual unbounded peak at critical levels is the sharper issue.

ZLD is a lower bound on all level bands; a finite $R_a$ is an upper
bound on all bands. Neither replaces the other as a matter of measure
theory. For example, with $U_I$ the uniform probability measure on
$I$, the measure
$\tfrac14U_{[-D,0]}+\tfrac14U_{[-1,0]}$ satisfies full-range ZLD with
$\kappa=1/4$, but has $\varphi_-=1/2$, $R\to1/4$ and bounded scalar.
Conversely, for fixed $0<\theta\le1$ and $p=D^{-2}$, the measure
\[
 \tfrac12\bigl((1-p)U_{[-\theta D/2,0]}+pU_{[-D,0]}\bigr)
\]
has scalar asymptotic to $\theta D/4$, but no uniform ZLD constant
on $[-\theta D,0]$: the lower part has density $p/(2D)$.
These are full-support measure countermodels to the proposed inference,
not distributions asserted for zeta. No implication, equivalence or
logical independence for the actual zero field is proved by them.

\begin{proposition}[The surviving weighted-cost implication]
\label{prop:v149-weighted-cost}
Let $\mathcal L_K$ be the finite subsets of $[-D,0]$ containing $-D$
and having at most $K\ge1$ elements. For $\mathcal L\in\mathcal L_K$
put $p_{\mathcal L}(t)=\max\{\ell\in\mathcal L:\ell\le t\}$ and define
\begin{equation}\label{eq:v149-qc}
 \mathrm{QC}_K(\widetilde m_a)
 =\inf_{\mathcal L\in\mathcal L_K}
       \int_{[-D,0]}(t-p_{\mathcal L}(t))\,d\widetilde m_a(t).
\end{equation}
This at-most-$K$ formulation avoids any need to assume strict spacing
or attainment of a minimum. Suppose a unit even source-admissible
$g\in\mathcal D(q_a)$ satisfies $q_a[g]\le Q$, $Q\ge0$, and
\begin{equation}\label{eq:v149-wlh}
 \mu_g(B)\ge c'\widetilde m_a(B)
 \quad\hbox{for every Borel }B\subset[-D,0],\qquad c'>0.
\end{equation}
Then every step minorant from \eqref{eq:v131-step-minorant} with at most
$K$ distinct values satisfies
\begin{equation}\label{eq:v149-weighted-cost-bound}
 \eta(m)\ge[c'\mathrm{QC}_K(\widetilde m_a)-Q]_+.
\end{equation}
No bounded density of $\widetilde m_a$ is needed. If it additionally has
a density bounded by a finite positive $R$, then
$\mathrm{QC}_K\ge\varphi_-^2/(2KR)$. For uniform density on $[-D,0]$,
$\mathrm{QC}_K=\varphi_-D/(2K)$, recovering v1.46 with $c=c'\varphi_-$.
\end{proposition}
\begin{proof}
The negative values of $m$ include $-D$ and form a set
$\mathcal L\in\mathcal L_K$. At a negative symbol value $t$ one has
$m\le p_{\mathcal L}(t)$; all discarded loss is nonnegative. Since
$\mu_g-c'\widetilde m_a$ is a finite positive measure, its integral
against the nonnegative Borel function $t-p_{\mathcal L}(t)$ is
nonnegative, by approximation by simple functions. Hence the total
packet loss is at least $c'\mathrm{QC}_K$. Subtract $q_a[g]\le Q$ as
in \eqref{eq:v146-gap-packet}. The infimum is nonincreasing in $K$
because $\mathcal L_K\subset\mathcal L_{K+1}$, so fewer levels cause
no gap in this argument.

Under the additional density bound, write $m_j$ for the mass in each
gap. Packing it nearest the gap's left endpoint at density at most
$R$ costs at least $m_j^2/(2R)$. Summing at most $K$ gaps and using
Cauchy--Schwarz gives $\varphi_-^2/(2KR)$. For uniform density, equal
gaps attain this value, proving the constant exactly.
\end{proof}

\paragraph{A finite alternative to the peak.}
For $\varphi_->0$ and $0<h\le D$, define the \emph{full} concentration
function $M(h)=\sup_{I\subset[-D,0],\,|I|=h}\widetilde m_a(I)>0$.
Then, without a bounded density,
\begin{equation}\label{eq:v149-concentration-cost}
 \mathrm{QC}_K\ge\frac{h\varphi_-^2}{2KM(h)}-\frac{h\varphi_-}{2}.
\end{equation}
Indeed, partition each level gap into consecutive blocks of length $h$.
Each block (including a final shorter one) has mass at most $M=M(h)$.
For gap mass $m=M(n+\theta)$, $n\ge0$ integer and $0\le\theta<1$, filling
blocks from the left minimizes the block-index cost and gives
\[
 \int_{\rm gap}(t-\ell)\,d\widetilde m_a
 \ge hM\left(\frac{n(n-1)}2+n\theta\right)
 \ge h\left(\frac{m^2}{2M}-\frac m2\right).
\]
The last inequality uses $\theta(1-\theta)\ge0$. Sum the gaps and apply
Cauchy--Schwarz to prove \eqref{eq:v149-concentration-cost}.
With uniform packet bounds $c'\ge c_0>0$, $Q\le Q_*$ and loss $\eta\le C$,
the correct rearrangement is
\begin{equation}\label{eq:v149-concentration-rearrangement}
 K\ge\frac{c_0h\varphi_-^2}
 {2M(h)(C+Q_*+c_0h\varphi_-/2)}.
\end{equation}
If $D$ is split into aligned bins of width $h=D/n$ and $b$ is their
largest \emph{exact} mass, then $b\le M(h)\le2b$, since any interval
of length $h$ meets at most two bins, up to massless endpoints.
The archived midpoint bin maximum is neither the full $M(h)$ nor a
certified upper bound for it. Replacing a denominator by that lower
estimate cannot certify a lower cost bound.

\paragraph{The one-level cost is an exact zero functional.}
Put $\mathcal Z_a^-=\max(-\mathcal Z_a,0)$, and keep the global
$D_a$ even when the head range does not attain it. Then
\begin{align}
 A_a:=\mathrm{QC}_1(\widetilde m_a)
 &=\frac1{X_a}\int_{\{0\le\xi\le X_a:\mathcal Z_a(\xi)<0\}}
                         (D_a+\mathcal Z_a(\xi))\,d\xi \notag\\
 &=\varphi_-(a)D_a-\frac1{X_a}\int_0^{X_a}\mathcal Z_a^-(\xi)\,d\xi
 =\int_0^{D_a}\left(\varphi_-(a)-\frac{F_{a,X_a}(t)}{X_a}\right)dt.
 \label{eq:v149-qc1-zero}
\end{align}
This follows directly from $p_{\{-D\}}=-D$ and the layer-cake identity
for the negative part. Thus $A_a\to\infty$ asks that the total
frequency-averaged distance of negative values from the deepest
value diverge. For example it suffices that
$\varphi_-D_a\to\infty$ and
$X_a^{-1}\int\mathcal Z_a^-\le(1-\delta)\varphi_-D_a$ for some fixed
$\delta>0$. It is a statement about values of the corrected zero sum,
not about the number of zeros or a prescribed zero spacing.

There is no conclusion here that $\varphi_-$ tends to a positive limit,
or to zero. Corollary~\ref{cor:v137-scalar-no-go} proves global depth
\emph{only}; it does not prove nested negative sets, increasing measure,
or any fraction for the chosen range. The four observed fractions do
not establish an asymptotic. The choice of $N(a)$ matters: at each fixed
window the negative set is bounded, so choosing an excessively large
prescribed head range can make its fraction arbitrarily small. No
uniform assertion over all choices of $N(a)$ is possible. These points
are independent of RH and do not themselves disprove a positive-fraction
statement for a suitably specified family.

\begin{proposition}[One-level growth does not imply fixed-level growth]
\label{prop:v149-qc1-insufficient}
There is no universal positive constant $\kappa$ such that
$\mathrm{QC}_K(\nu)\ge\kappa\mathrm{QC}_1(\nu)/K$ for all finite
absolutely continuous level measures $\nu$. In fact there are measures
on $[-D,0]$ of mass $1/2$, with support the whole interval, such that
$\mathrm{QC}_1\sim D/2$ but $\mathrm{QC}_2\to0$ as $D\to\infty$.
\end{proposition}
\begin{proof}
Take $D>3$, $h=D^{-1}$, $p=D^{-2}$ and $r=D^{-3}$, and put
\[
 \nu_D=\tfrac12\bigl((1-p-r)U_{[-1-h,-1]}
             +pU_{[-D,-D+h]}+rU_{[-D,0]}\bigr).
\]
The one-level cost is asymptotic to $D/2$, since almost all the mass
sits near $-1$ while the compulsory level is $-D$. The admissible two
levels $\{-D,-1-h\}$ give cost at most $h/4+rD/2$, hence
\[
 \mathrm{QC}_2(\nu_D)\le\frac1{4D}+\frac1{2D^2}\longrightarrow0.
\]
The bridge of mass $r/2$ makes the support all of $[-D,0]$. These are
ordinary absolutely continuous measures. They are not claimed to be
arithmetic level distributions; they disprove the proposed general
quantization inference.
\end{proof}

\begin{assumption}[Exact level quantization growth, QG]
\label{ass:v149-qg}
Fix a cofinal family and a prescribed rule $N(a)$, hence $X_a$.
For every fixed integer $K\ge1$, require
\begin{equation}\label{eq:v149-qg}
 \mathrm{QC}_K(\widetilde m_a)
 =\inf_{\mathcal L\in\mathcal L_K}\frac1{X_a}
   \int_{\{\mathcal Z_a<0\}\cap[0,X_a]}
       (\mathcal Z_a(\xi)-p_{\mathcal L}(\mathcal Z_a(\xi)))\,d\xi
 \longrightarrow\infty.
\end{equation}
Separately require source-admissible unit packets $g_a$ satisfying
\eqref{eq:v149-wlh} with $c'\ge c_0>0$ and $q_a[g_a]\le Q_*<\infty$, $Q_*\ge0$.
Both are named open inputs. Neither follows here from the four diagnostic
windows or from v1.37.
\end{assumption}
Under these hypotheses, a minorant family with bounded level count
$K(a)\le K_0$ has
$\eta(m_a)\ge[c_0\mathrm{QC}_{K_0}(\widetilde m_a)-Q_*]_+\to\infty$.
Thus a uniformly bounded loss forces $K(a)\to\infty$ (apply the same
argument to any subsequence with bounded $K$). This is a sufficient
obstruction, not a necessary condition for all possible proofs of one.
QG says that the corrected zero field's negative values cannot be
approximated from below by any fixed number of levels with bounded
mean error. It needs no bounded peak density.

A stronger quantitative alternative is \emph{quantization shape control}:
for some $\kappa>0$ independent of $a,K$, require
\begin{equation}\label{eq:v149-shape}
 A_a\to\infty,\qquad
 \mathrm{QC}_K(\widetilde m_a)\ge\frac{\kappa A_a}{K}
 \quad\hbox{for all }K\ge1\hbox{ on the cofinal family}.
\end{equation}
With the packet input this gives
$K\ge c_0\kappa A_a/(C+Q_*)$ for loss $\eta\le C$ and $C+Q_*>0$.
If $C+Q_*=0$, positive exact cost instead gives incompatibility.
Linear growth in $D_a$ additionally needs $A_a\ge d_0D_a$ for a fixed
$d_0>0$. Equation~\eqref{eq:v149-qc1-zero} by itself gives none of the
shape inequality: concentration in a few narrow level bands is precisely
the missing obstruction. The empirical relation
$\mathrm{QC}_K\approx\mathrm{QC}_1/(1.8K)$ is not that theorem.

\paragraph{Diagnostic and source constraints after review.}
The archived NS-28 code constrains 20 level bins for raw $N=256$ even
vectors, scans to $\xi=800$ independently of $2\pi N/L$, uses a sampled
depth, and computes a 200-bin surrogate cost with reconstruction levels
on bin edges. The reported finite peak is a peak of bin averages, and
only $K\le32$ costs are evaluated. At $\lambda=4,c'=1$, the saved
returned vector has minimum bin ratio $0.636<1$: its energy is not a
feasible upper bound even for the binned optimization. The dual number
alone supplies no such upper bound. Claims of exact WLH, cofinal
$Q\to0$, and measured costs at hundreds or thousands of levels are
therefore withdrawn as interpretations of those diagnostics; their
raw files are preserved.

For a raw unit vector $f$, projection to
$g=P_af/\sqrt{1-r^2}$ with $r=|\langle f,u_a\rangle|<1$ obeys
\eqref{eq:v147-projected-density}. If $E=q_a[f]$ and
$\delta=(2r+r^2)\|\mathsf W_{e^a}u_a\|$, the form expansion gives
\[
 \frac{E-\delta}{1-r^2}\le q_a[g]\le\frac{E+\delta}{1-r^2}.
\]
For $E>0$, the explicit sufficient checks $r^2\le1/4$ and
$\delta\le E/2$ would keep this within $[E/2,2E]$. No such overlap
check is archived for the LP minimizer; a small source residual alone
does not control the denominator or preserve the weighted level
constraints. The missing $\lambda=3$ computation is to recover that
minimizer in the centered source basis, enclose its overlap and source
tail, its complete energy and the source cross term, and recheck the
projected constraints. It is named, not run. Passing the finite bin
constraints would still not prove the all-Borel hypothesis.

\paragraph{Finding (NS-33).}
The peak-density scalar is the wrong general bound for an exact level
measure containing a negative critical point: it is vacuous there.
The valid alternative is the exact cost, reduced to the named open
input QG plus the admissible bounded-energy weighted packet input.
Replacing it by $\mathrm{QC}_1/K$ requires the additional shape
hypothesis \eqref{eq:v149-shape}; divergence of
$\varphi_-D_a-X_a^{-1}\int\beta_a^-$ alone is insufficient. No cofinal
input is proved unconditionally. All new local identities, implications
and countermodels require no RH. Odd packets satisfy the same measure
argument because their Fourier modulus is even, without the even-source
constraint. This is a limitation of the proposed bound and diagnostic
interpretation, not a failure of the physical form or the valid weighted
concentration criterion. No new window, tail metric, optimization or zero
enclosure is computed. The route remains blocked. G2 and RH remain open.


\subsection*{A positive arithmetic weight and the critical capacity budget (NS-34)}

The bounded-error target of Proposition~\ref{prop:v136-bounded-floor}
allows a fixed ordinary-norm error, but it does not automatically allow
relative slack in a ground-state transform. Here one explicit weight is
chosen from the archived Gaussian arithmetic radical. The resulting
identities are exact; the obstruction below concerns a weakened
\emph{sufficient bound}, not a negative direction of the Weil form.
No new window, tail metric or numerical experiment is used.

Put $\ell=\log2$, and use the unnormalized positive weight
\begin{equation}\label{eq:v150-weight}
 h(x)=e^{x/2}\sum_{n\ge1}H(ne^x),\qquad
 H(u)=\frac{4\pi}{\sqrt3}u^2(2\pi u^2-3)e^{-\pi u^2}.
\end{equation}
Thus $h(x)=f_\infty(e^x)$ from Lemma~\ref{lem:v115-source-mass}.
Its restriction to each $I_a=(-a,a)$ is the comparison weight; it is
\emph{not} identified with the actual repaired source $u_a$. Positive
scalar rescaling leaves the forms unchanged when $g$ and $I_h$ are
rescaled consistently. All displayed constants below use the
unnormalized weight \eqref{eq:v150-weight}.

\begin{lemma}[Explicit positive weight and normalization]
\label{lem:v150-weight}
The function $h$ is smooth, even and strictly positive, and has a positive
minimum on every closed finite interval. Write $c(x)=\cosh(x/2)$ and
$s(x)=\sinh(x/2)$. Then
\begin{equation}\label{eq:v150-pole-mass}
 I_h:=\int_{\mathbb R}c(x)h(x)\,dx=\frac1{\sqrt3},
 \qquad d\nu(x)=I_h^{-1}c(x)h(x)\,dx
\end{equation}
defines an even probability measure. For $x\ge0$,
\begin{equation}\label{eq:v150-weight-bounds}
 A_-e^{9x/2-\pi e^{2x}}\le h(x)\le A_+e^{9x/2-\pi e^{2x}},
 \quad A_-=\frac{4\pi(2\pi-3)}{\sqrt3},\quad
 A_+=\frac{16\pi^2}{\sqrt3}.
\end{equation}
In particular, for $x\ge\ell$,
\begin{equation}\label{eq:v150-weight-ratio}
 \frac{h(x-\ell)}{h(x)}\ge
 \frac{2\pi-3}{4\pi}\,2^{-9/2}
             \exp\!\left(\frac{3\pi}{4}e^{2x}\right).
\end{equation}
\end{lemma}
\begin{proof}
Poisson summation and the two zero moments of the Fourier-invariant
$H$ give evenness, as in Lemma~\ref{lem:v115-source-mass}.
For $x\ge0$ every summand is positive. The $n=1$ term gives the lower
bound. For the upper bound use
\[
 h(x)\le\frac{8\pi^2}{\sqrt3}e^{9x/2-\pi e^{2x}}
 \left(1+\sum_{n\ge2}n^4e^{-\pi(n^2-1)e^{2x}}\right)
 \le A_+e^{9x/2-\pi e^{2x}}.
\]
Indeed $4\log n\le n^2-1$ for $n\ge2$, and $\pi>3$ bounds the sum
by $\sum_{n\ge2}e^{-2(n^2-1)}\le\sum_{n\ge2}e^{-2n}<1$.
Evenness gives the same decay at the other end. Smoothness and strict
positivity then give the fixed-interval minimum. Taking the ratio of
the two bounds proves \eqref{eq:v150-weight-ratio}.

For $\Re z>1$, absolute convergence and a Gaussian integral give
\[
 \int_0^\infty H(u)u^{z-1}\,du
 =\frac{z(z-1)}{\sqrt3}\pi^{-z/2}\Gamma(z/2),
\]
\begin{equation}\label{eq:v150-mellin}
 \int_{\mathbb R}h(x)e^{(z-1/2)x}\,dx
 =\frac{z(z-1)}{\sqrt3}\pi^{-z/2}\Gamma(z/2)\zeta(z).
\end{equation}
The decay just proved makes the left side entire in $z$; the identity
continues through the removable singularities on the right. Letting
$z\to1$ gives $\int e^{x/2}h=1/\sqrt3$, and evenness gives
\eqref{eq:v150-pole-mass}. This is the classical completed-zeta kernel,
not a new zero-free function or a use of RH.
\end{proof}

\begin{proposition}[Exact full and windowed energy identities]
\label{prop:v150-critical-energy}
For $f=hg\in C_c^\infty(\mathbb R)$ define
\begin{align}
 \mathcal J_h[g]={}&\frac12\iint_{\mathbb R^2}
       \rho(|x-y|)h(x)h(y)|g(x)-g(y)|^2\,dx\,dy\notag\\
 &+\sum_{n\ge2}\frac{\Lambda(n)}{\sqrt n}
       \int_{\mathbb R}h(x)h(x+\log n)
                    |g(x+\log n)-g(x)|^2\,dx,
 \label{eq:v150-global-energy}
\end{align}
where $\rho(t)=e^{t/2}/(2\sinh t)$ as in v1.29. All displayed
nonnegative terms have a finite total for such $f$. Then
\begin{equation}\label{eq:v150-critical-identity}
 q[f]=\mathcal J_h[g]-\frac23\operatorname{Var}_{\nu}(g)
                 -2\left|\int_{\mathbb R}s(x)h(x)g(x)\,dx\right|^2,
 \qquad
 \operatorname{Var}_{\nu}(g)=\int|g|^2d\nu-\left|\int g\,d\nu\right|^2.
\end{equation}
For even $g$ the last square vanishes. For odd $g$ the mean in the
variance vanishes but the last square must be retained.

For the finite window let $\mathcal J_{a,h}$ be
\eqref{eq:v150-global-energy} with both endpoints of every edge inside
$I_a$, equivalently \eqref{eq:v129-transformed-energy}. Put
\begin{align}
 T_a(x)={}&\int_{\mathbb R\setminus I_a}\rho(|x-y|)h(y)\,dy\notag\\
 &+\sum_{n\ge2}\frac{\Lambda(n)}{\sqrt n}
 \bigl[1_{I_a^c}(x+\log n)h(x+\log n)
       +1_{I_a^c}(x-\log n)h(x-\log n)\bigr],\label{eq:v150-exterior}\\
 r_a(x)={}&\frac{T_a(x)-2I_hc(x)}{h(x)},\qquad x\in I_a.
 \label{eq:v150-potential}
\end{align}
For $f$ supported in $I_a$, with $g=0$ outside, the exact identities are
\begin{align}
 \mathcal J_h[g]&=\mathcal J_{a,h}[g]
                +\int_{I_a}\frac{T_a}{h}|f|^2\,dx,\label{eq:v150-edge-split}\\
 q_a[f]&=\mathcal J_{a,h}[g]+\int_{I_a}r_a|f|^2\,dx+\Pi[f],
 \quad \Pi[f]=2|\langle f,c\rangle|^2-2|\langle f,s\rangle|^2.
 \label{eq:v150-local-identity}
\end{align}
They extend to the physical form domain on each fixed window.
The potential $r_a$ is exactly $H_a h/h$ from v1.29, with the pole
removed; no exterior term has been discarded.
\end{proposition}
\begin{proof}
The known Gaussian arithmetic radical has zero pairing with every
compact smooth Weil test (the radical input of v1.35, originating in
\cite[Section~3]{ConnesConsani2023}). Equivalently, its full pointwise
archimedean-minus-prime expression is $-2I_hc$. This last statement
also follows distributionally by taking the v1.35 smooth cutoffs of
$h$: their complete operator residual tends to zero, while their pole
moments and the absolutely convergent prime and archimedean expressions
converge against each fixed compact test. Smooth local convergence then
gives the pointwise identity. No global positive realization is assumed.

Apply the elementary Picone identity from
Proposition~\ref{prop:v129-picone} to all archimedean and prime edges.
The potential is $-2I_hc/h$ and the full pole is $\Pi[f]$, by
Proposition~\ref{prop:v135-both-parities}. Thus
\[
 q[f]=\mathcal J_h[g]-2I_h\int ch|g|^2
               +2\left|\int chg\right|^2-2\left|\int shg\right|^2,
\]
which gives \eqref{eq:v150-critical-identity} because $I_h^2=1/3$.
The Gaussian decay gives absolute convergence of the prime terms when
one endpoint is in a fixed compact set; the continuous kernel's local
singularity is integrable after the squared difference. These facts
justify the edge identity without subtracting divergent positive sums.

Separating edges crossing $I_a^c$ gives \eqref{eq:v150-edge-split}
with $T_a\ge0$, hence \eqref{eq:v150-local-identity}.
For interior $x$, every prime shift with $n\ge e^{2a}$ is exterior,
so this is precisely the strict-cutoff finite form. Endpoint equalities
are irrelevant on the open interval. At each fixed $a$, $h$ and $1/h$
are bounded and smooth, while $c/h$ is bounded. On the physical core,
$\mathcal J_h=q_a+V_a$ where
$V_a[f]=2I_h\int(c/h)|f|^2-\Pi[f]$ is a bounded $L^2$ form.
The shifted form norms are therefore equivalent, and completion of the
squared-edge map identifies its nonnegative integral with this closure.
No uniform-in-$a$ multiplier bound is asserted. The exterior energy is
a nonnegative part of $\mathcal J_h$, so its integral, and then those
of $(r_a)_\pm$, are finite on the resulting form domain. This argument makes no
claim for noncompact $f/h$ with infinite separate transformed energies.
\end{proof}

\paragraph{The exact capacity demand, including the source.}
For $0\le\delta_a\le1$ and $\eta_a\ge0$, the proposed strict absorption test is
\begin{equation}\label{eq:v150-capacity-demand}
 \int_{I_a}(r_a)_-|f|^2\le
 (1-\delta_a)\mathcal J_{a,h}[f/h]+\Pi[f]+\eta_a\|f\|^2.
\end{equation}
Use even $f\perp u_a$ and all odd $f$, with $u_a$ the actual normalized
even source. In the even sector the constraint is
$\int_{I_a}h g\overline{u_a}=0$, not $\int g\,d\nu=0$.
In the odd sector the negative pole can instead be placed on the left
as $2|\langle f,s\rangle|^2$. Uniform finite $\eta_a$ in both sectors,
together with the archived vanishing source residual, would suffice
by Corollary~\ref{cor:v136-complement-floor}.
The form actually lower-bounded by this test is
\begin{equation}\label{eq:v150-budget-form}
 \mathcal B_{a,\delta}[f]
 =q_a[f]-\int_{I_a}(r_a)_+|f|^2
                  -\delta_a\mathcal J_{a,h}[f/h].
\end{equation}
This identity isolates both losses of the sufficient criterion.

\begin{proposition}[No fixed relative slack for the Gaussian arithmetic weight]
\label{prop:v150-capacity-slack}
Set $b=\ell/16$ and
\[
 \kappa=\frac{(2\pi-3)\log2}{128\pi}>0,\qquad
 \gamma=\frac{3\pi}{4}e^{-4\ell-2b}>0.
\]
There are $A<\infty$ and $a_0$, independent of $a$, such that every
valid \eqref{eq:v150-capacity-demand} on either the even source
complement or the odd sector, for $a\ge a_0$, must satisfy
\begin{equation}\label{eq:v150-slack-budget}
 \eta_a\ge\left[\delta_a\kappa\exp(\gamma e^{2a})-B_a\right]_+,
 \qquad B_a=A+(8a+1)e^a+2a.
\end{equation}
In particular no fixed $\delta_a\ge\delta_0>0$ permits a uniform
finite $\eta_a$. If $\eta_a\le C$, necessarily
\begin{equation}\label{eq:v150-slack-rate}
 \delta_a\le\frac{C+B_a}{\kappa}e^{-\gamma e^{2a}}
       =O(\lambda\log\lambda\,e^{-\gamma\lambda^2}).
\end{equation}
Also, retaining the coefficient one on every other transformed edge
but deleting the prime-$2$ energy alone forces
$\eta_a\ge[\kappa\exp(\gamma e^{2a})-B_a]_+$.
These are failures of the indicated lower comparisons, not of $q_a$.
\end{proposition}
\begin{proof}
Choose two fixed real even orthonormal smooth bumps
$\phi_1,\phi_2\in C_c^\infty(-b,b)$. Put $t=a-2\ell$ and
\[
 F_{j,t}^{\pm}(x)=2^{-1/2}\bigl(\phi_j(x-t)\pm\phi_j(x+t)\bigr).
\]
For large $a$ these give orthonormal pairs in the even and odd sectors,
respectively. The even two-dimensional span contains a unit vector
$f_a^{\rm ev}\perp u_a$, regardless of the source coefficients. Every
odd vector is source-orthogonal. All estimates below hold for every
unit vector in either span, hence include both admissible choices.

The support bands and their shifts by $\ell$ are disjoint. For $x$ in
the positive support band the edge to $x-\ell$ stays inside $I_a$ and
has $f(x-\ell)=0$; for the negative band use the reflected edge. The
prime-$2$ part $\mathcal J^{(2)}_{a,h}$ therefore satisfies
\begin{align*}
 \mathcal J_{a,h}[f/h]\ge\mathcal J^{(2)}_{a,h}[f/h]
 &\ge\frac{\log2}{\sqrt2}
       \inf_{t-b\le x\le t+b}\frac{h(x-\ell)}{h(x)}\|f\|^2\\
 &\ge\kappa\exp\!\left(\frac{3\pi}{4}e^{2(t-b)}\right)
 =\kappa\exp(\gamma e^{2a}),
\end{align*}
using \eqref{eq:v150-weight-ratio}. No cross term between the bumps
has been dropped: the endpoint opposite every selected nonzero endpoint
lies outside the support.

The original form on these same unit vectors has the elementary upper
bound $q_a[f]\le B_a$. For the archimedean multiplier
$\alpha(\xi)=\Re\psi(1/4+i\xi/2)-\log\pi$, one may take
\[
 A=2\sum_{j=1}^2\int_{\mathbb R}|\alpha(\xi)|
                      |\widehat\phi_j(\xi)|^2\,d\xi<\infty,
\]
since the translated parity factors have absolute value at most
$\sqrt2$. The prime correlation norm is at most $2P_\lambda$, and
\[
 P_\lambda\le 2a\sum_{n<e^{2a}}n^{-1/2}\le4ae^a.
\]
The even pole is at most $2\|c\|_{L^2(I_a)}^2
=2a+2\sinh a\le2a+e^a$; the odd pole is nonpositive and does not
increase this upper bound. Thus all terms, including mixed terms within
the chosen two-dimensional space, are covered by $B_a$.

From \eqref{eq:v150-budget-form} and $(r_a)_+\ge0$,
\[
 -\mathcal B_{a,\delta}[f]\ge
       \delta_a\mathcal J_{a,h}[f/h]-q_a[f]
 \ge\delta_a\kappa\exp(\gamma e^{2a})-B_a.
\]
The capacity demand requires $\mathcal B_{a,\delta}\ge-\eta_a$ on
these unit vectors, proving \eqref{eq:v150-slack-budget} and its
rearrangement. Deleting only the prime-$2$ energy subtracts
$\mathcal J^{(2)}_{a,h}$ in place of $\delta_a\mathcal J_{a,h}$ and
has the same lower obstruction. No eigenvalue or sign of the physical
form has been inferred from a negative comparison form.
\end{proof}

\begin{assumption}[Critical arithmetic energy comparison, CAE]
\label{ass:v150-cae}
For some finite $C$ independent of the window, prove on a cofinal family
\begin{equation}\label{eq:v150-cae}
 \mathcal J_h[f/h]\ge
 \frac23\operatorname{Var}_{\nu}(f/h)
       +2\left|\int s(x)f(x)\,dx\right|^2-C\|f\|^2
\end{equation}
for every even source-admissible $f$ and every odd $f$ in the complete
physical form domain. Both parity statements and all exterior edges
are required. This is a named open input.
\end{assumption}
By \eqref{eq:v150-critical-identity}, CAE is the existing uniform
complement-floor target in explicit positive-weight coordinates, not
an independent positivity theorem. The archived source residual and
v1.36 would then give the conditional RH implication. They do not
prove CAE. The unweakened local capacity demand with $\delta_a=0$
is a stronger sufficient condition because it also discards $(r_a)_+$;
its uniform error remains unproved as well.

\paragraph{Finding and limits.}
The explicit weight exists and its identities include the complete
arithmetic shifts and both poles. What had to change is the proposed
use of a generic strict absorption estimate: even one fixed fraction
of transformed energy, or the entire prime-$2$ channel, cannot be
sacrificed with a uniform ordinary error for this weight. Edge ratios
amplify that loss like $\exp(\gamma\lambda^2)$ on admissible tests.
The remaining critical comparison CAE is an equivalent form of the
missing arithmetic bound; this calculation has not made it easier
by proof. A different weight is not excluded. The complete capacity
route stays blocked, while the specified fixed-slack shortcut is closed.
The general ground-state method is classical
\cite{FrankSeiringer2008}; no worldwide novelty is claimed for it or
for the Gaussian completed-zeta kernel. All new statements here are
exact identities, proved implications or the named open input CAE.
No RH assumption, new window, tail metric or numerical certificate
enters. G2 and RH remain open.

\subsection*{Corrected zero field, quantization, and the critical energy target (NS-39)}
\label{sec:v151-input-comparison}
The three named inputs concern different quantities. ZLD is a lower
level-band measure, QG is growth of a one-sided quantization cost, and
CAE is a signed physical-form bound. The following proposition gives
what follows from their stated definitions. It does not assert logical
independence of unproved assertions about the actual zeta function.

\begin{proposition}[Zero-field potential and the implications actually supplied]
\label{prop:v151-input-comparison}
Let $a>0$, $I_a=(-a,a)$, $h$ be the positive weight of
\eqref{eq:v150-weight}, $h_a=1_{I_a}h$ on $\mathbb R$, and
$I_{h,a}=\int_{I_a}c(x)h(x)\,dx$, where $c(x)=\cosh(x/2)$.
Use the complete corrected field $\mathcal Z_a=\beta_a$ of
\eqref{eq:v147-corrected-zero-field}, including its endpoint, trivial
zeros and prescribed symmetric-height limit. Put
\[
 Y_a=\mathcal F^{-1}(\mathcal Z_a\widehat h_a),
\]
with the unitary Fourier transform. Let $u_a$ denote the actual normalized
even source. Then the following statements hold.
\begin{enumerate}[label=(\roman*)]
\item \emph{Exact identity, fixed window.} The product defining $Y_a$
is in $L^2(\mathbb R)$, and almost everywhere on $I_a$,
\begin{align}
 r_a(x)&=\frac{Y_a(x)-2I_{h,a}c(x)}{h(x)},
       \label{eq:v151-zero-potential}\\
 T_a(x)&=Y_a(x)+2(I_h-I_{h,a})c(x).
       \label{eq:v151-zero-exterior}
\end{align}
In particular, the global $I_h$ cannot replace $I_{h,a}$ in
\eqref{eq:v151-zero-potential}. For each $f$ in the complete physical
form domain, writing $g=f/h$ on $I_a$ and zero outside,
\begin{align}
 \mathcal J_h[g]-\tfrac23\operatorname{Var}_{\nu}(g)
              -2|\langle f,s\rangle|^2
 &=q_a[f]
 =\int_{\mathbb R}\mathcal Z_a(\xi)
                   |\widehat{\widetilde f}(\xi)|^2\,d\xi,
 \label{eq:v151-zero-energy}
\end{align}
where $s(x)=\sinh(x/2)$ and $\widetilde f$ is zero extension.
The identities retain all exterior edges and both parity poles.

\item \emph{Proved implication, identical cutoff rule.} At a fixed
window with $D_a>0$, suppose the band inequality in
\eqref{eq:v147-zld} holds with constants $\kappa>0$, $0<\theta\le1$ for the
prescribed $N(a)$ and $X_a=2\pi N(a)/L$. Then, for each $K\ge1$,
\begin{equation}\label{eq:v151-zld-qg}
 \mathrm{QC}_K(\widetilde m_a)
       \ge\frac{\kappa\theta^2D_a}{2K}.
\end{equation}
On a cofinal family, uniform $\kappa,\theta>0$ and $D_a\to\infty$
therefore imply the quantization-growth clause
\eqref{eq:v149-qg} of QG. For the actual symbol, the required global
depth growth is Corollary~\ref{cor:v137-scalar-no-go}. This does
\emph{not} supply the independent source-admissible, bounded-energy
packet clause of Assumption~\ref{ass:v149-qg}. At one window QG
has no growth assertion: only the finite cost inequality is applicable.

\item \emph{Exact equivalence; the open arithmetic input.} At any
fixed window and for any specified $C\ge0$, the CAE inequality is
equivalent to
\begin{equation}\label{eq:v151-zero-floor}
 \int\mathcal Z_a|\widehat{\widetilde f}|^2
       \ge-C\|f\|^2
\end{equation}
for every even $f\perp u_a$ and every odd $f$ in the complete form
domain. A window-dependent $C=D_a$ is automatic. On a cofinal family,
a single finite $C$ independent of $a$ is precisely CAE and precisely
the uniform complement-floor target. No additional signed estimate is
obtained by either change of coordinates. The source residual and
both parity requirements in Corollary~\ref{cor:v136-complement-floor}
are still the hypotheses for transferring that target to a full floor.

\item \emph{Limits of these implications.} On families with $D\to\infty$, the quantization-growth clause is
strictly less restrictive than uniform ZLD for general level measures, and
neither level-measure property alone forces a uniform signed physical
floor. The countermodels in the proof establish only these general
statements; they are not arithmetic counterexamples. CAE itself does
not refer to $N(a)$. For the actual symbol there is a prescribed integer head rule with
every fixed-$K$ cost tending to zero, whether or not CAE holds. This
refutes QG uniformly over all head rules; it also rules out uniform
ZLD on that diluting rule. CAE and the physical form are unchanged
by the choice of rule. No converse for a particular fixed
arithmetic rule, and no further implication between actual arithmetic
QG and CAE, is proved here.
\end{enumerate}
\end{proposition}
\begin{proof}
The compact function $h_a$ has bounded variation, so
$\widehat h_a(\xi)=O_a((1+|\xi|)^{-1})$. The finite prime sum and
archimedean asymptotic give
$\mathcal Z_a(\xi)=O_a(\log(2+|\xi|))$; hence
$\mathcal Z_a\widehat h_a\in L^2$. Equivalently, $Y_a$ is the
$L^2$ limit of the inverse transforms of
$e^{-\varepsilon|\xi|}\mathcal Z_a\widehat h_a$ as
$\varepsilon\downarrow0$.
Compression of the full symbol multiplier to $I_a$ is the physical
operator $H_a+\Pi$ in both parities, with
$\Pi=2|c\rangle\langle c|-2|s\rangle\langle s|$.
On the even $h_a$, $\Pi h_a=2I_{h,a}c$. The local jump differences
of $h/h=1$ vanish, so the v1.29 identity gives $H_a h=h r_a$.
Consequently $Y_a|_{I_a}=h r_a+2I_{h,a}c$. Combining this with
$h r_a=T_a-2I_hc$ proves (i). These are restricted $L^2$ identities;
endpoint singularities are permitted. The positive minimum of $h$
controls division only at each fixed window, not uniformly in $a$.
Equation~\eqref{eq:v151-zero-energy} combines the existing complete
symbol and energy identities and extends on their physical form domain.
No finite zero prefix has been substituted for $\mathcal Z_a$, and
no interchange of its height limit with the inverse transform or
quadratic integral is asserted. Such a substitution would require its
Perron remainder, not just the small correction $C_a$.

For (ii), ZLD says that $\widetilde m_a$ dominates
$(\kappa/D_a)\,dt$ on $[-\theta D_a,0]$ (first on intervals,
then on Borel sets). An admissible set of at most $K$ minorant levels
contains the compulsory level $-D_a$. Its gaps clipped to that interval
have total length $\theta D_a$ and number at most $K$. On a clipped
gap of length $\Delta$, its one-sided distance to the preceding
level has integral at least $\Delta^2/2$. Thus its cost is at least
\[
 \frac{\kappa}{2D_a}\sum\Delta^2
 \ge\frac{\kappa\theta^2D_a}{2K}.
\]
Taking the infimum gives \eqref{eq:v151-zld-qg}. Equation
\eqref{eq:v151-zero-energy} proves (iii); $\mathcal Z_a\ge-D_a$
and Plancherel give the automatic local bound. Neither statement
makes $\sup_aD_a$ finite.

For the first limitation in (iv), let $U_J$ denote the uniform
probability measure on $J$, and take $D>1$ and
\[
 \mu_D=\tfrac12\bigl[(1-D^{-2})U_{[-D,-D/2]}
                         +D^{-2}U_{[-D,0]}\bigr].
\]
It has mass $1/2$, full support and
$\mathrm{QC}_K(\mu_D)\ge(1-D^{-2})D/(8K)$ by the same gap
argument. Nevertheless its density on $(-D/2,0)$ is only
$1/(2D^3)$, contradicting every fixed positive ZLD constant on
$[-\theta D,0]$ for any fixed $\theta>0$ eventually. Thus QG need
not give uniform all-band control for general level measures.

For failure of a general implication to a signed floor, take the
nonarithmetic, even analytic multiplier $b_D(\xi)=\xi^2-D$ and
$X_D=\sqrt D$. Its normalized negative-level density is
\[
 \frac{1}{2\sqrt D\sqrt{D+t}}\ge\frac1{2D},
              \qquad -D<t<0.
\]
It satisfies ZLD with $\theta=1$, $\kappa=1/2$ and hence QG, yet
for every fixed nonzero odd compact smooth test $f$,
\[
 \int b_D|\widehat f|^2=\|f'\|^2-D\|f\|^2\longrightarrow-\infty.
\]
The test is admissible on all sufficiently large physical windows
and orthogonal to every even source. This is not the Weil multiplier
or its Gaussian-radical identity; it shows exactly why signed
energy control cannot be inferred from level-distribution data alone.

Finally put $\ell_a=|\{\beta_a<0\}\cap[0,\infty)|<\infty$.
For the actual field the compulsory one-level minorant gives
\begin{equation}\label{eq:v151-cutoff-dilution}
 0\le\mathrm{QC}_K(\widetilde m_a)
 \le\mathrm{QC}_1(\widetilde m_a)\le D_a\ell_a/X_a.
\end{equation}
A prescribed integer rule with $X_a\ge a(1+D_a\ell_a)$ makes
this at most $1/a$. Choosing such a rule changes only the frequency
averaging range, not CAE or the physical operator. Thus QG uniformly over every rule is false, whether or not CAE
holds. This does not prove a logical nonimplication from actual
arithmetic CAE, whose truth is unknown, or disprove QG for a
particular natural rule.
\end{proof}

\paragraph{Disposition.}
CAE is the bounded-floor target rewritten in positive-weight coordinates.
ZLD implies the QG growth clause on the same prescribed family; the
converse fails for general level measures. Neither hypothesis alone
supplies the missing signed arithmetic estimate. Whether a special
property of the actual corrected zero field yields any further
implication remains open. No assertion here proves CAE, G2 or RH.


\subsection*{Source-compressed Schatten refinement and two no-go tests}

There is a sharper scalar consequence of the concentration hierarchy,
but it also has a precise limitation. Retain the notation of
Proposition~\ref{prop:v131-concentration}. For a finite-measure
symmetric set $E$, write
\[
 C_E=\mathcal F_a^*1_E\mathcal F_a,
 \qquad A_E=P_aC_EP_a\big|_{P_a\mathcal H_a^{\rm ev}}.
\]

\begin{proposition}[Exact Schatten refinement and its structural limit]
\label{prop:v132-schatten}
The operator $A_E$ is a positive trace-class contraction and
\begin{align}
 \operatorname {tr}(A_E^2)
  ={}&\operatorname {tr}(C_E^2)-2\norm{C_Eu_a}^2
       +\ip{C_Eu_a}{u_a}^{,2},
 \label{eq:v132-hs-identity}\\
 \norm{A_E}
  \le{}&\min\left\{1,\operatorname {tr}A_E,
                  \sqrt{\operatorname {tr}(A_E^2)}\right\}.
 \label{eq:v132-hs-bound}
\end{align}
Equivalently, if
$k_\xi=P_a(2\pi)^{-1/2}\cos(\xi\,\cdot)$ is the
source-compressed even evaluation vector, then
\begin{equation}\label{eq:v132-hs-kernel}
 \operatorname {tr}(A_E^2)
   =\iint_{E\times E}|\ip{k_\xi}{k_\zeta}|^2\,d\xi\,d\zeta,
\end{equation}
with the symmetric-frequency normalization already included. Hence
\begin{equation}\label{eq:v132-hs-layercake}
 q_{W_\lambda}[f]\ge-\eta_a^{(2)}\norm f^2,
 \quad
 \eta_a^{(2)}:=\int_0^{D_a}
 \min\{1,\operatorname {tr}A_{E_a(t)},
       \sqrt{\operatorname {tr}(A_{E_a(t)}^2)}\}\,dt
\end{equation}
for $f=P_af$. This improves the trace specialization of
\eqref{eq:v131-layercake} and retains the actual source exactly.

Nevertheless, vanishing of $\eta_a^{(2)}$ is not a structural
consequence of nonnegativity plus an exact source. On the even
Dirichlet space $H_0^1(-a,a)$ consider
\[
 q_a^{\rm mod}[f]=a^2\norm{f'}^2-\frac{\pi^2}{4}\norm f^2,
 \qquad
 u_a^{\rm mod}(x)=a^{-1/2}\cos\frac{\pi x}{2a}.
\]
Then $q_a^{\rm mod}\ge0$ and $u_a^{\rm mod}$ is an exact normalized
null vector, but the coefficient obtained from
\eqref{eq:v132-hs-layercake} for the symbol
$a^2\xi^2-\pi^2/4$ is a strictly positive constant independent of
$a$.
\end{proposition}
\begin{proof}
Put $U=|u_a\rangle\langle u_a|$. Cyclicity of the trace gives
\[
 \operatorname {tr}(P_aC_EP_aC_EP_a)
 =\operatorname {tr}(C_E^2)-2\operatorname {tr}(UC_E^2)
   +\operatorname {tr}(UC_EUC_E),
\]
which is \eqref{eq:v132-hs-identity}. The evaluation-vector
factorization gives \eqref{eq:v132-hs-kernel}. For a positive
trace-class contraction, its largest eigenvalue is at most one, its
trace, and the Euclidean norm of its eigenvalue sequence. This proves
\eqref{eq:v132-hs-bound}; the layer-cake proof of
Proposition~\ref{prop:v131-concentration} then proves
\eqref{eq:v132-hs-layercake}.

For the model, the sharp Dirichlet Poincar\'e inequality gives
$q_a^{\rm mod}\ge0$ with equality precisely on the displayed even
source. Its negative levels are
\[
 E_a^{\rm mod}(t)=\left\{|\xi|<a^{-1}
       \sqrt{\pi^2/4-t}\right\},
 \qquad 0<t<\pi^2/4.
\]
Under the unitary dilation from $L^2(-a,a)$ to $L^2(-1,1)$, both the
source-compressed concentration operator and every displayed band
become independent of $a$. The compressed operator is nonzero for
every $t$ in a nonempty interval. Its Hilbert--Schmidt norm therefore
has a strictly positive, $a$-independent integral. Thus the sufficient
negative-only coefficient does not vanish although the complete form
is nonnegative and the source residual is zero.
\end{proof}

Rank-one compression also cannot remove an arbitrary second
concentration channel. Cauchy interlacing gives
\begin{equation}\label{eq:v132-interlace}
 \norm{P_aC_EP_a}\ge\lambda_2(C_E),
 \qquad
 \norm{P_aC_EP_a}_{\rm HS}^2
     \ge\sum_{j\ge2}\lambda_j(C_E)^2.
\end{equation}
This does not refute the arithmetic weighted criterion
\eqref{eq:v131-weighted-concentration}; it proves that replacing its
signed sum by any negative-only Schatten moment requires an additional
arithmetic collapse of all channels after the first.

A different tempting simplification is to smooth $\beta_a$ by a
positive frequency kernel while preserving the compact-support form
exactly. The following elementary obstruction rules that out.

\begin{proposition}[No nontrivial positive flat-top smoothing]
\label{prop:v132-flat-top}
Let $\kappa$ be a probability measure on $\mathbb R$ and let
$\chi(t)=\int e^{it\xi}\,d\kappa(\xi)$. If $\chi(t)=1$ on a nonempty
open interval containing zero, then $\kappa=\delta_0$ and
$\chi\equiv1$. Consequently no nontrivial positive convolution kernel
can leave every multiplier form on functions supported in $(-a,a)$
unchanged by having a flat top on the difference interval $[-2a,2a]$.
Any nontrivial form-exact flat-top frequency kernel must change sign.
\end{proposition}
\begin{proof}
For every $t$ in the interval,
\[
 0=1-\Re\chi(t)=\int(1-\cos(t\xi))\,d\kappa(\xi).
\]
The integrand is nonnegative. Fix a nonzero $t_0$ in the interval and
intersect the resulting full-measure zero sets for $t=t_0/n$ over all
sufficiently large integers $n$. If $\xi$ belongs to that intersection,
$t_0\xi/n$ is an integral multiple of $2\pi$ for every such $n$; for
large $n$ its absolute value is below $2\pi$, hence it is zero. Thus
$\xi=0$ and $\kappa$ is concentrated at zero.
\end{proof}

The exact form-preserving observation behind the proposition is still
useful: since the autocorrelation of the zero extension of $f$ is
supported in $[-2a,2a]$, convolution by a signed kernel whose
characteristic function equals one there does not change
$\int\beta_a|\widehat{\widetilde f}|^2$. But the proposition shows
that positivity cannot come from convex averaging. Establishing a
nonnegative signed flat-top extension is a new sign problem, not an
automatic regularization of the old one.

These results sharpen the candidate register. The Hilbert--Schmidt
identity is a useful continuation of G2.6-CAP-01, not a new mechanism;
the dilation model prevents it from being advertised as the missing
cofinal argument. Positive flat-top smoothing is ruled out. The
surviving project-new mechanism remains the signed weighted operator
inequality \eqref{eq:v131-weighted-concentration}, whose exact open
lemma is unchanged: construct an explicit finite or convergent
arithmetic hierarchy with ordinary error $\eta_a\to0$ along a cofinal
family. The odd sector remains separate, and no G2 sign gap is closed.

\subsection*{Exact concentration commutators and a cofinal obstruction}

There is another tempting simplification of the signed hierarchy: the
sets can be chosen nested, so one might hope that their concentration
operators become asymptotically commuting after source compression.
The complete product algebra identifies the exact leakage terms and
shows that such commutation is not a consequence of physical scaling.

Let $J=\mathcal F_a:\mathcal H_a^{\rm ev}\to L^2_{\rm ev}(\mathbb R)$,
put $\Pi=JJ^*$, and for a symmetric measurable set $E$ define
\[
 C_E=J^*1_EJ,
 \qquad H_E=(I-\Pi)1_EJ.
\]
Thus $H_E$ is the leakage out of the physical Paley--Wiener range.  It
is a truncated Wiener--Hopf/Hankel-type channel; no half-line model is
being silently identified with it.  Retain the normalized source $u_a$,
write $P_a=I-|u_a\rangle\langle u_a|$, and on
$P_a\mathcal H_a^{\rm ev}$ put
\[
 A_E=P_aC_EP_a,
 \qquad v_E=P_aC_Eu_a.
\]

\begin{proposition}[Toeplitz--leakage identities with the actual source]
\label{prop:v133-leakage}
For symmetric measurable $E,F$, all the following are exact bounded
operator identities:
\begin{align}
 C_EC_F&=C_{E\cap F}-H_E^*H_F,\label{eq:v133-product}\\
 [C_E,C_F]&=H_F^*H_E-H_E^*H_F,\label{eq:v133-commutator}\\
 A_EA_F&=P_aC_{E\cap F}P_a-P_aH_E^*H_FP_a
          -|v_E\rangle\langle v_F|,\label{eq:v133-source-product}\\
 [A_E,A_F]&=P_a(H_F^*H_E-H_E^*H_F)P_a
 +|v_F\rangle\langle v_E|-|v_E\rangle\langle v_F|.
 \label{eq:v133-source-commutator}
\end{align}
Here $|v\rangle\langle w|f=v\ip f w$, in the fixed first-slot-linear
convention.  If
\[
 L_Ef=(H_Ef,\ip f{v_E})
 \quad\hbox{on }P_a\mathcal H_a^{\rm ev},
\]
then
\begin{equation}
 A_EA_F=A_{E\cap F}-L_E^*L_F,
 \quad [A_E,A_F]=L_F^*L_E-L_E^*L_F,
 \quad A_E-A_E^2=L_E^*L_E.\label{eq:v133-joint-leakage}
\end{equation}
Consequently
\[
 \norm{[A_E,A_F]}
 \le2\sqrt{\norm{A_E-A_E^2}\norm{A_F-A_F^2}}
 \le\frac12.
\]
\end{proposition}
\begin{proof}
Insert $I=\Pi+(I-\Pi)$ between the commuting multipliers $1_E$ and
$1_F$.  This proves \eqref{eq:v133-product}; antisymmetrization proves
\eqref{eq:v133-commutator}.  Expanding the middle
$P_a=I-|u_a\rangle\langle u_a|$ gives
\eqref{eq:v133-source-product} with the displayed rank-one term, and
antisymmetrization gives \eqref{eq:v133-source-commutator}.  The direct
sum definition of $L_E$ proves \eqref{eq:v133-joint-leakage}.  Finally,
Cauchy--Schwarz for the positive block Gram
$(L_E^*L_F)_{E,F}$ gives the first norm bound; a positive contraction
satisfies $0\preceq A_E-A_E^2\preceq I/4$.
\end{proof}

The joint Gram in \eqref{eq:v133-joint-leakage} retains phases that are
lost by separate block norms.  That fact alone is not a sign mechanism:
the Gram occurs in product defects and an antisymmetric commutator, not
as a positive factorization of the signed Weil operator.  The following
exact stress test prevents nesting or reciprocal scale from being used
as a generic small-commutator argument.

\begin{proposition}[Actual-source reciprocal bands do not asymptotically commute]
\label{prop:v133-noncommute}
Let $a=\log\lambda$ and
\[
 E_a=[-1/(10a),1/(10a)],\qquad
 F_a=[-1/(5a),1/(5a)].
\]
For the normalized actual repaired source of
Lemma~\ref{lem:v115-source-mass},
\begin{equation}
 \liminf_{a\to\infty}
 \norm{[A_{E_a},A_{F_a}]}>9\cdot10^{-6}.
 \label{eq:v133-noncommute-lower}
\end{equation}
Thus generic cofinal approximate commutation fails even after removal of
the manuscript's actual source.
\end{proposition}
\begin{proof}
Scale $L^2_{\rm ev}(-1,1)$ unitarily to
$L^2_{\rm ev}(-a,a)$ by
$(U_af)(x)=a^{-1/2}f(x/a)$.  The two uncompressed concentration
operators become the fixed operators $C_{1/10}$ and $C_{1/5}$ on
$(-1,1)$.  Here $C_c$ has, on even functions, kernel
\[
 C_c(x,y)=\frac1\pi\int_0^c\cos(tx)\cos(ty)\,dt.
\]
Put $e_0=1/\sqrt2$ and
$e_2=\sqrt{5/8}(3x^2-1)$.  With
\[
 T_c=\frac c\pi|1\rangle\langle1|
 -\frac{c^3}{6\pi}
 (|x^2\rangle\langle1|+|1\rangle\langle x^2|),
\]
Taylor's theorem gives
\[
 \norm{C_c-T_c}\le\frac{2c^5}{15\pi},
 \qquad \norm{C_c}\le\frac{2c}{\pi}.
\]
Direct rank-one algebra gives
\[
 \ip{[T_c,T_d]e_0}{e_2}
 =\frac{4cd(d^2-c^2)}{9\sqrt5\,\pi^2}.
\]
For $c=1/10,d=1/5$, the commutator remainder is at most
\[
 \frac{272}{15\cdot10^6\pi^2}
 +\frac{256}{225\cdot10^{10}\pi^2}.
\]
Using only $9<\pi^2<10$ and $\sqrt5<9/4$ therefore yields
\[
 \norm{[C_{1/10},C_{1/5}]}
 >\frac1{84375}-\frac{272}{135000000}
 -\frac{256}{20250000000000}>9\cdot10^{-6}.
\]

It remains to retain the actual source.  By
\eqref{eq:v115-source-limit}, its normalized zero extension converges
strongly on the unscaled logarithmic line to a fixed nonzero vector.
After the preceding interval dilation, the normalized vectors $z_a$
converge weakly to zero in $L^2(-1,1)$: tightness controls the unscaled
tails, while absolute continuity of the $L^2$ integral controls every
fixed compact core after dilation.  Each $C_c$ is compact, so
$\norm{C_cz_a}\to0$.  If $Q_a=I-|z_a\rangle\langle z_a|$, then
\[
 \norm{Q_aC_cQ_a-C_c}
 \le2\norm{C_cz_a}+|\ip{C_cz_a}{z_a}|\longrightarrow0.
\]
The source-compressed commutator therefore converges in operator norm to
$[C_{1/10},C_{1/5}]$, proving \eqref{eq:v133-noncommute-lower}.
\end{proof}

This proposition concerns explicit interval bands, not the arithmetic
superlevel sets of $\beta_a$.  It does not falsify the signed hierarchy;
it proves that its commutators cannot be declared small from nesting,
reciprocal scaling, or the source limit alone.  Even exact commutation
would not supply weighted coverage: commuting contractions
$A_j=\theta_jI$ still make $-DI+\sum_jw_jA_j$ negative whenever
$\sum_jw_j\theta_j<D$.

For cofinite good sets $G_{a,j}$ put $B_{a,j}=\mathbb R\setminus G_{a,j}$
and define the complete joint analysis operator
\[
 Z_af=(\sqrt{w_{a,j}},1_{B_{a,j}}\mathcal F_af)_j,
 \qquad f\perp u_a.
\]
Then the remaining signed criterion is exactly
\begin{equation}
 \norm{Z_a}^2
 \le\sum_jw_{a,j}-D_a+\eta_a,
 \qquad \eta_a\longrightarrow0.\label{eq:v133-joint-channel}
\end{equation}
Its channel Gram retains all overlaps and the source before taking a
norm.  Equation~\eqref{eq:v133-joint-channel} is nevertheless
algebraically equivalent to \eqref{eq:v131-weighted-concentration}; it
is recorded as the exact surviving implementation target, not as a new
candidate or a proved estimate.

The product identity is the classical Toeplitz--Hankel mechanism; see
Brown--Halmos~\cite{BrownHalmos1964} and the product-formula review
\cite{Virtanen2022}.  The present even-interval and source-compressed
formulas were derived directly, with no transfer of a half-line
normalization.  Formula comparison also separates them from the v1.28
common-Gram proposal: the current Gram is automatic for concentration
\emph{product defects}, whereas v1.28 required a favorable factorization
of the actual Weil off-blocks.  Calling the former a new positivity
mechanism would be a relabeling.  The project-new result here is the
actual-source cofinal obstruction in
Proposition~\ref{prop:v133-noncommute}.  The arithmetic signed estimate
\eqref{eq:v133-joint-channel} and the odd sector remain open; no G2 sign
gap is closed.

\subsection*{Disjoint-channel persistence and optimal primitive transport}

The preceding commutator obstruction does not by itself exclude a
Cotlar--Stein estimate: a row sum may tolerate finitely many adjacent
interactions.  The next result identifies what survives for disjoint
channels and then tests a genuinely different scalar mechanism.

\begin{proposition}[Disjoint reciprocal channels retain a cofinal cross term]
\label{prop:v134-disjoint}
On the even Paley--Wiener range write
\[
 \Pi_a=\mathcal F_a\mathcal F_a^*,\qquad
 \Pi_{a,u}=\Pi_a-|\widehat{\widetilde u_a}\rangle
                    \langle\widehat{\widetilde u_a}|,
\]
where $u_a$ is the normalized repaired source.  Put
\[
 E_a=[-1/(10a),1/(10a)],\qquad
 F_a=[-3/(10a),-1/(5a)]\cup[1/(5a),3/(10a)].
\]
Then
\begin{equation}
 \liminf_{a\to\infty}
 \norm{1_{E_a}\Pi_{a,u}1_{F_a}}
 \ge \frac{73}{375\pi}>0.0619.                 \label{eq:v134-disjoint}
\end{equation}
Consequently disjointness and fixed-ratio reciprocal scaling do not
make the actual-source cross channels vanish.  Moreover, for any
atomization of the bad-set multiplier
$m_a=\sum_jw_{a,j}1_{B_{a,j}}$, the complete analysis operator satisfies
\[
 Z_a^*Z_a=P_a\mathcal F_a^*m_a\mathcal F_aP_a.
\]
Splitting or regrouping channels does not change the operator to be
bounded.
\end{proposition}
\begin{proof}
After the unitary frequency dilation, take
\[
 E=[-1/10,1/10],\qquad
 F=[-3/10,-1/5]\cup[1/5,3/10].
\]
Both have measure $1/5$.  Test the full Paley--Wiener projection on the
normalized even indicators $e_E=\sqrt5\,1_E$ and
$e_F=\sqrt5\,1_F$.  Its kernel is
\[
 \Pi(\xi,\eta)=\frac{\sin(\xi-\eta)}{\pi(\xi-\eta)}.
\]
For $\xi\in E$ and $\eta\in F$, one has
$|\xi-\eta|\le2/5$ and
\[
 \frac{\sin t}{t}\ge1-\frac{t^2}{6}\ge\frac{73}{75}.
\]
Therefore
\[
 \ip{1_E\Pi1_Fe_F}{e_E}
 \ge 5|E||F|\frac{73}{75\pi}
 =\frac{73}{375\pi}.
\]
The even compression gives the same pairing on even test vectors.
The source correction has norm
$\norm{1_{E_a}\widehat{\widetilde u_a}}
 \norm{1_{F_a}\widehat{\widetilde u_a}}=o(1)$:
the strong source limit \eqref{eq:v115-source-limit} and absolute
continuity of its $L^2$ integral control the shrinking bands.  This
proves \eqref{eq:v134-disjoint}.  The identity for $Z_a^*Z_a$ follows
directly by summing the channel adjoint products.
\end{proof}

Thus a Cotlar--Stein proof would still have to retain the coherent
adjacent interactions and prove the complete constant in
\eqref{eq:v133-joint-channel}.  In particular, $J$ identical copies of
one channel with weights $w/J$ have the same complete Gram as one
channel of weight $w$, although every individual diagonal norm tends
to zero.  A row-norm estimate that discards those cross terms is the
earlier block comparison in different coordinates, not a new
mechanism.

There is, however, a distinct way to retain favorable and unfavorable
values of the scalar symbol.  It uses signed primitive transport rather
than concentration operators.  The derivative estimate below is the
$L^2$ Bernstein inequality for the Paley--Wiener space; see
\cite{Bang1995}.

\begin{proposition}[Optimal signed-primitive/Bernstein criterion]
\label{prop:v134-primitive}
Let $\beta:\mathbb R\to\mathbb R$ be locally integrable, bounded below,
and tend to $+\infty$ at both ends.  Define
\begin{equation}
 \Delta(\beta)=
 \sup_{x<y}\left(-\int_x^y\beta(\xi)\,d\xi\right)_+ .
 \label{eq:v134-drawdown}
\end{equation}
For every $f\in L^2(-a,a)$ whose zero extension has Fourier transform
$F$, one has
\begin{equation}
 \int_{\mathbb R}\beta(\xi)|F(\xi)|^2\,d\xi
 \ge-a\Delta(\beta)\norm f^2.                 \label{eq:v134-primitive-bound}
\end{equation}
The coefficient is optimal within all scalar decompositions
$\beta=b+W'$ with $b\ge0$ and $W\in L^\infty$: after adding constants
to $W$,
\begin{equation}
 \inf\norm W_\infty=\frac12\Delta(\beta).     \label{eq:v134-optimal}
\end{equation}
Hence the arithmetic condition
\begin{equation}
 a\Delta(\beta_a)\longrightarrow0             \label{eq:v134-open-lemma}
\end{equation}
would prove the even estimate required by
Proposition~\ref{prop:v118-gap-free}, on the whole physical space and
therefore on the actual source complement, without a positive spectral
gap.

This scalar mechanism is not forced by positivity plus an exact source.
For the Dirichlet model of Proposition~\ref{prop:v132-schatten},
\[
 \beta_a^{\rm mod}(\xi)=a^2\xi^2-\frac{\pi^2}{4},
 \qquad
 \Delta(\beta_a^{\rm mod})=\frac{\pi^3}{6a}.
\]
Thus its optimal primitive error is the nonzero constant $\pi^3/6$,
although the physical form is nonnegative and its source is an exact
null vector.
\end{proposition}
\begin{proof}
Let $B(x)=\int_0^x\beta$ and
$M(x)=\sup_{t\le x}B(t)$.  The hypotheses make the maximum drawdown of
$B$ finite.  The function $M$ is locally absolutely continuous,
$b=M'\ge0$ almost everywhere, and
\[
 W=B-M+\frac12\Delta(\beta)
\]
satisfies $\beta=b+W'$ and
$\norm W_\infty\le\Delta(\beta)/2$.  Conversely, every decomposition
$\beta=b+W'$ with $b\ge0$ gives, for $x<y$,
\[
 W(x)-W(y)=\int_x^y(b-\beta)
 \ge-\int_x^y\beta.
\]
Hence the oscillation of $W$ is at least $\Delta(\beta)$, proving
\eqref{eq:v134-optimal}.

Now $F\in H^1(\mathbb R)$ and, in the unitary Fourier normalization,
\[
 \norm{F'}_2=\norm{xf}_2\le a\norm f.
\]
Distributional integration by parts, followed by Cauchy--Schwarz,
gives
\[
 \int\beta|F|^2
 =\int b|F|^2-\int W(|F|^2)'
 \ge-2\norm W_\infty\norm{F'}_2\norm F_2,
\]
which is \eqref{eq:v134-primitive-bound}.  For the model, direct
integration over its negative interval gives
\[
 \Delta(\beta_a^{\rm mod})
 =\int_{-\pi/(2a)}^{\pi/(2a)}
   \left(\frac{\pi^2}{4}-a^2\xi^2\right)d\xi
 =\frac{\pi^3}{6a}.
\]
\end{proof}

The primitive criterion is formulaically distinct from the Schur,
structured-inverse, block-metric, common-Gram, Picone and concentration
routes: it uses one-dimensional signed transport plus the physical
Paley--Wiener derivative bound, and introduces no auxiliary spectral
gap.  Its ingredients are classical, and worldwide novelty is not
claimed.  The exact drawdown optimization is new within this project.
It does not survive the present falsification threshold as a live
candidate. A historical non-rigorous direct scan, whose original script
and output are not archived, reported for the exact arithmetic symbol $a\Delta(\beta_a)$ approximately $2.5448,3.4002,4.9148,5.5052$
at $\lambda=3,4,5,6$, respectively.  These finite values prove neither
failure nor success of \eqref{eq:v134-open-lemma}; they show no trace of
the required decay.  The candidate register therefore marked the scalar
primitive route as rejected. Corollary~\ref{cor:v137-scalar-no-go}
below now proves divergence of this actual arithmetic budget, even
under the weaker bounded-error objective; a source-sensitive replacement
would need an essentially different estimate.  The live even-sector obligation remains the full
signed weighted concentration inequality
\eqref{eq:v131-weighted-concentration}; the odd sector remains open and
no G2 sign gap is closed.

\subsection*{Both parities and a quantitative growing radical block}

The signed concentration route can be formulated on the complete physical
space.  This extension does not infer the odd sign from the even sign.
Retain $a=\log\lambda$, $I=(-a,a)$, the unitary Fourier transform of
zero extensions, and the canonical closed operator $\mathsf W_\lambda$.
Its form domain is
\[
 \mathcal D(q_a)=\left\{f\in L^2(I):
 \int_{\mathbb R}\log(2+|\xi|)
       |\widehat{\widetilde f}(\xi)|^2d\xi<\infty\right\}.
\]

\begin{proposition}[The same exact symbol in both parities]
\label{prop:v135-both-parities}
For every complex $f\in\mathcal D(q_a)$,
\begin{equation}\label{eq:v135-full-symbol}
 q_a[f]=\int_{\mathbb R}\beta_a(\xi)
                    |\widehat{\widetilde f}(\xi)|^2d\xi,
\end{equation}
with exactly the symbol, prime cutoff and normalization in
\eqref{eq:v130-r-symbol}--\eqref{eq:v130-beta}.
Let $u_a$ be the actual normalized even source and let
$P_a=I-|u_a\rangle\langle u_a|$ now act on all of $L^2(I)$.
The weighted concentration criterion
\eqref{eq:v131-weighted-concentration}, with this full-space meaning,
is equivalent to its two parity blocks.  If their respective ordinary
negative errors are $\eta_a^{\rm ev}$ and $\eta_a^{\rm odd}$, and
$\norm{\mathsf W_\lambda u_a}\le\epsilon_a$, then
\begin{equation}\label{eq:v135-full-error}
 \mathsf W_\lambda\succeq-
 \max\left\{\eta_a^{\rm odd},
      \eta_a^{\rm ev}+\frac2{\sqrt3}\epsilon_a\right\}I.
\end{equation}
For a symmetric set $E$ of finite measure, the uncompressed parity
concentration traces are
\begin{equation}\label{eq:v135-parity-traces}
 \operatorname{tr}C_E^{\rm ev/odd}
 =\int_E\frac{a\mathbin{\pm}\sin(2a\xi)/(2\xi)}{2\pi}\,d\xi.
\end{equation}
Source removal subtracts $\int_E|\widehat{\widetilde u_a}|^2$ only
from the even trace.  In particular the complete compressed trace is
$a|E|/\pi-\int_E|\widehat{\widetilde u_a}|^2$.
\end{proposition}
\begin{proof}
The full pole has kernel $2\cosh((x-y)/2)$ and is
$2|c\rangle\langle c|-2|s\rangle\langle s|$.  The elementary
identity $2\cosh(t/2)-e^{|t|/2}=e^{-|t|/2}$ therefore gives
\[
 2|c\rangle\langle c|-2|s\rangle\langle s|-K_a^+=K_a^-.
\]
Apply the multiplier and digamma-recurrence calculation of
Proposition~\ref{prop:v130-remainder}; parity is no longer needed.
The differences from the archimedean form are bounded at fixed $a$,
so the identity extends from the smooth core to the displayed domain.
The symbol and all the symmetric level sets commute with reflection;
the even source projection is the identity on the odd sector.
Apply Proposition~\ref{prop:v118-gap-free} only to the even block and
take the worse of the two lower errors to obtain
\eqref{eq:v135-full-error}.  The squared norms of the cosine and sine
evaluation vectors give \eqref{eq:v135-parity-traces}.  There are no
cross-parity form terms, also for complex tests.
\end{proof}

The next result quantifies why a positive complement gap is an
unnecessarily strong objective.  The radical construction itself is
classical~\cite[Section 3]{ConnesConsani2023}.  The earlier research
log already used even derivatives of its Gaussian example to rule out
a fixed gap after removal of a fixed-rank block.  Here separated
translations give a uniformly invertible physical Gram and an explicit
\emph{growing-rank ordinary operator} estimate.  This is a continuation
of that obstruction, not a new positivity mechanism.

Use $h_\infty$ and $f_\infty$ from
Lemma~\ref{lem:v115-source-mass}, and put
\[
 k(x)=f_\infty(e^x),\qquad \phi=k/\norm k_2,\qquad
 M_2=\norm{(1+|x|)^2\phi}_2,\qquad
 D=\left\lceil4\pi\sqrt{M_2}\right\rceil.
\]
Thus $D$ is a fixed integer defined by the actual physical source,
not a parameter fitted at each window.  Fix a smooth $\vartheta$ with
$0\le\vartheta\le1$, $\vartheta(t)=1$ for $t\le-1$, and
$\vartheta(t)=0$ for $t\ge-1/2$.  For $a\ge2$ put
\[
 \chi_a(x)=\vartheta(|x|-a),\quad
 J_a=\left\lfloor\frac{a}{2D}\right\rfloor,\quad
 g_{a,j}(x)=\chi_a(x)\phi(x-jD)\quad(|j|\le J_a).
\]
All columns lie in $C_c^\infty(I)\subset\mathcal D(\mathsf W_\lambda)$.
Let $G_a$ be their synthesis map and $\Gamma_a=G_a^*G_a$.

\begin{proposition}[Growing physical blocks with a vanishing full residual]
\label{prop:v135-growing-radical}
There are constants $C,a_0$, depending only on the fixed source and
cutoff, such that, for $a\ge a_0$,
\begin{align}
 \Gamma_a&\succeq\tfrac12 I,
       &m_a:=\operatorname{rank}G_a&=2J_a+1,\label{eq:v135-gram}\\
 \norm{\mathsf W_\lambda P_{G_a}}
       &\le\varepsilon_a,
       &\varepsilon_a&:=Ca\exp(-\lambda/100).
       \label{eq:v135-block-residual}
\end{align}
Here $P_{G_a}=G_a\Gamma_a^{-1}G_a^*$ is the ordinary physical
orthogonal projection.  The same bound, after enlarging $C$ if
necessary, holds on the odd block of dimension $J_a$ spanned by
$(g_{a,j}-g_{a,-j})/\sqrt2$, $1\le j\le J_a$.
These are complete operators, with no Fourier truncation.
\end{proposition}
\begin{proof}
The Gaussian source formula, its Fourier invariance, and Poisson
summation give $k(-x)=k(x)\ne0$.  Differentiating its normally
convergent series on $x\ge0$, and then reflecting, gives, for each
fixed $r\ge0$,
\begin{equation}\label{eq:v135-gaussian-tail}
 |k^{(r)}(x)|\le C_r\exp\{-\tfrac\pi2 e^{2|x|}\}.
\end{equation}
Indeed a derivative only introduces polynomial factors in $e^x$ and
$ne^x$ into the Gaussian series; half its exponent absorbs them.
All translates are polarized global radicals: translation by $t$
corresponds to the even Schwartz source
$h_t(u)=e^{-t/2}h_\infty(e^{-t}u)$, which has zero value and zero
integral.  The Schwartz radical theorem applies to arbitrary compact
logarithmic tests, not just to even ones.  No positive global closed
operator is assumed here.

For $s\ne0$, splitting at $|x|=|s|/2$ and using the weighted norm
on whichever translated factor lies in a tail gives
\[
 |\ip{\phi(\cdot-s)}\phi|
 \le\frac{2M_2}{(1+|s|/2)^2}.
\]
Hence the sum of all off-diagonal absolute entries in any row of
the uncut lattice Gram is at most
\[
 4M_2\sum_{n\ge1}(1+nD/2)^{-2}
 \le\frac{8\pi^2M_2}{3D^2}\le\frac16.
\]
Every finite uncut Gram is therefore at least $5I/6$.

Write $r_{a,j}=(1-\chi_a)\phi(\cdot-jD)$ on the entire line.
For $|jD|\le a/2$, its support satisfies $|x|\ge a-1$, hence
$|x-jD|\ge a/2-1$.  Equation~\eqref{eq:v135-gaussian-tail} and
the bounded cutoff derivatives imply the uniform estimate
\begin{equation}\label{eq:v135-weighted-tail}
 \norm{r_{a,j}}_{H^1(\mathbb R)}+
 \sup_x e^{2|x|}\bigl(|r_{a,j}(x)|+|r'_{a,j}(x)|\bigr)
 \le\delta_a:=C\exp\{-\pi\lambda/(4e^2)\}.
\end{equation}
For clarity, with $X=e^{2|x-jD|}\ge\lambda/e^2$ one has
$e^{2|x|}\le\lambda X\le e^2X^2$; absorbing $X^2$ into half
the exponent in \eqref{eq:v135-gaussian-tail} proves the supremum
bound.  Its integrable envelope also proves the $H^1$ bound.
The cutoff Gram loss is positive semidefinite, and its trace is
at most $Cm_a\delta_a^2$, using the same estimate for the
uncut tails.  Since $m_a=O(a)$, this loss is eventually below $1/3$.
This proves \eqref{eq:v135-gram}.

It remains to bound the \emph{complete} residual, not merely its
quadratic value.  Denote the global explicit-formula distribution by
$\mathcal W$.  Radicality gives, on $I$,
\[
 \mathsf W_\lambda g_{a,j}=-\mathcal W r_{a,j}.
\]
This distributional equality has an $L^2(I)$ representative as follows.
The full-line archimedean multiplier
$\Re\psi(1/4+i\xi/2)-\log\pi$ is bounded by $C(1+|\xi|)$,
so its $L^2$ action on $r_{a,j}$ is at most $C\delta_a$.
The prime part includes \emph{every} $n\ge2$, also those beyond
$\lambda^2$.  Since $|x\pm\log n|\ge\log n-|x|$,
\[
 \sum_{n\ge2}\frac{\Lambda(n)}{\sqrt n}
   |r_{a,j}(x\pm\log n)|
 \le\delta_a e^{2|x|}
             \sum_{n\ge2}\frac{\Lambda(n)}{n^{5/2}}
 \le C\delta_a e^{2a}\qquad(x\in I).
\]
The series converges absolutely.  Both pole integrals are bounded
by $C\delta_a e^{|x|/2}$, using their kernels $e^{\pm(x-y)/2}$.
Thus
\[
 \norm{\mathsf W_\lambda g_{a,j}}_2
 \le C\sqrt a\,\lambda^2\delta_a
 \le C\sqrt a\,e^{-\lambda/100}.
\]
The last inequality absorbs a polynomial into the strictly stronger
exponent $\pi/(4e^2)>1/100$.  The operator domain is legitimate
because each $g_{a,j}$ is smooth and compactly supported inside $I$.
Now $G_a\Gamma_a^{-1/2}$ is an isometry, so
\[
 \norm{\mathsf W_\lambda P_{G_a}}
 =\norm{\mathsf W_\lambda G_a\Gamma_a^{-1/2}}
 \le\sqrt2\left(\sum_j
             \norm{\mathsf W_\lambda g_{a,j}}_2^2\right)^{1/2}.
\]
This is at most $Ca e^{-\lambda/100}$.  Reflection exchanges
indices $j$ and $-j$.  The antisymmetric coefficient embedding is
an isometry, so its Gram lower bound and residual estimate follow
from the same argument; it produces exactly $J_a$ odd columns.
\end{proof}

\begin{corollary}[A growing-rank obstruction to positive coercivity]
\label{cor:v135-coercivity-obstruction}
For every physical subspace $F_a\subset L^2(I)$ with
$\dim F_a<m_a$, there is a unit vector
$v_a\in F_a^\perp\cap\mathcal D(\mathsf W_\lambda)$ such that
\begin{equation}\label{eq:v135-complement-ceiling}
 \norm{\mathsf W_\lambda v_a}\le\varepsilon_a,
 \qquad |q_a[v_a]|\le\varepsilon_a.
\end{equation}
Consequently any positive coercivity constant on $F_a^\perp$ is
at most $\varepsilon_a$.  This excludes a uniform positive constant
or a positive polynomial-in-$\lambda^{-1}$ lower bound when
$\dim F_a=o(\log\lambda)$.
If all deep and plunge image columns in (G2.6) are even, their
low-concentration complement contains the complete odd block above,
regardless of the number of those columns; the same ceiling applies.
More generally it still applies if fewer than $J_a$ independent odd
constraints are added.
The spectral projection of $\mathsf W_\lambda$ onto
$[-2\varepsilon_a,2\varepsilon_a]$ has rank at least $m_a$,
including at least $J_a$ odd eigenvalues.
\end{corollary}
\begin{proof}
Dimension gives a nonzero vector in
$\operatorname{Ran}G_a\cap F_a^\perp$; normalize it and apply
\eqref{eq:v135-block-residual} and Cauchy--Schwarz.
Odd vectors are orthogonal to every even image, which proves the
physical (G2.6) assertion without identifying it with $u_a^\perp$.
If the indicated spectral projection had rank below $m_a$, a unit
vector in $\operatorname{Ran}G_a$ could be chosen orthogonal to its
range.  The spectral theorem would give
$\norm{\mathsf W_\lambda v_a}\ge2\varepsilon_a$, a contradiction.
The same proof on the odd block gives its count. Compact resolvent
is already established in Proposition~\ref{prop:v114-weil-core}.
\end{proof}

The result determines neither side of zero for these eigenvalues.
In particular it is compatible with RH and with the gap-free target.
For the explicitly constructed $P_{G_a}$,
Proposition~\ref{prop:v118-gap-free} still requires the independent
arithmetic estimate
\[
 q_a[f]\ge-\eta_a\norm f^2\quad
 (f\perp\operatorname{Ran}G_a),\qquad \eta_a\longrightarrow0.
\]
The signed concentration criterion on the whole actual rank-one
complement remains another live conditional target.  Neither estimate
follows from the small residual.  No source-to-physical norm transfer,
Fourier-cut limit, or sampler graph identity was used or modified.
The logarithmic diagonal, physical endpoints and graph defect remain
as previously stated.  No G2 sign gap is closed.


\subsection*{Dense radicals, ordinary resolvent collapse, and a bounded-error reduction}

All operators in this subsection are the complete physical Weil operators,
not the corrected logarithmic generators in a Weil-metric completion.
Put $H=L^2(\mathbb R,dx)$, $I_a=(-a,a)$, $\lambda=e^a$, and let
$J_a:L^2(I_a)\to H$ be zero extension. Inner products remain linear in
the first slot. Write $q_a$ for the canonical closed form of
$\mathsf W_{e^a}$. On compactly supported smooth tests it is the
restriction of a single support-consistent geometric form $q$.
The form domain and both pole signs are those of
Propositions~\ref{prop:v118-log-dirichlet} and
\ref{prop:v135-both-parities}. No Fourier cutoff is imposed.

\begin{lemma}[A total family of localized ordinary near-radicals]
\label{lem:v136-total-radicals}
Use the normalized source $\phi$ and smooth physical cutoff $\chi_a$
from Proposition~\ref{prop:v135-growing-radical}. Then
\[
 \mathcal R=\operatorname{span}\{\phi(\cdot-t):t\in\mathbb R\}
 \quad\hbox{is dense in }H.
\]
For every fixed $h\in\mathcal R$, $g_a=\chi_a h$ belongs to
$C_c^\infty(I_a)$ and satisfies
\begin{equation}\label{eq:v136-dense-residual}
 J_ag_a\longrightarrow h,\qquad
 \norm{\mathsf W_{e^a}g_a}_2\longrightarrow0.
\end{equation}
The even and odd reflected translate spans are dense in their respective
parity subspaces. These are ordinary-norm assertions; they assert neither
uniform approximation coefficients nor density in a global form norm.
\end{lemma}
\begin{proof}
The double-exponential estimate \eqref{eq:v135-gaussian-tail} makes
$\widehat\phi$ entire. It is not identically zero, so its real zeros
have measure zero. If $v\in H$ is orthogonal to every translate, the
inverse Fourier transform of the $L^1$ function
$\widehat v\,\overline{\widehat\phi}$ vanishes. Fourier uniqueness
gives $\widehat v\,\overline{\widehat\phi}=0$ almost everywhere,
and hence $v=0$. This is the elementary Hilbert-space version of the
classical translation-density principle~\cite{Wiener1932}.
Apply the bounded parity projections to obtain the two parity statements.

Write $h=\sum_{j=1}^r c_j\phi(\cdot-t_j)$ with fixed real $t_j$.
For sufficiently large $a$, every $|t_j|\le a/2$. The exterior-tail
proof of Proposition~\ref{prop:v135-growing-radical} uses only that
inequality, not membership in the spaced lattice. It therefore gives
\[
 \norm{\mathsf W_{e^a}(\chi_a\phi(\cdot-t_j))}_2
 \le C\sqrt a\,e^{-e^a/100}.
\]
The triangle inequality proves the residual assertion with a constant
depending on this fixed representation of $h$. Dominated convergence
proves the ordinary limit. In particular all global prime rows in the
exterior remainder, both poles and the logarithmic diagonal are retained.
\end{proof}

\begin{proposition}[Ordinary strong-resolvent collapse]
\label{prop:v136-resolvent-collapse}
Define the self-adjoint operator
\[
 A_a=\mathsf W_{e^a}\oplus0
 \quad\hbox{on }H=L^2(I_a)\oplus L^2(I_a^c),
 \qquad
 \mathcal D(A_a)=\mathcal D(\mathsf W_{e^a})\oplus L^2(I_a^c).
\]
Then, without a uniform lower bound or a sign hypothesis,
\begin{equation}\label{eq:v136-resolvent-collapse}
 (A_a-z)^{-1}\longrightarrow-z^{-1}I
 \quad\hbox{strongly for every }z\in\mathbb C\setminus\mathbb R.
\end{equation}
If $P_a$ is an orthogonal projection with range in $\mathcal D(A_a)$
and $\norm{A_aP_a}\to0$, the same conclusion holds for the operator
which deletes its block and assigns zero on its range. Precisely, define
\[
 A_a^{\perp}=A_a-A_aP_a-P_aA_a+P_aA_aP_a
 \quad\hbox{on }\mathcal D(A_a).
\]
The terms involving $P_aA_a$ mean their bounded extensions.
This applies to the actual repaired rank-one source and to the growing
block of Proposition~\ref{prop:v135-growing-radical}.
\end{proposition}
\begin{proof}
For $h\in\mathcal R$ use $g_a$ from Lemma~\ref{lem:v136-total-radicals},
regarded as a vector in $H$. The resolvent bound for self-adjoint
operators and the resolvent identity give
\begin{align*}
 \norm{(A_a-z)^{-1}h+z^{-1}h}
 &\le\left(\frac1{|\Im z|}+\frac1{|z|}\right)\norm{h-g_a}
       +\frac{\norm{A_ag_a}}{|z|\,|\Im z|}.
\end{align*}
This tends to zero. Uniform resolvent boundedness and density extend
the conclusion to every $h\in H$. For the deleted block, the difference
is a bounded self-adjoint perturbation of norm at most
$3\norm{A_aP_a}$. The resolvent identity bounds the new resolvent
difference by $3\norm{A_aP_a}/|\Im z|^2$.
\end{proof}

This concerns the ordinary $L^2$ realization of the Weil operator.
It does not identify a limit of the logarithmic generators in the
different Weil-metric spaces discussed, for example, in
\cite[Section 7]{Suzuki2026}. Nor does it give convergence of
$q_a[f]$ for fixed compactly supported $f$; those values are exactly
support-consistent already. Strong-resolvent convergence alone does
not control unbounded spectral moments or lower spectral edges.
Indeed, on $\ell^2$, the operators
$T_n=-n|e_n\rangle\langle e_n|$ converge to zero in this topology,
while $\inf\sigma(T_n)=-n$. Even replacing $n$ by $1$ leaves a
fixed negative edge. These are operator-topology countermodels,
not arithmetic Weil tests.

Even support consistency and a dense radical do not by themselves
fix the sign. On $c_{00}\subset\ell^2$ take
$q_\pm[x]=\pm|\sum_kx_k|^2$; its $n$-coordinate matrices are
$\pm\mathbf1_n\mathbf1_n^*$. Their common limiting radical,
the finite sequences of sum zero, is dense: for example
$e_1-m^{-1}\sum_{k=2}^{m+1}e_k\to e_1$.
Both zero-extended matrix families converge strong-resolvent to zero,
since their unit nonzero-eigenvalue direction
$n^{-1/2}\mathbf1_n$ tends weakly to zero. The positive family is
nonnegative and the negative family has lower edge $-n$.
The latter amplifies the form value $-1$ to Rayleigh value $-m$
on $m^{-1}\sum_{k=2}^{m+1}e_k$.

\begin{proposition}[A uniform finite lower floor suffices]
\label{prop:v136-bounded-floor}
Suppose that for a cofinal sequence $a_j\to\infty$ there is one
finite constant $C_*$, independent of $j$, such that
\begin{equation}\label{eq:v136-bounded-floor}
 q_{a_j}[f]\ge-C_*\norm f_2^2
 \quad\hbox{for all }f\in\mathcal D(q_{a_j}).
\end{equation}
Then $q[f]\ge0$ on every compact smooth test. Consequently every
complete finite-window form is nonnegative, and the Weil criterion
implies RH. Equivalently, if any compact smooth test has negative
Weil form, then
\begin{equation}\label{eq:v136-negative-divergence}
 \inf\sigma(\mathsf W_{e^a})\longrightarrow-\infty
 \qquad(a\longrightarrow\infty).
\end{equation}
This is a conditional reduction: no bound \eqref{eq:v136-bounded-floor}
is established here.
\end{proposition}
\begin{proof}
Suppose a compact smooth $f$ has $q[f]=-d<0$. Given $\varepsilon>0$,
choose a fixed $h\in\mathcal R$ with $\norm{f-h}<\varepsilon/2$.
Put $g_a=\chi_a h$ and $w_a=f-g_a$ for sufficiently large windows.
By support consistency and self-adjointness,
\[
 |q_a[w_a]-q[f]|
 \le (2\norm f+\norm{g_a})\norm{\mathsf W_{e^a}g_a}
 \longrightarrow0,\qquad
 \norm{w_a}\longrightarrow\norm{f-h}<\varepsilon/2.
\]
Thus eventually $q_a[w_a]<-d/2$ and $0<\norm{w_a}<\varepsilon$,
so the variational principle gives
\begin{equation}\label{eq:v136-negative-amplification}
 \inf\sigma(\mathsf W_{e^a})
 \le\frac{q_a[w_a]}{\norm{w_a}^2}
 <-\frac{d}{2\varepsilon^2}.
\end{equation}
Since $\varepsilon$ is arbitrary, this proves
\eqref{eq:v136-negative-divergence}, and contradicts every cofinal
finite lower floor. Therefore \eqref{eq:v136-bounded-floor} forces
$q[f]\ge0$ for every compact smooth $f$.
Compact smooth tests are a form core at each fixed window by
Proposition~\ref{prop:v118-log-dirichlet}, so complete local
positivity follows. The fixed-support Weil implication is the same
one used in Proposition~\ref{prop:v121-cofinal-rh}.
The argument gives no effective window threshold for a prescribed
$\varepsilon$: the finite translate coefficients and shifts can grow
as its ordinary approximation is refined.
\end{proof}

\begin{corollary}[Bounded-error version of the gap-free target]
\label{cor:v136-complement-floor}
Let $P_a$ be either the actual repaired source projection or the
growing radical projection, with its established
$\norm{\mathsf W_{e^a}P_a}\to0$. It is sufficient to prove, on a
cofinal family, the single uniform estimate
\begin{equation}\label{eq:v136-complement-target}
 q_a[f]\ge-C_*\norm f^2,
 \qquad f\in\mathcal D(q_a)\cap\operatorname{Ran}(I-P_a),
 \qquad C_*<\infty\text{ independent of }a.
\end{equation}
Both parity sectors must be included. In the signed concentration
criterion \eqref{eq:v131-weighted-concentration}, bounded cofinal
errors in both sectors therefore suffice; proving their decay to zero
is a stronger sufficient objective. This corollary does not identify
$C_a^{\rm low}$ with either displayed complement.
\end{corollary}
\begin{proof}
Proposition~\ref{prop:v118-gap-free} gives a full lower floor
$-C_*-(2/\sqrt3)\norm{\mathsf W_{e^a}P_a}$, which is uniformly
bounded along the family. Apply Proposition~\ref{prop:v136-bounded-floor}.
The even and odd blocks combine as in \eqref{eq:v135-full-error}.
For the physical deep/plunge decomposition, any missing block and
coupling still needs its own justified bound.
\end{proof}

\begin{proposition}[No persistent bounded positive comparison]
\label{prop:v136-positive-comparison}
Let $u_a$ be the actual ordinary-unit source, with $J_au_a\to\phi$
and $\norm{\mathsf W_{e^a}u_a}\to0$, and let
$Q_a=I-|u_a\rangle\langle u_a|$. Suppose
\[
 B_a=Q_aB_aQ_a\succeq0,\qquad\sup_a\norm{B_a}<\infty,
 \qquad
 q_a[f]\ge\ip{B_af}{f}-\eta_a\norm f^2\quad(f\perp u_a),
 \qquad\eta_a\to0,
\]
on the complete form domain. Then
$J_aB_aJ_a^*\to0$ strongly on $H$.
In particular a fixed nonzero positive bounded comparison on
$\phi^\perp$ cannot satisfy these assumptions.
If instead the comparison is imposed on the actual $C_a^{\rm low}$
and all its removed physical image columns are even, any uniformly
bounded positive comparison must tend strongly to zero on every
fixed odd vector. No identification of the two complements is used.
\end{proposition}
\begin{proof}
For $h\in\mathcal R$ set $f_a=Q_a(\chi_a h)$. The source assumptions
and Lemma~\ref{lem:v136-total-radicals} give
\[
 J_af_a\longrightarrow(I-|\phi\rangle\langle\phi|)h,
 \qquad\norm{\mathsf W_{e^a}f_a}\longrightarrow0.
\]
Consequently $\ip{B_af_a}{f_a}\to0$. If $M$ bounds $\norm{B_a}$,
positivity gives $\norm{B_af_a}^2\le M\ip{B_af_a}{f_a}\to0$.
The limiting projected translates are dense in $\phi^\perp$;
uniform boundedness extends the assertion to that subspace.
Compression and $J_au_a\to\phi$ give convergence on $\phi$ too.
For $C_a^{\rm low}$ use only odd reflected translates. Their odd
cutoffs belong to that physical complement exactly, and their span
is dense in the odd subspace. The same positive-square argument works
even when the comparison operator itself does not preserve parity.
\end{proof}

Uniform boundedness and the chosen topology matter: escaping rank-one
projections can have norm one and strong limit zero, while
$n^2|e_n\rangle\langle e_n|$ shows why testing a dense family without
a common norm bound is insufficient. A fixed closed nonnegative form
has a related obstruction only when all the projected cutoff tests
belong to its domain: lower semicontinuity then makes it vanish on
$\phi^\perp$, not necessarily on the source direction. This assertion
uses ordinary $L^2$ closedness and does not concern a Weil-metric
quotient completion.

The essential new project ingredient is totality of \emph{all fixed real
translates}, rather than a single uniformly separated growing lattice.
The classical translation-density and radical identities are not claimed
as new. These results strengthen the reduction and exclude persistent
positive comparisons; they do not supply an independent arithmetic
lower bound. The live signed concentration construction remains open,
now with the sufficient target \eqref{eq:v136-complement-target}.
No G2 sign gap is closed. The finite sampler, its explicit graph defect,
the physical endpoint and the existing Fourier scale factors are unchanged.

\subsection*{Why a bounded scalar primitive error cannot supply the finite floor}

The weaker target in Corollary~\ref{cor:v136-complement-floor} makes it
necessary to revisit the scalar primitive route: its earlier failure to
give an $o(1)$ error does not itself exclude an $O(1)$ error. The following
argument rules out that weaker scalar target for the \emph{actual}
arithmetic symbol. It does not rule out the signed concentration operator.
Keep $a=\log\lambda$, $T=2a$, the exact $\beta_a$ of
\eqref{eq:v130-beta}, and $\Delta_a=\Delta(\beta_a)$ from
\eqref{eq:v134-drawdown}.

\begin{lemma}[A positive probe with a negative truncated Fourier mass]
\label{lem:v137-probe}
Fix $\delta=1/100$ and an even nonnegative
$\kappa\in C_c^\infty(-\delta,\delta)$ with integral one. Put
\[
 w=\mathbf1_{[\pi,3\pi]}*\kappa,\qquad
 H(t)=\int_{\mathbb R}w(s)e^{-its}\,ds,\qquad
 K(u)=\int_{-1}^1H(t)e^{iut}\,dt.
\]
Then $w\ge0$, $\int w=2\pi$, $\norm{w'}_1=2$, and
\begin{equation}\label{eq:v137-probe-properties}
 H(1)=H(-1)=0,\qquad
 K(0)<-\frac1{3\pi},\qquad
 |K(u)|\le\frac{C_w}{1+u^2}.
\end{equation}
Here $K$ is real on the real line; no symmetry of $w$ about zero is
required.
\end{lemma}
\begin{proof}
The two translates in $w'(s)=\kappa(s-\pi)-\kappa(s-3\pi)$ have
disjoint supports, giving the stated variation. With the nonunitary
transform used only in this lemma,
\[
 H(t)=2e^{-2\pi it}\frac{\sin(\pi t)}{t}\widehat\kappa_{\rm nu}(t),
\]
where the quotient is interpreted continuously at zero. This proves
the endpoint zeros. Also
\begin{align*}
 \int_\pi^{3\pi}\frac{\sin s}{s}\,ds
 &=-\pi\int_0^\pi
       \frac{\sin u}{(\pi+u)(2\pi+u)}\,du
 \le-\frac1{3\pi}.
\end{align*}
Since $\sin s/s=\int_0^1\cos(ts)\,dt$ has derivative bounded by
$1/2$, convolution with $\kappa$ changes that integral by at most
$\pi\delta$. Therefore
\[
 K(0)=2\int_{\mathbb R}\frac{\sin s}{s}w(s)\,ds
 \le-\frac2{3\pi}+2\pi\delta<-\frac1{3\pi}.
\]
The last strict inequality follows from $6\pi^2<100$.
Twice integrating the definition of $K$ by parts, using $H(\pm1)=0$,
gives the claimed $u^{-2}$ decay; for example the numerator can be
bounded by $|H'(1)|+|H'(-1)|+\int_{-1}^1|H''|$.
Finally $H(-t)=\overline{H(t)}$, so $K(u)$ is real.
\end{proof}

\begin{proposition}[Conditional negative mass beside a critical zero]
\label{prop:v137-conditional-mass}
Assume RH. Fix any ordinate $\gamma_0$ of a nontrivial zero, of
multiplicity $m_0$, and use Lemma~\ref{lem:v137-probe}. Then
\begin{equation}\label{eq:v137-probe-limit}
 J_a:=\int_{\mathbb R}\beta_a(\xi)
              w\bigl(T(\xi-\gamma_0)\bigr)\,d\xi
 =m_0K(0)+O_{\gamma_0,w}(T^{-2})+E_T,
\end{equation}
where
\begin{equation}\label{eq:v137-exterior-error}
 |E_T|\le\frac{8\pi e^{-5T/2}}{5T(1-e^{-2T})}.
\end{equation}
Consequently, for all sufficiently large $a$,
\begin{equation}\label{eq:v137-conditional-lower}
 \Delta_a\ge\frac1{6\pi},\qquad
 \inf_{\xi\in\mathbb R}\beta_a(\xi)
      \le-\frac{a}{6\pi^2}.
\end{equation}
The threshold is not asserted effective here, and these displayed rates
are conditional on RH.
\end{proposition}
\begin{proof}
Write $\mathcal W$ for the complete geometric Weil distribution in
logarithmic coordinates. Under RH its spectral form is
$\mathcal W(t)=\sum_\gamma m_\gamma e^{i\gamma t}$ in the distributional
sense, with distinct real ordinates in this notation; this is the
classical explicit formula, not a new spectral assumption beyond RH
\cite[Theorem 2.1]{BurnolExplicit2000}.
The inverse-transform kernel of $\beta_a$ agrees with $\mathcal W$
on $|t|<T$. Outside that interval it is exactly
\[
 -\rho_*(|t|),\qquad
 \rho_*(t)=\frac{e^{-5t/2}}{1-e^{-2t}}\quad(t>0).
\]
Indeed this is the off-origin kernel of
$\Re\psi(5/4+i\xi/2)-\log\pi$. The prime and continuum terms in
\eqref{eq:v130-r-symbol} have support in $[-T,T]$; inside that interval
the identity
$-\rho_*(|t|)+e^{|t|/2}=-\rho(|t|)+2\cosh(t/2)$
retains the complete pole and archimedean terms. All distributions at
zero have the same normalization as \eqref{eq:v135-full-symbol}.

The nonunitary Fourier transform of the probe in \eqref{eq:v137-probe-limit}
is
\[
 V_T(t)=T^{-1}e^{-i\gamma_0t}H(t/T).
\]
The explicit formula applied to $\mathbf1_{[-T,T]}V_T$ therefore gives
the exact identity
\begin{equation}\label{eq:v137-zero-sum}
 J_a=\sum_\gamma m_\gamma K\bigl(T(\gamma-\gamma_0)\bigr)+E_T,
 \qquad
 E_T=-\int_{|t|>T}\rho_*(|t|)V_T(t)\,dt.
\end{equation}
There is no additional $2\pi$ in this identity: $V_T$ is the
nonunitary transform of the frequency test and $K$ is defined above
without a normalization factor. The ordinary physical Fourier transform
in \eqref{eq:v135-full-symbol} remains unitary.

For completeness the truncated test is legitimate. It is continuous,
compactly supported, smooth near zero and piecewise smooth elsewhere;
its endpoint values vanish because $H(\pm1)=0$. Its second distributional
derivative is a finite measure. Smooth mollification gives convergence
on the geometric side and Fourier decay $O((1+|\gamma|)^{-2})$
uniformly in the mollification parameter, for each fixed $T$.
The bound $N(Y)=O(Y\log(2+Y))$ then permits dominated convergence on
the RH zero sum. This extends the smooth explicit formula to this test.
If $e^T$ is a prime power, the endpoint atom pairs with zero; the strict
prime cutoff is therefore preserved exactly.

The ordinate $\gamma_0$ is isolated. Thus
\[
 \sum_{\gamma\ne\gamma_0}
       \frac{m_\gamma}{|\gamma-\gamma_0|^2}<\infty,
\]
and \eqref{eq:v137-probe-properties} bounds all other zero terms by
$O_{\gamma_0,w}(T^{-2})$. Since $|H(t)|\le2\pi$, integration of
$\rho_*$ bounds $E_T$ by \eqref{eq:v137-exterior-error}.
This proves \eqref{eq:v137-probe-limit}, and eventually
$J_a\le-1/(6\pi)$.

The scalar decomposition of Proposition~\ref{prop:v134-primitive}
gives $\beta_a=b_a+U_a'$ with $b_a\ge0$ and
$\norm{U_a}_\infty=\Delta_a/2$. Since the rescaled probe has variation
two, integration by parts yields
\[
 J_a\ge-\frac{\Delta_a}{2}
      \norm{(w(T(\cdot-\gamma_0)))'}_1=-\Delta_a.
\]
This proves the first inequality in \eqref{eq:v137-conditional-lower}.
Its mass is $2\pi/T=\pi/a$, which proves the second.
\end{proof}

\begin{corollary}[Unconditional scalar-budget obstruction]
\label{cor:v137-scalar-no-go}
For the exact arithmetic symbol of this manuscript,
\begin{equation}\label{eq:v137-scalar-divergence}
 a\Delta(\beta_a)\longrightarrow+\infty,
 \qquad
 \inf_{\xi\in\mathbb R}\beta_a(\xi)\longrightarrow-\infty
 \quad(a\longrightarrow\infty).
\end{equation}
In particular neither a bounded scalar primitive budget nor a bounded
pointwise negative error can supply the finite floor of
Proposition~\ref{prop:v136-bounded-floor}, even along a cofinal sequence.
No unconditional linear rate is claimed.
\end{corollary}
\begin{proof}
If $a_j\Delta(\beta_{a_j})$ were bounded on any sequence
$a_j\to\infty$, Proposition~\ref{prop:v134-primitive} would give a
uniform finite lower bound for the complete physical forms. The full
parity identity \eqref{eq:v135-full-symbol} and
Proposition~\ref{prop:v136-bounded-floor} would then imply RH.
But under RH, Proposition~\ref{prop:v137-conditional-mass} forces
$a_j\Delta(\beta_{a_j})\ge a_j/(6\pi)$ eventually, a contradiction.
Absence of a bounded cofinal subsequence is precisely the first limit
in \eqref{eq:v137-scalar-divergence}.
Likewise a cofinal pointwise lower bound $\beta_{a_j}\ge-C$ would
give a complete physical lower floor, imply RH, and contradict the
second estimate in \eqref{eq:v137-conditional-lower}.
This proves the second limit.
\end{proof}

This is not a negative physical Weil test. The nonnegative frequency
probe $w(T(\xi-\gamma_0))$ is not asserted to be the squared transform
of a vector supported in $(-a,a)$. Such an assertion would contradict
the RH positivity used in the conditional part of the argument.
The result identifies exactly the cost of discarding physical
time--frequency correlations. Both parities remain within the complete
symbol identity, but no projection, low-concentration identification,
source-to-physical transfer or sampler graph formula is changed.

The nearest project proposal is the v1.34 scalar primitive estimate,
retested under v1.36's weaker bounded-error objective. The new ingredient
is an endpoint-vanishing positive frequency probe, which gives an
absolutely convergent zero sum and a surviving negative cutoff side lobe.
This is a proved obstruction to that route, not a new positive mechanism.
The explicit formula and cutoff oscillations are classical. Nearby
pointwise-envelope barriers, for example \cite{Zhu2026}, bound the raw
prime comb; the present statement concerns the continuum-cancelled
symbol and its optimal signed primitive. No worldwide novelty is claimed.
The joint signed concentration estimate remains open: it must retain
the favorable levels jointly with their physical concentration operators
to obtain a uniform finite ordinary lower error. No G2 sign gap is closed.


\subsection*{What an arithmetic-location-insensitive certificate can and cannot mean}

There is no unrestricted theorem that every location-insensitive estimate
fails. The relevant information loss must be specified. Two precise
obstructions follow below: one concerns scalar density relaxations for
the actual multiplier, and one concerns its orientation relative to the
physical Fourier space. They do not identify all the earlier closures
with one hypothesis.

Put $a=\log\lambda$, $I_a=(-a,a)$, and let $\mathcal F_a$ be zero
extension followed by the unitary Fourier transform. Retain the
first-slot-linear convention and the exact $\beta_a$ in
\eqref{eq:v135-full-symbol}. Its form domain is
\[
 \mathcal D_a=\left\{f\in L^2(I_a):
 \int_{\mathbb R}\log(2+|\xi|)|\mathcal F_af(\xi)|^2d\xi<\infty\right\}.
\]
Indeed $\beta_a(\xi)=\log(2+|\xi|)+O_a(1)$; this is the same
complete physical form domain, not a new source norm. There is no
Fourier cutoff in this subsection. Either use the full space or use
the actual source complement $Q_aL^2(I_a)$, where
$Q_a=I-|u_a\rangle\langle u_a|$ and the established complete
residual $\norm{\mathsf W_{e^a}u_a}\to0$ is retained. Both parities
are included. No equality with $C_a^{\rm low}$ is asserted.

For an ordinary unit vector $f$, set
\[
 p_f(\xi)=\tfrac12\bigl(|\mathcal F_af(\xi)|^2+
                               |\mathcal F_af(-\xi)|^2\bigr).
\]
Evenness of $\beta_a$ gives $q_a[f]=\int\beta_ap_f$. The elementary
physical constraints are
\begin{equation}\label{eq:v139-density-constraints}
 p_f\ge0,\qquad \int p_f=1,\qquad
 \norm{p_f}_\infty\le\frac a\pi,\qquad
 \operatorname{TV}(p_f)\le2a.
\end{equation}
The height bound is Cauchy--Schwarz on an interval of length $2a$.
For $F=\mathcal F_af$, Plancherel gives
$\norm{F'}_2=\norm{xf}_2\le a$, and
$\norm{(|F|^2)'}_1\le2\norm{F'}_2\norm F_2\le2a$.
Reflection averaging preserves these bounds.

\begin{theorem}[Obstruction for relaxations admitting the cutoff probe]
\label{thm:v139-probe-relaxation}
Let $\mathcal R_a$ be a collection of even nonnegative probability
densities containing every $p_f$ from the full physical space, or
from its complete source complement as just specified. Assume that,
for all sufficiently large $a$, it also contains
\begin{equation}\label{eq:v139-probe-density}
 p_a^\star(\xi)=\frac a{2\pi}
 \left\{w\bigl(2a(\xi-\gamma_0)\bigr)
       +w\bigl(2a(-\xi-\gamma_0)\bigr)\right\},
\end{equation}
where $w$ is the fixed probe of Lemma~\ref{lem:v137-probe} and
$\gamma_0$ is any fixed critical-zero ordinate. If $B_a\ge0$ is a
sound relaxed error, meaning
\begin{equation}\label{eq:v139-relaxed-error}
 \int\beta_a p\ge-B_a\qquad(p\in\mathcal R_a),
\end{equation}
then $B_a\to+\infty$. In particular such a relaxation cannot supply
even a uniform finite lower floor along a cofinal family.

One explicit eligible relaxation is the set of all even $p\in
W^{1,1}(\mathbb R)$ satisfying \eqref{eq:v139-density-constraints}.
Thus retaining mass, height and total variation, with their displayed
physical constants, still loses information essential to this route.
Under RH the sharper conclusion $B_a\ge a/(6\pi^2)$ holds eventually;
that rate is not asserted unconditionally.
\end{theorem}
\begin{proof}
The unsymmetrized density $(a/\pi)w(2a(\xi-\gamma_0))$ has mass one,
height at most $a/\pi$, and variation $2a/\pi\le2a$.
Symmetrization preserves those inequalities and gives
\eqref{eq:v139-probe-density}. Compact smoothness makes its pairing
with $\beta_a$ finite. By evenness and
Proposition~\ref{prop:v137-conditional-mass}, under RH
\[
 \int\beta_ap_a^\star=\frac a\pi J_a
       \le-\frac a{6\pi^2}
\]
for all sufficiently large $a$. This proves the conditional rate.
If $B_{a_j}$ were bounded on any cofinal sequence, inclusion of all
physical densities would give a complete uniform lower floor. In
the source-complement case apply
Corollary~\ref{cor:v136-complement-floor}; in the full-space case
apply Proposition~\ref{prop:v136-bounded-floor}. Either implication
gives RH, and hence contradicts the just proved conditional rate.
The absence of a bounded cofinal subsequence proves the asserted
unconditional divergence.
\end{proof}

The quantifier over both parities is essential to this proof.
For separate parity errors it proves unconditional divergence of
their maximum, not unconditional divergence of each merely from
one sector's boundedness. The probe is not asserted to be a physical
squared transform or to satisfy the sharper source-dependent
pointwise evaluation bounds. In fact a nonzero compactly supported
frequency density cannot be such a squared transform: the entire
transform of a compact physical function cannot vanish on a real
open interval. The obstruction is precisely an excess of tests
in the relaxation, not a negative physical Weil test.

There is a particularly explicit location-insensitivity property.
For a real symbol $b$ bounded below with finite-measure sublevel sets,
define $b^\uparrow(s)$ to be its increasing rearrangement with respect
to Lebesgue measure. The strongest universal lower bound using only
unit mass and the height bound $a/\pi$ is
\begin{equation}\label{eq:v139-bathtub}
 \ell_a^{\rm cap}(b):=
 \inf_{\substack{p\ge0,\ \int p=1\\p\le a/\pi}}
        \int bp
 =\frac a\pi\int_0^{\pi/a} b^\uparrow(s)\,ds.
\end{equation}
This is unchanged by any measure-preserving rearrangement of $b$:
it retains the full distribution of favorable and unfavorable values
but discards their physical frequency positions. The formula is the
classical bathtub principle~\cite[Section 1.14]{LiebLoss2001}.
For clarity, choose a threshold $t$ whose lower sublevel has measure
at most $\pi/a$ and whose closed sublevel has measure at least
$\pi/a$, and fill that lowest region with density $a/\pi$, partially
filling the level $t$ if needed. For every competitor,
$(b-t)(p-p_*)\ge0$ pointwise; its integral and equal total masses
prove minimality and \eqref{eq:v139-bathtub}.
For even $b$, symmetrizing does not change the value.
Theorem~\ref{thm:v139-probe-relaxation} implies
\begin{equation}\label{eq:v139-cap-divergence}
 \ell_a^{\rm cap}(\beta_a)\longrightarrow-\infty.
\end{equation}
Any lower-bound rule valid for \emph{every} mass-and-height-admissible
density is at most \eqref{eq:v139-bathtub}. This is the exact scope of
the distribution-only assertion; it does not prohibit every possible
theorem whose input happens to include a histogram.

The scalar primitive criterion is included by a different verified
hypothesis: its decomposition $b=c+V'$, $c\ge0$,
$\norm V_\infty=\Delta(b)/2$, is valid against every nonnegative
unit density of variation at most $2a$. Thus its error $a\Delta(b)$
satisfies \eqref{eq:v139-relaxed-error}. Notice that $\Delta(b)$
itself is not rearrangement invariant. Calling it distribution-only
would be incorrect.

\begin{theorem}[Spectral-orientation obstruction with protected finite data]
\label{thm:v139-protected-orbit}
Fix $a$, let $M$ be multiplication by the exact $\beta_a$ on
$L^2(\mathbb R)$ with its multiplication-operator domain, and let
$S$ be any finite-dimensional frequency subspace. Let $\mathcal U_S$
consist of unitaries which fix $S$ pointwise, differ from the identity
by finite rank, and preserve $\mathcal D(M)$. Then
\begin{equation}\label{eq:v139-orbit-edge}
 \inf_{U\in\mathcal U_S}
 \inf_{\substack{f\in C_c^\infty(I_a),\ \norm f=1\\
                 \mathcal F_af\perp S}}
 \ip{U^*MU\mathcal F_af}{\mathcal F_af}
       =\operatorname*{ess\,inf}\beta_a.
\end{equation}
For any finite physical subspace $E$ with
$\mathcal F_aE\subset\mathcal D(M)$, one can take
$S\supset\mathcal F_aE+M\mathcal F_aE$. Every allowed unitary then
preserves the complete operator action on the protected columns:
\begin{equation}\label{eq:v139-protected-columns}
 U^*MU\mathcal F_ae=M\mathcal F_ae\qquad(e\in E).
\end{equation}
Consequently a sound universal compression bound insensitive to all
the allowed orientations, even after retaining these exact columns,
cannot exceed $\operatorname*{ess\,inf}\beta_a$. Along the actual
arithmetic family that upper limit tends to $-\infty$.
The result holds in either parity, with parity-preserving unitaries
and reflection-invariant protected data.
\end{theorem}
\begin{proof}
Write $m=\operatorname*{ess\,inf}\beta_a$. The lower bound in
\eqref{eq:v139-orbit-edge} is immediate from $M\succeq mI$.
Fix $r>m$. The spectral subspace $1_{\{\beta_a<r\}}L^2$ is
infinite dimensional because its support has positive Lebesgue measure;
its vectors lie in $\mathcal D(M)$ since $m\le\beta_a<r$ there.
Choose a unit vector $v$ in this subspace orthogonal to $S$.
Choose a unit smooth physical $f$ satisfying the finitely many
conditions $\mathcal F_af\perp S$, and put $h=\mathcal F_af$.
Such $f$ exists because the smooth physical space is infinite
dimensional. Its rapidly decreasing transform belongs to
$\mathcal D(M)$. On the space $\operatorname{span}\{h,v\}\subset
S^\perp$, choose a unitary carrying $h$ to $v$, and extend it by the
identity. Then $U-I$ has finite rank with range in $\mathcal D(M)$;
both $U$ and $U^*$ preserve that domain. Its Rayleigh value on $h$
is $\ip{Mv}v<r$. Let $r\downarrow m$.

If $S$ contains both protected columns and their images, $U$ fixes
each of them, which proves \eqref{eq:v139-protected-columns}.
For the constructed unitary, expansion of $U=I+K$ shows that
$U^*MU-M=K^*M+MK+K^*MK$ extends to a bounded finite-rank
self-adjoint perturbation: $MK$ is bounded and $K^*M$ is its
adjoint on the original domain. Thus the compressed forms also
keep their original domain. Evenness of $\beta_a$ leaves infinite
dimensional negative spectral subspaces in each parity, so the same
construction inside either sector proves the parity statement.
Finally apply Corollary~\ref{cor:v137-scalar-no-go}.
\end{proof}

For the actual finite Fourier source proxy, the additional domain
condition is automatic: its zero extension is piecewise smooth,
its transform is $O_a((1+|\xi|)^{-1})$, and multiplication by
$\beta_a=O_a(\log(2+|\xi|))$ remains square integrable. Protecting
it therefore preserves its exact compressed source residual, not
just the source norm or its Rayleigh value. Any finite number of
other domain columns can also be protected, even if their number
grows with $a$. This does not preserve the action on an infinite
physical block.

Spectral compression and orientation are classical concerns; see
Fan--Pall~\cite{FanPall1957}. The project-specific extension here is
the explicitly protected graph data and its application to the
actual divergent symbol. Most importantly, $U^*MU$ generally is
\emph{not} a multiplication operator. The theorem rules out bounds
required to be valid on that larger operator class. It does not
construct a rearranged prime sequence, a negative arithmetic Weil
test, or a no-go for methods retaining the multiplier's physical
support geometry. It is not a new finite-rank metric repair.

\begin{proposition}[Why mere location independence is too broad]
\label{prop:v139-location-counterexample}
Let $D_a>0$, choose any measurable set $E_a$ of measure
$\pi/(2a(D_a+1))$, and put
$b_a=1-(D_a+1)1_{E_a}$. Then for every physical $f$,
\[
 \int b_a|\mathcal F_af|^2\ge\tfrac12\norm f^2,
\]
independently of the location or shape of $E_a$, even if
$D_a\to\infty$ and hence $\operatorname*{ess\,inf}b_a\to-\infty$.
\end{proposition}
\begin{proof}
Use $|\mathcal F_af|^2\le(a/\pi)\norm f^2$ and integrate over
$E_a$. The subtracted contribution is at most
$(D_a+1)(a/\pi)|E_a|\norm f^2=\norm f^2/2$.
\end{proof}

This example is not the arithmetic symbol. It proves that neither
negative depth alone nor the informal phrase ``ignores arithmetic
locations'' implies the requested universal conclusion. The actual
symbol obstruction in Theorem~\ref{thm:v139-probe-relaxation} uses
the signed cutoff-probe pairing, and not just its minimum.
The proved scalar pointwise and primitive closures fall within the
displayed classes. The Schatten, signed block-norm, noncommuting
band, Cotlar, and flat-top results have additional geometric or
structural hypotheses; their recorded closures are not corollaries
of a single insensitivity assertion. Their precise coverage is
tabulated in \texttt{evidence/v139/adversarial\_review.md},
Sections~6--7 (the ten-route coverage table, NS-4b). The detailed
hypothesis-by-hypothesis analysis is retained in
\path{evidence/diag_ns4b_coverage/coverage.md}; it places one of the
ten recorded closures in the probe class and none in the protected-orbit
class as stated. The surviving joint signed
concentration operator retains information excluded by both
relaxations. Its cofinal finite-floor bound is still unproved.
All endpoint, prime-cut, logarithmic-diagonal, Fourier normalization
and sampler graph-defect conventions remain unchanged. No G2 sign
gap is closed, and no worldwide novelty or publication readiness is
asserted.


\subsection*{A fixed finite-floor experiment and its cutoff cost}

Proposition~\ref{prop:v136-bounded-floor} permits a fixed ordinary
negative error of any finite size. This does make some finite-window
certificates cheaper. It does not by itself remove the scalar-tail
cost in Proposition~\ref{prop:v125-cutoff-cost}.

\begin{proposition}[Additive shifts and the remaining cutoff barrier]
\label{prop:v138-shifted-floor}
In the orthonormal Fourier splitting of
Proposition~\ref{prop:v116-directional-schur}, let $\delta\ge0$ and
suppose $T\succeq D_g$, with increasing $g_n$ and
$g_{N+1}+\delta>0$. For a fixed finite solve $Z$ set
\[
 K_\delta=K_Z+\delta(I+Z^*Z),\qquad
 R_\delta=R-\delta Z.
\]
If $U_J$ encloses the remote residual beyond the support of $Z$, then
\begin{equation}\label{eq:v138-shifted-lower}
 \mathcal L_\delta:=K_\delta-
 \sum_{N<n\le J}\frac{R_{\delta,n}^*R_{\delta,n}}{g_n+\delta}
 -\frac{U_J}{g_{J+1}+\delta}\succeq0
 \quad\Longrightarrow\quad
 \mathsf W_\lambda\succeq-\delta I.
\end{equation}
The displayed sum is in one normalized parity basis. Verifying both
parities proves the inequality on the full complex form domain.
For the present scalar comparison with $c=0$, its prerequisite forces
\begin{equation}\label{eq:v138-shifted-cost}
 N+1>L\exp(M_{\phi,\lambda}-\delta),\qquad L=2\log\lambda.
\end{equation}
Consequently any bounded choice of $\delta$ retains
$\liminf_{\lambda\to\infty}\log((N+1)/L)/\lambda\ge1$.
A polynomial cutoff in this particular comparison would instead
require $\delta\ge M_{\phi,\lambda}-O(\log\lambda)$.
\end{proposition}
\begin{proof}
Apply the verified-solve identity to $\mathsf W_\lambda+\delta I$.
Its off-diagonal block is still $B$, its tail is $T+\delta I$, and
expansion gives the stated $K_\delta$ and $R_\delta$. Inverse order
gives $(T+\delta I)^{-1}\preceq(D_g+\delta I)^{-1}$.
The shifted remote rows beyond the solve support are unchanged.
The moment enclosure and completion of the square prove
\eqref{eq:v138-shifted-lower} on the closed form domain.
For the cost, positivity of the first scalar weight implies
$\log((N+1)/L)>M_{\phi,\lambda}-\delta$, since its other losses
are nonnegative. Use
$M_{\phi,\lambda}\ge\|\mathsf T_{{\rm pr},\lambda}\|$ and
Proposition~\ref{prop:v119-prime-scale}. This is a limitation of
the chosen scalar comparison, not of every matrix tail metric.
\end{proof}

Fresh outward evaluation of all window-dependent constants gives the
following smallest admitted integer $N\ge1$ with $g_{N+1}+\delta>0$:
\begin{center}
\begin{tabular}{rrrrr}
\toprule
$\lambda$ & $\delta=10^{-3}$ even & $\delta=10^{-3}$ odd
 & $\delta=8$ even & $\delta=8$ odd\\
\midrule
5&1501&7217&1&4\\
6&7481&35985&3&14\\
8&124317&598024&43&204\\
\bottomrule
\end{tabular}
\end{center}
Each threshold except the lower endpoint $N=1$ has a certified
negative predecessor. These are only remote-tail gates. They show
why a moderately sized constant, rather than $10^{-3}$, is the useful
finite experiment suggested by the bounded-floor reduction.

\begin{proposition}[Complete finite floors at three windows]
\label{prop:v138-three-floors}
For each $\lambda\in\{5,6,8\}$ and both reflection sectors,
\begin{equation}\label{eq:v138-three-floors}
 \mathsf W_\lambda\succeq-8I,\qquad
 -8\le\inf\sigma(\mathsf W_\lambda^{\pm})<10^{-16}.
\end{equation}
The lower estimate includes the complete infinite Fourier tail;
the upper estimate comes from a certified finite-support trial.
By physical support consistency, the lower estimate also holds for
every $1<\lambda\le8$. No cofinal lower estimate is asserted.
\end{proposition}
\begin{proof}[Complete interval certificate]
Use the canonical form and logarithmic domain of
Proposition~\ref{prop:v114-weil-core}, with ordinary physical
$L^2$ normalization and its established Fourier form core. Freeze
the same parameters at all three windows:
\[
 \delta=8,\quad N=256,\quad J=4096,\quad r=16,
 \quad\tau=1/10,\quad Z=0.
\]
The even head is $0{:}256$, the odd head $1{:}256$. There is no
conjugate-gradient solve or positive $10^{-8}D$ tail prerequisite.
In the remote-moment formula take $M=N$, so $G$ is the head embedding
and $G^*G=I$. The moment proof only requires $M<J$; a nonempty
intermediate solve block is unnecessary when $Z=0$.
Here $K_8=F+8I$, $R_8=B$, and all actual rows through $J$ are retained.
The entire remainder beyond $J$ is bounded by
Proposition~\ref{prop:v116-remote-moment}, including its leading
moment Gram and geometric remainder.

At each window recompute $L$, every strict-cutoff prime power,
the separate logarithmic diagonal, $M_{\phi,\lambda}$, the parity
tail constants and $B_*$. The coefficient enclosure retains $64$
archimedean series terms and the analytic infinite-series remainder.
The first runs at $160$ bits freeze inverse-Cholesky congruences as
exact IEEE dyadics. A fresh $256$-bit evaluation, using the identical
dyadics, verifies every row of
$C\mathcal L_8C^*$ is greater than $999/1000$ after subtraction
of the absolute off-diagonal row sum. The matrices $C$ are triangular
with nonzero diagonal. Thus every complete lower matrix is positive.
The code also checks exact finite-support Rayleigh quotients: each
is strictly positive and less than $10^{-16}$. These are upper
bounds for the true lower spectral edge, not signs of that edge.
The two parities give the full complex-space inequality.
Finally, zero extension of a test from a smaller physical interval
preserves its ordinary norm and Weil form. Apply the $\lambda=8$
bound and use the form core to obtain the assertion for $\lambda\le8$.
\end{proof}

The near-one congruence row margins are not invariant headroom.
For a comparable certificate quantity put
$\mathcal U_8=K_8-\mathcal L_8$ and
\[
 \mathfrak h_\lambda^\pm
 =1-\lambda_{\max}(\mathcal U_8,K_8).
\]
The same exact congruence verifies
$C\mathcal L_8C^*\succeq(999/1000)I$ and
$CK_8C^*\preceq M I$. Hence
$\mathfrak h_\lambda^\pm\ge(999/1000)/M$.
The following are downward-rounded lower bounds for this generalized
margin of the \emph{shifted} certificate:
\begin{center}
\begin{tabular}{rrrrr}
\toprule
$\lambda$ & even $g_{257}+8$ & odd $g_{257}+8$
 & even $\mathfrak h$ lower & odd $\mathfrak h$ lower\\
\midrule
5&6.2297&4.6549&0.8430&0.9190\\
6&4.6225&3.0457&0.7166&0.7620\\
8&1.8111&0.2302&0.2474&0.1022\\
\bottomrule
\end{tabular}
\end{center}
These margins quantify enclosure headroom with the same target and
construction. Their decrease is not evidence of a negative physical
eigenvalue; the actual lower edge remains bracketed only between
$-8$ and the small positive trial quotient. A worsening lower bound
cannot establish the divergence in
\eqref{eq:v136-negative-divergence}. A negative Rayleigh upper bound
would have a different logical role, and none was found here.

This is a continuation of the existing Schur/moment and shifted-inertia
routes, not a new arithmetic positivity mechanism. General certified
finite-window lower and trial upper bounds also appear in
\cite{Zhu2026}; no external numerical certificate is imported.
The actual unproved target remains one finite ordinary lower constant
on a cofinal family, on the full physical form or on a complement
meeting the residual hypotheses of Proposition~\ref{prop:v136-bounded-floor}.
The present bound covers every physical subspace at the listed windows
without identifying $C_a^{\mathrm{low}}$ with a rank-one complement.
All endpoint, scale and sampler graph-defect obligations are unchanged.
No growing-window G2 gap is closed.


\subsection*{Finite trial positivity and the exact omitted-tail correction}

The trial-head term in Proposition~\ref{prop:v120-structured-uniform}
must be distinguished from the exact Schur complement. A finite-support
certificate can settle the former without settling the latter.

\begin{proposition}[Finite-support trial forms and Galerkin Schur bounds]
\label{prop:v121-finite-trial}
Let the outer head be finite dimensional, let $T\succeq cI$ with
$c>0$ be a positive self-adjoint tail operator, and let $B$ be the
bounded head-to-tail coupling. For $X$ with range in $\mathcal D(T)$,
define $K_X,r$ as in Proposition~\ref{prop:v120-structured-uniform}
and put $S_\infty=F-B^*T^{-1}B$. Then
\begin{align}
 K_X-S_\infty
  &=(X-T^{-1}B)^*T(X-T^{-1}B)\succeq0,
       \label{eq:v121-square}\\
 S_\infty&=K_X-r^*T^{-1}r.\label{eq:v121-residual}
\end{align}
Let $P_M$ be the orthogonal projection onto a finite tail space in $\mathcal D(T)$, and
write $T_M=P_MT|_{\operatorname{Ran}P_M}$, $B_M=P_MB$ and
$S_M=F-B_M^*T_M^{-1}B_M$. If $X=P_MX$, then
\begin{equation}\label{eq:v121-finite-trial}
 K_X=\begin{bmatrix}I\\-X\end{bmatrix}^{*}
       \begin{bmatrix}F&B_M^*\\B_M&T_M\end{bmatrix}
       \begin{bmatrix}I\\-X\end{bmatrix}.
\end{equation}
In particular, $S_M\succeq0$ implies $K_X\succeq0$ for every such
$X$, and strict positivity transfers as well. This is the complete
operator's $K_X$, not a truncation of that trial form.

For the exact finite solve $X_M=T_M^{-1}B_M$, embedded in the complete
tail, one has $K_{X_M}=S_M$ and
\begin{equation}\label{eq:v121-galerkin-defect}
 \boxed{S_\infty=S_M-r_M^*T^{-1}r_M,
 \qquad r_M=B-TX_M.}
\end{equation}
If the finite tail spaces are nested and their union is a form core
for $T$, then $S_M\downarrow S_\infty$ in matrix norm.
\end{proposition}
\begin{proof}
Expanding the first square gives $K_X-F+B^*T^{-1}B$.
Since $r=-T(X-T^{-1}B)$, its value is $r^*T^{-1}r$,
which proves both identities. In \eqref{eq:v121-finite-trial}
every trial lies in the head plus the retained finite tail space,
so compression changes none of the quadratic terms. Finite square completion gives
\[
 K_X=S_M+(X-T_M^{-1}B_M)^*T_M(X-T_M^{-1}B_M).
\]
Thus $K_X\succeq S_M$, and the graph map $h\mapsto(h,-Xh)$ is
injective. Substitution of the exact finite solve proves
\eqref{eq:v121-galerkin-defect}; its retained residual rows vanish,
but its omitted rows need not vanish.

For each head vector $h$, the variational identity is
\[
 \ip{B^*T^{-1}Bh}{h}
 =\sup_{y\in\mathcal D(T^{1/2})}
 \{2\Re\ip{Bh}{y}-\norm{T^{1/2}y}^2\}.
\]
Restricting the supremum to $\operatorname{Ran}P_M$ gives
$\ip{B_M^*T_M^{-1}B_Mh}{h}$. Nestedness and form-core density make
these quantities increase to the full supremum. Polarization and
finite dimensionality of the head give matrix-norm convergence of
the Schur complements in the stated decreasing order.
\end{proof}

A finite list of positive $S_M$ therefore does not prove positivity
of $S_\infty$. Positivity for every member of a cofinal sequence at
a fixed window would imply $S_\infty\succeq0$, but strict positivity
would still require a positive limiting margin. The structured inverse
bounds concern this same $T$; a comparison metric or a shifted lower
operator does not redefine the target Schur complement.

\begin{remark}[The semidefinite range condition]\label{rem:v121-range}
For $T\succeq0$, the notation $T^{-1}B$ requires an additional
condition. A form version is available if
$\operatorname{Ran}B\subset\operatorname{Ran}T^{1/2}$. Define the
reduced inverse $C=T^{-1/2}B$, with values orthogonal to $\ker T$.
For $X$ with range in the form domain, use
$(T^{1/2}X)^*T^{1/2}X$ in the definition of $K_X$. Then
\[
 K_X-(F-C^*C)=(T^{1/2}X-C)^*(T^{1/2}X-C)\succeq0.
\]
Orthogonality to $\ker T$, or membership in the closure of its range,
does not suffice when the range is not closed. At the certified
$\lambda=4$ split there is no such issue:
Proposition~\ref{prop:v120-lambda4-tail} gives
$T\succeq10^{-8}D\succeq1.7940\cdot10^{-8}I$ on $Q_{16}$.
\end{remark}

\begin{proposition}[A certified local zero head-error term]
\label{prop:v121-alpha-zero}
For the actual $\lambda=4$ Weil form with outer head $E_{16}$,
$K_X\succ0$ in both parity sectors for $X=0$ and for the frozen
dyadic solves supported on modes $17,\ldots,256$ specified in
Appendix~\ref{app:v121-audit}. Consequently $\alpha_4=0$ is
admissible in Proposition~\ref{prop:v120-structured-uniform} for
each of these choices, on the complete operator.
\end{proposition}
\begin{proof}
Assemble the exact finite entries with the separate logarithmic
diagonal, normalized parity basis and $L=2\log4$. Outward interval
$LDL^*$ factorization of $K_X$ verifies all $17$ even pivots and
all $16$ odd pivots strictly positive, both for $X=0$ and for the
specified dyadics. The identical witnesses pass at $768$ and $896$
bits; Appendix~\ref{app:v121-audit} gives the files and gates.
Real symmetric positivity also proves positivity on the complex
parity spaces. Equation~\eqref{eq:v121-finite-trial} identifies these
matrices with the complete trial forms. No infinite Schur sign or
small residual is used in this argument.
\end{proof}

This result does not certify $\alpha_\lambda=0$ at untested windows
or for larger unsupported solves. Positivity of an exact Schur
complement would guarantee $K_X\succeq0$ for every trial; without
that premise the least admissible $\alpha$ can depend on $X$.
For example, $F=0$, $B=T=1$ gives $S=-1$, $K_0=0$ and $K_1=-1$.
Even $K_0\succ0$ need not help the signed target: its residual is
the entire coupling $B$.

With a proved $\alpha_\lambda=0$ along a growing family, the precise
remaining sufficient estimates would be
\begin{equation}\label{eq:v121-scaled-target}
 e_\lambda\longrightarrow0,
 \qquad t_\lambda/\sqrt{\mu_\lambda}\longrightarrow0.
\end{equation}
The second condition is stronger than $t_\lambda\to0$ when
$\mu_\lambda\to0$. Geometric convergence of an inverse approximation
improves an enclosure of $r_\lambda^*T_\lambda^{-1}r_\lambda$;
it does not force that residual energy to vanish for a fixed trial.
These growing-window estimates remain unproved. Failure of this
sufficient bound would not itself refute Weil positivity.


\subsection*{An inverse-weighted finite spectral certificate}
The complement estimate can also identify a single candidate without
dividing its entire residual by the smallest spectral gap.
The following finite statement is independent of any PSWF approximation.

\begin{proposition}[Inverse-weighted finite certificate]
\label{prop:v112-schur-angle}
Let $W$ be a finite Hermitian matrix, $u$ a unit vector, and $P$ the
orthogonal projection onto $u^\perp$. Put
\[
 \alpha=\ip{Wu}{u},\quad r=PWu,\quad
 C=PWP|_{u^\perp},\quad A=C-\alpha I.
\]
If $A\succeq\gamma I$ for some $\gamma>0$, then $W$ has a simple
lowest eigenvalue $\lambda_0\le\alpha$. Set
$q=\norm{A^{-1}r}$ and $E=\ip{A^{-1}r}{r}$. For the angle $\theta$
between $u$ and the lowest eigenspace,
\begin{equation}\label{eq:v112-angle}
 \tan\theta\le q,\qquad
 \sin\theta\le\frac{q}{\sqrt{1+q^2}},\qquad
 \alpha-E\le\lambda_0\le\alpha.
\end{equation}
If $W$ commutes with reflection and $u$ is even, its ground state is even.
For any approximate solve $z\in u^\perp$ one also has
\begin{equation}\label{eq:v112-solve}
 q\le\norm z+\frac{\norm{r-Az}}{\gamma}.
\end{equation}
\end{proposition}
\begin{proof}
Min--max gives $\lambda_0\le\alpha$ and
$\lambda_1\ge\inf\sigma(C)>\alpha$. Write a ground vector as $bu+h$,
where $h\perp u$ and $b\ne0$. The complementary equation gives
$h=-b(C-\lambda_0 I)^{-1}r$. With $t=\alpha-\lambda_0\ge0$, the
scalar equation is $t=\ip{(A+tI)^{-1}r}{r}$. Spectral calculus for
the same positive operator $A$ gives $t\le E$ and
$\norm{(A+tI)^{-1}r}\le\norm{A^{-1}r}$. These prove
\eqref{eq:v112-angle}. An odd ground state would be orthogonal to the
even $u$, contradicting the complement bound. Finally
$A^{-1}r=z+A^{-1}(r-Az)$ proves \eqref{eq:v112-solve}.
\end{proof}

For certified approximations $\widetilde A,\widetilde r$, suppose
$\norm{A-\widetilde A}\le\delta_A$,
$\norm{r-\widetilde r}\le\delta_r$ and
$\widetilde A\succeq\widetilde\gamma I$ with
$\widetilde\gamma>\delta_A$. Then the usable error budget is
\[
 q\le\norm z+
 \frac{\delta_r+\delta_A\norm z+
       \norm{\widetilde r-\widetilde A z}}
      {\widetilde\gamma-\delta_A}.
\]
For a fixed exact candidate and matrix error $\norm{W-\widetilde W}\le\eta$,
one may take $\delta_r=\eta$ and $\delta_A=2\eta$.
True-source approximation error is a separate input. The condition
$\gamma>0$ is sufficient, not necessary, for source-to-ground
convergence: in $W=\operatorname{diag}(0,d,1)$, the candidate
$u=\sqrt{1-\tau^2}e_0+\tau e_2$ has angle tending to zero but
$\gamma\le d-\tau^2<0$ when $0<d<\tau^2$.

Appendix~\ref{app:v112-certificate} certifies the hypotheses for one
explicit rational candidate at $\lambda=3$, $N=64$, using the actual
arithmetic Weil matrix. That certificate alone proves no uniform complement theorem or
exact-source error bound. The following argument supplies the latter
at its fixed parameters; endpoint and graph transfer remain separate.


\subsection*{From a finite prolate approximation to the exact source}
The preceding certificate treats its input vector as exact. To identify
that vector with a projection of the manuscript's source, the omitted
Legendre tail must also be controlled. Here this can be done at fixed
parameters, including pointwise values rather than only an $L^2$ error.

\begin{lemma}[Certified angular eigenfunction error]
\label{lem:v113-angular}
On $L^2(-1,1)$ define
\[
 e_\ell(x)=\sqrt{(2\ell+1)/2}\,P_\ell(x),\qquad
 De_\ell=\ell(\ell+1)e_\ell,
\]
and let
$\mathsf J_c=D+c^2X^2$, where $Xf(x)=xf(x)$ and
$\mathcal D(\mathsf J_c)=\mathcal D(D)$.
Restrict to the even subspace. Let $J_M$ be its compression to
$e_0,e_2,\ldots,e_{2M-2}$, let $b_M$ be the link from its last
coordinate to $e_{2M}$, and put $d_M=2M(2M+1)$.
For $x<d_M$, the number of eigenvalues of $\mathsf J_c$ below $x$
is between the negative inertia counts of
\begin{equation}\label{eq:v113-inertia}
 J_M-xI,\qquad
 J_M-xI-\frac{b_M^2}{d_M-x}E_{M-1,M-1}.
\end{equation}
If both matrices are nonsingular with the same count, that count is exact.

Suppose these counts certify exactly one eigenvalue $\nu$ in
$(\xi-g,\xi+g)$ and no other spectrum there, with $g>0$.
Let $\xi\in\mathbb R$, let $v$ be an exact normalized real vector in
this even finite span, and suppose
$\norm{(\mathsf J_c-\xi)v}\le\epsilon<g$.
For the normalized eigenfunction $\psi$ at $\nu$, with sign chosen
so that $\ip v\psi>0$, put $\delta=\sqrt2\epsilon/g$.
Then
\begin{align}
 |\nu-\xi|&\le\epsilon,& \norm{v-\psi}_2&\le\delta,\notag\\
 \norm{v-\psi}_\infty
 &\le 2\epsilon+(1+|\xi|+c^2)\delta.
 \label{eq:v113-pointwise}
\end{align}
The residual includes the first omitted coordinate $b_Mv_{M-1}$.
The gap $g$ is measured from $\xi$ to the other exact eigenvalues.
\end{lemma}
\begin{proof}
The diagonal operator $D$ is self-adjoint with compact resolvent;
$c^2X^2$ is bounded, positive and self-adjoint. Its angular realization
is the regular PSWF operator (the differential spectral parameter is
shifted by $c^2$ relative to the convention in \cite{DLMFPSWF}).
The omitted compression $T$ satisfies $T\succeq d_MI$. For $x<d_M$
its positive inverse gives the finite Schur complement
\[
 J_M-xI-b_M^2\ip{(T-xI)^{-1}e_{2M}}{e_{2M}}E_{M-1,M-1}.
\]
The scalar lies between $0$ and $(d_M-x)^{-1}$. Inertia and min--max
therefore prove \eqref{eq:v113-inertia}. Completing the square in the
quadratic form justifies the Schur argument for the unbounded $T$.

The spectral theorem places an eigenvalue within $\epsilon$ of $\xi$;
the certified interval identifies it as $\nu$. All other coefficients
of $v$ have squared mass at most $(\epsilon/g)^2$. Choosing positive
overlap gives $\norm{v-\psi}_2\le\sqrt2\epsilon/g$.
For $e=v-\psi$,
\[
 \mathsf J_ce=\xi e+(\xi-\nu)\psi+(\mathsf J_c-\xi)v,
\]
so $\norm{(I+D)e}_2\le2\epsilon+(1+|\xi|+c^2)\delta$.
Finally $|P_\ell(x)|\le1$ on $[-1,1]$ and
\[
 \sum_{\ell\ge0}\frac{(2\ell+1)/2}{[1+\ell(\ell+1)]^2}
 \le\frac12+\frac12\sum_{\ell\ge1}
             \left(\frac1{\ell^2}-\frac1{(\ell+1)^2}\right)=1.
\]
Cauchy--Schwarz gives uniform convergence of the Legendre expansion and
$\norm e_\infty\le\norm{(I+D)e}_2$, including both endpoint traces.
\end{proof}

Here is the source error budget used below. Set $\lambda=3$ and let
$\widehat h_j=\lambda^{-1/2}v_j(\cdot/\lambda)$ approximate
$h_{j,\lambda}$ for $j=0,4$. Let $d_j$ and $u_j$ bound their physical
$L^2$ and uniform errors. Define
\[
 \widehat a_0=-\int\widehat h_4,\quad
 \widehat a_4=\int\widehat h_0,\quad
 \widehat h^\circ=\widehat a_0\widehat h_0+
                  \widehat a_4\widehat h_4.
\]
Both this polynomial and the exact $h^\circ$ have zero integral.
With $d_{a_0}=\sqrt{2\lambda}\,d_4$ and
$d_{a_4}=\sqrt{2\lambda}\,d_0$, a valid uniform error is
\begin{equation}\label{eq:v113-source-budget}
 \eta=\sum_{j=0,4}
 \{(|\widehat a_j|+d_{a_j})u_j+
                d_{a_j}\norm{\widehat h_j}_\infty\}.
\end{equation}
Use the fixed repair with $b=3/4$,
$\phi(t)=\exp(1-1/(1-(t/b)^2))$ for $|t|<b$, zero outside, and
$\psi(t)=\phi(t)(1-Ct^2)$ where
$C=\int_0^b\phi/\int_0^b t^2\phi$.
Since $\phi\le1$ everywhere and $\phi\ge1/2$ on $[b/4,b/2]$,
$Cb^2\le128$ and $\norm\psi_\infty\le129$.
Consequently, with $Z=|\widehat h^\circ(0)|$,
\begin{equation}\label{eq:v113-repair-budget}
 \norm{\widetilde h-\widehat h^\circ}_\infty
 \le U:=\eta+129(Z+\eta).
\end{equation}
This bounds the entire exact repair without assuming a numerical
quadrature value for $C$.

Let $\widehat p$ be the inversion-even co-Poisson proxy of
$\widehat h^\circ$ on the same window. At almost every
$u\in[\lambda^{-1},\lambda]$, at most $\lambda/u$ summands occur,
so the physical factor $\sqrt u$ gives
$|\mathcal E(\widetilde h-\widehat h^\circ)(u)|
 \le\lambda^{3/2}U$.
Inversion averaging is an $L^2(du/u)$ contraction. Bessel's inequality
and the strict interior endpoint formula therefore give
\begin{align}
 \norm{P_Np_3-P_N\widehat p}_2&\le\sqrt{\lambda^3L}\,U,
       &L&=2\log3,\label{eq:v113-projection-budget}\\
 |B_3-\widehat B|&\le\tfrac12\left(\sqrt\lambda\,\eta+
                                     8\lambda^{-1/2}U\right),
 \label{eq:v113-endpoint-budget}
\end{align}
where
$\widehat B=\tfrac12\{\sqrt\lambda\,\widehat h^\circ(\lambda)
+\lambda^{-1/2}\sum_{m=1}^{8}\widehat h^\circ(m/\lambda)\}$.
The $m=9$ term is excluded from this lower interior trace; the repair
vanishes at the upper endpoint. These are fixed-parameter estimates,
not a growing-window spectral theorem.

\begin{proposition}[Exact-source angle and endpoint at one window]
\label{prop:v113-source-certificate}
Use the exact repaired source of Proposition~\ref{prop:v12-source}
with the bump just specified. At $\lambda=3,N=64$, let $c^*$ be the
unnormalized rational coefficient vector in
Proposition~\ref{prop:v112-finite-certificate}, and let $v^*$ be the
coefficient vector of the exact $P_{64}p_3$ in the same translated
Fourier basis. The interval certificate in
Appendix~\ref{app:v113-source} proves
\[
 \norm{v^*-c^*}_2<2.60\,10^{-36},\qquad
 \sin\angle(v^*,\text{ground state of }W_{3,64})<0.000216.
\]
It also proves
\begin{equation}\label{eq:v113-endpoint-failure}
 5.58669\,10^{-19}<B_3<5.58670\,10^{-19},\qquad
 1.374<\frac{|\delta_{64}(P_{64}p_3)-B_3|}{B_3}<1.375.
\end{equation}
Thus the small ordinary ground-state angle coexists with a relative
physical-endpoint error greater than one for this exact projected source.
\end{proposition}
\begin{proof}
The appendix certifies all error budgets above, including the full
infinite angular residuals, and the exact polynomial Fourier moments.
It gives the first bound and the displayed endpoint intervals.
For nonzero vectors $v^*,c^*$ with error at most $e<\norm{c^*}$,
\[
 \norm{P_{\mathbb Cv^*}-P_{\mathbb Cc^*}}
 =\sin\angle(v^*,c^*)\le\frac{e}{\norm{c^*}-e}.
\]
This is below $4\,10^{-36}$ here. The triangle inequality for rank-one
orthogonal projections and the sharper enclosed sine bound from
Proposition~\ref{prop:v112-finite-certificate} prove the angle claim.
No new positive-complement estimate for the perturbed candidate is
assumed. Equation~\eqref{eq:v113-endpoint-failure} concerns physical
endpoint evaluation, not the equal-coefficient functional without its
$L^{-1/2}$ factor.
\end{proof}

\section{Weighted G1 transfer and the arithmetic jump train}\label{sec:weighted}
\begin{proposition}[Cofinal finite G1 bound on a weighted sampler class]\label{prop:g1-weighted-sampler}
Use the explicit repaired source of Section~\ref{sec:g1}.
Let $R>3$, let $K\ge2N$, and write
\[
 x=\sum_{|m|\le K}x_mU_m\in E_K,
 \qquad
 \|x\|_{\mathcal H_{R,L}}
 :=\sup_{|m|\le K}
 \sqrt L\left(1+\frac{|m|}{L}\right)^R|x_m|.
\]
If $N\ge2L$, then the filtered proxy of
Proposition~\ref{prop:g1-vp-transfer} satisfies
\begin{equation}\label{eq:g1-weighted-sampler}
 \boxed{
 \frac{|\langle D_{\log}W_{\lambda,K}
 k^{\mathrm{VP}}_{\lambda,N},x\rangle|}{|B_\lambda|}
 \le C_R\|x\|_{\mathcal H_{R,L}}
 \left[
 \mathfrak r(\lambda)
 +\frac{\lambda\mathcal V_\lambda}{|B_\lambda|}
 \left\{
 \frac{L}{N}
 +\left(\frac{L}{N}\right)^{R-1}(1+\log(2N))
 \right\}
 \right].}
\end{equation}
The bound is uniform in the ambient cutoff $K$.  Thus, once the arithmetic
core cutoff $N$ has been fixed, enlarging the Hardy-sampler cutoff cannot
spoil this G1 estimate.
\end{proposition}
\begin{proof}
Put $e_N=k^{\mathrm{VP}}_{\lambda,N}-p_\lambda$.  Because the
Vallee--Poussin multiplier is one for $|n|\le N$, one has
$\widehat e_N(n)=0$ there, while for $|n|>N$ the BV estimate gives
\[
 |\widehat e_N(n)|
 \le |\widehat p_\lambda(n)|
 \le \frac{\sqrt L\,\mathcal V_\lambda}{2\pi|n|}.
\]
The coefficient hypothesis and $R>3$ imply, by comparison with Riemann sums,
\[
 \|D_{\log}x\|_2+\|D_{\log}^2x\|_2
 +|(D_{\log}x)(-a)|+|(D_{\log}x)(a)|
 \le C_R\|x\|_{\mathcal H_{R,L}}.
\]
Indeed the two $L^2$ sums reduce to
$L^{-1}\sum_m (|m|/L)^{2j}(1+|m|/L)^{-2R}$, $j=1,2$, and each endpoint
trace to $L^{-1}\sum_m (|m|/L)(1+|m|/L)^{-R}$.
Proposition~\ref{prop:g1-sobolev}, applied to $D_{\log}x$, therefore gives
\[
 |QW_\lambda(p_\lambda,D_{\log}x)|
 \le C_R|B_\lambda|\mathfrak r(\lambda)
       \|x\|_{\mathcal H_{R,L}}.
\]

For the projection error separate the off-diagonal and diagonal matrix
entries.  Lemma~\ref{lem:g1-high-low} gives, uniformly in $K$,
\[
 \begin{split}
 |QW_\lambda(e_N,D_{\log}x)|_{\mathrm{off}}
 \le{}& C_R\lambda\mathcal V_\lambda
 \|x\|_{\mathcal H_{R,L}}
 \sum_{|n|>N}\frac1{|n|}
 \sum_{m\ne n}\frac{a_m}{|n-m|},\\
 a_m:={}&\frac{|m|}{L}
       \left(1+\frac{|m|}{L}\right)^{-R}.
 \end{split}
\]
For nonzero integer $n-m$ one may replace $|n-m|^{-1}$ by
$2(1+|n-m|)^{-1}$.  For $|n|>N\ge2L$, split the inner sum into
$|m|\le|n|/2$, $|n|/2<|m|\le2|n|$, and $|m|>2|n|$.
Since $\sum_m a_m\le C_RL$, the first range is $O_R(L/|n|)$.
On the middle range,
$a_m\le C_R(|n|/L)^{1-R}$ and the harmonic sum contributes
$O(1+\log(2|n|))$.  On the last range the denominator is $\gg|m|$, and
summation of $m^{-R}$ gives $O_R((|n|/L)^{1-R})$.  Hence
\[
 \sum_{m\ne n}\frac{a_m}{|n-m|}
 \le C_R\left[
 \frac{L}{|n|}
 +\left(\frac{|n|}{L}\right)^{1-R}
  (1+\log(2|n|))\right],
\]
and therefore
\begin{equation}\label{eq:g1-weighted-offdiag}
 |QW_\lambda(e_N,D_{\log}x)|_{\mathrm{off}}
 \le C_R\lambda\mathcal V_\lambda\|x\|_{\mathcal H_{R,L}}
 \left[
 \frac{L}{N}
 +\left(\frac{L}{N}\right)^{R-1}(1+\log(2N))\right].
\end{equation}

The diagonal must be handled separately.  Corollary~\ref{cor:g1-complete-matrix},
the BV coefficient estimate, and the definition of $\mathcal H_{R,L}$ give
\[
 \begin{split}
 |QW_\lambda(e_N,D_{\log}x)|_{\mathrm{diag}}
 \le{}& C\frac{\mathcal V_\lambda}{L}
 \|x\|_{\mathcal H_{R,L}}\\
 &\times\sum_{N<|n|\le K}
 \bigl(\lambda+1+\log(2+|n|/L)\bigr)
 \left(1+\frac{|n|}{L}\right)^{-R}.
 \end{split}
\]
Comparison with the corresponding decreasing integral, using $N\ge2L$ and
$R>3$, yields
\begin{equation}\label{eq:g1-weighted-diag}
 |QW_\lambda(e_N,D_{\log}x)|_{\mathrm{diag}}
 \le C_R\mathcal V_\lambda\|x\|_{\mathcal H_{R,L}}
 \bigl(\lambda+1+\log(2N)\bigr)
 \left(\frac{L}{N}\right)^{R-1}.
\end{equation}
Because $\lambda\ge2$, \eqref{eq:g1-weighted-diag} is absorbed by the
second term on the right of \eqref{eq:g1-weighted-offdiag}.  Thus the
projection-error contribution has exactly the rate stated in
\eqref{eq:g1-weighted-sampler}, despite the logarithmic diagonal growth.
Finally
$\langle D_{\log}W_{\lambda,K}k^{\mathrm{VP}},x\rangle
=QW_\lambda(k^{\mathrm{VP}},D_{\log}x)$ because both finite vectors lie in
$E_K$.  Combining the two pieces proves \eqref{eq:g1-weighted-sampler}.
\end{proof}

\subsection*{Exact extraction of the arithmetic jump train}
The BV estimate above treats all variation of the prolate/co-Poisson proxy in
one lump.  That is unnecessarily expensive after normalization by the tiny
endpoint amplitude.  The actual discontinuities have a much more rigid
arithmetic form and can be separated exactly.

Let $y=\log(\lambda u)\in[0,L]$ and let $B_\lambda^{+}
=\sqrt\lambda\,\widetilde h_\lambda(\lambda^-)$ denote the upper source
boundary amplitude.  Put
\[
 \mathcal M_\lambda:=\{m\in\mathbb N:1<m<\lambda^2\}.
\]
The fixed smooth exact-moment repair used above is supported strictly inside
the source interval, so it creates no endpoint jumps.

\begin{lemma}[Exact arithmetic jump train]\label{lem:g1-jump-train}
There is a zero-mean periodic BV function $\jmath_\lambda$ on the logarithmic
circle such that its distributional derivative is
\begin{equation}\label{eq:g1-jump-measure}
 d\jmath_\lambda
 =\frac{B_\lambda^{+}}2
 \sum_{m\in\mathcal M_\lambda}m^{-1/2}
 \bigl(\delta_{\log m}-\delta_{L-\log m}\bigr).
\end{equation}
Moreover
\[
 c_\lambda:=p_\lambda-\jmath_\lambda
\]
is continuous and periodic, is piecewise smooth, and its first logarithmic
derivative has bounded variation.  If
\[
 t_n:=\frac{2\pi n}{L},\qquad
 \mathcal D_\lambda(t):=
 \sum_{m\in\mathcal M_\lambda}m^{-1/2}\sin(t\log m),
\]
then for every $n\ne0$,
\begin{equation}\label{eq:g1-jump-fourier-exact}
 \boxed{
 \widehat{\jmath_\lambda}(n)
 =-\frac{B_\lambda^{+}}{t_n\sqrt L}\,
 \mathcal D_\lambda(t_n).}
\end{equation}
Consequently, for all sufficiently large $\lambda$,
\begin{equation}\label{eq:g1-jump-fourier-trivial}
 \boxed{
 |\widehat{\jmath_\lambda}(n)|
 \le C|B_\lambda|\frac{\lambda\sqrt L}{|n|}.}
\end{equation}
Finally, if
\[
 \mathcal W_\lambda
 :=\bigl\|D_y^2c_\lambda\bigr\|_{\mathcal M(\mathbb T_L)}
\]
denotes the total variation of the distributional second derivative of the
continuous remainder, then
\begin{equation}\label{eq:g1-continuous-fourier}
 \boxed{
 |\widehat{c_\lambda}(n)|
 \le \frac{L^{3/2}\mathcal W_\lambda}{4\pi^2n^2},
 \qquad n\ne0.}
\end{equation}
All quantities in the statement are determined by the explicit source and
$\lambda$; no zeta zero is used.
\end{lemma}
\begin{proof}
For
$f_\lambda(u)=\mathcal E(\widetilde h_\lambda)(u)$, the summand with index
$m$ disappears as $u$ increases through $u=\lambda/m$.  Its one-sided jump is
therefore
\[
 -\left(\frac\lambda m\right)^{1/2}
  \widetilde h_\lambda(\lambda^-)
 =-\frac{B_\lambda^{+}}{\sqrt m}.
\]
In the $y$ coordinate this occurs at $L-\log m$.  The inverted function
$f_\lambda(u^{-1})$ has the opposite jump at $\log m$.  Since
$p_\lambda=\frac12(f_\lambda+Jf_\lambda)$ on the window, its complete
interior jump measure is exactly \eqref{eq:g1-jump-measure}.  The total mass
of that measure is zero, so it is the derivative of a periodic BV function;
fixing zero mean makes $\jmath_\lambda$ unique.  Subtracting it removes every
function jump.  The arithmetic sum has only finitely many thresholds and is
smooth between them, so $c_\lambda$ is continuous, piecewise smooth, and
$D_yc_\lambda$ is BV.

For $n\ne0$, periodic Stieltjes integration by parts gives
\[
 \widehat{\jmath_\lambda}(n)
 =\frac1{it_n\sqrt L}
 \int_0^Le^{-it_ny}\,d\jmath_\lambda(y).
\]
Because $e^{-it_nL}=1$, \eqref{eq:g1-jump-measure} yields
\[
 \int_0^Le^{-it_ny}\,d\jmath_\lambda(y)
 =\frac{B_\lambda^{+}}2
 \sum_{m\in\mathcal M_\lambda}m^{-1/2}
 \{m^{-it_n}-m^{it_n}\}
 =-iB_\lambda^{+}\mathcal D_\lambda(t_n),
\]
which is \eqref{eq:g1-jump-fourier-exact}.  The elementary estimate
\[
 |\mathcal D_\lambda(t)|
 \le\sum_{m<\lambda^2}m^{-1/2}\le2\lambda
\]
and $B_\lambda^{+}\asymp B_\lambda$ from
Corollary~\ref{cor:pswf-cross-tail} prove
\eqref{eq:g1-jump-fourier-trivial}.

For the continuous remainder there is no first-order boundary term.  A second
periodic Stieltjes integration by parts, now against the finite measure
$D_y^2c_\lambda$, gives
\[
 |\widehat{c_\lambda}(n)|
 \le\frac{\mathcal W_\lambda}{\sqrt L\,t_n^2}
 =\frac{L^{3/2}\mathcal W_\lambda}{4\pi^2n^2},
\]
proving \eqref{eq:g1-continuous-fourier}.
\end{proof}

The exact sine sum in \eqref{eq:g1-jump-fourier-exact} is potentially useful
in its own right: any unconditional exponential-sum improvement immediately
sharpens the cutoff below.  Even the trivial bound already removes the
factor $\mathcal V_\lambda/|B_\lambda|$ from the discontinuous arithmetic
part.

\begin{proposition}[Jump/continuous split for the weighted G1 transfer]\label{prop:g1-jump-split}
Let $R>3$, $K\ge2N$, $N\ge2L$, and let $x\in E_K$ be measured by
$\|x\|_{\mathcal H_{R,L}}$ as in
Proposition~\ref{prop:g1-weighted-sampler}.  Then
\begin{align}
 &\frac{
 |QW_\lambda((\mathsf V_N-I)p_\lambda,D_{\log}x)|
 }{|B_\lambda|}
 \notag\\
 &\qquad\le C_R\|x\|_{\mathcal H_{R,L}}
 \Bigg[
 \lambda^2\left\{
 \frac LN+
 \left(\frac LN\right)^{R-1}(1+\log(2N))
 \right\}
 \label{eq:g1-jump-split}\\
 &\hspace{43mm}
 +\frac{\lambda\mathcal W_\lambda}{|B_\lambda|}
 \left\{
 \left(\frac LN\right)^2+
 \left(\frac LN\right)^R(1+\log(2N))
 \right\}
 \Bigg].\notag
\end{align}
The first bracket is the entire contribution of the arithmetic function jumps
and has no small-boundary denominator.  The only surviving endpoint-normalized
regularity cost belongs to the continuous remainder and gains one additional
power of $N^{-1}$.
\end{proposition}
\begin{proof}
Write
$p_\lambda=\jmath_\lambda+c_\lambda$ and use that the
Vallee--Poussin multiplier is exactly one on $|n|\le N$.  For the jump part,
\eqref{eq:g1-jump-fourier-trivial} replaces the BV coefficient estimate used
in Proposition~\ref{prop:g1-weighted-sampler} by
\[
 |\widehat{(\mathsf V_N-I)\jmath_\lambda}(n)|
 \le C|B_\lambda|\frac{\lambda\sqrt L}{|n|},
 \qquad |n|>N.
\]
Repeating the off-diagonal convolution estimate there therefore replaces
$\lambda\mathcal V_\lambda/|B_\lambda|$ by $C\lambda^2$.  The diagonal
bound in Corollary~\ref{cor:g1-complete-matrix} is absorbed by the same
second term because $\lambda\ge2$.

For the continuous remainder use \eqref{eq:g1-continuous-fourier}.  With
\[
 a_m=\frac{|m|}{L}\left(1+\frac{|m|}{L}\right)^{-R},
\]
the off-diagonal part is bounded by
\[
 C_R\lambda L\mathcal W_\lambda
 \|x\|_{\mathcal H_{R,L}}
 \sum_{|n|>N}\frac1{n^2}
 \left[
 \frac L{|n|}+
 \left(\frac{|n|}{L}\right)^{1-R}
 (1+\log(2|n|))
 \right].
\]
Comparison with decreasing integrals gives
\[
 C_R\lambda\mathcal W_\lambda
 \|x\|_{\mathcal H_{R,L}}
 \left[
 \left(\frac LN\right)^2+
 \left(\frac LN\right)^R(1+\log(2N))
 \right].
\]
Using the diagonal estimate from
Corollary~\ref{cor:g1-complete-matrix} with the $n^{-2}$ coefficient bound
produces no larger term.  Division by $|B_\lambda|$ proves
\eqref{eq:g1-jump-split}.
\end{proof}

\begin{corollary}[The arithmetic jumps are resolvable at near-Slepian scale]\label{cor:g1-jump-natural-scale}
Fix $R>3$ and put
\begin{equation}\label{eq:g1-jump-cutoff}
 N_\lambda^{\rm jump}
 :=\left\lceil
 C_R\frac{\lambda^2L}{\mathfrak r(\lambda)}
 \right\rceil.
\end{equation}
For $C_R$ sufficiently large, the jump contribution in
\eqref{eq:g1-jump-split} is $O_R(\mathfrak r(\lambda))$ uniformly for every
ambient cutoff $K\ge2N_\lambda^{\rm jump}$.  Moreover
\begin{equation}\label{eq:g1-jump-near-slepian}
 \boxed{
 N_\lambda^{\rm jump}=\lambda^{2+o(1)}.}
\end{equation}
At the same cutoff, the continuous-remainder contribution is
\begin{equation}\label{eq:g1-continuous-at-jump-scale}
 O_R\!\left(
 \frac{\mathcal W_\lambda}{|B_\lambda|}
 \frac{\mathfrak r(\lambda)^2}{\lambda^3}
 \right)
\end{equation}
(up to a smaller $R$-dependent logarithmic term).  Thus the concrete
zero-independent condition
\begin{equation}\label{eq:g1-continuous-subgate}
 \boxed{
 \frac{\mathcal W_\lambda}{|B_\lambda|}
 \ll_R\frac{\lambda^3}{\mathfrak r(\lambda)} }
\end{equation}
would upgrade the filtered G1 transfer to the natural
$\lambda^{2+o(1)}$ Fourier scale.  Condition
\eqref{eq:g1-continuous-subgate} is \emph{not} proved here; it is a newly
isolated regularity subgate rather than an RH assumption.
\end{corollary}
\begin{proof}
The first term in \eqref{eq:g1-jump-split} is made
$O_R(\mathfrak r(\lambda))$ by \eqref{eq:g1-jump-cutoff}; the second
high-order term is smaller because $R>3$ and $\mathfrak r(\lambda)<1$ for
large $\lambda$.  Since
\[
 \log\frac1{\mathfrak r(\lambda)}=O(\sqrt L+\log L)=o(L)
\]
while $\lambda^2=e^L$, one has
$\log N_\lambda^{\rm jump}=L+o(L)$, proving
\eqref{eq:g1-jump-near-slepian}.  Finally
$L/N_\lambda^{\rm jump}\asymp
\mathfrak r(\lambda)/\lambda^2$; substitution in the continuous term of
\eqref{eq:g1-jump-split} gives
\eqref{eq:g1-continuous-at-jump-scale}.  Requiring it to be
$O_R(\mathfrak r(\lambda))$ gives
\eqref{eq:g1-continuous-subgate}.
\end{proof}


\subsection*{A polynomial cutoff for the literal source}
The small physical endpoint need not force an exponentially large Fourier
cutoff for this fixed two-mode source.  The additional ingredient is an
exterior norm estimate for a finite entire extension of the arithmetic
sum.  It controls the endpoint derivative jets without bounding the whole
interior variation by the endpoint amplitude.

\begin{lemma}[Polynomial approximation with a growing derivative order]
\label{lem:v114-entire-polynomial}
Fix a compact real interval $I$ of positive length. Suppose
\[
 F(z)=\int_{-\Omega}^{\Omega}e^{izt}\,d\mu(t),\qquad
 \Omega\le A_0c,\quad \norm\mu\le c^{A_1}.
\]
For each fixed $H>0$, there is $A>0$ such that a polynomial $T$ of degree
$D=\lceil Ac\rceil$ satisfies, simultaneously for $0\le j\le\lceil c\rceil$,
\begin{equation}\label{eq:v114-polynomial-error}
 \norm{(F-T)^{(j)}}_{L^\infty(I)}
       \le e^{-Hc}(Cc)^j.
\end{equation}
For all sufficiently large $c$, the constants are independent of $c$ and $j$. Moreover,
\[
 \norm T_\infty\le C(D+1)\norm T_2,\qquad
 \norm{T^{(j)}}_\infty\le(CD^2)^j\norm T_\infty
\]
on $I$.
\end{lemma}
\begin{proof}
Take the degree-$D$ Taylor polynomial about the midpoint of $I$ under the
finite measure. If $d$ is the half-length of $I$, its differentiated
remainder is at most
\[
 \norm\mu\,\Omega^j e^{\Omega d}
       \frac{(\Omega d)^{D+1-j}}{(D+1-j)!}.
\]
The inequality $k!\ge(k/e)^k$, with $A$ sufficiently large, proves
\eqref{eq:v114-polynomial-error} for the entire stated range of $j$.
The first polynomial inequality follows from normalized Legendre expansion,
$|P_k|\le1$, and $\sum_{k=0}^D(2k+1)=(D+1)^2$. The second is Markov's
inequality iterated on the derivatives, with the fixed interval scaling.
Real and imaginary parts give the same bounds for complex polynomials.
\end{proof}

\begin{lemma}[Endpoint jets of the finite arithmetic sum]
\label{lem:v114-endpoint-jets}
Put $c=2\pi\lambda^2$, $r=\lceil c\rceil$, and let $p_\lambda^\circ$
be the inversion-even proxy before the fixed value-at-zero repair.
If $J_j$ is the sum of absolute jumps of its $j$th ordinary logarithmic
derivative, including the jump at the identified periodic endpoint, then
\begin{equation}\label{eq:v114-jet-bound}
 J_j\le C|B_\lambda^+|\lambda^3(C\lambda^4)^j,
          \qquad 1\le j<r.
\end{equation}
Its zeroth-order periodic boundary jump is zero. On its smooth pieces,
for fixed constants $A,C,\rho>0$,
\begin{equation}\label{eq:v114-smooth-remainder}
 \norm{(p_\lambda^\circ)^{(r)}_{\rm ac}}_1
       \le CL\lambda^3 r!\rho^{-r}e^{Ac}.
\end{equation}
\end{lemma}
\begin{proof}
Write
\[
 g(x)=\sqrt\lambda h_\lambda^\circ(\lambda x)
     =\sum_{j=0,4}a_{j,\lambda}\psi_{j,c}(x),\qquad g(1)=B_\lambda^+.
\]
The exact finite Fourier eigenrelation used in
Lemma~\ref{lem:v12-endpoint-energy} represents $g$ by a measure supported
on $[-c,c]$ with total variation $O(\lambda)$; this is also the
bandlimited extension in \cite[\S30.15]{DLMFPSWF}.
Proposition~\ref{prop:endpoint-radial} and endpoint noncancellation give
$|g(x)|\le C|B_\lambda^+|$ for $x\ge1$.

Set $M=\lceil2\lambda^2\rceil$ and
$G(v)=\sum_{m=1}^M g(mv/\lambda^2)$. Its measure has bandwidth
$2\pi M=O(c)$ and total variation $O(\lambda^3)$. For $1/2\le v\le1$,
\[
 G(v)=\lambda v^{-1/2}f_\lambda^\circ(v/\lambda)
       +\sum_{\substack{m\le M\\mv\ge\lambda^2}}
                       g(mv/\lambda^2)
\]
almost everywhere. The last sum is $O(|B_\lambda^+|\lambda^2)$.
Corollary~\ref{cor:pswf-cross-tail}, which includes the unrepaired source,
and $dv/v=du/u$ therefore give
\[
 \norm G_{L^2(1/2,1)}\le C|B_\lambda^+|\lambda^2.
\]
Since $|B_\lambda^+|\asymp c^{11/4}e^{-c}$, apply
Lemma~\ref{lem:v114-entire-polynomial} with fixed sufficiently large $H$.
Its polynomial inequalities imply
\begin{align*}
 |G^{(j)}(1)|&\le C|B_\lambda^+|\lambda^4(C\lambda^4)^j,\\
 \sup_{1\le x\le3}|g^{(j)}(x)|
       &\le C|B_\lambda^+|(C\lambda^4)^j,
          \qquad 0\le j\le r.
\end{align*}
For the second bound approximate $g$ on $[1,3]$ and use its already
bounded sup norm, so the extra Legendre degree factor is unnecessary.

At the strict interior lower trace $v=1+$, subtract from $G$ precisely the
terms with $m\ge\lambda^2$, including equality when $\lambda^2$ is an
integer, and multiply by $\sqrt v/\lambda$. Their arguments lie in
$[1,3]$, their number is $O(\lambda^2)$, and differentiation costs at most
$3^j$. Thus the logarithmic trace jets are bounded by
$C|B_\lambda^+|\lambda^3(C\lambda^4)^j$.
To justify the conversion uniformly in $j$, use
$(v\partial_v)^j=\sum_k S(j,k)v^k\partial_v^k$ and
$S(j,k)\le\binom jk j^{j-k}$; hence
$\sum_k S(j,k)A^k\le(A+j)^j$. Here $j=O(\lambda^2)$ and
$A=C\lambda^4$. For the product factor, $|(\sqrt v)^{(k)}|_{v=1}\le k!$ and
$\sum_{k=0}^j\binom jk k!A^{j-k}\le2A^j$ when $A\ge2j$;
thus its cost is uniform in $j$ as well.
At the upper trace only the first summand survives, giving the stronger
bound without $\lambda^3$.

Each interior derivative jump is a source endpoint jet times
$m^{-1/2}$, with the reflected sign. Summing $m<\lambda^2$ costs
$O(\lambda)$; the preceding trace estimates control the periodic jump.
This proves \eqref{eq:v114-jet-bound}. Inversion symmetry gives the zero
function jump at the periodic boundary.
Finally extend each smooth piece using its fixed finite collection of
entire summands. On a complex logarithmic disk of fixed small radius
$\rho$, every active source argument has imaginary part bounded by a
fixed constant. The Fourier representation bounds the extension by
$C\lambda^3e^{Ac}$. Cauchy's estimate, integrated over total length $L$,
proves \eqref{eq:v114-smooth-remainder}. No asymptotic remainder was
differentiated, and the smooth compact repair is not used in this step.
\end{proof}

\begin{proposition}[Polynomial physical-endpoint recovery]
\label{prop:v114-polynomial-endpoint}
For the source of Proposition~\ref{prop:v12-source}, repaired as in
Corollary~\ref{cor:v12-repair}, set
\[
 N_\lambda=\left\lceil\lambda^8(1+L)\right\rceil,
       \qquad k_\lambda=P_{N_\lambda}p_\lambda.
\]
For all sufficiently large $\lambda$,
\begin{align}
 \frac{|\delta_{N_\lambda}(k_\lambda)-B_\lambda|}{|B_\lambda|}
       &\le\frac C\lambda,\label{eq:v114-polynomial-endpoint}\\
 |\widehat p_\lambda(n)|
       &\le C|B_\lambda|\frac{\lambda\sqrt L}{|n|},
            \qquad |n|>N_\lambda.\label{eq:v114-high-tail}
\end{align}
This is an asymptotic estimate with no certified numerical threshold or
optimized exponent. The vector is the literal ordinary Fourier projection;
no endpoint shell is added.
\end{proposition}
\begin{proof}
First omit the compact repair. The periodic proxy is BV and continuous
at the identified endpoint, so the Dirichlet--Jordan theorem identifies
its symmetric Fourier sum there with its common one-sided trace. We may
therefore estimate the actual endpoint error by its symmetric Fourier tail.
In the jump train of
Lemma~\ref{lem:g1-jump-train}, the endpoint tail contributed by a threshold
$m$ is, up to sign,
\[
 \frac{B_\lambda^+}{\pi\sqrt m}
       \sum_{n>N}\frac{\sin(2\pi n\log m/L)}n.
\]
Dirichlet summation and the bounded sine-series partial sums give
\[
 \left|\sum_{n>N}\frac{\sin(n\theta)}n\right|
       \le C\min\left(1,\frac1{N|\sin(\theta/2)|}\right).
\]
Separate the largest integer $m_0<\lambda^2$: its contribution is
$O(|B_\lambda^+|/\lambda)$ even when $\lambda^2$ approaches an integer.
For every remaining $m$, $\lambda^2-m\ge1$. Splitting at $m=\lambda$
and $m=\lambda^2/2$ gives
\[
 \sum_{\substack{1<m<\lambda^2\\m\ne m_0}}
 \frac{m^{-1/2}}{|\sin(\pi\log m/L)|}
       \le C\lambda L(1+\log\lambda).
\]
Near the upper end use
$\log(\lambda^2/m)\ge(\lambda^2-m)/\lambda^2$ and sum the harmonic
denominator. Thus the entire function-jump endpoint error divided by
$|B_\lambda^+|$ is at most
\begin{equation}\label{eq:v114-resonance-bound}
 C\left\{\lambda^{-1}+\frac{\lambda L(1+\log\lambda)}N\right\}.
\end{equation}
This retains, rather than assumes away, near-threshold arithmetic resonance.

Integrate by parts $r=\lceil c\rceil$ times on the smooth pieces, retaining
the jumps of every derivative and the periodic derivative jumps. With
$t_n=2\pi n/L$, the $j$th jump term in a normalized Fourier coefficient
has magnitude at most $J_j/(\sqrt L\,|t_n|^{j+1})$. Its physical-endpoint
tail is at most
\[
 \frac{2J_j}{L}\sum_{n>N}|t_n|^{-j-1}
       \le \frac{C J_j}{j}\left(\frac L{2\pi N}\right)^j.
\]
By Lemma~\ref{lem:v114-endpoint-jets}, summing $1\le j<r$ gives
$C|B_\lambda^+|\lambda^7L/N$ whenever $C\lambda^4L/N<1/2$.
The final smooth-piece remainder, after endpoint summation and division by
$|B_\lambda^+|$, is bounded by
\begin{equation}\label{eq:v114-analytic-tail}
 C\lambda^A(1+L)^Ae^{Ac}r!
                 \left(\frac{CL}{N}\right)^{r-1}.
\end{equation}
Here $A,C$ are fixed and independent of $r,\lambda$.
At $N=N_\lambda$, $r!\le(C\lambda^2)^r$, so the logarithm of this bound
is at most $-6c\log\lambda+O(c+\log\lambda+\log L)$.
It is smaller than every fixed inverse power of $\lambda$.
Combining this with \eqref{eq:v114-resonance-bound} proves the first
claim for $p_\lambda^\circ$.

The same expansion gives, at every $|n|>N_\lambda$,
\[
 |\widehat p_\lambda^\circ(n)|
 \le \frac{C|B_\lambda^+|}{\sqrt L}
       \left(\frac\lambda{|t_n|}+\frac{\lambda^7}{|t_n|^2}\right)
       +|R_n|.
\]
The second term and remainder are absorbed by the first. For the remainder,
its ratio to $|B_\lambda^+|\lambda/(\sqrt L|t_n|)$ decreases with $|n|$
and is already smaller than every fixed inverse power at $N_\lambda$.
This proves \eqref{eq:v114-high-tail} before repair.

Finally let $e=p_\lambda-p_\lambda^\circ$ and
$a=h_\lambda^\circ(0)$. Direct counting gives
$\norm e_\infty\le C\lambda^{3/2}|a|$.
The Fourier Lebesgue constant is $O(\log(N+2))$, so
\[
 \frac{|(P_N-I)e(0)|}{|B_\lambda|}
 \le C\lambda^{3/2}\log(N+2)\frac{|a|}{|B_\lambda|}
       =O_A(\lambda^{-A})
\]
along this polynomial diagonal, by \eqref{eq:v12-zero-value}.
For the coefficient assertion a direct variation estimate is stronger:
if the fixed repair bump is supported in $[-b,b]$, then
\[
 \int_{1/\lambda}^{\lambda}|D_{\log}\mathcal E(\psi)(u)|\frac{du}{u}
 \le \sum_{m\le b\lambda}m^{-1/2}
   \int_0^b x^{-1/2}\bigl(\tfrac12|\psi(x)|+x|\psi'(x)|\bigr)\,dx
 \le C\sqrt\lambda.
\]
Inversion averaging preserves this bound and has no periodic function
jump. Hence $|\widehat e(n)|\le C|a|\sqrt{\lambda L}/|n|$, which is
negligible. Use $B_\lambda^+\asymp B_\lambda$ to finish.
\end{proof}

\begin{corollary}[Weak G1 on the same polynomial-cutoff vector]
\label{cor:v114-polynomial-g1}
Let $R>3$, $N=N_\lambda$, $K\ge N$, and $x\in E_K$, measured by
the weighted norm of Proposition~\ref{prop:g1-weighted-sampler}.
For all sufficiently large $\lambda$,
\begin{align}
 \frac{|\ip{D_{\log}W_{\lambda,K}k_\lambda}{x}|}
      {|\delta_K(k_\lambda)|}
 &\le C_R\norm x_{\mathcal H_{R,L}}
 \left[\mathfrak r(\lambda)+\lambda^2
 \left\{\frac LN+\left(\frac LN\right)^{R-1}
                   (1+\log(2N))\right\}\right].
 \label{eq:v114-polynomial-g1}
\end{align}
The bound is uniform in $K$ and tends to zero on bounded weighted test
families. This is a fixed-source weak pairing statement, not control
in operator norm on a growing source block.
\end{corollary}
\begin{proof}
The coefficients of $(P_N-I)p_\lambda$ vanish for $|n|\le N$.
Insert \eqref{eq:v114-high-tail} into the absolute convolution proof
of Proposition~\ref{prop:g1-weighted-sampler}: on the only frequencies
used there, $\mathcal V_\lambda$ is replaced by $C|B_\lambda|\lambda$.
Both its off-diagonal estimate and its separately treated logarithmic
diagonal therefore give the displayed $\lambda^2$ error. The sums
converge absolutely for $R>3$; the fixed-window domain identification
in Proposition~\ref{prop:v114-weil-core} also justifies the passage
from finite Fourier sums to $p_\lambda$.
The continuum term is Proposition~\ref{prop:g1-sobolev}. Finally
$|\delta_K(k_\lambda)|\ge(1-C/\lambda)|B_\lambda|$ and, with the
first-slot-linear convention,
$\ip{D_{\log}W_{\lambda,K}k_\lambda}{x}
=QW_\lambda(k_\lambda,D_{\log}x)$.
Unlike the filtered vector, $P_Np_\lambda$ lies in $E_N$, so $K\ge N$
suffices.
\end{proof}

These results assert no positive Weil complement along this diagonal,
a small inverse-weighted spectral residual, or control of the corrected
sampler's graph defect. Those obligations remain separate. The fixed
orders $0,4$ are essential; no uniform claim over a growing prolate block
is made.

\subsection*{Use on the centered and corrected sampler}
The weighted estimate is unchanged by the factors $(-1)^n$, because its
proof bounds coefficient magnitudes. For any fixed finite root family,
sufficient Hardy smoothing therefore controls the centered sampler and
any fixed finite number of its derivatives. The source input is now
proved in Section~\ref{sec:g1}; constants may depend on the test family.
It supplies no uniformity for a derivative order growing with $L$.

For the later even endpoint shell, the zeroth residual pairing is preserved
exactly by parity: an even candidate $k$ has odd $D_{\log}Wk$ and hence
zero pairing with that shell. Higher residual jets do not all have the
same parity, so no automatic preservation is claimed for every jet.
A fixed finite Gramian made from pairings that tend to zero also tends to
zero, by summing their squared magnitudes. Such a Gramian cannot by itself
supply a uniformly positive metric on that root space.

\begin{lemma}[A legitimate common-cutoff selection]\label{lem:diagonal-selection}
Suppose a normalized G1 error is at most $\epsilon$ for every
$N\ge N_0(L,\epsilon)$ and every ambient cutoff $K\ge2N$.
Suppose finitely many additional errors are at most $\epsilon$ for every
$K\ge K_0(L,N,\epsilon)$. Then one can choose a common cutoff satisfying
all these bounds. The same statement holds with nested eventual thresholds
$N\ge N_0$, $K\ge K_0(N)$ when that is what has actually been proved.
\end{lemma}
\begin{proof}
Choose an admissible core cutoff first, then choose the ambient cutoff above
all the finitely many thresholds. Apply this with a positive tolerance
sequence tending to zero. Each error must be normalized by the actual
endpoint or scale used in the target expression before taking thresholds.
Merely knowing that two sets of acceptable cutoffs are unbounded would not
suffice: two unbounded sets can be disjoint.
\end{proof}
The explicit corrected-sampler rule below can be enlarged to meet additional
proved eventual thresholds: its error bounds still vanish when $K$ grows
beyond that rule. A formula can be independent of the zero list while a
proof of its convergence on a fixed hypothetical window has window-dependent
constants. Neither observation proves a missing graph-pairing estimate.
\section{A consistent finite sampling problem}\label{sec:sampling}

The immediate question is whether a finite sampler can simultaneously retain
the Cauchy evaluation, use the actual CCM endpoint functional, and converge
in the same finite Weil form. The graph identity required by the wider
G1/G2 argument is a fourth property. It is kept separate here because it
does not follow from the other three.

Fix $L\ge2$, $\lambda=e^{L/2}$, $A=L/2$, and
\[
 t_n=2\pi n/L,\qquad s_n=\tfrac12+it_n.
\]
In the coordinate $x=\log(\lambda u)\in[0,L]$, use
\[
 V_n(x)=L^{-1/2}e^{it_nx},\qquad D_LV_n=t_nV_n.
\]
Interval vectors are extended by zero outside the interval when evaluating the
global Weil form. Write $W_{L,N}$ for its compression to $|n|\le N$.
This is the same compression previously denoted $W_{\lambda,N}$, with $\lambda=e^{L/2}$; $\Q=QW$ denotes the same global Weil form in the fixed first-slot-linear convention. The coordinate change is unitary. The Hermitian product is linear in its first slot. The physical endpoint is
\begin{equation}\label{v1:eq:delta}
 \delta_{L,N}\Big(\sum a_nV_n\Big)=L^{-1/2}\sum a_n,
 \qquad \eta_{L,N}=\sqrt L\,\delta_{L,N}.
\end{equation}
This is the equal-coefficient endpoint of the original matrix, not an
interior evaluation. The basis and matrix realization are those of
\cite{ConnesConsaniMoscovici2025}; its opposite sesquilinear convention is translated here.

To remove the Fourier-sign ambiguity, define
\begin{equation}\label{v1:eq:transform}
 H_F(s)=\int_{\mathbb R}g_F(y)e^{-(s-1/2)y}\,\dd y,
 \qquad
 g_F(y)=\frac1{2\pi}\int_{\mathbb R}H_F(\tfrac12+it)e^{ity}\,\dd t.
\end{equation}
Thus the critical-line transform has the usual negative Fourier sign.
Reversing this sign reverses physical tails and Fourier indices, but does
not remove the half-period character identified below.

For the manuscript's Hardy output put
\begin{equation}\label{v1:eq:hardy}
 \begin{split}
 H_F(s)&=\kappa(s)F(s),\\
 \kappa(s)&=\frac{\sigma(\sigma+a)^M}{\sigma-1}
       \frac{s-1}{(s+a)^M(\sigma-s)},
 \qquad \E(F)=\sigma F(\sigma).
 \end{split}
\end{equation}
Here $F$ ranges over a fixed finite-dimensional space of root functions
$\zeta(s)/(s-\rho)^j$, with the denominator canceled by a genuine zeta
zero of order at least $j$. For the one-sided application all these zeros
have real part greater than $1/2$. This is a conditional model: their
existence is not asserted. The additional output $H=\kappa\zeta$ will
provide an unconditional numerical test of the sampling mechanism.

Choosing $M$ sufficiently large gives uniform polynomial vertical decay
of any fixed order needed below on the relevant bounded strip. Contour
shifts, after cancellation at $s=1$ and at the selected zeros, imply the
following useful facts. They also follow directly by Fourier inversion:
\begin{equation}\label{v1:eq:tails}
 g_F(y)=\E(F)e^{qy}\quad(y<0),\qquad q=\sigma-\tfrac12,
 \qquad |g_F^{(j)}(y)|\le C e^{-ry}\quad(y>0),\quad j\le2,
\end{equation}
where any fixed $r<a+1/2$ can be used after sufficient smoothing.
Constants are uniform on a finite-dimensional unit sphere. In particular,
$g_F(0)=\E(F)$. To verify the negative tail, close the inverse contour to
the right for $y<0$. The only pole crossed is $\sigma$, whose residue is
$-\sigma F(\sigma)$; the contour orientation gives \eqref{v1:eq:tails}.
For the positive tail shift left to $\Re s=1/2-r>-a$.
Using a rate strictly below $a+1/2$ avoids hiding the polynomial associated
with the order-$M$ pole at $-a$.

Both transfer and the uncentered nonconvergence theorem use $q,r>1/2$.
These margins are available for every fixed $\sigma>1$ and $a>0$, after
choosing $r<a+1/2$ and taking $M$ sufficiently large. No conclusion below
requires moving these parameters with $L$.

In the convention \eqref{v1:eq:transform}, the first-slot-linear zero formula is
\begin{equation}\label{v1:eq:zeros}
 \Q(f,g)=\sum_{\rho\in\Z}
 H_f(\rho)\overline{H_g(\rho^\sharp)},
 \qquad \rho^\sharp=1-\overline\rho,
\end{equation}
with nontrivial zeros counted with multiplicity. The symmetries of the
zero set make this consistent with the Fourier form of the explicit
formula. All sums used below are absolutely convergent. The standard
bound $N(T)=O(T\log T)$ implies
$\sum_{\rho}(1+|\Im\rho|)^{-2}<\infty$.

\section{The two samplers are different finite vectors}

Let $P_Lg(y)=\sum_{k\in\mathbb Z}g(y+kL)$. Its Fourier series is
$L^{-1}\sum_n H(s_n)e^{it_ny}$. Define
\begin{align}
 \Ju_{L,N}F&=L^{-1/2}\sum_{|n|\le N}H_F(s_n)V_n,
       \label{v1:eq:raw}\\
 \Jc_{L,N}F&=L^{-1/2}\sum_{|n|\le N}(-1)^nH_F(s_n)V_n.
       \label{v1:eq:centered}
\end{align}
The superscripts mean uncentered and centered, respectively.
For fixed $L$, their full Fourier limits on $[0,L]$ are
\[
 u_{L,F}(x)=P_Lg_F(x),\qquad
 c_{L,F}(x)=P_Lg_F(x-A).
\]
Translation of $c_{L,F}$ back to $[-A,A]$ gives the centered fundamental
representative. Translation invariance of the global form therefore
applies to $c_{L,F}$; it does not identify $u_{L,F}$ with that representative.

\begin{proposition}[Endpoint and interior evaluations]\label{v1:prop:endpoints}
With Fourier truncations taken sufficiently far,
\[
 \delta(\Ju F)\longrightarrow \E(F),\qquad
 \delta(\Jc F)\longrightarrow0\quad(L\to\infty).
\]
On the centered sampler the evaluation at $x=A$, rather than at the cut,
converges to $\E(F)$.
\end{proposition}
\begin{proof}
For fixed $L$, Fourier convergence gives
$\delta(\Ju F)\to P_Lg_F(0)$ and
$\delta(\Jc F)\to P_Lg_F(A)$.
Equation \eqref{v1:eq:tails} gives
$P_Lg_F(0)=\E(F)+O(e^{-\min(q,r)L})$ and
$P_Lg_F(A)=O(e^{-\min(q,r)L/2})$.
Evaluation of \eqref{v1:eq:centered} at $A$ cancels $(-1)^n$.
\end{proof}

The diagonal unitary $P=\operatorname{diag}((-1)^n)$ commutes with $D_L$.
A pure change of basis must replace $W$ by $PWP$ and the boundary covector
by its transformed covector. Applying $P$ to the samples while retaining
the original $W$ and $\delta$ is instead a different physical vector.
That is precisely the change used in \eqref{v1:eq:centered}.

\section{Centered Weil transfer with the boundary terms retained}

\begin{lemma}[A sufficient strip estimate]\label{v1:lem:strip}
Suppose $|g^{(j)}(y)|\le C e^{-\mu|y|}$ for $j\le2$, with $\mu>1/2$.
Let $h_L=1_{[-A,A]}P_Lg$ and let $H_L$ be its transform. Then, uniformly
for $0\le\Re s\le1$,
\begin{equation}\label{v1:eq:stripbound}
 |H_L(s)-H(s)|\le
 \frac{C e^{-(\mu-1/2)L/2}}{1+|\Im s|}.
\end{equation}
The estimate includes the jumps of the zero-extended representative.
\end{lemma}
\begin{proof}
Set $e_L=h_L-g$ and $d=\Re s-1/2$.
Outside $[-A,A]$ use the tails of $g$; inside use the nonzero periodic
images. Integrating the resulting geometric series, including their first
derivatives, gives
\[
 \norm{e_L e^{-d\,\cdot}}_1+
 \operatorname{Var}(e_L e^{-d\,\cdot})
 \le C e^{-(\mu-1/2)A},\qquad |d|\le1/2.
\]
The variation includes both endpoint jumps, whose weighted sizes obey
the same bound. The $L^1$ bound handles bounded ordinates. One integration
by parts in the sense of a finite derivative measure handles large
ordinates, giving \eqref{v1:eq:stripbound}.
\end{proof}

Arbitrary inverse powers of the ordinate are not claimed. A nonzero
endpoint jump normally produces a $1/t$ term. That decay is sufficient:
each Weil summand is a product of two transforms.

\begin{theorem}[Centered form transfer]\label{v1:thm:center}
On any fixed finite-dimensional family satisfying the preceding decay
bounds and $|H_F(s)|\le C(1+|\Im s|)^{-1}$ on the strip,
\[
 \Q(c_{L,F},c_{L,G})\longrightarrow\Q(g_F,g_G)
\]
uniformly on bounded sets of $F,G$. The same conclusion holds for
$\ip{W_{L,N}\Jc_{L,N}F}{\Jc_{L,N}G}$ with a sufficiently large common
Fourier cutoff. On a one-sided root space the self-form limit is zero.
For reflected top-root outputs, the nonzero global cross-form also transfers.
\end{theorem}
\begin{proof}
Let $\epsilon_L=e^{-(\mu-1/2)L/2}$. In \eqref{v1:eq:zeros}, replace
each transform by $H_F+E_{L,F}$. Lemma~\ref{v1:lem:strip} bounds each difference
by $C\epsilon_L/(1+|\Im\rho|)$, and the unchanged transform has the same
ordinate decay without $\epsilon_L$. Hence the difference of the forms is
at most
$C(\epsilon_L+\epsilon_L^2)\sum_\rho(1+|\Im\rho|)^{-2}$.
Translation of $h_L$ to $[0,L]$ does not change the paired form.

For completeness, Fourier truncation can be justified directly, without
inferring convergence of these particular projections from abstract core
density. The correct CCM matrix bounds are
\begin{equation}\label{v1:eq:matrix}
 |W_{nm}|\le C\lambda/|n-m|\ (n\ne m),\qquad
 |W_{nn}|\le C\{\lambda+1+\log(2+|n|/L)\}.
\end{equation}
The diagonal uses $2(1-y/L)\cos(t_ny)$, not a limit of the off-diagonal
kernel. Near zero its archimedean numerator is bounded by
$C\{\min(1,t_n^2y^2)+y\}$; integration against $O(1/y)$ gives the
logarithm in \eqref{v1:eq:matrix}. The prime and pole pieces are $O(\lambda)$.

If $|H_F(s_n)|\le C(1+|n|/L)^{-R}$ with $R\ge2$, then the normalized
coefficients $a_n$ have
\[
 \sum_n|a_n|\le C\sqrt L,\qquad
 \sum_{|n|>N}|a_n|\le C\sqrt L(1+N/L)^{1-R}.
\]
The absolute double matrix series therefore converges. Bounding off-diagonal
entries by $C\lambda$ and treating the logarithmic diagonal separately gives,
for $N\ge L$,
\begin{equation}\label{v1:eq:galerkin}
 |\Q(c_{L,F},c_{L,G})-
 \ip{W_{L,N}\Jc_{L,N}F}{\Jc_{L,N}G}|
 \le C\lambda L(1+N/L)^{1-R}.
\end{equation}
The constants are uniform on the fixed family. This also identifies the
matrix series with the semilocal form, using its closed realization.

For a one-sided root space, $H_F$ is nonzero at at most finitely many zeros,
all on one side, so each paired summand in its self-form is zero. If $F$ and
$G$ are the top root functions at $\rho$ and $\rho^\sharp$, their cross-form is
\[
 m_\rho\,\kappa(\rho)\frac{\zeta^{(m_\rho)}(\rho)}{m_\rho!}
 \overline{\kappa(\rho^\sharp)
 \frac{\zeta^{(m_\rho)}(\rho^\sharp)}{m_\rho!}},
\]
which is nonzero. This is the first-slot-linear convention.
\end{proof}

\section{The uncentered sampler has an explicit cut defect}

The failure of the original proof is more than an omitted phase in its
notation. It is possible to compute the form created by the different cut.
In this section assume $q,r>1/2$, and abbreviate $E=\E(F)$,
$H=H_F$, $z_\rho=\rho-1/2$. Let $g_+=1_{(0,\infty)}g$ and
$g_-=1_{(-\infty,0)}g$. Set
\[
 f_L=g_++g_-(\,\cdot-L).
\]
The negative half has been moved by $L$. Since
$H_-(s)=E/(\sigma-s)$, its full transform is
\begin{equation}\label{v1:eq:cuttransform}
 \widehat f_L(s)=H(s)+
       E\frac{e^{-(s-1/2)L}-1}{\sigma-s}.
\end{equation}

\begin{lemma}[Comparison with the actual periodic cut]\label{v1:lem:uncut}
For $u_L=1_{[0,L]}P_Lg$ and some $\nu>0$,
\[
 \Q(u_L,u_L)=\Q(f_L,f_L)+O(e^{-\nu L}).
\]
One may take $\nu=\min(q,r)-1/2>0$.
\end{lemma}
\begin{proof}
Write $s=1/2+d+it$, $|d|\le1/2$, and $p_0=\min(q,r)$.
On $[0,L]$ the negative images sum exactly to
$E e^{q(x-L)}/(1-e^{-qL})$. Their transform differs from the transform
of $g_-(\,\cdot-L)$ by exactly
\[
 \frac{E e^{-qL}}{1-e^{-qL}}
       \frac{e^{-(s-1/2)L}-1}{\sigma-s}.
\]
The extra positive images are $O(e^{-rL}e^{-rx})$ on $[0,L]$; the removed
positive tail lies in $x>L$. Their weighted $L^1$ norms and variations,
including endpoint jumps, give the same type of bound. Consequently
\[
 |\widehat u_L(s)-\widehat f_L(s)|
 \le \frac{C e^{-p_0L}(1+e^{-dL})}{1+|t|},\qquad
 |\widehat f_L(s)|\le \frac{C(1+e^{-dL})}{1+|t|}.
\]
It matters that the partner has real displacement $-d$. Products in the
form error contain
$(1+e^{-dL})(1+e^{dL})\le4e^{L/2}$, rather than two independent
worst-case factors. Expanding the zero sum therefore bounds the error by
$C e^{-(p_0-1/2)L}\sum_\rho(1+|\Im\rho|)^{-2}$, as required.
\end{proof}

For the one-sided root family, let $S$ be the finite set of zeros at which
$H$ does not vanish. Sums over $S$ below include their multiplicities.
For $H=\kappa\zeta$, take $S$ empty.

\begin{proposition}[Cut-defect formula]\label{v1:prop:cut}
Under the stated hypotheses,
\begin{align}
 \Q(u_L,u_L)={}&\Q(g,g)
 +2\Re\sum_{\rho\in S}
 \frac{H(\rho)\overline E}{q+z_\rho}
          (e^{z_\rho L}-1)\notag\\
 &+2|E|^2\sum_{\rho\in\Z}
       \frac{1-e^{z_\rho L}}{q^2-z_\rho^2}
 +O(e^{-\nu L}).\label{v1:eq:cutdefect}
\end{align}
Both displayed sums are well defined without RH; the second is absolutely
convergent for each fixed $L$.
\end{proposition}
\begin{proof}
Insert \eqref{v1:eq:cuttransform} into \eqref{v1:eq:zeros}. The two mixed sums
are complex conjugates after replacing $\rho$ by $\rho^\sharp$.
The product of the two added terms has numerator
$2-e^{z_\rho L}-e^{-z_\rho L}$ and denominator $q^2-z_\rho^2$.
The symmetry $\rho\mapsto1-\rho$ makes the two exponential sums equal,
giving \eqref{v1:eq:cutdefect}. Lemma~\ref{v1:lem:uncut} handles the actual cut.
Absolute convergence follows from $|\Re z_\rho|<1/2$ and the quadratic
denominator at large ordinates.
\end{proof}

\begin{theorem}[Failure of the claimed full-scale raw transfer]\label{v1:thm:rawfailure}
Suppose $E\ne0$, and $H$ vanishes at every critical-line zero and at all
but finitely many other nontrivial zeros. Under the preceding decay
hypotheses, $\Q(u_L,u_L)$ does not converge to any finite limit as real
$L\to\infty$. In particular it cannot converge to $\Q(g,g)=0$ for a
one-sided root output with nonzero $E$.
\end{theorem}
\begin{proof}
Choose a critical-line zero $\rho_0=1/2+i\gamma_0$ of multiplicity $m_0$,
with $\gamma_0\ne0$. Such zeros exist unconditionally by Hardy's theorem
\cite{Hardy}. The coefficient of $e^{i\gamma_0L}$ in the infinite sum
in \eqref{v1:eq:cutdefect} is
\begin{equation}\label{v1:eq:residue}
 -\frac{2|E|^2m_0}{q^2+\gamma_0^2}\ne0.
\end{equation}
The finite mixed terms have exponents strictly off the imaginary axis,
so do not contribute at this exponent.

Here is an argument that does not assume away possible off-line zeros.
Take the Laplace transform in $L\ge L_0$ of \eqref{v1:eq:cutdefect}, initially
for $\Re p>1/2$. Its series of rational-exponential terms continues
meromorphically across the imaginary axis: the denominators at large
ordinates are $O(|\Im\rho|^3)$, giving locally normal convergence away
from the discrete poles. The exponentially decaying error is analytic
in $\Re p>-\nu$. At $p=i\gamma_0$ the continued transform has residue
\eqref{v1:eq:residue}; the lower integration limit does not change the residue.

If $\Q(u_L,u_L)$ had a finite limit, it would be bounded and its Laplace
transform would be holomorphic for $\Re p>0$. All apparent poles in that
half-plane would therefore have to cancel. Moreover bounded convergence,
after subtracting the finite limit, would give
\[
 \lim_{\epsilon\downarrow0}\epsilon
 \int_{L_0}^\infty e^{-(\epsilon+i\gamma_0)L}
           \Q(u_L,u_L)\,\dd L=0.
\]
For the constant part the same limit is zero since $\gamma_0\ne0$.
But the meromorphic expression gives the nonzero residue
\eqref{v1:eq:residue}. This is a contradiction.
\end{proof}

\begin{remark}[Exact scope]
The theorem concerns the full real-scale limit. It does not exclude
specially chosen subsequences $L_j$ on which the form happens to approach
zero. Nor does it exclude different sampling weights or a changed
matrix. It applies to Galerkin approximations whenever their error relative
to this continuum cut tends to zero; choosing a larger faithful Fourier
cutoff cannot remove the continuum defect.

For any nonempty hypothetical one-sided root window, a basis vector has
$E=\sigma\zeta(\sigma)/(\sigma-\rho)^j\ne0$. Thus the original full-scale
transfer and the original unphased sampling convention are incompatible
in that window. This is a conditional diagnosis, not an assertion that an
off-line zero exists or an unconditional disproof of an implication with
a potentially empty premise. The output $H=\kappa\zeta$ is an unconditional
example of the same failure in the larger analytic output class.

On $\ker\E$, the negative half in \eqref{v1:eq:tails} vanishes, and this cut
obstruction disappears. That codimension-one restriction does not recover
the whole nonempty root window.
\end{remark}

\section{A single finite object with form transfer and exact endpoint}

For an integer $K\ge1$ use the ambient space $|n|\le2K$, and define the
even endpoint shell
\begin{equation}\label{v1:eq:shell}
 c_{L,K}=\frac{\sqrt L}{2K}\sum_{K<|n|\le2K}V_n.
\end{equation}
There are exactly $2K$ terms. Consequently
\begin{equation}\label{v1:eq:shellnorms}
 \delta(c_{L,K})=1,\qquad
 \norm{c_{L,K}}^2=\frac L{2K},\qquad
 \norm{D_Lc_{L,K}}^2=\frac{\pi^2}{3L}(14K+9+K^{-1}).
\end{equation}
Define a linear functional and a corrected sampler by
\begin{equation}\label{v1:eq:repair}
 d_{L,K}(F)=\E(F)-\delta(\Jc_{L,2K}F),\qquad
 \Jr_{L,K}F=\Jc_{L,2K}F+c_{L,K}d_{L,K}(F).
\end{equation}

\begin{theorem}[Three-property compatibility]\label{v1:thm:repair}
On a fixed finite-dimensional Hardy-output family as above, choose
\[
 K(L)=\lceil\lambda^4L^4\rceil,\qquad \lambda=e^{L/2},
\]
and choose $M$ sufficiently large to give $R\ge2$ in
\eqref{v1:eq:galerkin}. The same finite matrix, sampler, and endpoint obey
\begin{align}
 \delta_{L,2K}(\Jr_{L,K}F)&=\sigma F(\sigma)
                       &&\text{exactly},\label{v1:eq:exactendpoint}\\
 \norm{\Jr_{L,K}F-\Jc_{L,2K}F}&\longrightarrow0,\label{v1:eq:normrepair}\\
 \ip{W_{L,2K}\Jr_{L,K}F}{\Jr_{L,K}G}
            &\longrightarrow\Q(g_F,g_G),\label{v1:eq:formrepair}\\
 \ip{\Jr_{L,K}F}{\Jr_{L,K}G}
            &\longrightarrow\ip{g_F}{g_G}.\label{v1:eq:hilbertrepair}
\end{align}
The limits are uniform on bounded subsets of the fixed family. The
construction requires no zero list in the matrix, shell, cutoff rule, or
formula for the endpoint correction. It is not an economical cutoff claim.
\end{theorem}
\begin{proof}
The endpoint identity follows immediately from \eqref{v1:eq:shellnorms}.
Uniform decay of the critical-line samples gives
$\norm{\Jc F}\le C\norm F$, and the corresponding Riemann sum converges
to $\norm{g_F}^2$. The same coefficient bound and the endpoint result give
$\norm{d_{L,K}}\le C$ for large $L$ and the stated cutoffs. Therefore
$\norm{\Jr F-\Jc F}\le C\sqrt{L/K}\norm F$.

The Schur row-sum estimate applied to the correct matrix bounds yields
\begin{equation}\label{v1:eq:opnorm}
 \norm{W_{L,2K}}\le C\lambda(1+\log(4K+1)).
\end{equation}
Expansion of the two corrected arguments bounds the change in their
form pairing by
\begin{equation}\label{v1:eq:formprice}
 C\lambda(1+\log(4K+1))
       \left(\sqrt{L/K}+L/K\right)\norm F\norm G.
\end{equation}
For the displayed cutoff this is $O(\lambda^{-1}L^{-1/2})$.
The centered Fourier-truncation error in \eqref{v1:eq:galerkin} is at most
$O(\lambda^{-3}L^{-2})$ when $R=2$, and smaller for larger $R$.
Theorem~\ref{v1:thm:center} proves \eqref{v1:eq:formrepair}.
The ordinary norm correction tends to zero, so the Riemann-sum limit also
proves \eqref{v1:eq:hilbertrepair}.
\end{proof}

If $F\mapsto g_F$ is injective on the fixed finite-dimensional space, the
last limit gives a uniform lower ordinary-norm bound there. In particular,
on a one-sided root space the corrected self-form tends to zero, and adding
any fixed positive scalar to $W$ gives restricted coercivity for sufficiently
large cutoffs. The nonzero reflected cross-form also transfers to this
corrected sampler. These are legitimate consequences of this particular
construction; they do not identify its graph transport with the old one.

\section{The graph identity is the remaining compatibility cost}

Suppose the root-space operator satisfies
\[
 BF=(s-\tfrac12)F-\eta_{\mathrm B}(F)\zeta.
\]
The centered samples satisfy the exact coefficientwise identity
\begin{equation}\label{v1:eq:centergraph}
 iD_L\Jc-\Jc B=z^{\mathrm c}_{L,2K}\eta_{\mathrm B},\qquad
 z^{\mathrm c}_{L,2K}=L^{-1/2}\sum_{|n|\le2K}
        (-1)^n\kappa(s_n)\zeta(s_n)V_n.
\end{equation}
For the corrected sampler, however,
\begin{equation}\label{v1:eq:graphdefect}
 iD_L\Jr-\Jr B=z^{\mathrm c}_{L,2K}\eta_{\mathrm B}+\mathcal R_{L,K},
 \qquad
 \mathcal R_{L,K}F=iD_Lc_{L,K}\,d_{L,K}(F)
                         -c_{L,K}\,d_{L,K}(BF).
\end{equation}
This is an exact rank-at-most-two defect. Parity gives
$\ip{D_Lc_{L,K}}{c_{L,K}}=0$, so
\begin{equation}\label{v1:eq:graphnorm}
 \norm{\mathcal R_{L,K}F}^2
 =|d_{L,K}(F)|^2\frac{\pi^2(14K+9+K^{-1})}{3L}
       +|d_{L,K}(BF)|^2\frac L{2K}.
\end{equation}
For $\E(F)\ne0$, this diverges along the displayed construction.
Smallness of a correction in ordinary norm or Weil form therefore does
not make it small in the graph identity.

\begin{proposition}[Rigidity and a general trace obstruction]\label{v1:prop:graph}
For a fixed one-sided off-line root space, the equation
$iD_LJ-JB=z\eta_{\mathrm B}$ has at most one solution $J$ for each prescribed $z$.
Moreover, any family of corrections $T_L$ with $\norm{T_L}\to0$ and
$\delta(T_LF)\to\E(F)\ne0$ for some fixed $F$ has
$\norm{iD_LT_LF-T_LBF}\to\infty$.
\end{proposition}
\begin{proof}
All eigenvalues of $B$ have positive real part, whereas those of $iD_L$
are purely imaginary. Row by row, the homogeneous equation is
$\ell_n(it_nI-B)=0$, and $it_nI-B$ is invertible. This proves uniqueness.

For a periodic $H^1$ function $b$ on $[0,L]$, choose a point where
$|b|^2$ does not exceed its average and integrate the derivative of
$|b|^2$. Cauchy--Schwarz gives
\begin{equation}\label{v1:eq:trace}
 |b(0)|^2\le L^{-1}\norm b^2+2\norm b\,\norm{D_Lb}.
\end{equation}
Apply this to $b=T_LF$. Its norm tends to zero while its endpoint tends
to a nonzero value, so $\norm{D_LT_LF}\to\infty$.
Since $B$ is bounded on the fixed finite-dimensional space,
$\norm{T_LBF}\to0$. The reverse triangle inequality proves the claim.
\end{proof}

There is also a useful exact description at the weaker pairing level.
Here $\eta$ denotes the boundary covector $\eta_{L,2K}$, distinct from
the root-space functional $\eta_{\mathrm B}$. In physical scaling the finite commutator is
\[
 [D_L,W]=|\beta\rangle\eta- |\eta\rangle\beta^*,
 \qquad \eta(c_{L,K})=\sqrt L,\qquad \beta^*(c_{L,K})=0,
\]
where $\beta=\beta^{\log}$ includes the physical factor $2\pi/L$ and is odd and $c_{L,K}$ is even. Here the rank-one notation
means $|a\rangle b^*:x\mapsto a\ip{x}{b}$, and the vector representing
$\eta$ is understood. Therefore for every finite test $v$,
\begin{equation}\label{v1:eq:paireddefect}
 \begin{split}
 \ip{W\mathcal R_{L,K}F}{v}
 ={}&i\,d_{L,K}(F)\ip{D_LWc_{L,K}}{v}
       -i\sqrt L\,d_{L,K}(F)\ip{\beta}{v}\\
    &-d_{L,K}(BF)\ip{Wc_{L,K}}{v}.
 \end{split}
\end{equation}
Thus even where the first and last terms can be made small, the boundary
commutator term has not disappeared. No unproved asymptotic value for this
term is used in this manuscript.

One part of G1 is preserved particularly simply. The original finite matrix
preserves parity; if the G1 candidate $k$ is even, $D_LWk$ is odd. Hence
\begin{equation}\label{v1:eq:g1parity}
 \ip{D_LWk}{\Jr F}=\ip{D_LWk}{\Jc F}.
\end{equation}
This is exact and requires no estimate on the shell derivative. Establishing
the centered right-hand G1 bound still requires the independently audited
continuum G1 inputs. Equation \eqref{v1:eq:g1parity} does not cancel
\eqref{v1:eq:paireddefect} in the later graph calculation.

\section{The rebuilt finite metric and reflected-partner argument}\label{sec:g2}
Throughout this section $J=\Jr_{L,K}$, $W=W_{L,2K}$, $D=D_L$,
$z=z^{\mathrm c}_{L,2K}$, and $\mathcal R=\mathcal R_{L,K}$ are the
objects of \eqref{v1:eq:repair} and \eqref{v1:eq:graphdefect} on the
\emph{same} finite space. Write $\eta_\partial=\sqrt L\delta$ and
$\beta=\beta^{\log}$. Then
\[
 [D,W]=|\beta\rangle\eta_\partial-|\eta_\partial\rangle\beta^*.
\]
The rank-one convention is $|u\rangle v^*:x\mapsto u\ip{x}{v}$.
In particular the root functional $\eta_{\mathrm B}$ is not identified
with either finite endpoint functional.
Theorem~\ref{thm:v15-finite-metric} gives the full finite-core metric
budget without a global closed-operator assumption. The source-annihilating
projection of Proposition~\ref{prop:v15-projector} is a separate positive
benchmark on this same finite object, not a replacement for $W$ in the
identities below or a claim of small metric defect. Proposition~\ref{prop:v16-projected-cut}
also shows why inserting that sharp band projection into a Weil pairing
does not automatically preserve the earlier holomorphic-strip transfer.
Proposition~\ref{prop:v17-pair-edge} further shows that reflection alone
does not cancel the leading edge term; control of the entire zero sum is
still required.
Proposition~\ref{prop:v18-full-projection} supplies a different,
full-space projection that kills the source and both shell directions
while preserving ordinary coercivity and Weil transfer through a proved
full-strip product bound. Its exact commutator remains in
\eqref{eq:v18-full-budget}. Corollary~\ref{cor:v18-shift} shows why
adding a fixed positive scalar to its projected Weil form still fails
the small-defect criterion. Neither result licenses meromorphic
resolvent transfer or changes the endpoint of $J$.
With the stronger fixed tail margin $p>1$,
Proposition~\ref{prop:v19-action} now controls $WSJF$ in ordinary norm
by an explicit finite exponential sum. Corollary~\ref{cor:v19-square}
audits the globally positive candidate $SW^2S$: top-root mass grows,
while lower jets at multiple zeros lose unscaled coercivity. The full
squared-metric commutator is still required.

\begin{proposition}[Restricted positive shifts after corrected transfer]\label{prop:v1-shift}
On any fixed one-sided root space, for every fixed $a_0>0$ there is
$c>0$ such that, for sufficiently large $L$,
\[
 \ip{(W+a_0I)JF}{JF}\ge c\|F\|^2.
\]
This is restricted coercivity, not positivity on the full finite space.
\end{proposition}
\begin{proof}
The output map $F\mapsto g_F$ is injective, so its Gram matrix has a
positive minimum eigenvalue on the fixed space. Theorem~\ref{v1:thm:repair}
transfers that ordinary Gram matrix, while the Weil Gram matrix tends to
zero by one-sided isotropy. The positive scalar term therefore dominates.
\end{proof}
If a negative compact Weil test is used to bound the lowest finite form
eigenvalue away from zero under failure of RH, its particular Fourier
approximants must also converge in the closed form. For a smooth compact
test this follows from the correct matrix bound and rapid coefficient
decay. This optional calibration is not needed for the proposition.
In either formulation coercivity alone does not imply a small Lyapunov defect.

\begin{proposition}[Lyapunov identity with the graph defect retained]\label{prop:v1-lyapunov}
Let $BF=\mu F$, $d=\Re\mu>0$, $\|F\|=1$, and put
\[
 x=JF,\quad e=\eta_{\mathrm B}(F),\quad
 y=e z+\mathcal RF,\quad b=\eta_\partial(x).
\]
For every self-adjoint finite matrix $S$,
\begin{equation}\label{eq:v1-general-metric}
 d\ip{Sx}{x}
 =-\Re\ip{Sy}{x}-\frac{i}{2}\ip{[D,S]x}{x}.
\end{equation}
For $T=W+a_0I$ this becomes
\begin{equation}\label{eq:v1-residual}
 2d\ip{Tx}{x}=-2\Re\ip{T(ez+\mathcal RF)+ib\beta}{x}.
\end{equation}
If $\ip{Tx}{x}\ge c$ and $\|x\|\le C$, then
\[
 \|T(ez+\mathcal RF)+ib\beta\|\ge dc/C.
\]
\end{proposition}
\begin{proof}
The exact graph equation is $iDx-\mu x=y$. Pair it in both slots
with $Sx$. Self-adjointness gives
$2d\ip{Sx}{x}=-i\ip{[D,S]x}{x}-2\Re\ip{Sy}{x}$.
For $S=T$ the scalar part commutes with $D$. The rank-two commutator
has expectation $b\ip{\beta}{x}-\overline{b\ip{\beta}{x}}$,
which gives \eqref{eq:v1-residual}. Apply Cauchy--Schwarz for the bound.
\end{proof}
For this sampler $b=\sqrt L\,\mathscr E(F)$ exactly. The vector
$\mathcal RF$ cannot be omitted on the ground that $J-\Jc$ tends to zero
in ordinary norm: \eqref{v1:eq:graphnorm} proves the opposite behavior
for its derivative. On a left root window the signs reverse, with the
same absolute-value conclusion.

\begin{proposition}[Candidate comparison and scalar cancellation]\label{prop:v1-candidate}
For an even candidate $k$ with $a=\eta_\partial(k)\ne0$, put
$r=DWk$, $q=ib/a$, and $m=y-qDk$. Then
\begin{align}
 T y+ib\beta&=T m+q(r+a_0Dk),\label{eq:v1-candidate}\\
 d\ip{Wx}{x}&=-\Re\ip{Wm}{x}-\Re\bigl(q\ip{r}{x}\bigr).
 \label{eq:v1-unshifted}
\end{align}
More generally, a self-adjoint correction commuting with $D$ satisfies
$d\ip{Sx}{x}=-\Re\ip{S y}{x}$ exactly.
\end{proposition}
\begin{proof}
Parity gives $[D,W]k=a\beta$, hence $WDk=r-a\beta$.
Since $qa=ib$, expansion proves the first identity. Use
Proposition~\ref{prop:v1-lyapunov} with $S=W$ for the second.
The final statement is \eqref{eq:v1-general-metric} with zero commutator.
\end{proof}
Thus scalar coercivity balances its own source and graph terms. It does
not independently exclude an off-line root. For a noncommuting positive
$S$, the exact combined quantity
\[
 \mathfrak X_S=\Re\ip{S(ez+\mathcal RF)}{x}
       +\frac i2\ip{[D,S]x}{x}=-d\ip{Sx}{x}
\]
still needs an independent upper estimate before its lower bound can be
used in a contradiction. With $S=W^2$, $[D,W^2]=[D,W]W+W[D,W]$ has
rank at most four. Form smallness alone does not give norm smallness.

\begin{proposition}[Reflected top-root lower bound for the corrected sampler]\label{prop:v1-reflected}
Suppose $\rho$ is an off-line zero of exact multiplicity $m$ and
$\Re\rho>1/2$. Set $\rho^\sharp=1-\overline\rho$,
$F_+=\zeta/(s-\rho)^m$, and $F_-=\zeta/(s-\rho^\sharp)^m$.
On the common finite construction, with $x=JF_+$ and $x_-=JF_-$,
\[
 \ip{Wx}{x}\to0,\qquad \ip{Wx}{x_-}\to C_{\rho,m}\ne0,
 \qquad \liminf\|Wx\|>0,
\]
where the first-slot-linear constant is
\begin{equation}\label{eq:v1-crossconstant}
 C_{\rho,m}=m\,\kappa(\rho)\frac{\zeta^{(m)}(\rho)}{m!}
 \overline{\kappa(\rho^\sharp)
                   \frac{\zeta^{(m)}(\rho^\sharp)}{m!}}.
\end{equation}
\end{proposition}
\begin{proof}
The functional equation and conjugation preserve zero multiplicity, so
$\rho^\sharp$ also has exact order $m$. The top root output is nonzero
only at its selected zero in the global value-level zero sum. Thus the
global cross form is exactly \eqref{eq:v1-crossconstant}; its factors
are nonzero. Both outputs meet the holomorphic strip hypotheses of
Theorem~\ref{v1:thm:repair}, which transfers the self and cross forms.
Ordinary norms of $x_-$ are bounded. The inequality
$\|Wx\|\|x_-\|\ge|\ip{Wx}{x_-}|$ proves the lower bound.
\end{proof}
For $m>1$ this is a bound on the top generalized vector. The genuine
eigenvector $\zeta/(s-\rho)$ vanishes at every zero when $m>1$, and the
value-level Weil formula supplies no analogous bound on it. No claim
of coercivity on an entire multiple root block follows from this proposition.

\subsection*{Boundary-annihilating tests: the exact surviving identity}
Write $\mu=\rho-1/2$, $A_\rho=iD-\mu I$, and $x_j=JF_{\rho,j}$.
With $x_0=0$, the corrected graph relation gives
\[
 A_\rho x_m=u_m,\qquad
 u_m=x_{m-1}+\eta_{\mathrm B}(F_{\rho,m})z+\mathcal R F_{\rho,m}.
\]
In the canonical chain $\eta_{\mathrm B}(F_{\rho,1})=1$ and the
functional vanishes on higher chain vectors. Thus the old radical
predecessor now has the additional term $\mathcal R F_{\rho,m}$.
Define $v^\pm=(A_\rho^*)^{-1}x_\pm$, with $x_+=x_m$.
The inverse has norm at most $1/d$, where $d=\Re\rho-1/2$.
If $\delta(v^+)\ne0$, set
\[
 \alpha=\delta(v^-)/\delta(v^+),\qquad
 \widetilde v=v^- -\alpha v^+,\qquad
 \widetilde y=x_- -\alpha x_m=A_\rho^*\widetilde v.
\]
These are finite definitions, not assertions of a bounded limiting ratio.
For $b_m=\eta_\partial(x_m)$, commutation gives exactly
\begin{equation}\label{eq:v1-boundarytest}
 \ip{Wu_m+ib_m\beta}{\widetilde v}
       =\ip{Wx_m}{\widetilde y}.
\end{equation}
Indeed $A_\rho Wx_m=Wu_m+i[D,W]x_m$, and
$\eta_\partial(\widetilde v)=0$ removes the other boundary channel.
If the ratio $\alpha$ is bounded along a chosen sequence, the right side
tends to $C_{\rho,m}$ by Proposition~\ref{prop:v1-reflected} and
conjugate-linearity in the second slot. Corollary~\ref{cor:v11-ratio}
below shows that this bounded-ratio premise actually fails for the
corrected sampler. Even a hypothetical bounded-ratio construction would
also require, before isolating the beta term,
\begin{equation}\label{eq:v1-needed-pole-transfer}
 \ip{W(x_{m-1}+\eta_{\mathrm B}(F_{\rho,m})z)}{\widetilde v}
       +\ip{W\mathcal R F_{\rho,m}}{\widetilde v}\longrightarrow0.
\end{equation}
Neither term is disposed of by the holomorphic form-transfer theorem.
The first involves a potentially meromorphic adjoint-resolvent test; the
second is the new graph correction. Equation~\eqref{v1:eq:paireddefect}
shows its explicit boundary contribution.

\paragraph{What is withdrawn.}
The previous pole-aware transfer proof controlled the singular zero but
did not bound the product error over all zeros. Its conclusion is not used
here. A valid replacement must give a summable majorant for the paired
transforms on the entire strip, including cut jumps and slow reflected
tails, on the actual finite test family. Consequently the old nonzero
limit for $\sqrt L\ip{\beta}{\widetilde v}$ is not an established result
of this version. The exact arithmetic formula for that quantity remains
valid for every finite test, as shown in Section~\ref{sec:arithmetic}.

\subsection*{The endpoint ratio and a limitation on fixed test spaces}
The bounded-ratio premise above can now be decided for the corrected
sampler. Fix a hypothetical zero $\rho=1/2+\mu$, where
$\mu=d+i\gamma$ and $0<d<1/2$, and set
$p=\rho^\sharp=1/2-\overline\mu$. Let $\mathcal F$ be a fixed
finite-dimensional space of Hardy inputs satisfying the holomorphy,
tail and vertical-decay hypotheses of Sections~\ref{sec:sampling}--\ref{sec:g2}.
One may include any fixed finite set of genuine root functions and
$\zeta$ itself, choosing the smoothing order sufficiently large.
Fix any norm on $\mathcal F$. Use the common cutoff
\[
 K\ge\left\lceil e^{2L}L^4\right\rceil,
 \qquad J=\Jr_{L,K},\qquad
 V_LG=(A_\rho^*)^{-1}JG.
\]
The letter $V_L$ here denotes a resolvent map; $V_n$ still denotes a
Fourier basis vector. Choose $r$ as in \eqref{v1:eq:tails} and put
$\nu=\min(q,r,1)>1/2$. All constants below may depend on the fixed
space and on $\rho$, but not on $L$, $K$, or a bounded input.

\begin{lemma}[Endpoint of the corrected adjoint resolvent]
\label{lem:v11-resolvent-endpoint}
Uniformly for $G\in\mathcal F$,
\begin{equation}\label{eq:v11-endpoint-asymptotic}
 \delta V_LG
 =-H_G(p)\frac{e^{-\overline\mu L/2}}{1-e^{-\overline\mu L}}
       +O(e^{-\nu L/2})\|G\|.
\end{equation}
Consequently the rescaled endpoint functional satisfies, in the dual
norm of this fixed space,
\begin{equation}\label{eq:v11-rescaled-endpoint}
 \ell_L(G):=e^{\overline\mu L/2}(1-e^{-\overline\mu L})\delta V_LG
 =-H_G(p)+O(e^{-(\nu-d)L/2})\|G\|.
\end{equation}
\end{lemma}
\begin{proof}
The multiplier of $(A_\rho^*)^{-1}$ is
$(-it-\overline\mu)^{-1}$. Thus the full centered Fourier series is
the periodization of the inverse transform of
\[
 H_{v,G}(s)=-\frac{H_G(s)}{s-p}.
\]
This transform is used only for Fourier inversion, not in a Weil
zero-sum formula. Its possible pole at $p$ has residue $-H_G(p)$.
Contour shifts with the same decay bounds as \eqref{v1:eq:tails} give
\[
 h_G(y)=
 \begin{cases}
 -\E(G)(\sigma-p)^{-1}e^{qy},&y<0,\\
 -H_G(p)e^{-\overline\mu y}+O(e^{-ry})\|G\|,&y>0.
 \end{cases}
\]
Summing these tails at the half-period points $(j+1/2)L$ gives
\begin{align*}
 \delta V_L^{\mathrm c,\infty}G
 ={}&-H_G(p)\frac{e^{-\overline\mu L/2}}{1-e^{-\overline\mu L}}\\
 &-\frac{\E(G)}{\sigma-p}
          \frac{e^{-qL/2}}{1-e^{-qL}}+O(e^{-rL/2})\|G\|.
\end{align*}
If $|H_G(1/2+it)|\le C(1+|t|)^{-R}\|G\|$, $R\ge2$,
the omitted Fourier endpoint tail is at most
$C(1+K/L)^{-R}\|G\|$, because the resolvent supplies one further
inverse power. The endpoint correction has coefficient
$d_{L,K}(G)=\E(G)-\delta\Jc_{L,2K}G=O(\|G\|)$.
For its shell $c=c_{L,K}$, pairing the positive and negative indices
gives the exact identity
\begin{equation}\label{eq:v11-shell-resolvent}
 \begin{split}
 \delta(A_\rho^*)^{-1}c
 &=\frac1{2K}\sum_{K<|n|\le2K}\frac1{-it_n-\overline\mu}\\
 &=-\frac{\overline\mu}{K}
       \sum_{n=K+1}^{2K}\frac1{t_n^2+\overline\mu^{\,2}}
   =O_\mu(L^2/K^2).
 \end{split}
\end{equation}
For large $L$, $K/L$ exceeds every fixed multiple of $|\mu|$,
which justifies the last bound despite the complex denominator.
Both cutoff errors are $O(e^{-4L}L^{-6})$ at the stated minimum
cutoff and decrease upon enlarging it. They are absorbed by
$O(e^{-\nu L/2})$. This proves \eqref{eq:v11-endpoint-asymptotic};
multiplication by the displayed scale proves
\eqref{eq:v11-rescaled-endpoint}.
\end{proof}

\begin{lemma}[Uniform lower norm bound for the resolvent map]
\label{lem:v11-resolvent-gram}
There are $c,C>0$ such that for all sufficiently large $L$,
\[
 c\|G\|\le\|V_LG\|\le C\|G\|\qquad(G\in\mathcal F).
\]
\end{lemma}
\begin{proof}
The Gram form of the centered resolvent samples is a Riemann sum
converging, uniformly on the fixed finite-dimensional space, to
\[
 \frac1{2\pi}\int_{\mathbb R}
       \frac{|H_G(1/2+it)|^2}{|-it-\overline\mu|^2}\,dt.
\]
This form is positive definite: a zero integral forces the analytic
output, and hence $G$, to vanish. Its minimum eigenvalue is positive.
The shell changes the resolvent vector by at most
$d^{-1}|d_{L,K}(G)|\sqrt{L/(2K)}=o(\|G\|)$.
Finite-dimensional compactness proves both bounds.
\end{proof}

\begin{corollary}[Divergence of the proposed endpoint ratio]
\label{cor:v11-ratio}
For the top-root pair in Proposition~\ref{prop:v1-reflected}, put
$v^\pm=V_LF_\pm$. Wherever $\delta v^+\ne0$ and $L$ is sufficiently
large, the proposed ratio obeys
\[
 \left|\frac{\delta v^-}{\delta v^+}\right|
       \ge c e^{(\nu-d)L/2}.
\]
Thus it cannot be bounded along any sequence tending to infinity on
which it is defined. At zeros of the denominator the proposed formula
for the test is undefined. The vectors
$v^- -(\delta v^-/\delta v^+)v^+$ have unbounded ordinary norm.
\end{corollary}
\begin{proof}
Genuine zero cancellation gives
\[
 H_{F_+}(p)=0,\qquad
 H_{F_-}(p)=\kappa(p)\frac{\zeta^{(m)}(p)}{m!}\ne0.
\]
The preceding endpoint lemma bounds the denominator by $Ce^{-\nu L/2}$
and the numerator below by $ce^{-dL/2}$. The inputs $F_+$ and $F_-$
are linearly independent, and Lemma~\ref{lem:v11-resolvent-gram}
then gives norm divergence of the displayed combination.
\end{proof}

\begin{theorem}[Loss of the reflected signal on a fixed test space]
\label{thm:v11-fixed-tests}
Include $F_+$ in $\mathcal F$ and let $G_L\in\mathcal F$.
Suppose $v_L=V_LG_L$ is bounded in ordinary norm and $\delta v_L=0$.
Then
\begin{equation}\label{eq:v11-signal-loss}
 \ip{WJF_+}{A_\rho^*v_L}
       =\ip{WJF_+}{JG_L}\longrightarrow0.
\end{equation}
This conclusion uses holomorphic form transfer only. No pole-aware
Weil transfer is assumed.
\end{theorem}
\begin{proof}
Lemma~\ref{lem:v11-resolvent-gram} makes $G_L$ bounded.
Equation~\eqref{eq:v11-rescaled-endpoint} and $\delta v_L=0$
imply $H_{G_L}(p)\to0$. On the other hand, the global value-level
zero formula and the top-root cancellation give
\[
 \Q(g_{F_+},g_G)=mH_{F_+}(\rho)\overline{H_G(p)}.
\]
Theorem~\ref{v1:thm:repair} transfers this identity uniformly for
bounded $G\in\mathcal F$. Substitution of $G_L$ proves the result.
Equivalently, for every bounded family of inputs,
\begin{equation}\label{eq:v11-endpoint-cross-relation}
 \ip{WJF_+}{JG_L}
   =-mH_{F_+}(\rho)\overline{\ell_L(G_L)}+o(1).
\end{equation}
\end{proof}

In particular, normalizing the divergent two-vector test does not retain
the desired nonzero cross limit. The endpoint constraint suppresses the
very reflected evaluation detected by the Weil form. This is a limitation
of a fixed finite-dimensional input space with the specified resolvent
and sampler. It does not exclude test spaces that grow or change with
$L$, nor other ways to handle the boundary channel.

\paragraph{A remaining alternative and its cost.}
A finite test can instead be corrected directly:
\[
 w_L=v^- -\delta(v^-)c_{L,K},\qquad \delta w_L=0.
\]
It is bounded and differs from $v^-$ by a norm-vanishing vector, but
$A_\rho^*w_L=x_- -\delta(v^-)A_\rho^*c_{L,K}$ is no longer the
image of a bounded input in the same fixed test space. The exact price is
\[
 \ip{WJF_+}{A_\rho^*w_L}
 =\ip{WJF_+}{x_-}
  -\overline{\delta(v^-)}\ip{WJF_+}{A_\rho^*c_{L,K}}.
\]
For this particular large shell,
$\|A_\rho^*c_{L,K}\|^2=\|Dc_{L,K}\|^2+
|\mu|^2\|c_{L,K}\|^2$ by symmetry. Consequently its graph
correction has size comparable to
$e^{-dL/2}\sqrt{K/L}$, which diverges at the specified cutoff.
A different endpoint correction or a smaller auxiliary shell inside the
same ambient space may change that cost. Either choice requires an
independent bound for the displayed Weil pairing and the complete graph
source. The present calculation establishes no such bound.


\subsection*{A parity benchmark that uses ordinary norm transfer only}
Let $\mathsf P V_n=V_{-n}$ and $P_+=(I+\mathsf P)/2$.
Then $P_+\succeq0$ and $[D,P_+]=D\mathsf P$.
For the fixed one-sided root family the even projection of the Hardy
output is injective. To see this, a vanishing even part implies
$\kappa(s)F(s)+\kappa(1-s)F(1-s)=0$ by analytic continuation.
Writing $F=\zeta R_F$ and using the functional equation gives a
meromorphic identity whose rational principal parts lie in disjoint
right and left windows. The coefficients at each pole on the right
must vanish, so $F=0$. Finite-dimensional compactness gives a positive
lower norm bound. Ordinary Riemann-sum convergence transfers it to
$\Jc$; the centered phase commutes with reflection, and the norm-small
endpoint correction preserves the bound for $J$.
Thus $P_+$ is a bounded noncommuting positive benchmark. Its exact
Lyapunov source is $P_+(ez+\mathcal RF)$, however, and the new graph
term is still present. Positivity by itself supplies no small upper bound.

\section{The arithmetic boundary observable}\label{sec:arithmetic}
\begin{proposition}[Exact arithmetic realization of the CCM boundary channel]\label{prop:beta-arithmetic-channel}
Let $L=2\log\lambda$, $t_n=2\pi n/L$, and let
$\mathfrak b_\lambda(n)$ for $n\in\mathbb Z$ denote the real odd coefficient in the CCM
rank--two commutator before physical rescaling, so that
\begin{equation}\label{eq:beta-bcoeff-expansion}
 \beta^{\log}_{\lambda,N}
 =\frac{2\pi}{L}\sum_{|n|\le N}\mathfrak b_\lambda(n)V_n.
\end{equation}
With
\[
 \rho_\infty(y)=\frac{e^{y/2}}{e^y-e^{-y}},
 \qquad
 \mathcal V(y)=2(e^{y/2}-1)
 -\sum_{2\le k\le e^y}\frac{\Lambda(k)}{\sqrt k},
 \qquad 0\le y\le L,
\]
where $\mathcal V$ is exactly the centered prime cumulative already used in
\eqref{eq:centered-prime-cumulative}, one has the finite, zero-independent
identity
\begin{align}
 \mathfrak b_\lambda(n)
 ={}&\frac{32L\sinh^2(L/4)\,n}
          {L^2+16\pi^2n^2}
 +\frac1\pi\int_0^L\sin(t_ny)\rho_\infty(y)\,dy \notag\\
 &+\frac1\pi\sum_{1<k\le e^L}
   \frac{\Lambda(k)}{\sqrt k}\sin(t_n\log k)                                      \label{eq:beta-bcoeff-direct}\\
 ={}&\boxed{\frac{t_n}{\pi}\int_0^L
       \mathcal V(y)\cos(t_ny)\,dy
 +\frac1\pi\int_0^L
   \bigl(\rho_\infty(y)-e^{-y/2}\bigr)\sin(t_ny)\,dy.}             \label{eq:beta-bcoeff-V}
\end{align}
Equivalently, if
\[
 \mathcal E_L(t)
 :=\sum_{1<k\le e^L}\frac{\Lambda(k)}{\sqrt k}\sin(t\log k)
   -\int_0^L e^{y/2}\sin(ty)\,dy,
\]
then at the lattice frequencies
\begin{equation}\label{eq:beta-bcoeff-E}
 \boxed{
 \mathfrak b_\lambda(n)
 =\frac{\mathcal E_L(t_n)}{\pi}
  +\frac1\pi\int_0^L\sin(t_ny)\rho_\infty(y)\,dy
  -\frac{4t_n(1-\lambda^{-1})}{\pi(1+4t_n^2)}.}
\end{equation}
Thus the same centered pole--prime discrepancy that drives G1 also generates
the boundary vector appearing in the graph calculation; no nontrivial-zero data occur in
\eqref{eq:beta-bcoeff-direct}--\eqref{eq:beta-bcoeff-E}.

More precisely, write an arbitrary finite vector as
\[
 v=L^{-1/2}\sum_{|n|\le N}h_nV_n,
 \qquad
 H_v(y)=\frac1L\sum_{|n|\le N}h_ne^{it_ny},
 \qquad
 A_v(y)=\overline{H_v(-y)}-\overline{H_v(y)}.
\]
Then, for the first-slot-linear convention used throughout this manuscript,
\begin{equation}\label{eq:beta-arithmetic-trace}
 \boxed{
 i\sqrt L\,\langle\beta^{\log}_{\lambda,N},v\rangle
 =\int_0^L\mathcal V(y)A_v'(y)\,dy
  +\int_0^L\bigl(\rho_\infty(y)-e^{-y/2}\bigr)A_v(y)\,dy.}
\end{equation}
Since $A_v$ is $L$-periodic and odd, one also has the exact folded form
\begin{align}
 i\sqrt L\,\langle\beta^{\log}_{\lambda,N},v\rangle
 ={}&\int_0^{L/2}
   \bigl(\mathcal V(y)+\mathcal V(L-y)\bigr)A_v'(y)\,dy \notag\\
 &+\int_0^{L/2}
 \Bigl[
  \rho_\infty(y)-e^{-y/2}
  -\rho_\infty(L-y)+e^{-(L-y)/2}
 \Bigr]A_v(y)\,dy.                                      \label{eq:beta-arithmetic-folded}
\end{align}
\end{proposition}
\begin{proof}
Connes--Consani--Moscovici show that the truncated Weil matrix has
$\tau_{ij}=(b_i-b_j)/(i-j)$ off the diagonal and that its commutator with the
integer Fourier generator is the corresponding rank--two matrix
\cite{ConnesConsaniMoscovici2025}.  Their point contribution
\[
 W_{0,2}(V_n,V_m)
 =\frac{32L\sinh^2(L/4)(L^2-16\pi^2mn)}
 {(L^2+16\pi^2m^2)(L^2+16\pi^2n^2)}
\]
therefore contributes
$32L\sinh^2(L/4)n/(L^2+16\pi^2n^2)$ to $b_n$.
For $n\ne m$, the kernel already used in Lemma~\ref{lem:g1-high-low} is
\[
 q(U_n,U_m)(y)
 =\frac{\sin(t_my)-\sin(t_ny)}{\pi(n-m)}.
\]
The signs in the semilocal Weil decomposition then give the positive prime
sine sum and the positive archimedean sine integral in
\eqref{eq:beta-bcoeff-direct}.  This proves the first identity directly from
the finite Weil matrix.

The point term also has the integral representation
\[
 \frac{32L\sinh^2(L/4)n}{L^2+16\pi^2n^2}
 =-\frac1\pi\int_0^L
   (e^{y/2}+e^{-y/2})\sin(t_ny)\,dy.
\]
In Stieltjes form,
\[
 d\mathcal V(y)=e^{y/2}\,dy
 -\sum_{1<k\le e^L}\frac{\Lambda(k)}{\sqrt k}\,\delta_{\log k}(dy).
\]
Substitution in \eqref{eq:beta-bcoeff-direct} gives
\[
 \mathfrak b_\lambda(n)
 =-\frac1\pi\int_0^L\sin(t_ny)\,d\mathcal V(y)
 +\frac1\pi\int_0^L
   (\rho_\infty(y)-e^{-y/2})\sin(t_ny)\,dy.
\]
The boundary term in Stieltjes integration by parts vanishes because
$t_nL=2\pi n$, giving \eqref{eq:beta-bcoeff-V}.  Formula
\eqref{eq:beta-bcoeff-E} follows from
\[
 \int_0^Le^{-y/2}\sin(t_ny)\,dy
 =\frac{4t_n(1-\lambda^{-1})}{1+4t_n^2}.
\]

Finally, \eqref{eq:beta-bcoeff-expansion} and the first-slot-linear inner
product give
\[
 \sqrt L\,\langle\beta^{\log}_{\lambda,N},v\rangle
 =\frac{2\pi}{L}\sum_{|n|\le N}
   \mathfrak b_\lambda(n)\overline{h_n}.
\]
But
\[
 A_v(y)=\frac{2i}{L}\sum_{|n|\le N}
 \overline{h_n}\sin(t_ny),
 \qquad
 A_v'(y)=\frac{2i}{L}\sum_{|n|\le N}
 t_n\overline{h_n}\cos(t_ny),
\]
which proves \eqref{eq:beta-arithmetic-trace}.  Periodicity gives
$A_v(L-y)=-A_v(y)$ and $A_v'(L-y)=A_v'(y)$; splitting at $L/2$ proves
\eqref{eq:beta-arithmetic-folded}.
\end{proof}
\begin{proposition}[The present absolute PNT remainder is quantitatively insufficient for the G2 boundary channel]\label{prop:pnt-beta-barrier}
Fix $T>0$.  The unconditional estimate
\eqref{eq:centered-prime-cumulative}, inserted absolutely into the exact
formula \eqref{eq:beta-bcoeff-V}, gives constants $c_T,C_T>0$ such that
\begin{equation}\label{eq:pnt-beta-barrier}
 \boxed{
 \sup_{|t_n|\le T}|\mathfrak b_\lambda(n)|
 \le C_T\bigl(1+\lambda e^{-c_T\sqrt L}\bigr).}
\end{equation}
In particular this already-audited PNT input, used only by absolute values,
does not imply that the boundary observable
$\sqrt L\langle\beta^{\log},v\rangle$ on a varying test family tends to zero.  The
extra factor $e^{-L/2}=\lambda^{-1}$ that closes the G1 pole--prime estimate
comes from the special co-Poisson cross-tail kernel and is absent from the
periodic Hardy boundary kernel in \eqref{eq:beta-arithmetic-trace}.
This records the limitation of this absolute-value estimate. It is not a lower bound on the actual error, nor a proof that signed cancellation is impossible.
\end{proposition}
\begin{proof}
The second integral in \eqref{eq:beta-bcoeff-V} is $O_T(1)$: near zero
$\rho_\infty(y)=(2y)^{-1}+O(1)$ while
$\sin(t_ny)=O_T(y)$, and for $y\ge1$ the kernel decays exponentially.
For the first integral, \eqref{eq:centered-prime-cumulative} gives
\[
 |t_n|\int_1^L|\mathcal V(y)|\,dy
 \ll_T\int_1^Le^{y/2-c\sqrt y}\,dy
 \ll_T\lambda e^{-c_T\sqrt L},
\]
after splitting at $L/4$ and decreasing $c_T$ if necessary.  The bounded
initial interval is absorbed into the constant.  This proves
\eqref{eq:pnt-beta-barrier}.
\end{proof}
\section{Analytic filters, moving phases, and the reflected pole}\label{sec:filters}
The following continuum identities remain useful diagnostics. They do
not identify the endpoint or graph behavior of the corrected finite
sampler without a further transfer argument.

\subsection*{The exact meromorphic normal form}
For an off-line zero $\rho$ of order $m$, the critical-line adjoint
resolvent multiplier is $-(s-\rho^\sharp)^{-1}$. Define
\[
 H_{v^+}(s)=-\frac{\kappa(s)\zeta(s)}{(s-\rho)^m(s-\rho^\sharp)},
 \qquad
 H_{v^-}(s)=-\frac{\kappa(s)\zeta(s)}{(s-\rho^\sharp)^{m+1}}.
\]
Their continuum values at the uncentered cut follow by the right-hand
Cauchy pole at $\sigma$. They give
\[
 \alpha_\infty=
 \left(\frac{\sigma-\rho}{\sigma-\rho^\sharp}\right)^m,
 \qquad
 H_{\widetilde v}(s)=\frac{\kappa(s)\zeta(s)}{s-\rho^\sharp}
 \left[\frac{\alpha_\infty}{(s-\rho)^m}
              -\frac1{(s-\rho^\sharp)^m}\right].
\]
The bracket vanishes at $\sigma$, canceling the Hardy Cauchy pole.
At $\rho$ the zeta zero cancels the apparent pole. At $\rho^\sharp$
there remains exactly one simple pole, with residue
\begin{equation}\label{eq:v1-pole-residue}
 -\kappa(\rho^\sharp)\frac{\zeta^{(m)}(\rho^\sharp)}{m!}\ne0.
\end{equation}
This is an algebraic statement about the uncentered continuum transform.
For the corrected finite resolvent tests, Corollary~\ref{cor:v11-ratio}
proves that the endpoint ratio instead diverges in magnitude; it cannot
tend to this finite $\alpha_\infty$.

Adding a transform holomorphic at $\rho^\sharp$ cannot change this
residue. If such a correction $w$ is cross-invisible to the top root,
then, wherever the paired form is well-defined,
\[
 QW(g_{\rho,m},A_\rho^*(c v^-+w))
       =\overline c\,C_{\rho,m}.
\]
The conjugate on $c$ is required by the first-slot-linear convention.
This corrects the opposite convention used in the previous proof.

\begin{proposition}[The bounded-filter tradeoff]\label{prop:v1-filter}
Let $\Phi_L$ be bounded and holomorphic on $\Re s>1/2$ and admit
continuation near $\rho^\sharp$. Put
$M_L=\|\Phi_L\|_{H^\infty(\Re s>1/2)}$.
If $|\Phi_L(\rho)\overline{\Phi_L(\rho^\sharp)}|\ge c_0>0$, then
\[
 |\Phi_L(\rho^\sharp)|\ge c_0/M_L.
\]
In particular residue damping by $e^{-\epsilon L}$ while retaining
that product forces $M_L\ge c_0e^{\epsilon L}$.
\end{proposition}
\begin{proof}
Since $\rho$ is in the right half-plane, $|\Phi_L(\rho)|\le M_L$.
The conclusion follows by dividing the product lower bound by $M_L$.
\end{proof}
To interpret the product as a global Weil cross form, impose in addition
global transform admissibility and the strip/tail estimates required for
that form and its transfer. Local continuation at two points is not
sufficient. A boundary-annihilation ratio also requires
$\Phi_L(\sigma)\ne0$. Under these extra conditions the cross constant
is multiplied by $\Phi_L(\rho)\overline{\Phi_L(\rho^\sharp)}$ and
the residue by $\Phi_L(\rho^\sharp)$.

The absolute PNT envelope applied to a reflected image of size
$|\Phi_L(\rho^\sharp)|e^{-dL}$ yields an upper estimate with scale
$|\Phi_L(\rho^\sharp)|e^{(1/2-d)L-c\sqrt L}$.
Making this particular estimate tend to zero by amplitude damping
alone calls for exponentially small residue and hence an exponentially
large available Hardy norm. The existing G1 estimate pays uniform
Sobolev constants and does not cover that regime. This is a limitation
of this proof strategy. A large upper bound on an error does not prove
that the actual error is large, and the calculation does not exclude
signed cancellation or a new estimate that avoids that norm cost.

\subsection*{Moving phases and the omitted periodic image}
For real $\tau$ the multiplier
$H_F^{(\tau)}(s)=e^{-\tau(s-1/2)}H_F(s)$ translates the inverse
transform to $g_F(y-\tau)$. Global cross forms are invariant, because
the two reflected factors cancel. The uncentered Cauchy value at zero
is $e^{-q\tau}\mathscr E(F)$ for $\tau>0$.
These are exact continuum statements, and the coefficientwise graph
identity changes its source by the same phase.

The earlier claimed periodic kernel bound proportional to $e^{-\tau/2}$
omitted the nearest reflected-pole image. A pole at real displacement
$-d$, $0<d<1/2$, contributes a term on the scale
$e^{-d(L-\tau)}$ at the opposite end of the period. For example,
$d=0.1$ and $\tau=0.26L$ give $e^{-0.074L}$, larger than
$e^{-0.13L}$. Restricting $\tau<L/2$ does not repair that comparison.

The historical threshold comparison
\[
 \frac{d}{d+q}<\frac1{2q+1}\quad\Longleftrightarrow\quad d<1/2
 \qquad(q>0)
\]
is algebraically correct. It proves a conflict only if both proposed
asymptotic thresholds have first been justified, with remainders small
relative to the relevant periodic image. An absolute $O(e^{-cL})$
remainder is not enough if the image decays faster. The previous
unconditional exclusion of the whole moving-phase route is therefore
withdrawn. Moving $\sigma$ likewise requires uniform smoothing, cutoff,
and transfer estimates; the scalar relation
$\sigma F(\sigma)\to\eta_{\mathrm B}(F)$ on a fixed root space is
not such a uniform theorem.

\section{Candidate changes, exact matching, and their scope}\label{sec:matching}
\subsection*{Projective candidate invariance is an exact finite identity}
For any candidate $k$ with $a=\eta_\partial(k)\ne0$, parity is not
needed for the full commutator identity
\[
 W\frac{Dk}{a}=\frac{DWk}{a}-\beta
             +\frac{\beta^*(k)}{a}\eta_\partial.
\]
When paired with a fixed test $v$ having $\eta_\partial(v)=0$, the
last term vanishes. For two such candidates,
\[
 \ip{W(Dk_1/a_1-Dk_2/a_2)}{v}
 =\ip{DWk_1/a_1-DWk_2/a_2}{v}.
\]
Consequently a candidate change cannot alter that fixed paired mismatch
when both normalized G1 pairings tend to zero. This does not assert a
nonzero limit, and it does not cover changes to the tests, matrix, or
sampler. Replacing the prolate candidate by an endpoint shell preserves
this exact identity but does not import a disputed G2 lower limit.

\subsection*{An exactly matching sampler}
On an algebraic one-sided root space put
\[
 R_F(s)=F(s)/\zeta(s),\qquad
 R_{BF}(s)=(s-1/2)R_F(s)-\eta_{\mathrm B}(F).
\]
For an even finite candidate $k$ with physical endpoint $B_k\ne0$,
let $z^{(k)}=iDk/B_k$ and define
\begin{equation}\label{eq:v1-exactmatcher}
 [J^{(k)}F]_n=[z^{(k)}]_nR_F(s_n).
\end{equation}
Multiplication of the rational graph equation by $[z^{(k)}]_n$ proves
\[
 iDJ^{(k)}-J^{(k)}B=z^{(k)}\eta_{\mathrm B}.
\]
With $x=J^{(k)}F$ and $q=i\eta_\partial(x)/\eta_\partial(k)$,
\[
 \eta_{\mathrm B}(F)z^{(k)}-qDk
 =\frac{iDk}{B_k}\bigl(\eta_{\mathrm B}(F)-\delta(x)\bigr).
\]
This removes a vector mismatch algebraically. Endpoint recovery is only
one remaining requirement: injectivity, uniform ordinary norm bounds,
Weil transfer, and weak G1 on these tests also need proof. No coercivity
follows from the displayed coefficient identity alone.

\subsection*{Endpoint recovery is a periodic Laplace problem}
Let $k$ be a trigonometric polynomial on $[0,L]$, with
$k(0)=k(L)=B_k$, and $\Re\alpha>0$. Put
\[
 E_{L,k}(\alpha)=\frac{\alpha}{B_k(1-e^{-\alpha L})}
                   \int_0^Le^{-\alpha y}k(y)\,dy.
\]
Solving $(\partial_y-\alpha)x=k'/B_k$ with periodic boundary condition
gives $x(0)=1-E_{L,k}(\alpha)$: multiply the equation by
$e^{-\alpha y}$, integrate, and integrate its right side by parts.
Since $\partial_\alpha(\partial_y-\alpha)^{-1}
=(\partial_y-\alpha)^{-2}$, higher Jordan coordinates satisfy
\[
 \eta_{\mathrm B}(F_{\rho,j})-\delta(J^{(k)}F_{\rho,j})
 =\frac{\partial_\alpha^{j-1}E_{L,k}(\alpha)}{(j-1)!},
 \qquad \rho=1/2+\alpha.
\]
This formula is exact for each finite candidate.

For the endpoint shell $k=c_{L,K}$ of \eqref{v1:eq:shell}, the
coefficient of the first adapted root is
$\sqrt L(2K)^{-1}it_n/(it_n-\alpha)$ on its $2K$ modes.
If $K/L\to\infty$, its physical endpoint tends to one while its squared
ordinary norm is $O(L/K)\to0$. Higher root endpoints have the corresponding
vanishing limits and their ordinary masses also vanish. Thus near-perfect
endpoint matching can coexist with collapse of the whole sampled norm.
The shell exact matcher is not a coercive transport.

\subsection*{Global divisibility and endpoint-jet expansions}
For a smooth exact-moment source $h$ supported up to $\lambda$, the
Mellin identity is $M_0(Eh)=\zeta M_0h$. Set
$a=\log\lambda$, $\varphi(y)=e^{y/2}h(e^y)$ for $y\le a$, and
$B^+=\varphi(a)\ne0$. With the positive Mellin convention used in
the Sonine sections,
\[
 it\,M_0h(1/2+it)/B^+
 =e^{ita}-\frac1{B^+}\int_{-\infty}^a\varphi'(y)e^{ity}\,dy.
\]
This is integration by parts. The negative Fourier convention of
Section~\ref{sec:sampling} represents the same formula after $y\mapsto-y$;
the signs must be changed together. If $\varphi'\in L^1$, the bulk
integral tends to zero for fixed $\lambda$ and $|t|\to\infty$.
Thus the global zeta quotient cancels but a nondecaying endpoint character
remains. If $\varphi'\in L^2$, Plancherel gives bulk squared norm
$\|\varphi'\|_2^2/|B^+|^2$.

For $r\ge1$, assuming one-sided derivatives, vanishing lower traces,
and $\varphi^{(r)}\in L^1$, repeated integration gives
\[
 it\,M_0h(1/2+it)/B^+
 =e^{ita}\sum_{j=0}^{r-1}
   \frac{(-1)^j\varphi^{(j)}(a)}{B^+(it)^j}
   +O\!\left(\frac{\|\varphi^{(r)}\|_1}
                       {|B^+||t|^{r-1}}\right).
\]
This is a fixed-cutoff high-frequency expansion. The terms $(it)^{-j}$
cannot be subtracted through $t=0$ without a low-frequency prescription.
On the lattice $t_n=2\pi n/L$, $L=2a$, the character is $(-1)^n$.
Its appearance is consistent with the centered/unphased distinction and
does not license treating the two samplers as one vector.

\subsection*{Shell estimates and the limited scalar-shift obstruction}
The corrected logarithmic diagonal is sufficient for shell weak-pairing
estimates. For the explicit shell and a bounded ordinary-norm test,
\[
 |\ip{Wc_{L,K}}{v}|
 \le C\lambda(1+\log(4K+1))\sqrt{L/K}\,\|v\|.
\]
For a Hardy test with adequate weighted coefficient decay, the weighted
convolution estimate controls $\ip{DWc}{v}=\ip{Wc}{Dv}$ similarly.
One may enlarge a common ambient cutoff to make these shell pairings
small, but \eqref{v1:eq:paireddefect} still contains its separate beta
term. The old shell calculation is not a proof of graph compatibility.

Finally a scalar-shift shortcut has a limited, elementary obstruction.
If on the same normalized candidate a full residual is $o(1)$, while a
proposed shift has $|\alpha|\ge\alpha_0>0$ and normalized derivative
norm bounded below, the reverse triangle inequality prevents the shifted
full residual from also tending to zero. Weak fixed-band G1 supplies
neither of those full-norm premises. The exact cancellation in
Proposition~\ref{prop:v1-candidate} is the appropriate statement for the
weaker paired argument.

\section{The remaining fixed-space criterion and research goals}\label{sec:goals}
The original relative boundary criterion is unchanged. If, for every
fixed nontrivial zero and prime $p$,
\[
 \frac{\|P_RD_p(v_\rho^{pR}-v_\rho^R)\|}{\|v_\rho^R\|}
 \longrightarrow0,
\]
then contractivity of $T_{R,p}$ and \eqref{eq:boundary} give
$|p^{1/2-\rho}|\le1$. Reflection gives RH. The critical-line rate
derived from Assumption~\ref{ass:evaluators} is consistent with this
criterion; it does not establish the off-line estimate. Equivalently,
the exact diagonal defect asks for an independent arithmetic sign on
each zero evaluator. Full positivity on their span is an unnecessarily
strong sufficient condition.
Proposition~\ref{prop:v110-normalized-shell} now removes a technical
prerequisite from this alternative: the Fredholm remainder and its signed
moment correction vanish after evaluator normalization without invoking
Assumption~\ref{ass:evaluators}. The precise outstanding target is
\eqref{eq:v110-sign-target} at the single prime $2$. The elementary
off-line upper bound suffices for this reduction; a new arithmetic
argument must still orient the leading shell imbalance. The reduction
does not infer that sign from the already known prime eigenvalue.
Proposition~\ref{prop:v111-adjoint} now writes the arithmetic constraint
as the local distributional equation $\mathscr Ay=A_y+B_y/t$ on the
expanded source interval. The integral-corrected cutoff converges with
the explicit bound in \eqref{eq:v111-adjoint-error}; the raw cutoff can
instead diverge in norm. Corollary~\ref{cor:v111-gram} justifies fixed
finite source-Gram calculations, not inversion of a limiting source
Gram operator. The next step is a signed quadratic estimate exploiting
this equation together with both Sonine support conditions. Neither
linear source orthogonality alone nor the power response
\eqref{eq:v111-power-response} supplies that estimate.

The weak G1 source chain is now closed: Theorem~\ref{thm:v14-radical}
supplies compact-source radicality and the semilocal domain passage;
Proposition~\ref{prop:v12-source} and Corollary~\ref{cor:v12-repair}
give the normalized source and exact repair; the subsequent tail,
endpoint and correlation estimates give \eqref{eq:G1-continuum} and
\eqref{eq:G1-sobolev}. The deterministic finite transfer in
Theorem~\ref{thm:g1-boundary-corrected} is consequently unconditional
for this source and a fixed physical test band.

Three distinct research tasks remain for the finite-Weil route. First,
seek genuinely uniform estimates for a growing deep prolate block, or
the separately stated continuous-remainder improvement needed for an
economical near-Slepian cutoff. Neither is required for the proved weak
G1 limit, and neither follows from its fixed-source domain passage.
Second, estimate the
corrected graph pairings on a new, fully defined reflected test family.
The finite object is fixed by \eqref{v1:eq:repair}, but the former
bounded endpoint ratio is ruled out by Corollary~\ref{cor:v11-ratio}.
Theorem~\ref{thm:v11-fixed-tests} further shows that bounded resolvent
tests from any fixed finite-dimensional input space lose the reflected
cross signal under exact endpoint annihilation. A viable replacement must
address that obstruction, for example through a changing test space or
a direct endpoint correction with its additional Weil pairing controlled.
For such a replacement, the full graph source, the relevant meromorphic
transform errors, and any entire zero-sum estimate remain obligations.
Third, prove a positive
metric comparison whose Lyapunov defect tends to zero independently of
the assumed off-line roots. Restricted scalar coercivity is already a
consequence of form transfer and cannot serve as that independent input.
Section 19 now sharpens this task: Theorem~\ref{thm:v15-finite-metric}
supplies a finite-core criterion, and Proposition~\ref{prop:v15-projector}
constructs a positive, uniformly coercive metric that kills the zeta source
and endpoint-shell graph terms exactly. Its remaining commutator has the
strictly positive variance limit in Proposition~\ref{prop:v15-variance};
no independent vanishing signed pairing follows. Removing more source
derivatives by bounded projections loses coercivity by
Proposition~\ref{prop:v15-krylov}. A sharp fixed-band Weil correction
faces a separate transfer cost: Proposition~\ref{prop:v16-projected-cut}
proves exponential off-line transform growth after source projection.
Proposition~\ref{prop:v17-pair-edge} computes the reflected-pair term:
unless its explicit endpoint product vanishes, it remains exponentially
large and oscillatory. Any proposed cancellation must therefore be checked
over the entire zero sum or on the actual finite prime/pole/archimedean
form; bounded ordinary norms and pairwise reflection do not supply it.
There is now a transfer-preserving alternative:
Proposition~\ref{prop:v18-full-projection} removes the source and the
two shell directions on the full expanding Fourier space. It proves a
uniform summable strip-product estimate, fixed-core coercivity and the
unchanged global Weil pullback. The next metric target is its remaining
signed commutator, or an independently justified arithmetic correction
to it. Reusing the same projected Weil form plus a fixed scalar is
excluded by Corollary~\ref{cor:v18-shift}: the Weil pullback tends to
zero on the one-sided root space and cannot cancel the nonzero scalar
metric budget. A proposed correction must be tested against the full
relative criterion before any positivity claim is made.
Proposition~\ref{prop:v19-action} now supplies a stronger diagnostic:
at tail rate $p>1$, the actual finite Weil action converges in norm to
the projected finite exponential sum of zero values. Consequently
the positive square $SW^2S$ has the top-root growth and lower-jet
collapse in Corollary~\ref{cor:v19-square}; it does not furnish the
missing signed estimate. Further metrics should be checked for both
the relative defect and multiplicity-sensitive coercivity, with any
scale change accounted for explicitly.

The fixed-space route has separate foundational obligations: prove the
quantitative evaluator estimates and give a complete closed-operator
realization if the Riesz-window formulation is retained. These are not
prerequisites for the finite-core exclusion criterion in
Theorem~\ref{thm:v15-finite-metric}, nor for the direct complete-Weil
route below. Bypassing them does not establish them for the Riesz
formulation. The shell-sign, uniform metric, large-support
spectral-overlap, and Schur-complement formulations are alternative
targets; none has been shown to follow from weak G1.

\begin{proposition}[The direct cofinal Weil route]\label{prop:v121-cofinal-rh}
Suppose $\lambda_j\to\infty$ and the complete, unshifted semilocal
Weil forms satisfy
\[
 QW_{\lambda_j}(f,f)\ge-\epsilon_j\norm f^2
 \quad(f\in\mathcal D(QW_{\lambda_j})),
 \qquad \epsilon_j\ge0,\quad\epsilon_j\to0.
\]
Then the Weil form is nonnegative on every smooth compactly supported
test function, and RH follows by Weil's positivity criterion.
\end{proposition}
\begin{proof}
Fix such a test $f$. For all sufficiently large $j$, its multiplicative
support lies in $[\lambda_j^{-1},\lambda_j]$. Extension by zero
preserves both its ordinary norm and the global Weil value, so
\[
 QW(f,f)=QW_{\lambda_j}(f,f)\ge-\epsilon_j\norm f^2.
\]
Sending $j$ to infinity proves nonnegativity. Apply the compact-support
Weil criterion recalled in \cite{ConnesConsaniMoscovici2025}.
\end{proof}

Consequently the full cofinal conclusion of
Proposition~\ref{prop:v120-structured-uniform}, if its hypotheses and
vanishing error are proved, would already yield RH by this direct route.
The conditional Riesz realization and evaluator assumptions belong to
other routes and are not additional premises of
Proposition~\ref{prop:v121-cofinal-rh}. A fixed-window certificate,
a bound on a finite compression alone, or an incomplete transport
estimate does not meet this proposition. The implication isolates the
strength of the missing estimate; it does not establish that estimate.

% NS-17: determinacy_insert.tex (qualified integration)
The gap isolated by Proposition~\ref{prop:v121-cofinal-rh} can be
viewed as a continuation and identification problem. In the screw-function
realization of Suzuki~\cite{Suzuki2026}, let $g$ be the Weil screw
function, $a=\log\lambda$, and $\lambda_a=\inf\sigma(A_a)$ as in
\eqref{eq:g2-monotone}. A certified window
$\mathsf W_\lambda\succeq0$ gives $\lambda_a\ge0$ and the associated
local positive kernel. The continuation theory of
Kre\u{\i}n--Langer~\cite{KreinLanger2014} distinguishes the existence of
positive extensions of local data from their uniqueness and from the
identification of a prescribed extension. Here the prescribed global
function is the arithmetic function $g$; Suzuki's global screw-function
criterion identifies its positivity with RH. A local certificate does
not establish that global identification. In particular we do not infer
that RH is equivalent to uniqueness of an extension of one fixed-window
kernel, or that uniqueness alone identifies the extension with $g$.
Any determinacy route must state and prove that additional identification.

Truncated moment problems and the Toeplitz--Hankel determinant theory
of Basor--Ehrhardt~\cite{BasorEhrhardt2002} provide a useful comparison,
not a further positivity theorem for this operator. The Fourier-basis
split in \texttt{evidence/diag\_circle\_split/} includes divided-difference
commutator terms; it is not an identification with the standard
single-symbol matrices $T(\phi)\pm H(\phi)$ covered by that reference.
Those numerical comparisons remain diagnostics. The cofinal estimate in
Proposition~\ref{prop:v121-cofinal-rh} is still a sufficient arithmetic
input of RH strength, not a consequence of extension terminology or a
finite list of positive sections. No G2 gap or RH claim follows.




Proposition~\ref{prop:v112-schur-angle} isolates an inverse-weighted
residual that can be far smaller than the full residual divided by the
smallest gap. Proposition~\ref{prop:v112-finite-certificate} certifies
one finite case. Proposition~\ref{prop:v113-source-certificate} now
identifies the exact projected source there, including its angular tail,
and proves a relative physical-endpoint mismatch greater than one.
Proposition~\ref{prop:v114-polynomial-endpoint} now recovers the exact
source endpoint, with relative error $O(\lambda^{-1})$ at the unchanged
polynomial Fourier cutoff. Corollary~\ref{cor:v114-polynomial-g1}
puts weak G1 on this same finite vector, uniformly over the ambient
Fourier cutoff in the weighted test norm. These results supply no
practical certified threshold
at $\lambda=3$, and no endpoint comparison to an actual ground vector.
Proposition~\ref{prop:v114-weil-core} closes the semilocal Weil
operator-core issue and establishes fixed-window graph convergence for
BV sources. This is distinct from the open arithmetic realization in
Section~18. Proposition~\ref{prop:v115-sharp-tail} sharpens the omitted-tail
estimate to $N=256$ at $\lambda=3$. Proposition~\ref{prop:v115-absolute-residual}
additionally proves a small full ordinary residual and Rayleigh value
for the actual source and its polynomial projection. It gives no
endpoint-normalized operator estimate and no lowest-state overlap.
Proposition~\ref{prop:v116-window-positive} now closes the signed
head/complement test at $\lambda=3$ for both parity sectors and the
entire closed form. Positive weighted prime bounds and the increasing
tail diagonal are essential to its verified inverse-action estimate;
the residual's directional moments control every omitted row.
Proposition~\ref{prop:v117-ground-order} now proves that the complete
operator at this window has a simple even ground, with
$0<\mu_0<3.644\cdot10^{-38}$ and $\mu_1>10^{-36}$.
Corollary~\ref{cor:v117-ground-overlap} bounds the ordinary sine angle
of the exact normalized $P_{64}p_3$ by $0.01931$. The stronger
$10^{-34}$ threshold for all other even states is essential here.
The earlier $N=64$ finite-matrix angle bound and this complete-operator
bound are different statements.

The remaining G2 spectral task is a complete lower bound tending to
zero along unbounded windows, or sufficient quantitative source-to-ground
overlap on such a family. Proposition~\ref{prop:v121-cofinal-rh}
shows why this is RH-strength: even a cofinal sequence of complete
lower bounds $-o(1)$ implies nonnegativity of every compact test.
A finite list of certified windows cannot replace that estimate.
Proposition~\ref{prop:v118-gap-free} removes an unnecessary positive-gap
requirement from one route. The already proved rank-one source residual
reduces the target to asymptotic nonnegativity of its entire orthogonal
complement. A negative compact test would persist in that complement;
the required arithmetic sign is still unproved. Growing source blocks
also need an independent uniform operator-residual bound.
Proposition~\ref{prop:v121-complement-cancellation} rules out the
claim that substantial cancellation is confined to the single source:
an exactly source-orthogonal finite test has non-prime energy greater
than $0.4$ but positive total energy below $3\cdot10^{-31}$.
This does not disprove complementary positivity; it prevents using
basis or random-vector samples as a uniform dominance argument.

The unsigned-prime shortcut is now excluded by
Propositions~\ref{prop:v119-prime-essential}--\ref{prop:v119-prime-scale}:
every finite-dimensional source or Fourier-head removal retains the
full prime norm, asymptotic to $\lambda$. The rank-two pole term cannot
make their combined ordinary operator norm smaller either. The next
estimate must therefore use the archimedean energy, directional inverse
weights or signed correlations; enlarging a finite source block alone
cannot make that unweighted arithmetic norm small. This is an
obstruction to a sufficient norm bound, not to full Weil positivity.

Proposition~\ref{prop:v120-lambda4-tail} now proves a signed lower
bound on the entire $\lambda=4$ tail $|n|>16$, retaining its exact
archimedean diagonal. Both complex parity sectors are included.
The effective head $F-B^*T^{-1}B$ has $33$ dimensions.
Proposition~\ref{prop:v125-even-complete} certifies its entire
$17$-dimensional even sector and hence the complete even form.
Proposition~\ref{prop:v126-full-window} now completes the
$16$-dimensional odd sector and proves $\varepsilon_4=0$.
The remaining target is complete ordinary-error control along an
unbounded window family, not another finite compression at this
already certified window.
Corollary~\ref{cor:v126-local-error} gives a refined sufficient
target: bound the complementary correction below its trial energy,
then control the unresolved block with the physical factor
$\norm{K_\lambda^{1/2}P_{E_\lambda}V_\lambda^{-1}}^2$
retained. Its error convergence is not established. Proposition~\ref{prop:v121-alpha-zero} now certifies the
raw head and specified finite-support trial forms $K_X$ positive in
both parity sectors. This gives local $\alpha_4=0$, without an
infinite-tail transfer for $K_X$. It still does not suffice for the
exact Schur sign because the nonnegative residual inverse correction
must be subtracted, as in \eqref{eq:v121-galerkin-defect}.
This leaves a separate growing-window target: a corresponding ordinary
lower error tending to zero. Proposition~\ref{prop:v120-weighted-error}
shows that a dimensionless weighted error must be multiplied by its
arithmetic scale before it can serve that purpose.

Proposition~\ref{prop:v120-structured-inverse} now reuses the signed
factorization to bound the actual inverse-tail correction, including
all remote rows through \eqref{eq:v120-structured-finite-rows}.
Proposition~\ref{prop:v120-structured-uniform} identifies the sufficient
ordinary error $\alpha_\lambda+e_\lambda^2+
 t_\lambda^2/\mu_\lambda$. Its operator norms concern the entire
growing head; individual source or column bounds do not suffice.
Proposition~\ref{prop:v120-inverse-polynomial} additionally gives
convergent upper bounds for the complete far inverse, with certified
constants at $\lambda=4$, sharpened to $20$ and $90$ by
Proposition~\ref{prop:v122-parity-inverse}. The complete far energy
now fits below the trial budget on one certified head direction,
by Proposition~\ref{prop:v122-directional-far}. This is not the full
inverse correction. Proposition~\ref{prop:v123-first-majorant-fails}
now proves that its first-polynomial scalar-head bound fails after
the complete mixed term is included. Corollary~\ref{cor:v123-updated-first-fails}
also excludes its exact positive head update with the saved initial
certificate. Proposition~\ref{prop:v124-positive-direction}
proves a complete positive value on the saved head direction.
Proposition~\ref{prop:v125-even-complete} extends this to the entire
even head through frozen energy-scaled trials and a simultaneous
complete Gram bound. Proposition~\ref{prop:v126-full-window}
completes the odd sector by a different finite trial in one
certified direction and the whole complementary bound. Thus
$\mathsf W_4\succeq0$ exactly at this window.
Proposition~\ref{prop:v140-ground4} additionally proves its simple
even ground, and Corollary~\ref{cor:v140-real-zeros} gives the
real-zero ground transform. Proposition~\ref{prop:v142-ground-zero}
locates one simple complete-ground zero within $1/10$ of $\gamma_1$,
without identifying it with a Riemann zero. The primary
next target is the cofinal ordinary-error estimate, rather than a
longer finite list of certified windows.

Proposition~\ref{prop:v125-cutoff-cost} identifies a structural cost:
with the present scalar far majorant, its inner cutoff must grow
exponentially in $\lambda$. The fresh threshold table gives the
window dependence explicitly. A viable family must either control
the associated head dimensions, conditioning and residual Grams at
that cost, or improve the signed far estimate before constructing
the growing inner heads. Finite-only scaling diagnostics in the
research log do not provide the missing complete witnesses.

The historical rigid transfer test reports failure of the unchanged
$10^{-8}D$ tail metric at $\lambda=5$ in both parities; its original
witnesses and verifiers are not archived
(Proposition~\ref{prop:v127-metric-fails}). Its complete certificate
margin is unavailable. Conditional on the reported counterwitness
bounds, changing remote cutoffs alone cannot repair that prerequisite. A successful head or metric adaptation rule
remains to be proved. Proposition~\ref{prop:v127-rank-barrier} rules
out polynomial-rank repairs of the old unsigned diagonal comparison;
Proposition~\ref{prop:v127-block-metric} identifies a genuinely signed
infinite-block alternative, without claiming its hypotheses hold.
The actual $\lambda=5$, $N=\lambda^2$ test in
Proposition~\ref{prop:v128-dyadic-fails} reports failure of this dyadic
partition under every admissible smaller block metric and every
positive scalar row reweighting, but its original evidence is not archived.
The reported finite norm-comparison obstruction would settle that failure; it does not settle the Weil sign.
The physical G2.6 spaces now require finite image Grams, since
Proposition~\ref{prop:v127-nonclosable} excludes bounded transport
through the fixed point-value repair. The precise remaining ordinary
complementary target is \eqref{eq:v127-weighted-complement}.

The trial-head hypothesis is not yet verified along an unbounded
window family. With $\alpha_\lambda=0$, the remaining target is
$e_\lambda\to0$ and $t_\lambda/\sqrt{\mu_\lambda}\to0$, retaining
all ordinary-energy and dimension factors. Inverse approximation
convergence alone does not establish those limits. The complete
sign at one window closes a local obligation, not G2 or RH.

The successful one-sided bound does not require a norm contraction:
Proposition~\ref{prop:v120-norm-counterexample} certifies that the
latter fails at the very same cutoff. Nor can the whole weighted
remainder be estimated by an infinite Hilbert--Schmidt sum, by
Proposition~\ref{prop:v120-not-schatten}. Finite-column residual Grams
remain valid. The logarithmic weighted-tail asymptotic in
Proposition~\ref{prop:v120-weighted-tail-asymptotic} is strictly a
fixed-window statement and supplies no uniform joint-limit obstruction.

The physical continuum-symbol alternative is now sharply delimited.
Proposition~\ref{prop:v135-both-parities} extends its exact representation
to both parities; the odd weighted sign estimate remains an independent
block of the same target.  Proposition~\ref{prop:v135-growing-radical}
proves a growing complete near-radical block, and
Corollary~\ref{cor:v135-coercivity-obstruction} excludes positive
polynomial gaps on the specified complements.  Neither result fixes the
sign of any near-zero eigenvalue.
The sharper reduction in Corollary~\ref{cor:v136-complement-floor}
now makes a uniform $O(1)$ signed lower error sufficient along a cofinal
family; $o(1)$ remains sufficient but is not necessary for that proof.
This applies only to a complete complementary form with the established
small operator residual, including both parities. The arithmetic bound
is still unproved. Ordinary strong-resolvent convergence cannot supply it:
Proposition~\ref{prop:v136-resolvent-collapse} proves that convergence to
zero unconditionally. Proposition~\ref{prop:v136-positive-comparison}
also rules out a nonzero bounded positive comparison persisting strongly
on the source complement with vanishing error. It does not exclude
moving metrics, signed estimates, or other Hilbert-space realizations.

The scalar primitive and scalar minimum cannot supply that finite
constant: Corollary~\ref{cor:v137-scalar-no-go} proves their cofinal
divergence for the exact symbol. This excludes those simplifications,
not the signed concentration operator or all matrix-valued comparisons.
The remaining estimate must retain the physical correlations between
favorable and unfavorable levels.

The original even-sector reduction remains as follows.
Propositions~\ref{prop:v130-continuum-cancellation} and
\ref{prop:v130-remainder} reduce it to the exact arithmetic symbol
$\beta_a$ on the Paley--Wiener image of the physical window. The live
project-new mechanism is the signed source-compressed weighted
concentration inequality \eqref{eq:v131-weighted-concentration}.
Proposition~\ref{prop:v132-schatten} sharpens its negative-only scalar
consequence but proves, by an exact scale-invariant model, that no such
negative-only Schatten coefficient is forced to vanish even under
complete positivity and an exact null source.
Proposition~\ref{prop:v132-flat-top}
also excludes nontrivial positive form-exact smoothing. Thus the next
lemma must retain favorable and unfavorable arithmetic levels jointly;
neither a second-moment estimate nor a renamed one-negative-square
claim supplies it. The cofinal signed inequality and the odd sector
remain open.

% NS-17: no_prior_art_note.tex
For four of these closures the recorded search of the arXiv and journal literature
as of 2026-09-21 found no specific published prior art: the failure of
norm contraction (Proposition~\ref{prop:v120-norm-counterexample}), the
Schatten and Hilbert--Schmidt no-go
(Propositions~\ref{prop:v120-not-schatten} and~\ref{prop:v132-schatten}),
the failure of reciprocal-band approximate commutation
(Proposition~\ref{prop:v133-noncommute}), and the unconditional
divergence of the scalar signed-primitive budget
(Corollary~\ref{cor:v137-scalar-no-go}); the cutoff cost of
Proposition~\ref{prop:v125-cutoff-cost}, by contrast, belongs to a
published family and is cited as such there.

\paragraph{Diagnostic concentration geometry.}
The validated-symbol experiments in
\nolinkurl{evidence/diag_true_symbol/} concern finite-compression ground
vectors and midpoint quadrature, not certified complete-ground densities.
At $\lambda=3,4$ they place approximately half the Fourier mass on
$\{\beta_a<0\}$, with negative and positive level energies of about
$\mp0.090$ and $\mp0.226$. The quadrature resolves their cancellation
only to about $10^{-16}$; the much smaller $10^{-38}$ and $10^{-75}$
scales are the separately computed finite Rayleigh values.
The plotted mass beyond $\xi\simeq20$ is negligible at the diagnostic
resolution, not proved zero. In lattice units $u=\xi L/(2\pi)$, the
observed alternating symbol train and the ground's outward spread
with its first nearby zero at $u\simeq\gamma_1L/(2\pi)$ are diagnostic
context for a future signed concentration estimate. The zero-frequency
value is the archimedean constant minus the weighted strict-cutoff
remainder $2[\sum_{n<\lambda^2}\Lambda(n)/\sqrt n-2(\lambda-1)]$,
as in \eqref{eq:v146-amplitude}; it is not determined by
$\psi(\lambda^2)-\lambda^2$, nor an exact amplitude for every lattice
cell. These observations supply no bound on the source-orthogonal
complement and trigger no new experiment here.


The continuum endpoint is now settled by
Theorem~\ref{thm:v118-eigen-endpoint-zero}: every complete eigenfunction
has continuous zero physical trace. Hence a nonzero endpoint comparison
to the repaired source is false, with exact relative mismatch one.
Corollary~\ref{cor:v118-ground-fejer} supplies a fixed-window filtered
endpoint-decay rate. It gives no sharp-cut or joint growing-window rate
for finite-compression ground vectors. Those finite endpoint estimates,
the sampler graph defect and full-strip normalization remain separate.
The next spectral work should retain these distinctions while seeking
an independent signed complementary lower bound, rather than imposing
a nonzero continuum ground trace. G2 and RH remain open.


\paragraph{What is established by the sampler correction.}
The uncentered full-scale transfer in the previous version cannot be used
under the stated nonzero-endpoint hypotheses. Centering repairs the form
transfer; an explicit shell then restores the original physical endpoint
on the same finite matrix. Ordinary norm coercivity and the reflected
top-root cross form survive. The graph identity gains an explicit defect
with unavoidable derivative growth. Strong graph convergence fails for
norm-vanishing endpoint repairs, but that fact alone does not settle the
weaker arithmetic pairing needed for the paper's objective. The new
resolvent calculation identifies a second obstruction: on a fixed test
space the rescaled endpoint converges to minus the reflected evaluation,
which is exactly the evaluation needed for the nonzero cross form.
Killing that endpoint with bounded coefficients therefore kills that
particular G2 signal.

\paragraph{Claim discipline.}
This manuscript claims neither RH nor a completed positivity transfer.
The pole-aware transfer and the earlier nonzero critical beta limit have
been removed from the established chain. The moving-phase exclusion is
restricted to its proved algebra, pending relative alias estimates.
Unspecified earlier source-quasimode, partner-blind, and certified
nonorthogonality calculations are not used as theorems. Numerical checks
below include exploratory illustrations, a rigorous finite-matrix
certificate and its exact-source identification, complete fixed-window
sign certificates, finite-trial positivity and an exact-source
cancellation certificate. These do not settle the remaining uniform
analytic obligations. A publication draft would need the listed inputs
proved or its scope narrowed to the results that are already supported.

\section{Signed comparison obstructions and complete-space selection (v1.52)}
\label{sec:v152-route-selection}
The new material has three classifications: exact identities and proved
implications with stated hypotheses; two certified computations at an
already studied window; and explicit open inputs. No diagnostic is used
as a lower bound. A blocked parent route is not declared impossible
merely because one of its sufficient comparisons fails.

The direct uniform ordinary lower floor remains the weakest sufficient
cofinal target established here. CAE and unrestricted exact adaptive
concentration express that target in other coordinates. The additional
CCM problem is selection of a specific complete ground profile, with
entire-transform control. The following results remove specified
shortcuts and identify the estimates still missing; they do not supply
a new cofinal signed floor, G2 or RH.



\subsection*{Fresh certificates for two existing comparison failures}

The following computations use the existing $\lambda=5$ window and
the already stipulated metrics and partition. They introduce neither a
new window nor a replacement tail metric. The frozen vectors are new;
the original v1.27--v1.28 evidence is still not archived. In particular,
the historical displayed witness ratios and four-block tables are not
claimed to have been recovered or replayed.

\begin{proposition}[Fresh failure of the stipulated $10^{-8}D$ tail comparison]
\label{prop:v152-tail-comparison}
At $\lambda=5$, in each parity there is an exact dyadic vector supported
on $17\le n\le128$ with strictly positive Weil energy, but with
\[
 0<\frac{v^*\mathsf W_5v}{v^*Dv}
 <\begin{cases}7\cdot10^{-18},&\text{even},\\
               12\cdot10^{-18},&\text{odd}.
   \end{cases}
\]
Here $D=\operatorname{diag}(a_n)$ is the stipulated positive
archimedean metric, not the complete diagonal of $\mathsf W_5$.
Consequently $Q_{16}\mathsf W_5Q_{16}\succeq10^{-8}Q_{16}DQ_{16}$
is false in both parities. No negative Weil direction is asserted.
\end{proposition}
\begin{proof}[Certified verification]
The integer coordinates and common denominator $2^{60}$ are frozen in
\path{evidence/ns44_metric_replay/witnesses.json}. With those identical
vectors, fresh Arb coefficient enclosures at $320$ and $448$ bits,
including the analytic $96$-term series remainder, prove positive norm,
positive $D$ energy, positive Weil energy, and
$v^*(\mathsf W_5-10^{-8}D)v<0$ in each parity.
The respective ratio balls are contained in the displayed rational
intervals; rounded display values are
$6.4235244156530\cdot10^{-18}$ and
$1.1127922753304\cdot10^{-17}$.
The shifted energies are respectively approximately
$-2.4267798261306\cdot10^{-8}$ and
$-2.4549167289288\cdot10^{-8}$; the signs are decided by the
archived balls, not these rounded values.

A separate expansion into the normalized $+n,-n$ Fourier coordinates
verifies the same signs and overlaps the parity assembly. This checks
signs and indexing independently while sharing the analytic coefficient
enclosures; it is not a second derivation of the kernel. A further
reviewer replay and direct four-sign double sum reproduce the gates.
All coefficients are freshly assembled with an empty temporary cache.
The float proposal is used only to choose frozen vectors; no midpoint
solve enters a proof gate. The complete ball strings, source hashes,
verifier and independent review are archived and indexed in
\path{evidence/v152/}.
\end{proof}

\begin{proposition}[Fresh obstruction for the existing dyadic block comparison]
\label{prop:v152-block-comparison}
For the $\lambda=5$, $N=25$ dyadic partition, the signed block
norm-comparison criterion fails in both parities for every admissible
choice $0<M_j\preceq T_{jj}$ and every positive scalar row weighting.
The first three blocks
\[
 I_0=\{26,\ldots,50\},\quad I_1=\{51,\ldots,100\},\quad
 I_2=\{101,\ldots,200\}
\]
already suffice. This closes this partition's sufficient comparison,
not all signed block methods or the sign of the complete Weil form.
\end{proposition}
\begin{proof}[Certified pairings and exact comparison argument]
For each of the three pairs in each parity, separate frozen dyadic
vectors $x,y$ satisfy, at both $320$ and $448$ bits,
\[
 X=x^*T_{jj}x>0,\qquad Y=y^*T_{kk}y>0,\qquad
 |x^*T_{jk}y|^2-H_{jk}^2XY>0.
\]
The zero-diagonal symmetric rational matrices $H$ have entries
\[
 (H_{01},H_{02},H_{12})=
 \begin{cases}(997,568,481)/1000,&\text{even},\\
              (997,484,612)/1000,&\text{odd}.
 \end{cases}
\]
The same fresh coefficient and independent signed-index checks are
used for both energies and the cross pairing. The exact Rayleigh
quotients on $(1,1,1)$ are
\[
 \frac{341}{250}=1.364>1\quad\text{and}\quad
 \frac{2093}{1500}>1,
\]
respectively. The frozen witnesses and strict interval gates are in
\path{evidence/ns45_block_replay/}; an independent rerun gives a
byte-identical certificate. Pair witnesses need not be one common
physical vector because they lower-bound separate operator norms.

If admissible positive metrics exist, then each $T_{jj}$ is positive
definite. The verified quotients lower-bound the normalized coupling
norms for $T$, and Lemma~\ref{lem:v128-comparison-obstruction} gives
$\beta_{jk}(M)\ge\beta_{jk}(T)>H_{jk}$. If the metrics do not exist,
the criterion already fails its prerequisite. Thus no unconditional
full-block inertia assertion is needed. A positive row weighting making
every comparison row sum below one would, by diagonal similarity and
the row-sum norm, force the spectral radius of $H$ below one, contrary
to its exact Rayleigh quotient. An infinite continuation adds only
nonnegative norm terms, so it cannot repair these first three rows.
This proves precisely the stated comparison failure.
\end{proof}


\subsection*{Prime-channel retention for the explicit Gaussian weight}

% NS-43 working proof fragment. Exact identities and proved implications only.
% Not a manuscript/version change. No numerical certificate is asserted.

\paragraph{Notation and the comparison being ruled out.}
Write $C_0=(2\pi-3)/(4\pi)$ and retain the unnormalized Gaussian weight
$h$ of \eqref{eq:v150-weight}. For a prime power $n=p^k\ge2$, put
$\ell_n=\log n$, $w_n=\Lambda(n)/\sqrt n$, and
\[
 J^{(n)}_{a,h}[g]
 =w_n\int_{-a}^{a-\ell_n}h(x)h(x+\ell_n)
                 |g(x+\ell_n)-g(x)|^2\,dx
\]
when $\ell_n<2a$, and set it to zero otherwise. The global channel
$J^{(n)}_h$ has the same integral over $\mathbb R$. With $g=f/h$ on
$I_a$ and zero outside, $J^{(n)}_h\ge J^{(n)}_{a,h}$.

For $0\le\varepsilon_{n,a}\le1$, a coefficientwise loss is
$L_a[g]=\sum_n\varepsilon_{n,a}J^{(n)}_h[g]$ in the global comparison,
or the corresponding sum of interior channels in the local comparison.
Any additional discarded nonnegative energy can be included in $L_a$.
The two comparisons at issue are, respectively,
\begin{align*}
 J_h[g]-L_a[g]&\ge\tfrac23\operatorname{Var}_\nu(g)
       +2|\langle f,s\rangle|^2-\eta_a\|f\|^2,\\
 \int(r_a)_-|f|^2&\le J_{a,h}[g]-L_a[g]+\Pi[f]
                                  +\eta_a\|f\|^2.
\end{align*}
They require, respectively, $q_a[f]-L_a[g]\ge-\eta_a\|f\|^2$
and $q_a[f]-\int(r_a)_+|f|^2-L_a[g]\ge-\eta_a\|f\|^2$.
Thus both imply $\eta_a\ge L_a[g]-q_a[f]$ on every admissible unit test.
No compensating change in the potential, variance, pole or retained
coefficients is part of this statement.

\begin{proposition}[Every fixed nonzero prime-power channel has a critical coefficient]
\label{prop:ns43-fixed-channel}
Fix a prime power $n\ge2$, let $b=(\log2)/16$, and put
\[
 \kappa_n=C_0\Lambda(n)n^{-5},\qquad
 \gamma_n=\pi(1-n^{-2})e^{-6b}.
\]
Choose two fixed real even orthonormal functions
$\phi_1,\phi_2\in C_c^\infty(-b,b)$, and define
\[
 A_b=2\sum_{j=1}^2\int_{\mathbb R}
    |\Re\psi(1/4+i\xi/2)-\log\pi|\,|\widehat\phi_j(\xi)|^2\,d\xi,
 \qquad B_{a,b}=A_b+(8a+1)e^a+2a.
\]
For every sufficiently large $a$, each of the even actual-source
complement and the odd sector contains a unit smooth test satisfying
\[
 J^{(n)}_{a,h}[f/h]\ge\kappa_n\exp(\gamma_n e^{2a}),
 \qquad q_a[f]\le B_{a,b}.
\]
Consequently either comparison above, if valid on the indicated sector
with $\eta_a\ge0$, must obey
\[
 \eta_a\ge
 [\varepsilon_{n,a}\kappa_n\exp(\gamma_n e^{2a})-B_{a,b}]_+.
\]
In particular, deleting this channel, or decreasing its coefficient by
any fixed positive amount along a cofinal subfamily, is incompatible
with a uniform finite ordinary-norm error. This concerns the comparison,
not the sign of $q_a$ or the unmodified coefficient-one comparison.
\end{proposition}
\begin{proof}
The existing two-sided bounds for $h$ give, for $x\ge\ell_n$,
\[
 \frac{h(x-\ell_n)}{h(x)}
 \ge C_0 n^{-9/2}
             \exp\!\bigl(\pi(1-n^{-2})e^{2x}\bigr).
\]
The restriction $x\ge\ell_n$ is required here: both weight bounds
are being applied at nonnegative arguments.

Put $t=a-2b$ and
$F_{j,t}^{\pm}=2^{-1/2}(\phi_j(\cdot-t)\pm\phi_j(\cdot+t))$.
For $a>3b+\ell_n$, their support bands
$B_+\subset(a-3b,a-b)$ and $B_-=-B_+$ are disjoint and inside $I_a$.
The two $F_{j,t}^+$ are orthonormal; a unit vector in their complex
span can be chosen exactly orthogonal to the actual $u_a$, by the
kernel of the single linear functional $f\mapsto\langle f,u_a\rangle$.
This argument does not identify $u_a$ with $h$, approximate its
coefficients or require a source error estimate. The two odd functions
are orthonormal and every vector in their span is source-orthogonal.

For every unit vector in either span, an endpoint $x\in B_+$ has
$x-\ell_n\ge0$, remains inside the window, and is outside both support
bands, since $\ell_n\ge\log2>2b$. Select the edge from $x-\ell_n$ to
$x$, and use its reflected inward edge for $x\in B_-$. These selected
edges are distinct; no selected edge has two nonzero endpoints. Hence
\[
 J^{(n)}_{a,h}[f/h]
 \ge w_n\inf_{x\in(a-3b,a-b)}\frac{h(x-\ell_n)}{h(x)}
             \int_{B_+\cup B_-}|f(x)|^2dx
 \ge\kappa_n\exp\!\bigl(\pi(1-n^{-2})e^{2(a-3b)}\bigr).
\]
All interference between $\phi_1$ and $\phi_2$ is retained inside
$|f|^2$; its integral is the ordinary norm, exactly one.

The archimedean contribution is at most $A_b$: translation/reflection
multipliers have modulus at most $\sqrt2$, and Cauchy--Schwarz in the
two-dimensional coefficient vector gives the displayed bound uniformly
for every unit mixture. Prime correlations have total absolute value
at most $2\sum_{m<e^{2a}}\Lambda(m)/\sqrt m\le8ae^a$.
The even pole is at most $2a+2\sinh a\le2a+e^a$; the odd pole is
nonpositive. These estimates give $q_a[f]\le B_{a,b}$ including all
cross terms. The preceding comparison identities finish the proof.
Exterior edges were retained in those identities; the lower bound uses
only interior edges, so adding the exterior part cannot weaken it.
\end{proof}

\begin{proposition}[Uniform obstruction for moving omitted channels away from the endpoint]
\label{prop:ns43-moving-channel}
Fix $0<\theta<1$. Let
\[
 b_\theta=\frac{\min(\log2,-\log\theta)}{16},\qquad
 d_\theta=\min(\log2,-\log\theta-6b_\theta),
\]
\[
 \gamma_\theta=\pi(1-e^{-2d_\theta})e^{-6b_\theta}>0,
 \qquad c_\theta=C_0(\log2)\theta^{-1/2}e^{9b_\theta/2}>0.
\]
For all sufficiently large $a$, the same two parity spaces built with
$b=b_\theta$, $t=a-2b_\theta$, satisfy, for every unit vector and
\emph{simultaneously for every prime power} $2\le n\le\theta e^{2a}$,
\[
 J^{(n)}_{a,h}[f/h]\ge
 c_\theta e^{-11a/2}\exp(\gamma_\theta e^{2a}).
\]
Thus either comparison, with $\eta_a\ge0$, and an omitted channel
$n=n(a)$ in this range, even when it varies with $a$, forces
\[
 \eta_a\ge
 [\varepsilon_{n(a),a}c_\theta e^{-11a/2}
             \exp(\gamma_\theta e^{2a})-B_{a,b_\theta}]_+.
\]
The same statement holds on any cofinal subfamily satisfying these
hypotheses. No conclusion is asserted merely because an omitted
channel exists at each window if its ratio $n(a)/e^{2a}$ tends to one.
\end{proposition}
\begin{proof}
Here $d_\theta\ge10b_\theta>2b_\theta$. For $x\in(a-3b,a-b)$ and
$\log2\le\ell_n\le2a+\log\theta$, one has
\[
 \ell_n\ge d_\theta,\qquad 2x-\ell_n\ge d_\theta,
 \quad\text{hence}\quad |x-\ell_n|\le x-d_\theta.
\]
For large $a$ these inward endpoints lie inside the window and outside
both support bands, since
$|x-\ell_n|\le a-b-d_\theta<a-3b$. Unlike the fixed-channel proof,
$x-\ell_n$ need not be nonnegative. Use evenness of $h$ and apply its
bounds to $|x-\ell_n|$ and $x$, rather than incorrectly applying the
previous ratio formula past its range. They give
\begin{align*}
 \frac{h(x-\ell_n)}{h(x)}
 &\ge C_0\exp\!\left(\tfrac92(|x-\ell_n|-x)
                   +\pi(e^{2x}-e^{2|x-\ell_n|})\right)\\
 &\ge C_0e^{-9(a-b)/2}
                \exp\!\left(\pi(1-e^{-2d_\theta})e^{2(a-3b)}\right).
\end{align*}
Also $w_n\ge(\log2)\theta^{-1/2}e^{-a}$. The same distinct selected
inward edges and norm-one integration prove the claimed bound. The
source kernel and the original-form upper budget are exactly as above,
with fixed $b_\theta$ independent of $a$ and $n$.
\end{proof}

\paragraph{Truncations actually covered.}
With every retained channel coefficient at most one and all other
signed terms unchanged, the following shortcuts require unbounded error:
(i) any fixed finite set of retained prime-power channels; (ii) any
fixed finite set of retained primes, even retaining all of their powers;
(iii) a uniformly bounded number of retained channels (or primes),
although their labels may vary with the window; and (iv) a cutoff
retaining only $n\le M(a)$ with $M(a)=o(e^{2a})$.
For (iii), take more fixed prime powers (respectively primes) than the
retained count: one is omitted, and all lie in the uniform range for
large $a$. For (iv), if $M(a)\ge1$, the first power of $2$ exceeding
$M(a)$ is at most $2M(a)$ and is eventually at most $\theta e^{2a}$;
if $M(a)<1$, take $n=2$ instead. The same proof covers
$\limsup M(a)e^{-2a}<1/2$ by choosing $\theta$ between twice that
limsup and one. No prime-gap theorem is being used.

A fixed coefficient loss in any of those omitted channels is covered.
A loss tending sufficiently rapidly to zero is only subject to the
necessary displayed rate; it is not excluded categorically. Merely
saying ``finitely many channels at each window'' is insufficient:
the full local sum itself is finite. A cutoff approaching the strict
endpoint, compensating changes in the signed diagonal/pole, coefficients
larger than one in other channels, and a different positive weight are
not ruled out by these propositions.

\paragraph{One specific surviving signed, cofinal next input.}
This is a proposed \emph{open input on a restricted test family}, not
a theorem and not a restatement of CAE on all tests. For fixed real even
$\phi_1,\phi_2\in C_c^\infty(-b,b)$ as above, set
\[
 R_{jk}(u)=\int\phi_j(z+u)\phi_k(z)\,dz,\qquad
 C_j=\int\cosh(x/2)\phi_j(x)\,dx,
\]
and consider the explicit real symmetric two-by-two matrix
\[
 E_{jk}(t)=e^t C_jC_k-
       \sum_{n\ge2}\frac{\Lambda(n)}{\sqrt n}
                          R_{jk}(\log n-2t),\qquad t>2b.
\]
Only $e^{2t-2b}<n<e^{2t+2b}$ contribute. The candidate next lemma is
\[
 \sup_{t\ge t_0}\|E(t)\|_{\mathrm{op}}<\infty
 \quad\text{for this specified pair of bumps.}
\]
This asks for a signed, smoothly weighted prime--continuum cancellation
at cofinal multiplicative locations, not separate absolute estimates
of Gaussian transformed channels.

Its limited consequence is checkable. Let
$D_{jk}=q(\phi_j,\phi_k)$ and
$C_{jk}(t)=q(\phi_j(\cdot-t),\phi_k(\cdot+t))$ for the polarized
physical form. For $a=t+2b$ the complete prime and pole formula gives
\[
 C_{jk}(t)=E_{jk}(t)+e^{-t}C_jC_k+A_{jk}(t),
 \qquad A_{jk}(t)=O_b(e^{-t}),
\]
where the last term is the off-diagonal archimedean form. Indeed its
kernel on disjoint supports is $-\rho(|x-y|)$, so the decay follows
from $\rho(s)=O(e^{-s/2})$. The pole is
$2C_jC_k\cosh t=(e^t+e^{-t})C_jC_k$; precisely one orientation of
the prime shift connects the separated supports, giving the sum above
without an extra factor two. The parity matrices are exactly
$D+C(t)$ and $D-C(t)$, with all cross terms present. The candidate
bound therefore supplies a uniform floor only on these two-dimensional
parity spaces, including their actual-source-orthogonal even vectors.
It does not cover arbitrary physical tests, does not establish CAE,
and is not asserted to be weaker than RH or presently provable.


\subsection*{A restricted prime-discrepancy input can already carry the full zero obstruction}
\label{sec:ns43-bump-strength}
The two-dimensional test family in the channel analysis is concrete,
but its uniform discrepancy estimate must not be advertised as a
weaker arithmetic foothold without checking its strength.

\begin{proposition}[What bounded translated-bump discrepancy excludes]
\label{prop:ns43-bump-strength}
Let $\phi\in C_c^\infty(-b,b)$ be real and even, put
\[
 P(s)=\int e^{sx}\phi(x)\,dx,\qquad
 R(v)=\int\phi(x+v)\phi(x)\,dx,\qquad C=P(1/2),
\]
and define
\[
 F(u)=e^{u/2}C^2-
        \sum_{n\ge2}\frac{\Lambda(n)}{\sqrt n}R(\log n-u).
\]
If $F$ is bounded for $u\ge u_0$, every nontrivial zero $\rho$ of
$\zeta$ with $\Re\rho>1/2$ must obey $P(\rho-1/2)=0$.
There are explicit nonnegative even smooth compactly supported probes
$\phi$ for which all zeros of $P$ lie on the imaginary axis. For
such a probe, the asserted boundedness would exclude every zero off
the critical line. This is a conditional implication, not a proved
bound for $F$ and not a new RH claim. For an arbitrary preassigned
pair of bumps, the matrix input forces common transform zeros at
any off-line zero; no conclusion without that qualification is asserted.
\end{proposition}
\begin{proof}
The correlation is supported in $[-2b,2b]$ and satisfies
$\int e^{sv}R(v)\,dv=P(s)^2$ by evenness. For $\Re s>1/2$,
absolute convergence of the von Mangoldt Dirichlet series gives
\begin{equation}\label{eq:ns43-bump-laplace}
 \int_{u_0}^{\infty}e^{-su}F(u)\,du
 =\frac{C^2}{s-1/2}+P(s)^2\frac{\zeta'}{\zeta}(s+1/2)+H(s),
\end{equation}
where
\begin{align*}
 H(s)={}&C^2\frac{e^{-(s-1/2)u_0}-1}{s-1/2}\\
 &+\sum_{\log n\le u_0+2b}\frac{\Lambda(n)}{\sqrt n}
       \int_{-\infty}^{u_0}e^{-su}R(\log n-u)\,du.
\end{align*}
This correction is entire: the quotient is removable and only
finitely many compactly supported integrals occur. The usual
identity $-\zeta'/\zeta(s)=\sum\Lambda(n)n^{-s}$ is used only in
its half-plane of absolute convergence (\href{https://dlmf.nist.gov/27.4.E12}{DLMF 27.4.12}). The pole at $s=1/2$ in
\eqref{eq:ns43-bump-laplace} cancels because $C=P(1/2)$.
Boundedness of $F$ makes the left side holomorphic on $\Re s>0$.
Uniqueness of meromorphic continuation therefore forces the residue
at $s=\rho-1/2$ to vanish. That residue is
$\operatorname{mult}(\rho)P(\rho-1/2)^2$, proving the first assertion.
For a matrix of two bumps apply this argument to each diagonal;
all off-line zeros must then be common zeros of their transforms.

Here is a zero-independent choice of probe. Fix $0<B<b$ and set
$b_k=B/(\zeta(2)k^2)$. Let $\mu$ be the distribution of
$\sum_{k\ge1}U_k$, where the $U_k$ are independent uniform variables
on $[-b_k,b_k]$. The sum converges absolutely since $\sum b_k=B$.
Its measure is nonnegative, even, supported in $[-B,B]$, and has
entire bilateral transform
\[
 P_0(s)=\prod_{k\ge1}\frac{\sinh(b_ks)}{b_ks}.
\]
The product converges uniformly on compact sets, since each factor
is $1+O(b_k^2)$ there and $\sum b_k^2<\infty$. It is nonzero away
from the zeros of its factors. Those zeros are exactly imaginary
points $i\pi m/b_k$, $m\ne0$.
For every integer $M>0$,
\[
 |P_0(i\xi)|\le
 \left(\prod_{k=1}^M b_k^{-1}\right)|\xi|^{-M},
\]
since all remaining sinc factors have modulus at most one. Thus
Fourier inversion gives a $C^\infty$ density supported in $[-B,B]$.
It is nonnegative and nonzero because it is the density of $\mu$;
normalizing it in $L^2$ does not change its transform zeros. This is
the required $\phi\in C_c^\infty(-b,b)$. If an orthonormal pair is
desired, take the second even smooth bump supported in
$\{B<|x|<b\}$ and normalize it; disjoint supports give orthogonality.
For the first probe no zero with $\Re(\rho-1/2)>0$ can be masked.
The functional-equation symmetry supplies the corresponding exclusion
on the other side of the critical line, conditional on the unproved
boundedness of $F$.
\end{proof}

For the proposed matrix in the channel analysis, $F(u)=E_{11}(u/2)$
for this first probe. Hence even this restricted uniform estimate can
already contain the full off-line-zero obstruction. Nothing above
proves that estimate or asserts the converse. The useful outcome of
this check is a limit on the advertised strength of a proposed next
lemma, not a new positivity argument.


% NS-43 concentration subtask. Analytic working fragment, not a version claim.
% All conclusions are explicitly scoped. No numerical or RH claim.
\subsection*{Adaptive concentration: completeness is not an estimate (NS-43)}
\label{ns43-conc-sec}

\paragraph{Objects and source constraint.}
Fix $a>0$. Let $\mathcal H_a=L^2(-a,a)$, let $\mathcal F_a$ mean
zero extension followed by the unitary Fourier transform, and let
$\beta_a$ be the exact complete symbol, including its endpoint and pole
terms. Put $D_a=\| (\beta_a)_-\|_\infty$.
Use either of the closed spaces
\[
 S_a^{\rm ev}=\{f\in L^2_{\rm ev}(-a,a):\langle f,u_a\rangle=0\},
 \qquad S_a^{\rm odd}=L^2_{\rm odd}(-a,a),
\]
where $u_a$ is the actual normalized repaired source.
In this subsection $S_a$ denotes one of these spaces and $P_{S_a}$
its orthogonal projection. All operator lower bounds act on the
\emph{whole} $S_a$, not on a head or on individual pencil directions.
For an even Borel set $G$ put
\[
 C_{a,S}(G)=P_{S_a}\mathcal F_a^*1_G\mathcal F_aP_{S_a}|_{S_a}.
\]
The complete form is
$q_a[f]=\int\beta_a|\mathcal F_af|^2$, with form domain given by
$\int(\beta_a+D_a)|\mathcal F_af|^2<\infty$.

\begin{proposition}[Exact adaptive hierarchy and the original lower edge]
\label{ns43-conc-prop-adaptive}
Suppose $\beta_a$ is continuous, even, bounded below, and tends to
$+\infty$ at both frequency ends. Assume also that
$S_a\cap\mathcal D(q_a)$ contains a nonzero vector.
These hypotheses hold for the manuscript's fixed-window symbol:
$\beta_a(\xi)=O_a(\log(2+|\xi|))$ puts every compactly supported
smooth physical test in the form domain. In the even sector, a
nonzero linear combination of two independent even smooth tests
can satisfy the one source constraint; in the odd sector any
nonzero odd smooth test suffices. For integers $n\ge1$ define
\begin{align*}
 \delta_n&=2^{-n},\qquad J_n=n2^n,\\
 m_{a,n}(\xi)&=-D_a+
 \delta_n\sum_{j=1}^{J_n}
       1_{\{\beta_a(\xi)+D_a\ge j\delta_n\}}.
\end{align*}
These are finite step minorants with positive weights, and
$m_{a,n}\uparrow\beta_a$ pointwise. Their exact compressed lower
edges obey
\begin{equation}\label{ns43-conc-eq-edge-limit}
 \lambda_{a,n}:=\inf_{\substack{f\in S_a\\\|f\|=1}}
   \int m_{a,n}|\mathcal F_af|^2
 \ \uparrow\
 \lambda_a:=\inf_{\substack{f\in S_a\cap\mathcal D(q_a)\\\|f\|=1}}
   q_a[f].
\end{equation}
Consequently the following two existence statements on any prescribed
cofinal family are equivalent:
\begin{enumerate}[label=(\alph*)]
\item the complete even source complement and complete odd sector
have one common finite lower floor $-C$, independent of $a$;
\item one may choose finite indices $n_{\rm ev}(a),n_{\rm odd}(a)$
so that the corresponding exact signed concentration sums
\[
 -D_a I_{S_a}+\delta_n\sum_{j=1}^{J_n}
 C_{a,S}\bigl(\{\beta_a+D_a\ge j\delta_n\}\bigr)
\]
have a common finite lower floor independent of $a$.
\end{enumerate}
More precisely, (a) yields a floor $-C-\varepsilon$ in (b) for any
fixed $\varepsilon>0$; (b) yields its stated floor in (a).
This equivalence supplies no effective choice of $n(a)$ and no
arithmetic lower estimate for any of the sums.
\end{proposition}
\begin{proof}
The formula is lower dyadic rounding of
$\min(\beta_a+D_a,n)$, minus $D_a$. Refinement of the grid and
increase of the cap prove monotonicity and pointwise convergence.
Each $m_{a,n}$ is bounded. Thus its compressed form is a bounded
self-adjoint operator, and $\lambda_{a,n}\le\lambda_a$.
The numbers increase to some finite $\ell\le\lambda_a$; finiteness
above follows by testing one fixed unit vector in
$S_a\cap\mathcal D(q_a)$, and below by $-D_a$.

Choose unit $f_n\in S_a$ with
$\int m_{a,n}|\mathcal F_af_n|^2\le\lambda_{a,n}+1/n$.
The integrals against $m_{a,n}+D_a$ are uniformly bounded by some
$B<\infty$. Given $R\ge1$, choose $M_R$ with
$\beta_a(\xi)+D_a\ge2R$ for $|\xi|>M_R$.
For sufficiently large $n$, lower rounding gives
$m_{a,n}+D_a\ge R$ there. Hence
\[
 \int_{|\xi|>M_R}|\mathcal F_af_n|^2\le B/R.
\]
Restriction of $\mathcal F_a$ to $[-M_R,M_R]$ is a compact map:
its integral kernel has finite square integral on the product of
the two bounded intervals. This fact and the uniform frequency-tail
bound make $(\mathcal F_af_n)$, after discarding finitely many
terms for each $R$, relatively compact in $L^2(\mathbb R)$.
Take a strongly convergent subsequence. Its physical limit $f$ is
in the closed space $S_a$ and has norm one.

For each fixed $k$, boundedness of $m_{a,k}$ and monotonicity give
\[
 \int m_{a,k}|\mathcal F_af|^2
 =\lim_n\int m_{a,k}|\mathcal F_af_n|^2\le\ell.
\]
Monotone convergence applied after adding $D_a$ gives
$f\in\mathcal D(q_a)$ and $q_a[f]\le\ell$.
Thus $\lambda_a\le\ell$, proving \eqref{ns43-conc-eq-edge-limit}.
Choose the two finite indices separately at each window to within
$\varepsilon$ of their limits. Conversely a pointwise minorant's
operator lower bound transfers immediately to $q_a$.
\end{proof}

\paragraph{What this closes.}
The proposal ``allow arbitrary adaptive levels and assert that their
exact signed operator sums have a uniform lower floor'' is an
equivalent restatement of the unknown complement-floor target.
It is not a weaker arithmetic input. This closes only that claimed
logical shortcut. It does not close a method that proves a specific
operator comparison from additional arithmetic structure.
The finite level count is allowed to grow without restriction here.
Therefore the statement does not evade a proved obstruction by
silently retaining bounded complexity.

\begin{proposition}[All scalar set bounds can still lose the signed floor]
\label{ns43-conc-prop-scalar}
Even exact upper and lower concentration bounds for every measurable
set do not in general determine the signed operator lower edge.
This failure occurs already for commuting positive operator-valued
measures with three outcomes, after an exact rank-one source has
been removed.
\end{proposition}
\begin{proof}
On $\mathbb C^3$ define two three-outcome positive operator-valued
measures by making outcome operator $C_i$ diagonal, with its three
diagonal entries the $i$th column of the following arrays:
\[
 A=\frac1{10}
 \begin{pmatrix}0&7&3\\3&0&7\\7&3&0\end{pmatrix},
 \qquad
 B=\frac1{10}
 \begin{pmatrix}7&0&3\\3&7&0\\0&3&7\end{pmatrix}.
\]
Each row sums to one, so in both models $\sum_iC_i=I$.
For every singleton $E$, the exact lower and upper scalar bounds are
$p(E)=0$, $c(E)=7/10$. For every two-element set they are
$p(E)=3/10$, $c(E)=1$. Empty and full sets have their usual values.
These list \emph{all} subsets; hence all scalar data coincide.

For symbol values $b=(-1/2,1/2,3/2)$, the exact signed operators have
diagonal entries
\[
 \sum_i b_i C_i^A=\operatorname{diag}(4/5,9/10,-1/5),
 \qquad
 \sum_i b_i C_i^B=\operatorname{diag}(1/10,1/5,6/5).
\]
Thus their bottoms are $-1/5$ and $1/10$ despite identical scalar
concentration data for every set.

The usual scalar certification rule cannot recover the positive
bottom in model $B$, even after optimizing over all sets and weights.
Indeed the probability vector $z=(7/10,3/10,0)$ satisfies every
scalar constraint $p(E)\le z(E)\le c(E)$ in model $B$.
For any pointwise minorant
$\alpha+\sum_jw_j1_{G_j}\le b$ with $w_j\ge0$, one therefore has
\[
 \alpha+\sum_jw_jp(G_j)
 \le\alpha+\sum_jw_jz(G_j)\le\sum_i b_i z_i=-1/5.
\]
This upper bound on the certifiable lower bound is attained by
$b=-1/2+1_{\{2,3\}}+1_{\{3\}}$.
Multiplying $b$ by any $t>0$ gives an exact best scalar lower bound
$-t/5$, while the actual model-$B$ operator is at least $(t/10)I$.
The scalar loss therefore diverges along $t\to\infty$ even though
the true signed floor is nonnegative.

To add an exact source, take an orthogonal direct sum with one unit
vector $u$ and give the three outcome operators on that vector the
values $(1/2,1/2,0)$. Their signed sum annihilates $u$ and has no
cross terms with $u^\perp$. Compression off this same exact source
returns the two arrays above. The countermodel is finite and
nonarithmetic: it does not assert that either array arises from
physical Paley--Wiener concentration or from the corrected zero field.
\end{proof}

\paragraph{Relation to the existing obstructions.}
The negative-only and primitive scalar estimates discard signed
physical information; v1.37 closes its stated pointwise and primitive
uniform-budget proposals. The signed operators in
Proposition~\ref{ns43-conc-prop-adaptive} retain the positive regions
and all cross terms, so those scalar no-go proofs do not apply to them.
This is a distinction in hypotheses, not a proof that their lower
edges meet the target.

The QG growth clause in Assumption~\ref{ass:v149-qg}, together with
its \emph{separate} source-admissible unit packet satisfying the
all-Borel domination and uniform upper energy assumptions, gives
\[
 \eta(m_a)\ge
 [c_0\,\mathrm{QC}_{K(a)}(\widetilde m_a)-Q_*]_+.
\]
It obstructs a bounded level count under those two open inputs.
For every fixed bound $K_0$, the cost diverges, so a successful
bounded-loss family would have $K(a)\to\infty$.
Neither QG alone, a finite-bin packet test, nor the conditional
obstruction proves failure of unrestricted adaptive comparison.
In particular there is no contradiction with
Proposition~\ref{ns43-conc-prop-adaptive}: its indices are unrestricted
and its existence is conditional on the original floor.

The odd sector requires no source projection. In the even sector,
the actual source must be removed before every concentration or
packet comparison. For a raw unit $f$ with
$r=|\langle f,u_a\rangle|<1$, its normalized projection $g$ obeys
the already proved pointwise and energy estimates
\[
 |\mathcal F_ag(\xi)|\ge
 \frac{[|\mathcal F_af(\xi)|-r|\mathcal F_au_a(\xi)|]_+}
 {\sqrt{1-r^2}},\qquad
 |(1-r^2)q_a[g]-q_a[f]|
 \le(2r+r^2)\|\mathsf W_{e^a}u_a\|.
\]
Thus raw packet coverage or a raw pencil ratio cannot be substituted
for the source-compressed hypotheses. The repaired-source residual
is used only in the subsequent complement-to-full-floor transfer;
it does not supply the missing signed comparison or packet coverage.

\paragraph{Disposition and missing input.}
The full signed concentration route remains open. The proposed
unrestricted exact hierarchy is classified as an \emph{exact
reformulation}; completeness of optimization from individual scalar
set bounds is a \emph{proved-false general simplification}, with no
arithmetic impossibility claim. A productive next input must estimate
joint signed operators using actual arithmetic structure, with a
cofinal uniform error on both complete parity obligations. No such
estimate is proved here. QG is an obstruction-investigation input,
not a positivity mechanism; CAE is already the same floor in another
coordinate system. These statements add no numerical window, tail
metric, certificate, G2 conclusion, or RH conclusion.


% NS-43 analytic working note. No numerical computation or certificate.
% Scope: additional complete-space inputs for CCM step (b).

\subsection*{CCM selection requires a relative estimate, not another absolute residual}

\paragraph{Status.}
The fixed-window simple ground and its real-zero transform do not identify
Riemann's $\Xi$. The known source limit and complete residual are useful
prerequisites. The new input would be a cofinal \emph{complete spectral
selection estimate}, with a rate compatible with the transform topology.
The propositions below are conditional elementary functional analysis;
none establishes the hypotheses for the arithmetic Weil family.

Put $H=L^2(\mathbb R)$, $I_a=(-a,a)$ and let $J_a$ denote zero extension.
Let $A_a$ be the complete lower-bounded self-adjoint even-sector operator
on $L^2(I_a)$, with compact resolvent. If the global ground is used,
its evenness and ordering below the odd sector must also be supplied.
In the following selection lemmas assume the selected ground is simple,
with unit vector $v_a$, eigenvalue $\mu_a$, and complete next-level
separation $g_a>0$:
\[
 q_a[f]-\mu_a\|f\|^2\ge g_a\|f\|^2
 \quad(f\in\mathcal D(q_a),\ f\perp v_a).
\]
This is an ordinary complete spectral separation, not an LDL pivot,
finite-compression gap, or a source-complement estimate inferred from
separate diagonal blocks.

Let $h$ be the explicit positive Gaussian arithmetic radical of
\eqref{eq:v150-weight}, and $\phi=h/\|h\|_2$. Use exactly the fixed
boundary profile of \ref{prop:v135-growing-radical}: fix one smooth
$\vartheta$ with $0\le\vartheta\le1$, $\vartheta(t)=1$ for $t\le-1$,
and $\vartheta(t)=0$ for $t\ge-1/2$. For $a\ge2$ set
$\chi_a(x)=\vartheta(|x|-a)$. This is smooth (it is constant near zero),
even, supported in $I_a$, equal to one on $[-a+1,a-1]$, and has
uniformly bounded derivatives of every fixed order. Define
\[
 p_a=\chi_a\phi/\|\chi_a\phi\|_2,\qquad
 \alpha_a=q_a[p_a],\qquad E_a=\alpha_a-\mu_a\ge0.
\]
For large $a$, these vectors lie in the complete operator domain.
The double-exponential tail from v1.35 implies, for every fixed $R\ge0$,
\[
 \int_{\mathbb R}e^{R|x|}|J_ap_a(x)-\phi(x)|\,dx\longrightarrow0.
\]
The vector $p_a$ is the normalized $j=0$ column of the growing block.
Thus \ref{prop:v135-growing-radical}, specifically
\eqref{eq:v135-block-residual}, gives the complete estimate
\[
 \|A_ap_a\|\le\varepsilon_a:=Ca\exp(-e^a/100),\qquad
 |\alpha_a|\le\varepsilon_a,
\]
with constants for this fixed source and profile. No residual estimate
is asserted for arbitrary cutoffs whose derivatives may grow with $a$.
The cutoff $p_a$ here is explicit and is not silently identified with
the different repaired prolate source $u_a$.

\begin{proposition}[A complete relative selection estimate]
\label{prop:ns43-ccm-rayleigh-selection}
Choose the phase of $v_a$ so that $\langle p_a,v_a\rangle\ge0$.
If $E_a/g_a\to0$, then $\|J_av_a-\phi\|_2\to0$.
Define
\[
 C_R(a)=\begin{cases}
 \sqrt{(e^{2Ra}-1)/(2\pi R)},&R>0,\\
 \sqrt{a/\pi},&R=0.
 \end{cases}
\]
If, in addition, for every fixed $R>0$,
\[
 C_R(a)\sqrt{E_a/g_a}\longrightarrow0,
 \tag{S1}
\]
then the unitary Fourier transforms of $J_av_a$ converge locally
uniformly on $\mathbb C$ to $\widehat\phi$.
\end{proposition}
\begin{proof}
Write $p_a=\kappa_a v_a+e_a$, where $e_a\perp v_a$ and
$\kappa_a\ge0$. The spectral projection preserves the form domain,
and the cross term is exactly zero because $v_a$ is an eigenvector.
Consequently
\[
 E_a=q_a[e_a]-\mu_a\|e_a\|^2\ge g_a\|e_a\|^2,
 \quad \kappa_a^2=1-\|e_a\|^2,
 \quad \|p_a-v_a\|^2=2(1-\kappa_a)\le2E_a/g_a.
\]
This proves the ordinary-norm conclusion. For $|\Im z|\le R$,
Cauchy--Schwarz gives
\[
 |\widehat{J_a(v_a-p_a)}(z)|
 \le (2\pi)^{-1/2}\left(\int_{-a}^a e^{2R|x|}\,dx\right)^{1/2}
       \|v_a-p_a\|
 \le C_R(a)\sqrt{2E_a/g_a}.
\]
The weighted cutoff convergence above controls
$\widehat{J_ap_a}-\widehat\phi$ on the same strip. Combining the
bounds proves the assertion. No compactness hypothesis is hidden:
(S1) supplies a direct quantitative topology upgrade.
\end{proof}

\begin{proposition}[A residual-to-separator version]
\label{prop:ns43-ccm-residual-selection}
Let $p_a\in\mathcal D(A_a)$ be unit, put
$\alpha_a=\langle A_ap_a,p_a\rangle$ and
$\rho_a=\|(A_a-\alpha_a)p_a\|$. Suppose an independent complete
second-eigenvalue lower bound $\beta_a>\alpha_a$ is established,
with eigenvalues counted with multiplicity. Then the ground is simple and
\[
 \|v_a-p_a\|\le\frac{\sqrt2\,\rho_a}{\beta_a-\alpha_a}.
\]
Thus the cofinal condition
\[
 C_R(a)\frac{\rho_a}{\beta_a-\alpha_a}\longrightarrow0
 \quad\text{for every fixed }R>0 \tag{S1r}
\]
selects the entire transform of the explicit cutoff radical.
\end{proposition}
\begin{proof}
The variational upper bound is $\mu_a\le\alpha_a<\beta_a$, so
exactly one eigenvalue lies below $\beta_a$ and it is simple.
On its orthogonal complement the spectrum of $A_a-\alpha_a$ lies
above $\beta_a-\alpha_a$. Spectral expansion gives
$\rho_a^2\ge(\beta_a-\alpha_a)^2\|e_a\|^2$.
The preceding phase/norm identity and Fourier estimate complete the proof.
\end{proof}
This formulation does not require a lower ground-energy estimate matching
$\alpha_a$ to high relative precision. The missing arithmetic input is
the complete cofinal separator at the correct relative scale. A finite
matrix's second eigenvalue supplies no such lower bound without its full
tail and coupling comparison. An even-sector separator still needs the
separate odd ordering if the real-zero theorem requires the global ground.

\paragraph{A weaker alternative to the evaluator rate.}
Ordinary selection $E_a/g_a\to0$ also gives locally uniform entire
convergence if the selected grounds satisfy the independently proved
weighted tightness condition
\[
 \lim_{K\to\infty}\sup_a\int_{|x|>K}e^{R|x|}|J_av_a(x)|\,dx=0
 \quad\text{for every }R>0. \tag{S2}
\]
Indeed the integral on $[-K,K]$ converges by ordinary $L^2$ convergence
and Cauchy--Schwarz, while (S2) and the known tail of $\phi$ control
the exterior. A sufficient explicit condition for (S2) is
$\sup_a\int e^{2(R+1)|x|}|J_av_a|^2<\infty$ for each $R$.
No such ground-tail statement follows from the current residual bound.

\begin{proposition}[Complete evaluator-dual alternative]
\label{prop:ns43-ccm-evaluator-selection}
In the preceding setting suppose $E_a/g_a\to0$. Suppose further that,
for every fixed $R>0$, a complete complement estimate is established:
\[
 \sup_{|z|\le R}|\widehat{J_ae}(z)|^2
 \le D_{a,R}\bigl(q_a[e]-\mu_a\|e\|^2\bigr)
 \quad(e\in\mathcal D(q_a),\ e\perp v_a),
 \qquad D_{a,R}E_a\to0. \tag{S3}
\]
Then the same locally uniform entire convergence follows.
\end{proposition}
\begin{proof}
For the orthogonal error above,
$\widehat p_a=\kappa_a\widehat v_a+\widehat e_a$ and
$\kappa_a\to1$. Thus, on the disk,
\[
 |\widehat v_a-\widehat p_a|
 \le \kappa_a^{-1}\sqrt{D_{a,R}E_a}
       +(\kappa_a^{-1}-1)|\widehat p_a|\to0.
\]
The explicit cutoff transforms are uniformly bounded on each disk
and converge to $\widehat\phi$. This proves the assertion.
\end{proof}
The shift by $\mu_a$ and the complete complement in (S3) are part of
its hypotheses. The fixed-window v1.43 evaluator comparison for one
real zero is not a cofinal bound of this kind, nor a bound on every
complex compact set. Its scalar resolution floor is not a no-go
result for (S1)--(S3).

\begin{proposition}[A source-relative block condition sufficient for selection]
\label{prop:ns43-ccm-source-separation}
Let $p_a\in\mathcal D(A_a)$ be unit, $P_a=|p_a\rangle\langle p_a|$,
$Q_a=I-P_a$, $\alpha_a=\langle A_ap_a,p_a\rangle$, and
$r_a=Q_aA_ap_a$. Suppose the \emph{complete compressed form} satisfies
\[
 q_a[f]\ge(\alpha_a+\delta_a)\|f\|^2
 \quad(f\in\mathcal D(q_a)\cap p_a^\perp),\qquad \delta_a>0.
 \tag{S4}
\]
Then $A_a$ has exactly one eigenvalue below $\alpha_a+\delta_a$,
it is simple, and, for its phase-aligned ground $v_a$,
\[
 \alpha_a-\|r_a\|^2/\delta_a\le\mu_a\le\alpha_a,
 \qquad \|v_a-p_a\|\le\sqrt2\,\|r_a\|/\delta_a.
\]
For the explicit cutoff source, the cofinal rate
$C_R(a)\|r_a\|/\delta_a\to0$ for every $R>0$ therefore implies
entire-transform selection.
\end{proposition}
\begin{proof}
The restriction to $p_a^\perp$ is a closed lower-bounded form,
because subtracting its component along the fixed operator-domain
vector preserves the form domain continuously. The min--max principle
places the second eigenvalue at least $\alpha_a+\delta_a$; the trial
$p_a$ puts the lowest at most $\alpha_a$. For $f=tp_a+y$, $y\perp p_a$,
\[
 q_a[f]-\alpha_a\|f\|^2
 =q_a[y]-\alpha_a\|y\|^2+2\Re(t\langle A_ap_a,y\rangle)
 \ge\delta_a\|y\|^2-2|t|\|r_a\|\|y\|
 \ge-\|r_a\|^2|t|^2/\delta_a.
\]
This handles the full cross term and proves the lower edge bound.
Writing $v_a=tp_a+y$, the weak projected eigen-equation is
$(Q_aA_aQ_a-\mu_a)y=-tr_a$. Its inverse has norm at most
$1/(\alpha_a+\delta_a-\mu_a)\le1/\delta_a$, so
$\|y\|\le|t|\|r_a\|/\delta_a$. Phase alignment then gives
$\|v_a-p_a\|^2=2(1-|t|)\le2\|y\|^2$.
\end{proof}
Condition (S4) is a new, source-specific signed spectral estimate,
not an algebraic reformulation of the bounded-floor target. It asks for information beyond a floor: it separates one named source
from its whole complement and controls the relative error. No reverse
implication from RH or from a uniform floor to (S4) is claimed. In this
arithmetic family, with the known vanishing source residual, a successful
cofinal condition of this strength gives a uniform floor in the sector
where it is imposed. To use v1.36 one must impose it on the full space
or separately control the odd sector by a uniform floor; those additional
hypotheses would imply RH. Even-sector selection does not silently supply
the odd estimate. It is not an unconditional shortcut.

For clarity, the floor consequence itself needs no relative-rate
condition: in the displayed cross-term formula,
$2|t|\|y\|\le |t|^2+\|y\|^2$ and $\delta_a>0$ give
\[
 q_a[tp_a+y]\ge(\alpha_a-\|r_a\|)(|t|^2+\|y\|^2).
\]
Thus the known $\alpha_a\to0$ and $\|r_a\|\to0$ give the asserted
sector floor from (S4) alone. The stronger relative-rate condition
is needed for selection of the Gaussian profile and its transform,
not for this floor deduction.

The v1.35 even reflected near-radical block has dimension $J_a+1$,
so eventually it contains a unit vector orthogonal to $p_a$. Its complete
residual is at most $\varepsilon_a$, hence (S4) forces
$\delta_a\le\varepsilon_a-\alpha_a\le\varepsilon_a+|\alpha_a|$.
The explicit cutoff source also has $|\alpha_a|=O(\varepsilon_a)$.
This rules out a uniform positive $\delta_a$ or a polynomial lower scale.
A source-selecting separation must shrink and must still dominate the
source error in the relative sense above. The existence of a growing
near-zero block does not logically exclude such a hierarchy. Neither
that hierarchy nor its required rate is in the archive.

\paragraph{Identification of the limit.}
By \eqref{eq:v150-mellin}, with
$\Xi(z)=\xi(1/2+iz)$ and the even functional equation,
\[
 \widehat\phi(z)=\frac{2}{\sqrt{6\pi}\,\|h\|_2}\,\Xi(z).
\]
In particular $\widehat\phi(0)>0$. Once either selection criterion is
proved, $\widehat v_a(0)\ne0$ eventually and
\[
 \frac{\widehat v_a(z)}{\widehat v_a(0)}
 \longrightarrow\frac{\Xi(z)}{\Xi(0)}
 \quad\text{locally uniformly on }\mathbb C.
\]
The scalar is fixed by the stated $L^2$ normalization and then by the
value at zero; phase/scale and translations must not be left free.
If the real-zero-ground-transform theorem's complete hypotheses hold
cofinally as well, this convergence rules out a nonreal zero of $\Xi$:
a small closed disk about such a zero, disjoint from the real axis and
with zero-free boundary, would by uniform convergence and the argument
principle eventually contain the same positive number of zeros of
$\widehat v_a$, a contradiction. This is a conditional implication,
not a new RH proof. Cofinal simplicity/evenness and applicability of
that theorem are themselves additional to the single $\lambda=4$ result.

\subsection*{An exact abstract nonselection model}

The next construction tests precisely the proposed \emph{abstract}
premises: complete compact-resolvent operators on increasing intervals;
simple even nonnegative ground functions; entire ground transforms with only real
zeros; monotone positive ground energies tending to zero; a fixed
candidate with vanishing full residual; dense fixed tests with vanishing
residual; ordinary strong-resolvent collapse; and even uniformly bounded
support of the ground functions. It does \emph{not} impose the
support-consistent arithmetic Weil form, the prime identities, the
CvS distribution representation, or a positivity-improving semigroup.
Its ground functions are nonnegative, not strictly positive on all of
the increasing interval. It cannot disprove an arithmetic
selection theorem using those additional structures.

\begin{proposition}[Abstract nonselection despite fast absolute residuals]
\label{prop:ns43-ccm-nonselection}
There is a family satisfying all the abstract properties just listed,
including uniformly tight ground supports and the explicit Gaussian
cutoff candidate, whose phase-fixed normalized ground transforms have
two different subsequential limits. The absolute residual rate can be
made arbitrarily fast by a positive scalar rescaling.
\end{proposition}
\begin{proof}
Let $b=1_{[-1/4,1/4]}$ and let $g_1=b*b*b*b/\|b*b*b*b\|_2$.
Then $g_1$ is nonnegative and even, lies in $H^2(\mathbb R)$, and is
supported in $[-1,1]$. Put $g_2(x)=2^{-1/2}g_1(x/2)$, supported in
$[-2,2]$. Both are unit vectors, distinct and not scalar multiples.
Their entire transforms are nonzero constants times
$\operatorname{sinc}(z/4)^4$ and $\operatorname{sinc}(z/2)^4$,
respectively, where $\operatorname{sinc}z=\sin z/z$. Every zero is real.
Their normalized transforms have distinct zero sets.

For integers $n\ge3$, let $D_n=-d^2/dx^2$ be the Dirichlet Laplacian
on $H_n=L^2(-n,n)$, with domain $H^2(-n,n)\cap H_0^1(-n,n)$.
Its unit first eigenvector is
$h_n(x)=n^{-1/2}\cos(\pi x/(2n))$.
Let $j(n)=1$ for even $n$ and $2$ for odd $n$. Define
\[
 w_n=\frac{h_n-g_{j(n)}}{\|h_n-g_{j(n)}\|},\qquad
 U_n=I-2|w_n\rangle\langle w_n|,\qquad
 A_n=n^{-1}U_n(I+D_n)U_n.
\]
The denominator is nonzero; its square is
$2-2\langle h_n,g_{j(n)}\rangle\to2$.
Thus $U_n$ is an even-parity-preserving unitary involution with
$U_nh_n=g_{j(n)}$. Because $w_n\in\mathcal D(D_n)$, it preserves
that domain. The complete $A_n$ is positive, self-adjoint with compact
resolvent, and its simple even ground is exactly $g_{j(n)}$ with
\[
 \mu_n=\frac1n\left(1+\frac{\pi^2}{4n^2}\right)\downarrow0,
 \qquad \lambda_{2,n}-\mu_n=\frac{3\pi^2}{4n^3}>0.
\]
All ground transforms have only real zeros, but the normalized ground
functions and their normalized entire transforms alternate between two
different limits. Positivity and a fixed positive phase do not repair this.

For each fixed $f\in C_c^\infty(\mathbb R)$ and all sufficiently large
$n$, $f\in\mathcal D(D_n)$. Moreover
$\|(I+D_n)w_n\|$ is uniformly bounded: the $g_j$ have fixed finite
$H^2$ norms, while $\|D_nh_n\|=\pi^2/(4n^2)$ and the denominators
are bounded away from zero. Consequently
\[
 \|A_nf\|\le\frac1n\left(\|f-f''\|
        +2\|f\|\,\|(I+D_n)w_n\|\right)\longrightarrow0.
\]
The same estimate holds for the fixed candidate $g_1$. It also holds
uniformly for the explicit Gaussian cutoff candidate $p_n$ used above:
its $H^2$ norms are uniformly bounded, its cutoff normalization stays
away from zero, and $\|p_n\|=1$. Thus $\|A_np_n\|=O(1/n)$ while
$p_n\to\phi$ in every exponentially weighted $L^1$ norm. Even the
known explicit Gaussian candidate and its transform identity do not
select the alternating grounds of this abstract operator family.
Let $\widetilde A_n=A_n\oplus0$ on $L^2(\mathbb R)$ by zero extension.
For every nonreal $z$ and every such fixed $f$,
\[
 \|(\widetilde A_n-z)^{-1}f+z^{-1}f\|
 \le\frac{\|\widetilde A_nf\|}{|z||\Im z|}\to0.
\]
Density of $C_c^\infty$ and the uniform self-adjoint resolvent bound
extend this to all of $H$. Thus the ordinary strong-resolvent limit
is zero. This model meets every listed abstract premise, including
weighted tightness of the grounds, yet supplies no unique selection.
Its ground multiplicity is one at every $n$; the failure is not caused
by a degenerate eigenspace at a fixed window.
The scalar $1/n$ can be replaced throughout by any positive sequence
$\theta_n\downarrow0$. The eigenvalues and residuals are rescaled,
the ground functions are unchanged, and the residual estimate becomes
$O(\theta_n)$. Thus even an arbitrarily prescribed rapid absolute
residual rate cannot replace the missing relative estimate. This
rescaling does not restore arithmetic support consistency.
\end{proof}


\paragraph{Ordinary norm is also the wrong transform topology by itself.}
Even if a selection estimate gives $f_a\to\phi$ in $L^2$, locally
uniform entire convergence still needs a rate or a tail condition.
For any fixed $R>0$, add to the positive cutoff source a positive even
unit bump supported near $\pm(a-1)$ with coefficient $e^{-Ra/2}$,
and renormalize. The resulting compactly supported unit functions
converge to $\phi$ in $L^2$, while their transforms at $iR$ grow like
a positive constant times $e^{Ra/2}$. This is a separate topology
counterexample; it does not claim the real-zero property or the
arithmetic ground equation.

\subsection*{A restriction on renormalized profile selection}

A possible alternative is to renormalize the shrinking spectrum and
prove a nonzero profile form with a unique selected direction.
The ordinary zero-resolvent limit cannot provide this profile.
\emph{Scope reset for this subsection only:} $q_a$ now denotes the
\emph{full both-parity} physical Weil form on $L^2(I_a)$, not merely
the even-sector form used for the preceding selection statements.
Here $J_a$ is zero extension into the full $H=L^2(\mathbb R)$;
$\alpha_a$ remains the Rayleigh value of the even cutoff $p_a$.
Thus the use of all real translates and their density in full $H$
is exactly the setting of \ref{lem:v136-total-radicals}.
Here is a precise obstruction to choosing too coarse a scale.
\begin{proposition}[A coarse rescaled liminf cannot select one radical]
\label{prop:ns43-ccm-coarse-profile}
Let $s_a>0$ and suppose
\[
 \varepsilon_a/s_a\to0,\quad |\alpha_a|/s_a\to0,
 \qquad \varepsilon_a=Ca e^{-e^a/100},
\]
where each fixed finite translate combination may have its own constant
$C$, as in v1.35--v1.36. Suppose a nonnegative lower-semicontinuous
quadratic form $\mathcal E$ on $H$ satisfies the proposed strong
liminf inequality
\[
 \mathcal E(f)\le\liminf_a
 \frac{q_a[f_a]-\alpha_a\|f_a\|^2}{s_a}
 \quad\text{whenever }J_af_a\to f.
 \tag{S5}
\]
Then $\mathcal E\equiv0$ on $H$.
\end{proposition}
\begin{proof}
For every fixed $f$ in the dense translated-radical span, insert its
smooth cutoff $f_a$. The complete residual implies the right side is
zero, so $\mathcal E(f)=0$. Lower semicontinuity and density then imply
$\mathcal E\equiv0$ on $H$. Therefore no such coarse rescaled limit
can have one-dimensional kernel $\mathbb C\phi$.
\end{proof}
This is a scoped implication about (S5), not a prohibition on finer
scales, different compactness mechanisms, or arithmetic ground selection.

Finally, in the actual support-consistent Weil family, write
$\mu_a^{\mathrm{full}}=\inf\sigma(\mathsf W_{e^a})
=\min(\mu_a^{\mathrm{even}},\mu_a^{\mathrm{odd}})$ on the full
$L^2(I_a)$. A uniform lower bound on this \emph{complete, both-parity}
ground energy is already the bounded-floor
target of v1.36, and $\mu_a^{\mathrm{full}}\to0$ is equivalent to
full both-parity positivity
via monotonicity and the known source upper bound. Those must not be
rebranded as weak new selection assumptions. Neither conclusion alone
has been shown here to identify the Gaussian ground profile. The
abstract countermodel establishes no arithmetic nonimplication.


\section{Fine translated-radical blocks and the rank cost of positive coercivity (v1.53)}
\label{sec:v153-fine-block}
This section strengthens the existing near-zero multiplicity obstruction.
It does not prove a lower bound for the complete Weil form. The change
from v1.35 is essential: the translate centers now remain in a fixed
physical interval and their spacing shrinks. This improves the complete
boundary residual to an exponential in $-\lambda^2$, while the finite
Gram is poorly conditioned. An elementary interpolation estimate pays
that conditioning cost explicitly. A growing source block must therefore
be assessed as a block, with its full Gram and operator residual.

The resulting complete near-zero subspace has dimension asymptotic to
$\lambda^2/(20000\log\lambda)$, in both parities together. It rules out
a positive uniform or polynomial gap after removal of
$o(\lambda^2/\log\lambda)$ directions. This is a failure of that positive
coercivity objective, not a negative Weil direction. The weaker uniform
negative-floor target remains open; its gap-free transfer was already
proved in \ref{prop:v118-gap-free} and \ref{cor:v136-complement-floor}.
No new numerical window, Fourier cutoff or tail metric is introduced.

% NS-46 analytic proof draft. No numerical or cofinal sign claim.
\subsection*{Finite fine-lattice Grams without an infinite-lattice assumption}

\begin{lemma}[An elementary finite-arc Gram lower bound]
\label{lem:ns46-finite-arc}
Let $\phi\in L^2(\mathbb R)$ have unitary Fourier transform satisfying
$|\widehat\phi(\xi)|^2\ge c_0>0$ on $[-r,r]$, with $r>0$.
Let $m\ge2$, $\delta>0$ and $0<r\delta\le\pi/2$.
The Gram matrix of the $m$ translates
$\phi(\cdot-t_0-j\delta)$, $0\le j<m$, obeys
\begin{equation}\label{eq:ns46-gram-bound}
 \Gamma_{m,\delta}\succeq\gamma_{m,\delta}I,
 \qquad
 \gamma_{m,\delta}:=\frac{c_0r}{m}
       \left(\frac{r\delta}{4\pi e}\right)^{2(m-1)}.
\end{equation}
No positive lower bound for an infinite-lattice Gram is assumed.
The displayed bound is allowed to deteriorate with both $m$ and $\delta$.
\end{lemma}
\begin{proof}
For coefficients $c_j$ put $P(z)=\sum_{j=0}^{m-1}c_jz^j$ and
$\theta_0=r\delta$, $\ell=\theta_0/m$. The uncut physical norm is at least
\[
 \frac{c_0}{\delta}\int_{-\theta_0}^{\theta_0}
                      |P(e^{i\theta})|^2\,d\theta.
\]
The sign of the phase and the common shift do not affect this inequality.
Write $I_P$ for the integral. For each $0\le k<m$, the interval
\[
 B_k=[-\theta_0+2k\ell,-\theta_0+(2k+1)\ell]
\]
has length $\ell$, and lies inside the integration arc. Choose
$\theta_k\in B_k$ with
$|P(e^{i\theta_k})|\le\sqrt{I_P/\ell}$, which is possible by averaging.
For $j\ne k$ the angles satisfy
\[
 |\theta_j-\theta_k|\ge |j-k|\ell,
 \qquad |\theta_j-\theta_k|\le2\theta_0\le\pi.
\]
Thus, with $z_k=e^{i\theta_k}$,
$|z_j-z_k|\ge(2\ell/\pi)|j-k|$.

Lagrange interpolation at these distinct nodes is exact for $P$.
On the full unit circle its numerator factors have modulus at most two,
so
\begin{align*}
 |P(z)|
 &\le\sqrt{I_P/\ell}
   \sum_{k=0}^{m-1}
   \frac{2^{m-1}}{(2\ell/\pi)^{m-1}k!(m-1-k)!}\\
 &=\sqrt{I_P/\ell}\,
   \frac{(2\pi/\ell)^{m-1}}{(m-1)!}.
\end{align*}
Parseval on the full circle bounds $\sum_j|c_j|^2$ by the square
of this supremum. Therefore
\[
 I_P\ge\ell
 \left((m-1)!\left(\frac{\ell}{2\pi}\right)^{m-1}\right)^2
 \sum_j|c_j|^2.
\]
The elementary factorial bound
$(m-1)!\ge((m-1)/e)^{m-1}$ and $(m-1)/m\ge1/2$ give
\[
 (m-1)!\left(\frac{r\delta}{2\pi m}\right)^{m-1}
 \ge\left(\frac{r\delta}{4\pi e}\right)^{m-1}.
\]
Multiplying by $c_0/\delta$ proves the lemma, for complex coefficients
as well as real ones. The result uses only a finite polynomial and a
single fixed Fourier interval; it makes no assertion about the zeros
of the transform on any infinite arithmetic progression.
\end{proof}

\paragraph{Application data for the actual source.}
The Gaussian arithmetic radical in v1.35 is real, even, smooth and
nonnegative, and not zero, as made explicit in v1.50. Its normalized
version $\phi$ has $\widehat\phi(0)>0$. Continuity therefore supplies
fixed $r,c_0>0$ as in the lemma. Their existence is analytic; no numerical
value or certified effective threshold is asserted here.

For $a=\log\lambda$ sufficiently large, take
\[
 J_a=\left\lfloor\frac{\lambda^2}{40000a}\right\rfloor,
 \qquad m_a=2J_a+1,\qquad \delta_a=1/J_a,
 \qquad t_{a,j}=j\delta_a\quad(-J_a\le j\le J_a).
\]
All centers are in the fixed interval $[-1,1]$.
The lemma applies after a common shift to the indices $0,\ldots,m_a-1$.
For sufficiently large $a$,
\[
 m_a\le\frac{\lambda^2}{10000a},\qquad
 \log\left(\frac{4\pi e}{r\delta_a}\right)\le3a.
\]
Consequently its uncut Gram lower bound satisfies
\begin{equation}\label{eq:ns46-gram-asymptotic}
 \gamma_a\ge\frac{c_0r}{m_a}
       \exp\left(-\frac{6\lambda^2}{10000}\right).
\end{equation}
This is a lower bound for the smallest ordinary Gram eigenvalue,
not for a Weil eigenvalue. In particular the conditioning is explicitly
paid for, not silently assumed uniform.


% Astra-2 / NS-46: complete residual, cutoff loss, exact parity and rank.
% Companion finite-arc Gram estimate is owned by the coordinator.
\subsection*{Dense finite radical blocks: complete residual and rank accounting}
\label{ns46-a2-sec-block}

\paragraph{Scope and input.}
All operators below are the complete physical $\mathsf W_{e^a}$,
with the domain, both pole terms, all prime rows and fixed smooth
cutoff $\chi_a$ of Proposition~\ref{prop:v135-growing-radical}.
Write $\phi=h/\|h\|_2$ for the even Gaussian arithmetic radical.
Choose fixed $r,c_0>0$ such that $|\widehat\phi(\xi)|^2\ge c_0$
for $|\xi|\le r$; these exist because $\widehat\phi(0)>0$.
Lemma~\ref{lem:ns46-finite-arc} supplies the following
uncut Gram lower bound for $m\ge2$ equally spaced translates,
spacing $d>0$, whenever $rd\le\pi/2$:
\begin{equation}\label{ns46-a2-eq-arc-input}
 K_m\succeq\gamma_m I,\qquad
 \gamma_m=\frac{c_0r}{m}
       \left(\frac{rd}{4\pi e}\right)^{2(m-1)}.
\end{equation}
The results here use this proved explicit bound as their input.
No lower bound for an infinite lattice with shrinking spacing is
being assumed. This is an analytic construction, not an interval
computation or a new finite-window certificate.

\begin{lemma}[Uniform complete tails for a fixed range of centers]
\label{ns46-a2-lem-fixed-centers}
Fix $T>0$. For $|t|\le T$, let
$g_{a,t}=\chi_a\phi(\cdot-t)$ and
$r_{a,t}=(1-\chi_a)\phi(\cdot-t)$. There is a constant $C_T$,
independent of $a,t$, such that, for all sufficiently large $a$,
\begin{align}
 \|r_{a,t}\|_{H^1}
  +\sup_x e^{2|x|}(|r_{a,t}(x)|+|r'_{a,t}(x)|)
 &\le C_Te^{-c_Te^{2a}},\notag\\
 \int_{|x|\ge a-1}|\phi(x-t)|^2\,dx
 &\le C_T^2e^{-2c_Te^{2a}},
 \quad c_T=\frac\pi4e^{-2(T+1)},
 \label{ns46-a2-eq-tails}\\
 \|\mathsf W_{e^a}g_{a,t}\|_2
 &\le C_T\sqrt a\,e^{2a}e^{-c_Te^{2a}}.
 \label{ns46-a2-eq-column-residual}
\end{align}
\end{lemma}
\begin{proof}
On the tail support $|x|\ge a-1$, one has
$|x-t|\ge a-1-T$. Put $X=e^{2|x-t|}$, so
$X\ge e^{2a-2(T+1)}$ and $e^{2|x|}\le e^{2T}X$.
The derivative bounds \eqref{eq:v135-gaussian-tail} and the uniformly
bounded derivatives of the fixed cutoff bound the weighted pointwise
tails by $C_TXe^{-\pi X/2}\le C_Te^{-\pi X/4}$.
Integrating the squared Gaussian bound after the substitution
$X=e^{2|x-t|}$ proves the $H^1$ and full tail-mass estimates, with
the same weaker exponent. Constants are uniform for all $|t|\le T$.

Polarized radicality gives, distributionally on $I_a$,
$\mathsf W_{e^a}g_{a,t}=-\mathcal W r_{a,t}$.
The archimedean multiplier is bounded by $C(1+|\xi|)$, so its
$L^2$ contribution is controlled by the $H^1$ tail. For $x\in I_a$,
each of the two full prime sums is bounded by
\[
 C_Te^{-c_Te^{2a}}e^{2|x|}
         \sum_{n\ge2}\frac{\Lambda(n)}{n^{5/2}}.
\]
This includes prime rows beyond $e^{2a}$; there is no finite-prime
truncation of the residual. Both pole integrals are bounded by
$C_Te^{-c_Te^{2a}}e^{|x|/2}$. Integration on $I_a$ gives
\eqref{ns46-a2-eq-column-residual}. Each $g_{a,t}$ is smooth and
compactly supported inside $I_a$, hence in the complete operator domain.
\end{proof}

\begin{proposition}[Gram loss and residual amplification]
\label{ns46-a2-prop-transfer}
Take any $m$ distinct centers in $[-T,T]$ whose uncut synthesis
Gram satisfies $K\succeq\gamma I$, $\gamma>0$.
Let $G$ synthesize the cutoff columns, $\Gamma=G^*G$, and
$\tau_a=C_Te^{-c_Te^{2a}}$. If
$m\tau_a^2\le\gamma/2$, then
\begin{equation}\label{ns46-a2-eq-transfer}
 \Gamma\succeq\gamma I/2,\qquad
 \operatorname{rank}G=m,\qquad
 \|\mathsf W_{e^a}P_G\|
 \le \sqrt{2m/\gamma}\,
         C_T\sqrt a\,e^{2a}e^{-c_Te^{2a}}.
\end{equation}
Here $P_G=G\Gamma^{-1}G^*$ is the ordinary orthogonal projection.
\end{proposition}
\begin{proof}
The Gram loss is positive semidefinite and
\[
 \operatorname{tr}(K-\Gamma)
 =\sum_j\int(1-\chi_a^2)|\phi(x-t_j)|^2\,dx
 \le m\tau_a^2.
\]
The full tail mass in \eqref{ns46-a2-eq-tails} is essential here:
one cannot replace $1-\chi_a^2$ by $(1-\chi_a)^2$ as an upper bound.
Thus $\Gamma\succeq\gamma I/2$. Since $G\Gamma^{-1/2}$ is an
isometry onto its range, the complete column residuals give
\[
 \|\mathsf W_{e^a}P_G\|
 =\|\mathsf W_{e^a}G\Gamma^{-1/2}\|
 \le\sqrt{2/\gamma}
       \left(\sum_j\|\mathsf W_{e^a}g_{a,t_j}\|^2\right)^{1/2}.
\]
This proves the displayed bound without a midpoint inverse or a
quadratic-value substitute for the operator residual.
\end{proof}

\begin{corollary}[A block of order $\lambda^2/\log\lambda$]
\label{ns46-a2-cor-large-block}
Put $\lambda=e^a$ and, for sufficiently large $a$, choose
\[
 J=\left\lfloor\frac{\lambda^2}{40000a}\right\rfloor,
 \quad m=2J+1,\quad d=1/J,\quad t_j=j/J\quad(-J\le j\le J).
\]
Using \eqref{ns46-a2-eq-arc-input}, the cutoff synthesis has
rank $m$ and its ordinary projection satisfies
\begin{equation}\label{ns46-a2-eq-large-residual}
 \|\mathsf W_{e^a}P_G\|\le\eta_a,
 \qquad \eta_a=C\exp(-\lambda^2/2000).
\end{equation}
The even and odd parts have exact dimensions $J+1$ and $J$,
and inherit the same residual bound. All constants and the eventual
threshold are independent of $a$; no numerical threshold is certified.
\end{corollary}
\begin{proof}
Use $T=1$. Then $c_T=\pi/(4e^4)>1/1000$.
Absorbing the polynomial prefactor in
\eqref{ns46-a2-eq-column-residual} gives a column residual at most
$Ce^{-\lambda^2/1000}$, and the full tail norm has the same bound.
For large $a$, $rd\le\pi/2$, $m\le\lambda^2/(10000a)$ and
$\log(4\pi e/(rd))\le3a$. Hence
\[
 \gamma_m^{-1/2}\le C\sqrt m\,
       \exp(3ma)\le C\sqrt m\exp(3\lambda^2/10000).
\]
In particular $m\tau_a^2/\gamma_m\to0$, so
Proposition~\ref{ns46-a2-prop-transfer} applies. Its resulting
bound is at most
$Cm\exp(-7\lambda^2/10000)\le C\exp(-\lambda^2/2000)$.

Reflection sends the $j$th cutoff column to the $-j$th column.
The coefficient embeddings consisting of the central column and
$(e_j+e_{-j})/\sqrt2$, or of $(e_j-e_{-j})/\sqrt2$, are
isometries of dimensions $J+1$ and $J$. The full Gram lower bound
makes both restricted syntheses injective. Their images are exactly
the two parity parts and their projections lie under $P_G$.
\end{proof}

\paragraph{Exact source accounting.}
An even source imposes no condition on the $J$ odd columns. In the
even block its orthogonality condition removes at most one dimension:
the resulting dimension is $J$ if its projection into that block is
nonzero, and $J+1$ otherwise. The lower dimension $J$ is enough for
the conclusion below; no identification of the cutoff column with
the actual source or unproved nonzero-overlap claim is needed.

\begin{corollary}[Necessary removal rank for positive coercivity]
\label{ns46-a2-cor-rank}
For every physical subspace $F_a$ with $\dim F_a<m$, there is a
unit $v_a\in\operatorname{Ran}G\cap F_a^\perp$ with
\[
 \|\mathsf W_{e^a}v_a\|\le\eta_a,
 \qquad |q_a[v_a]|\le\eta_a.
\]
Thus every positive coercivity constant on $F_a^\perp$ is at most
$\eta_a$. In particular removing $o(\lambda^2/\log\lambda)$
dimensions cannot leave a fixed positive or positive polynomial-in-
$\lambda^{-1}$ gap. The complete spectral projection onto
$[-2\eta_a,2\eta_a]$ has rank at least $2J+1$, with at least
$J+1$ even and $J$ odd eigenvalues counted with multiplicity.
\end{corollary}
\begin{proof}
Dimension gives the indicated unit vector, and the complete residual
gives its form bound. If the spectral projection had smaller rank,
a unit vector in the block orthogonal to its range would have
operator norm at least $2\eta_a$, contradicting
\eqref{ns46-a2-eq-large-residual}. Apply the same argument to each
parity block. The complete self-adjoint operators have compact
resolvent, as already established in the manuscript.
\end{proof}
With an actual even source constraint already imposed, fewer than
$J$ additional independent even constraints still leave such a vector;
fewer than $J$ odd constraints do so in the odd sector. More precisely
the relevant number is the rank of the constraint functionals restricted
to the indicated block. No identification with a sampler image or a
finite Fourier head is made.

\paragraph{What remains missing.}
The block construction supplies many complete near-zero directions and
their full residual; it supplies neither their energy signs nor a lower
bound on the infinite-dimensional orthogonal complement. The remaining
signed estimate is exactly
\[
 q_a[y]\ge-C_*\|y\|^2\quad
 (y\in\mathcal D(q_a)\cap\operatorname{Ran}(I-P_G)),
 \qquad C_*<\infty\ \hbox{uniformly on a cofinal family},
\]
in both parities. Together with the residual, the existing gap-free
transfer in v1.18/v1.36 would give a full uniform floor. Rewriting this
estimate in block coordinates is not new arithmetic input. Increasing
the block's rank does not bound negative eigenvalues remaining outside
it. The result is a stronger necessary rank obstruction for a positive
gap strategy; it does not exclude gap-free block methods, produce a
negative Weil direction, prove G2 or RH, or close the signed-floor route.


\section{Kernel exclusion, paired heat transport and arithmetic approximation (v1.54)}
\label{sec:v154-new-shapes}
This checkpoint tests three different intermediate mechanisms. It proves an
exact complete-domain nullvector reduction and a scoped regularity countermodel,
finite heat-zero transport identities with an explicitly conditional comparison
obstruction, and an unconditional arithmetic block-gain inequality. The missing
inputs are, respectively, arithmetic injectivity, suitable full heat transport,
and asymptotic residual correlation. None is supplied by changing coordinates.
The countermodels and implications below do not produce a negative Weil
direction. No new window or tail metric is computed; G2 and RH remain open.

\subsection*{The complete nullvector equation without an $H^1$ assumption (NS-51)}

\paragraph{Classification and scope.}
The next statements are exact identities and proved implications, with a
named open arithmetic injectivity input. They concern the complete
fixed-window operator, not a finite head. The new regularity countermodel
rules out one generic inference, not injectivity of the arithmetic
operator. No new window or tail metric is computed. G2 and RH remain open.

Put $I_a=(-a,a)$, and let $E_a$ denote zero extension from $I_a$.
Write $V_a=\mathcal D(q_a)$ for the complete physical form domain:
\[
 V_a=\left\{f\in L^2(I_a):
  \int_{\mathbb R}\log(2+|\xi|)|\widehat{E_af}(\xi)|^2\,d\xi<\infty\right\}.
\]
Let $A_a$ be its self-adjoint form operator. These are the unshifted,
complete objects of Proposition~\ref{prop:v118-log-dirichlet}; no source
projection is imposed. Let $\mathfrak g$ be Suzuki's real even screw
function, normalized by $\mathfrak g(0)=0$ and, for $t>0$,
\begin{align}
 \mathfrak g(t)={}&-4(e^{t/2}+e^{-t/2}-2)
  +\sum_{n<e^t}\frac{\Lambda(n)}{\sqrt n}(t-\log n)
  -\frac t2\bigl(\psi(1/4)-\log\pi\bigr)\notag\\
 &-\frac14\left\{\Phi(1,2,1/4)
             -e^{-t/2}\Phi(e^{-2t},2,1/4)\right\}.
 \label{eq:ns51-screw}
\end{align}
Here $\Phi(z,s,v)=\sum_{j\ge0}z^j/(j+v)^s$ for $|z|<1$,
with its convergent value at $(1,2,1/4)$. Including a prime-power
endpoint instead makes no difference in this formula, since its ramp
vanishes at that endpoint. Suzuki's distributional identity is
$A_a\varphi=(-\mathfrak g''*E_a\varphi)|_{I_a}$ for
$\varphi\in C_c^\infty(I_a)$~\cite[Sections 2.5 and 3]{Suzuki2026}.
All convolutions below are local convolutions with a compactly supported
function or distribution; no global temperedness of $\mathfrak g$ is
assumed.

\begin{proposition}[Affine screw potential criterion on the complete domain]
\label{prop:ns51-affine-potential}
For $f\in L^2(I_a)$ set
\[
 F_f(x)=\int_{-a}^a\mathfrak g(x-y)f(y)\,dy,\qquad
 U_f(x)=\int_{-a}^a\mathfrak g'(x-y)f(y)\,dy.
\]
Then $F_f\in C^1(\mathbb R)$ and $F_f'=U_f\in C(\mathbb R)$.
For $f\in V_a$ the following conditions are equivalent:
\begin{enumerate}[label=(\roman*)]
 \item $f\in\mathcal D(A_a)$ and $A_af=0$;
 \item $F_f$ is affine on $I_a$;
 \item $U_f$ is constant on $I_a$.
\end{enumerate}
Let $\Pi_{\rm aff}$ be the ordinary $L^2(I_a)$ projection onto
$\operatorname{span}\{1,x\}$ and $P_0$ the projection onto mean-zero
functions. The operators
\[
 \mathcal C_a f=(I-\Pi_{\rm aff})F_f|_{I_a},\qquad
 \mathcal T_a f=P_0 U_f|_{I_a}
\]
are compact on $L^2(I_a)$, and
\begin{equation}\label{eq:ns51-kernel-reduction}
 \ker A_a=V_a\cap\ker\mathcal C_a=V_a\cap\ker\mathcal T_a.
\end{equation}
For even $f$ the affine potential is constant and $U_f=0$; for odd
$f$ it is $b x$ and $U_f=b$. The constant $b$ must not be discarded.
\end{proposition}
\begin{proof}
The screw function is locally absolutely continuous. Its derivative
has only a logarithmic singularity at zero and finitely many bounded
jumps on every compact interval. Indeed differentiation of
\eqref{eq:ns51-screw} gives, almost everywhere for $t>0$,
\begin{equation}\label{eq:ns51-screw-derivative}
 \mathfrak g'(t)=-4\sinh(t/2)
 +\sum_{n<e^t}\frac{\Lambda(n)}{\sqrt n}
 -\frac12(\psi(1/4)-\log\pi)
 -\frac12e^{-t/2}\Phi(e^{-2t},1,1/4),
\end{equation}
and odd reflection gives its value for $t<0$.
In particular $\mathfrak g'\in L^2_{\rm loc}(\mathbb R)$.
Translation continuity in local $L^2$ and Cauchy--Schwarz show that
$U_f$ is continuous, locally uniformly in $x$. The distributional
identity $F_f'=U_f$ then proves $F_f\in C^1$.

For $f\in V_a$ and $\varphi\in C_c^\infty(I_a)$, the core identity
and symmetry give, in the first-slot-linear convention,
\[
 q_a(f,\varphi)=\langle f,A_a\varphi\rangle
 =-\int_{I_a}F_f(x)\overline{\varphi''(x)}\,dx
 =\int_{I_a}U_f(x)\overline{\varphi'(x)}\,dx.
\]
A nullvector therefore has $F_f''=0$ as a distribution on $I_a$,
which is equivalent to the displayed affine and constant conditions.
Conversely these conditions annihilate $q_a(f,\varphi)$ on the smooth
form core. The form norm, after a fixed-window positive shift, makes
$q_a(f,\cdot)$ continuous. Core density therefore gives
$q_a(f,v)=0$ for every $v\in V_a$. The representation theorem for
closed forms implies $f\in\mathcal D(A_a)$ and $A_af=0$.
The kernels $\mathfrak g(x-y)$ and $\mathfrak g'(x-y)$ are square
integrable on $I_a^2$; for the latter the squared kernel norm is
\[
 \int_{-2a}^{2a}(2a-|t|)|\mathfrak g'(t)|^2\,dt<\infty.
\]
Thus their integral operators are Hilbert--Schmidt, and bounded
projections preserve compactness. Evenness of $\mathfrak g$ gives
the two parity conclusions.
\end{proof}

\paragraph{What changes in the derivative reduction.}
For $f\in H_0^1(I_a)$ one can write the familiar equation
$G_a Df=0$ with $D=i\,d/dx$ and the ordinary mean-zero screw operator
$G_a$. For general $f\in V_a$, $D(E_af)$ is only guaranteed to be a
compactly supported $H^{-1}$ distribution. Integration by parts in
this distributional sense gives
\[
 \mathfrak g*D(E_af)=i\,\mathfrak g'*E_af=iU_f.
\]
This is the extension on this particular derivative range; it does
not assert a bounded extension of $G_a$ to every $H^{-1}$ distribution.
It retains any boundary contributions of zero extension. Therefore
$\mathcal T_af=0$ is meaningful without making $Df$ an ordinary
$L^2$ vector. Suzuki's shifted derivative-completion and generalized
eigenvalue construction already addresses this distinction
\cite[Sections 8.3--8.5]{Suzuki2026}; the proposition gives its
zero-eigenvalue question as a local continuous integral equation.
It is a derived formulation, not a newly discovered RH equivalence.

\begin{proposition}[A zero eigenvalue does not supply generic $H_0^1$ regularity]
\label{prop:ns51-no-generic-bootstrap}
There is an interval $I_r$, a bounded real self-adjoint rank-one
operator $K$ on $L^2(I_r)$, and a nonzero even
$\tau\in\mathcal D(\tfrac12L_\Delta^{\rm Dir})$ such that
\[
 (\tfrac12L_\Delta^{\rm Dir}+K)\tau=0,
 \qquad \tau\notin H_0^1(I_r).
\]
The operator $K$ is bounded on every $L^p(I_r)$, $1\le p\le\infty$,
and $\tau$ has the continuous zero extension and logarithmic boundary
rate appearing in Theorem~\ref{thm:v118-eigen-endpoint-zero}.
This is a countermodel for inference from the principal operator and
bounded perturbation alone, not a counterexample for the arithmetic
$\mathcal K_\lambda$.
\end{proposition}
\begin{proof}
Choose $0<r<e^{-\gamma_{\rm E}}$. The one-dimensional specialization
of \cite[Theorem 1.2]{HLSLogBoundary} supplies the positive even torsion
function $\tau$ with $L_\Delta^{\rm Dir}\tau=1$ and two-sided bounds
$\tau(x)\asymp\ell(r-|x|)^{1/2}$. Its weak equation with right side
$1\in L^2(I_r)$ puts it in the operator domain. Set
\[
 M=\int_{I_r}\tau(x)\,dx>0,\qquad
 Kf=-\frac{1}{2M}\left(\int_{I_r}f(x)\,dx\right)\mathbf1_{I_r}.
\]
This is bounded, real, self-adjoint and rank one, and $K\tau=-1/2$.
If $\tau\in H_0^1(I_r)$, absolute continuity and Cauchy--Schwarz
would give $|\tau(x)|\le\sqrt{r-|x|}\,\|\tau'\|_2$ near either
endpoint. This contradicts its logarithmic lower bound. The two-sided
boundary estimate also supplies the stated continuous extension.
\end{proof}

\paragraph{Named open input: affine-potential injectivity (API).}
For every finite $a>0$, require
\begin{equation}\label{eq:ns51-api}
 f\in V_a,\quad (\mathfrak g*E_af)|_{I_a}\in
 \operatorname{span}\{1,x\}\quad\Longrightarrow\quad f=0.
\end{equation}
The parity versions must use the constants specified above. No
injectivity estimate for the actual arithmetic kernel is proved here.
Injectivity on all of $L^2(I_a)$ would be sufficient but is a stronger
statement than the one established as equivalent to the physical
kernel condition; no regularity upgrade is being assumed.

\begin{proposition}[Sign continuation with an explicit missing input]
\label{prop:ns51-sign-continuation}
Assume the complete fixed-window realizations above, compact resolvent,
continuity of their lowest eigenvalue in $a$, positivity for sufficiently
small $a$, and API for every finite $a>0$. Then every $A_a$ is strictly
positive at fixed $a$. Consequently the complete compact-test Weil
form is nonnegative and Weil's criterion implies RH. No common positive
spectral gap is required. This is a conditional implication; API is open.
\end{proposition}
\begin{proof}
The first four structural inputs, excluding API, are supplied by the
fixed-window theory and \cite[Theorems 1.3--1.4]{Suzuki2026}.
If the continuous lowest eigenvalue ever became nonpositive, it would
be zero at some finite intermediate window. Compact resolvent makes
that lower edge an eigenvalue. Proposition~\ref{prop:ns51-affine-potential}
then contradicts API. Every smooth compact test lies in a finite window,
so its energy is nonnegative. Apply the classical Weil criterion.
\end{proof}

\paragraph{Scope of the finding.}
The unjustified $H_0^1$ shortcut has been replaced by a valid weak
integral formulation. The generic bootstrap is false, while its truth
for a hypothetical arithmetic nullvector has not been decided. The
remaining obstacle is uniqueness for this specific signed arithmetic
integral equation, not a proved negative direction and not a required
uniform inverse bound. Complete even and odd sectors have both been
retained; no source-complement argument is substituted. Growing
near-zero blocks are consistent with fixed-window injectivity.


% NS-52. Exact finite identities and a conditional obstruction; no RH assumption.
% Standalone fragment: requires amsmath, amssymb, amsthm, hyperref and
% proposition/proof environments. No manuscript version is assigned here.
\subsection*{NS-52: paired heat-zero transport and its first obstruction}

\paragraph{Normalization and claim scope.}
Use the conventional de Bruijn--Newman family
\[
 H_t(z)=\int_0^\infty e^{tu^2}\Phi(u)\cos(zu)\,du,
 \qquad H_0(z)=\tfrac18\xi(\tfrac12+iz/2),
\]
as in Polymath15, equations (1)--(4) and Proposition 3.1
(\url{https://arxiv.org/html/1904.12438#S3}).
Put $Z_t(w)=8H_t(2w)$ and $\kappa=1/4$. Then
\begin{equation}\label{eq:ns52-normalization}
 Z_0(w)=\Xi(w):=\xi(\tfrac12+iw),\qquad
 \partial_t Z_t=-\kappa\partial_w^2 Z_t.
\end{equation}
Thus the heat coefficient in the ordinary zeta-ordinate coordinate is
$1/4$, not $1$. All zero sums below count multiplicity. No assertion
about the locations of the zeros of $Z_0$ is assumed.

For $f,g\in C_c^\infty(-a,a)$ define the nonunitary transform and
its analytic sesquilinear product by
\[
 F_f(w)=\int_{-a}^a f(x)e^{-iwx}\,dx,\qquad
 F_g^*(w)=\overline{F_g(\overline w)},\qquad
 \Psi_{f,g}(w)=F_f(w)F_g^*(w).
\]
The second factor is an entire function of $w$. Replacing it by
$\overline{F_g(w)}$ off the real axis would give the wrong form.
With $\mathcal Z_t$ the zero multiset of $Z_t$, put
\begin{equation}\label{eq:ns52-zero-form}
 q_{a,t}(f,g)=\sum_{w\in\mathcal Z_t}\Psi_{f,g}(w).
\end{equation}
At $t=0$, $\rho=1/2+iw$ and
$\rho^\sharp=1-\overline\rho=1/2+i\overline w$ identify this
exactly with the manuscript's first-slot-linear Weil zero formula:
$q_{a,0}(f,g)=q_a(f,g)$. There is no $2\pi$ factor with the
displayed nonunitary transform. Using the unitary transform instead
would put $2\pi$ in front of the sum. The factor $8$ in $Z_t$ does
not change any zero or its multiplicity. The completed function has
no poles; no additional pole term is to be added to this zero-side
definition. This is not an assertion that the undeformed geometric
prime and pole terms can be heat-evolved separately.

For completeness, the sum in \eqref{eq:ns52-zero-form} is absolutely
convergent on this smooth core for every $t\ge0$, locally uniformly
in $t$. The de Bruijn strip theorem (Polymath15, Theorem 3.2) and
the unconditional initial critical strip give
$|\operatorname{Im} w|\le1/2$. For $0\le t\le T$, Gaussian
majorization of the super-exponentially decreasing $\Phi$ gives
$|Z_t(w)|\le C_T e^{C_T|w|^2}$ and $Z_t(0)\ge Z_0(0)>0$.
Jensen's formula therefore bounds the zero count in $|w|\le R$
by $C_T(1+R)^2$, uniformly in $t$. Repeated integration by parts
gives arbitrary inverse powers of $|\operatorname{Re}w|$ for both
test transforms on the fixed strip. This proves convergence and
uniformly small tails. Local contour integrals count any finite
cluster continuously through collisions, so the complete form is
continuous in $t$ on the smooth core. This argument does not justify
termwise differentiation of the infinite sum or claim a closed
deformed form on the original complete form domain.
If all zeros at a fixed $t_0$ are real, then
\[
 q_{a,t_0}[f]=\sum_{w\in\mathcal Z_{t_0}}|F_f(w)|^2\ge0.
\]
For example the classical de Bruijn theorem supplies such $t_0\ge1/2$.

\begin{proposition}[Finite-cluster transport, including a collision]
\label{prop:ns52-cluster}
Let $Z(t,w)$ be holomorphic near $J\times\overline\Omega$, where
$J$ is an interval of real times and $\Omega$ is a bounded domain
with piecewise smooth positively oriented boundary $\Gamma$.
Assume $Z$ has no zero on $\Gamma$ for $t\in J$ and satisfies
$\partial_t Z=-\kappa Z_{ww}$ with constant $\kappa>0$.
For an entire $\Psi$, define the finite sum $S_\Gamma(t)$ of
$\Psi$ over the zeros in $\Omega$, counted with multiplicity.
Then $S_\Gamma$ is differentiable through all internal collisions and
\begin{align}
 S_\Gamma(t)&=\frac1{2\pi i}\int_\Gamma
                   \Psi(w)\frac{Z_w}{Z}(t,w)\,dw,
                   \label{eq:ns52-contour}\\
 S_\Gamma'(t)&=\frac{\kappa}{2\pi i}\int_\Gamma
                   \Psi'(w)\frac{Z_{ww}}{Z}(t,w)\,dw.
                   \label{eq:ns52-contour-derivative}
\end{align}
At a time with $m$ simple cluster zeros $w_1,\ldots,w_m$, set
$A(w)=Z(t,w)/\prod_{j=1}^m(w-w_j)$ and $b=A'/A$ locally in
the cluster. Then
\begin{equation}\label{eq:ns52-divided-difference}
 S_\Gamma'
 =2\kappa\sum_{1\le j<k\le m}
       \frac{\Psi'(w_j)-\Psi'(w_k)}{w_j-w_k}
   +2\kappa\sum_{j=1}^m b(w_j)\Psi'(w_j).
\end{equation}
If the cluster consists at a particular time of one zero $c$ of
order $m$, write $Z(t,w)=(w-c)^m A(w)$ with $A(c)\ne0$. Its
derivative at that time is exactly
\begin{equation}\label{eq:ns52-multiple}
 S_\Gamma'=
 \kappa m(m-1)\Psi''(c)
       +2\kappa m\frac{A'(c)}{A(c)}\Psi'(c).
\end{equation}
In particular, a double zero contributes
$2\kappa\Psi''(c)+4\kappa(A'/A)(c)\Psi'(c)$.
\end{proposition}
\begin{proof}
The argument principle proves \eqref{eq:ns52-contour}. The nonvanishing
boundary allows differentiation under its compact integral. Since
\[
 \partial_t(Z_w/Z)=-\kappa\partial_w(Z_{ww}/Z),
\]
integration by parts around $\Gamma$ gives
\eqref{eq:ns52-contour-derivative}. At a simple zero,
$w_j'=\kappa Z_{ww}(w_j)/Z_w(w_j)$, and factorization gives
\[
 w_j'=2\kappa\sum_{k\ne j}\frac1{w_j-w_k}
                         +2\kappa b(w_j).
\]
Pair the $j,k$ and $k,j$ contributions to $\sum_jw_j'\Psi'(w_j)$.
This proves \eqref{eq:ns52-divided-difference}. At an order-$m$ zero,
\[
 \frac{Z_{ww}}Z=\frac{m(m-1)}{(w-c)^2}
       +\frac{2m}{w-c}\frac{A'}A+\frac{A''}A.
\]
The residue after multiplying by $\Psi'$ is precisely
\eqref{eq:ns52-multiple}. No labeling or derivative of a colliding
individual root is needed.
\end{proof}

\paragraph{Complete symmetries and the term that must not be omitted.}
For real $t$, $Z_t$ is real entire and even. Choose the union of
cluster contours invariant under conjugation and $w\mapsto-w$,
counting each distinct zero once with its actual multiplicity.
Then the sum defines a Hermitian sesquilinear contribution and the
same identities apply. A double collision at a nonzero real $c$
has a reflected double collision at $-c$; both must be included
in the symmetry-complete contribution. At $c=0$ one would count
one cluster, not duplicate it (the actual $H_t$ has no imaginary-axis
zero). There is no RH assumption here.

For an isolated pair $w_\pm=c\pm d$, the internal term is
$2\kappa[\Psi'(w_+)-\Psi'(w_-)]/(w_+-w_-)$.
It converges to $2\kappa\Psi''(c)$ as $d\to0$.
The two divergent individual velocities therefore cancel exactly.
The remaining $b=A'/A$ term retains the interaction with the rest
of the entire function; it does not vanish by symmetry. The local
analytic factor suffices for the proof, so no additional rearrangement
of an infinite velocity sum is used here. For a full isolated double
cluster the limiting value is the double-zero formula above, not
merely $2\kappa\Psi''(c)$.

\begin{proposition}[The elementary estimate still grows with support]
\label{prop:ns52-support}
For $f,g\in L^2(-a,a)$ and $n\ge0$,
\begin{equation}\label{eq:ns52-pw}
 |\Psi_{f,g}^{(n)}(w)|
 \le 2a(2a)^n e^{2a|\operatorname{Im}w|}\|f\|_2\|g\|_2.
\end{equation}
Consequently, at a double collision with
$|\operatorname{Im}c|\le Y$ and $|(A'/A)(c)|\le B$, the bound
\begin{equation}\label{eq:ns52-support-cost}
 |S_\Gamma'|
 \le16\kappa(a^3+Ba^2)e^{2aY}\|f\|_2\|g\|_2
\end{equation}
holds. The $a^3$ dependence of a single real collision contribution
cannot be replaced by a constant for all compact smooth tests.
Specifically, suppose a double collision occurs at real $c$, and
choose a real odd $h\in C_c^\infty(-1,1)$ with $\|h\|_2=1$ and
$M_1=\int u h(u)\,du\ne0$. For
$f_a(x)=a^{-1/2}h(x/a)e^{icx}$ one has exactly
\begin{equation}\label{eq:ns52-collision-witness}
 \|f_a\|_2=1,\qquad
 S_\Gamma'[f_a]=4\kappa a^3 M_1^2.
\end{equation}
The occurrence of such a collision in the nonnegative-time arithmetic
family is not asserted.
\end{proposition}
\begin{proof}
Cauchy--Schwarz gives
$|F_f^{(j)}(w)|\le a^j\sqrt{2a}e^{a|\operatorname{Im}w|}\|f\|_2$.
Leibniz's formula proves \eqref{eq:ns52-pw}; use $n=1,2$ in
\eqref{eq:ns52-multiple} to obtain \eqref{eq:ns52-support-cost}.
For the witness, $F_{f_a}(c+v)=\sqrt a F_h(av)$, so
$F_{f_a}(c)=0$, $\Psi_{f_a,f_a}'(c)=0$, and
$\Psi_{f_a,f_a}''(c)=2a^3M_1^2$. Apply the double-zero formula.
This is a lower bound on a conditional individual cluster's derivative,
not a lower bound on the complete arithmetic derivative: other clusters
remain present.
\end{proof}

\paragraph{Passage to the whole time-dependent sum.}
Absolute convergence of \eqref{eq:ns52-zero-form} at each time
does not imply convergence of its differentiated cluster series.
One sufficient hypothesis is an exhausting family of finite zero
multisets tracked with multiplicity on $[0,t_0]$, with piecewise
contours as necessary, whose endpoint sums converge to the complete
forms, and whose integrated derivatives satisfy
\[
 \int_0^{t_0}S_n'(t)[f]\,dt\le C\|f\|_2^2
\]
for every compact smooth $f$, with one $C$ independent of the
exhaustion and its support. An integrable bound
$S_n'(t)[f]\le B(t)\|f\|_2^2$, $\int_0^{t_0}B(t)\,dt<\infty$,
would suffice. Full symmetry makes these diagonal quantities real.
These are hypotheses, not consequences of pair cancellation. They
would imply $q_a[f]\ge q_{a,t_0}[f]-C\|f\|_2^2$ by the fundamental
theorem of calculus followed by the endpoint limit. The next result
shows that this particular comparison has an additional obstruction.

\begin{proposition}[An unmatched real zero obstructs uniform
same-test comparison]\label{prop:ns52-unmatched}
Assume the archived cutoff-radical lemma: there exists
$\phi\in\mathcal S(\mathbb R)$ such that, for every fixed finite
linear combination $h$ of its real translates, smooth cutoffs
$\chi_a h\in C_c^\infty(-a,a)$ satisfy
\[
 \chi_a h\longrightarrow h\text{ in }L^2\text{ and }L^1,
 \qquad q_a[\chi_a h]\longrightarrow0.
\]
Assume $q_{a,t_0}$ is the same-test zero form
\eqref{eq:ns52-zero-form}, all its zeros are real, and it has a
zero $c$ of multiplicity $m_c\ge1$ with $F_\phi(c)\ne0$.
Then no finite $C\ge0$ independent of $a$ can satisfy
\begin{equation}\label{eq:ns52-rejected-comparison}
 q_a[f]\ge q_{a,t_0}[f]-C\|f\|_2^2
 \quad(f\in C_c^\infty(-a,a))
\end{equation}
on a cofinal family of windows.
\end{proposition}
\begin{proof}
Fix $T>0$. The Schwartz assumption implies
\[
 B_T:=\sum_{k\in\mathbb Z}
       |\langle\phi,\phi(\cdot-kT)\rangle|<\infty.
\]
For each fixed positive integer $M$ put
\[
 h_M=M^{-1/2}\sum_{j=1}^M e^{icjT}\phi(\cdot-jT).
\]
Counting equal differences of shifts gives
$\|h_M\|_2^2\le B_T$, whereas
$F_{h_M}(c)=\sqrt M F_\phi(c)$.
Since all deformed zeros are real,
$q_{a,t_0}[f]\ge m_c|F_f(c)|^2$ for every smooth test.
Apply \eqref{eq:ns52-rejected-comparison} to $\chi_a h_M$ and
let $a\to\infty$ along the prescribed cofinal family, keeping $M$
fixed. The $L^1$ convergence controls the single real Fourier
evaluation. Thus
\[
 m_c M|F_\phi(c)|^2\le C\|h_M\|_2^2\le CB_T.
\]
Now let $M\to\infty$, a contradiction. No cutoff estimate uniform
in $M$, global closedness, or lower semicontinuity of a deformed
form is used.
\end{proof}

The cutoff input is Lemma~\ref{lem:v136-total-radicals}: its ordinary
norm convergence and vanishing complete operator residual imply the
displayed form limit. Schwartz decay and dominated convergence supply
the additional $L^1$ limit. The source used there has Fourier transform
a nonzero constant times $\Xi$, by \eqref{eq:v150-mellin} and the
source normalization, so $F_\phi(c)\ne0$ means precisely that this
real deformed zero is not an original Xi zero. An unmatched zero
at a specified $t_0$ has not been certified or proved in this
checkpoint. The proposition states its hypothesis explicitly; it
does not assert an unconditional arithmetic closure. It also does
not assert any negative direction of the original Weil form.

\paragraph{Disposition.}
The finite algebraic test succeeds: multiplicities and complete
pairing remove the inverse-gap singularity. What had to change is
the normalization, the analytic off-real sesquilinear product, and
retention of the exterior-cluster logarithmic derivative. What is
still missing is a complete signed transport estimate with support
independent error. Moreover the proposed same-test positive-time
comparison is impossible whenever there is an unmatched real
deformed zero. A broader heat approach would therefore need a
different comparison or a justified transport of tests that respects
the original radical; no such transport is supplied. These are
finite exact identities, proved implications with explicit hypotheses,
and a named open transport input. No new window, tail metric,
numerical certificate, G2 assertion, or RH assertion is made.


% NS-53. Exact identities, proved implications, and an explicitly open input.
% No finite numerical approximation is offered as convergence evidence.
\subsection{NS-53: residualized integer-dilation blocks}
\label{sec:ns53-nb-blocks}

Work in the real Hilbert space $H=L^2((0,\infty),dx)$, with
$\chi=\mathbf 1_{(0,1)}$ and $\rho_n(x)=\{1/(nx)\}$ for integers $n\geq1$.
Real coefficients lose nothing for this real target; the complex version
uses conjugate transposes. Put
\[
 V_N=\operatorname{span}\{\rho_1,\ldots,\rho_N\},\qquad
 r_N=\chi-P_N\chi,\qquad E_N=\|r_N\|^2=d_N^2,
\]
where $P_N$ is the orthogonal projection onto $V_N$.
The strong Nyman--Beurling criterion is the established equivalence
$E_N\to0\Longleftrightarrow\mathrm{RH}$; see B\'aez-Duarte,
\href{https://arxiv.org/abs/math/0202141}{Theorem~1.1}.
No particular upper rate is required. Burnol's
\href{https://doi.org/10.1006/aima.2001.2066}{lower-bound theorem}
does not require a matching upper bound. This checkpoint proves no
convergence statement and makes no G2 or RH claim.

\begin{proposition}[Exact arithmetic entries and finite independence]
\label{prop:ns53-nb-entries}
Let $G_{mn}=\langle\rho_m,\rho_n\rangle$, $b_n=\langle\chi,\rho_n\rangle$,
and $\kappa=\log(2\pi)-\gamma$, where $\gamma$ is Euler's constant.
Then
\begin{align}
 G_{mn}&=\int_0^\infty\{y/m\}\{y/n\}\,\frac{dy}{y^2},
 & b_n&=\frac{\log n+1-\gamma}{n},
 &G_{nn}&=\frac{\kappa}{n}.\label{eq:ns53-nb-entries}
\end{align}
In particular, writing $a_j=\lfloor j/m\rfloor$ and
$d_j=\lfloor j/n\rfloor$, the following is an exact convergent series:
\begin{equation}
 G_{mn}=\frac1{mn}+\sum_{j=1}^{\infty}
 \left[\frac1{mn}-\left(\frac{d_j}{m}+\frac{a_j}{n}\right)
 \log\left(1+\frac1j\right)+\frac{a_jd_j}{j(j+1)}\right].
 \label{eq:ns53-nb-gram-series}
\end{equation}
Each bracket is a nonnegative cell integral. Every Gram matrix belonging
to finitely many distinct positive integers is positive definite.
Moreover $E_N>0$ for every finite $N$.
\end{proposition}

\begin{proof}
Substitute $y=1/x$ and integrate on $[j,j+1)$; both floor functions
are constant on that interval. This gives the integral and series formulas.
The change $y=nt$ gives the diagonal scaling and
\[
 b_n=\frac1n\left(\log n+\int_1^\infty\frac{\{t\}}{t^2}\,dt\right).
\]
The integral truncated at an integer $K$ is $1+\log K-H_K$.
Similarly, integrating $\{t\}^2/t^2$ on the first $M+1$ unit cells gives
\[
 2M+2-H_{M+1}-2M\log(M+1)+2\log(M!).
\]
Stirling's formula gives its limit $\kappa$. For finite independence,
a vanishing linear combination of $\{y/n\}$ is piecewise affine.
At the smallest index with nonzero coefficient its jump cannot be
cancelled by a smaller nonzero coefficient, a contradiction. If
$\chi=\sum_{n\leq N}c_n\rho_n$, comparison of the jumps at every positive
integer in the $y$ variable gives
$\sum_{n\mid j}c_n=-\mathbf 1_{j=1}$, with $c_n=0$ for $n>N$.
M\"obius inversion would force $c_n=-\mu(n)$ for every $n$, contradicted
by a prime greater than $N$. A finite-dimensional span is closed, so
$E_N>0$.
\end{proof}

\begin{proposition}[Exact block improvement, including all cross terms]
\label{prop:ns53-nb-block-gain}
Fix integers $1\leq N<M$. Let $G$ be the $N$ by $N$ old Gram matrix,
$C=(G_{mn})_{1\leq m\leq N,\,N<n\leq M}$, and
$D=(G_{mn})_{N<m,n\leq M}$.
Write $b=(b_m)_{m\leq N}$, $b_B=(b_n)_{N<n\leq M}$, and define
\begin{equation}
 c=G^{-1}b,\qquad h=b_B-C^{\mathsf T}c,\qquad
 S=D-C^{\mathsf T}G^{-1}C.
 \label{eq:ns53-nb-schur}
\end{equation}
Then $S$ is positive definite, and
\begin{equation}
 E_N-E_M=h^{\mathsf T}S^{-1}h
 =\sup_{w\ne0}\frac{|w^{\mathsf T}h|^2}{w^{\mathsf T}Sw}.
 \label{eq:ns53-nb-gain}
\end{equation}
The gain is zero if and only if $h=0$; strict finite improvement is not
asserted for every block. The coefficients for the enlarged optimum are
$(c-G^{-1}CS^{-1}h,S^{-1}h)$. In addition,
\begin{equation}
 E_N-E_M\ \geq\
 \frac{\sum_{n=N+1}^M|b_n-\sum_{m\leq N}c_mG_{mn}|^2}
 {\kappa\sum_{n=N+1}^M1/n}
 \ \geq\
 \frac{\|h\|_2^2}{\kappa\log(M/N)}.
 \label{eq:ns53-nb-trace-gain}
\end{equation}
These are finite, unconditional identities and inequalities, not estimates
of the numerator as $N$ grows.
\end{proposition}

\begin{proof}
Set $z_n=(I-P_N)\rho_n$ for $N<n\leq M$. Its Gram matrix is $S$,
and $\langle r_N,z_n\rangle=h_n$. Finite independence shows that these
$z_n$ are independent. The orthogonal decomposition
$V_M=V_N\mathbin{\oplus}\operatorname{span}\{z_n\}$ gives
\eqref{eq:ns53-nb-gain} and the coefficient update. Since
$0<S\leq D$, the largest eigenvalue of $S$ is at most
$\operatorname{tr}D=\kappa\sum_{n=N+1}^M1/n$. Therefore
$h^{\mathsf T}S^{-1}h\geq\|h\|_2^2/\operatorname{tr}D$.
The decreasing function $1/x$ gives
$\sum_{n=N+1}^M1/n\leq\log(M/N)$.
\end{proof}

\begin{remark}[Conditioning is retained]
Replacing $S$ by $D$ in the exact identity, or using $b_B$ instead of $h$,
would discard the old--new interaction. A small eigenvalue of $S$ can
make the exact gain much larger than the trace lower bound; it causes no
missing inverse error in the analytic lower bound. For a proposed real
block vector $w$, its exact gain after optimizing its scalar coefficient
is $|w^{\mathsf T}h|^2/(w^{\mathsf T}Sw)$. If
$\sum_{n=N+1}^M w_n/n=0$, the unprojected combination
$\sum w_n\rho_n$ vanishes on $x>1/(N+1)$. Its residualized combination
need not have this support, and its norm is still $w^{\mathsf T}Sw$.
\end{remark}

\begin{proposition}[Exact divisor cells and the raw M\"obius-prefix penalty]
\label{prop:ns53-nb-divisor-cells}
For arbitrary real coefficients $c_1,\ldots,c_N$, put
\[
 A=\sum_{n\leq N}\frac{c_n}{n},\qquad
 B_j=1+\sum_{n\leq N}c_n\left\lfloor\frac jn\right\rfloor
     =1+\sum_{k=1}^j\ \sum_{\substack{n\mid k\\n\leq N}}c_n.
\]
For $r=\chi-\sum c_n\rho_n$, one has $r(1/y)=-Ay$ on $(0,1)$
and $r(1/y)=B_j-Ay$ on $[j,j+1)$, $j\geq1$. Consequently
\begin{equation}
 \|r\|^2=A^2+\sum_{j=1}^{\infty}
 \left[A^2-2AB_j\log\left(1+\frac1j\right)
          +\frac{B_j^2}{j(j+1)}\right].
 \label{eq:ns53-nb-divisor-energy}
\end{equation}
The brackets are nonnegative cell integrals; they are not to be split
into three separately summed series. For the raw choice $c_n=-\mu(n)$,
\begin{equation}
 \left\|\chi+\sum_{n\leq N}\mu(n)\rho_n\right\|^2
 \geq (N+1)\left|\sum_{n\leq N}\frac{\mu(n)}n\right|^2.
 \label{eq:ns53-nb-mobius-penalty}
\end{equation}
\end{proposition}

\begin{proof}
Expand each fractional part. The energy identity is the integral of the
squared affine function on each unit cell. For the M\"obius choice,
$B_j=0$ for $1\leq j\leq N$, by
$\sum_{n\mid k}\mu(n)=\mathbf 1_{k=1}$. The interval $(0,N+1)$ alone
then contributes the displayed lower bound.
\end{proof}

\begin{remark}[What the arithmetic lemma does and does not supply]
Equation~\eqref{eq:ns53-nb-mobius-penalty} gives an explicit necessary
condition for this particular unsmoothed coefficient choice to converge:
$\sqrt N\sum_{n\leq N}\mu(n)/n\to0$. It is not a lower bound for the
optimal $E_N$. The divergence of the raw natural approximation is already
recorded in B\'aez-Duarte's Introduction; no novelty is claimed for its
failure. Smoothing or adaptive coefficients change the problem and must
be evaluated through the full residual and Gram matrix.
\end{remark}

\begin{proposition}[A sufficient block-correlation input at any rate]
\label{prop:ns53-nb-any-rate}
Fix an integer $q\geq2$, $N_j=q^j$, and real numbers
$0\leq\alpha_j<1$ for $j\geq j_0$, with $\sum_j\alpha_j=\infty$.
Use the actual optimal coefficients $c^{(N_j)}=G_{N_j}^{-1}b_{N_j}$.
Suppose the following \emph{residual block-correlation input} holds:
\begin{equation*}
 \sum_{n=N_j+1}^{qN_j}
 \left|\frac{\log n+1-\gamma}{n}
       -\sum_{m\leq N_j}c_m^{(N_j)}G_{mn}\right|^2
 \geq\kappa\log q\,\alpha_j E_{N_j}.
 \tag{RBC}\label{eq:ns53-nb-rbc}
\end{equation*}
Then $E_N\to0$. For example, the hypothetical bound (RBC) with
$\alpha_j=1/(2(j+1))$ would give
$E_{q^j}=O(j^{-1/2})$, hence $E_N=O((\log N)^{-1/2})$.
This slower-than-optimal rate would already meet the established criterion.
\end{proposition}

\begin{proof}
Equation~\eqref{eq:ns53-nb-trace-gain} gives
$E_{N_{j+1}}\leq(1-\alpha_j)E_{N_j}$. Iteration and
$1-t\leq e^{-t}$ prove convergence. Monotonicity supplies intermediate
$N$. For the example, summing $1/(2(j+1))$ yields the stated rate.
\end{proof}

\paragraph{Named open input and disposition.}
(RBC) is \emph{unproved}; it is a specific sufficient arithmetic estimate,
not asserted equivalent to RH or known to follow from RH. It may be
stronger than needed because the trace bound can lose the small
Gram eigenvalues that amplify the exact gain. The weaker exact target is
a nonsummable sequence of relative gains
$h_j^{\mathsf T}S_j^{-1}h_j/E_{N_j}$, which is equivalent to convergence
and is not itself a new theorem. The arithmetic work still missing is
a lower estimate for the corrected correlations, or for an explicitly
constructed block quotient, for the \emph{actual optimal residuals}.
Entry formulas, a finite collection of positive gains, and a finite list of distances
do not supply it. No such asymptotic gain estimate is proved here.
The missing step is the asymptotic gain estimate; this establishes
neither failure of (RBC) nor failure of the approximation object. No new Weil window, tail metric, G2 or RH
claim is introduced.

\section{Local uniqueness tests for the complete arithmetic operator (v1.55)}
\label{sec:v155-local-uniqueness}

The affine-potential criterion of Proposition~\ref{prop:ns51-affine-potential}
asks for the absence of a nullvector satisfying the equation on the
\emph{whole} window. Here we test three possible mechanisms for that
uniqueness. The complete arithmetic operator fails a local-patch
unique-continuation property, and its semigroup does not preserve the
ordinary cone of nonnegative functions. Neither statement gives a
nullvector on the whole window or a negative diagonal Weil direction.
An exact inward-dilation identity isolates the arithmetic overlap
terms still required at a hypothetical first zero. API, the uniform
floor, G2 and RH remain open. No numerical window or tail metric is added.

% NS-55: exact and scoped. Local-patch witnesses are not full nullvectors.
\subsection*{Unique continuation: the arithmetic shifts change the theorem}

Use $I_a=(-a,a)$, its complete operator $A_a$, and
\[
 \ell=\log2,\quad w_2=\frac{\log2}{\sqrt2},\quad
 \rho(t)=\frac{e^{-t/2}}{1-e^{-2t}},\quad
 R(t)=2\cosh(t/2)-\rho(t)\quad(t>0).
\]
If $f$ vanishes on an open $J\subset I_a$ separated from its support,
the exact operator there is
\begin{equation}\label{eq:ns55-off-support-operator}
 (A_af)(x)=\int_{I_a}R(|x-y|)f(y)\,dy
 -\sum_{1<n<e^{2a}}\frac{\Lambda(n)}{\sqrt n}
       \{E_af(x-\log n)+E_af(x+\log n)\}.
\end{equation}
This follows from Proposition~\ref{prop:v118-log-dirichlet}: local
diagonal terms vanish, and the logarithmic kernel and its correction
combine into $-\rho$. Equivalently it is the off-origin distribution
identity for $-\mathfrak g''$. In particular the pole term is retained.

\begin{proposition}[Failure of local open-set unique continuation]
\label{prop:ns55-local-ucp-counterexample}
For every $a>\tfrac12\log2$, there are a nonempty $U\Subset I_a$
and a nonzero real $f\in\mathcal D(A_a)$ with
\[
 f|_U=0,\qquad (A_af)|_U=0.
\]
The second equality is pointwise and distributional on $U$ only.
No complete nullvector or negative Weil direction is asserted.
\end{proposition}
\begin{proof}
Set
\[
 \delta_0=\min\{\ell/8,\tfrac14\log(3/2)\},\qquad
 M_0=2\cosh((\ell+2\delta_0)/2)+\rho(\ell-2\delta_0).
\]
Choose $0<\epsilon<\min\{a-\ell/2,\delta_0,w_2/(4M_0)\}$ and put
\[
 U=(\ell/2-\epsilon,\ell/2+\epsilon),\quad V=U-\ell,\quad
 W=(-\epsilon,\epsilon).
\]
Their closures are disjoint and inside $I_a$. Choose any nonzero real
$h\in C_c^\infty(W)$ and define on $L^2(U)$
\[
 (Tv)(x)=\int_U R(x-t+\ell)v(t)\,dt,\qquad
 H(x)=\int_W R(x-y)h(y)\,dy .
\]
Since $|R(t)|\le M_0$ for $|t-\ell|\le2\delta_0$, Schur's estimate
gives $\|T\|_{2\to2}\le2\epsilon M_0<w_2/2$. Define
$v=(w_2I-T)^{-1}H$ by its convergent Neumann series. The equation
$w_2v=H+Tv$ makes $v$ smooth on a neighborhood of $\overline U$:
the smooth kernels and their derivatives are bounded on that
neighborhood times the fixed integration intervals.

Set $f(y)=v(y+\ell)$ on $V$, $f=h$ on $W$, and zero elsewhere.
This is nonzero. On $U$, positive prime shifts lie beyond the support.
For $n\ge3$, $2\epsilon<\log(3/2)$ puts $x-\log n$ left of $V$.
The only nonzero translated value can therefore be $n=2$, equal to
$v(x)$. Equation~\eqref{eq:ns55-off-support-operator} gives
$A_af=-w_2v+Tv+H=0$ on $U$.

The zero extension is compactly supported and piecewise smooth with
finitely many jumps, hence belongs to $H^s(\mathbb R)$ for
$0<s<1/2$. Its full-line logarithmic multiplier lies in $L^2$:
$\log^2|\xi|\ll_s1+|\xi|^{2s}$ at high frequency, and its transform
is bounded near zero, where $\log^2|\xi|$ is integrable. The form
representation therefore puts $f$ in
$\mathcal D(\tfrac12L_\Delta^{\rm Dir})$. Bounded perturbation
gives $f\in\mathcal D(A_a)$. This is the complete operator, not a head.
\end{proof}

\begin{proposition}[Parity and finitely many linear constraints]
\label{prop:ns55-parity-ucp-counterexample}
For every $a>\log2$ and $\sigma\in\{1,-1\}$, there are a nonempty
$U\Subset(0,a)$ and an infinite-dimensional real linear subspace
$\mathcal L_\sigma\subset\mathcal D(A_a)$ such that
\[
 f(-x)=\sigma f(x),\qquad f=A_af=0\text{ on }U\cup(-U)
 \quad(f\in\mathcal L_\sigma).
\]
One may impose any prescribed finite number of real linear conditions
on $\mathcal L_\sigma$ and still obtain nonzero witnesses. In particular,
source orthogonality does not restore this local unique-continuation
statement. This assertion remains about a local patch, not the full kernel.
\end{proposition}
\begin{proof}
Put $\mathcal P_a=\{\log n:1<n<e^{2a},\ \Lambda(n)>0\}$.
Choose $v_0\in(0,a-\ell)$ avoiding the finitely many equations
$2v_0+\ell\in\mathcal P_a$, and let $u_0=v_0+\ell$.
Choose $w_0\in(0,v_0/2)$ with
$u_0-w_0,u_0+w_0\notin\mathcal P_a$.

Choose an initial radius $\epsilon_0>0$ such that the intervals
$U_{\epsilon_0},V_{\epsilon_0},W_{\epsilon_0}$ about $u_0,v_0,w_0$
and their reflections have disjoint closures inside $I_a$.
Choose it smaller if necessary so differences from $U_{\epsilon_0}$
to $V_{\epsilon_0},-V_{\epsilon_0},W_{\epsilon_0},-W_{\epsilon_0}$
meet $\mathcal P_a$ only in the intended $U-V$ difference $\ell$.
There are only finitely many forbidden central differences, all
separated from $\mathcal P_a$, so this is possible. Set
\[
 M_\sigma=\sup_{x,t\in\overline{U_{\epsilon_0}}}
  |R(x-t+\ell)+\sigma R(x+t-\ell)|<\infty .
\]
Choose $0<\epsilon\le\epsilon_0$ with
$2\epsilon M_\sigma<w_2/2$, and write $U,V,W$ for the smaller intervals.
For arbitrary $h\in C_c^\infty(W)$, define
\begin{align*}
 (T_\sigma v)(x)&=\int_U
       \{R(x-t+\ell)+\sigma R(x+t-\ell)\}v(t)\,dt,\\
 H_\sigma(x)&=\int_W
       \{R(x-y)+\sigma R(x+y)\}h(y)\,dy .
\end{align*}
All arguments are positive. Solve
$(w_2I-T_\sigma)v=H_\sigma$ by the same norm bound.
Set $f(y)=v(y+\ell)$ on $V$, $f=h$ on $W$, reflect with sign $\sigma$,
and put it zero elsewhere. The same smoothness and domain proof applies.
The sole nonzero prime-shift term on $U$ is $-w_2v$, and the integrals
give $T_\sigma v+H_\sigma$, hence $A_af=0$ there. Parity gives the
reflected equality. The linear map $h\mapsto f$ is injective since
$f|_W=h$, so its image is infinite dimensional. Its intersection with
the kernels of finitely many linear functionals is still infinite
dimensional. A complex orthogonality condition on real witnesses is
at most two real linear conditions.
\end{proof}

\paragraph{The logarithmic theorem and its missing hypotheses.}
Chen--Hauer--Weth~\cite[Theorem 1.7]{ChenHauerWeth2023} prove that
$u=L_\Delta u=0$ on the \emph{same} nonempty open set implies $u=0$
on $\mathbb R$, for $u\in L^1((1+|x|)^{-1}dx)$.
See also Harrach--Lin--Weth~\cite[Proposition 4.1]{HarrachLinWeth2025}.
Compactly supported $L^2$ functions satisfy the integrability hypothesis.
But on an interior open set the arithmetic equation gives
$L_\Delta E_af=-2\mathcal K_a f$; $f=0$ does not annihilate this
nonlocal remainder. The preceding counterexamples show the difference
for the actual operator. Replacing it by a local multiplication
potential would change the problem.

Outside the window, only $E_af=0$ is prescribed; neither
$L_\Delta E_af=0$ nor the full arithmetic equation is prescribed.
For example, at $x>a$ one has
\[
 (L_\Delta E_af)(x)=-\int_{-a}^a\frac{f(y)}{x-y}\,dy,
\]
which generally is nonzero. Endpoint trace zero supplies no open
interior vanishing set or missing exterior equation.

\begin{proposition}[Uniqueness with a prime-shift-free exterior-of-support cell]
\label{prop:ns55-prime-free-cell}
Let $f\in V_a$ have essential support $S\subset[-a,a]$, with
$b=\sup S<a$. Suppose a nonempty open interval $J\Subset(b,a)$ obeys
\[
 J\cap\bigcup_{s\in\mathcal P_a}(S+s)=\varnothing,\qquad
 A_af=0\text{ on }J\text{ in distributions}.
\]
Then $f=0$. The reflected left-hand statement also holds.
Here $A_af$ is understood weakly if only $f\in V_a$ is initially given.
\end{proposition}
\begin{proof}
Assume $S$ nonempty. Positive shifts lie beyond $S$; the displayed
separation removes negative shifts. Thus
\[
 \mathcal R_f(x)=\int_{I_a}R(x-y)f(y)\,dy=0\quad(x\in J).
\]
This is real analytic for all $x>b$, so vanishes on that half-line.
Put
$M_-=\int e^{-y/2}f(y)\,dy$ and
$M_j=\int e^{(2j+1/2)y}f(y)\,dy$.
The locally uniformly convergent identity
$\rho(t)=\sum_{j\ge0}e^{-(2j+1/2)t}$ for $t>0$ gives
\[
 0=e^{-x/2}\mathcal R_f(x)
   =M_- -\sum_{j=1}^\infty M_j e^{-(2j+1)x}\quad(x>b).
\]
The $j=0$ term cancels exactly one pole contribution. Since
$|M_j|\le\|f\|_1e^{(2j+1/2)b}$, the last series is a power
series in $z=e^{-x}$ with radius at least $e^{-b}$.
Taking $x\to\infty$ and then its coefficients gives $M_-=0$
and $M_j=0$ for every $j\ge1$.

Push the finite complex measure $e^{y/2}f(y)\,dy$ forward by
$t=e^{2y}$. It is supported in $[e^{-2a},e^{2b}]$, away from zero,
and all its positive integer moments vanish. Polynomials without a
constant term are dense among continuous functions on this interval:
approximate the desired function divided by $t$ and multiply by $t$.
Thus the measure is zero, proving $f=0$.
Only the separated integral component was analytically continued,
not the whole arithmetic equation with its shifts.
\end{proof}

\paragraph{Scope of the remaining gap.}
The last proposition requires actual support separation, not smallness
or zero endpoint trace. A separate first-crossing support argument forces
a first-crossing nullvector to reach both window endpoints; hence it
cannot supply the last proposition's exterior-of-support cell inside
the window. The local counterexamples are not first-crossing nullvectors.
They disprove the direct transplantation of the logarithmic open-set
unique-continuation theorem to the full shifted operator, while the
prime-free-cell theorem is a restricted arithmetic uniqueness result.
Complete affine-potential injectivity, G2 and RH remain open.

% NS-55 first-crossing subtask. Exact identities and scoped countermodel.
% No H^1 assumption, shape derivative, numerical window, or RH assumption.
\subsection*{First-crossing geometry on the complete form domain}

Let $V_a=\mathcal D(q_a)$ be the complete form domain from NS-51,
with unitary Fourier transform of the zero extension understood.
Let $\mu(a)=\inf\sigma(A_a)$. A \emph{first-zero window} means
\begin{equation}\label{eq:ns55fc-first}
 \mu(a_*)=0,\qquad \mu(b)>0\quad(0<b<a_*).
\end{equation}
This is a conditional setup, not an assertion that such a window exists.
If a zero or negative lower edge occurs, continuity and small-window
positivity yield such a first window. The required continuity is
Suzuki's Theorem 1.3, with its complete-domain scaling argument in
Section 4 of version 2
(\url{https://arxiv.org/html/2606.09096v2#S4}).

\begin{proposition}[Support saturation and exact inward variation]
\label{prop:ns55fc-geometry}
Assume \eqref{eq:ns55fc-first}, and let $0\ne f\in\ker A_{a_*}$.
Then the essential support of its zero extension has diameter $2a_*$;
its extreme support points are $-a_*$ and $a_*$. For $0<r<1$ set
\[
 (U_rf)(x)=r^{-1/2}f(x/r),
\]
where $f$ has already been extended by zero. Then $U_rf\in V_{ra_*}$,
$\|U_rf\|_2=\|f\|_2$, and
\begin{equation}\label{eq:ns55fc-null-variation}
 q_{ra_*}[U_rf]
 =q_{a_*}[U_rf-f]
 \ge\mu(ra_*)\|f\|_2^2>0.
\end{equation}
All identities are valid on the complete form domain. Inward dilation
preserves each parity sector.
\end{proposition}
\begin{proof}
The global compact-support form is invariant under simultaneous
translation, and its restrictions are support consistent. If the
essential support diameter were smaller than $2a_*$, a translation
would place $f$ in $V_b$ for some $b<a_*$. Its unchanged form value
zero contradicts $q_b\ge\mu(b)\|\cdot\|_2^2$.

The Fourier identity $\widehat{U_rf}(\xi)=\sqrt r\,\widehat f(r\xi)$
and comparison of $\log(2+|\xi|/r)$ with $\log(2+|\xi|)$ prove
membership in $V_{ra_*}$. Since $f$ is an operator nullvector,
$q_{a_*}(f,v)=0$ for every $v\in V_{a_*}$ by the representation
theorem for closed forms. Expand $q_{a_*}[U_rf-f]$; both cross terms
and $q_{a_*}[f]$ vanish. Support consistency and the variational
principle on $I_{ra_*}$ give \eqref{eq:ns55fc-null-variation}.
No derivative of $r\mapsto U_rf$ is taken.
\end{proof}

The endpoint conclusion means that an argument requiring an empty
exterior support strip \emph{inside} $I_{a_*}$ cannot be applied to
a first-zero nullvector without a new arithmetic step. It does not
mean that $f$ has a nonzero trace: the archived continuous zero
extension and logarithmic endpoint decay are fully compatible with
support saturation.

\begin{proposition}[Exact dilation defect with both poles and all primes]
\label{prop:ns55fc-dilation}
For $f\in V_a$, let
\[
 R_f(t)=\int_{\mathbb R}f(x+t)\overline{f(x)}\,dx,
 \quad P_f(r)=\int_{-a}^a f(x)\cosh(rx/2)\,dx,
 \quad S_f(r)=\int_{-a}^a f(x)\sinh(rx/2)\,dx,
\]
and set $\alpha(\xi)=\operatorname{Re}\psi(1/4+i\xi/2)-\log\pi$.
For $0<r<1$ define
\begin{align}
 \mathcal A_r[f]&=\int_{\mathbb R}
    [\alpha(\xi/r)-\alpha(\xi)]|\widehat f(\xi)|^2\,d\xi,
       \label{eq:ns55fc-arch}\\
 \mathcal P_r[f]&=2r\bigl(|P_f(r)|^2-|S_f(r)|^2\bigr)
                  -2\bigl(|P_f(1)|^2-|S_f(1)|^2\bigr),\notag\\
 \mathcal R_r[f]&=2\sum_{1<n<e^{2a}}\frac{\Lambda(n)}{\sqrt n}
          \operatorname{Re}\{R_f(\log n/r)-R_f(\log n)\}.
          \notag
\end{align}
Then
\begin{equation}\label{eq:ns55fc-defect}
 q_{ra}[U_rf]-q_a[f]
       =\mathcal A_r[f]+\mathcal P_r[f]-\mathcal R_r[f].
\end{equation}
For nonzero $f$,
\begin{equation}\label{eq:ns55fc-arch-bound}
 0<\mathcal A_r[f]\le5\log(1/r)\|f\|_2^2.
\end{equation}
In the even sector $S_f(r)=0$; in the odd sector $P_f(r)=0$.
Neither pole sign may be dropped in the full-space identity.
\end{proposition}
\begin{proof}
The exact complete formula, obtained from
\eqref{eq:v129-jump-decomposition} and
Proposition~\ref{prop:v135-both-parities}, is
\[
 q_a[f]=\int\alpha(\xi)|\widehat f(\xi)|^2\,d\xi
  +2|P_f(1)|^2-2|S_f(1)|^2
  -2\sum_{1<n<e^{2a}}\frac{\Lambda(n)}{\sqrt n}
                        \operatorname{Re}R_f(\log n).
\]
This extends from the smooth core by the form norm: at fixed $a$
the pole and finite prime operators are bounded, and the logarithmic
multiplier defines the principal closed form. Change variables in
the Fourier integral. Directly,
$R_{U_rf}(t)=R_f(t/r)$ and
$P_{U_rf}(1)=\sqrt r P_f(r)$, with the same formula for $S$.
The correlation $R_f$ is continuous, and $R_f(t)=0$ for
$|t|\ge2a$. Thus the dilated prime sum can use the fixed upper
cutoff $e^{2a}$: every added term past $e^{2ra}$ is exactly zero,
including equality. This proves \eqref{eq:ns55fc-defect}, retaining
every changed support overlap rather than differentiating the prime
cutoff or discarding boundary terms.

For $s_k=k+1/4$ and $t=|\xi|/2$, the convergent digamma series gives
\[
 \xi\alpha'(\xi)
 =\sum_{k\ge0}\frac{2t^2s_k}{(s_k^2+t^2)^2}>0\quad(\xi\ne0).
\]
For fixed $t>0$, $F_t(s)=2t^2s/(s^2+t^2)^2$ is nonnegative and
unimodal. Comparing its unit-spaced samples on either side of its
maximum with their adjacent integrals gives
\[
 \sum_{k\ge0}F_t(k+1/4)
 \le\int_{1/4}^\infty F_t(s)\,ds
               +2\sup_{s\ge1/4}F_t(s)\le1+4=5;
\]
the last estimate uses $F_t(s)\le1/(2s)$. Integrate this logarithmic
derivative from $|\xi|$ to $|\xi|/r$. Strict increase away from the
origin and the fact that a nonzero $L^2$ transform cannot be supported
at one point prove \eqref{eq:ns55fc-arch-bound}.
\end{proof}

At a first-zero nullvector, the exact necessary inequality is therefore
\begin{equation}\label{eq:ns55fc-necessary}
 \mathcal R_r[f]
 \le\mathcal A_r[f]+\mathcal P_r[f]
             -\mu(ra_*)\|f\|_2^2
 <\mathcal A_r[f]+\mathcal P_r[f],\qquad 0<r<1.
\end{equation}
There is no contradiction: inward scaling has positive energy.
An independent arithmetic estimate forcing the reverse non-strict
inequality for even one $r<1$ at every proposed first-zero nullvector
would exclude such a nullvector. No such estimate is derived here.
Merely naming that reverse inequality restates a sufficient exclusion
test; it does not make an arithmetic lemma out of the variational identity.

The form domain does not in general make $R_f$ differentiable at the
prime shifts. Its Fourier representation controls a logarithmic moment,
not $\int|\xi||\widehat f(\xi)|^2d\xi$. Consequently dividing
\eqref{eq:ns55fc-defect} by $1-r$ and passing to a differentiated
virial identity would require an additional regularity or limit
argument. The finite identity requires neither. Strong continuity of
dilations in the logarithmic form topology, and hence continuity of
the defect, follows by smooth approximation and locally uniform
dilation bounds in that topology; it supplies no differentiability rate.


% NS-55 prime-shift lane: exact identities and scoped operator obstructions.
% Uses v118 physical realization and v154 affine-potential notation.
\subsection*{Prime-shift primitives and the limits of transfer arguments}

\paragraph{Scope.}
Let $I_a=(-a,a)$, $L=2a$, and use the complete operator $A_a$,
form domain $V_a$ and screw function $\mathfrak g$ of
Proposition~\ref{prop:ns51-affine-potential}. Write
$\ell_n=\log n$ and $w_n=\Lambda(n)/\sqrt n$. Terms with
$w_n=0$ are omitted. The following are exact identities and proved
obstructions to particular inferences. They prove neither an arithmetic
nullvector nor affine-potential injectivity (API). Both pole signs are
retained. No new window or tail metric is computed; G2 and RH remain open.

\begin{proposition}[The exact two-sided primitive equation]
\label{prop:ns55-prime-primitives}
Remove the prime ramps from the screw function, writing on $|t|<L$
\[
 \mathfrak g(t)=\mathfrak g_0(t)+
       \sum_{1<n<e^L}w_n(|t|-\ell_n)_+.
\]
Thus $\mathfrak g_0$ includes the archimedean and both pole terms.
For $f\in L^2(I_a)$ define the clamped primitive and total integral
\[
 H_f(x)=\int_{-\infty}^x E_af(y)\,dy,
 \qquad M_f=\int_{I_a}f(y)\,dy.
\]
Then, at every $x\in I_a$,
\begin{equation}\label{eq:ns55-prime-primitive}
 U_f(x)=(\mathfrak g_0'*E_af)(x)+
 \sum_{1<n<e^L}w_n
       \{H_f(x-\ell_n)+H_f(x+\ell_n)-M_f\}.
\end{equation}
In particular API's equation $U_f=b$ is exactly this two-sided
shifted-primitive equation, subject to $f\in V_a$ and the parity constants
of Proposition~\ref{prop:ns51-affine-potential}. Differentiation is valid
in distributions and gives
\begin{equation}\label{eq:ns55-prime-differentiated}
 (\mathfrak g_0''*E_af)|_{I_a}
 +\sum_{1<n<e^L}w_n\{E_af(\cdot-\ell_n)+E_af(\cdot+\ell_n)\}|_{I_a}=0.
\end{equation}
No $H^1$ regularity of $f$ is required or inferred.
\end{proposition}
\begin{proof}
The derivative of $(|t|-\ell)_+$ is
$\mathbf1_{t>\ell}-\mathbf1_{t<-\ell}$ almost everywhere. Its
convolution with $E_af$ is
$H_f(x-\ell_n)-\{M_f-H_f(x+\ell_n)\}$.
The local $L^2$ property of $\mathfrak g_0'$ and compact support of
$E_af$ make its convolution continuous, as in the affine-potential
criterion. The finite sum is continuous too. Finally
$H_f'=E_af$ in distributions. Membership of this antiderivative in
$H^1_{\rm loc}$ does not imply $f\in H^1$.
\end{proof}

\paragraph{Why this is not a propagation rule.}
The prime part couples both $x-\ell_n$ and $x+\ell_n$.
More importantly, $\mathfrak g_0''*E_af$ remains a nonlocal term.
Knowing $f=0$ on a subinterval does not make that term zero there;
values of $f$ throughout the interval still enter. The logarithmic
singularity also lies on the convolution diagonal, so analyticity
of the kernel away from its singularities does not imply analyticity
of the potential across the support. No identity theorem can be
applied to that unsupported analytic continuation. Additional
prime-free exterior hypotheses can change this conclusion; none
is imposed in the present statement.

\begin{proposition}[The full semigroup is not positivity preserving]
\label{prop:ns55-nonpositive-semigroup}
Let $r>1$ be the unique real solution of $r^3-r-1=0$, and suppose
$L>\log r$. There are real nonnegative
$f,g\in C_c^\infty(I_a)$ with disjoint supports such that
\begin{equation}\label{eq:ns55-positive-cross}
 q_a(f,g)>0.
\end{equation}
Consequently the complete semigroup $e^{-sA_a}$ is not positivity
preserving on the full real $L^2(I_a)$ space with its ordinary
nonnegative cone: for these tests and all sufficiently small $s>0$,
\[
 \langle e^{-sA_a}f,g\rangle<0.
\]
This asserts no negative diagonal Weil energy and makes no separate
claim about order preservation after a parity restriction.
\end{proposition}
\begin{proof}
For disjoint supports the diagonal terms in the operator contribute
nothing. Away from prime-power separations, the off-diagonal integral
kernel of Proposition~\ref{prop:v118-log-dirichlet} is
\[
 J(t)=2\cosh(t/2)-\rho(t),\qquad
 \rho(t)=\frac{e^{t/2}}{e^t-e^{-t}},\quad t>0.
\]
The logarithmic Laplacian contributes $-1/(2t)$ and its bounded
correction contributes $-(\rho(t)-1/(2t))$, giving $-\rho(t)$;
the pole kernel gives the displayed $2\cosh(t/2)$.
For $u=e^t>1$,
\[
 J(t)=\frac{u^3-u-1}{\sqrt u\,(u^2-1)}.
\]
The numerator is strictly increasing on $(1,\infty)$, changes sign
once, and is positive precisely when $t>\log r$.
Choose a separation $t_0\in(\log r,L)$ outside the finite set
$\{\log n:1<n<e^L,\ \Lambda(n)\ne0\}$. Two sufficiently small
disjoint intervals inside $I_a$, at separation $t_0$, have all
cross separations in the positive region of $J$ and avoid every
prime-power separation. Choose nonzero nonnegative smooth bumps
$f,g$ supported in these intervals. All translated prime cross
terms vanish and the remaining double integral is strictly positive.
The tests belong to $\mathcal D(A_a)$, and $A_a$ is self-adjoint
and bounded below at this fixed window. Strong differentiability
of its semigroup on the operator domain gives
\[
 \langle e^{-sA_a}f,g\rangle
 =\langle f,g\rangle-s\langle A_af,g\rangle+o(s)
 =-s q_a(f,g)+o(s),
\]
which is negative for small positive $s$.
\end{proof}

\paragraph{Meaning of the order obstruction.}
A standard positive-kernel or maximum-principle transfer cannot simply
be attributed to the complete arithmetic operator. This does not
exclude a more specialized uniqueness argument, positivity of a
particular eigenfunction, a different cone, or API itself.
Positivity of the quadratic form is distinct from preservation of
pointwise positivity by its semigroup. Moreover conjugation by a
strictly positive multiplier $h$ with $h,h^{-1}\in L^\infty(I_a)$
cannot repair this: multiplication by $h$ and by $h^{-1}$ both
preserve the nonnegative cone, so positivity of the conjugated
semigroup would imply positivity of the original one. This concerns
the full signed operator, not a separate positive edge form.


\paragraph{What remains after these tests.}
The whole-window API input has not been weakened by a proved arithmetic
estimate. The local counterexamples rule out one proposed uniqueness
principle for the actual operator, while the nonpositive semigroup rules
out using its ordinary cone as a maximum-principle mechanism. The
prime-free-cell result needs support separation that a first-zero
nullvector lacks. The finite dilation identity preserves the precise
signed prime overlaps and pole changes; their required estimate remains
unproved. This is a failure of these proposed mechanisms, not a
counterexample to whole-window injectivity or Weil positivity.

The full NS-55 evidence also records the norm discontinuity of an
entering compressed prime shift, its continuous screw-ramp smoothing,
and a valid fixed-window compact Fredholm reduction. The latter leaves
the physical coupling point potentially exceptional; discreteness is
not exclusion. These supporting arguments are archived in
\path{evidence/ns55_kernel/}; no global analyticity, additional endpoint
regularity, new window, tail metric, G2 or RH result is asserted.

\section{Finite prime-weight controls and stable lower floors (v1.56)}
\label{sec:v156-prime-weight-controls}

\paragraph{Classification and scope.} Analytic consequences of the pinned project's complete-domain
and total near-radical results. No claim of literature novelty, new numerical
certificate, arithmetic injectivity, uniform floor, G2, or RH proof.
The perturbed forms below are controls; the original arithmetic operator remains unchanged.

\subsection*{Inputs and conventions}
The standalone checkpoint uses v1.55 as its source baseline.
Write $H=L^2(\mathbb R)$, $I_a=(-a,a)$, and $P_a$ for its ordinary support projection.
The complete Weil form $q_a$ and operator $A_a$ have their original
archimedean part, both pole terms, and all prime-power contributions.
Their compact smooth cores are restrictions of one support-consistent form $q$.
The common fixed-window form domain is
\[
 V_a=\{f\in L^2(I_a):\int_\mathbb R\log(2+|\xi|)
           |\widehat{E_af}(\xi)|^2\,d\xi<\infty\}.
\]
Here $E_a$ is zero extension and the Fourier transform is unitary.
The inner product is linear in the first slot.

We use exactly the following archived inputs, without rerunning their proofs
or numerical certificates:
\begin{itemize}
\item Lemma~\ref{lem:v136-total-radicals}: a dense reflected-translate
span $\mathcal R\subset H$ and even cutoffs give, for every fixed
$h\in\mathcal R$, $g_a\in C_c^\infty(I_a)$ with
$E_ag_a\to h$ and $\norm{A_ag_a}\to0$.
The parity-projected spans are dense in the even and odd subspaces.
For the quantitative corollary we additionally use v1.53,
Lemma~\ref{ns46-a2-lem-fixed-centers}: for a fixed range of translate
centers, the complete cutoff residual is at most
$C\sqrt a\,e^{2a}\exp(-c e^{2a})$, with $C,c>0$ independent of $a$.
\item The complete $A_a$ have compact resolvent, are strictly positive
for sufficiently small $a$, and have continuous lowest eigenvalue.
The standard fixed-window realization is a logarithmic principal form
plus a bounded perturbation.
\item Proposition~\ref{prop:v136-bounded-floor}: a common finite
ordinary lower floor on cofinal windows implies RH. If RH fails,
$\inf\sigma(A_a)\to-\infty$.
\end{itemize}
All analytic statements below inherit these hypotheses. A bounded perturbation
does not inherit the actual near-radical property of $A_a$.

\subsection*{A general bounded-perturbation test}
\begin{lemma}[Detection on the total near-radical span]\label{lem:ns57-bounded-detection}
Let $B$ be any bounded self-adjoint operator on $H$, and define
\[
 q_a^B[f]=q_a[f]+\langle BE_af,E_af\rangle.
\]
If $B$ has a negative direction, then $q^B$ has a negative compact
smooth test. If $B$ preserves parity and has a negative direction
in each parity, so does $q^B$. Conversely, nonnegativity of $q^B$
on all compact smooth tests forces $B\ge0$.
\end{lemma}
\begin{proof}
For each fixed $h\in\mathcal R$ the archived approximants satisfy
\[
 |q_a[g_a]|\le\norm{A_ag_a}\norm{g_a}\longrightarrow0,
 \qquad q_a^B[g_a]\longrightarrow\langle Bh,h\rangle.
\]
If $\langle Bv,v\rangle<0$, ordinary density and boundedness of $B$
give a fixed $h\in\mathcal R$ with $\langle Bh,h\rangle<0$.
For sufficiently large $a$, $g_a$ is the required compact test.
Apply the parity-projected density for either sector. The converse
follows by taking the same limit and using continuity on $H$.
This uses neither global form-norm density nor any uniform bound
on the coefficients in the translate approximation.
When $B$ preserves real functions, taking real parts of approximants
gives real compact tests for real negative directions.
\end{proof}

\begin{theorem}[Bounded-floor stability and exact compensation]\label{thm:ns57-bounded-floor-stability}
For fixed bounded self-adjoint $B$, the existence of a uniform finite
cofinal floor for $q_a^B$ is equivalent to that for $q_a$, hence to RH
under the archived inputs. More sharply, for any real $c$,
\[
 q^B[f]+c\norm f^2\ge0\quad(f\in C_c^\infty(\mathbb R))
 \quad\Longleftrightarrow\quad
 \bigl[\mathrm{RH}\ \hbox{and}\ B+cI\ge0\bigr].
\]
Writing $\beta=\inf\sigma(B)$, the full lowest energies satisfy
\[
 \lim_{a\to\infty}\inf\sigma(A_a+E_a^*BE_a)
 =\begin{cases}\beta,&\text{if RH holds},\\-\infty,&\text{if RH fails}.
 \end{cases}
\]
\end{theorem}
\begin{proof}
The forms differ by at most $\norm B\norm f^2$, so each finite
cofinal floor transfers to the other. The archived bounded-floor
implication and Weil criterion give the equivalence with RH.
Nonnegativity after adding $cI$ therefore implies RH, and the lemma
applied to $B+cI$ forces its positivity. The converse is immediate.

Under RH, $q_a\ge0$ implies the lower bound $\beta$. Given
$\epsilon>0$, approximate a unit vector with $B$-Rayleigh quotient
less than $\beta+\epsilon$ by a nonzero fixed $h\in\mathcal R$,
and then by its cutoff $g_a$. The displayed limits in the lemma,
also after normalization, give a limsup at most $\beta+2\epsilon$.
Let $\epsilon\downarrow0$. If RH fails, the archived divergence
of the original bottom and the bound $\norm B$ give divergence
of the perturbed bottom. No arithmetic sign is established by this dichotomy.
\end{proof}

\subsection*{Changing finitely many actual prime-shift weights}
Let $S$ be a nonempty finite set of distinct prime powers, and let
real changes $\delta_n$ not all be zero. Write
\[
 \ell_n=\log n,\qquad w_n=\Lambda(n)/\sqrt n,\qquad
 (T_tf)(x)=f(x+t),\qquad R_f(t)=\langle T_tf,f\rangle.
\]
Only the coefficient $w_n$ is changed to $w_n+\delta_n$ for $n\in S$.
The new form is exactly
\[
 q_\delta[f]=q[f]-2\sum_{n\in S}\delta_n\operatorname{Re}R_f(\ell_n)
            =q[f]+\langle B_\delta f,f\rangle,
 \qquad B_\delta=-\sum_{n\in S}\delta_n(T_{\ell_n}+T_{-\ell_n}).
\]
This retains both poles, all original prime locations, and the original
archimedean term. Its perturbation is bounded by
$\norm{B_\delta}\le2\sum|\delta_n|$. Its Fourier symbol is
\[
 b_\delta(\xi)=-2\sum_{n\in S}\delta_n\cos(\ell_n\xi).
\]

\begin{lemma}[Both signs without independence of prime logarithms]\label{lem:ns57-symbol-signs}
The continuous even function $b_\delta$ takes both positive and negative
values. Consequently $B_\delta$ has negative directions in both parity sectors.
\end{lemma}
\begin{proof}
Distinct positive frequencies suffice for the elementary Cesaro identities
\[
 \lim_{R\to\infty}\frac1{2R}\int_{-R}^R b_\delta(\xi)\,d\xi=0,
 \qquad
 \lim_{R\to\infty}\frac1{2R}\int_{-R}^R b_\delta(\xi)^2\,d\xi
       =2\sum_{n\in S}\delta_n^2>0.
\]
They follow by integrating each cosine and each product of two cosines.
If $b_\delta\ge0$, then $b_\delta^2\le\norm{b_\delta}_\infty b_\delta$
contradicts these limits. Apply the same argument to $-b_\delta$.
Choose Fourier bumps supported in a negative open interval and its
reflection. A real even Fourier bump produces a real even physical
test; a purely imaginary odd Fourier bump produces a real odd test.
Both have negative $B_\delta$ energy. These are ordinary $L^2$
directions; the preceding lemma supplies the compact smooth Weil tests.
No rational independence, equidistribution rate, or prime number theorem
is required.
\end{proof}

\begin{corollary}[Arbitrarily small finite changes create negative tests]\label{cor:ns57-finite-weight-negative}
For every nonzero finite change $\delta$, $q_\delta$ has negative compact
smooth tests in both parity sectors. This holds for arbitrarily small
$\sum|\delta_n|$, including changes that keep all modified weights
strictly positive. It does not assert a negative test for $q$ itself.
For one changed coefficient, $\inf\sigma(B_\delta)=-2|\delta|$.
Thus, conditionally on RH, the perturbed full lowest energy tends
to $-2|\delta|$.
\end{corollary}

\subsection*{Actual first-zero controls, rather than only matrices}
\begin{theorem}[A finite first-zero window for every nonzero finite change]\label{thm:ns57-first-zero-controls}
Each perturbed family $q_{\delta,a}$ starts strictly positive at small
windows and has a finite first-zero window. The full form is nonnegative
at that first window and has a nonzero operator nullvector. Its support
reaches both endpoints. Each parity sector also has a finite first-zero
window, possibly different from the full first window.
\end{theorem}
\begin{proof}
If $2a\le\min_{n\in S}\ell_n$, every added correlation is zero, so the
original small-window positivity applies. The negative compact tests
above give finite larger windows with negative lowest eigenvalue.

For completeness, continuity survives the finite changes despite the
operator-norm discontinuity of a compressed shift at prime entry.
Scale $I_a$ unitarily to $(-1,1)$. Relative to a fixed reference
archimedean form with compact resolvent, the rescaled form is
$t[u]+\langle K(a)u,u\rangle$ on a common form domain.
The reference can be shifted so $t\ge\norm{\cdot}^2$.
The archimedean multiplier differences are norm-continuous by the
bound $|\xi\alpha'(\xi)|\le5$ proved in Proposition~\ref{prop:ns55fc-dilation}; the pole terms are continuous
finite-rank forms. On a compact range of $a$, there are only finitely
many prime shifts. Each compressed translation by $\ell_n/a$ is
strongly continuous, including where it becomes zero at the endpoint,
and uniformly bounded. Thus $K(a)$ is locally uniformly bounded,
self-adjoint, and strongly continuous.

For $a_j\to a$, a fixed trial vector gives upper semicontinuity of
the ground energy. Normalized ground vectors have uniformly bounded
$t$-energy. Compactness gives an ordinary strongly convergent
subsequence. Uniform boundedness and strong continuity of $K(a_j)$
give convergence of its quadratic values on that subsequence; lower
semicontinuity of $t$ gives the reverse inequality. This proves
continuity, also on the closed parity subspaces.

Continuity between the positive and negative windows gives a first
zero, and compact resolvent makes it an eigenvalue. The perturbed
global form is still translation invariant and support consistent.
If a full first-zero nullvector had support diameter less than $2a_*$,
translate it into a smaller window, contradicting strict positivity
there. Thus both extreme support points equal the endpoints.
For a parity-sector first zero, reflection symmetry of the support
gives the same conclusion without translating out of the sector.
The statement concerns essential support, not a nonzero boundary trace.
\end{proof}

At each fixed window, the elementary variational perturbation estimate is
\[
 |\mu_\delta(a)-\mu_0(a)|\le\norm{B_\delta}
                    \le2\sum_{n\in S}|\delta_n|.
\]
It holds in each parity as well. Thus these controls do not invalidate
finite-window positivity margins; their negative tests may occur later.

\begin{corollary}[A one-sided sensitivity rate for a fixed perturbation direction]\label{cor:ns57-fixed-direction-sensitivity}
Fix a nonzero finite coefficient pattern $\delta^0$ and put
$\delta=\epsilon\delta^0$, $\epsilon>0$. There exist constants
$C_0>0$ and $0<\epsilon_0<1$, depending on that fixed pattern, such
that each parity has a negative test in a window with
\[
 \lambda=e^a\le C_0\sqrt{\log(1/\epsilon)}
       \qquad(0<\epsilon<\epsilon_0).
\]
In particular the first-zero parameter $\lambda_*=e^{a_*}$ has this upper bound.
This is not a matching lower bound or a numerically effective threshold.
\end{corollary}
\begin{proof}
In either parity fix one finite translate combination $h$ with
$\langle B_{\delta^0}h,h\rangle=-\kappa<0$, as above.
All its translate centers lie in a fixed bounded interval. The
v1.53 complete-residual estimate, the bounded norm of $g_a=\chi_a h$,
and absorption of polynomial factors give constants $C,c>0$ with
\[
 |q_a[g_a]|\le C\exp(-c\lambda^2)
\]
for all sufficiently large $a$. Norm convergence and boundedness of
$B_{\delta^0}$ give, after increasing this fixed threshold,
$\langle B_{\delta^0}g_a,g_a\rangle\le-\kappa/2$.
Therefore
\[
 q_{\epsilon\delta^0,a}[g_a]
 \le C\exp(-c\lambda^2)-\epsilon\kappa/2<0
 \quad\hbox{if}\quad
 \lambda^2\ge c^{-1}\log\frac{4C}{\epsilon\kappa}.
\]
Take $\epsilon_0$ small enough to exceed the fixed window threshold
and absorb the additive constant in this logarithm. Choose the larger
constant for the two parity witnesses. No translate coefficients or
centers depend on $\epsilon$; none is asserted uniform over different
perturbation patterns.
\end{proof}

The screw potential is changed by precisely
\[
 \mathfrak g_\delta(t)=\mathfrak g(t)+
                 \sum_{n\in S}\delta_n(|t|-\ell_n)_+.
\]
The proof of Proposition~\ref{prop:ns51-affine-potential} applies to this finite modification:
nullvectors are exactly the complete-domain functions for which
$\mathfrak g_\delta*E_af$ is affine throughout $I_a$.
These controls therefore retain the global affine-potential equation
in its modified form, not just a finite matrix resemblance.
They lose the original arithmetic near-radical identity:
for cutoff translates the perturbed residual generally approaches
$B_\delta h$, rather than zero.

\subsection*{What this changes in the proposed geometry task}
A proof excluding the first arithmetic nullvector must use a property
that fails in these perturbed families, such as the exact arithmetic
radical identities. Translation and reflection geometry, finite-window
compactness, both-endpoint support, and an affine-potential reformulation
alone do not distinguish them. This is a test of those hypotheses,
not a closure of conformal geometry or of API.

A proof of a uniform finite lower floor has a different standard:
bounded coefficient changes preserve that objective. A method robust
under these changes can still prove RH through the original floor
criterion. It would be incorrect to reject such a method merely
because the altered forms are sometimes negative.

The preliminary bounded-multiplier identity supplies no new sign.
For a real symmetric matrix $K$, real vectors $f,m$, and ordinary
quadratic form $q_K$, exact expansion gives
\[
 q_K[mf]-\langle Kf,m^2f\rangle
 =-\frac12\sum_{i,j}K_{ij}f_if_j(m_i-m_j)^2.
\]
At a nullvector the middle term vanishes, but the right-hand side
retains signed coefficients. The analogous form identity requires
the proper domain argument and supplies no automatic positivity
mechanism. We use its finite algebra only as a sign check, not as
an unproved infinite-dimensional theorem.


\paragraph{Evidence and unresolved input.}
The twelve exact Maxima checks in \texttt{evidence/v156/} verify finite
algebra only. The complete-domain implications above rely on the cited
analytic lemmas, not on these checks. No new numerical window, tail metric,
effective threshold, actual arithmetic lower bound, API, G2 or RH proof is
obtained. Integration and local audit details are archived with this version;
no independent review is claimed.

\section{The translated-radical boundary test and its exact limitation (NS-58)}
\label{sec:ns58-boundary}

\paragraph{Classification and outcome.}
The statements below are exact identities and proved implications from
the archived radical and complete-domain inputs. The proposed boundary
test is equivalent to the original nullvector equation; it supplies no
independent sign or uniqueness estimate. This completes the bounded
NS-58 attempt at its stated stop condition. API, the uniform finite
floor, G2 and RH remain open. No numerical window is added.

Use the real even normalized Gaussian arithmetic radical $\phi$ of
Lemma~\ref{lem:v136-total-radicals}, its translates
$\phi_t(x)=\phi(x-t)$, and $I_a=(-a,a)$. The inputs used here are
global distributional radicality $\mathcal W\phi_t=0$, the
double-exponential bounds for every derivative of $\phi$, and the
nonzero entire Fourier transform of $\phi$. These are the archived
v1.35--v1.36 inputs, not a new radical construction. Let $E_a$ denote
zero extension, $P_a$ restriction to $I_a$, $V_a$ the full physical
form domain, and $A_a$ its self-adjoint operator. No source projection
or $H^1$ hypothesis is imposed. Inner products are linear in the first slot.

\subsection*{A sharp cut with every exterior contribution retained}

Write
\[
 u_t=\phi_t|_{I_a},\qquad r_t=\mathbf1_{I_a^c}\phi_t,
 \qquad \ell_n=\log n,\quad w_n=\frac{\Lambda(n)}{\sqrt n},
 \qquad \rho(s)=\frac{e^{-s/2}}{1-e^{-2s}}\quad(s>0).
\]
Define the two exterior pole moments
\[
 C_t=\int_{I_a^c}\cosh(y/2)\phi_t(y)\,dy,\qquad
 S_t=\int_{I_a^c}\sinh(y/2)\phi_t(y)\,dy.
\]
For almost every $x\in I_a$ put
\begin{align}
 R_{a,t}(x)={}&-\int_{I_a^c}\rho(|x-y|)\phi_t(y)\,dy
       +2\cosh(x/2)C_t-2\sinh(x/2)S_t\notag\\
 &-\sum_{n\ge2}w_n\{r_t(x-\ell_n)+r_t(x+\ell_n)\}.
 \label{eq:ns58-exterior-row}
\end{align}
The prime sum is infinite: the exterior rows with $n\ge e^{2a}$
must not be removed. Values assigned at the two cut points do not
affect any of the following $L^2$ identities.

\begin{proposition}[Exact boundary pairing on the complete form domain]
\label{prop:ns58-boundary-identity}
For each fixed $a>0$ and $t\in\mathbb R$,
$u_t\in\mathcal D(A_a)$ and $R_{a,t}\in L^2(I_a)$, and
\begin{equation}\label{eq:ns58-row-action}
 A_a u_t=-R_{a,t}.
\end{equation}
For every $f\in V_a$ the full exterior pairing is therefore
\begin{equation}\label{eq:ns58-boundary-pairing}
 \mathcal B_a f(t):=\langle f,R_{a,t}\rangle
                  =-q_a(f,u_t).
\end{equation}
Each integral and series obtained by inserting
\eqref{eq:ns58-exterior-row} into this pairing is absolutely convergent.
\end{proposition}
\begin{proof}
The zero extension of $u_t$ is piecewise smooth with two possible
jumps. Its Fourier transform is bounded by
$C_{a,t}(1+|\xi|)^{-1}$, using integration by parts on $I_a$.
Consequently it belongs to $H^s(\mathbb R)$ for $0<s<1/2$, and applying the
full archimedean multiplier
$\alpha(\xi)=\operatorname{Re}\psi(1/4+i\xi/2)-\log\pi$
to that zero extension gives an $L^2$ function. The finite interior prime and pole operators are
bounded. The form representation, or
Proposition~\ref{prop:v118-log-dirichlet}, therefore gives
$u_t\in\mathcal D(A_a)$ without requiring a zero boundary trace.

On the open interval, $r_t$ vanishes. Thus the local diagonal terms
in $\mathcal W r_t$ vanish, the archimedean off-diagonal kernel is
$-\rho$, and the pole kernel is $2\cosh((x-y)/2)$. This gives
exactly \eqref{eq:ns58-exterior-row}, including both pole signs.
Near an endpoint, the archimedean integral is bounded in absolute
value by
\[
 C_{a,t}\{1+|\log(a-x)|+|\log(a+x)|\}.
\]
Indeed $\rho(s)=O(s^{-1})$ at zero, $\rho$ decays at infinity,
and $\phi_t$ is bounded and rapidly decreasing. This bound is in
$L^2(I_a)$. Also $|r_t(y)|\le C_t' e^{-2|y|}$, so uniformly on
$I_a$ the prime sum is bounded by
\[
 2C_t'e^{2a}\sum_{n\ge2}\frac{\Lambda(n)}{n^{5/2}}<\infty.
\]
Both pole integrals converge absolutely. These estimates give the
$L^2$ claim and, by Cauchy--Schwarz with $f\in L^2(I_a)$,
absolute convergence of all paired terms, including the double integral.

The global identity $\mathcal W\phi_t=0$ and the decomposition
$\phi_t=E_a u_t+r_t$ give \eqref{eq:ns58-row-action} in distributions
on $I_a$. Both sides are in $L^2(I_a)$, hence are equal there.
The representation theorem for the closed form gives
$q_a(f,u_t)=\langle f,A_a u_t\rangle$, proving the pairing formula.
No globally closed Weil operator or ordinary global form norm is used.
\end{proof}

\subsection*{Testing every translate recovers exactly the kernel equation}

\begin{proposition}[No additional exclusion from the vanishing pairing alone]
\label{prop:ns58-equivalence}
For $f\in V_a$,
\begin{equation}\label{eq:ns58-equivalence}
 f\in\ker A_a
 \quad\Longleftrightarrow\quad
 \mathcal B_a f(t)=0\quad\hbox{for every }t\in\mathbb R.
\end{equation}
For $\sigma\in\{1,-1\}$ and $f(-x)=\sigma f(x)$ the same
equivalence holds using only
\[
 u_t^\sigma=\tfrac12(u_t+\sigma u_{-t}),\qquad
 R_{a,t}^\sigma=\tfrac12(R_{a,t}+\sigma R_{a,-t}).
\]
For these parity rows the exterior odd pole moment vanishes when
$\sigma=1$, and the exterior even pole moment vanishes when
$\sigma=-1$. The remaining odd pole has its negative sign.
\end{proposition}
\begin{proof}
The forward implication follows from
\eqref{eq:ns58-boundary-pairing}. For the converse, define a linear
distribution on smooth test functions on the line by
\[
 D_f(\psi)=q_a(f,\overline{\psi}|_{I_a}).
\]
It is supported in $[-a,a]$ and has finite order. To check this rather
than assume density in an unspecified global form norm, integration
by parts gives
\[
 \|\psi|_{I_a}\|_{V_a}
 \le C_a\bigl(\|\psi\|_{L^\infty(I_a)}
                +\|\psi'\|_{L^\infty(I_a)}\bigr),
\]
where the norm on $V_a$ is the ordinary norm plus the logarithmically
weighted Fourier norm of its zero extension. The endpoint jumps are
included in the Fourier bound. Form continuity then bounds $D_f$
by the same seminorm, multiplied by a constant depending on $f$.
In particular $D_f$ is a compactly supported distribution of order
at most one, acting also on noncompact smooth functions by a cutoff.

The assumed vanishing and the reality and evenness of $\phi$ give
\[
 (D_f*\phi)(t)=D_f(\phi_t)=q_a(f,u_t)=0\quad(t\in\mathbb R).
\]
Fourier transformation of this convolution is legitimate for a
compactly supported distribution and a Schwartz function. It gives
$\widehat D_f\,\widehat\phi=0$. The first factor is a smooth
function on the real line; the second is a nonzero entire function
with isolated real zeros. Hence $\widehat D_f=0$ first off those
zeros and then everywhere by continuity. Fourier uniqueness implies
$D_f=0$. In particular $q_a(f,v)=0$ for every
$v\in C_c^\infty(I_a)$. This is a form core, so form continuity
and the representation theorem give $f\in\mathcal D(A_a)$ and
$A_af=0$.

Reflection sends $u_t$ to $u_{-t}$ and commutes with $A_a$.
For a vector of parity $\sigma$, pairing against $u_t^\sigma$
is the same as pairing against $u_t$. The preceding proof therefore
applies in either sector. The assertions about pole moments follow
by integrating odd functions over the symmetric exterior set.
\end{proof}

This result does not assert that the exterior rows form a total family
in $L^2(I_a)$. Proving the needed separation of all nonzero $V_a$
vectors would itself prove the still-open kernel exclusion. Ordinary
translation density of $\phi$ cannot be transferred through the
unbounded operator $A_a$ while silently assuming that exclusion.
Nor has an endpoint derivative of $f$ been introduced: the formula
contains actual exterior interactions, not an assumed local boundary flux.

\subsection*{The altered-weight controls acquire an explicit forcing term}

Use NS-57's real finite coefficient change $\delta\ne0$ and the
bounded self-adjoint whole-line operator
\[
 B_\delta=-\sum_n\delta_n(T_{\ell_n}+T_{-\ell_n}),\qquad
 A_a^\delta=A_a+P_a B_\delta E_a,
\]
where $(T_s h)(x)=h(x-s)$. Put
\[
 R_{a,t}^\delta=R_{a,t}+P_aB_\delta r_t,
 \qquad d_{a,t}^\delta=P_aB_\delta\phi_t.
\]
Thus the added terms in the first expression are
$-\sum_n\delta_n\{r_t(x-\ell_n)+r_t(x+\ell_n)\}$, and the
same expression with $\phi_t$ replaces $r_t$ in the forcing term.

\begin{proposition}[Forced pairing and a necessarily visible defect]
\label{prop:ns58-forced-pairing}
For every $f\in V_a$,
\begin{equation}\label{eq:ns58-forced-pairing}
 \langle f,R_{a,t}^\delta\rangle-
 \langle f,d_{a,t}^\delta\rangle=-q_a^\delta(f,u_t).
\end{equation}
Vanishing of the left side for all $t$ is equivalent to
$f\in\ker A_a^\delta$, also in either parity sector.
For every nonzero $f\in L^2(I_a)$ the forcing profile
\[
 F_\delta f(t)=\langle f,d_{a,t}^\delta\rangle
\]
is not identically zero. In particular every nonzero altered-weight
nullvector has a nonzero forcing profile. Nevertheless this profile
has no uniform positive $L^2$ lower bound over the whole unit sphere
of either fixed-window parity space.
\end{proposition}
\begin{proof}
The exact decomposition $E_a u_t+r_t=\phi_t$ gives
\[
 A_a^\delta u_t=-R_{a,t}^\delta+d_{a,t}^\delta,
\]
which proves \eqref{eq:ns58-forced-pairing}. The equivalence proof
above applies with $q_a^\delta$ in place of $q_a$: the added operator
is bounded and preserves the form domain and its continuity estimates.

Self-adjointness gives
$F_\delta f(t)=\langle B_\delta E_a f,\phi_t\rangle$.
If this vanishes for every $t$, the ordinary translation-density
lemma implies $B_\delta E_a f=0$. Its multiplier is
\[
 b_\delta(\xi)=-2\sum_n\delta_n\cos(\xi\ell_n).
\]
By Lemma~\ref{lem:ns57-symbol-signs} this is not identically zero.
It is real analytic, so its real zero set has measure zero; its
multiplication operator is injective on $L^2(\mathbb R)$.
Thus $f=0$, a contradiction. The same argument uses reflected
translates in either parity sector, since $B_\delta$ commutes
with reflection.

For clarity this injectivity supplies no uniform inverse estimate.
With the unitary Fourier convention, Plancherel gives exactly
\[
 \|F_\delta f\|_{L^2(dt)}^2
   =2\pi\int_{\mathbb R}|b_\delta(\xi)|^2
                |\widehat{E_a f}(\xi)|^2|\widehat\phi(\xi)|^2\,d\xi.
\]
Its integral kernel is $B_\delta\phi(x-t)$ for $x\in I_a$,
so its squared Hilbert--Schmidt norm is
$2a\|B_\delta\phi\|_2^2<\infty$. It is compact, also on
each infinite-dimensional parity space. An orthonormal sequence
there converges weakly to zero, and its image under this compact
map converges to zero in norm. This precludes a positive lower
bound on the whole unit sphere. No lower bound restricted to the
finite-dimensional kernel of a particular altered operator is excluded.
\end{proof}

\paragraph{Disposition and remaining input.}
The original family has a homogeneous boundary pairing; the altered
families have a forced one. The forcing is necessarily visible, so
the original arithmetic radical identity really does distinguish the
controls. But homogeneity for the original family is already equivalent
to its null equation. Neither that distinction nor ordinary density
establishes the missing sign. The first-zero assumptions add the
previous support saturation and variational inequalities, but no
independent bound on the signed expression above has been obtained.
NS-58 therefore stops as specified, with the boundary test fully
derived and its logical strength identified. The whole API and
bounded-floor routes remain open. Finite exact algebra checks in
\texttt{evidence/v157/} check signs and the forcing bookkeeping only;
the complete-domain arguments are the analytic proofs above.


\section{Nyman--Beurling difference budgets and their arithmetic limits (NS-59)}
\label{sec:v158-nb-difference}

\paragraph{Status.}
This section gives exact identities, proved implications, and separately
labelled finite certificates. It supplies an unconditional asymptotic upper
bound for a complete block Gram matrix. It does not supply the lower
correlation estimate needed for Nyman--Beurling convergence. An almost
scale-invariant bound for the \emph{unprojected} difference Gram is equivalent
to Lindel\"of; every fixed logarithmic version of that bound is false.
Neither assertion is transferred to the projected Gram or the particular
residual direction without an additional argument. No priority claim is made
for these consequences of the cited classical zeta estimates.

Use the Hilbert space, $\rho_n$, $P_N$, $r_N$, $E_N$, $h$ and $S$ from
Section~\ref{sec:ns53-nb-blocks}. In this section the new block is
$N<n\leq 2N$, and its indices always retain their integer values.
For $n\geq2$ put
\begin{equation}
 \delta_n=\rho_n-\frac{n-1}{n}\rho_{n-1},\qquad
 H_k=\sum_{j=1}^k\frac1j,\qquad
 L_N=\sum_{n=N+1}^{2N}\frac{H_{n-1}}{n^2}.
 \label{eq:v158-difference-def}
\end{equation}

\begin{proposition}[An elementary adjacent-dilation budget]
\label{prop:v158-adjacent-budget}
For every integer $n\geq2$,
\begin{equation}
 \frac{H_{n-1}}{n^2}-\frac{n-1}{n^3}
 \leq \norm{\delta_n}_H^2
 \leq\frac{H_{n-1}}{n^2}.
 \label{eq:v158-adjacent-norm}
\end{equation}
In particular $\norm{\delta_n}_H^2=(\log n+O(1))/n^2$ and
$L_N\leq H_{2N-1}/(2N)$.
\end{proposition}
\begin{proof}
Write $a=n-1$ and $t=1/x$. The function
\[
 \delta_n(1/t)=\frac{a}{n}\left\lfloor\frac{t}{a}\right\rfloor
              -\left\lfloor\frac{t}{n}\right\rfloor
\]
has period $an$. On $[ka,kn)$, $1\leq k\leq a$, its value is
$1-k/n$; on $[kn,(k+1)a)$, $1\leq k\leq a-1$, its value is $-k/n$.
It vanishes on $[0,a)$. Thus, including the zero-length last interval,
\begin{align*}
 \int_0^{an}\frac{|\delta_n(1/t)|^2}{t^2}\,dt
 &=\sum_{k=1}^a\left\{
 \frac{(n-k)^2}{k a n^3}
 +\frac{k(a-k)}{(k+1)a n^3}\right\}\\
 &=\frac{H_a}{n^2}-\frac{a}{n^3}.
\end{align*}
Periodicity and the displayed values give $|\delta_n(1/t)|\leq a/n$ everywhere.
The remaining nonnegative integral is at most
$(a/n)^2\int_{an}^{\infty}t^{-2}\,dt=a/n^3$.
Finally $\sum_{N<n\leq2N}n^{-2}\leq\int_N^{2N}u^{-2}\,du=1/(2N)$.
\end{proof}

Let $T_N$ be the upper bidiagonal matrix whose column indexed by $n$
has diagonal entry $1$ and, when $n>N+1$, preceding entry $-(n-1)/n$.
The first column represents $(I-P_N)\delta_{N+1}=(I-P_N)\rho_{N+1}$;
the old $\rho_N$ term is not discarded before projection. Define
\begin{equation}
 \mathcal D_N=(\langle \delta_m,\delta_n\rangle)_{N<m,n\leq2N},\quad
 K_N=T_N^*S T_N,\quad g=T_N^*h.
 \label{eq:v158-difference-gram}
\end{equation}
Here $\mathcal D_N$ is the raw difference Gram and $K_N$ is its
orthogonally projected Gram. In coordinates,
$g_n=h_n-(n-1)h_{n-1}/n$, with $h_N=0$.

\begin{proposition}[Complete difference-block gains]
\label{prop:v158-difference-gain}
The following identities and inequalities retain all old/new cross terms:
\begin{align}
 E_N-E_{2N}&=g^*K_N^{-1}g, & 0<K_N&\leq\mathcal D_N,
 \label{eq:v158-exact-difference-gain}\\
 E_N-E_{2N}&\geq\frac{\norm{g}_2^2}{L_N},&
 E_N-E_{2N}&\geq
 \frac{\norm{g}_2^4}{g^*K_Ng}\quad(g\ne0).
 \label{eq:v158-difference-bounds}
\end{align}
Replacing the old denominator by $L_N$ while keeping the old numerator
$\norm{h}_2^2$ is not justified by these formulas.
\end{proposition}
\begin{proof}
The invertible matrix $T_N$ changes the projected spanning family, so the
first identity follows from Proposition~\ref{prop:ns53-nb-block-gain}.
Orthogonal projection is contractive, hence $K_N\leq\mathcal D_N$.
Proposition~\ref{prop:v158-adjacent-budget} bounds
$\operatorname{tr}\mathcal D_N$ by $L_N$. The first lower bound follows
from $K_N\leq L_N I$, and the second is the variational gain formula
applied to the coefficient direction $g$. If $g=0$, the gain is zero.
\end{proof}

\begin{proposition}[Mellin multiplier and an unconditional block norm]
\label{prop:v158-mellin-budget}
For complex coefficients $a=(a_n)_{N<n\leq2N}$ set
\[
 F_a(u)=\sum_{n=N+1}^{2N}\frac{a_n}{n}\mathbf1_{(n-1,n]}(u),
 \qquad q_a(v)=e^{v/2}F_a(e^v).
\]
With $\widehat q(t)=\int_{\mathbb R}q(v)e^{-itv}\,dv$,
\begin{equation}
 a^*\mathcal D_Na
 =\frac1{2\pi}\int_{\mathbb R}
       |\zeta(\tfrac12+it)|^2|\widehat q_a(t)|^2\,dt,
 \qquad
 \norm{q_a}_2^2=\sum_{N<n\leq2N}\frac{|a_n|^2}{n^2}.
 \label{eq:v158-mellin-identity}
\end{equation}
If $0<\theta<1/2$ and
$|\zeta(1/2+it)|\leq A_\theta(1+|t|)^\theta$ for all real $t$, then
\begin{equation}
 \norm{\mathcal D_N}_{\rm op}
 \leq C_\theta A_\theta^2 N^{-2+2\theta},
 \qquad \norm{K_N}_{\rm op}\leq\norm{\mathcal D_N}_{\rm op}.
 \label{eq:v158-zeta-budget}
\end{equation}
Consequently, for every $\varepsilon>0$,
\begin{equation}
 \norm{\mathcal D_N}_{\rm op}\ll_\varepsilon N^{-5/3+\varepsilon}.
 \label{eq:v158-weyl-budget}
\end{equation}
This is an analytic bound with an unspecified constant, not a numerical
certificate for a finite cutoff or a claim of the best available exponent.
\end{proposition}
\begin{proof}
For $0<\Re s<1$ the usual fractional-part Mellin identity gives
$\mathcal M\rho_n(s)=-n^{-s}\zeta(s)/s$. Also
\[
 n^{-s}-\frac{n-1}{n}(n-1)^{-s}
 =\frac{1-s}{n}\int_{n-1}^n u^{-s}\,du.
\]
On $s=1/2+it$ the factor $(1-s)/s$ has modulus one. Mellin Plancherel
therefore gives~\eqref{eq:v158-mellin-identity}, including the entire
frequency integral. The norm identity follows by changing variables
$u=e^v$.

Here is the needed inverse estimate, including the discontinuities of
$q_a$. Its support is $[\log N,\log(2N)]$, partitioned into intervals
$I_n=(\log(n-1),\log n]$ of lengths between $1/(2N)$ and $1/N$.
On each interval $q_a(v)$ is a constant times $e^{v/2}$. For
$0<\theta<1/2$,
\begin{equation}
 \int_{\mathbb R}(1+|t|)^{2\theta}|\widehat q_a(t)|^2\,dt
 \leq C_\theta N^{2\theta}\norm{q_a}_2^2.
 \label{eq:v158-inverse-sobolev}
\end{equation}
To check this, use the Fourier identity between the homogeneous
$H^\theta$ seminorm and
$\iint |q(v)-q(w)|^2|v-w|^{-1-2\theta}\,dv\,dw$.
On a single cell, $q'=q/2$ bounds the integral by a constant times its
$L^2$ mass. For different cells, including the zero exterior, use
$|q(v)-q(w)|^2\leq2(|q(v)|^2+|q(w)|^2)$.
For $v\in I=(b,c)$, integration outside $I$ gives
$((v-b)^{-2\theta}+(c-v)^{-2\theta})/(2\theta)$.
Since $|q|^2$ varies by at most $e^{|I|}\leq2$ on a cell, its integral
against these endpoint powers is at most
$C_\theta |I|^{-2\theta}\int_I|q|^2$.
The powers are integrable exactly for $2\theta<1$. Summing gives
\eqref{eq:v158-inverse-sobolev}; thus no $H^1$ regularity is assumed.
Combining with~\eqref{eq:v158-mellin-identity} proves the operator bound.
The classical Weyl estimate
$\zeta(1/2+it)\ll_\delta(1+|t|)^{1/6+\delta}$
from~\cite[Section 2]{HuxleyIvic2007} gives
\eqref{eq:v158-weyl-budget} by taking $\theta=1/6+\varepsilon/2$
for small $\varepsilon$; larger $\varepsilon$ follow at once.
\end{proof}

\begin{proposition}[A sharp raw norm would already encode Lindel\"of]
\label{prop:v158-lindelof-budget}
The assertion
\begin{equation}
 \text{for every }\varepsilon>0,\qquad
 \norm{\mathcal D_N}_{\rm op}\ll_\varepsilon N^{-2+\varepsilon}
 \quad(N=2^j)
 \label{eq:v158-near-scale-budget}
\end{equation}
is equivalent to the Lindel\"of hypothesis. Moreover, for every fixed
$B\geq0$,
\begin{equation}
 \norm{\mathcal D_N}_{\rm op}
 \ne O\left(\frac{(\log N)^B}{N^2}\right)
 \quad\text{even when }N\text{ is restricted to powers of two}.
 \label{eq:v158-no-log-budget}
\end{equation}
Both statements concern the raw Gram $\mathcal D_N$ uniformly over
\emph{all} coefficient vectors. They assert neither such an equivalence
nor such an obstruction for $K_N$ or for $g^*K_Ng$.
\end{proposition}
\begin{proof}
Lindel\"of and Proposition~\ref{prop:v158-mellin-budget} imply
\eqref{eq:v158-near-scale-budget}. For the converse we first show that,
with an absolute constant and all sufficiently large $N$,
\begin{equation}
 \int_{\tau-1}^{\tau+1}|\zeta(\tfrac12+it)|^2\,dt
 \leq C N^2\norm{\mathcal D_N}_{\rm op}
 \qquad(|\tau|\leq\sqrt N).
 \label{eq:v158-local-zeta-test}
\end{equation}
Fix a nonzero nonnegative smooth function $w$ compactly supported in
$(0,\log2)$. Its Fourier transform satisfies
$|\widehat w(t)|\geq\cos(\log2)\int w>0$ for $|t|\leq1$.
Put $v_n=\log(n-1/2)$ and choose
\[
 a_n=n e^{-v_n/2}w(v_n-\log N)
                      e^{i\tau(v_n-\log N)}.
\]
The associated $q_a$ approximates
$q(v)=w(v-\log N)e^{i\tau(v-\log N)}$ with
\[
 \norm{q_a-q}_{L^1}\leq C_w\frac{1+|\tau|}{N},
 \qquad \norm{a}_2^2\leq C_w N^2.
\]
These estimates follow cell by cell from the mean value theorem and
$|v-v_n|\leq C/N$; the constants are independent of $N$ and $\tau$.
For $|\tau|\leq\sqrt N$, the Fourier error tends uniformly to zero.
Hence $|\widehat q_a(t)|$ is bounded below by a fixed positive constant
on $[\tau-1,\tau+1]$ for all sufficiently large $N$.
Equation~\eqref{eq:v158-mellin-identity} proves
\eqref{eq:v158-local-zeta-test}. Complex coefficients are harmless:
$\mathcal D_N$ is real symmetric, and its operator norm is the same
over real or complex vectors.

The classical short-interval inequality
\cite[Eq. (2.1)]{HuxleyIvic2007} is
\begin{equation}
 |\zeta(\tfrac12+iT)|^2\ll\log T
 \left(1+\int_{T-\log^2T}^{T+\log^2T}
                      |\zeta(\tfrac12+it)|^2\,dt\right)
 \quad(T\geq T_0).
 \label{eq:v158-short-interval}
\end{equation}
Choose a power of two $N$ with $4T^2\leq N<8T^2$.
For large $T$ the integration interval can be covered by
$O(\log^2T)$ intervals in~\eqref{eq:v158-local-zeta-test}.
Consequently
\begin{equation}
 |\zeta(\tfrac12+iT)|^2
 \ll\log T+(\log T)^3 N^2\norm{\mathcal D_N}_{\rm op}.
 \label{eq:v158-budget-to-zeta}
\end{equation}
Assertion~\eqref{eq:v158-near-scale-budget}, with arbitrarily small
exponents, now implies Lindel\"of.
If the big-$O$ estimate excluded in~\eqref{eq:v158-no-log-budget}
held, the same formula would give
$|\zeta(1/2+iT)|^2\ll(\log T)^{B+3}$.
This contradicts the unconditional large values in
\cite[Theorem 1]{Soundararajan2008}.
\end{proof}

\paragraph{Certified finite comparison; no extrapolation.}
The exact rational-cotangent formula for the Gram entries is taken from
\cite[Proposition 89]{BaezDuarteAutocorrelation2003}. With
$V(p,q)=\sum_{k=1}^{q-1}\{kp/q\}\cot(\pi k/q)$, $d=(m,n)$ and
$\kappa=\log(2\pi)-\gamma$, its scaling to the current conventions is
\begin{equation}
 G_{mn}=\frac\kappa2\left(\frac1m+\frac1n\right)
 +\frac{n-m}{2mn}\log\frac mn
 -\frac{\pi d}{2mn}\{V(m/d,n/d)+V(n/d,m/d)\}.
 \label{eq:v158-cot-gram}
\end{equation}
All entries, cross terms and solves in \texttt{evidence/v158/nb\_certify.py}
are evaluated as Arb balls, independently at 256 and 384 bits, for the
same matrices through index 512. The program also checks
\eqref{eq:v158-adjacent-norm} for $2\leq n\leq512$ as a finite replay
of the analytic estimate. At the displayed blocks, the following numbers
are outward-rounded enclosing intervals; all ratios are relative to $E_N$:
\begin{center}
\begin{tabular}{rcccc}
\toprule
$N$ & actual gain & raw-trace bound & difference budget & difference direction\\
\midrule
128 & $(.14752,.14753)$ & $(.0003503,.0003505)$
    & $(.0004286,.0004289)$ & $(.08380,.08381)$\\
256 & $(.10289,.10290)$ & $(.0002135,.0002137)$
    & $(.0001328,.0001331)$ & $(.05955,.05956)$\\
\bottomrule
\end{tabular}
\end{center}
The raw-trace bound is less than $1/400$ of the actual gain at both
blocks. At $N=256$ the smaller adjacent-difference denominator produces
an even smaller trace bound after transforming the numerator correctly.
The full directional quotient performs better on these two finite blocks;
this is not an asymptotic lower bound. The floating scan through index
2048 is separately labelled a numerical illustration in
\texttt{evidence/diag\_ns59\_nb\_budget/}.

\paragraph{What remains missing.}
The unconditional bound~\eqref{eq:v158-weyl-budget} supplies a denominator,
so it gives $E_N-E_{2N}\geq C_\varepsilon^{-1}
N^{5/3-\varepsilon}\norm{g}_2^2$ for a suitable constant.
To use it for convergence one still needs, for example, a cofinal lower
bound $\norm{g}_2^2\geq C_\varepsilon N^{-5/3+\varepsilon}
\alpha_j E_N$ at $N=2^j$, with $0\leq\alpha_j<1$ and
$\sum_j\alpha_j=\infty$. No such estimate is proved here.
The nonsummability of the \emph{exact} relative gain is simply the known
convergence criterion, not an independent arithmetic input.
Uniform logarithmic control of $\mathcal D_N$ is ruled out, while control
restricted to the actual residual direction remains open. The raw RBC
input of Section~\ref{sec:ns53-nb-blocks} is neither disproved by the
finite losses nor established by the new denominator estimates.
RH, G2 and the uniform Weil floor remain open.


\section{Explicit old-space Euler corrections and the grid cost (NS-60)}
\label{sec:ns60-euler-corrections}

\paragraph{Status and scope.}
The preceding raw-norm obstruction does not apply automatically after
optimal projection. Here we test one explicit family of old-space
corrections. We derive its exact grid error and prove that its full
coefficient norm still cannot obey a fixed-logarithmic $N^{-2}$ bound,
even with suitably growing Euler products. This closes only that uniform
comparison for this correction family. It does not close the optimal
projected comparison, tapered non-Euler mollifiers, or a bound for the
actual residual direction. All results below are analytic; the separate
finite scan is a numerical illustration, not proof evidence.

Extend the difference notation by $\delta_1=\rho_1$. For a finitely
supported integer step function $F$, with value $F_n$ on $(n-1,n]$, set
\begin{equation}
 \Psi F=\sum_{n\geq1}nF_n\delta_n.
 \label{eq:ns60-step-map}
\end{equation}
The triangular change of basis gives
$\operatorname{span}\{\delta_1,\ldots,\delta_N\}=V_N$.
The Mellin identity in Proposition~\ref{prop:v158-mellin-budget} extends
to the first cell because $\int_0^1u^{-s}\,du=(1-s)^{-1}$.

For an integer $d\geq2$, define the old-grid averaging operator
\begin{equation}
 (J_dF)(u)=\int_{m-1}^{m}F(dv)\,dv
 =\frac1d\sum_{d(m-1)<n\leq dm}F_n
 \quad(m-1<u\leq m).
 \label{eq:ns60-cell-average}
\end{equation}
If $F$ is supported in $(N,2N]$, then $J_dF$ is supported in $(0,N]$.
For a squarefree positive integer $Q$ put
\begin{equation}
 \Phi_QF=F+\sum_{\substack{d\mid Q\\d>1}}\mu(d)J_dF,
 \qquad
 \mathcal C_{N,Q}=\text{Gram matrix of }a\longmapsto\Psi\Phi_QF_a,
 \label{eq:ns60-euler-correction}
\end{equation}
where $F_a$ is the new-block step function in
Proposition~\ref{prop:v158-mellin-budget}. When $Q=1$, this is the raw
difference Gram $\mathcal D_N$.

\begin{proposition}[A valid correction, with the grid error retained]
\label{prop:ns60-grid-correction}
The correction lies in $V_N$ and has the explicit old difference
coefficients
\begin{equation}
 (B_{N,Q}a)_m
 =m\sum_{\substack{d\mid Q\\d>1}}\frac{\mu(d)}d
       \sum_{\substack{N<n\leq2N\\d(m-1)<n\leq dm}}\frac{a_n}{n}.
 \label{eq:ns60-correction-matrix}
\end{equation}
Therefore
\begin{equation}
 0<K_N\leq\mathcal C_{N,Q},\qquad
 \mathcal C_{N,Q}
 =\mathcal D_N+X^*B_{N,Q}+B_{N,Q}^*X+B_{N,Q}^*G_\delta B_{N,Q},
 \label{eq:ns60-full-correction-gram}
\end{equation}
where $G_\delta$ is the old difference Gram and
$X=(\langle\delta_m,\delta_n\rangle)_{m\leq N<n\leq2N}$.
Every old/new cross term is present.
Equivalently,
\[
 \mathcal C_{N,Q}-K_N
 =(B_{N,Q}+G_\delta^{-1}X)^*G_\delta
                         (B_{N,Q}+G_\delta^{-1}X).
\]

If $F$ is constant on each interval $(jQ,(j+1)Q]$ on the positive
axis, then $J_dF=F(d\,\cdot)$ for $d\mid Q$. On this compatible
subspace,
\begin{equation}
 \norm{\Psi\Phi_QF}_H^2
 =\frac1{2\pi}\int_{\mathbb R}|\zeta(\tfrac12+it)|^2
       \prod_{p\mid Q}|1-p^{-1/2-it}|^2
       |\widehat q_F(t)|^2\,dt,
 \qquad q_F(v)=e^{v/2}F(e^v).
 \label{eq:ns60-compatible-euler}
\end{equation}
For general $F$, the missing term in this simplification is exactly
\begin{equation}
 R_QF=\sum_{\substack{d\mid Q\\d>1}}\mu(d)
                   \{J_dF-F(d\,\cdot)\}.
 \label{eq:ns60-grid-defect}
\end{equation}
Its Mellin transform and its cross term cannot be omitted.
\end{proposition}
\begin{proof}
Equation~\eqref{eq:ns60-cell-average} gives
\eqref{eq:ns60-correction-matrix}, and the old support gives membership
in $V_N$. Applying $I-P_N$ to $\Psi\Phi_QF_a$ therefore gives the
same projected new vector as applying it to $\Psi F_a$.
Contractivity proves $K_N\leq\mathcal C_{N,Q}$; expanding the squared
norm gives the complete Gram formula.

For compatible $F$, every breakpoint of $F(d\,\cdot)$ is an integer
because $d\mid Q$. Averaging leaves this function unchanged. Also
\[
 \int_0^\infty F(du)u^{-s}\,du
       =d^{s-1}\int_0^\infty F(u)u^{-s}\,du,
 \qquad
 \sum_{d\mid Q}\mu(d)d^{s-1}=\prod_{p\mid Q}(1-p^{s-1}).
\]
On $\Re s=1/2$, the last product has the modulus of
$\prod_{p\mid Q}(1-p^{-s})$. Mellin Plancherel gives
\eqref{eq:ns60-compatible-euler}. This is a finite product identity;
no convergence of an Euler product in the critical strip is assumed.
For general $F$ the transformed
quantity is instead
\[
 \prod_{p\mid Q}(1-p^{s-1})\widehat q_F(t)
                           +\widehat q_{R_QF}(t).
\]
Its absolute square contains both the defect energy and the mixed term.
All these finite step functions have finite weighted Mellin energy,
by the fractional Sobolev argument already proved, so the displayed
identities are Hilbert-space identities rather than formal expansions.
\end{proof}

The grid defect is already nonzero when $N=2$, $Q=2$ and
$F=\mathbf1_{(2,3]}$: then
$F(2u)=\mathbf1_{(1,3/2]}(u)$ whereas
$J_2F=\tfrac12\mathbf1_{(1,2]}$.
Their ordinary $L^2(du)$ distance squared is $1/4$.
This elementary example concerns the grid functions, not a negative
Weil direction or a convergence obstruction.

\begin{lemma}[The smallest Euler multiplier for a bounded grid product]
\label{lem:ns60-euler-minimum}
Let $m_Q=\prod_{p\mid Q}(1-p^{-1/2})$, with $m_1=1$.
For $L=\max(3,\log Q)$,
\begin{equation}
 0\leq-\log m_Q\leq C\frac{\sqrt L}{\log L}
 \label{eq:ns60-euler-minimum}
\end{equation}
with an absolute constant. Consequently, for $Q\leq N$ and sufficiently
large $N$,
$m_Q^{-2}\leq\exp(C\sqrt{\log N}/\log\log N)$.
\end{lemma}
\begin{proof}
For primes $p\geq2$,
$-\log(1-p^{-1/2})\leq C_0p^{-1/2}$.
The elementary prime-counting upper bound
$\pi(x)\ll x/\log x$ (for example
\cite[Corollary 1, Eq. (3.6)]{RosserSchoenfeld1962}) and partial
summation give $\sum_{p\leq L}p^{-1/2}\ll\sqrt L/\log L$.
For primes dividing $Q$ that exceed $L$, use
\[
 \sum_{\substack{p\mid Q\\p>L}}p^{-1/2}
 \leq\frac1{\sqrt L\log L}\sum_{p\mid Q}\log p
 \leq\frac{\sqrt L}{\log L}.
\]
The last inequality also holds when $\log Q<3$. Increasing the absolute
constant covers bounded $L$ and proves the result.
\end{proof}

\begin{proposition}[Growing Euler-grid corrections do not give a logarithmic norm]
\label{prop:ns60-growing-euler-obstruction}
Fix $0\leq\alpha<1$. For each dyadic $N$, let $Q_N$ be any squarefree
positive integer with $Q_N\leq N^\alpha$. For every fixed $B\geq0$,
\begin{equation}
 \norm{\mathcal C_{N,Q_N}}_{\rm op}
 \ne O\left(N^{-2}(\log N)^B\right)
 \quad\text{as }N\to\infty\text{ through powers of two}.
 \label{eq:ns60-growing-no-log}
\end{equation}
The excluded bound is uniform over all new coefficient vectors for the
explicit correction~\eqref{eq:ns60-euler-correction}. It is not a lower
bound for $K_N$, and it does not exclude such a bound restricted to the
actual residual vector $g$.
\end{proposition}
\begin{proof}
Use the fixed nonnegative smooth $w$ from the proof of
Proposition~\ref{prop:v158-lindelof-budget}. Approximate
\[
 F(u)=u^{-1/2}w(\log(u/N))e^{i\tau\log(u/N)}
\]
by its midpoint values on the complete $Q$-cells contained in $(N,2N]$,
and set the other cells to zero. Call the result $F_{N,Q,\tau}$ and set
$a_n=n(F_{N,Q,\tau})_n$. Because $w$ has compact support away from the
two endpoints, for sufficiently large $N$ and $Q\leq N^\alpha$ no
nonzero part is lost in a partially intersecting endpoint cell.
The mean value theorem gives, uniformly in this range,
\begin{equation}
 \norm{q_{F_{N,Q,\tau}}
       -w(\,\cdot-\log N)e^{i\tau(\,\cdot-\log N)}}_{L^1}
 \leq C_w\frac{Q(1+|\tau|)}N,
 \qquad \norm{a}_2^2\leq C_wN^2.
 \label{eq:ns60-packet-grid-error}
\end{equation}
Thus whenever $Q(1+|\tau|)/N$ tends uniformly to zero, the Fourier
transform is bounded below by a fixed positive constant on
$[\tau-1,\tau+1]$. On this compatible packet,
\eqref{eq:ns60-compatible-euler} and
$\prod_{p\mid Q}|1-p^{-1/2-it}|\geq m_Q$ imply
\begin{equation}
 \int_{\tau-1}^{\tau+1}|\zeta(\tfrac12+it)|^2\,dt
 \leq C m_Q^{-2} N^2\norm{\mathcal C_{N,Q}}_{\rm op}.
 \label{eq:ns60-corrected-local-test}
\end{equation}

Suppose the excluded big-$O$ bound held. For large $T$, choose a power
of two $N$ with
$T^{2/(1-\alpha)}\leq N<2T^{2/(1-\alpha)}$.
Then $Q_N(1+2T)/N\leq(1+2T)/T^2\to0$, uniformly for
$|\tau|\leq2T$. Cover the interval in
\eqref{eq:v158-short-interval} by $O(\log^2T)$ unit intervals and use
\eqref{eq:ns60-corrected-local-test}. Lemma~\ref{lem:ns60-euler-minimum}
gives
\[
 |\zeta(\tfrac12+iT)|^2
 \ll \log T+
 \exp\left(C_\alpha\frac{\sqrt{\log T}}{\log\log T}\right)
 (\log T)^{B+3}.
\]
The logarithm of the right side is
$O_\alpha(\sqrt{\log T}/\log\log T)+(B+3)\log\log T$.
It is $o(\sqrt{\log T/\log\log T})$, contradicting the unconditional
large values in~\cite[Theorem 1]{Soundararajan2008}.
No relation between the choices $Q_N$ at different $N$ is required.
\end{proof}

\paragraph{What changed and what remains open.}
Downsampling does give an admissible old-space correction, but on a
general new vector it requires a cell average. Replacing that average
by an exact Euler multiplier loses~\eqref{eq:ns60-grid-defect} and its
mixed term. Even on a compatible family where the replacement is exact,
known zeta large values defeat the stated uniform norm goal whenever
$Q_N\leq N^\alpha$ for fixed $\alpha<1$.
Since $K_N\leq\mathcal C_{N,Q_N}$, this failure of the explicit upper
comparison does not give a lower obstruction for the optimal projection.
The range $Q_N=N^{1-o(1)}$, tapered non-Euler coefficients and
actual-residual-only estimates are outside the proved exclusion.

The separately labelled scan in
\texttt{evidence/diag\_ns60\_nb\_corrections/} compares fixed Euler
products $Q=2,6,30$ and tapered M\"obius lengths $4,16,64$ with the full
projection on fixed blocks $32\leq N\leq512$. It tests the correction
proposal only; its finite reductions and selected-direction values are
not asymptotic bounds. The remaining RH-relevant input is still a
cofinal gain estimate for the actual optimal residuals. No such estimate,
NB convergence, G2 or RH proof is supplied.

\section{One Mellin integration: an elementary block budget (NS-61)}
\label{sec:ns61-smoothed-nb}

\paragraph{Status and prior art.}
The $q=2$ weighted Nyman--Beurling problem already appears in
\cite{Ehm2024}; it is not a new RH criterion. We prove its equivalence
here with the tail space stated, and use adjacent differences to obtain
an unconditional $N^{-2}$ block budget. The numerator is transformed
along with the norm. No asymptotic lower correlation estimate, NB
convergence, G2 or RH proof follows from this upper budget.

Let $H=L^2(0,\infty;dx)$ and retain $\rho_n$, $\chi$, $V_N$ and
$\kappa=\log(2\pi)-\gamma$ from Section~\ref{sec:ns53-nb-blocks}.
Define the bounded Mellin convolution
\[
 (Jf)(x)=\int_0^1f(x/u)\,\frac{du}{u},\qquad
 \sigma_n=J\rho_n,\qquad \tau=J\chi=(-\log x)\mathbf1_{(0,1)}.
\]
On the Mellin line $\Re s=1/2$, $J$ has multiplier $1/s$, so
$\norm J=2$ and $J$ is injective. Set
\[
 W_N=\operatorname{span}\{\sigma_1,\ldots,\sigma_N\},\quad
 R_N=\tau-\Pi_N\tau,\quad E_N^{(2)}=\norm{R_N}^2,
\]
where $\Pi_N$ is the orthogonal projection onto $W_N$.
The superscript is a smoothing label; $E_N^{(2)}$ is already a
squared norm, and $\norm\tau^2=2$.

\begin{proposition}[The smoothed target still detects an off-line zero]
\label{prop:ns61-rh-equivalence}
Every finite $E_N^{(2)}$ is positive, and
\[
 E_N^{(2)}\longrightarrow0\quad\Longleftrightarrow\quad\mathrm{RH}.
\]
If $\zeta(\beta+i\gamma_*)=0$ with $1/2<\beta<1$, put
$z=\beta+i\gamma_*$. Then, for every $N$,
\begin{equation}
 E_N^{(2)}\geq
 \frac{|z|^{-4}}{(2\beta-1)^{-1}+|1-z|^{-2}}>0.
 \label{eq:ns61-zero-obstruction}
\end{equation}
\end{proposition}
\begin{proof}
The original strong NB criterion and boundedness of $J$ give
$E_N^{(2)}\leq4E_N$, proving the forward approximation under RH.
For the converse, the functions $\tau$ and $\sigma_n$ belong to the
closed space
\[
 \mathcal W=L^2(0,1)\oplus
       \operatorname{span}\{x^{-1}\mathbf1_{(1,\infty)}\}.
\]
Indeed $\sigma_n(x)=1/(nx)$ for $x>1$. On this space the functional
\[
 L_z(f)=\int_0^1 f(x)x^{z-1}\,dx+\frac{c}{1-z},
 \qquad f(x)=c/x\quad(x>1),
\]
is bounded with norm squared $(2\beta-1)^{-1}+|1-z|^{-2}$.
The absolutely convergent Mellin identities on $0<\Re s<1$ are
\[
 \mathcal M\sigma_n(s)=-\frac{n^{-s}\zeta(s)}{s^2},
 \qquad \mathcal M\tau(s)=\frac1{s^2}.
\]
Thus $L_z(\sigma_n)=0$ and $L_z(\tau)=z^{-2}$, which proves
\eqref{eq:ns61-zero-obstruction}. If RH fails, the functional equation
supplies a zero on the right of the critical line. This proves the
converse without assuming a bounded inverse for $J$.
Injectivity of $J$ also transfers finite linear independence and
$\chi\notin V_N$ from the original finite problem, so $E_N^{(2)}>0$.
\end{proof}

\begin{proposition}[A unitary average identity and an elementary trace budget]
\label{prop:ns61-average-budget}
Put $U=J-I$ and $h=U\rho_1$. The operator $U$ is unitary on $H$,
$\norm h^2=\kappa$, and $h$ is nonnegative and supported in $(0,1]$.
For $n\geq2$ define
\[
 \epsilon_n=\sigma_n-\frac{n-1}{n}\sigma_{n-1}=J\delta_n.
\]
Then, as a Bochner integral in $H$,
\begin{equation}
 n\epsilon_n=\int_{n-1}^{n}h(u\,\cdot)\,du
       =U\int_{n-1}^{n}\rho_1(u\,\cdot)\,du,
 \qquad
 \norm{\epsilon_n}^2\leq
       \frac{\kappa}{n^2}\log\frac n{n-1}.
 \label{eq:ns61-average}
\end{equation}
In particular $\epsilon_n$ is nonnegative and supported in
$(0,1/(n-1)]$, with
\begin{equation}
 \int_0^\infty x\epsilon_n(x)\,dx
             =\frac{\pi^2}{24n^2(n-1)}.
 \label{eq:ns61-positive-first-moment}
\end{equation}
If $\mathcal D_N^{(2)}$ is the raw Gram of
$\epsilon_{N+1},\ldots,\epsilon_{2N}$, then
\begin{equation}
 \norm{\mathcal D_N^{(2)}}_{\rm op}
 \leq\operatorname{tr}\mathcal D_N^{(2)}\leq \mathfrak b_N,
 \qquad \mathfrak b_N=\kappa\sum_{n=N+1}^{2N}
          \frac{\log(n/(n-1))}{n^2}
 \leq\frac{3\kappa}{8N^2}.
 \label{eq:ns61-trace-budget}
\end{equation}
Moreover $N^2\operatorname{tr}\mathcal D_N^{(2)}\to3\kappa/8$.
These are unconditional analytic bounds and a trace asymptotic, not
an asymptotic formula for the residual gain.
\end{proposition}
\begin{proof}
The Mellin multiplier of $U$ is $(1-s)/s$, which has modulus one on
the critical line and is nonzero there. Hence $U$ is unitary.
Writing
\[
 A(z)=\int_0^z\frac{\{v\}}v\,dv,
 \qquad \sigma_u(x)=A(1/(ux))\quad(u>0),
\]
gives $u\partial_u\sigma_u=-\rho_u$ almost everywhere. Therefore
$\partial_u(u\sigma_u)=\sigma_u-\rho_u=h(u\,\cdot)$.
On compact positive $u$ intervals this is also the strong integral
identity: dilations are strongly continuous in $H$, and integrating
the displayed derivative recovers the pointwise formula and an
$H$-integrable majorant. Since $\sigma_1(x)=\rho_1(x)=1/x$ for
$x>1$, $h$ has the asserted support.
For $z\in[k,k+1)$ with $k\geq1$, integration in the definition of
$A$ gives $A(z)=z-k\log z+\log(k!)$, hence
$h(1/z)=k-k\log z+\log(k!)>0$.
Indeed monotonicity of $\log u$ gives
$\int_1^{k+1}\log u\,du\leq\log((k+1)!)$, or
$\log(k!)\geq k\log(k+1)-k$; use $z<k+1$.
Positivity of $\epsilon_n$ follows from its average representation.
The compactly supported $h$ has Mellin transform
$-(1-s)\zeta(s)/s^2$, which extends analytically to $\Re s>0$.
At $s=2$ it gives $\int xh(x)\,dx=\pi^2/24$. Integrating the
average identity against $x\,dx$ proves
\eqref{eq:ns61-positive-first-moment}.
The Cauchy--Schwarz inequality for a unit-length Bochner integral gives
\[
 n^2\norm{\epsilon_n}^2\leq
 \int_{n-1}^{n}\norm{h(u\,\cdot)}^2\,du
 =\kappa\int_{n-1}^{n}\frac{du}{u}.
\]
For $n-1\leq u\leq n$, $1/(n^2u)\leq u^{-3}$, whence
\[
 \mathfrak b_N\leq\kappa\int_N^{2N}u^{-3}\,du
       =\frac{3\kappa}{8N^2}.
\]
Finally, dilation continuity gives
$n^3\norm{\epsilon_n}^2\to\kappa$: rescale $x=y/n$ in
\eqref{eq:ns61-average}, then $h((u/n)y)\to h(y)$ strongly,
uniformly for $n-1\leq u\leq n$. Summation on $N<n\leq2N$
and $N^2\sum n^{-3}\to3/8$ prove the trace asymptotic.
\end{proof}

The associated exact Gram identity, which retains all off-diagonal
entries, is
\begin{equation}
 \langle\epsilon_m,\epsilon_n\rangle
 =\frac1{mn}\int_{m-1}^{m}\int_{n-1}^{n}
        \langle\rho_u,\rho_v\rangle\,du\,dv.
 \label{eq:ns61-averaged-full-gram}
\end{equation}
Thus smoothing converts the singular adjacent-difference budget into
an average of the original complete Gram kernel. It does not diagonalize
that kernel or remove the old-space projection.

\begin{proposition}[Complete smoothed gain and its still-missing numerator]
\label{prop:ns61-gain}
Form the old Gram $G^{(2)}$, old/new block $C^{(2)}$, new Gram
$D^{(2)}$, and target loads $b^{(2)}$ using $\sigma_n$ and $\tau$.
Set
\[
 c=(G^{(2)})^{-1}b^{(2)},\quad
 h^{(2)}=b_{\rm new}^{(2)}-(C^{(2)})^*c,\quad
 S^{(2)}=D^{(2)}-(C^{(2)})^*(G^{(2)})^{-1}C^{(2)}.
\]
For the same triangular new-coordinate matrix $T_N$ as in
Proposition~\ref{prop:v158-difference-gain}, put
$g^{(2)}=T_N^*h^{(2)}$ and $K_N^{(2)}=T_N^*S^{(2)}T_N$.
Then
\begin{equation}
 0<K_N^{(2)}\leq\mathcal D_N^{(2)},\qquad
 E_N^{(2)}-E_{2N}^{(2)}
 =(g^{(2)})^*(K_N^{(2)})^{-1}g^{(2)}
 \geq\frac{\norm{g^{(2)}}_2^2}{\mathfrak b_N}
 \geq\frac{8N^2}{3\kappa}\norm{g^{(2)}}_2^2.
 \label{eq:ns61-gain}
\end{equation}
Consequently the cofinal estimate
\begin{equation}
 N^2\norm{g_{2^j}^{(2)}}_2^2
 \geq\frac{3\kappa}{8}\alpha_j E_{2^j}^{(2)},
 \quad N=2^j,\quad 0\leq\alpha_j<1,\quad
 \sum_j\alpha_j=\infty
 \label{eq:ns61-open-correlation}
\end{equation}
for all sufficiently large $j$ would prove RH. This lower estimate
is an open input; no necessity claim is made for it.
\end{proposition}
\begin{proof}
The finite Gram matrices are positive definite by injectivity of $J$.
Orthogonal projection, the complete Schur complement and the invertible
coordinate change give the exact identities as in NS-53.
The raw trace budget bounds the largest eigenvalue of the projected
Gram. The resulting contraction
$E_{2^{j+1}}^{(2)}\leq(1-\alpha_j)E_{2^j}^{(2)}$ proves the
conditional conclusion. In particular, the unsmoothed residual
correlation cannot be inserted in~\eqref{eq:ns61-gain}.
\end{proof}

\begin{proposition}[Explicit loads do not supply a projected lower bound]
\label{prop:ns61-loads}
Use the standard convention
$\zeta(1+z)=z^{-1}+\gamma-\gamma_1z+O(z^2)$ for $\gamma_1$.
Then
\begin{equation}
 b_n^{(2)}=\langle\tau,\sigma_n\rangle
 =\frac{P(\log n)}n,
 \quad P(v)=\tfrac12v^2+(2-\gamma)v+3-2\gamma-\gamma_1.
 \label{eq:ns61-loads}
\end{equation}
In difference coordinates, for $n\geq2$,
\[
 \langle\tau,\epsilon_n\rangle
 =\frac{P(\log n)-P(\log(n-1))}{n}
 \sim\frac{\log n+2-\gamma}{n^2}.
\]
For the actual residual the exact expression is instead
\begin{equation}
 g_n^{(2)}=\frac{P(\log n)-P(\log(n-1))}{n}
       -\sum_{m\leq N}c_m\langle\sigma_m,\epsilon_n\rangle,
 \quad N<n\leq2N.
 \label{eq:ns61-retained-cross}
\end{equation}
The sign and size of the second term are not controlled here.
\end{proposition}
\begin{proof}
Since $\sigma_1(x)=1/x$ for $x>1$, the transform
$-\zeta(s)/s^2-1/(1-s)$ is the Mellin transform of the restriction
to $(0,1)$. Its value and derivative at $s=1$ give
\[
 \int_0^1\sigma_1(x)\,dx=2-\gamma,\qquad
 \int_0^1(-\log x)\sigma_1(x)\,dx=3-2\gamma-\gamma_1.
\]
Change variables $u=nx$ in the load integral and split at $u=1$.
The tail contribution is $(\log n)^2/(2n)$, giving
\eqref{eq:ns61-loads}. Taking the adjacent difference and then
subtracting the old projection gives the remaining identities.
The loads agree, with the sign convention for $\sigma_n$ reversed,
with the $q=2$ formula in~\cite{Ehm2024}.
\end{proof}

\paragraph{The exact point at which the lower-bound attempt stops.}
Every new $\epsilon_n$ is supported in $(0,1/N]$, so the numerator
only sees the small-$x$ part of $R_N$. The elementary formula
\[
 A(z)=z-k\log z+\log(k!),\qquad k=\lfloor z\rfloor,
\]
and Stirling's estimate imply, uniformly for $z\geq1$,
$|A(z)-\tfrac12\log(2\pi z)|\leq3/z$; a complete elementary
remainder is given in Section~\ref{sec:ns61-gram-remainder}.
Thus for $0<x\leq1/N$,
\begin{equation}
 \begin{aligned}
 R_N(x)&=a_N\log(1/x)+\ell_N+\mathcal R_N(x),\\
 a_N&=1-\tfrac12\sum_{m\leq N}c_m,\qquad
 \ell_N=\tfrac12\sum_{m\leq N}c_m\log\frac m{2\pi},\\
 |\mathcal R_N(x)|&\leq3x\sum_{m\leq N}m|c_m|.
 \end{aligned}
 \label{eq:ns61-residual-endpoint}
\end{equation}
Writing $\Delta_n f=f(n)-f(n-1)$, the exact first two moments of
$\epsilon_n$ give
\begin{equation}
 g_n^{(2)}=
 \frac{a_N\{\tfrac12\Delta_n(\log^2)+(2-\gamma)\Delta_n\log\}
              +\ell_N\Delta_n\log}{n}
 +\mathcal E_{N,n},
 \quad N<n\leq2N,
 \label{eq:ns61-endpoint-correlation}
\end{equation}
where
\[
 |\mathcal E_{N,n}|\leq
       \frac{\pi^2\sum_{m\leq N}m|c_m|}{8n^2(n-1)}.
\]
The two leading moments follow from $\int h=1$ and
$\int(-\log x)h(x)\,dx=2-\gamma$, obtained by taking the value
and derivative of $-(1-s)\zeta(s)/s^2$ at $s=1$.
For the error, use nonnegativity and the exact first moment
\eqref{eq:ns61-positive-first-moment}. Neither a lower bound for the displayed
coefficient combination relative to $E_N^{(2)}$, nor a small enough
bound for $\sum m|c_m|$, has been obtained. Replacing these optimized
coefficients by the target loads would discard
\eqref{eq:ns61-retained-cross}.

The improvement is therefore a proved elementary denominator and an
explicit description of the remaining arithmetic numerator. It bypasses
the raw unsmoothed norm obstruction by changing both the functions and
the target. It does not bypass the zero obstruction, certify a lower
correlation, or establish convergence. This is the stop condition for
NS-61 unless an independent estimate for the optimized residual is found.

\subsection{A complete remainder for the smoothed Gram computation}
\label{sec:ns61-gram-remainder}

This subsection supplies the analytic tail used by the finite Arb replay;
the diagnostic truncations do not enter that replay. Write
$\widetilde B_j(x)=B_j(\{x\})$ and let $R_1,R_2,S_2$ be the
functions in~\cite[Theorem 2.1]{Ehm2024}. For $r=p/q\geq1$ in lowest
terms the Gram formula used is
\[
 G^{(2)}_{m,n}=\frac{K_2+K_1\log(n/m)+\tfrac14\log^2(n/m)}n
                  +\frac{S_2(n/m)}m\quad(m\leq n),
\]
with
\[
 K_1=\frac{\log(2\pi)-\gamma+2}{2},\qquad
 K_2=(1-\gamma/2)\log(2\pi)+\tfrac14\log^2(2\pi)
       +\frac{\pi^2}{48}-\frac{\gamma^2}{4}-\gamma-\gamma_1+\frac32.
\]
The first $K$ terms of $S_2(r)=\sum_{k\geq1}R_2(kr)$ are evaluated
by the finite formula in that theorem, retaining balls throughout.
The rest is enclosed by the following lemma.

\begin{lemma}[Bernoulli tail with an explicit remainder]
\label{lem:ns61-r2-tail}
For an integer $M\geq4$, let $c_2=c_3=0$ and, for $j\geq4$, let
\[
 c_j=\frac{2j-3-(j-1)H_{j-1}}{j(j-1)},\qquad
 P_M(x)=\sum_{j=4}^{M+2}c_j\frac{\widetilde B_j(x)}{x^j}.
\]
Put
\[
 a_{M+1}=2(M+2)(M+1)c_{M+2}-(M+1)Mc_{M+1},\qquad
 a_{M+2}=-(M+2)(M+1)c_{M+2}.
\]
For any constants $U_j\geq\sup_{0\leq x\leq1}|B_j(x)|$,
\begin{equation}
 |R_2(x)-P_M(x)|\leq
 \frac{U_M}{M^2(M-1)x^M}
 +\frac{|a_{M+1}|U_{M+1}}{M(M+1)x^{M+1}}
 +\frac{|a_{M+2}|U_{M+2}}{(M+1)(M+2)x^{M+2}}.
 \label{eq:ns61-r2-remainder}
\end{equation}
The rational constants
$U_j=2j!\,(1+1/(j-1))/6^j$, $j\geq2$, are valid.
For $r=p/q$ the Bernoulli part of the infinite sum is exactly
\begin{equation}
 \sum_{k>K}\frac{\widetilde B_j(kr)}{(kr)^j}
 =p^{-j}\sum_{a=1}^{q}
 B_j\left(\left\{\frac{(K+a)p}{q}\right\}\right)
 \zeta\left(j,\frac{K+a}{q}\right).
 \label{eq:ns61-hurwitz-tail}
\end{equation}
The sum of the three remainders in~\eqref{eq:ns61-r2-remainder}
is bounded by replacing each $x^{-l}$ by
$r^{-l}\zeta(l,K+1)$.
\end{lemma}
\begin{proof}
From $R_1(x)=-\int_x^\infty\widetilde B_1(t)t^{-2}\,dt$,
successive integrations by parts give the exact finite formula
\[
 R_1(x)=\sum_{j=2}^{M}\frac{\widetilde B_j(x)}{jx^j}
            -\int_x^\infty\frac{\widetilde B_M(t)}{t^{M+1}}\,dt.
\]
The remainder is at most $U_M/(Mx^M)$. Also
\cite[Proposition 5.2, Eq. (40)]{Ehm2024} gives
$-(x^2R_2'(x))'=R_1(x)$ and the inverse formula
\[
 (\mathcal If)(x)=\int_x^\infty f(t)\left(\frac1t-\frac1x\right)dt.
\]
If $|f(t)|\leq At^{-l}$ with $l>1$, then
$|\mathcal If(x)|\leq A x^{-l}/(l(l-1))$.
For $\mathcal L=-(x^2\partial_x)'$ the coefficient of
$\widetilde B_k(x)/x^k$ in $\mathcal LP_M$ is
\[
 -(k+2)(k+1)c_{k+2}+2(k+1)kc_{k+1}-k(k-1)c_k.
\]
It equals $1/k$ for $2\leq k\leq M$; the two remaining coefficients
are $a_{M+1},a_{M+2}$. Periodic Bernoulli polynomials of orders at
least four are twice continuously differentiable across integers, so
no point masses are omitted in this differentiation. Apply
$\mathcal I$ to $R_1-\mathcal LP_M$ and the preceding three bounds.
The possible homogeneous terms $A+B/x$ vanish because both $R_2$
and $P_M$ are $O(x^{-4})$. This proves~\eqref{eq:ns61-r2-remainder}.

The absolutely convergent Bernoulli Fourier series
\cite[Eq. (24.8.3)]{DLMFBernoulliFourier} gives
$\sup|B_j|\leq2j!\zeta(j)/(2\pi)^j$.
Using $2\pi>6$ and $\zeta(j)\leq1+1/(j-1)$ proves the rational
upper bound. Finally write $k=K+a+q\ell$ and sum over $\ell\geq0$.
The periodic numerator is constant in each such progression, giving
\eqref{eq:ns61-hurwitz-tail}. Absolute convergence justifies every
rearrangement and the sum of the absolute remainder bounds.
\end{proof}

\paragraph{An independent positive-cell convention check.}
On a cell where $k=\lfloor t/m\rfloor$,
$A(t/m)=t/m-k\log t+k\log m+\log(k!)$.
Products of two such expressions divided by $t^2$ are integrated
exactly between consecutive multiples of $m$ and $n$; the initial
cell has constant integrand $1/(mn)$.
For $z=k+r\geq1$, $0\leq r<1$, the elementary Stirling remainder
\cite[\S5.11]{DLMFGammaBounds}
$0<\log(k!)-(k+1/2)\log k+k-\tfrac12\log(2\pi)<1/(12k)$
gives
\[
 \left|A(z)-\tfrac12\log(2\pi z)\right|
 \leq\frac{13}{12k}\leq\frac{13}{6z}<\frac3z.
\]
Indeed $0\leq r-k\log(1+r/k)\leq r^2/(2k)$ and
$0\leq\tfrac12\log(1+r/k)\leq r/(2k)$.
For $T\geq\max(m,n)$ put $L_m=\tfrac12\log(2\pi T/m)$.
The remaining physical integral is within
\[
 \frac3{T^2}\left[m\left(\frac{L_n}2+\frac18\right)
                 +n\left(\frac{L_m}2+\frac18\right)\right]
            +\frac{3mn}{T^3}
\]
of $(L_mL_n+(L_m+L_n)/2+1/2)/T$.
This separate integral enclosure checks normalization and several
off-diagonal entries of the series formula; it is not used to sharpen
the block estimates.

\paragraph{Certified finite checks, not a cofinal estimate.}
The complete smoothed Grams through index 64 were enclosed at 256 and
384 bits with the entire series tail retained. A separate replay doubles
the finite series cutoff from 128 to 256 while keeping the same matrix.
All 180 enclosure comparisons across the precision and cutoff replays pass; the minimum internal accuracies are
99, 99 and 123 bits. Seven independent positive-cell Gram integrals
also pass at both precisions.
\begin{center}
\begin{tabular}{rccc}
\toprule
$N$ & exact relative gain & trace-budget relative bound & bound/gain\\
\midrule
16 & $0.2372848\ldots$ & $0.003293361\ldots$ & $0.01387935\ldots$\\
32 & $0.1891280\ldots$ & $0.003566353\ldots$ & $0.01885681\ldots$\\
\bottomrule
\end{tabular}
\end{center}
The trace lower estimate is below $1/50$ of the exact gain on both
blocks. The endpoint remainder budget is more than 2000 times the norm
of its leading term there, so that particular triangle lower bound is
vacuous, despite every individual remainder gate passing.
These are finite statements, not a failure of
\eqref{eq:ns61-open-correlation} or an asymptotic claim.

\section{Smoothing does not repair the canonical sharp interpolant (NS-62)}
\label{sec:ns62-sharp-interpolant}

\paragraph{Status and scope.}
The previous section changes the approximation norm and proves a smaller
upper budget. It is natural to ask whether simply smoothing the canonical
M\"obius interpolant now supplies convergence. It does not. The following
localized dual test extends the classical sharp-interpolant obstruction
in~\cite[Section 4]{BaezDuarte2000Arithmetic} through Mellin integration.
The conclusion concerns this explicit coefficient rule, including its
tail correction. It does not concern the optimized coefficients in
Proposition~\ref{prop:ns61-gain}.

Define
\[
 M(x)=\sum_{k\leq x}\mu(k),\quad
 s_N=\sum_{k\leq N}\frac{\mu(k)}k,\quad
 \eta_N=s_N-\frac{M(N)}N
       =\sum_{k<N}\frac{M(k)}{k(k+1)},
\]
and let
\begin{equation}
 B_N=\sum_{k=1}^N\mu(k)\rho_k-Ns_N\rho_N.
 \label{eq:ns62-natural-interpolant}
\end{equation}
The sign is chosen to match the classical notation: $B_N=-\chi$
on $x>1/N$ (up to endpoints). The sum is zero on $x>1$.
The smoothed target error is $J^r(\chi+B_N)$.

\begin{proposition}[A localized witness survives any number of integrations]
\label{prop:ns62-smoothed-obstruction}
For every integer $N\geq1$ and integer $r\geq0$,
\begin{equation}
 \norm{J^r(\chi+B_N)}_H
 \geq\frac{|N\eta_N-k_{N,r}|}{\sqrt N}
 \geq\frac{|N\eta_N-2|}{\sqrt N}
            -\frac{4\pi^2}{3^{r+2}N^{5/2}},
 \label{eq:ns62-laguerre-lower}
\end{equation}
where
\[
 k_{N,r}=\frac N{r!}\int_0^{1/N}
       \frac{\sin(2\pi x)}{\pi x}
                    \left(\log\frac1{Nx}\right)^r dx.
\]
In particular, for every sequence of nonnegative integers $r_N$,
$\norm{J^{r_N}(\chi+B_N)}_H$ does not tend to zero.
For growing $r_N$ the target also changes with $N$; this assertion
is only a failure of these particular approximants, not a new RH
criterion for varying norms.
The nonconvergence is for the full sequence of integers $N$. No
exclusion of a selected subsequence, including the dyadic subsequence,
is asserted without a corresponding lower bound for $\sqrt N|\eta_N|$
on that subsequence.
\end{proposition}
\begin{proof}
Use the dilation-invariant unitary $\mathcal V$ whose critical-line
Mellin multiplier, away from the discrete zero set, is
\[
 v(s)=\frac{s\zeta(1-s)}{(1-s)\zeta(s)}.
\]
Its modulus is one on $\Re s=1/2$; define it arbitrarily with modulus
one on the measure-zero exceptional set. Mellin identities and the
functional equation give
\[
 \mathcal V\rho_k(x)=\frac{\{kx\}}{kx},\qquad
 \mathcal V\chi(x)=k(x):=\frac{\sin(2\pi x)}{\pi x}.
\]
These are also the defining identities of the classical unitary in
\cite[Section 4]{BaezDuarte2000Arithmetic}. Both $\mathcal V$ and $J$
are Mellin multipliers, so they commute. For $0<x<1/N$,
\begin{equation}
 \mathcal V(\chi+B_N)(x)=k(x)+M(N)-Ns_N=k(x)-N\eta_N.
 \label{eq:ns62-unitary-small-cell}
\end{equation}

Let $L_r(z)=\sum_{j=0}^r(-1)^j\binom rj z^j/j!$ be the ordinary
Laguerre polynomial and put
\[
 \psi_{N,r}(x)=\sqrt N\,
       L_r\left(\log\frac1{Nx}\right)\mathbf1_{(0,1/N)}(x).
\]
Laguerre orthogonality \cite[Table 18.3.1]{DLMFLaguerre} gives
$\norm{\psi_{N,r}}=1$. The adjoint is
$(J^*f)(x)=x^{-1}\int_0^x f(y)\,dy$. Repeated integration, or
Rodrigues' formula and integration by parts, gives the complete
identity
\begin{equation}
 (J^*)^r\psi_{N,r}(x)=
 \frac{(-1)^r\sqrt N}{r!}
       \left(\log\frac1{Nx}\right)^r\mathbf1_{(0,1/N)}(x).
 \label{eq:ns62-adjoint-laguerre}
\end{equation}
There is no exterior term: at the $j$th integration, $j\leq r$,
the potential exterior coefficient is a Laguerre moment of a
polynomial of degree less than $r$, hence zero. On the interior,
the polynomial identity follows by integrating monomials against
$e^{-z}$; equivalently $(1-\partial_z)^{-r}L_r(z)=(-1)^rz^r/r!$.
This also covers $r=0$ directly.

Taking the inner product with $\psi_{N,r}$, moving $J^r$ to the
adjoint, and using~\eqref{eq:ns62-unitary-small-cell} yields
\[
 \left|\left\langle J^r\mathcal V(\chi+B_N),\psi_{N,r}\right\rangle\right|
       =\frac{|k_{N,r}-N\eta_N|}{\sqrt N}.
\]
The unitary norm and Cauchy--Schwarz prove the first inequality.
For real $x$, $|k(x)-2|\leq4\pi^2x^2/3$. Since
\[
 \int_0^{1/N}x^2\left(\log\frac1{Nx}\right)^r dx
        =\frac{r!}{3^{r+1}N^3},
\]
we have $|k_{N,r}-2|\leq4\pi^2/(3^{r+2}N^2)$, proving the second.

For completeness, $\eta_N\ne o(N^{-1/2})$ follows from the known
existence of a critical-line zeta zero. Indeed let
$\eta(x)=\int_1^xM(t)t^{-2}\,dt$. For $\Re s>0$,
\[
 \int_1^\infty\eta(x)x^{-s-1}\,dx
        =\frac1{s(s+1)\zeta(s+1)}.
\]
The right side has a genuine pole at $s=\rho-1$ for any
critical-line zero $\rho$. If $\eta(x)=o(x^{-1/2})$, the integral
would be analytic on $\Re s>-1/2$, and the usual Abelian boundary
estimate would make $(\sigma+1/2)$ times its value at
$\sigma+i\Im\rho$ tend to zero as $\sigma\downarrow-1/2$.
This contradicts a pole of any positive order at that point.
Also $\eta(N)=\eta_N$ and
$\eta(x)-\eta(\lfloor x\rfloor)=O(1/x)$, so the discrete
little-$o$ assertion would imply the continuous one. Thus
$\limsup_N\sqrt N|\eta_N|>0$, whereas the error term in
\eqref{eq:ns62-laguerre-lower} tends to zero uniformly in $r\geq0$.
This proves the stated nonconvergence for every choice of $r_N$.
No assumption about zeros off the critical line was made.
\end{proof}

\begin{proposition}[Positive smoothed differences still require signed coefficients]
\label{prop:ns62-positive-cone}
Every $\epsilon_n$, $n\geq2$, is nonnegative. Nevertheless $\tau$
is not in the $H$-closure of the cone of finite nonnegative linear
combinations of $\sigma_1,\epsilon_2,\epsilon_3,\ldots$.
This is a restriction on that coefficient cone, not on their signed
linear span or on the Weil form.
\end{proposition}
\begin{proof}
Nonnegativity was proved in Proposition~\ref{prop:ns61-average-budget}.
Its average identity also gives strict positivity on
$(1/n,1/(n-1))$.

Triangular telescoping gives
\begin{equation}
 -JB_N=\sum_{n=2}^{N}n s_{n-1}\epsilon_n.
 \label{eq:ns62-pointwise-series}
\end{equation}
It equals $\tau$ on $x>1/N$, because $-B_N=\chi$ there and
$Jf(x)$ uses only values at arguments at least $x$.
For $N=6$ the coefficient of $\epsilon_6$ is
$6s_5=6(1-1/2-1/3-1/5)=-1/5$.

Restrict now to $x>1/6$. Every $\epsilon_n$ with $n\geq7$
vanishes there. The restrictions of
$\sigma_1,\epsilon_2,\ldots,\epsilon_6$ are linearly independent:
first use $x>1$ for $\sigma_1$, then the successive intervals
$(1/n,1/(n-1))$ for $n=2,\ldots,6$.
Their finite-dimensional nonnegative cone is closed, but the unique
representation of the restricted target has the negative coefficient
$-1/5$. Thus the restricted target has positive distance from this
cone. Restriction is contractive, proving the full-space claim.
\end{proof}

\paragraph{Consequences for the next attempt.}
The positive shapes of the smoothed differences do not make the
coefficient problem positive. The canonical pointwise interpolation
coefficients are explicit, but their partial sums still fail in norm
after smoothing. The optimized residual criterion from NS-61 remains
open and is unaffected by this exclusion. A next attempt must obtain
an independent estimate for those adaptive residuals, or a different
coefficient rule with a complete norm bound; visual smoothness,
pointwise exactness away from zero, and a small raw block norm do not
supply it.

\section{Fixed smoothing and a changing-norm control (NS-63)}
\label{sec:ns63-changing-norm}

\paragraph{Why this check is needed.}
NS-61 uses one fixed extra Mellin integration. Increasing the smoothing
order with the approximation size changes the problem. The following
control shows that relative errors can then vanish with just one fixed
dilation and a fixed coefficient, unconditionally. It is not convergence
in a fixed NB norm.

For an integer $r\geq0$ put
\[
 \tau_r=J^r\chi=\frac{(-\log x)^r}{r!}\mathbf1_{(0,1)},\quad
 \sigma_{n,r}=J^r\rho_n,\quad
 T_r=\norm{\tau_r}^2=\binom{2r}{r}.
\]
Let $E_{N,r}$ be the squared distance from $\tau_r$ to the span of
$\sigma_{1,r},\ldots,\sigma_{N,r}$, so $E_{N,1}=E_N^{(2)}$.

\begin{proposition}[Keep the smoothing order fixed in the RH criterion]
\label{prop:ns63-fixed-order}
For each fixed $r$, $E_{N,r}\to0$ if and only if RH. If
$z=\beta+i\gamma_*$ is a zero with $1/2<\beta<1$, then
\[
 E_{N,r}\geq
 \frac{|z|^{-2(r+1)}}{(2\beta-1)^{-1}+|1-z|^{-2}}.
\]
The lower obstruction itself depends on $r$. After division by $T_r$
it tends to zero for every such $z$, since $|z|>1/2$; it gives no
fixed positive obstruction for the relative error as $r\to\infty$.
\end{proposition}
\begin{proof}
Boundedness of $J^r$ gives $E_{N,r}\leq4^rE_N$ under the original
strong NB criterion. Every $\sigma_{n,r}$ has the same tail $1/(nx)$
for $x>1$, whereas $\tau_r$ has zero tail. Thus the bounded functional
$L_z$ in Proposition~\ref{prop:ns61-rh-equivalence} applies unchanged.
Its value on the target is $z^{-(r+1)}$ and it annihilates the finite
span. This proves the lower bound and the converse for fixed $r$.
\end{proof}

\begin{proposition}[One fixed atom has vanishing relative error as the norm changes]
\label{prop:ns63-one-atom}
Let $z_0=\zeta(1/2)$, $c_0=-1/z_0$, and
\[
 D_0=\frac{\zeta'(1/2)}{\zeta(1/2)}
     =\frac12\left(\gamma+\frac\pi2+\log(8\pi)\right)>0.
\]
Unconditionally, as $r\to\infty$ through integers,
\begin{equation}
 \frac{\norm{\tau_r-c_0\sigma_{1,r}}^2}{T_r}
    =\frac{D_0^2}{4(2r-1)}+O(r^{-2})
    \sim\frac{D_0^2}{8r}\longrightarrow0.
 \label{eq:ns63-relative-control}
\end{equation}
Nevertheless the absolute squared error for this fixed coefficient
diverges:
\begin{equation}
 \norm{\tau_r-c_0\sigma_{1,r}}^2
      \sim\frac{D_0^2}{8\sqrt\pi}\frac{4^r}{r^{3/2}}
      \longrightarrow\infty.
 \label{eq:ns63-absolute-control}
\end{equation}
Moreover $(J/2)^rf\to0$ in $H$ for every fixed $f\in H$.
Neither this normalized convergence nor
\eqref{eq:ns63-relative-control} tests RH.
\end{proposition}
\begin{proof}
The alternating-series expression for zeta shows $z_0<0$: its
numerator $\sum_{n\geq1}(-1)^{n-1}n^{-1/2}$ is positive and its
denominator $1-\sqrt2$ is negative. Thus $c_0$ is well-defined.
Mellin Plancherel writes the relative squared error as the average of
\[
 H(t)=\left|1-\frac{\zeta(1/2+it)}{z_0}\right|^2
\]
under the probability density proportional to
$(1/4+t^2)^{-(r+1)}dt$. The exact second and fourth moments of this
density are
\[
 \mathbb E_r(t^2)=\frac1{4(2r-1)},\qquad
 \mathbb E_r(t^4)=\frac3{16(2r-1)(2r-3)}
\]
for $r\geq1$ and $r\geq2$, respectively. They follow directly
from the beta integral. Near zero, conjugation symmetry and Taylor
expansion give $H(t)=D_0^2t^2+O(t^4)$.
Away from any fixed neighborhood of zero the normalized integral is
exponentially small in $r$. To justify this without an unproved zeta
estimate, use
\[
 \left|\frac{\zeta(1/2+it)}{1/2+it}\right|
 \leq\int_0^\infty\rho_1(x)x^{-1/2}\,dx<4,
\]
or the integrability of the original Mellin square; either controls
the entire exterior integral. The displayed moments then prove
\eqref{eq:ns63-relative-control}. Differentiating the zeta functional
equation at $1/2$, and using
$\psi(1/2)=-\gamma-2\log2$ \cite[Eq. (5.4.13)]{DLMFGammaBounds},
gives the stated value of $D_0$.
Finally $T_r=\binom{2r}{r}\sim4^r/\sqrt{\pi r}$ proves
\eqref{eq:ns63-absolute-control}.

For the last assertion the Mellin multiplier of $(J/2)^r$ is
$(2s)^{-r}$. Its modulus is at most one on $\Re s=1/2$ and tends
to zero for every $t\ne0$. Dominated convergence in the Mellin
$L^2$ norm applies to every fixed $f$. The exceptional point $t=0$
has measure zero.
\end{proof}

\paragraph{Finite control and its full tail.}
The companion Arb calculation integrates the analytic product
\[
 \frac{(1-\zeta(s)/z_0)(1-\zeta(1-s)/z_0)}
             {\pi T_r\{s(1-s)\}^{r+1}},\qquad s=\tfrac12+it,
\]
over $0\leq t\leq A$. It uses reflection $1-s$, not complex
conjugation of the integration variable, so the quadrature callback
is meromorphic. The preceding elementary bound supplies the complete
positive tail bound
\begin{equation}
 \frac1{\pi T_r}\left\{
       \frac{2A^{-2r-1}}{2r+1}
       +\frac{32A^{-2r+1}}{z_0^2(2r-1)}\right\},\qquad r\geq1.
 \label{eq:ns63-quadrature-tail}
\end{equation}
The fixed-order interval results illustrate the control, not the
asymptotic proof or any RH assertion.

The useful research problem is therefore the fixed-order cofinal
correlation estimate in~\eqref{eq:ns61-open-correlation}. Further
smoothing cannot be credited as a convergence result without retaining
one fixed norm and its target. The unsolved arithmetic estimate remains
the one for the optimized coefficients as $N$ increases.

\section{An explicit geometric separator for the positive coefficient cone (NS-64)}
\label{sec:ns64-cone-separator}

The coefficient-cone exclusion in Proposition~\ref{prop:ns62-positive-cone}
has a small, explicit Hilbert-space witness. It quantifies that particular
failure without imposing any restriction on signed NB approximation.
Let $I=(1/6,\infty)$, put $p_1=\sigma_1|_I$ and
$p_n=\epsilon_n|_I$ for $2\leq n\leq6$, and write
\[
 G_I=(\langle p_m,p_n\rangle_{L^2(I)})_{m,n=1}^6,\quad
 \beta=(0,2,3/2,2/3,5/6,-1/5)^T,\quad
 d=G_I^{-1}e_6.
\]
The matrix $G_I$ is positive definite by the nested-support independence
proved in NS-62. In particular $d_6>0$.
The exact local target identity is $\tau|_I=\sum\beta_np_n$.

\begin{proposition}[A quantitative separation, confined to one coefficient cone]
\label{prop:ns64-cone-distance}
Define $\psi=\sum_{n=1}^6d_np_n$ on $I$ and extend it by zero outside
$I$. Then
\[
 \langle\psi,\sigma_1\rangle=0,\qquad
 \langle\psi,\epsilon_n\rangle=\mathbf1_{\{n=6\}}\quad(n\geq2),
 \qquad \langle\psi,\tau\rangle=-\frac15,
 \qquad \norm\psi^2=d_6.
\]
Consequently, if $\mathcal C_+$ denotes the closed cone generated by
$\sigma_1,\epsilon_2,\epsilon_3,\ldots$ with nonnegative real
coefficients, then
\begin{equation}
 \operatorname{dist}_H(\tau,\mathcal C_+)
 \geq\frac1{5\sqrt{d_6}}>0.00228323.
 \label{eq:ns64-positive-gap}
\end{equation}
The exact distance from $\tau|_I$ to the six-generator local cone is
$1/(5\sqrt{d_6})$; the interval replay encloses it between
$0.00228323766118601628$ and $0.00228323766118601629$.
The local equality is not asserted for the full-space distance.
\end{proposition}
\begin{proof}
The defining equation $G_Id=e_6$ gives the first six pairings and
$\norm\psi^2=d^TG_Id=d_6$. Every $\epsilon_n$, $n\geq7$, vanishes
on $I$. The target pairing follows from $\beta_6=-1/5$.
Thus $\langle\psi,f-\tau\rangle\geq1/5$ for every $f\in\mathcal C_+$;
continuity and Cauchy--Schwarz give the first inequality.

For the local equality, put
\[
 c=\beta+\frac{d}{5d_6},\qquad w=\frac{\psi}{5d_6}.
\]
Here $c_6=0$ exactly. The complete interval calculation below certifies
$c_1,\ldots,c_5>0$. Hence $\tau|_I+w=\sum c_np_n$ belongs to the
local cone. Its distance from $\tau|_I$ is $\norm w=1/(5\sqrt{d_6})$,
attaining the dual lower bound. Only the signs of these five finite
coefficients and the displayed decimal enclosure use certified
computation; the separating identity and formula for its bound are exact.
\end{proof}

\paragraph{Complete finite calculation, with no omitted tail.}
Under $t=1/x$, the domain $I$ becomes $0<t<6$ and its measure is
$dt/t^2$. Each $\sigma_n$ becomes $A(t/n)$. On the unit cell
$j<t<j+1$, $k=\lfloor j/n\rfloor$ and
\[
 A(t/n)=t/n-k\log t+k\log n+\log(k!).
\]
On $0<t<1$ the product $A(t/m)A(t/n)/t^2$ is exactly $1/(mn)$;
this cell includes the entire original $x>1$ tail. On the remaining
five cells the products have elementary primitives, retained in the
Maxima and Arb scripts. The triangular difference transformation gives
$G_I$ with every cross term present.

Arb at 256 and 384 bits, using certified preconditioned solves, gives
\[
 c\approx(0.00005096281,\ 1.99987424,\ 1.49456836,\
            0.68321323,\ 0.76905516,\ 0)^T,
\]
and
\[
 \frac1{25d_6}=0.000005213174217458189678\ldots.
\]
The figures in the vector are illustrative rounded displays of the
archived enclosures. All five positivity gates are strict, including
the small first coefficient. The two-precision comparison, complete
Gram and coefficient enclosures, and 39 exact Maxima checks are in
\texttt{evidence/v159/ns64/}. This is a separation of a specified
positive coefficient cone, not negative Weil energy, a lower bound for
the signed best-approximation error, or an obstruction to RH.

\section{Ordinary averaging does not remove the canonical obstruction (NS-65)}
\label{sec:ns65-cesaro}

The preceding results leave different coefficient rules open. One simple
repair of the sharp canonical rule can be tested exactly: ordinary
Ces\`aro averaging over its index. The conclusion below concerns a fixed
number of these averages. Logarithmic tapers, a growing number of averages,
and optimized coefficients are outside its scope.

For a sequence $(f_N)$ of vectors or scalars define
\[
 (Cf)_N=\frac1N\sum_{n=1}^Nf_n,\qquad
 B_N^{[q]}=(C^qB)_N,\qquad
 H_N^{[q]}=(C^q a)_N,\quad a_n=n\eta_n,
\]
where $B_n$ and $\eta_n$ are the canonical objects of NS-62 and
$q\geq0$ is a fixed integer. The averaging weights are nonnegative,
sum to one, and only involve indices $n\leq N$.

\begin{proposition}[A fixed number of Ces\`aro averages leaves a critical-zero pole]
\label{prop:ns65-cesaro-obstruction}
For every fixed integer $q\geq0$ and every sequence of nonnegative
integers $r_N$,
\[
 \norm{J^{r_N}(\chi+B_N^{[q]})}_H\not\longrightarrow0
 \quad(N\to\infty\text{ through all integers}).
\]
More precisely, for every $N$ and every $r\geq0$,
\begin{equation}
 \norm{J^r(\chi+B_N^{[q]})}_H
 \geq\frac{|H_N^{[q]}-k_{N,r}|}{\sqrt N}
 \geq\frac{|H_N^{[q]}-2|}{\sqrt N}
             -\frac{4\pi^2}{3^{r+2}N^{5/2}},
 \label{eq:ns65-averaged-witness}
\end{equation}
and $H_N^{[q]}\ne o(\sqrt N)$. This is full-sequence nonconvergence;
no selected-subsequence exclusion or assertion about signed optimal
approximation follows.
\end{proposition}
\begin{proof}
For $0<x<1/N$, every index $n\leq N$ obeys the same small-cell
identity from NS-62. Linearity and normalization of the weights give
\[
 \mathcal V(\chi+B_N^{[q]})(x)=k(x)-H_N^{[q]}.
\]
Apply the same unit Laguerre test $\psi_{N,r}$ and its complete
adjoint identity. This proves~\eqref{eq:ns65-averaged-witness},
including its remainder uniform in $r$.

It remains to prove the scalar obstruction. Put
\[
 \eta(x)=\int_1^x\frac{M(t)}{t^2}\,dt,\qquad
 A_0(x)=x\eta(x),\qquad
 A_{q+1}(x)=\frac1x\int_1^x A_q(u)\,du.
\]
For $\Re s>1$, absolute convergence and Fubini give
\begin{equation}
 \int_1^\infty A_q(x)x^{-s-1}\,dx
   =\frac1{(s+1)^q(s-1)s\zeta(s)}.
 \label{eq:ns65-cesaro-mellin}
\end{equation}
For $q=0$ this is the Mellin identity for $\eta$ with its variable
shifted by one. Each continuous averaging contributes exactly the
factor $(s+1)^{-1}$; there is no discarded endpoint term.

We now connect this continuous identity to the discrete averages.
Directly $A_0(x)=O(x\log(2x))$ and, almost everywhere,
\[
 A_0'(x)=\eta(x)+M(x)/x
         =\sum_{n\leq x}\mu(n)/n=O(\log(2x)).
\]
Inductively $A_q(x)=O_q(x\log(2x))$ and
$A_{q+1}'(x)=(A_q(x)-A_{q+1}(x))/x=O_q(\log(2x))$.
The functions are locally absolutely continuous. Comparing each
unit-cell integral with its endpoint value therefore gives
\[
 \frac1N\sum_{n\leq N}A_q(n)-\frac1N\int_1^N A_q(x)\,dx
       =O_q(\log(2N)).
\]
Since $A_0(N)=a_N$, induction yields
\begin{equation}
 H_N^{[q]}=A_q(N)+O_q(\log(2N)).
 \label{eq:ns65-discrete-continuous}
\end{equation}
Thus $H_N^{[q]}=o(\sqrt N)$ would imply $A_q(x)=o(\sqrt x)$
also between integers, by the displayed derivative bounds.

Under that hypothetical little-$o$ estimate the Mellin integral in
\eqref{eq:ns65-cesaro-mellin} is analytic for $\Re s>1/2$.
At a known critical-line zero $\rho=1/2+i\gamma_*$, its usual
Abelian boundary estimate gives
\[
 \lim_{\sigma\downarrow1/2}(\sigma-1/2)
       \int_1^\infty A_q(x)x^{-\sigma-i\gamma_*-1}\,dx=0.
\]
For example, split the integral at a fixed large $X$, use
$|A_q(x)|\leq\varepsilon\sqrt x$ on the tail, then let
$\sigma\downarrow1/2$ and $\varepsilon\downarrow0$.
But the meromorphic continuation on the right side of
\eqref{eq:ns65-cesaro-mellin} has a genuine pole at $\rho$:
none of $\rho+1$, $\rho-1$ or $\rho$ is zero, and its numerator
is one. A pole of any positive order contradicts the boundary limit.
Therefore $\limsup_N|H_N^{[q]}|/\sqrt N>0$.
Together with~\eqref{eq:ns65-averaged-witness}, whose remainder tends
to zero uniformly in $r\geq0$, this proves the stated nonconvergence.
\end{proof}

\paragraph{Interpretation and stop condition.}
Ordinary averaging attenuates a critical-zero contribution by
$(\rho+1)^{-q}$, but a fixed $q$ cannot annihilate it. This is an
analytic obstruction for the specified averaged sharp rule, not a
numerical observation or a priority claim about summability methods.
The general summability viewpoint is already part of
\cite[Section 3]{BaezDuarteStrong}. The useful open task remains an
independent estimate for the optimized signed residual, or a genuinely
different coefficient rule with a complete norm bound. No such lower
correlation or convergence estimate is supplied here.

\section{A weaker canonical target: growth rather than convergence (NS-66)}
\label{sec:ns66-canonical-growth}

Strong convergence of the full canonical sequence is excluded, but
that does not exclude slow growth or bounded subsequences of its norms.
One Mellin integration makes a complete Abel-summation estimate
available. The resulting growth exponent is an exact reformulation
of the zero-location question, not an independently established RH bound.

For a fixed integer $r\geq1$ define
\[
 D_{N,r}=\norm{J^r(\chi+B_N)}_H,\qquad
 \beta_* =\sup\{\Re z:\zeta(z)=0,\ 0<\Re z<1\}.
\]
Known critical-line zeros give $1/2\leq\beta_*\leq1$.

\begin{proposition}[Complete growth bounds for the smoothed canonical error]
\label{prop:ns66-growth-bounds}
If $z=\beta+i\gamma_*$ is a zero with $1/2<\beta<1$, then
\begin{equation}
 D_{N,r}\geq\frac{\sqrt{2\beta-1}}{|z|^{r+1}}N^{\beta-1/2}
 \quad\text{for every }N\geq1.
 \label{eq:ns66-zero-growth}
\end{equation}
Conversely, suppose for some $1/2\leq\theta<1$ and constant $C$
that $|M(u)|\leq Cu^\theta$ for all $u\geq1$. Then
\begin{equation}
 D_{N,1}\leq\sqrt2+
  \frac{C\norm{\sigma_1}}{1-\theta}N^{\theta-1/2}
  +C\sqrt\kappa\int_1^N u^{\theta-3/2}\,du,
 \qquad D_{N,r}\leq2^{r-1}D_{N,1}.
 \label{eq:ns66-mertens-upper}
\end{equation}
The integral is $(N^{\theta-1/2}-1)/(\theta-1/2)$ for
$\theta>1/2$, and $\log N$ for $\theta=1/2$.
The Mertens bound in this implication is a hypothesis, not a result
proved here.
\end{proposition}
\begin{proof}
Exact exterior interpolation gives $\chi+B_N=0$ on $x>1/N$.
Since $Jf(x)$ only uses arguments at least $x$, $J^r(\chi+B_N)$
is supported in $(0,1/N]$. Its Mellin value at $z$ is
$z^{-(r+1)}$: the zeta factor annihilates every basis contribution.
Cauchy--Schwarz on that support gives
\[
 |z|^{-(r+1)}\leq D_{N,r}
 \left(\int_0^{1/N}x^{2\beta-2}\,dx\right)^{1/2}
 =D_{N,r}\frac{N^{1/2-\beta}}{\sqrt{2\beta-1}},
\]
proving~\eqref{eq:ns66-zero-growth}.

Under the stated Mertens hypothesis, partial summation makes
$\sum\mu(n)/n$ convergent and identifies its value as
$\lim_{s\downarrow1}1/\zeta(s)=0$. Hence
\[
 \eta_N=-\int_N^\infty M(u)u^{-2}\,du,\qquad
 |\eta_N|\leq\frac{C}{1-\theta}N^{\theta-1}.
\]
The strong derivative $\partial_u\sigma_u=-\rho_u/u$ gives the
complete Bochner summation-by-parts identity
\begin{equation}
 JB_N=-N\eta_N\sigma_N+
           \int_1^N\frac{M(u)}u\rho_u\,du.
 \label{eq:ns66-bochner-abel}
\end{equation}
There is no lower-endpoint contribution, since $M(1^-)=0$.
Use $\norm{\sigma_N}=N^{-1/2}\norm{\sigma_1}$,
$\norm{\rho_u}=\sqrt{\kappa/u}$ and $\norm{J\chi}=\sqrt2$.
The triangle inequality yields~\eqref{eq:ns66-mertens-upper},
including all coefficient and tail-correction terms. Boundedness
of $J^{r-1}$ gives the final assertion.
\end{proof}

\begin{corollary}[The canonical norm-growth exponent records the rightmost zeros]
\label{cor:ns66-growth-exponent}
For every fixed integer $r\geq1$ the following limit exists and equals
\begin{equation}
 \lim_{N\to\infty}\frac{\log(1+D_{N,r})}{\log N}
                  =\beta_*-\frac12.
 \label{eq:ns66-exponent}
\end{equation}
In particular, RH is equivalent to the subpolynomial bound
\[
 \text{for every }\varepsilon>0,\qquad
 D_{N,r}=O_{r,\varepsilon}(N^\varepsilon).
\]
Even a bounded cofinal subsequence of $D_{N,r}$ would suffice to
prove RH. No such bounded subsequence or unconditional subpolynomial
bound is supplied, and boundedness is not claimed necessary under RH.
\end{corollary}
\begin{proof}
For each zero to the right of the line,
\eqref{eq:ns66-zero-growth} bounds the lower limit by
$\beta-1/2$. Taking the supremum gives $\beta_*-1/2$; when
$\beta_*=1/2$, the same lower limit is at least zero trivially.

If $\beta_*<1$, choose $\beta_*<\alpha<\theta<1$.
The classical zero-free-half-plane implication, stated precisely
in~\cite[Lemma 2.1]{BaezDuarteStrong}, gives convergence of
$\sum\mu(n)n^{-\theta}$. Partial summation of its bounded partial
sums gives $M(x)=O_\theta(x^\theta)$.
Proposition~\ref{prop:ns66-growth-bounds} then bounds the upper
limit by $\theta-1/2$. Let $\theta\downarrow\beta_*$.
This includes RH by choosing $\theta>1/2$ arbitrarily close to the
line, without claiming a bound at the endpoint $\theta=1/2$.

If $\beta_*=1$, the direct bounds $|s_N|\leq1+\log N$ and
$\norm{\sigma_n}=\norm{\sigma_1}/\sqrt n$ give
$D_{N,1}=O(\sqrt N\log(2N))$, so the upper limit is at most
$1/2$. This matches the preceding lower bound.
Thus~\eqref{eq:ns66-exponent} holds in every case.
The functional equation identifies $\beta_*=1/2$ with RH.
Finally~\eqref{eq:ns66-zero-growth} diverges along every cofinal
subsequence if any zero lies to the right of the line, proving the
bounded-subsequence implication.
\end{proof}

\paragraph{What the weaker target changes.}
For this explicit coefficient rule, norm convergence is false, while
subpolynomial norm growth is equivalent to RH. These are compatible
statements. The exact growth reformulation does not itself improve
our knowledge of $\beta_*$ or supply the missing Mertens estimate.
It gives a second precise research target: a direct subpolynomial
bound for the complete smoothed canonical norm, or a bounded cofinal
subsequence. Increasing the number of Mellin integrations changes
constants in~\eqref{eq:ns66-zero-growth}; the corollary is stated
only for one fixed order. No inference from a changing normalization
or finite table is used.

\section{Keeping canonical cancellation inside an exact cell norm (NS-67)}
\label{sec:ns67-canonical-cells}

\paragraph{Status.}
The canonical growth criterion in Section~\ref{sec:ns66-canonical-growth}
is still open. Here the complete once-smoothed error is evaluated before
any absolute-value estimate is taken. This gives exact arithmetic cell
identities, a smaller elementary tail bound, and a finite growing-interval
version of the same growth target. It does not provide the missing
subpolynomial estimate. No finite table proves a cofinal bound.

Set $e_N=J(\chi+B_N)$ and $F_N(t)=e_N(1/t)$, $t>0$.
Write
\[
 b_N=Ns_N,\qquad c_{N,n}=\mu(n)-b_N\mathbf1_{n=N},\qquad
 q_N(j)=\sum_{\substack{n\mid j\\n\leq N}}\mu(n).
\]
The letter $b_N$ in this section denotes the canonical boundary
coefficient, not a target-load vector. The identity
$\sum_{n\leq N}c_{N,n}/n=0$ is exact.

\begin{proposition}[Complete physical cells]
\label{prop:ns67-cells}
We have $F_N(t)=0$ for $0<t\leq N$, and, for every integer $j\geq N$,
\begin{equation}
 F_N(t)=f_{N,j}+a_{N,j}\log(t/j),\quad j\leq t\leq j+1,
 \qquad
 a_{N,j}=1-\sum_{n\leq N}\mu(n)\left\lfloor\frac jn\right\rfloor
                   +b_N\left\lfloor\frac jN\right\rfloor.
 \label{eq:ns67-cell}
\end{equation}
The endpoint values are generated by
\[
 f_{N,N}=0,\qquad
 f_{N,j+1}=f_{N,j}+a_{N,j}\log(1+1/j).
\]
Equivalently, putting $d_N(j)=q_N(j)-b_N\mathbf1_{N\mid j}$,
\[
 a_{N,j}-a_{N,j-1}=-d_N(j),\qquad j>N.
\]
The initial slope is $a_{N,N}=b_N$, and therefore
\begin{equation}
 F_N(t)=b_N\log(t/N)
          -\sum_{N<m\leq t}d_N(m)\log(t/m),\qquad t\geq N.
 \label{eq:ns67-divisor-response}
\end{equation}
For $h_j=\log(1+1/j)$ and
$I_k(h)=\int_0^h u^ke^{-u}\,du$, the complete cell energy is
\begin{equation}
 C_{N,j}=\int_j^{j+1}\frac{F_N(t)^2}{t^2}\,dt
 =\frac{f_{N,j}^2 I_0(h_j)+2f_{N,j}a_{N,j}I_1(h_j)
                       +a_{N,j}^2 I_2(h_j)}j.
 \label{eq:ns67-cell-energy}
\end{equation}
In particular $D_{N,1}^2=\sum_{j=N}^\infty C_{N,j}$.
Every signed divisor contribution and every cross term remains present.
\end{proposition}
\begin{proof}
The support assertion is the exact exterior interpolation in NS-62.
For $t\geq1$ the complete expression is
\[
 F_N(t)=\log t+\sum_{n\leq N}c_{N,n}A(t/n),
 \qquad A(z)=z-k\log z+\log(k!),\quad k=\lfloor z\rfloor.
\]
On a unit cell all floors are constant. The coefficient of $t$ is
$\sum c_{N,n}/n=0$, and the coefficient of $\log t$ is precisely
$a_{N,j}$. Continuity of $A$ at integers gives the endpoint recurrence.
The floor jumps give the divisor expression; at the initial endpoint
$\sum_{n\leq N}\mu(n)\lfloor N/n\rfloor=1$, so $a_{N,N}=b_N$.
The change $t=je^u$ proves~\eqref{eq:ns67-cell-energy}.
\end{proof}

\begin{lemma}[An elementary complete Stirling remainder]
\label{lem:ns67-remainder}
For every real $z\geq1$,
\begin{equation}
 \left|A(z)-\tfrac12\log(2\pi z)\right|\leq\frac1{6z}.
 \label{eq:ns67-remainder}
\end{equation}
\end{lemma}
\begin{proof}
Let $R(z)=A(z)-\tfrac12\log(2\pi z)$ and
$\widetilde B_2(z)=\{z\}^2-\{z\}+1/6$.
Then $R'(z)=(\{z\}-1/2)/z$ almost everywhere.
Stirling's formula gives $R(z)\to0$ as $z\to\infty$.
Since $\widetilde B_2$ is continuous at integers and its derivative
is $2(\{z\}-1/2)$ off the integers, integration by parts gives
\[
 R(z)=\frac{\widetilde B_2(z)}{2z}
             -\frac12\int_z^\infty\frac{\widetilde B_2(u)}{u^2}\,du.
\]
Both terms are controlled by $|\widetilde B_2|\leq1/6$, proving
\eqref{eq:ns67-remainder}. There is no unestimated periodic tail.
\end{proof}

Define the exact signed leading coefficients and the absolute remainder
mass by
\[
 A_N=1+\tfrac12\sum_{n\leq N}c_{N,n}=1-\tfrac12N\eta_N,
 \qquad B_N^\flat=\tfrac12\sum_{n\leq N}c_{N,n}\log\frac{2\pi}{n},
 \qquad L_N=\sum_{n\leq N}n|c_{N,n}|.
\]
The superscript distinguishes $B_N^\flat$ from the canonical function $B_N$.
Then for $t\geq N$,
\begin{equation}
 F_N(t)=A_N\log t+B_N^\flat+R_N^\flat(t),\qquad
 |R_N^\flat(t)|\leq\frac{L_N}{6t}.
 \label{eq:ns67-full-tail}
\end{equation}
For $T\geq N$ put $v_N(T)=A_N\log T+B_N^\flat$ and
\[
 P_N(T)=\frac{v_N(T)^2+2A_Nv_N(T)+2A_N^2}{T},\qquad
 r_N(T)=\frac{L_N}{6\sqrt{3T^3}}.
\]
The entire tail obeys
\begin{equation}
 \left|\int_T^\infty\frac{F_N(t)^2}{t^2}\,dt-P_N(T)\right|
       \leq2\sqrt{P_N(T)}\,r_N(T)+r_N(T)^2.
 \label{eq:ns67-tail-energy}
\end{equation}
This follows by integrating the squared linear logarithm exactly and
applying Cauchy--Schwarz to its cross term with the bounded remainder.
The right side includes the remainder square; no cancellation with it
is presumed.

\begin{proposition}[The polynomial cutoff retains the growth target]
\label{prop:ns67-polynomial-cutoff}
Let
\[
 Q_N=\left(\sum_{j=N}^{N^2-1}C_{N,j}\right)^{1/2},\qquad N\geq2.
\]
Unconditionally,
\begin{equation}
 0\leq D_{N,1}-Q_N\ll\log(2N).
 \label{eq:ns67-cutoff-error}
\end{equation}
Consequently
\[
 \lim_{N\to\infty}\frac{\log(1+Q_N)}{\log N}=\beta_*-\frac12,
\]
and RH is equivalent to $Q_N=O_\varepsilon(N^\varepsilon)$ for every
$\varepsilon>0$. This is an exact finite growing-interval reformulation,
not a proven bound for $Q_N$.
\end{proposition}
\begin{proof}
The elementary divisor identity gives
\[
 b_N=1+\sum_{n\leq N}\mu(n)\{N/n\},\qquad |b_N|\leq N+1.
\]
Hence $\sum|c_{N,n}|\leq2N+1$, $|A_N|\leq N+3/2$,
$|B_N^\flat|\ll N\log(2N)$, and $L_N\ll N^2$.
At $T=N^2$, therefore, $P_N(T)^{1/2}=O(\log(2N))$ and
$r_N(T)=O(1/N)$. Equation~\eqref{eq:ns67-tail-energy} bounds the
complete tail norm by $O(\log(2N))$. Disjoint-support orthogonality
gives $D_{N,1}^2=Q_N^2+\text{tail energy}$ and thus
\eqref{eq:ns67-cutoff-error}.
Apply the exponent theorem~\ref{cor:ns66-growth-exponent}. When
$\beta_*=1/2$, $0\leq Q_N\leq D_{N,1}$ gives exponent zero.
When $\beta_*>1/2$, subtracting an $O(\log N)$ error does not change
any positive power lower bound coming from a zero to the right of the
line. Take the supremum over those zeros and combine with the upper
bound $Q_N\leq D_{N,1}$.
\end{proof}

\paragraph{Complete boundary/integral comparison.}
Write $w_N=N\eta_N$ and
$I_N=\int_1^N M(u)\rho_u\,du/u$.
Then the exact Abel identity is $e_N=\tau+I_N-w_N\sigma_N$.
Once the full energy and the mixed integrals have been enclosed,
\[
 \|\tau+I_N\|^2=D_{N,1}^2+2w_N\langle e_N,\sigma_N\rangle
                          +\frac{w_N^2}{N}\|\sigma_1\|^2,
\]
\[
 \|I_N\|^2=\|\tau+I_N\|^2+2
       -2\{\langle e_N,\tau\rangle+w_N\langle\sigma_N,\tau\rangle\}.
\]
These are comparisons of the complete terms, not of their absolute
integrands. The implementation integrates both mixed terms on every
cell and retains their complete tails. The identities make the
cancellation measurable without discarding it; they do not bound it
uniformly in $N$.

\begin{proposition}[The two complete Abel terms have a uniform conditioning bound]
\label{prop:ns67-abel-conditioning}
Put $c_\sigma=\|\sigma_1\|$ and $k_N=k_{N,1}$ from
Proposition~\ref{prop:ns62-smoothed-obstruction}. With
$U_N=\tau+I_N$ and $V_N=w_N\sigma_N$, we have
\begin{equation}
 \|V_N\|\leq c_\sigma D_{N,1}+\frac{2c_\sigma}{\sqrt N},\qquad
 D_{N,1}\leq\|U_N\|+\|V_N\|
 \leq(1+2c_\sigma)D_{N,1}+\frac{4c_\sigma}{\sqrt N}.
 \label{eq:ns67-conditioning}
\end{equation}
Moreover define $\widetilde V_N=(w_N-k_N)\sigma_N$ and
$\widetilde U_N=U_N-k_N\sigma_N$. Then
$e_N=\widetilde U_N-\widetilde V_N$ and
\begin{equation}
 \|\widetilde V_N\|\leq c_\sigma D_{N,1},\qquad
 D_{N,1}\leq\|\widetilde U_N\|+\|\widetilde V_N\|
                \leq(1+2c_\sigma)D_{N,1}.
 \label{eq:ns67-centered-conditioning}
\end{equation}
Thus cancellation between these two complete vectors cannot change the
norm-growth exponent. Their internal arithmetic cancellation remains
unestimated. This does not justify replacing the integral $I_N$ by
its absolute integrand.
\end{proposition}
\begin{proof}
The localized unitary witness already proved in NS-62 gives
$|w_N-k_N|/\sqrt N\leq D_{N,1}$.
The defining integral of $k_N$ has absolute value at most $2$, since
$|\sin(2\pi x)/(\pi x)|\leq2$ and
$N\int_0^{1/N}\log(1/(Nx))\,dx=1$.
Use $\|\sigma_N\|=c_\sigma/\sqrt N$ and the triangle inequality.
The centered inequalities follow with no additive error; both centered
vectors differ by exactly $e_N$, so no cross term is omitted.
\end{proof}

\paragraph{Complete finite checks.}
The physical cell implementation evaluates $N=8,16,64,512,4096$ at
256 and 384 bits, with the entire tail beyond $T=1024N$ enclosed by
\eqref{eq:ns67-tail-energy}. It also retains the mixed tails with
$\sigma_N$ and $\tau$. Independent complete Gram evaluations at
$N=8,16,64$ lie inside all three corresponding physical enclosures.
Doubling the physical cutoff at $N=16,512$ gives overlapping full
energies and mixed terms. The replay verifies $175$ cross-precision
quantities, $28$ cutoff comparisons and $9$ independent Gram comparisons.
In all five cases the centered ratio in
\eqref{eq:ns67-centered-conditioning} lies below the certified
constant $1+2c_\sigma<4.617$. These are complete finite norms;
no bounded cofinal subsequence is inferred.

\paragraph{What changed in the proposed cancellation route.}
The large finite ratios from separately bounding the two original Abel
vectors can include a harmless $N^{-1/2}$ component. Removing its explicit
coefficient $k_N$ makes their comparison uniformly controlled by the
actual error. Therefore the proposed search for a power-saving hidden
solely in cancellation between those two \emph{complete} vectors does
not supply a weaker arithmetic target. The possible useful cancellation
is inside the integral, or inside the exact divisor-cell response
\eqref{eq:ns67-divisor-response}. No bound for either signed quantity
is proved here. This is an obstruction to that specific decomposition
argument, not a failure of the canonical subpolynomial-growth target.

\section{A positive renewal equation retains the arithmetic forcing (NS-68)}
\label{sec:ns68-renewal}

\paragraph{Status.}
The first multiplicative interval of the canonical response gives a
particularly simple signed arithmetic quantity. Its exact renewal
operator is positive and contractive in every strictly positive power
weight. This controls inversion of the operator but does not improve
the growth exponent of its forcing. The required arithmetic estimate
is still RH-equivalent. The argument below is a direct calculation;
no claim of priority for these reformulations is made.

For real $x>0$ put
\[
 M(x)=\sum_{n\leq x}\mu(n),\quad s(x)=\sum_{n\leq x}\frac{\mu(n)}n,
 \quad p(x)=x s(x)-M(x),
\]
where all three are zero for $x<1$. Then $p$ is continuous, including
at integers, and $p(x)=x\eta(x)$ for $x\geq1$.
Let $v=\log2$ and define
\begin{equation}
 f(x)=x s(x)v+\sum_{x<n\leq2x}\mu(n)\log\frac{2x}{n}.
 \label{eq:ns68-forcing}
\end{equation}
In particular $f(x)=0$ for $x<1/2$.

\begin{proposition}[Exact first-interval response and renewal identity]
\label{prop:ns68-renewal}
For each integer $N\geq1$,
\begin{equation}
 f(N)=F_N(2N)=e_N(1/(2N)).
 \label{eq:ns68-canonical-sample}
\end{equation}
For every real $x>0$,
\begin{equation}
 f(x)=v p(x)+\int_x^{2x}\frac{M(u)}u\,du
     =p(2x)-(1-v)p(x)-\int_x^{2x}\frac{p(u)}u\,du.
 \label{eq:ns68-renewal-identity}
\end{equation}
The coefficients in the corresponding average have total mass
$(1-v)+v=1$, and both are nonnegative.
\end{proposition}
\begin{proof}
On $N<m<2N$, every proper divisor of $m$ is at most $N$, so
$q_N(m)=-\mu(m)$. The endpoint $m=2N$ has zero logarithmic weight.
Equation~\eqref{eq:ns67-divisor-response} therefore gives
\eqref{eq:ns68-canonical-sample}. It also holds for $N=1$ directly.
Partial summation gives
\[
 \sum_{x<n\leq2x}\mu(n)\log\frac{2x}{n}
       =\int_x^{2x}\frac{M(u)-M(x)}u\,du,
\]
including integer endpoints, since the coefficient vanishes at $2x$.
This proves the first identity. The continuous, piecewise differentiable
function $p$ satisfies $u p'(u)-p(u)=M(u)$ almost everywhere. Integrate
$p'(u)-p(u)/u$ from $x$ to $2x$ to get the second identity. These
identities also apply when the interval crosses $1$.
\end{proof}

\begin{proposition}[Weighted causal inversion, with no power gain]
\label{prop:ns68-weighted-inversion}
For $\theta>0$ define
\[
 q_\theta=(1-v)2^{-\theta}+\frac{1-2^{-\theta}}\theta<1.
\]
Write $P_\theta(t)=e^{-\theta t}p(e^t)$ and
$G_\theta(t)=e^{-\theta t}f(e^{t-v})$, both zero for $t<0$.
Then
\begin{equation}
 P_\theta(t)=(1-v)2^{-\theta}P_\theta(t-v)
       +\int_0^v e^{-\theta u}P_\theta(t-u)\,du+G_\theta(t).
 \label{eq:ns68-causal}
\end{equation}
For every $1\leq a\leq\infty$ and every $T>0$,
\begin{equation}
 (1-q_\theta)\|P_\theta\|_{L^a(0,T)}
 \leq\|G_\theta\|_{L^a(0,T)}
 \leq(1+q_\theta)\|P_\theta\|_{L^a(0,T)}.
 \label{eq:ns68-two-sided}
\end{equation}
These bounds also apply on $(0,\infty)$ whenever either side is
finite. In particular, a bound $|f(x)|\leq Cx^\theta$ on $x>0$
implies
\begin{equation}
 |p(x)|\leq\frac{2^{-\theta}C}{1-q_\theta}\,x^\theta.
 \label{eq:ns68-sup-transfer}
\end{equation}
No smaller power follows from this renewal identity for general forcing.
\end{proposition}
\begin{proof}
Substitute $x=e^{t-v}$ in~\eqref{eq:ns68-renewal-identity}. The
convolution kernel in~\eqref{eq:ns68-causal} is the positive measure
\[
 K_\theta=(1-v)2^{-\theta}\delta_v
              +e^{-\theta u}\mathbf1_{(0,v)}(u)\,du.
\]
Its mass is $q_\theta$. At $\theta=0$ the mass is one; for
$\theta>0$ all its mass is strictly damped at positive delays, so
$q_\theta<1$. Extend a restriction to $(0,T)$ by zero. Causality
makes its convolution agree with the original convolution on $(0,T)$.
Young's inequality and the triangle inequality give both sides of
\eqref{eq:ns68-two-sided}. The finite-interval estimate followed by
$T\to\infty$ justifies the assertion without assuming an unknown
infinite norm is finite. For the supremum bound,
$|G_\theta|\leq2^{-\theta}C$.

For an explicit scope check, take the nonarithmetic continuous function
$p_\theta(x)=(x^\theta-1)\mathbf1_{x\geq1}$ and define its forcing
by the renewal identity. For $x\geq1$ that forcing is exactly
\[
 f_\theta(x)=\left\{2^\theta-(1-v)
                  -\frac{2^\theta-1}{\theta}\right\}x^\theta
            =2^\theta(1-q_\theta)x^\theta.
\]
The coefficient is positive. Thus the inverse preserves this power,
which excludes a generic power improvement by positivity alone. This
is a test of the operator class, not a counterexample concerning the
actual M\"obius coefficients.
\end{proof}

\begin{proposition}[The single sample retains every nontrivial-zero obstruction]
\label{prop:ns68-scalar-criterion}
The following two statements are equivalent to RH:
\[
 \forall\varepsilon>0,\quad p(x)=O_\varepsilon(x^{1/2+\varepsilon}),
 \qquad
 \forall\varepsilon>0,\quad f(N)=O_\varepsilon(N^{1/2+\varepsilon}).
\]
The second assertion is only over integer $N$. No such bound is
proved here.
\end{proposition}
\begin{proof}
RH implies $M(x)=O_\delta(x^{1/2+\delta})$. The complete tail
formula for $\eta$ in NS-66 then bounds $p$ at the same power.
Conversely the Mellin identity
\begin{equation}
 \int_1^\infty p(x)x^{-s-1}\,dx
             =\frac1{(s-1)s\zeta(s)},\qquad \Re s>1,
 \label{eq:ns68-p-transform}
\end{equation}
continues analytically to $\Re s>1/2$ if all the stated power bounds
hold. A zero there would give a genuine pole, a contradiction.
The apparent singularity at $s=1$ is removable. The functional equation
then gives RH.

The renewal identity and~\eqref{eq:ns68-sup-transfer} transfer these
power bounds in both directions between $p$ and $f$ for real $x$.
For the discrete-to-continuous step, the exact divisor identity gives
$|s(x)|\leq1+1/x$ for $x\geq1$. Off the discrete breakpoints,
\[
 f'(x)=v s(x)+\frac{M(2x)-M(x)}x,
 \qquad |f'(x)|\leq2v+2<4.
\]
The sums defining $f$ match continuously at integers and half-integers:
a term entering at $2x=n$ has coefficient zero, and a term leaving
at $x=n$ is exactly replaced by the change in $xs(x)v$.
Therefore $f$ is globally $4$-Lipschitz on $[1,\infty)$.
Its integer values suffice for every displayed power bound. Compact
initial intervals change only the constant.
\end{proof}

For an independent analytic check, termwise integration on $\Re s>1$
gives the exact Mellin response
\begin{equation}
 \int_0^\infty f(x)x^{-s-1}\,dx
  =\frac{H(s)}{\zeta(s)},\qquad
 H(s)=\frac{v}{s-1}+\frac{2^s-1-sv}{s^2}.
 \label{eq:ns68-forcing-transform}
\end{equation}
The factor $H$ has no zero on $\Re s>0$ (and has a pole at $1$).
Indeed the unweighted probability kernel has Laplace transform
\[
 \widehat K(s)=(1-v)2^{-s}+\frac{1-2^{-s}}s,
 \qquad H(s)=\frac{2^s\{1-\widehat K(s)\}}{s(s-1)}.
\]
For $\sigma=\Re s>0$, positivity of that kernel gives
$|\widehat K(s)|\leq q_\sigma<1$. Hence the numerator never
vanishes in this half-plane. The apparent pole of $H/\zeta$ at
$s=1$ is removable, as in~\eqref{eq:ns68-p-transform}.
Thus no nontrivial-zeta-zero contribution is removed by this scalar
sampling kernel. This is the same weighted contraction in Mellin
coordinates, not an additional arithmetic estimate.

\begin{corollary}[Critical weighted total energy cannot be finite]
\label{cor:ns68-critical-energy}
Put $q=q_{1/2}$. For every $X\geq1$,
\begin{equation}
 (1-q)^2\int_1^X\eta(x)^2\,dx
 \leq\frac12\int_{1/2}^{X/2}\frac{f(x)^2}{x^2}\,dx
 \leq(1+q)^2\int_1^X\eta(x)^2\,dx.
 \label{eq:ns68-finite-energy-comparison}
\end{equation}
Both integrals diverge as $X\to\infty$, unconditionally. Finiteness
of this particular critical weighted total energy is therefore not a
possible sufficient estimate to pursue. This does not exclude a
logarithmic or subpolynomial growth bound for these integrals, and is
not a claim that the canonical norm $D_{N,1}$ diverges.
\end{corollary}
\begin{proof}
Use~\eqref{eq:ns68-two-sided} with $a=2$, $\theta=1/2$ and
$T=\log X$, and change variables. Here
$\|P_{1/2}\|^2=\int_1^X\eta(x)^2\,dx$ and
$\|G_{1/2}\|^2=\tfrac12\int_{1/2}^{X/2}f(x)^2x^{-2}\,dx$.

If $\int_1^\infty\eta(x)^2dx$ were finite, Cauchy--Schwarz would
make~\eqref{eq:ns68-p-transform} analytic on $\Re s>1/2$ and give,
for $s=1/2+\delta+i\gamma$,
\[
 \left|\int_1^\infty p(x)x^{-s-1}dx\right|
       \leq\frac{\|\eta\|_2}{\sqrt{2\delta}}.
\]
At any known critical-line zero $1/2+i\gamma$ its meromorphic
continuation has a pole of integer order at least one. Along the
horizontal approach that grows at least as $\delta^{-1}$, contradicting
the displayed bound. The positive lower factor in
\eqref{eq:ns68-finite-energy-comparison} transfers divergence to the
forcing energy. No assumption on simplicity or RH is needed.
\end{proof}

\paragraph{Finite verification.}
Exact Maxima checks retain integer and noninteger endpoints, the physical
canonical sample, the Mellin response and both energy primitives.
Arb replays at 256 and 384 bits verify the complete renewal identity
and finite energy comparison for
$N=16,64,512,4096,65536,1048576$, using every M\"obius coefficient
through $2N$. There is no omitted tail in these finite identities.
The critical kernel mass and supremum-transfer constant satisfy
\[
 .8027641<q_{1/2}<.8027642,\qquad
 3.58508<\frac{2^{-1/2}}{1-q_{1/2}}<3.58509.
\]
The archive matches $74$ cross-precision quantities. No growth bound
is inferred from these six values.

\paragraph{Stop condition.}
The averaging geometry yields an explicit inverse bound, not a new
arithmetic estimate. A square-root-plus-epsilon estimate for its signed
forcing is exactly strong enough for RH and is still absent. Bounding
all critical weighted energy by a finite constant is actually false.
The remaining permissible objectives are growth estimates with an
explicit rate; neither positivity of the renewal kernel nor the finite
checks supply them.

\section{The optimized endpoint remainder is not a small correction on the tested blocks (NS-69)}
\label{sec:ns69-actual-remainder}

\paragraph{Scope.}
The coefficients in this section minimize the actual smoothed NB error.
They are not canonical M\"obius coefficients. We retain the complete
Gram, the old-space projection and the signed arithmetic remainder.
The finite certificates show a limitation of the two-moment model;
they neither prove an asymptotic obstruction to that model nor supply
the missing cofinal correlation estimate.

Use $R_N=\tau-\sum_{m\leq N}c_m\sigma_m$, $E_N=\|R_N\|^2$,
and $g_n=\langle R_N,\epsilon_n\rangle$, $N<n\leq2N$.
The $c_m$ solve the complete normal equations. Write
\[
 L_{N,n}=\frac{a_N\{\tfrac12\Delta_n(\log^2)
                  +(2-\gamma)\Delta_n\log\}
                   +\ell_N\Delta_n\log}{n},
 \qquad g_n=L_{N,n}+r_{N,n},
\]
with the exact $a_N,\ell_N$ from
\eqref{eq:ns61-residual-endpoint}. Thus $L$ denotes the endpoint
model and $r$ its signed error, rather than an independent substitute
for the actual correlation vector.

\begin{proposition}[Complete signed remainder and improved absolute allowance]
\label{prop:ns69-remainder}
With the complete arithmetic function $S_2$ from Ehm's Gram formula,
\begin{equation}
 r_{N,n}=-\sum_{m\leq N}\frac{c_m}{m}
       \left\{S_2(n/m)-\frac{n-1}{n}S_2((n-1)/m)\right\}.
 \label{eq:ns69-signed-remainder}
\end{equation}
For every $n>N$,
\begin{equation}
 |r_{N,n}|\leq
 \frac{\pi^2\sum_{m\leq N}m|c_m|}{144n^2(n-1)}.
 \label{eq:ns69-remainder-bound}
\end{equation}
This improves the earlier elementary allowance by a factor of $18$.
It does not assert that the allowance is small relative to $L$ or $g$.
\end{proposition}
\begin{proof}
Subtract the two complete Gram columns at $n,n-1$ with weights
$1,-(n-1)/n$. Their $K_2$ terms cancel exactly. Combining their
remaining polynomial terms with the exact target-load difference
gives $L_{N,n}$, while their $S_2$ terms give
\eqref{eq:ns69-signed-remainder}. This retains every old coefficient.
Alternatively, Lemma~\ref{lem:ns67-remainder} gives the physical
error bound
$|R_N(x)-a_N\log(1/x)-\ell_N|\leq x\sum m|c_m|/6$
on the support of $\epsilon_n$. Its nonnegativity and exact moment
$\int x\epsilon_n(x)dx=\pi^2/[24n^2(n-1)]$ give
\eqref{eq:ns69-remainder-bound}. Both derivations apply to the
complete functions, with no finite physical cutoff.
\end{proof}

Let $H_N$ be the complete projected Gram of the new differences, so
the full improvement is $\Delta E_N=g^T H_N^{-1}g=E_N-E_{2N}$.
For any nonzero real column $v$, restricting the new correction to
that column gives the exact improvement
\begin{equation}
 Q_N(v)=\frac{(v^Tg)^2}{v^TH_Nv}\leq\Delta E_N.
 \label{eq:ns69-selected-gain}
\end{equation}
For a two-column matrix $W$ it gives
$(W^Tg)^T(W^TH_NW)^{-1}(W^Tg)$. We use the fixed columns
\[
 W_{n,0}=\frac{N^2\Delta_n\log}{n},\qquad
 W_{n,1}=\frac{N^2\{\tfrac12\Delta_n(\log^2)
                                      -\log N\Delta_n\log\}}{n}.
\]
Their span contains $L$ exactly. They are independent for $N\geq2$:
the ratio of the second to the first is
$\tfrac12\log(n(n-1))-\log N$, strictly increasing in $n$.
The old projection is subtracted before evaluating every quotient.

\paragraph{Certified finite comparisons.}
Arb calculations at 256 and 384 bits use the complete Bernoulli/Hurwitz
tail enclosure of Section~\ref{sec:ns61-gram-remainder}. All solves
retain ball errors. The table gives outward decimal intervals; the
archived JSON files retain the narrower balls. The final two columns
are fractions of the full improvement, not relative errors.
\begin{center}\small
\begin{tabular}{rcccc}
\toprule
$N$ & $\|r\|_2/\|L\|_2$ & $\cos(L,g)$ & two-column fraction
 & $Q_N(g)/\Delta E_N$\\
\midrule
16 & $(.6074,.6076)$ & $(.8477,.8479)$ & $(.1204,.1205)$ & $(.2332,.2333)$\\
32 & $(.9250,.9252)$ & $(.7982,.7983)$ & $(.2036,.2037)$ & $(.3258,.3259)$\\
64 & $(2.0203,2.0205)$ & $(.5966,.5968)$ & $(.0806,.0807)$ & $(.2723,.2724)$\\
128 & $(3.1748,3.1750)$ & $(.4592,.4594)$ & $(.0330,.0331)$ & $(.2946,.2948)$\\
\bottomrule
\end{tabular}\end{center}
Here $\|\cdot\|_2$ is the Euclidean norm of the finite correlation
vector, while $E_N$ is the physical Hilbert-space squared error.
The exact vector identity
$\|g\|_2^2=\|L\|_2^2+2L^Tr+\|r\|_2^2$ is checked;
no mixed term is omitted. At $N=64$ the optimized logarithmic
coefficient satisfies
$-.000554<a_N<-.000553$, whereas it is positive at the other
three displayed sizes. Positivity of this coefficient is therefore
not available for all optimized blocks.

The smaller absolute allowance is still loose: its vector norm divided
by the actual remainder norm lies between $216$ and $541$ on these
four blocks. But the failure to treat the remainder as small is not
only a failure of that allowance: the actual remainder norm exceeds
the model norm at $N=64,128$. Conversely, the mixed term is positive
at those two sizes, so this is not evidence of universal destructive
cancellation. At the isolated far index $n=16N$, the relative scalar
remainders are below $.017$ in all four cases. This finite far-field
observation provides no uniform or cofinal estimate.

\paragraph{What had to change.}
A lower bound obtained by treating the leading endpoint terms as the
whole residual is unsupported. Even allowing both endpoint directions
captures only about $3.3\%$ of the full block improvement at $N=128$.
The actual-correlation direction captures about $29.5\%$ there, and
must retain its arithmetic remainder. These are finite comparisons,
not a convergence theorem. The open input is still an independently
proved, nonsummable relative contraction along a cofinal sequence.

\section{The elementary Gram background requires a structural correction (NS-70)}
\label{sec:ns70-background}

\paragraph{Scope.}
The endpoint model in NS-69 comes from the elementary part of Ehm's
complete $q=2$ Gram formula. Its scalar correction is small in some
entries, but that does not make it a small quadratic-form perturbation.
We first identify and repair a mismatch in diagonal regularity. Even
after that repair, a uniform comparison on all coefficient vectors is
excluded. None of these statements is a lower bound or obstruction
for the actual optimized residual direction.

For positive real $u,v$ let $G(u,v)=\langle\sigma_u,\sigma_v\rangle$.
With $K_1,K_2$ from Section~\ref{sec:ns61-gram-remainder}, define
\[
 K(u,v)=\frac{K_2+K_1|\log(u/v)|+\tfrac14\log^2(u/v)}{\max(u,v)},
 \qquad S(u,v)=G(u,v)-K(u,v).
\]
At integer indices this is the full arithmetic remainder of the Ehm
formula. Set $d=K_2-2K_1$. The constants are certified below; in
particular $d>0$ and $K_1>3/2$.

\begin{proposition}[A small entrywise correction cancels the diagonal cusp]
\label{prop:ns70-cusp}
Both $G$ and $K$ are positive definite on every finite collection of
distinct positive dilations. For the adjacent coefficient vector
corresponding to
$\epsilon_n=\sigma_n-(n-1)\sigma_{n-1}/n$, write $G[\epsilon_n]$
and $K[\epsilon_n]$ for its two quadratic forms. Then
\begin{equation}
 G[\epsilon_n]\sim\frac\kappa{n^3},\qquad
 K[\epsilon_n]\sim\frac d{n^2},\qquad
 \frac{S[\epsilon_n]}{K[\epsilon_n]}\longrightarrow-1.
 \label{eq:ns70-cusp-asymptotic}
\end{equation}
Consequently no fixed $c>0$ can satisfy $G\geq cK$ on all finite
integer coefficient vectors. In particular, no relative remainder
bound $|S[a]|\leq\delta K[a]$ with a fixed $\delta<1$ can hold
on that class. This follows from local regularity and does not use
an assumption about the location of zeta zeros.
\end{proposition}
\begin{proof}
Write $r=\log(v/u)\geq0$. Normalizing by $\sqrt{uv}$ turns $K$
into the translation kernel
\[
 \phi(r)=e^{-|r|/2}(K_2+K_1|r|+r^2/4).
\]
Its Fourier density, with
$\phi(r)=(2\pi)^{-1}\int\widehat\phi(t)e^{itr}dt$, is
\begin{equation}
 w_K(t)=\frac{d t^4+(K_2/2-3/2)t^2
                         +K_2/16+K_1/8+1/8}{(t^2+1/4)^3}>0.
 \label{eq:ns70-background-symbol}
\end{equation}
This is obtained by differentiating
$\int e^{-a|r|}e^{-itr}dr=2a/(a^2+t^2)$ twice in $a$ and then
putting $a=1/2$. The three numerator coefficients are positive by
the recorded constant enclosures. Hence $K$ is strictly positive
definite; $G$ is already a Gram of independent dilations.

For $h_n=\log(n/(n-1))$ direct evaluation gives
\begin{equation}
 n^2K[\epsilon_n]
 =K_2-2K_1(n-1)h_n-\tfrac12(n-1)h_n^2\longrightarrow d.
 \label{eq:ns70-exact-adjacent-background}
\end{equation}
NS-61 gives $n^3G[\epsilon_n]\to\kappa$ and its complete upper
bound. Thus the ratio $G[\epsilon_n]/K[\epsilon_n]$ tends to zero,
proving all assertions. The cusp is explicit:
$\phi'(0^+)=K_1-K_2/2=-d/2$, while the actual smoothed Gram
has the additional order of cancellation in
\eqref{eq:ns70-cusp-asymptotic}.
\end{proof}

\paragraph{The cusp correction is exactly diagonal in adjacent differences.}
For $B(u,v)=1/\max(u,v)$, direct evaluation gives
\[
 B[\epsilon_m,\epsilon_n]=\frac{\mathbf1_{m=n}}{n^2},
 \qquad B[\sigma_j,\epsilon_n]=0\quad(j<n).
\]
Indeed for $j<n$ the right difference is
$1/n-((n-1)/n)/(n-1)=0$, including $j=n-1$.
The diagonal difference equals $1/n^2$. Consequently the repair
subtracts exactly $d\,\operatorname{diag}(n^{-2})$ from the raw
new-difference Gram and leaves every elementary old/new cross column
unchanged. In particular it cannot improve the endpoint model $L$
in NS-69: the correction cancels from those correlations identically.
Its role is to repair the energy assigned to correction directions.
The old model Gram also changes, so a projected model must still
recompute its complete Schur complement.

\begin{proposition}[Repairing the cusp still leaves the zeta-weight comparison]
\label{prop:ns70-repaired-background}
The corrected elementary kernel
\begin{equation}
 K_0(u,v)=K(u,v)-\frac d{\max(u,v)}
 =\frac{2K_1+K_1|\log(u/v)|+\tfrac14\log^2(u/v)}{\max(u,v)}
 \label{eq:ns70-repair}
\end{equation}
is strictly positive definite and has no first-derivative cusp.
It satisfies
\[
 K_0[\epsilon_n]\sim\frac{K_1-1/2}{n^3},\qquad
 \frac{G[\epsilon_n]}{K_0[\epsilon_n]}
                \longrightarrow\frac\kappa{K_1-1/2}.
\]
Nevertheless there are no constants $c>0$ or $C<\infty$ such that,
on all finite integer coefficient vectors, respectively
\[
 G\geq cK_0\qquad\text{or}\qquad G\leq CK_0.
\]
Both failures are unconditional. They exclude uniform all-coefficient
comparisons, not selected-direction estimates, finite-block inequalities
with size-dependent constants, or the optimized NB criterion.
\end{proposition}
\begin{proof}
The Fourier density of the repaired kernel is
\begin{equation}
 w_0(t)=\frac{(K_1-3/2)t^2+K_1/4+1/8}{(t^2+1/4)^3}>0.
 \label{eq:ns70-repaired-symbol}
\end{equation}
This proves strict positivity and the cancellation of the cusp.
In~\eqref{eq:ns70-exact-adjacent-background} replace $K_2$ by
$2K_1$ and use
$(n-1)h_n=1-1/(2(n-1))+O(n^{-2})$ and
$(n-1)h_n^2=1/(n-1)+O(n^{-2})$.
This gives the displayed adjacent asymptotic.

The complete Gram has Mellin density
\[
 w_G(t)=\frac{|\zeta(1/2+it)|^2}{(t^2+1/4)^2}.
\]
For clarity, a uniform quadratic-form inequality on all finite integer
coefficients would force the corresponding pointwise inequality between
$w_G$ and $w_0$. Here is the full localization argument. Both densities
are continuous, even and integrable. Their homogeneous kernels extend
to all positive real dilation parameters. Replace any finite set of real
parameters $u_j$ by integers nearest to $D u_j$ and multiply both forms
by $D$; kernel continuity lets $D\to\infty$. Thus an integer inequality
extends to all real parameters. Riemann sums then extend it to compactly
supported real coefficient densities in $r=\log u$. Choose long real
cosine packets in $r$. Their squared Fourier transforms form approximate
identities at the paired frequencies $\pm t_0$. Continuity near the
frequencies and integrability away from them give the pointwise inequality.
This reasoning retains all quadratic cross terms and uses arbitrary
finite real coefficients, not an arithmetic coefficient rule.

A known critical-line zero makes $w_G(t_0)=0$ while $w_0(t_0)>0$,
so the lower comparison is impossible. For the upper comparison,
\[
 \frac{w_G(t)}{w_0(t)}
 =|\zeta(1/2+it)|^2\,
   \frac{t^2+1/4}{(K_1-3/2)t^2+K_1/4+1/8}.
\]
The rational factor tends to the positive constant $(K_1-3/2)^{-1}$.
Known unbounded critical-line values, for example
\cite[Theorem 1]{Soundararajan2008}, therefore exclude any finite $C$.
The lower failure uses zeros on the line, not hypothetical off-line zeros.
Neither failure says that the actual Gram, or an original Weil form, has
negative energy.
\end{proof}

\paragraph{Complete constant and adjacent-vector certificates.}
Arb at 256 and 384 bits encloses every numerator coefficient in
\eqref{eq:ns70-background-symbol} and
\eqref{eq:ns70-repaired-symbol} strictly above zero, with
$.0091604<d<.0091605$. Complete two-vector Gram comparisons give
\begin{center}\small
\begin{tabular}{rcc}
\toprule
$n$ & $G[\epsilon_n]/K[\epsilon_n]$
    & $G[\epsilon_n]/K_0[\epsilon_n]$\\
\midrule
256 & $(.36207,.36209)$ & $(1.11185,1.11187)$\\
1024 & $(.11987,.11989)$ & $(1.11423,1.11425)$\\
4096 & $(.03260,.03262)$ & $(1.11498,1.11500)$\\
\bottomrule
\end{tabular}\end{center}
The true quadratic energy stays positive. The Bernoulli/Hurwitz
remainder is retained; a doubled series cutoff gives overlapping
values for all $34$ constants and quantities. The asymptotic assertions
follow from the proofs, not from this finite table.

\paragraph{What had to change.}
The elementary background had the wrong local regularity; removing
$d/\max(u,v)$ repairs that mismatch and recovers the correct adjacent
power $n^{-3}$. This repair is real but insufficient: the complete
arithmetic factor is still neither uniformly above nor uniformly below
the repaired background. A useful next bound must keep the particular
residual direction or supply an additional arithmetic estimate. The
finite endpoint remainder in NS-69 is not silently discarded.

\section{Using the repaired kernel to select a direction (NS-71)}
\label{sec:ns71-selected-preconditioner}

\paragraph{Question and scope.}
NS-70 excludes a uniform comparison over every coefficient vector.
It does not exclude using the repaired kernel to choose one direction,
then measuring that direction with the true complete Gram. The finite
prediction tested here is that such a choice captures more of the
optimized improvement than the unmodified correlation vector. A good
finite direction is a computational finding, not a cofinal estimate.

On the new block $N<n\leq2N$, let $D_N$ and $H_N$ denote the raw
and old-space-projected actual difference Grams. Let $D_{0,N}$ and
$H_{0,N}$ be the corresponding Grams for the repaired kernel $K_0$.
All four matrices are strictly positive definite. In defining $H_{0,N}$,
both the old projection and the new block use $K_0$ consistently;
this is not the actual Hilbert-space projection.
For the actual optimized correlation vector $g$, consider
\[
 v_{\rm raw}=D_{0,N}^{-1}g,\qquad
 v_{\rm projected}=H_{0,N}^{-1}g.
\]
Both are only rules for selecting coefficients. Every reported energy
uses $H_N$ or $D_N$, with the complete arithmetic remainder retained.

\begin{proposition}[Selected-direction improvement and the precise open input]
\label{prop:ns71-selected-input}
For either positive model matrix $A_N\in\{D_{0,N},H_{0,N}\}$ put
\[
 v=A_N^{-1}g,\quad m_N=g^TA_N^{-1}g,\quad
 C_N=\frac{v^TH_Nv}{v^TA_Nv},\quad
 C_N^{\rm raw}=\frac{v^TD_Nv}{v^TA_Nv}.
\]
Whenever $g\ne0$, all denominators are positive, and
\begin{equation}
 E_N-E_{2N}\geq Q_N(v)=\frac{m_N}{C_N}
                  \geq\frac{m_N}{C_N^{\rm raw}}>0.
 \label{eq:ns71-selected-ratio}
\end{equation}
If along $N_j=2^jN_0$ independently justified estimates give
\[
 \frac{m_{N_j}}{E_{N_j}}\geq\alpha_j>0,
 \qquad C_{N_j}\leq B_j,
 \qquad \sum_j\frac{\alpha_j}{B_j}=\infty,
\]
then $E_{N_j}\to0$ and RH follows by the fixed smoothed NB criterion.
The same statement holds using $C_{N_j}^{\rm raw}$ in place of $C_{N_j}$.
No such infinite sequence of estimates is proved here.
\end{proposition}
\begin{proof}
Since $A_Nv=g$, $v^Tg=v^TA_Nv=m_N$. Substitute in
\eqref{eq:ns69-selected-gain}. Orthogonal projection gives
$H_N\leq D_N$, proving the second inequality with all old-space
cross terms accounted for. The assumed bounds yield
$E_{2N_j}\leq(1-\alpha_j/B_j)E_{N_j}$.
The exact gain ensures $0<\alpha_j/B_j\leq1$; if equality occurs,
the error vanishes already. Otherwise the divergent sum makes the
product of contraction factors tend to zero. Monotonicity extends
convergence from the dyadic sequence to all $N$. The criterion in
Proposition~\ref{prop:ns61-rh-equivalence} then applies.
This uses a selected scalar ratio, not a uniform matrix comparison.
\end{proof}

For its geometric meaning, let $u=H_N^{-1/2}g$ and
$z=H_N^{1/2}v$. Then the exact fraction of the full improvement is
\[
 \frac{Q_N(v)}{E_N-E_{2N}}=
       \frac{(u^Tz)^2}{\|u\|_2^2\|z\|_2^2}.
\]
Thus the finite experiment measures an angle in the actual projected
energy, including the arithmetic terms. A Euclidean angle or a small
entrywise kernel remainder would not be a substitute.

\paragraph{Certified finite outcome.}
The table gives outward intervals for the fraction of the full
improvement $E_N-E_{2N}$, using the actual projected energy throughout.
\begin{center}\small
\begin{tabular}{rccc}
\toprule
$N$ & $v=g$ & $v=D_{0,N}^{-1}g$ & $v=H_{0,N}^{-1}g$\\
\midrule
16 & $(.2332,.2333)$ & $(.7167,.7169)$ & $(.8390,.8391)$\\
32 & $(.3258,.3259)$ & $(.8580,.8581)$ & $(.8252,.8254)$\\
64 & $(.2723,.2724)$ & $(.9593,.9594)$ & $(.9577,.9578)$\\
128 & $(.2946,.2948)$ & $(.9729,.9730)$ & $(.9781,.9782)$\\
256 & $(.3840,.3842)$ & $(.9597,.9598)$ & $(.9705,.9707)$\\
\bottomrule
\end{tabular}\end{center}
Both rules outperform $v=g$ at every checked size. The projected
rule captures more than $97\%$ of the full improvement at $N=128,256$;
neither efficiency is asserted monotone. At $N=256$ it reduces the
squared residual by a fraction in $(.08081,.08083)$, compared with
the full optimum in $(.08326,.08328)$. The corresponding numerator
and selected ratio in Proposition~\ref{prop:ns71-selected-input}
satisfy
\[
 .7668< m_{256}/E_{256}<.7669,\qquad
 9.4884<C_{256}<9.4886.
\]
These two finite bounds make explicit where the remaining arithmetic
problem sits; no size-independent bound is inferred.

All five sizes were evaluated at 256 and 384 bits with the full
Bernoulli/Hurwitz Gram tail. A doubled series cutoff at $N=256$
replays the largest problem. The verifier matches $115$ quantities
between precisions, $23$ at the doubled cutoff, and $16$ overlapping
quantities against the separate NS-69 implementation. The minimum
recorded accuracy is $88$ bits. All model and actual solves retain
their ball enclosures. The Maxima checks concern the exact quotient
and projection algebra, not the missing arithmetic estimates.

\paragraph{Stop condition.}
The model inversion is explicit and numerically useful on the checked
blocks. The arithmetic numerator $m_N/E_N$ and the selected true-energy
ratio still need cofinal estimates. Finite efficiency, by itself,
does not control either quantity as $N\to\infty$. NS-70's obstruction
to all-coefficient comparisons is respected throughout; no new
zero-free region, uniform Weil floor, G2 or RH proof follows.

\section{Local differential geometry of the repaired model (NS-72)}
\label{sec:ns72-local-geometry}

\paragraph{Scope.}
The repaired kernel in NS-70 has a positive local differential-energy
representation. This makes its selected numerator uniformly comparable
to an explicit discrete curvature energy, and supplies two directions
that can be selected in a linear number of scalar operations once the
actual correlation vector is known. These are statements about the
model geometry. Every arithmetic gain below is measured with the
complete actual Gram. No cofinal arithmetic estimate is supplied.

Put
\[
 a_0=K_1-\tfrac32>0,\qquad b_0=K_1/4+\tfrac18=a_0/4+\tfrac12,
 \qquad \alpha_0=\sqrt{b_0/a_0}>\tfrac12.
\]
The normalized stationary kernel is
\[
 \phi_0(r)=e^{-|r|/2}(2K_1+K_1|r|+r^2/4),\qquad
 K_0(u,v)=\frac{\phi_0(\log u-\log v)}{\sqrt{uv}}.
\]
Its complete Fourier density is
\begin{equation}
 w_0(t)=\frac{a_0t^2+b_0}{(t^2+1/4)^3}
       =\frac{a_0}{(t^2+1/4)^2}
                       +\frac{1/2}{(t^2+1/4)^3}.
 \label{eq:ns72-positive-split}
\end{equation}

\begin{proposition}[A causal factor and a local positive energy]
\label{prop:ns72-local-energy}
With Fourier convention $\widehat f(t)=\int f(r)e^{-itr}dr$, define
\[
 h_0(r)=\sqrt{a_0}\,e^{-r/2}
       \left\{r+\frac{\alpha_0-1/2}{2}r^2\right\}
       \mathbf1_{r\geq0}.
\]
Then $h_0\geq0$ and
\[
 \widehat h_0(t)=\frac{\sqrt{a_0}(\alpha_0+it)}{(1/2+it)^3},
 \qquad \phi_0(r)=\int_{\mathbb R}h_0(u+r)h_0(u)\,du.
\]
Let $\mathcal H_0$ be the real Hilbert space with norm
\[
 \|f\|_{\mathcal H_0}^2=
       \frac1{2\pi}\int_{\mathbb R}\frac{|\widehat f(t)|^2}{w_0(t)}\,dt.
\]
As a set it is exactly $H^2(\mathbb R)$, and $\phi_0(r-\cdot)$
is its point-evaluation representer. If $z\in H^3(\mathbb R)$ is the
unique solution of $z'+\alpha_0z=f$, then
\begin{equation}
 \|f\|_{\mathcal H_0}^2
 =\frac1{a_0}\left\{
 \|z'''\|_2^2+\frac34\|z''\|_2^2
               +\frac3{16}\|z'\|_2^2+\frac1{64}\|z\|_2^2\right\}.
 \label{eq:ns72-local-functional}
\end{equation}
In particular, with $L_0=-d^2/dr^2+1/4$,
\begin{equation}
 \frac1{4b_0}\|L_0f\|_2^2\leq\|f\|_{\mathcal H_0}^2
                 \leq\frac1{a_0}\|L_0f\|_2^2.
 \label{eq:ns72-sobolev-comparison}
\end{equation}
All these are whole-line statements; no boundary or remote tail is omitted.
\end{proposition}
\begin{proof}
The elementary transforms of $r e^{-r/2}$ and $r^2e^{-r/2}/2$ on
$r\geq0$ give the displayed causal factor. Its squared modulus is
exactly~\eqref{eq:ns72-positive-split}, proving the autocorrelation
identity. In particular this factorization includes the complete model,
not just its large-frequency term.

The positive rational weight $1/w_0$ is comparable to $(1+t^2)^2$.
Point evaluation is continuous because $\int w_0(t)dt<\infty$;
Fourier inversion gives its stated representer. The multiplier
$(\alpha_0+it)^{-1}$ is an isomorphism from $H^2$ to $H^3$.
It is also the causal convolution with
$e^{-\alpha_0r}\mathbf1_{r\geq0}$; thus there is no unspecified
homogeneous solution. Plancherel gives
$\|f\|_{\mathcal H_0}^2=a_0^{-1}\|(d/dr+1/2)^3z\|_2^2$.
Expanding $(t^2+1/4)^3$ gives~\eqref{eq:ns72-local-functional}.
Equivalently one may use density of compact smooth functions in $H^3$;
no integration-by-parts boundary term survives that limit.
Finally
\[
 \frac1{w_0(t)}=(t^2+1/4)^2\frac{t^2+1/4}{a_0t^2+b_0},
 \qquad \frac1{4b_0}\leq\frac{t^2+1/4}{a_0t^2+b_0}\leq\frac1{a_0},
\]
which proves~\eqref{eq:ns72-sobolev-comparison}.
\end{proof}

\paragraph{The corresponding differential equation.}
For prescribed values $f(r_j)=y_j$, the minimum-norm interpolant is
$f=\sum_j\lambda_j\phi_0(r-r_j)$. In terms of its $z$, it satisfies
in distributions
\begin{equation}
 \frac1{a_0}L_0^3z
     =\sum_j\lambda_j\{\alpha_0\delta_{r_j}-\delta'_{r_j}\}.
 \label{eq:ns72-local-equation}
\end{equation}
Thus between sample points $L_0^3z=0$: its pieces are combinations
of $e^{r/2}$ and $e^{-r/2}$ times quadratic polynomials.
The whole-line Sobolev condition fixes the decaying tails.
The derivatives through order three are continuous, while the jumps are
$[z^{(4)}]_{r_j}=a_0\lambda_j$ and
$[z^{(5)}]_{r_j}=-a_0\alpha_0\lambda_j$.
These signs follow directly from~\eqref{eq:ns72-local-equation}.
This is a positive local model problem with explicit matching conditions;
its coefficients contain no zeta-zero information.

\begin{proposition}[The actual selected numerator is a model interpolation energy]
\label{prop:ns72-interpolation}
Fix $N\geq2$ and the actual optimized correlations $g_n$,
$N<n\leq2N$, from NS-71. Put $r_j=\log j$ and define
\begin{equation}
 y_j=0\quad(j\leq N),\qquad
 y_j=\frac1{\sqrt j}\sum_{k=N+1}^j k g_k\quad(N<j\leq2N).
 \label{eq:ns72-data}
\end{equation}
Then the numerator for the projected repaired model is exactly
\begin{equation}
 m_N=g^TH_{0,N}^{-1}g
       =\min\{\|f\|_{\mathcal H_0}^2:
                          f(r_j)=y_j,\ 1\leq j\leq2N\}.
 \label{eq:ns72-interpolation-numerator}
\end{equation}
This equality holds for every real vector $g$, not just the arithmetic one.
For the actual vector, $y_j=\sqrt j\langle R_N,\sigma_j\rangle$.
\end{proposition}
\begin{proof}
Write $q_j=\langle R_N,\sigma_j\rangle$.
Old-space orthogonality gives $q_j=0$ for $j\leq N$.
The exact difference relation is
$g_j=q_j-(j-1)q_{j-1}/j$; telescoping gives
$q_j=j^{-1}\sum_{k=N+1}^jkg_k$, with no endpoint error.

Let $\Phi_{ij}=\phi_0(r_i-r_j)$ and
$\mathsf D=\operatorname{diag}(\sqrt1,\ldots,\sqrt{2N})$.
The full model Gram is $\mathsf D^{-1}\Phi\mathsf D^{-1}$.
The invertible change from old atoms and new atoms to old atoms and
new differences sends the load data to $(0,g)$.
The complete block-inverse identity therefore gives
$y^T\Phi^{-1}y=g^TH_{0,N}^{-1}g$.
For completeness, the minimum-norm interpolant lies in the span of
the evaluation representers: subtracting its component orthogonal to
that span preserves all values and lowers its squared norm.
Its coefficients are $\Phi^{-1}y$, proving the remaining equality.
Every old/new cross term is retained in this argument.
\end{proof}

For arbitrary data~\eqref{eq:ns72-data}, write
\[
 h_j=r_{j+1}-r_j,\qquad
 d_i=\frac{y_{i+1}-y_i}{h_i}-\frac{y_i-y_{i-1}}{h_{i-1}}
       \quad(2\leq i\leq2N-1).
\]
Let $T$ be the positive tridiagonal matrix indexed by those interior
points, with
\[
 T_{ii}=\frac{h_{i-1}+h_i}{3},\qquad T_{i,i+1}=\frac{h_i}{6}.
\]
Define the complete bending energy and the local curvature sum by
\[
 \mathcal C_N=d^TT^{-1}d,\qquad
 \mathcal W_N=\sum_{i=2}^{2N-1}\frac{d_i^2}{h_{i-1}+h_i}.
\]
They vanish only when $g=0$: zero curvature would make all data affine
in $r$, and the two old zero values force that affine function to vanish.

\begin{proposition}[Uniform comparison with discrete curvature, within the model]
\label{prop:ns72-curvature}
For every $N\geq2$ and every real $g$,
\begin{equation}
 2\mathcal W_N\leq\mathcal C_N\leq6\mathcal W_N,
 \qquad
 \frac{\mathcal C_N}{4b_0}\leq m_N
                         \leq\frac{84}{a_0}\mathcal C_N.
 \label{eq:ns72-curvature-comparison}
\end{equation}
The constants are independent of $N$; the upper constant is deliberately
not optimized. These are uniform comparisons for the model numerator,
not comparisons between the model Gram and the actual arithmetic Gram.
\end{proposition}
\begin{proof}
Let $\psi_i$ be the standard tent function, equal to one at $r_i$
and zero at $r_{i-1},r_{i+1}$, with support between those two endpoints.
Integration by parts on each interval gives
$d_i=\int f''\psi_i$ for any $H^2$ interpolant. The complete Gram of
the tents is exactly $T$. Minimizing $\|f''\|_{L^2(r_1,r_{2N})}^2$
therefore gives $d^TT^{-1}d$. Conversely a minimizing second derivative
in the tent span integrates, with its affine part chosen to match the
first two data, to the natural cubic spline through every sample.
Thus this is an equality for the finite interval, not merely an
uncertified spline approximation. Moreover the off-diagonal row sum
of $T$ is at most half its diagonal. Applying
$2|uv|\leq u^2+v^2$ to its cross terms shows
\[
 \tfrac16\operatorname{diag}(h_{i-1}+h_i)
 \leq T\leq\tfrac12\operatorname{diag}(h_{i-1}+h_i),
\]
and inversion proves the first pair of bounds.
The left side of~\eqref{eq:ns72-sobolev-comparison} bounds every
whole-line interpolant below by $\mathcal C_N/(4b_0)$.

For the uniform upper bound let $F$ be the natural spline on
$[r_1,r_{2N}]$, and set $E=\|F''\|_2^2=\mathcal C_N$ there.
All old nodal values vanish. On an old interval of length $h\leq\log2<1$,
$F'$ has mean zero. Elementary one-dimensional estimates give
$\|F'\|_2^2\leq h^2\|F''\|_2^2$ and
$\|F\|_2^2\leq h^4\|F''\|_2^2$ on that interval.
The derivatives at $r_1$ and $r_N$ have squared magnitudes at most
$E/3$: integrate $F''$ against the endpoint linear weight on an
adjacent zero-data interval. The new interval $[r_N,r_{2N}]$ has
length $\log2<1$, starts at value zero, and hence
\[
 \|F'\|_{L^2(r_N,r_{2N})}^2,\quad
 \|F\|_{L^2(r_N,r_{2N})}^2,\quad
 |F(r_{2N})|^2,\quad |F'(r_{2N})|^2\ \leq\frac83 E.
\]
Combining old and new intervals yields
$\|F'\|_2^2,\|F\|_2^2\leq11E/3$ on the full finite interval.
For any interval $I$ define the positive local functional
\[
 \mathcal J_I(F)=\int_I\left(|F''|^2+\tfrac12|F'|^2
                                      +\tfrac1{16}|F|^2\right)dr.
\]
Consequently $\mathcal J_{[r_1,r_{2N}]}(F)\leq49E/16<4E$.
Only after summing over the complete extension below do we use
$\mathcal J_{\mathbb R}(F)=\|L_0F\|_2^2$. On a finite piece the
two expressions can differ by an endpoint term; none is discarded.

Extend to the left over an interval of length one by the cubic with
zero value and derivative at its outer endpoint and value zero,
derivative $F'(r_1)$ at $r_1$. Its local $\mathcal J$ energy is
$(6833/1680)|F'(r_1)|^2<2E$.
On the right use $y(1-3t^2+2t^3)+w(t-2t^2+t^3)$,
$0\leq t\leq1$, where $y=F(r_{2N})$, $w=F'(r_{2N})$,
and then extend by zero. Its local $\mathcal J$ energy matrix has entries
\[
 A_{00}=7069/560,\quad A_{01}=20339/3360,\quad A_{11}=6833/1680.
\]
They are positive and $A_{00}+2A_{01}+A_{11}<29$.
Since $y^2,w^2,|yw|\leq8E/3$, the right extension costs at most
$232E/3$. The resulting compactly supported function is $C^1$
and belongs to $H^2(\mathbb R)$; jumps in its second derivative
are allowed. Its total $L_0$ energy is less than
$(4+2+232/3)E<84E$.
The upper side of~\eqref{eq:ns72-sobolev-comparison} proves the result.
The extensions explicitly control both tails and all joining conditions.
\end{proof}

\paragraph{Two local coefficient rules.}
Let $\mathsf S_N$ be the exact linear map $g\mapsto d$ defined above,
and put $D_h=\operatorname{diag}(h_{i-1}+h_i)$. The choices
\begin{equation}
 v_{\rm spline}=\mathsf S_N^TT^{-1}\mathsf S_Ng,
 \qquad v_{\rm diagonal}=\mathsf S_N^TD_h^{-1}\mathsf S_Ng
 \label{eq:ns72-local-directions}
\end{equation}
have exact numerators $g^Tv_{\rm spline}=\mathcal C_N$ and
$g^Tv_{\rm diagonal}=\mathcal W_N$.
Both directions are nonzero for $g\ne0$.
Cumulative sums construct $y$, local differences construct $d$, a
positive-pivot tridiagonal solve computes $T^{-1}d$, and weighted suffix
sums apply the transpose. Thus direction selection takes $O(N)$ scalar
operations after the actual $g$ has been supplied. Computing that
arithmetic vector and its true energy is a separate task.

For either direction the actual gain remains
\[
 Q_N(v)=\frac{(g^Tv)^2}{v^TH_Nv}.
\]
A cofinal lower bound for $\mathcal C_N/E_N$ (or $\mathcal W_N/E_N$)
and an upper bound for the corresponding selected true cost, giving a
nonsummable relative contraction, would suffice for RH by the same
argument as NS-71. The geometric equivalence does not establish these
arithmetic bounds. In particular no inequality in
\eqref{eq:ns72-curvature-comparison} includes the actual Gram $H_N$.

\paragraph{Certified finite outcome.}
The following outward intervals enclose each direction's fraction of the
full actual projected gain. Direction selection uses the displayed rule
without fitting its coefficients to the actual Gram.
\[
\begin{array}{c|ccc}
 N&\text{repaired model}&\text{spline curvature}&\text{diagonal curvature}\\ \hline
16&(.8390,.8391)&(.8529,.8530)&(.9308,.9310)\\
32&(.8252,.8254)&(.8334,.8335)&(.9040,.9041)\\
64&(.9577,.9578)&(.9602,.9603)&(.9783,.9785)\\
128&(.9781,.9782)&(.9795,.9796)&(.9852,.9854)\\
256&(.9705,.9707)&(.9725,.9727)&(.9763,.9764)
\end{array}
\]
The simpler diagonal rule outperforms the repaired-kernel solve at all
five checked sizes. This does not assert monotonicity or asymptotic
optimality. The uniform numerator minorant $\mathcal C_N/(4b_0)$ captures
between $6.19\%$ and $6.26\%$ of the exact model numerator on these blocks;
its constant is conservative. Every complete actual Gram, tail and
old-space correction is retained.

Arb runs at 256 and 384 bits match all 140 recorded scalar enclosures.
Doubling the Gram-series cutoff at $N=256$ matches another 28 quantities;
25 overlapping quantities match the published NS-71 calculation.
The minimum recorded accuracy is 73 bits. At $N=16$ an independent
dense solve also encloses all entries of the tridiagonal solution.
The 41 exact Maxima checks include the causal factor, actual Green-function
jumps, normalization, tent Grams and the full extension energies.
No finite computation supplies the missing cofinal arithmetic bound.

\section{Testing explicit arithmetic coefficient templates (NS-73)}
\label{sec:ns73-mobius-templates}

\paragraph{Question and scope.}
The successful curvature rule depends on the actual optimized residual.
We test whether three prescribed M\"obius templates can replace that
information on the same finite blocks. This is a finite control, not an
asymptotic exclusion of M\"obius coefficients. Logarithmic tapers have
established prior motivation: Bettin--Conrey--Farmer's Theorem~1
in~\cite{BettinConreyFarmer2013} proves an optimal $q=1$ asymptotic under
RH and an additional inverse-zeta-derivative moment hypothesis. Neither
hypothesis is used here. Our $q=2$ projected block test is a different
problem from their full canonical approximant.

For $p=0,1,2$, prescribe physical new-atom coefficients
\[
 a^{(p)}_n=-\mu(n)\log^p(2N/n),\qquad N<n\leq2N,
\]
where $p=0$ means the constant taper, including at $n=2N$.
No coefficient is fitted when selecting these three individual templates.
Define their adjacent-difference coordinates by
\begin{equation}
 v^{(p)}_n=n\sum_{k=n}^{2N}\frac{a^{(p)}_k}{k}.
 \label{eq:ns73-coordinate-change}
\end{equation}
Direct backward telescoping gives
\[
 \sum_{n=N+1}^{2N}v^{(p)}_n\varepsilon_n
 =\sum_{n=N+1}^{2N}a^{(p)}_n\sigma_n
                  -\frac{N}{N+1}v^{(p)}_{N+1}\sigma_N.
\]
Thus applying the complete old-space orthogonal complement gives exactly
the intended physical new direction; the boundary term lies in the old
space and is removed, not silently dropped before projection.

\begin{proposition}[Complete finite template and span gains]
\label{prop:ns73-span}
For any real new-atom templates converted by
\eqref{eq:ns73-coordinate-change}, let $V$ collect their difference
coordinate vectors, $\Gamma=V^TH_NV$, and $\ell=V^Tg$.
If the projected templates are independent, then $\Gamma>0$ and
\[
 \max_{v\in\operatorname{ran}V\setminus\{0\}}
       \frac{(g^Tv)^2}{v^TH_Nv}=\ell^T\Gamma^{-1}\ell.
\]
For any nonzero direction $w$, the fraction of its actual projected
energy lying in that span is
\[
 \frac{(V^TH_Nw)^T\Gamma^{-1}(V^TH_Nw)}{w^TH_Nw}\in[0,1].
\]
These formulae keep every cross term and the complete old-space correction.
The maximizing span coefficients are fitted using the actual Gram; this
span control is stronger than any fixed combination of the three templates.
\end{proposition}
\begin{proof}
The coordinate relation follows from
$v_n-nv_{n+1}/(n+1)=a_n$, with $v_{2N+1}=0$.
The remaining claims are finite-dimensional orthogonal projection in the
positive $H_N$ metric. Equivalently apply Cauchy--Schwarz to
$\Gamma^{1/2}b$ and $\Gamma^{-1/2}\ell$. This also covers $\ell=0$,
when every gain is zero. The second expression is the usual squared norm
of the metric projection, divided by the original squared norm.
\end{proof}

\paragraph{Certified finite comparison.}
The outward intervals below enclose fractions of the same full arithmetic
block gain. The last column permits the best actual-Gram combination of
all three M\"obius templates, not merely their individual scalings.
\[
\begin{array}{c|cccc}
N& p=0&p=1&p=2&\text{three-template span}\\ \hline
16&(.2134,.2135)&(.4692,.4694)&(.1828,.1829)&(.6023,.6024)\\
32&(.0000683,.0000684)&(.6743,.6745)&(.5652,.5654)&(.7627,.7628)\\
64&(.9222,.9223)&(.6407,.6409)&(.3966,.3967)&(.9512,.9514)\\
128&(.2942,.2943)&(.4138,.4139)&(.2908,.2909)&(.8903,.8904)\\
256&(.1023,.1025)&(.3390,.3392)&(.4338,.4339)&(.8366,.8367)
\end{array}
\]
The curvature rule exceeds even this fitted span at all five sizes.
Thus the three proposed templates do not explain its finite performance.
The comparison does not rule out other arithmetic templates, larger
families, or any asymptotic M\"obius approximation theorem.
Two-precision and doubled-cutoff enclosures, the positive restricted Gram
and all cross terms are retained in the archived evidence. Ten exact
Maxima checks verify the coordinate and span algebra. No new cofinal
estimate results from this test.

\section{A same-Gram control with a positive error floor (NS-74)}
\label{sec:ns74-same-gram-control}

\paragraph{Scope.}
The finite efficiency of a direction is a fraction of the available
improvement, not an asymptotic convergence test. The following explicit
altered family preserves every complete Gram entry, can approximate any
fixed finite list of the original loads arbitrarily closely, and has an
unconditional nonzero error floor. It is a control on that inference.
It supplies neither an off-line zeta zero nor an error floor for the
original arithmetic family.

Fix $1/2<\beta<1$, set $a=\beta-1/2$, and put
\[
 B_\beta(s)=\frac{s-\beta}{s-(1-\beta)},\qquad
 b_\beta=\frac{1-\beta}{\beta},\qquad
 c_\beta=\frac{2\beta-1}{\beta^2}.
\]
Let $\mathcal U:L^2(0,\infty;dx)\to L^2(\mathbb R;dt)$ be
$(\mathcal Uf)(t)=e^{-t/2}f(e^{-t})$. On this latter space define
\[
 (V_\beta u)(t)=u(t)-2a\int_{-\infty}^t
                                      e^{-a(t-r)}u(r)\,dr,
 \qquad U_\beta=\mathcal U^{-1}V_\beta\mathcal U.
\]
The convolution is a bounded $L^2$ operator, since its kernel is in
$L^1$. Its Fourier multiplier is $(it-a)/(it+a)$, of modulus one,
so $U_\beta$ is unitary. Define the altered atoms
$\sigma_n^{(\beta)}=U_\beta\sigma_n$, keeping the original target
$\tau$ and norm fixed.

\begin{proposition}[Identical geometry, altered loads, and an explicit floor]
\label{prop:ns74-same-gram}
For every $m,n\geq1$ the complete altered Gram satisfies
\[
 \langle\sigma_m^{(\beta)},\sigma_n^{(\beta)}\rangle=G(m,n).
\]
The target loads are exactly
\begin{equation}
 b_n^{(\beta)}=\langle\tau,\sigma_n^{(\beta)}\rangle
 =b_\beta b_n-\frac{c_\beta}{n}
                    \left(\log n+2-\gamma+\frac1\beta\right).
 \label{eq:ns74-loads}
\end{equation}
For every finite coefficient vector, including optimized coefficients,
\begin{equation}
 \left\|\tau-\sum_n q_n\sigma_n^{(\beta)}\right\|_2^2
 \geq F_\beta:=\frac{(2\beta-1)(1-\beta)^2}{\beta^6}>0.
 \label{eq:ns74-floor}
\end{equation}
All actual and elementary-model Gram matrices, old-space Schur
complements, and geometric preconditioners are unchanged. The loads
and residuals are changed.
\end{proposition}
\begin{proof}
Unitarity proves the Gram equality on the complete space. In logarithmic
coordinates $\mathcal U\tau=t e^{-t/2}\mathbf1_{t>0}$.
Evaluating the adjoint anti-causal convolution gives
\[
 U_\beta^*\tau(x)=
 \begin{cases}
 b_\beta(-\log x)-c_\beta,&0<x<1,\\
 -c_\beta x^{-\beta},&x>1.
 \end{cases}
\]
Both pieces are retained. Since $\sigma_n(x)=1/(nx)$ for $x>1$ and
\[
 \int_0^1\sigma_n(x)\,dx=\frac{\log n+2-\gamma}{n},
 \qquad \int_1^\infty x^{-\beta}\sigma_n(x)\,dx=\frac1{n\beta},
\]
the load formula follows. The first integral follows from Fubini in
$\sigma_n=J\rho_n$: its interior contribution is
$(\log n+1-\gamma)/n$ and its exterior contribution is $1/n$.

All original atoms lie in the closed tail space
\[
 \mathcal T=L^2(0,1)\oplus
              \operatorname{span}\{x^{-1}\mathbf1_{x>1}\}.
\]
Projecting the exterior piece of $U_\beta^*\tau$ onto that one-dimensional
tail leaves squared distance
\[
 c_\beta^2\left(\frac1{2\beta-1}-\frac1{\beta^2}\right)
 =F_\beta.
\]
The original approximation span is contained in $\mathcal T$, so its
distance is at least this value. Applying the unitary $U_\beta$ proves
\eqref{eq:ns74-floor}. No assumption about density of the original span
in $\mathcal T$, or about RH, is made.
\end{proof}

\paragraph{The corresponding explicit dual obstruction.}
The altered atoms still have exterior tail $b_\beta/(nx)$: integrate the
causal convolution against $e^{t/2}/n$ for $t<0$.
Their Mellin transforms, for $1-\beta<\Re s<1$, are
\[
 \widehat{\sigma_n^{(\beta)}}(s)
     =-B_\beta(s)\frac{\zeta(s)}{s^2n^s}.
\]
This identity follows either by absolutely convergent convolution in a
smaller strip and continuation, or directly from the physical formula.
Thus the deliberately inserted zero at $s=\beta$ annihilates every atom.
On $\mathcal T$, evaluation at $\beta$ is represented by
\[
 w_\beta(x)=x^{\beta-1}\mathbf1_{0<x<1}
          +\frac{x^{-1}}{1-\beta}\mathbf1_{x>1},
 \quad
 \|w_\beta\|_2^2=\frac1{2\beta-1}+\frac1{(1-\beta)^2}.
\]
Since $\langle\tau,w_\beta\rangle=1/\beta^2$, Cauchy--Schwarz
recovers precisely $F_\beta$. This is a zero of the inserted factor,
not an asserted real or nonreal zero of $\zeta$.

\begin{corollary}[Finite efficiency cannot distinguish convergence]
\label{cor:ns74-finite-indistinguishability}
For any fixed finite set of block sizes with nonzero original gain,
the altered optimized energies, full gains, and either curvature
rule's gain fractions converge to their original values as
$\beta\downarrow1/2$. Nevertheless every fixed $\beta>1/2$ has the
positive floor~\eqref{eq:ns74-floor} at all sizes. Consequently no
strict finite collection of successful gain-fraction inequalities,
together with the shared Gram geometry, implies convergence to zero.
\end{corollary}
\begin{proof}
The loads~\eqref{eq:ns74-loads} converge entrywise to $b_n$.
At each fixed size all Gram matrices are unchanged and positive
definite. Hence their solved coefficients, projected loads and errors
are continuous in $\beta$. Each curvature direction is a fixed linear
map of those loads, and the positive gain denominators stay nonzero
near the limit. Continuity proves the statement for each block and
therefore for a finite list. The floor is strictly positive for every
fixed parameter. There is no interchange of the limits $N\to\infty$
and $\beta\downarrow1/2$.
\end{proof}

\paragraph{Certified control.}
For $\beta=1/2+10^{-6}$ at $N=256$, the complete squared residual is
in $(.00004614,.00004615)$, while the unconditional floor is in
$(.00003199,.00003200)$. Thus more than $69.34\%$ of this finite error
cannot be removed by any further approximation in the altered family.
Nevertheless its diagonal curvature direction captures a fraction in
$(.9763,.9764)$ of the available full block improvement, essentially
the same finite efficiency as the original family. Its relative squared
error reduction is only in $(.02491,.02493)$, compared with the original
fraction in $(.08129,.08130)$: the distinction between gain fraction and
error reduction matters.

The parameter $\beta=1/2+10^{-10}$ gives the further exact positive floor
$F_\beta\in(3.1999,3.2001)\cdot10^{-9}$ and still closer finite statistics.
Both parameters and the original control are checked at all five sizes,
at 256/384 bits and with the doubled cutoff at $N=256$. Sixteen exact
Maxima checks cover the complete norm, tail, multiplier, dual functional
and limiting identities. The original family's limiting error remains
unknown; the positive floors in this paragraph belong only to the
explicitly altered families.

\section{Divisor-defect feedback from the actual residual (NS-75)}
\label{sec:ns75-divisor-feedback}

\paragraph{Question.}
The fixed M\"obius templates in NS-73 do not use the optimized old
coefficients. We now test a correction that does: cancel the next block
of derivative jumps of the actual residual. This yields an exact
arithmetic construction, but removing those local jumps is not a
norm-descent theorem.

Write the complete residual as
$R_N=\tau-\sum_{n\leq N}c_n\sigma_n$, with the actual optimized $c_n$,
and put
\[
 \eta_c=\sum_{n\leq N}\frac{c_n}{n},\qquad
 r_N(m)=\mathbf1_{m=1}+\sum_{\substack{n\mid m\\n\leq N}}c_n.
\]
The following statements also hold for arbitrary fixed real old coefficients.

\begin{proposition}[Complete divisor-log cells and derivative jumps]
\label{prop:ns75-divisor-cells}
For every $x>0$,
\begin{equation}
 R_N(x)=-\frac{\eta_c}{x}
       +\sum_{m\leq1/x}r_N(m)\log\frac1{mx}.
 \label{eq:ns75-divisor-cells}
\end{equation}
At $x=1/m$ the function is continuous and the right-minus-left jump of
its first derivative is $m r_N(m)$. In distributions on $(0,\infty)$,
\begin{equation}
 R_N''=\frac1{x^2}\sum_{m\leq1/x}r_N(m)
           -\frac{2\eta_c}{x^3}
           +\sum_{m\geq1}m r_N(m)\delta_{1/m}.
 \label{eq:ns75-jump-measure}
\end{equation}
The atomic sum is locally finite away from zero. The complete exterior
tail is $R_N(x)=-\eta_c/x$ for $x>1$.
\end{proposition}
\begin{proof}
The exact identity
$A(z)=z-\sum_{k\leq z}\log(z/k)$ gives
\[
 \sum_{n\leq N}c_n\sigma_n(x)
 =\frac{\eta_c}{x}-\sum_{m\leq1/x}
       \left(\sum_{\substack{n\mid m\\n\leq N}}c_n\right)
       \log\frac1{mx}.
\]
For each fixed $x$ both sums are finite, so this rearrangement requires
no limiting interchange. The target is exactly the term $m=1$ in the
logarithmic sum. A term entering at $1/m$ has value zero there, proving
continuity. Between endpoints, differentiation gives
$R_N'=\eta_c/x^2-x^{-1}\sum_{m\leq1/x}r_N(m)$.
The stated jump and second-derivative distribution follow, including
the target endpoint $m=1$.
\end{proof}

\paragraph{An exact local cancellation step.}
For $N<n\leq2N$ choose physical new coefficients
\[
 a_n^{(p)}=-r_N(n)\log^p(2N/n),\qquad p=0,1,2,
\]
with the constant taper convention for $p=0$.
With $p=0$, append these coefficients while holding all old coefficients
fixed. For $N<m\leq2N$, every proper divisor of $m$ is at most $N$.
Consequently the updated divisor defect is exactly
$r_N(m)+a_m^{(0)}=0$ throughout this next block. The old defects are
unchanged. Other cells, the regular part of~\eqref{eq:ns75-jump-measure}
and both tails remain in the norm calculation. No smallness follows
from cancellation of the displayed atomic part alone.

If the old coefficients happen to be $c_n=-\mu(n)$, then
$r_N(m)=\mu(m)$ on this next block, so the constant rule specializes
to NS-73's sharp M\"obius template. This is a specialization, not an
assumption about the optimized old coefficients.

For comparison with the optimum we also convert each $a^{(p)}$ to
difference coordinates using~\eqref{eq:ns73-coordinate-change}, apply
the complete old-space complement, and optimize its scalar step.
The exact gain and the best three-direction span gain are then those
of Proposition~\ref{prop:ns73-span}. This old-space reoptimization
can change the local jumps again; we do not claim it preserves the
raw cancellation property. We separately record the complete error
of the held-old, unit-sized step:
\[
 \frac{\|R_N-\sum_{n=N+1}^{2N}a_n^{(p)}\sigma_n\|_2^2}{E_N}
 =1+\frac{(a^{(p)})^TG_{\rm new,new}a^{(p)}
                 -2(g^Tv^{(p)})}{E_N}.
\]
Old-space orthogonality justifies the numerator equality here;
the raw physical Gram in this expression is not replaced by a
projected or model Gram.

\paragraph{Certified finite outcome.}
The complete 256/384-bit replays and doubled-cutoff replay at $N=256$
give the following strict enclosures for the fraction of full block gain.
The single cancellation direction is optimally scaled \emph{after} old-space
projection; its scalar is not constrained positive.
\begin{center}
\begin{tabular}{rccc}
\hline
$N$&Cancellation direction&Best three-direction span&Diagonal curvature\\
\hline
16&(.3110,.3111)&(.5836,.5837)&(.9308,.9310)\\
32&(.4292,.4294)&(.8086,.8087)&(.9040,.9041)\\
64&(.6954,.6956)&(.9063,.9065)&(.9783,.9784)\\
128&(.0029,.0030)&(.7232,.7233)&(.9852,.9854)\\
256&(.0010,.0011)&(.5602,.5603)&(.9763,.9764)\\
\hline
\end{tabular}
\end{center}
At $N=256$ the span retains $56.02048\%$ of full gain, versus
$97.63068\%$ for curvature. The raw held-old cancellation step increases
the complete squared error by a factor in $(1495.54,1495.55)$.
Indeed all fifteen held-old unit steps tested here increase squared error.
The cancellation direction's correlation at $N=256$ is negative, so its
best scalar step reverses the prescribed sign. Merely cancelling derivative
jumps is therefore not a descent principle even on these certified blocks.
This is a finite failure of the proposed unit step and finite underperformance
of its span, not an asymptotic exclusion of every divisor-feedback scheme.
Eleven exact Maxima checks cover the identities and coordinate conversion;
all finite gains retain complete tails, the true old projection and all
three-direction cross terms.

\section{The continuous relaxation and its exact inner-factor defect (NS-76)}
\label{sec:ns76-inner-distance}

\paragraph{Scope and prior work.}
Burnol's quantitative Nyman formula identifies the $q=1$ continuous
approximation defect through the off-critical Blaschke product
\cite[Theorem 1.2 and Lemma 2.4]{BurnolNyman2001}. We apply the same
classical Hardy-space mechanism to this project's $q=2$ target, keeping
its full exterior tail. The calculation gives an exact distance for
\emph{continuous} dilations and a lower bound for integer dilations.
No unconditional equality of the two nonzero distances is assumed.
This is an explicit reformulation of the unknown arithmetic inner
factor, not a new zero-free estimate.

Let $\mathcal T$ be the complete tail space from NS-74, and let
$\mathbb H^2$ denote the Hardy space of $\Re s>1/2$, with boundary norm
$(2\pi)^{-1}\int_{\mathbb R}|f(1/2+it)|^2dt$.
Multiplication on the Mellin boundary by $(s-1)/s$ induces a unitary
map $W$ on the full $L^2$ space. It maps $\mathcal T$ onto
$L^2(0,1)$, whose Mellin image is $\mathbb H^2$.
Indeed in logarithmic coordinates it is
\[
 (Wu)(t)=u(t)-\int_{-\infty}^t e^{-(t-r)/2}u(r)\,dr.
\]
For $t<0$ a tail vector is $C e^{t/2}$ and cancels exactly in this
formula. Conversely the adjoint applied to a function supported on
$t\geq0$ has on $t<0$ a constant multiple of $e^{t/2}$.
Thus the map is onto, and the full tail has been accounted for.

The transformed target and first atom have Mellin transforms
\[
 T(s)=\frac{s-1}{s^3},\qquad
 F(s)=-\frac{(s-1)\zeta(s)}{s^3}.
\]
Write $B$ for the off-critical Blaschke product, normalized by $B(1)>0$:
\[
 B(s)=\prod_{\substack{\zeta(\rho)=0\\\Re\rho>1/2}}
 \frac{s-\rho}{s-(1-\overline\rho)}
 \frac{1-\overline\rho}{\rho}
 \left|\frac{\rho}{1-\rho}\right|.
\]
Multiplicities are included. Burnol's factorization proves that the
inner factor of $(s-1)\zeta(s)/s^2$ is exactly $B$: analytic
continuation excludes a finite-boundary singular factor, and the
real-axis asymptotic excludes an exponential factor. The additional
factor $1/s$ is outer, so $F$ has the same inner factor. The product
and its derivatives converge locally in $\Re s>1/2$ away from its
zeros; $B(1)$ and its first two derivatives are real by conjugation
symmetry.

\begin{proposition}[Exact distance for the complete continuous family]
\label{prop:ns76-continuous-distance}
Let $\mathcal V_{\rm cont}$ be the closed span in the original Hilbert
space of $\sigma_u$, for all real $u\geq1$. Its image under $W$ and
Mellin transformation is $B\mathbb H^2$.
Put $b=B(1)$, $d=B'(1)$ and $e=B''(1)$. Then
\begin{equation}
 D_{\rm cont}^2:=\operatorname{dist}(\tau,\mathcal V_{\rm cont})^2
 =2-2b^2+2bd-d^2-\frac{e^2}{4}.
 \label{eq:ns76-continuous-distance}
\end{equation}
This expression is nonnegative and at most every integer finite
squared error $E_N$. It vanishes exactly when RH holds.
The formula contains unknown off-critical zero information and
provides no unconditional upper bound tending to zero.
\end{proposition}
\begin{proof}
The normalized real dilations of $\sigma_1$ become the causal
translation semigroup applied to $F$. The Beurling--Lax theorem gives
that its closed generated subspace is the inner factor $B$ times
$\mathbb H^2$. This use of the full continuous semigroup is the reason
we do not replace it by the integer set without a separate argument.

For $\Re z>1/2$, let $k_z(s)=1/(s+\overline z-1)$ be the evaluation
kernel. Multiplication $M_B$ is an isometry and
$M_B^*k_z=\overline{B(z)}k_z$.
Differentiating at real $z=1$ is legitimate in the Hilbert norm,
as is seen directly from the integrable kernel derivatives. Since
$T=-\partial_z k_z|_1-\tfrac12\partial_z^2 k_z|_1$, it follows that
\[
 M_B^*T=-\frac{d+e/2}{s}+\frac{b+d}{s^2}-\frac{b}{s^3}.
\]
The complete boundary Gram for $(1/s,1/s^2,1/s^3)$ is
\[
 \begin{pmatrix}1&1&1\\1&2&3\\1&3&6\end{pmatrix};
\]
its entries follow by integrating the Laplace representatives
$e^{-t/2}t^{j-1}/(j-1)!$ on $t>0$.
Therefore $\|M_B^*T\|^2=2b^2-2bd+d^2+e^2/4$.
The orthogonal projection is $B M_B^*T$, with that same norm, and
$\|T\|^2=\|\tau\|^2=2$. This proves the distance identity and
nonnegativity without any sign assumption on $d$ or $e$.

The integer span is contained in the continuous span, giving
$E_N\geq D_{\rm cont}^2$. If RH holds, the product is empty and
$(b,d,e)=(1,0,0)$. Conversely, any off-critical zero $\rho$ gives a
zero of $B$ in this half-plane, whereas $T(\rho)\ne0$ since
$\rho\notin\{0,1\}$. Continuous evaluation in $\mathbb H^2$ then
prevents $T$ from belonging to the closed subspace $B\mathbb H^2$,
so its distance is positive.
\end{proof}

\begin{corollary}[The fixed component invisible to every block gain]
\label{cor:ns76-invisible-component}
Let $P_{\rm cont}$ and $P_N$ be the complete projections onto the
continuous and finite integer spaces. Then exactly
\[
 R_N=(I-P_{\rm cont})\tau+(P_{\rm cont}\tau-P_N\tau),
 \qquad
 E_N=D_{\rm cont}^2+\|P_{\rm cont}\tau-P_N\tau\|^2.
\]
Every new correlation $g_n$ and every selected-direction gain depend
only on the second component. The first is orthogonal to every real
dilation, hence also to each projected new integer atom.
\end{corollary}
\begin{proof}
The finite integer space is a subspace of $\mathcal V_{\rm cont}$,
so $P_NP_{\rm cont}=P_N$. The two displayed components are
orthogonal. Taking inner products with any old-projected new atom
annihilates the first, proving the final statement.
\end{proof}

\begin{proposition}[The integer and continuous closures are strictly different]
\label{prop:ns76-strict-closure}
Let $\mathcal V_{\rm int}$ be the closed span of the integer atoms.
Then $\mathcal V_{\rm int}\subsetneq\mathcal V_{\rm cont}$,
unconditionally. More explicitly,
\[
 \operatorname{dist}(\sigma_{3/2},\mathcal V_{\rm int})^2>.00111713.
\]
This is a bound for $\sigma_{3/2}$, not for the target $\tau$.
\end{proposition}
\begin{proof}
On $I=(1/2,1)$, the first integer atom is
$\sigma_1=1/x+\log x$ and every integer atom with $n\geq2$
is $1/(nx)$. Their restricted closed span lies in the closed
two-dimensional space $\operatorname{span}\{1/x,\log x\}$.
In contrast, $\sigma_u$ for $1<u<2$ has a nonzero derivative jump
at the interior point $x=1/u$; it cannot equal an element of that
analytic two-dimensional span, even almost everywhere. This already
proves strict inclusion.

For the quantitative bound, put $u=3/2$, $L=\log2$, $c=\log u$ and
$h=L-c$. The complete local Gram, loads and squared norm are
\[
 G_I=\begin{pmatrix}1&-L^2/2\\-L^2/2&1-L-L^2/2\end{pmatrix},\qquad
 b_I=\begin{pmatrix}
 u^{-1}-h^2/2\\
 (c+2-L^2/2)/u+c(L+1)/2-(L^2+2L+2)/2
 \end{pmatrix},
\]
\[
 q_I=u^{-2}-h^2/u+2/u-h^2/2-h-1.
\]
These follow by splitting only at $x=1/u$, with
$\sigma_u=1/(ux)+\log(ux)$ on the left and $1/(ux)$ on the right.
The Gram is positive definite. Its exact local projection error
$q_I-b_I^TG_I^{-1}b_I$ lies in $(.00111713,.00111714)$ by independent
256/384-bit enclosures; Maxima checks all complete elementary integrals.
Restriction of a norm to $I$ can only decrease it, so this local
minimum bounds the full-space distance below. No omitted tail is
set to zero in the asserted full-space inequality.
\end{proof}

\paragraph{What the calculation does and does not resolve.}
The $q=2$ continuous defect is now expressed through three values of the
classical unknown inner factor. It is not estimated away. In particular,
the integer limiting error may contain an additional approximation
component. The closures are strictly different, but that fact alone
does not decide whether their two distances from this particular target
are different. No equality of positive target distances is supplied.
The original sufficient cofinal contraction remains valid, because it
would force the entire $E_N$ to zero, including this fixed component.
Replacing $E_N$ by $E_N-D_{\rm cont}^2$ would remove exactly the part
that must be ruled out and would not prove RH.

\section{How a high zero tail can hide the continuous defect (NS-77)}
\label{sec:ns77-continuous-tail}

\paragraph{Scope.}
We give a complete upper bound on the continuous defect in NS-76 using
positive sums over hypothetical off-critical zeros. This quantifies a
limitation of finite geometric evidence. It is not an upper bound on
$E_N$ or its integer limit, not a new verified zero height, and not a
proof that the defect is exactly zero.

Suppose every off-critical zero in the upper half-plane has height
$\gamma>H\geq1$. Index the conjugate pairs by $\rho=\beta+i\gamma$,
with $1/2<\beta<1$, once per pair and with multiplicity, and put
\[
 S_2=\sum_\rho\frac{2\beta-1}{\gamma^2},\qquad
 S_4=\sum_\rho\frac{2\beta-1}{\gamma^4}.
\]
These sums converge by the usual zero-counting bound. Empty sums are zero.

\begin{proposition}[A complete positive zero-tail budget]
\label{prop:ns77-tail-budget}
For the continuous $q=2$ distance,
\begin{equation}
 0\leq D_{\rm cont}^2\leq
       \min\left\{2,\frac83S_2^3+5S_4\right\}.
 \label{eq:ns77-tail-budget}
\end{equation}
If the total upper-half-plane zero count satisfies
$N(t)\leq A t\log t$ for all $t\geq H\geq100$, then
\begin{equation}
 D_{\rm cont}^2\leq
 \frac{\frac{64}{3}A^3(\log H+1)^3
                  +20A(\log H/3+1/9)}{H^3}.
 \label{eq:ns77-height-bound}
\end{equation}
No tail prefix is substituted for the complete zero sums.
\end{proposition}
\begin{proof}
For one pair write $a=\beta$, $c=1-\beta$, $\delta=a-c=2\beta-1$,
so $a+c=1$. Its real-normalized inner factor at $1$ is
$B_\rho(1)=(\gamma^2+c^2)/(\gamma^2+a^2)$.
Set $j_\rho=-\log B_\rho(1)$ and
$\ell_\rho=(\log B_\rho)'(1)$. Exactly,
\[
 j_\rho=\int_{c^2}^{a^2}\frac{dt}{\gamma^2+t},\qquad
 \ell_\rho=-\frac{2\delta(\gamma^2-ac)}
                        {(\gamma^2+a^2)(\gamma^2+c^2)}.
\]
Thus $0\leq j_\rho\leq\delta/\gamma^2$ and
\[
 \ell_\rho+\frac{2\delta}{\gamma^2}
 =\frac{2\delta\{\gamma^2(a^2+c^2+ac)+a^2c^2\}}
        {\gamma^2(\gamma^2+a^2)(\gamma^2+c^2)}
 \leq\frac{17\delta}{8\gamma^4},
\]
using $a^2+c^2+ac\leq1$, $a^2c^2\leq1/16$, and $\gamma\geq1$.
In particular $\ell_\rho+2j_\rho\leq17\delta/(8\gamma^4)$.
Absolute convergence permits summing these inequalities and the
logarithmic derivatives. With
\[
 J=-\log B(1)\in[0,S_2],\qquad \ell=B'(1)/B(1),
 \qquad \epsilon=\ell+2J,
\]
we have $\epsilon\leq17S_4/8$. There is no lower-sign assumption on
$\epsilon$.

Let $\ell_2=(\log B)''(1)$, a real number. Substituting
$b=e^{-J}$, $d=b\ell$, $e=b(\ell^2+\ell_2)$ into
\eqref{eq:ns76-continuous-distance} and discarding only the
nonnegative square that is subtracted gives
\begin{align*}
 D_{\rm cont}^2
 &\leq 2-e^{-2J}(2-2\ell+\ell^2)\\
 &=2-e^{-2J}(2+4J+4J^2)
       +(2+4J)e^{-2J}\epsilon-e^{-2J}\epsilon^2.
\end{align*}
The first term on the last line is zero at $J=0$ and has derivative
$8J^2e^{-2J}$, so it is at most $8J^3/3$.
Also $(2+4J)e^{-2J}\leq2$ for $J\geq0$. If $\epsilon\leq0$ its
linear contribution is nonpositive; otherwise it is at most
$17S_4/4<5S_4$. This proves~\eqref{eq:ns77-tail-budget}.

For the second assertion, $0<2\beta-1<1$ and partial summation give,
for $k=2,4$,
\[
 S_k\leq k\int_H^\infty\frac{N(t)}{t^{k+1}}\,dt
 \leq kA H^{1-k}
             \left(\frac{\log H}{k-1}+\frac1{(k-1)^2}\right).
\]
The endpoint at infinity vanishes and dropping the nonpositive lower
endpoint term only enlarges the bound. Inserting these complete
integrals proves~\eqref{eq:ns77-height-bound}.
\end{proof}

\paragraph{Calibration using published finite-height information.}
The bound of Hasanalizade--Shen--Wong
\cite[Corollary 1.2]{HasanalizadeShenWong2021} implies
$N(t)\leq t\log t$ for $t\geq100$. Indeed its main term is at most
$(t/6)\log t$, and its remainder is at most
$.362\log t+10$; since $\log100>4$, their sum is less than
$(t/6+3)\log t\leq t\log t$ on this range.
Platt--Trudgian's published interval verification
\cite[Theorem 1]{PlattTrudgian2021} supplies the absence of off-critical
zeros up to $H=3\cdot10^{12}$. Taking $A=1$ in the analytic bound
therefore gives the explicit consequence
\begin{equation}
 0\leq D_{\rm cont}^2<2.08\cdot10^{-32}.
 \label{eq:ns77-published-calibration}
\end{equation}
Only the final scalar evaluation is replayed here at 256/384 bits;
the cited large zero verification is an external published input,
not a new computation or an archived verification by this project.
This bound is deliberately conservative.

A small positive continuous defect remains compatible with this
inequality. Moreover the integer and continuous closures are strictly
different by Proposition~\ref{prop:ns76-strict-closure}, and no matching
upper bound on the integer limiting error has been supplied.
Thus the calibration explains why the relaxed problem can look
extremely favorable without proving the original approximation target,
a new zero-free region, G2 or RH.

\section{A complete norm from reciprocal-integer samples (NS-78)}
\label{sec:ns78-sampling}

\paragraph{Purpose.}
NS-75 identifies the signed divisor defects but shows that cancelling
individual jumps need not decrease the norm. Here we retain the entire
residual through its reciprocal-integer values, with constants independent
of the number and size of the coefficients. This is an exact geometric
reduction and a complete physical-space replay, not the missing arithmetic
estimate. The infinite sample sum is not replaced by a finite prefix.

For arbitrary real $\alpha,c_1,\ldots,c_N$, put
\[
 R=\alpha\tau-\sum_{n\leq N}c_n\sigma_n,\qquad
 \eta=\sum_{n\leq N}\frac{c_n}{n},\qquad
 z_m=R(1/m),\qquad w_m=\frac1{m(m+1)}.
\]
Define the complete nonnegative arithmetic quantity
\begin{equation}
 \mathcal S(c)=\eta^2+
       \sum_{m=1}^\infty w_m(z_m^2+z_{m+1}^2).
 \label{eq:ns78-sampling-functional}
\end{equation}
Its samples are explicitly
\[
 z_m=-\eta m+\sum_{k\leq m}
 \left(\alpha\mathbf1_{k=1}+\sum_{\substack{n\mid k\\n\leq N}}c_n\right)
                    \log(m/k).
\]
All cancellation inside each displayed sample is retained before squaring.

\begin{proposition}[Uniform complete sampling comparison]
\label{prop:ns78-sampling}
For every real $\alpha$ and finite real coefficient vector,
\begin{equation}
 \frac1{16}\mathcal S(c)\leq\|R\|_2^2
                 \leq\frac{21}{16}\mathcal S(c).
 \label{eq:ns78-sampling-comparison}
\end{equation}
Both sides include the exterior tail and every reciprocal cell. The
infinite series in~\eqref{eq:ns78-sampling-functional} converges.
\end{proposition}
\begin{proof}
Change variables to $t=1/x$ and write $r(t)=R(1/t)$, with measure
$dt/t^2$. On $0<t<1$, $r(t)=-\eta t$, so the complete exterior
energy is exactly $\eta^2$. By NS-75, on $m\leq t\leq m+1$,
\[
 r(t)=-\eta t+D_m\log t-L_m,
\quad D_m=\sum_{k\leq m}\delta_k,
\quad L_m=\sum_{k\leq m}\delta_k\log k.
\]
Here $\delta_k=\alpha\mathbf1_{k=1}+\sum_{n\mid k,\,n\leq N}c_n$
is the signed divisor defect.
Let $h_m=\log(1+1/m)$, $\theta=\log(t/m)/h_m$, and set
\[
 p(t)=(1-\theta)z_m+\theta z_{m+1},\qquad
 b(t)=\theta-(t-m).
\]
Exactly $r(t)=p(t)+\eta b(t)$ on every cell.
Under $t=m e^{h_m\theta}$ the measure is
$(h_m/m)e^{-h_m\theta}d\theta$. Its density divided by $w_m$
is between $m h_m$ and $(m+1)h_m$, hence between $1/2$ and $2$.
Also
\[
 \int_0^1((1-\theta)a+\theta d)^2d\theta
       =\frac{a^2+ad+d^2}{3}
       \in\left[\frac{a^2+d^2}{6},\frac{a^2+d^2}{2}\right].
\]
Thus, writing $V=\int_1^\infty p(t)^2dt/t^2$ and
$S=\sum_{m\geq1}w_m(z_m^2+z_{m+1}^2)$, we get
$S/12\leq V\leq S$, initially for finite unions of cells.

Concavity of $\log(m+v)$ for $0\leq v\leq1$, and its second
derivative bound, give
\[
 0\leq b(m+v)\leq\frac{v(1-v)}{2m^2h_m}
       \leq\frac1{8m^2h_m}\leq\frac1{4m}.
\]
Consequently
\[
 \Gamma:=\int_1^\infty b(t)^2\frac{dt}{t^2}
 \leq\frac1{16}\sum_{m\geq1}\frac{w_m}{m^2}
 \leq\frac1{16},
\]
since $\sum w_m=1$. The original finite combination lies in $L^2$,
so $p=r-\eta b$ is square integrable. Monotone passage through the
finite-cell estimates now proves finiteness of $S$ and their complete
versions; no convergence of $S$ was assumed.

Cauchy--Schwarz and Young's inequality give
\[
 |2\eta\langle p,b\rangle|
       \leq\eta^2/4+4\Gamma V\leq(\eta^2+V)/4.
\]
Using the exact identity
$\|R\|_2^2=\eta^2+V+2\eta\langle p,b\rangle+\eta^2\Gamma$,
we obtain
\[
 \frac34(\eta^2+V)\leq\|R\|_2^2
 \leq\frac{21}{16}\eta^2+\frac54V
 \leq\frac{21}{16}(\eta^2+V).
\]
Inserting $S/12\leq V\leq S$ proves the proposition.
\end{proof}

\paragraph{The arithmetic input in explicit coordinates.}
Specialize to $\alpha=1$. By the fixed-order smoothed Nyman--Beurling
criterion in NS-61, RH is
equivalent to the existence of coefficient vectors, of unbounded allowed
length, with $\mathcal S(c)\to0$. For the actual optimized coefficients,
this is equivalent to $E_N\to0$. Boundedness of $\mathcal S(c)$ alone
does not suffice; it is not the bounded-Weil-floor criterion.
The local geometric comparison is unconditional, but no sequence with
vanishing complete $\mathcal S$ has been produced. A finite collection
of small samples cannot replace its infinite sum or its cofinal limit.

The case $\alpha=0$ also bounds the complete energy of any finite
correction vector. Its comparison quantity still contains all signed
arithmetic samples; this is not an all-vector comparison with the
smooth repaired model excluded in NS-70.

\paragraph{Exact inversion and the remaining arithmetic constraint.}
The samples also recover the coefficients without a Gram solve. Put
$D_0=0$. Continuity at every reciprocal knot gives
\[
 \eta=-z_1,\qquad
 D_m=\frac{z_{m+1}-z_m+\eta}{\log(1+1/m)},\qquad
 \delta_m=D_m-D_{m-1}.
\]
M\"obius inversion then gives, extending $c_n=0$ beyond $N$,
\begin{equation}
 c_n=\sum_{d\mid n}\mu(n/d)\delta_d-\alpha\mu(n).
 \label{eq:ns78-sample-inverse}
\end{equation}
For $n>N$ the right-hand side must therefore vanish. These infinitely
many arithmetic compatibility conditions distinguish admissible sample
sequences from arbitrary small knot values.

In particular, norm convergence of residuals with $\alpha=1$ forces
$c_n\to-\mu(n)$ at each fixed index $n$. Indeed
$|z_m|\leq4\sqrt{m(m+1)}\,\|R\|_2$ follows from the complete
sampling comparison, while $|\eta|\leq\|R\|_2$ follows from the
exterior tail. The finite formulas above then force $D_m,\delta_m\to0$
for every fixed $m$, and~\eqref{eq:ns78-sample-inverse} gives the claim.
This necessary coefficientwise limit supplies neither a uniform growing-index
bound nor a norm-convergent sharp M\"obius truncation; the latter was
excluded in NS-62/66. No interchange of those two limits is made.

\subsection{Complete cell integrals and explicit tails}
\label{subsec:ns78-tails}
For $m\geq1$, put $a=\log m$, $d=\log(m+1)$ and
\[
 I_{1,m}=\frac{a+1}{m}-\frac{d+1}{m+1},\qquad
 I_{2,m}=\frac{a^2+2a+2}{m}-\frac{d^2+2d+2}{m+1}.
\]
The exact cell contribution is
\begin{equation}
 \int_m^{m+1}\frac{r(t)^2}{t^2}dt
 =\eta^2+D_m^2I_{2,m}+L_m^2w_m
 -\eta D_m(d^2-a^2)+2\eta L_m(d-a)-2D_mL_mI_{1,m}.
 \label{eq:ns78-exact-cell}
\end{equation}
The signed cross terms are evaluated before accumulating cell energies.

We also need a complete tail, since $m$ extends indefinitely even when
$c$ has finite support. Write
\[
 u=\alpha-\tfrac12\sum c_n,\qquad
 v=\tfrac12\sum c_n\log n-\tfrac12\log(2\pi)\sum c_n,
 \qquad D=\tfrac16\sum n|c_n|.
\]
For $t\geq N$,
\begin{equation}
 r(t)=u\log t+v+e(t),\qquad |e(t)|\leq D/t.
 \label{eq:ns78-complete-asymptotic}
\end{equation}
Indeed $A'(z)=\{z\}/z$ almost everywhere and Stirling's formula at
integer endpoints identifies the limiting constant
$A(z)-\tfrac12\log z\to\tfrac12\log(2\pi)$.
Using the periodic Bernoulli polynomials,
\[
 A(z)-\tfrac12\log(2\pi z)
 =-\int_z^\infty\frac{B_1(\{t\})}{t}dt
 =\frac{B_2(\{z\})}{2z}
       -\frac12\int_z^\infty\frac{B_2(\{t\})}{t^2}dt.
\]
The first integral is an improper convergent integral; the second is
absolutely convergent. Since $|B_2|\leq1/6$, the remainder is at most
$1/(6z)$. The periodic $B_2$ is continuous at integer joins, so no
endpoint mass is discarded. Summing this bound gives
\eqref{eq:ns78-complete-asymptotic}.

Fix an integer $M\geq N$ and put $q=u\log M+v$. The complete
main-tail energy and remainder budget are
\[
 T_M=\frac{q^2+2uq+2u^2}{M},\qquad
 U_M=\frac{D^2}{3M^3}.
\]
Thus the actual integral from $M$ to infinity differs from $T_M$ by
at most $2\sqrt{T_MU_M}+U_M$. For the infinite sample tail, set
$Q=|q|+(D+|u|)/M$. Then
\begin{equation}
 \sum_{m\geq M}w_m(z_m^2+z_{m+1}^2)
 \leq \frac{2(Q^2+2|u|Q+2u^2)}{M}.
 \label{eq:ns78-sampling-tail}
\end{equation}
To verify this last bound, on each $m\leq t\leq m+1$ both endpoint
values have absolute value at most $Q+|u|\log(t/M)$:
use~\eqref{eq:ns78-complete-asymptotic} and
$\log((m+1)/t)\leq1/M$. Integrating the square with measure $dt/t^2$
and summing all cells proves~\eqref{eq:ns78-sampling-tail}.
These are bounds on the complete tails, not estimates based on a zero
prefix or an unmeasured remainder.

\paragraph{Independent complete physical-space check.}
For the actual optimized coefficients at $N=16,64,256$, the interval
calculation accumulates every cell through $M=65536$ and encloses the
entire remaining physical and sample tails by the formulas above.
At both 256 and 384 bits, the resulting complete physical norm encloses
the independently evaluated full Gram norm. The complete ratios satisfy
\[
 \mathcal S(c_{16})/E_{16}\in(3.323,3.325),\qquad
 \mathcal S(c_{64})/E_{64}\in(2.9002,2.9004),\qquad
 \mathcal S(c_{256})/E_{256}\in(2.644,2.646).
\]
The analytic physical-tail uncertainty divided by $E_N$ is below
$1.64\cdot10^{-6}$, $5.20\cdot10^{-6}$ and $2.38\cdot10^{-4}$,
respectively. Replaying $N=256$ with $M=131072$ and the doubled complete
Gram-kernel cutoff reduces that tail ratio below $4.17\cdot10^{-5}$.
All solves retain Arb balls. The precise finite part, main tail and
remainder radius are archived separately; the printed combined decimal
balls may be wider due to decimal serialization. Thirty exact Maxima
checks cover the cell, mass, interpolation and tail identities.
This is an independent norm evaluation on finite coefficient vectors,
not a cofinal estimate of their sample sums.

\section{Balancing the exterior tail of divisor feedback (NS-79)}
\label{sec:ns79-balanced-feedback}

The sampling norm isolates the exterior coefficient $\eta$. We therefore
test whether removing its change repairs the failed NS-75 step. This is
an explicitly modified correction, not a reinterpretation of that test.
Let $a_n^{(p)}=-r_N(n)\log^p(2N/n)$ be the three original new-block
physical directions, with the constant convention for $p=0$. Define
\begin{equation}
 \widetilde a_n^{(p)}=a_n^{(p)}\quad(N<n<2N),\qquad
 \widetilde a_{2N}^{(p)}=-2N\sum_{n=N+1}^{2N-1}\frac{a_n^{(p)}}n.
 \label{eq:ns79-balance}
\end{equation}

\begin{proposition}[Exact tail balance and its scope]
\label{prop:ns79-balance}
Each raw correction $f_p=\sum_{n=N+1}^{2N}\widetilde a_n^{(p)}\sigma_n$
vanishes for $x>1/(N+1)$. Thus the held-old step preserves $\eta$.
For $p=0$ it cancels the new derivative jumps for $N<m<2N$, leaving
the jump at $m=2N$ unconstrained. It is not an exact cancellation of
the entire next block.

The original three-direction span augmented by the physical endpoint
atom $\sigma_{2N}$ equals the balanced three-direction span augmented
by that atom. This identity persists after the complete old-space
projection. No assertion of norm descent follows from these identities.
\end{proposition}
\begin{proof}
The weighted coefficient sum in~\eqref{eq:ns79-balance} is zero.
For $x>1/(N+1)$ every new atom equals $1/(nx)$, proving vanishing
and tail preservation. In the open block each proper divisor is old,
so the NS-75 jump argument applies at every index except the modified
last one. More explicitly, for $N<t<2N$ the held-old updated residual
has one formula $-\eta t+D_N\log t-L_N$, with $D_N,L_N$ from NS-78;
no new jumps occur inside this open interval.

The difference between each original and balanced physical vector is
a multiple of the same endpoint coordinate. Hence augmenting either
span by that coordinate gives the same space; applying any linear
old-space projection preserves this equality. In adjacent-difference
coordinates the endpoint atom has vector $v_n=n/(2N)$, modulo its
old boundary term. The transformation is retained exactly in the
finite calculation.
\end{proof}

Old-space reoptimization can change $\eta$ again. The projected gains
below therefore do not preserve the raw tail balance. Conversely,
repeating only raw tail-preserving steps with a fixed nonzero initial
$\eta$ can never reduce squared error below $\eta^2$, because the
exterior tail remains unchanged. This observation does not exclude
initialization with $\eta=0$ or algorithms that reoptimize old coefficients.

\paragraph{Certified finite outcome.}
At 256/384 bits, with a doubled complete-kernel cutoff at $N=256$,
the fractions of full projected gain have the following strict enclosures.
\begin{center}
\begin{tabular}{rccc}
\hline
$N$&Balanced three-span&Endpoint-augmented four-span&Diagonal curvature\\
\hline
16&(.5960,.5961)&(.8542,.8543)&(.9308,.9310)\\
32&(.6580,.6581)&(.8240,.8241)&(.9040,.9041)\\
64&(.8978,.8979)&(.9192,.9193)&(.9783,.9784)\\
128&(.8809,.8810)&(.8810,.8811)&(.9852,.9854)\\
256&(.5996,.5997)&(.7118,.7119)&(.9763,.9764)\\
\hline
\end{tabular}
\end{center}
Thus removing the tail change improves the $N=256$ three-span fraction
from $56.02\%$ to $59.97\%$, and the full four-span reaches $71.18\%$;
both remain below curvature's $97.63\%$. The raw untapered unit step
at $N=256$ still increases squared error by a factor in $(53.07,53.09)$,
although this is smaller than NS-75's factor above $1495$.
Thirteen of the fifteen raw balanced unit steps increase error. The two
exceptions are the linear and quadratic tapers at $N=128$, whose error
ratios lie in $(.9643,.9645)$ and $(.9890,.9893)$. This is neither a
uniform-descent theorem nor an exclusion of all balanced feedback.
Sixteen exact Maxima checks cover endpoint compensation, coordinate
conversion and the complete exterior floor. The four-dimensional span
keeps all cross terms; no mixed term is discarded to claim positivity.

\section{A uniform preconditioner that retains arithmetic (NS-80)}
\label{sec:ns80-arithmetic-preconditioner}

The complete sampling comparison also bounds every finite Gram matrix.
Its comparator retains the arithmetic values $A(m/n)$ at \emph{all}
integer knots; it is not the smooth model from NS-70. This gives a
uniform finite-block efficiency theorem, but neither a fast inverse
nor the cofinal lower bound on available gain needed for RH.

Define the complete arithmetic sampling Gram
\begin{equation}
 Q_{ij}=\frac1{ij}+\sum_{m=1}^\infty\frac{
 A(m/i)A(m/j)+A((m+1)/i)A((m+1)/j)}{m(m+1)}.
 \label{eq:ns80-sampling-gram}
\end{equation}
The series converges by NS-78 and Cauchy--Schwarz. For any finite
coefficient vector, its quadratic form is exactly the complete
sampling quantity for the corresponding pure atom combination.
Thus on every finite index set,
\begin{equation}
 \frac1{16}Q\preceq G\preceq\frac{21}{16}Q.
 \label{eq:ns80-full-comparison}
\end{equation}
Both matrices are positive definite: zero physical norm forces zero
samples and exterior coefficient, and the exact inversion in NS-78
then forces every coefficient to vanish.

\begin{proposition}[Uniform efficiency with the full arithmetic comparator]
\label{prop:ns80-uniform-efficiency}
Use the ordered basis consisting of the old $N$ atoms and the new
adjacent differences. Let $H_G,H_Q$ be the respective Schur complements
of the old blocks in their complete Grams. Then
\[
 \frac1{16}H_Q\preceq H_G\preceq\frac{21}{16}H_Q.
\]
For any nonzero \emph{actual} residual correlation vector $g$, set
$v=H_Q^{-1}g$. Its optimally scaled true gain satisfies
\begin{equation}
 \frac{(g^Tv)^2}{v^TH_Gv}
 \geq\frac{21}{121}\,g^TH_G^{-1}g.
 \label{eq:ns80-gain-fraction}
\end{equation}
This is uniform in $N$ and $g$. All old projections and all cross terms
are included. It is a fraction of the full available block gain, not
of the current residual energy.
\end{proposition}
\begin{proof}
The full comparison survives the invertible change to adjacent
coordinates. For each new vector $v$, the Schur quadratic form is the
minimum of the complete quadratic form over old coordinates. Taking
those minima on both sides preserves the same two constants; the old
minimizers need not agree.

More generally put $a=1/16$, $b=21/16$ and
$M=H_Q^{-1/2}H_GH_Q^{-1/2}$. Its eigenvalues lie in $[a,b]$.
With $s=H_Q^{-1/2}g$, normalize its squared spectral coordinates to
probability weights and write $x=\mathbb E\lambda$,
$y=\mathbb E\lambda^{-1}$. For $a\leq\lambda\leq b$,
\[
 \lambda+\frac{ab}{\lambda}\leq a+b,
 \qquad x+aby\leq a+b,
 \qquad xy\leq\frac{(a+b)^2}{4ab}.
\]
The final inequality follows from $(x-aby)^2\geq0$.
The ratio of the selected gain to the full gain is precisely
$1/(xy)$, hence is at least $4ab/(a+b)^2=21/121$.
This is the usual elementary Kantorovich argument, displayed here
to make every hypothesis and normalization explicit.
\end{proof}

\paragraph{Repeated steps within a fixed finite block.}
Start at $x_0=0$ and minimize the exact block quadratic
$\Phi(x)=x^TH_Gx-2g^Tx$. At each step use the current residual
$r_k=g-H_Gx_k$, direction $p_k=H_Q^{-1}r_k$, and the exact line search
\[
 x_{k+1}=x_k+\frac{r_k^Tp_k}{p_k^TH_Gp_k}p_k,
\]
with termination if $r_k=0$. The remaining block gain is
$r_k^TH_G^{-1}r_k$, so Proposition~\ref{prop:ns80-uniform-efficiency}
reduces it by at least $21/121$ each time. After $k$ steps the captured
fraction of full gain is therefore at least
\[
 1-(100/121)^k.
\]
In particular twenty ideal steps capture more than $97.79\%$, uniformly
in block size. This statement assumes application of the complete
$H_Q^{-1}$ and the true $H_G$; no implementation or complexity bound for
that infinite-sample comparator is asserted. Truncating its sum without
a certified tail would invalidate the theorem's premises.

For completeness, minimizing the full sampling objective over the first
$N$ coefficients gives a residual with squared error at most $21E_N$:
apply both bounds of NS-78 before and after the two minimizations.
This uniform approximation factor is also not a decay estimate for $E_N$.

\paragraph{The remaining obstruction.}
Even perfect capture of every block gain does not prove that the limiting
error is zero. The full block gains telescope to $E_N-E_\infty$, and
NS-76 exhibits a component invisible to all of them. A lower estimate
for available gain relative to the entire current error remains essential.
The present theorem preserves arithmetic in its comparator and supplies a
uniform finite-block efficiency guarantee; it does not provide that lower
estimate, a fast computable replacement for $H_Q$, a bound for the earlier
smooth curvature rule, G2 or RH.

\section{Polynomial conditioning and an effective finite comparator (NS-81)}
\label{sec:ns81-effective-sampling}

\paragraph{Scope.}
The coefficient inverse in NS-78 makes the finite arithmetic comparator
in NS-80 effective. We prove an ordinary positive-Gram conditioning bound
and an explicitly sufficient finite sample cutoff. The cutoff is extremely
large; this is not a practical solver or a cofinal error estimate. None of
these bounds concerns a signed Weil operator or its spectral floor.

Let $H_N=\sum_{n=1}^N1/n$ in this section only, and let $G_N$ be the
ordinary Gram of $\sigma_1,\ldots,\sigma_N$. For a pure combination
$f=\sum c_n\sigma_n$, define the finite prefix form
\[
 P_N(c)=\eta^2+\sum_{m=1}^N\frac{f(1/m)^2+f(1/(m+1))^2}{m(m+1)},
 \qquad \eta=\sum_{n\leq N}c_n/n.
\]
It includes the sample at $1/(N+1)$. The same notation $P_N$ denotes
its matrix, and put
\[
 C_N=48H_NN^3(N+2).
\]

\begin{proposition}[A finite-prefix coefficient bound]
\label{prop:ns81-conditioning}
For every real coefficient vector,
\begin{equation}
 \|c\|_{\ell^2}^2\leq C_N P_N(c),\qquad
 \lambda_{\min}(G_N)\geq\frac1{16C_N}.
 \label{eq:ns81-conditioning}
\end{equation}
Writing $K_2=\|\sigma_1\|_2^2$, one also has
\[
 \lambda_{\max}(G_N)\leq K_2H_N,\qquad
 \operatorname{cond}_2(G_N)\leq768K_2H_N^2N^3(N+2)
                         =O(N^4\log^2(2N)).
\]
This is an ordinary spectral condition bound in the physical atom
coordinates, not a lower bound on approximation error.
\end{proposition}
\begin{proof}
Use $R=-f$ and $\alpha=0$ in NS-78; signs of the samples do not affect
$P_N$. Since $m\log(1+1/m)\geq1/2$, the exact inverse gives
\[
 \sum_{m=1}^N\frac{D_m^2}{m^2}
 \leq12\sum_{m=1}^N(z_m^2+z_{m+1}^2+\eta^2)
 \leq12N(N+2)P_N(c).
\]
Here each cell weight is at least $1/[N(N+1)]$ and $\eta^2\leq P_N$.
With $D_0=0$ and $\delta_m=D_m-D_{m-1}$,
\[
 \sum_{m=1}^N\frac{\delta_m^2}{m^2}
 \leq4\sum_{m=1}^N\frac{D_m^2}{m^2}
 \leq48N(N+2)P_N(c).
\]
M\"obius inversion and weighted Cauchy--Schwarz imply
\[
 |c_n|^2\leq\left(\sum_{d\mid n}d\right)
                  \sum_{d\mid n}\frac{\delta_d^2}{d}
 \leq nH_N\sum_{d\mid n}\frac{\delta_d^2}{d}.
\]
After summing and using $\sum_{k\leq y}k\leq y^2$ for $y\geq1$,
\[
 \sum_{n\leq N}|c_n|^2
 \leq H_N\sum_{d\leq N}\delta_d^2\sum_{k\leq N/d}k
 \leq H_NN^2\sum_{d\leq N}\frac{\delta_d^2}{d^2}
 \leq C_NP_N(c).
\]
The complete sampling form dominates its finite prefix, and NS-78
therefore gives $c^TG_Nc\geq P_N(c)/16$. This proves the ordinary
minimum-eigenvalue bound. The maximum is at most the trace,
$\operatorname{tr}G_N=K_2H_N$, by dilation and positivity. Division
of these two bounds proves the condition estimate.
\end{proof}

\begin{proposition}[Explicit complete tail and a finite comparator]
\label{prop:ns81-finite-comparator}
Let $Q^{<M}$ be the complete arithmetic sampling Gram of NS-80 truncated
to cells $1\leq m<M$, retaining the exterior term, where $M\geq N+1$.
Set
\[
 a_{N,M}=\tfrac12\log(2\pi M)+\frac{N+3}{6M},\qquad
 B_N(M)=\frac{2N}{M}\left(a_{N,M}^2+a_{N,M}+\tfrac12\right).
\]
Then the entire discarded sample tail satisfies
\begin{equation}
 0\preceq Q-Q^{<M}\preceq B_N(M)I
                  \preceq C_NB_N(M)Q^{<M}.
 \label{eq:ns81-tail-matrix}
\end{equation}
In particular the explicit integer choice
\begin{equation}
 M=2^{16}(N+1)^6
 \label{eq:ns81-explicit-cutoff}
\end{equation}
satisfies $C_NB_N(M)<1$, and hence
\[
 Q^{<M}\preceq Q\preceq2Q^{<M},\qquad
 \frac1{16}Q^{<M}\preceq G_N\preceq\frac{21}{8}Q^{<M}.
\]
The same constants survive any old-space Schur complement. For an
$N$-to-$2N$ block use~\eqref{eq:ns81-explicit-cutoff} with the total
atom count $2N$, not the old count $N$.
\end{proposition}
\begin{proof}
For $t\geq M\geq N$, the complete Bernoulli remainder in NS-78 gives
\[
 0\leq A(t/n)\leq\tfrac12\log(2\pi t)+\frac{N}{6M}
 \quad(1\leq n\leq N).
\]
On a cell $m\leq t\leq m+1$ with $m\geq M$, either endpoint's
atom value is bounded above by
$a_{N,M}+\tfrac12\log(t/M)$, since
$\log((m+1)/t)\leq1/M$. For a unit coefficient vector,
$\sum|c_n|\leq\sqrt N$, so both squared endpoint values are bounded
by $N(a_{N,M}+\tfrac12\log(t/M))^2$.
Integrating with $dt/t^2$ on every remaining cell gives $B_N(M)$,
including the complete tail to infinity. Finally
$Q^{<M}\succeq P_N\succeq I/C_N$, proving the last comparison.

For the displayed cutoff put $v=N+1\geq2$. Elementary bounds give
\[
 H_N\leq1+\log v\leq3v^{1/4},\qquad
 a_{N,M}\leq7+3\log v\leq13v^{1/4}.
\]
Here $\log v\leq2v^{1/4}$ follows by maximizing
$\log v/v^{1/4}$, whose maximum is $4/e<2$.
Also $\tfrac12\log(2\pi2^{16})<6.5$ and
$(N+3)/(6M)<1/2$, giving the bound for $a_{N,M}$.
Since $N^4(N+2)\leq v^5$, we obtain
\[
 C_NB_N(M)
 \leq\frac{96}{2^{16}}v^{-1}
       (3v^{1/4})\left(\frac{365}{2}v^{1/2}\right)
 =\frac{52560}{65536}v^{-1/4}<1.
\]
The full Gram comparison follows from NS-80. Taking old-coordinate
minima preserves the constants exactly as there.
\end{proof}

\paragraph{What is now finite, and what remains open.}
The finite comparator has projected condition ratio at most $42$.
The elementary NS-80 argument gives one-step capture at least
$168/1849$ of full available block gain and forty ideal steps capture
\[
 1-(41/43)^{80}>.9778.
\]
This is an effective finite-dimensional prescription; it does not make
that prescription fast. For a block with old size $256$, the total atom
count is $512$, and the sufficient cutoff exceeds $10^{21}$ cells.
No such enormous matrix is assembled in this update. The cutoff is a
conservative existence bound, not a necessary cutoff or a complexity
lower bound on every algorithm.

The scalar inequalities and positivity of the small finite-prefix matrices
are replayed with certified arithmetic. Those checks do not evaluate the
large finite comparator. Polynomial conditioning removes a numerical
resolution concern for these ordinary Grams, but leaves the missing
cofinal arithmetic decay completely unresolved. The original Weil floor,
G2 and RH remain open.

\paragraph{Finite validation.}
At $N=16,64,256$, the first $N$ cells and their final endpoint are
assembled exactly with Arb logarithms and integer factorial identities.
Their positive definiteness is supplied by the coefficient inverse.
A certified solve for $P_N^{-1}$ gives
$\operatorname{tr}(P_N^{-1})<C_N$ at each tested size, independently
confirming the displayed prefix floor; no midpoint solve is used.
The two 256/384-bit replays agree. At $N=256$ the analytic ordinary
Gram eigenvalue lower bound exceeds $4.91\cdot10^{-14}$, and the scalar
complete-tail ratio at~\eqref{eq:ns81-explicit-cutoff} is below $.01925$.
The latter is an evaluation of the proved tail bound, not an assembly
of the comparator's $18883333817345572864$ cells. Eighteen exact Maxima
checks validate the constants, primitives and rational step guarantees.

\appendix
\section{Reproducible finite checks}\label{app:numerics}

\subsection*{Complete first zero and archived resolution (v1.43)}
The incremental directory \texttt{evidence/v143/} contains the evaluator
comparison, $304$ earlier-zero exclusion intervals, resolution-threshold
certificates and saved $1024/1280$-bit reports. The unchanged v1.42
verifier and complete v1.25/v1.26 tail bounds are inherited by reference.
A second implementation checks the dual solve, exact dyadic coverage,
serialized signs, resolution algebra and cross-precision agreement.
The finite $N=120$ benchmark is separately labelled. No new external
audit or complete signed-discrepancy certificate is claimed.

\subsection*{One local complete-ground zero (v1.42)}
The directory \texttt{evidence/v142/} contains the energy-projection
proof, the $1024/1280$-bit even-separator and root gates, dependency
hashes, and a second implementation using principal determinants and
termwise cosine integrals. All $4097$ exact trial coefficients are
retained. The original complete residual bounds remain inherited
inputs; only their small final gates are replayed. The root interval
is $(\gamma_1-1/10,\gamma_1+1/10)$. No earlier-zero enumeration,
discrepancy sign, high-accuracy CCM transfer, or cofinal conclusion
is claimed. This pass has no new external audit.

\subsection*{Complete $\lambda=4$ ground ordering and real zeros (v1.40)}
The directory \texttt{evidence/v140/} contains the shifted-resolvent
proof, outward interval verifier, $1024/1280$-bit gate reports,
reproduction instructions and independent adversarial review.
The reports bind all inherited small-gate inputs by SHA256; full
original assemblies and witnesses remain under \texttt{evidence/v126/}.
The new computation recomputes the exact frozen trials' ordinary Grams,
replays the complete even margin and odd complement/directional gates,
and certifies the shifted lower-matrix inertia. It does not claim a new
full residual assembly. The finite even Rayleigh witness supplies the
negative shifted direction needed for the complete inertia conclusion.
The external real-zero implication uses the actual closed operator and
its core, not the finite compression or an endpoint approximation.

\subsection*{Scoped meta-obstruction and exact adversarial models (v1.39)}
The directory \texttt{evidence/v139/} contains the complete new proof,
its dependency and novelty report, an independent adversarial review,
and exact-rational model checks. The checks preserve a nonzero source
residual under a spectral swap and verify location-independent positive
examples; they are not arithmetic Weil certificates. The two theorems
are analytic and do not depend on floating-point diagnostics.
The complete source is compiled with \texttt{manuscript/build.sh};
the temporary ignored PDF is used only for the build gate and is not
part of delivery. The manuscript and evidence remain LaTeX/text only.

\subsection*{Complete additive-floor certificates (v1.38)}
The directory \texttt{evidence/v138/} contains the exact protocol,
unchanged archived coefficient assembly, the shifted verifier,
compressed frozen dyadic witnesses, initial and same-witness higher
precision certificates, fresh cutoff thresholds, and adversarial review.
The $256$-bit replay checks every complete lower-matrix row and every
trial quotient using Arb enclosures; floating eigensolvers and Cholesky
only propose exact dyadic witnesses. It separately verifies rational
upper bounds for the transformed head energy, giving the generalized
margin bounds in Proposition~\ref{prop:v138-three-floors}. No external
sign certificate, weighted positive tail metric or RH input is used.
The original v1.38 source was checked in draft mode. The current
complete manuscript uses the actual build gate described above.

\subsection*{Actual arithmetic scalar-budget obstruction (v1.37)}
The relocated directory \path{evidence/v137/} contains the proof,
candidate and literature audit, independent adversarial review, and a
standard-library probe diagnostic. Lemma~\ref{lem:v137-probe} proves its
constants analytically; the script separately checks the rational gate
$6(22/7)^2<100$, the negative probe mass and isolated-atom scaling.
Its floating-point quadrature is diagnostic, not an interval certificate.
Proposition~\ref{prop:v137-conditional-mass} uses RH only for its displayed
zero-sum asymptotic. Corollary~\ref{cor:v137-scalar-no-go} removes that
assumption through the v1.36 finite-floor implication and proves the two
divergent limits unconditionally. The strict prime endpoint, full exterior
archimedean term and all Fourier factors were checked independently.
No finite-window positivity calculation, numerical zero list, or PDF is
added in this revision.

\subsection*{Dense radicals and the finite-floor criterion (v1.36)}
Lemma~\ref{lem:v136-total-radicals},
Propositions~\ref{prop:v136-resolvent-collapse},
\ref{prop:v136-bounded-floor}, \ref{prop:v136-positive-comparison}, and
Corollary~\ref{cor:v136-complement-floor} are analytic results. The directory
\path{evidence/v136/} contains their proof excerpt, candidate and
literature audit, independent adversarial review, and exact countermodel
checks. The exact-rational script in that directory checks the support-consistent forms
$q_\pm[x]=\pm|\sum x_k|^2$. It verifies dense-radical approximation,
negative Rayleigh amplification, and the exact resolvent-error formula.
These models show why radical density and strong-resolvent convergence
alone do not determine sign. They are not arithmetic Weil certificates.
No finite-window numerical positivity result is added in this revision.

\subsection*{Growing radical blocks and the full parity symbol (v1.35)}
The analytic proofs are Propositions~\ref{prop:v135-both-parities} and
\ref{prop:v135-growing-radical}, and
Corollary~\ref{cor:v135-coercivity-obstruction}.  The accompanying directory
\path{evidence/v135/} contains the proof excerpt, novelty and
scope report, independent adversarial review, and
\texttt{check\_growing\_radical.py}.  Its double-precision diagnostics check
the Gaussian source moments and Poisson symmetry, and the complete pole
cancellation on complex mixed-parity tests at $\lambda=3,4,5$.
Omitting the negative odd pole produces a visible nonzero error in each
case.  A numerical estimate suggests $D=15$; this value is not certified
or used by the theorem, which uses its exact norm-based definition.
No growing-window numerical residual or finite-window sign certificate is
claimed.  The result uses every global prime row in the tail estimate,
including those above $\lambda^2$, and applies the ordinary Gram
normalization before claiming an operator bound.



The accompanying file \texttt{check\_sampler.py} constructs the real symmetric CCM matrix with the
separate logarithmic diagonal, using prime-power sums and numerical
quadrature only for the smooth archimedean integrals. The test is
$H(s)=\kappa(s)\zeta(s)$ with $\sigma=2$, $a=4$, $M=6$. Its global
zero-side form is exactly zero and $E=2\zeta(2)=3.289868133696453$.
No hypothetical zero location is inserted. At a fixed cutoff $N=512$:
\begin{center}
\begin{tabular}{rrrr}
\toprule
$L$ & Uncentered form & Centered form & Centered endpoint\\
\midrule
4 & 0.763119 & $1.174\,10^{-3}$ & 0.129293\\
6 & 1.036652 & $1.085\,10^{-4}$ & 0.033565\\
8 & 0.669071 & $4.000\,10^{-6}$ & 0.008042\\
10 & 1.337515 & $4.078\,10^{-7}$ & 0.001816\\
\bottomrule
\end{tabular}
\end{center}
For $L=6$, endpoint correction is exact to the displayed floating-point
precision, while its form price decreases and derivative price increases:
\begin{center}
\begin{tabular}{rrrr}
\toprule
$N=2K$ & Corrected form & Correction norm & Derivative norm\\
\midrule
128 & 1.400023 & 0.705010 & 72.5381\\
256 & 0.906192 & 0.498518 & 102.3284\\
512 & 0.492218 & 0.352505 & 144.5331\\
\bottomrule
\end{tabular}
\end{center}
These moderate cutoffs do not yet make the corrected form small. The
existence and rates in Theorem~\ref{v1:thm:repair} come from the proof, not
from extrapolating this table. Increasing the quadrature order on a
representative matrix changed entries by less than $2\,10^{-12}$.
The script uses 30-digit zeta evaluations followed by double-precision
matrix arithmetic; the checks are not interval-certified bounds.

\subsection*{A rational check of the resolvent endpoint asymptotic}
The file \texttt{check\_resolvent\_endpoints.py} tests the sign, half-period
phase, shell contribution and residue scale in
Lemma~\ref{lem:v11-resolvent-endpoint} using explicit rational functions:
\[
 H_-(s)=\frac{(\sigma+a)^M}{(\sigma-s)(s+a)^M},
 \qquad H_+(s)=(s-p)H_-(s).
\]
Here $\sigma=2$, $a=4$, $M=6$, $\mu=0.2+1.3i$ and
$p=1/2-\overline\mu$. These are model transforms, not zeta root
functions. Their Cauchy tail coefficients are $E_-=1$ and
$E_+=\sigma-p$. The centered samples, their adjoint resolvents and
the endpoint shell correction are summed explicitly in Fourier pairs.
For numerical tractability this check uses
$K=\lceil e^{L/2}L^2\rceil$, much smaller than the theorem cutoff.
It does not compute a Weil form.

Write $e_\pm=\delta V_LH_\pm$, interpreting $V_L$ here as the same
sampling and resolvent operations applied directly to these transforms.
The predicted reflected limit is
\[
 S_L:=\frac{e^{\overline\mu L/2}(1-e^{-\overline\mu L})e_-}
                  {-H_-(p)}\longrightarrow1.
\]
The output, rounded from double-precision arithmetic, is:
\begin{center}
\begin{tabular}{rrrr}
\toprule
$L$ & $K$ & $|S_L-1|$ & $|e_-/e_+|$\\
\midrule
4  & 119    & $7.347\,10^{-2}$ & 4.426\\
6  & 724    & $5.792\,10^{-3}$ & 35.261\\
8  & 3495   & $9.638\,10^{-4}$ & 282.874\\
10 & 14842  & $2.276\,10^{-4}$ & 1869.438\\
12 & 58094  & $7.847\,10^{-5}$ & 5886.940\\
14 & 214941 & $1.872\,10^{-5}$ & 24924.379\\
\bottomrule
\end{tabular}
\end{center}
This is a reproducible check of the analytic mechanism and conventions,
not an interval certificate, a verification on hypothetical zeta zeros,
or evidence for RH. The theorem and its application to genuine root
functions depend on the contour and uniformity arguments in the text.
\subsection*{A high-precision check of the normalized prolate source}
The script \texttt{check\_prolate\_source.py} constructs the even
Legendre--Galerkin matrix for
$-\partial_x((1-x^2)\partial_x)+c^2x^2$ on $[-1,1]$.
It takes the first and third even eigenvectors, corresponding to indices
$0$ and $4$, and normalizes them in $L^2(-1,1)$. The compressed Fourier
eigenvalue is recovered from the ratio of the integral to the central
value, with the factor $\sqrt{c/(2\pi)}$ retained. Thus the source
coefficients are those in \eqref{eq:v12-coefficients}, not an imposed
endpoint fit.

Let $S_\lambda$ denote the noncancellation ratio on the left of
\eqref{eq:v12-noncancellation}, and let
$Z_\lambda=|h_\lambda^\circ(0)|/|B_\lambda^+|$. With 50 even Legendre
modes and 85-decimal-digit arithmetic the rounded output is:
\begin{center}
\begin{tabular}{rrrr}
\toprule
$c$ & $b_4/(c^2b_0)$ & $S_\lambda$ & $Z_\lambda$\\
\midrule
10 & 6.40490  & 1.00164349 & $1.92330\,10^{-2}$\\
20 & 9.94110  & 1.00029407 & $4.00922\,10^{-6}$\\
30 & 10.99058 & 1.00012034 & $3.99136\,10^{-10}$\\
40 & 11.51074 & 1.00006514 & $3.10538\,10^{-14}$\\
\bottomrule
\end{tabular}
\end{center}
Repeating $c=40$ with 65 modes and 105-digit arithmetic preserves the
displayed figures. The script also reports the concentration loss divided
by Fuchs's leading term and the ratio $b_n^2/[c(1-\vartheta_n)]$.
Increased precision is important because concentration losses become
exponentially small. This is a finite spectral approximation with no
certified truncation or roundoff bound; the analytic proof, not this
table, establishes the asymptotic estimates. No zeta zero is inserted.

\subsection*{A compact-source primitive convention check}
The script \texttt{check\_bv\_radical\_passage.py} tests the powers and
sign in the primitive mechanism of Theorem~\ref{thm:v14-radical}.
Take $h(u)=u^2-\frac53u^4$ on $[-1,1]$, with zero extension. It has
the two exact moments and $h(1^-)=-2/3\ne0$. If $N=\lfloor1/u\rfloor$
and $S_j(N)=\sum_{n=1}^N n^j$, an explicit logarithmic primitive is
\[
 P(u)=\frac25u^{5/2}S_2(N)-\frac{10}{27}u^{9/2}S_4(N)
       +\frac4{135}\zeta(\tfrac12,N+1),
\]
where the last function is the Hurwitz zeta function. Away from a cut,
$uP'(u)=\mathcal E(h)(u)$; the jump cancellation in this expression
agrees with continuity of the primitive. At 70-digit working precision,
testing $1/u=N+\theta$ for $N=10,100,1000,10000$ and
$\theta=0.2,0.5,0.8$ gives maximum derivative discrepancy less than
$3.3\,10^{-69}$. The sampled values of $P(u)/u^{3/2}$ remain bounded.
This is a toy-source convention check, not an interval certificate or
a numerical proof of the uniform primitive bound, radicality, G1, or RH.
Those conclusions use the analytic arguments in the text.


\subsection*{Finite-core metric and projection checks}
The script \texttt{check\_section19\_metric.py} checks the exact budget
\eqref{eq:v15-budget}, source/shell annihilation,
\eqref{eq:v15-variance}, and the new projected graph source
\eqref{eq:v15-projected-graph}. It uses the rational source samples
\[
 z_n=\frac{(-1)^n}{\sqrt L(2+it_n)^3},\qquad t_n=2\pi n/L,
\]
with a length-two Jordan chain at $\mu=0.23+0.7i$ and a further
eigenvector at $\nu=0.37+1.4i$. These parameters are model eigenvalues,
not zeta zeros. The tests include a complex endpoint-shell correction and
the first-slot-linear convention. At $T=3$ the rounded results are
\begin{center}
\begin{tabular}{rrrr}
\toprule
$L$ & $\lambda_{\min}(J^*SJ)$ & $\|[D,S]\|$
 & $|\mathcal L_{11}|/G_{11}$\\
\midrule
8 & $2.55980\,10^{-4}$ & 0.976536 & 0.460000\\
16 & $8.40732\,10^{-4}$ & 0.997413 & 0.460000\\
32 & $8.95576\,10^{-4}$ & 1.007014 & 0.460000\\
\bottomrule
\end{tabular}
\end{center}
The nine algebra-check discrepancies are below $2.1\,10^{-15}$;
the source-polynomial remainder identity is also checked on a grid.
The final column equals $2\Re\mu$, as the exact identity predicts:
source annihilation does not make the relative Lyapunov defect small.
No Weil matrix or zeta zeros enter this model. The calculations are
floating-point convention checks, not positivity certificates or proofs
of any limiting estimate; Section 19 supplies the analytic arguments.

\subsection*{Sharp-band transform convention checks}
The script \texttt{check\_band\_cut.py} uses the rational source
$Z(t)=(2+it)^{-3}$ and $R(t)=(it-\mu)^{-1}$ with
$\mu=0.23+0.7i$. It subtracts the exact sampled source projection,
sets $T=3$, and takes $L=2\pi N/3$, $N=8,16,32,64,128,256$.
The second-difference identity in Lemma~\ref{lem:v16-band-transform}
is checked to error below $3\,10^{-15}$, and direct physical integration
at the smallest size agrees to below $10^{-14}$.
At $N=128$ the ordinary squared norm is approximately $0.0156032$,
while the transform magnitudes at $s=1$ and $s=0$ are approximately
$1.088\,10^{24}$ and $6.525\,10^{23}$, respectively. These large
values are consistent with the proved cut mechanism, not evidence about
zeta zeros or the sign of a Weil form. The experiment evaluates no Weil
matrix and supplies no certified numerical error bound. It also checks
the transform product at $v=\mu$ and $v^\sharp=-\overline\mu$ against
the edge expression in Proposition~\ref{prop:v17-pair-edge}. At $N=256$
the relative discrepancy is approximately $0.0148$ and the real pair sum
is approximately $1.414\,10^{43}$; this is a rational-model convention
check, not a zeta-zero or whole-Weil-form computation.

\subsection*{Full-space projection and holomorphic strip checks}
The file \texttt{check\_full\_projection.py} tests
Proposition~\ref{prop:v18-full-projection} with the rational model
\[
 Z(v)=\frac{(v-\mu)^2(v-\nu)}{(v+2)^8(2-v)},\qquad
 \mu=0.23+0.7i,\quad\nu=0.37+1.4i.
\]
The model root outputs $Z(v)/(v-\mu)$, $Z(v)/(v-\mu)^2$ and
$Z(v)/(v-\nu)$ are holomorphic on the closed critical strip;
no meromorphic-output assumption is hidden in this check. At
$L=8,12,16,24,32$ it uses the tractable cutoff $K=4L^2$, not the
theorem's exponential cutoff. Eight finite algebra checks, including
the complete graph source, both shell directions and the commutator
budget, have discrepancies below $4.2\,10^{-14}$.
On a 25-point grid in the strip, the largest ordinate-weighted transform
error decreases from approximately $0.02043$ at $L=8$ to
$4.121\,10^{-6}$ at $L=32$. The model first-eigenvector relative
Lyapunov defect remains $0.46=2\Re\mu$. No zeta zero or Weil matrix
enters this experiment. These are floating-point convention checks,
not uniform strip certificates or evidence for RH.

\subsection*{An arithmetic-matrix check of the Weil action}
The script \texttt{check\_weil\_action.py} uses the actual finite
arithmetic matrix from \texttt{check\_sampler.py}, including the
separate logarithmic diagonal, and the first known critical-line zero
$\rho_1\simeq1/2+14.1347251417i$. The transform is
$\kappa(s)\zeta(s)/(s-\rho_1)$ with $\sigma=2$, $a=4$, $M=6$.
The root location and derivative are evaluated at 40-digit precision,
then samples and matrices use double precision. This is a numerical
simple-root check, not a new multiplicity certificate. After the
full-space projection, it compares $WSJF$ with the exact finite
Fourier projection of
$\kappa(\rho_1)\zeta'(\rho_1)e^{(\rho_1-1/2)y}$.
At the moderate cutoff $K=256$, $N=2K=512$:
\begin{center}
\begin{tabular}{rrr}
\toprule
$L$ & Absolute action error & Relative action error\\
\midrule
4 & $7.9105\,10^{-4}$ & $5.7219\,10^{-2}$\\
6 & $2.3776\,10^{-4}$ & $1.4045\,10^{-2}$\\
8 & $4.1955\,10^{-5}$ & $2.1459\,10^{-3}$\\
10 & $1.5078\,10^{-5}$ & $6.8992\,10^{-4}$\\
\bottomrule
\end{tabular}
\end{center}
The script also checks $K=128$, the target Fourier coefficient by
independent physical quadrature, and the last matrix action with 400
additional archimedean quadrature nodes. These finite cutoffs do not
meet the theorem's exponential rule. No hypothetical off-line zero
is inserted, no interval certification is claimed, and the table is
not evidence for RH. It checks the action's sign, centered phase and
physical normalization using an actual arithmetic matrix.

\subsection*{Signed Fredholm moment check}
The accompanying \texttt{check\_fredholm\_shell.py} uses Gauss--Legendre
Nystr\"om matrices for the actual kernel $2\cos(2\pi xt)$ on $(0,a)$.
It compares the two-channel inverse, the signed $A_a$ formula and the
constant-kernel resummation, using complex test profiles at
$a=0.2,0.1,0.05,0.025$ and two quadrature orders. It separately checks
that reversing one profile reverses the mixed term. The profiles are
arbitrary interval data, not certified Sonine vectors or arithmetic
evaluators. These floating-point checks test signs, scale factors and
the remainder formula only; they do not test the arithmetic shell sign,
certify a uniform estimate, or supply evidence for RH.

\subsection*{Regularized co-Poisson adjoint check}
The script \texttt{check\_copoisson\_adjoint.py} uses
$\phi(t)=(t-1)^3(2-t)^3$ on $[1,2]$, extended by zero, and
$h=\mathcal D\phi$. Both source moments vanish exactly. This toy
source has sufficient piecewise smoothness for the cutoff estimates;
it is not asserted to be a smooth arithmetic-core vector or an evaluator.
At 45-digit precision, independent physical and adjoint integrals are
compared for $y(x)=e^{-x}$ and $M=4,8,16,32$. The weak limit is checked
against its convergent zeta-coefficient series. A separate power-response
check uses $w=0.7+1.1i$, not a zeta zero, and $M$ up to 1024. These
checks include $y(x)=e^{-(1+0.4i)x}$ to test the first-slot conjugation.
The
non-certified computations check counterterm signs and scales, not
the shell inequality or RH. Results are recorded in the supplied JSON.

\subsection*{A certified finite Weil matrix and candidate angle}
\label{app:v112-certificate}
The following finite result is certified by interval arithmetic; the
earlier numerical illustrations remain non-certified. Set
$\lambda=3$, $L=2\log3$, $N=64$ and retain all $129$ indices
$-64\le n\le64$ in the translated $[0,L]$ basis $U_n$.
The matrix $W$ is the same arithmetic Weil matrix used above, with all
prime powers $2,3,4,5,7,8,9$, the pole term and the separately computed
logarithmic diagonal. No list of zeta zeros is used.

The candidate is defined exactly, rather than by an accuracy claim:
the $65$ decimal strings in \texttt{g2\_finite\_candidate.json} are
interpreted as rational numbers $c_0,\ldots,c_{64}$, with
$c_{-n}=c_n$, and $u=c/\norm c$. The file's SHA-256 is
\begin{center}\ttfamily\small
6d326cafbe715752f89b90cd92eb1d5a68\par
ab0b78befb6afe074e797372b1a18e
\end{center}
The two lines concatenate to the identifier. The strings were generated
from a $105$-digit, $70$-even-mode approximation of the repaired source,
but the conclusions below depend only on their exact rational values.
This candidate definition alone does not certify proximity to the true
PSWF coefficients; Appendix~\ref{app:v113-source} supplies that bound.

\begin{proposition}[Computer-assisted finite result]
\label{prop:v112-finite-certificate}
For this exact $W$ and rationally specified normalized candidate,
the complement in Proposition~\ref{prop:v112-schur-angle} satisfies
\[
 C-\alpha I\succ10^{-34}I,\qquad
 \norm{(C-\alpha I)^{-1}r}<0.000216.
\]
The matrix is positive definite and has a simple even ground state, with
\begin{equation}\label{eq:v112-certified-bottom}
 3.64\,10^{-38}<\lambda_0<5.32\,10^{-38},
 \qquad \sin\theta<0.000216.
\end{equation}
\end{proposition}
\begin{proof}[Certificate and verification]
The supplied \texttt{certify\_g2\_finite.py} uses
Python-FLINT~0.9.0 and FLINT~3.6.0 ball arithmetic
\cite{PythonFlint2026} at $512$ bits. It encloses all archimedean
integrals, uses both exact reflection sectors, and constructs the
complement by the Householder transformation determined by the exact
normalized candidate. Orthogonality belongs to that underlying exact
matrix, not to arbitrary independent choices from its entrywise balls.
Interval $LDL^T$ elimination gives $64$ strictly positive pivot
enclosures for each of
\[
 C_{\rm even}-\alpha I-10^{-34}I,
 \qquad W_{\rm odd}-\alpha I-10^{-34}I.
\]
An enclosed LU solve gives $q<0.000216$ and verifies, by strict
interval comparisons,
\[
 \alpha<5.32\,10^{-38},\qquad
 \alpha-\ip{(C-\alpha I)^{-1}r}{r}>3.64\,10^{-38}.
\]
Proposition~\ref{prop:v112-schur-angle} proves the claims.

The independent program \texttt{certify\_g2\_series.py} repeats the
certificate at $768$ bits using the analytic series below, with $128$
terms and explicit tail bounds. It uses a rational, nonorthonormal
complement in the full $129$-dimensional space. For the unnormalized
rational vector $c$, let $s=c^Tc$ and choose $c_k\ne0$. Columns
$Q_j=e_j-(c_j/c_k)e_k$, $j\ne k$, span $c^\perp$ exactly. Put
\[
 G=Q^TQ,\quad H=Q^T(W-\alpha I)Q,\quad b=Q^TWc.
\]
All $128$ interval pivots of $H-10^{-34}G$ are positive.
For the enclosed solution $Hz=b$, the physical quantities are
$q^2=z^TGz/s$ and $E=b^Tz/s$. The Gram matrix $G$ is retained.
This independently verifies the same three strict inequalities.
All $8321$ parity-block entries also have overlapping enclosures in the
two independently assembled matrices. The supplied JSON files retain
the pivot and result intervals, input hash, versions and precision.
No approximate floating-point solve is used for either certificate.
\end{proof}

For reproducibility, write $\rho_\infty(y)=e^{y/2}/(2\sinh y)$,
$t_n=2\pi n/L$ and let $\psi_0,\psi_1$ denote digamma and trigamma.
The first calculation removes the apparent singularity at zero using
$\operatorname{sinc}z=\sin z/z$ and
$\sinh y/y=\operatorname{sinc}(iy)$. The second uses the following
exact formulas for the sine integral and the \emph{complete}
archimedean diagonal:
\begin{align*}
 S_n&:=\int_0^L\rho_\infty(y)\sin(t_ny)\,dy
 =-\tfrac12\Im\psi_0(z_n)
   -\sum_{k\ge0}e^{-a_kL}\frac{t_n}{a_k^2+t_n^2},\\
 A_n^\infty&:=\int_0^L\rho_\infty(y)
  \{2(1-y/L)\cos(t_ny)-2e^{-y/2}\}\,dy\\
 &\hspace{15mm}+\log(4\pi)+\gamma_{\!E}+\log\tanh(L/2)\\
 &=\log\pi-\Re\psi_0(z_n)-\frac{\Re\psi_1(z_n)}{2L}
   +\frac2L\sum_{k\ge0}e^{-a_kL}
          \frac{a_k^2-t_n^2}{(a_k^2+t_n^2)^2},
\end{align*}
where $a_k=2k+\tfrac12$ and $z_n=\tfrac14-it_n/2$.
To verify them, expand $\rho_\infty(y)=\sum_{k\ge0}e^{-a_ky}$.
Integrate the convergent combinations
$\rho_\infty\sin(t_ny)$,
$\rho_\infty(\cos(t_ny)-e^{-y/2})$, and
$y\rho_\infty\cos(t_ny)$ on $(0,\infty)$ and subtract their tails.
The identity $t_nL=2\pi n$ simplifies the endpoint terms. The odd
exponential tail cancels $\log\tanh(L/2)$, while
$\psi_0(1/2)=-\gamma_{\!E}-2\log2$ leaves $\log\pi$.
The digamma difference must not be split into two divergent sums.
After retaining $k<K$, the absolute tail bounds are
\[
 R_S\le\frac{e^{-a_KL}}{2a_K(1-e^{-2L})},\qquad
 R_A\le\frac{2e^{-a_KL}}{La_K^2(1-e^{-2L})}.
\]
These follow from $|t|/(a^2+t^2)\le1/(2a)$ and
$|a^2-t^2|/(a^2+t^2)^2\le1/a^2$. At $\lambda=3$ the weights are
the exact rationals $3^{-4k-1}$; both tails are enclosed symmetrically.

\paragraph{Scope and endpoint limitation.}
This certificate concerns one finite matrix and one explicitly frozen
candidate. It gives no lower bound on the omitted Fourier complement,
no growing-window ordering theorem. Its fixed-source approximation
prerequisite is addressed in Appendix~\ref{app:v113-source}. The preparatory non-certified source pilot gives relative
physical-endpoint errors about $1.375$ at $N=64$ and $0.1046$ at
$N=4096$ for $\lambda=3$; the latter calculation concerns source
coefficients only, not a matrix of that dimension. Small ordinary
ground-state angle therefore cannot substitute for endpoint recovery.
An endpoint-shell correction would still require its derivative and
graph terms. The full-strip transform cost and an independent uniform
arithmetic estimate remain open. Neither G2 nor RH follows from this
finite certificate.


\subsection*{Certified exact-source identification and endpoint error}
\label{app:v113-source}
The additional verifier \texttt{certify\_pswf\_source.py} upgrades the
source comparison at this one window from numerical agreement to a
computer-assisted proof. Its input \texttt{pswf\_rational\_input.json}
contains two $70$-coordinate exact decimal rational vectors and two
exact decimal rational approximate eigenvalues. Its SHA-256 is
\begin{center}\ttfamily\small
ac9dfa8c03f54a326e3ba7af59917ac435e\par
63d41e978310301e6df6e53e83076
\end{center}
The two lines concatenate. Approximate eigensolver output is used only
to select this input; its accuracy is not assumed. Python-FLINT~0.9.0 /
FLINT~3.6.0 at $768$ bits verifies the following steps.

In the normalized Legendre basis,
$Xe_\ell=A_\ell e_{\ell+1}+A_{\ell-1}e_{\ell-1}$ with
$A_\ell=(\ell+1)/\sqrt{(2\ell+1)(2\ell+3)}$ and $A_{-1}=0$.
The even Jacobi diagonal and next diagonal are
\[
 2k(2k+1)+c^2(A_{2k}^2+A_{2k-1}^2),\qquad
 c^2A_{2k}A_{2k+1},\qquad c=18\pi.
\]
The omitted diagonal lower bound is $140\cdot141=19740$.
For mode zero the two bounding matrices in
\eqref{eq:v113-inertia} both have inertia $0$ at $\xi_0-1$ and
$1$ at $\xi_0+1$. For mode four their counts are $2$ and $3$.
Every interval $LDL^T$ pivot has a certified nonzero sign.
Thus $g=1$ is a valid separation from each $\xi_j$ to all other
exact even eigenvalues. The full residuals include the first omitted
link, which dominates here. The following rounded numbers are strict
upper bounds, not error estimates inferred from precision settings:
\begin{center}
\begin{tabular}{@{}lrr@{}}
\toprule
Quantity & Mode $0$ & Mode $4$\\
\midrule
Full residual $\epsilon$ & $5.913\,10^{-56}$ & $4.796\,10^{-52}$\\
Angular $L^2$ error & $8.362\,10^{-56}$ & $6.783\,10^{-52}$\\
Angular uniform error & $2.723\,10^{-52}$ & $2.509\,10^{-48}$\\
\bottomrule
\end{tabular}
\end{center}
The approximate physical central values exceed $1.187$ and $0.716$,
respectively, by interval comparison with their errors, fixing the true
sign convention $h_{j,3}(0)>0$.
The error budget gives
\[
 \eta<1.725\,10^{-48},\quad
 |\widehat h^\circ(0)|<2.611\,10^{-39},\quad
 U<3.368\,10^{-37}.
\]
The exact moment repair is bounded by $\norm\psi_\infty\le129$;
no uncertified evaluation of its defining integral ratio is inserted.

The polynomial Fourier coefficients are also evaluated with enclosures.
If $\widehat h^\circ(v)=\sum_r d_r(v/\lambda)^{2r}$ and
$t_n=2\pi n/L$, the translated coefficients of its inversion average
are exactly
\begin{equation}\label{eq:v113-mellin-polynomial}
 \widehat c_n=\frac1{\sqrt L}\Re\sum_r d_r
 \frac{\displaystyle\sqrt\lambda\sum_{m=1}^{9}m^{-1/2-it_n}
       -\displaystyle\lambda^{-4r-1/2}\sum_{m=1}^{9}m^{2r}}
      {2r+1/2+it_n}.
\end{equation}
To check all scale factors, integrate each monomial on
$-\log\lambda<x<\log(\lambda/m)$ with the sole co-Poisson factor
$e^{x/2}$. Translation from the centered interval contributes
$(-1)^n$, which cancels $e^{\pm it_n\log\lambda}$ at the integration
limits. Taking the real part performs inversion averaging. At $m=9$
the interval has zero length, so its two terms cancel exactly; this
integral endpoint convention differs from the strict trace in
\eqref{eq:v113-endpoint-budget}.

The interval polynomial coefficient vector differs from the frozen
rational vector by less than $6.422\,10^{-40}$ in ordinary norm.
Adding \eqref{eq:v113-projection-budget} gives a bound below
$2.595\,10^{-36}$, in particular the slightly looser strict threshold
$2.60\,10^{-36}$ in Proposition~\ref{prop:v113-source-certificate}.
The normalized projector distance is below $4\,10^{-36}$; no division
by the tiny Weil spectral gap is used in this final source replacement.
The saved interval sine bound is less than $0.000215893$.

Direct Legendre evaluation at $1$ and $m/9$, $1\le m\le8$, together
with \eqref{eq:v113-endpoint-budget}, encloses $B_3$.
Physical evaluation of the projection is
$L^{-1/2}(c_0+2\sum_{n=1}^{64}c_n)$, and its propagated error is
at most $\sqrt{129/L}$ times the coefficient error. Strict interval
comparisons verify \eqref{eq:v113-endpoint-failure}. Thus the endpoint
mismatch is now proved for the exact repaired source, not just observed
for a numerical approximation.

The JSON output records the input hashes, inertia counts, full residual
and pointwise bounds, coefficient error, physical endpoint and relative
error intervals. The earlier frozen-candidate certificate remains a
separate dependency. These two source modes use the canonical regular
angular realization; no closed arithmetic operator assumption is needed
for this fixed-source computation. This source certificate alone supplies
no omitted \emph{Weil} Fourier-complement bound, growing-window sign,
endpoint-shell graph estimate or full-strip normalization. G2 remains open.

\subsection*{Scale check for the explicit omitted-tail constant}
The new operator-core and polynomial-endpoint arguments are analytic
proofs, not extrapolations from numerical data. A small additional
interval calculation evaluates the explicit constant in
\eqref{eq:v114-tail-constant} at $\lambda=3$. Using 256-bit Arb
quadrature of the regularized zero-mode archimedean integral gives
\[
 17.69849<C_3^{\rm tail}<17.69851,\qquad
 \Gamma_{3,50\,000\,000}>1.0797.
\]
The supplied script \texttt{g2\_deep\_next/check\_tail\_constant.py}
and its output record the complete intervals. No matrix of that size
was assembled: only the scalar constant was enclosed. This bound
illustrates the cost of the unsigned arithmetic estimate. It does not
certify the intervening low modes, their Schur correction, or positivity
of the whole window. In particular the earlier $N=64$ certificate is
not extended to the continuum by this calculation.

\subsection*{Sharper tail certificate and ordinary residual scope}
The script \texttt{g2\_signed\_tail/check\_sharp\_tail\_constants.py}
evaluates the exact elementary constants of
Proposition~\ref{prop:v115-sharp-tail} with 256-bit Arb intervals.
Its output gives
\[
 \Gamma^{\rm all}_{3,256}>0.0148329,\qquad
 \Gamma^{\rm all}_{3,512}>0.7085077,\qquad
 \Gamma^{\rm even}_{3,256}>1.5867887.
\]
The interval arithmetic certifies these scalar evaluations of analytic
bounds. That v1.15 calculation did not certify a head matrix or its
inverse-action correction; the complete certificate below now does so
at the same physical window. The ordinary residual theorem is an analytic
asymptotic estimate with no claimed effective threshold at $\lambda=3$.
It concerns the same exact source and literal Fourier projection as
the endpoint theorem. An independent internal adverse review checked
the norm normalization, low-output duality, high-output matrix bound,
graph tail, all scale factors and the remaining lack of ground overlap.

\subsection*{Complete fixed-window certificate and its reproduction}
Proposition~\ref{prop:v116-window-positive} is certified by
\texttt{g2\_schur\_directional/certify\_infinite\_schur.py}, using
the coefficient assembly in \texttt{assembly.py} and the scalar check
\texttt{certify\_weighted\_prime\_tail.py}. The bundle contains both
exact dyadic inverse-action witnesses, all positive LDL pivot enclosures,
and 768-bit and 896-bit replay reports. Its reproduction guide gives
the complete commands and analytic dependencies. The frozen witnesses
are unchanged between the two precision runs.

All finite output rows through $J=4096$ are enclosed, and the moment
inequality bounds every row beyond $J$. Thus the result is not an
extrapolation from finite eigenvalues. The finite Fourier entries through
$N=64$ were also checked against the earlier independent interval
integration: all 8,321 even and odd entries overlap their reference
enclosures. This consistency check supplements, but does not replace,
the analytic coefficient identities and series-remainder proof.

The unshifted result certifies complete fixed-window positivity.
The shifted extension below also establishes ground-state ordering.
Neither certificate alone supplies an endpoint comparison or a bound
for unbounded windows. The separate boundary theorem now determines
the continuum eigenfunction trace to be zero. Internal adversarial review checked the inverse order, dyadic
freezing residual, parity factors, remote Gram and form-domain passage;
this is not external peer review or proof-assistant verification.
Exploratory failures of weaker sufficient bounds are preserved in the
research log and are not presented as negative Weil-form results.

\subsection*{Shifted inertia and complete ground-overlap reproduction}
Proposition~\ref{prop:v117-ground-order} uses
\texttt{g2\_ground\_order/certify\_shifted\_inertia.py} with
$\texttt{shift}=10^{-34}$ in the even sector and $10^{-36}$ in the
odd sector. The bundle retains the signed pivot reports at 768 and
896 bits. It reuses the unchanged exact witnesses from
\texttt{g2\_schur\_directional}; their identity is checked by hash.
The script shifts both the actual logarithmic diagonal and all inverse
comparison weights, recomputes the middle residual and retains the
full remote moment bound. A negative lower-matrix pivot is not itself
a negative direction of the complete operator: the proved trial quotient
supplies that direction for the shifted even operator.

The separate \texttt{certify\_trial\_reuse.py} checks the exact rational
trial's full-form Rayleigh quotient, constructs the finite-supported
correction and imports the proved exact-source error from
Proposition~\ref{prop:v113-source-certificate}. Its corrected Rayleigh
enclosure is contained in
\[
 (3.64339965699233062,3.64339965699233064)\cdot10^{-38}.
\]
The separate scalar script also verifies the conservative constants
displayed in the theorem.
The complete ground-overlap conclusion uses the independently proved
even-sector spectral threshold. It does not reuse the finite matrix's
gap as if it were a gap for the complete operator.
All reproducibility commands and unsuccessful exploratory tests are
recorded separately in the guide and research log.

\subsection*{Analytic endpoint and complement checks}
The proof of Theorem~\ref{thm:v118-eigen-endpoint-zero} is analytic.
It does not use a numerical endpoint extrapolation. Its three gates are
the exact physical kernel identity with scalar $-\log(2\pi)$, a common
closed-form core, and the simultaneous $L^2/L^\infty$ resolvent argument
that makes the boundary theorem applicable. The archived independent
reviews check each gate, the retained pole signs and the distinction
between a restricted logarithmic Laplacian and a spectral logarithm.
The script \texttt{check\_log\_kernel\_normalization.py} independently
compares the original and decomposed constant-mode archimedean diagonal
at $\lambda=1.2,2,3,10$, with 90-decimal working precision. Agreement is
a diagnostic only; the scalar integral in the proof establishes the
normalization exactly.

The Fej\'er bound retains the physical $L^{-1/2}$ Fourier factor and
uses a positive kernel of total mass one. It does not replace the
paper's literal Fourier cutoff in any earlier formula.
Proposition~\ref{prop:v118-graph-trace} explicitly shows why graph
convergence does not transfer the physical trace.
The complement estimate in Proposition~\ref{prop:v118-gap-free}
is an exact Hilbert-space inequality with a sharp two-dimensional
example. Its use of the rank-one source invokes the proved ordinary
operator residual; its arithmetic complementary-sign hypothesis remains
open. These checks close the continuum endpoint question and improve
the formulation of the growing-window target, without closing G2.

\subsection*{Source-weighted capacity candidate audit}
The original \texttt{g2\_capacity\_candidate} companion directory is
\emph{not archived}: its candidate register, novelty comparison,
adversarial review and recomputation outputs are absent. The complete
analytic proof of Proposition~\ref{prop:v129-capacity} is printed in this
manuscript and does not depend on that missing review. The finite
projected-source values following it were historically reported as
recomputed from the available $\lambda=3$, $N=64$ interval midpoints;
the recomputation script and output are not archived. Those values are
falsification diagnostics, not interval certificates or proof inputs.
The historical review was described as checking the archimedean constant,
prime orientation factors, exterior killing terms, pole sign and the
two-coordinate counterexample. This is a description of an unavailable
review, not a claim of current independent verification. The arithmetic
capacity estimate, its cofinal error and the odd sector remain open.

\subsection*{Weighted-concentration and Schatten no-go audit}
The original \texttt{g2\_weighted\_concentration} and
\texttt{g2\_schatten\_no\_go} evidence sets are \emph{not archived}.
Only \path{evidence/v131/g2_weighted_concentration_report.md} survives,
with its recovery provenance. The original scripts, outputs, candidate
register and complete adversarial reviews are absent. The analytic
Hilbert--Schmidt, interlacing, dilation and positive flat-top arguments
are proved in the manuscript and do not rely on those diagnostics.

The original \texttt{certify\_beta4\_negative.py} and its output are
\emph{not archived}. The historical account reports exact rational
interval arithmetic proving $\beta_{\log4}(1)<-3/5$; that numerical
claim cannot presently be replayed. Proposition~\ref{prop:v132-global-index}
now explicitly separates the reported enclosure from the analytic
negative-index implication conditional on it. The implication concerns
unrestricted Schwartz tests, not the compact physical window.

Finite cosine-compression and trapezoidal-flat-top tests were reported
only for $\lambda=3,4,5$. Their original scripts and outputs are not
archived; they were diagnostics, not certificates or cofinal evidence.
The signed weighted-superlevel implication remains analytic and
conditional; its cofinal arithmetic constants remain open.

\subsection*{Continuum-cancellation and local no-go audit}
The original \texttt{g2\_continuum\_cancellation} proof report and
\texttt{test\_continuum\_symbol.py}, including its saved output, are
\emph{not archived}. The complete analytic proofs of
Propositions~\ref{prop:v130-continuum-cancellation}--\ref{prop:v130-local-no-go}
are printed in this manuscript. Kernel factorization, the positive
Fourier multiplier and the six-bump algebraic argument require no
numerical diagnostic or absent companion review.
The missing script was reported to scan $\lambda=3,4,5$ as a
falsification diagnostic, not a sign certificate. The exact normal
form retains the physical support constraint; arbitrary frequency
localization is not substituted for it.

\subsection*{Prime norm recurrence and growing-window scale}
Propositions~\ref{prop:v119-prime-essential}--\ref{prop:v119-prime-scale}
are analytic. The archived proof and adverse review distinguish a
finite-codimension norm obstruction from the still-open signed estimate.
The script \texttt{check\_prime\_norm.py} checks the positive trial
$\phi=e^{-x/2}+e^{-(L-x)/2}$ using
\[
 \|\phi\|_2^2=2(1-\lambda^{-2})+2L/\lambda,
 \qquad
 C_\phi(s)=2e^{-s/2}(1-e^{-(L-s)})
             +2(L-s)\lambda^{-1}\cosh(s/2).
\]
Thus its prime Rayleigh quotient is
$2\sum_m w_m C_\phi(\log m)/\|\phi\|_2^2$.
Independent piecewise physical quadrature at $\lambda=2,3$ agrees
with the autocorrelation formula to below $10^{-50}$ at 60-decimal
precision. At $\lambda=300$ the trial quotient divided by $\lambda$
is approximately $1.026889$; this is illustrative, not a norm certificate
or an asymptotic proof.

The same script tests phase recurrence at $\lambda=2$. Modulation by
Fourier index $381074$ has normalized $P_{64}$ head norm about
$9.52\cdot10^{-13}$, while its prime Rayleigh quotient differs from
the unmodulated value by about $4.7\cdot10^{-14}$.
These floating-point checks illustrate the analytic mechanism only.
They give no negative complete-Weil direction, because the
archimedean logarithmic energy grows with the modulation frequency.

\subsection*{A signed energy certificate at $\lambda=4$}
The directory \texttt{g2\_weighted\_signed} in the cumulative bundle
contains the proof, internal adverse review, exact dyadic witnesses
and interval reports for Proposition~\ref{prop:v120-lambda4-tail}.
\texttt{assembly\_general.py} encloses the actual digamma/trigamma
series with explicit exponential remainders. The separate pole diagonal
is retained. \texttt{certify\_weighted\_tail.py} evaluates the
verified-solve lower matrix, including every row up to $4096$ and the
infinite moment remainder, and verifies every Gershgorin row after an
exact dyadic congruence. The saved solve and congruence witnesses are
reused at higher precision; their origin is not a proof assumption.

\texttt{certify\_weighted\_counterexample.py} verifies the finite
dyadic positive direction disproving the two-sided norm contraction.
\texttt{weighted\_pilot.py} is only a floating-point exploration:
its small negative values at larger windows are unresolved numerical
roundoff, not negative Weil directions. The analytic arguments for
weighted compactness, non-Schatten behavior, fixed-window norm decay
and ordinary-error scaling do not rely on these computations.
The primary-source comparison with \cite{Zhu2026} is contextual;
no external certificate or conjectured decay law is imported.

\subsection*{Structured low-mode inverse bounds and their scope}
The low-mode investigation retains the exact $\lambda=4$ matrix,
$E_{16}$ and the previously certified complete tail. The scripts
\texttt{certify\_low\_schur.py} and
\texttt{structured\_low\_schur.py} evaluate the coarse and structured
inverse bounds, respectively. All original arithmetic diagonals,
normalized parity factors and infinite residual moment bounds remain.
The saved tail witnesses are used as exact dyadics; their hashes and
previously certified congruence margins are checked before reuse.
The outer residual belongs to $\mathsf W_4$, whereas the reused inner
factors belong to $\mathsf W_4-10^{-8}D$. Inverse order justifies
that distinction.

These earlier experiments did not certify the low Schur sign. Failed
lower enclosures and their exact diagnostic witnesses are kept in the
research log, not asserted as negative Weil directions. The script \texttt{certify\_far\_inverse.py} verifies the
strict far-inverse constants at 320 and 384 bits; the proof includes
the upper archimedean diagonal bound and the exact pole scale.
The analytic
results of Propositions~\ref{prop:v120-structured-inverse} and
\ref{prop:v120-structured-uniform} are proved directly above and
do not rely on the success of a numerical positivity gate.



\subsection*{Signed-block test: original evidence not archived}
\label{app:v128-test}
The original \texttt{g2\_block\_metric} directory is \emph{not archived}.
Its frozen protocol, numerical proposals, exact dyadic pair and congruence
witnesses, $320/448$-bit outputs, \nolinkurl{certify_block_failure.py}
and independent signed-index verifier are absent. The manuscript preserves
only their historical description: $N=\lambda^2=25$, blocks $26{:}50$,
$51{:}100$, $101{:}200$, $201{:}400$, and actual signed block metrics.

The historical account reports positive-block Gershgorin gates and six
normalized coupling lower bounds per parity, followed by integer
Rayleigh checks. The displayed integer arithmetic is reproducible from
the text; its asserted relation to the actual blocks is not replayable
without the missing witnesses and verifiers. These are reported gates,
not currently available certificates. If those gates hold, additional
nonnegative norm terms cannot repair the finite comparison obstruction.
No complete $\lambda=5$ sign certificate or positivity margin follows.

\subsection*{Rigid $\lambda=5$ transfer and the stopped certificate}
\label{app:v127-transfer}
The original \texttt{g2\_lambda5\_transfer} directory is \emph{not
archived}: the protocol, dyadic counterwitnesses, $320/448$-bit outputs
and obstruction reports are absent. What follows is surviving historical
metadata, not an available certificate. The parameter record was reported
to precede the tests and to have SHA256
\begin{quote}\small\nolinkurl{6de75f0e5121fb4a452de7f0518418820c2acb113b9c02915c450f42875dff1e}.\end{quote}
Its tail start is $16$, $c=10^{-8}$, with even $(N,M,J,r)=(512,1024,4096,48)$
and odd $(1536,2048,4096,64)$, remote Young parameter $10^{-6}$,
$96$ archimedean series terms and target $\mu=0.9999999999$.
The proposed inner solve is the same finite middle solve, dyadic
freezing, inverse-Cholesky proposal and outward Gershgorin verification.

The outer protocol keeps even head $0{:}16$, odd head $1{:}16$, seed
support $256$, inverse-Cholesky trial-energy scaling, support $4096$,
$16$ CG steps per column at $512$ candidate bits and the freshly
certified triangular tail preconditioner. It retains $J=65536$, moment
order $64$, even Young parameter $1/10$, odd complementary parameter
$1/20$ and threshold $1/1000$. The odd directional rule freezes a
proposed top generalized direction and uses support $8192$, $24$ steps,
and the same square-root optimization of its scalar Young parameter.
Those outer steps were not executed at $\lambda=5$, because the
stipulated tail prerequisite was reported to fail.

The absent counterwitness files have historically recorded hashes
\begin{quote}\small
Even: \nolinkurl{747e401533af131a98d5ccac241c5145c87d34245efaf3b8f1d6015f18926976}.\\
Odd: \nolinkurl{105c49f88128480f8d471e614ba91da589d61d87169d692e0a49c55ce2deaa6b}.
\end{quote}
\nolinkurl{falsify_fixed_tail_metric.py} and
\nolinkurl{audit_l5_tail_metric_expanded.py} are \emph{not archived}.
The historical account describes complete finite-support quadratic
replay and an independent signed-index assembly using the same analytic
coefficient enclosures and Arb arithmetic. The asserted agreement is
not currently replayable; it was a parity/sign check, not an independent
analytic kernel derivation. None of the recorded hashes above has been
revalidated against its absent original in NS-38.
No remote enclosure is necessary for the finite-support negative
shifted quadratics. No new $\lambda=8$ computation is part of this run.

The manuscript's additional no-go results are analytic proofs; they
are not inferred from the certificate table. The full first-slot-linear
convention, diagonal, physical endpoint and sampler graph defect are
unchanged. Neither the stopped $\lambda=5$ construction nor the block
metric criterion supplies growing-window G2 or RH.

\subsection*{Odd complement and complete fixed-window sign}
\label{app:v126-complement}
The bundle adds \texttt{g2\_odd\_complement}.  The script
\texttt{direction\_complement\_certificate.py} binds the complete
trial ingredients to their exact dyadic witness, forms the exact
$K$-dependent complement projection with outward arithmetic, and
checks all $15$ complement pivots.  The script
\texttt{quantify\_head\_error.py} tests the ordinary-energy target
using the actual physical head $V$, rather than discarding its
congruence factor.  The complete residual ingredients come from
\texttt{simultaneous\_trial.py}; they include every row beyond
$J=65536$ through the moment majorant.  Their coordinate pivots
are not ordinary spectral gaps.

The script \texttt{directional\_tail\_repair.py} reconstructs the
same frozen $8192$-support directional trial at $1024$ and $1280$
bits and bounds its complete inverse correction. The success
certificate combines the exact rational thresholds $429/1000$
and $1/1000$ to give the relative matrix margin $428/1000$.
The physical head is checked both by outward dyadic arithmetic
and by a separate exact rational calculation. The additional
complement calculation at $896$ bits reuses the previously
certified $768$-bit ingredient enclosures; it is not represented
as a fresh evaluation of those ingredients.

The proof and implementations have been checked adversarially.
Earlier unsuccessful repairs, the intermediate ordinary error
$10^{-43}$, and the limits of these tests are retained in the
research log. No accuracy assumption for a suggested eigenvector
or convergence theorem for conjugate gradients is used. This
internal computer-assisted certificate does not constitute
independent external verification or a growing-window theorem.

\subsection*{Simultaneous even certificate and window scaling}
\label{app:v125-simultaneous}
The cumulative bundle adds \texttt{g2\_simultaneous}. Its main
\texttt{simultaneous\_trial.py} verifier reconstructs the frozen
$17$-column even trial and evaluates all complete residual Grams;
\texttt{complete\_margin.py} certifies the exact rational generalized
margin. Both use the same witness at $768$ and $896$ bits.
The optional \texttt{joint\_remote\_gate.py} retains one common
remote Gram and its cross block. The complete inner tail certificate
and its original exact witnesses are retained unchanged.

The initial odd trials at support $4096$ and the unscaled even
trial are retained as inconclusive experiments in the log. Those
earlier trials are not sign certificates.
All trial heads are checked for full rank; midpoint proposals never
replace outward evaluation in a proof gate. The recorded $LDL^*$
pivots are coordinate pivots, not ordinary spectral gaps.

\texttt{window\_scaling\_probe.py} freshly evaluates the analytic
far-cut thresholds at $320$ and $384$ bits. Its separate finite
Galerkin study uses windows $3,4,5,6,8$, heads through $\lambda^2$
and fixed retained/test cuts $256/384$. It records the finite
Schur correction, dimensionless margins and conditioning. Modes
beyond $384$ are absent from that diagnostic; it is not an execution
of the complete-tail pipeline at the larger windows. Every finite
tail used in its elimination must pass its own positive interval
gate. Complete window-specific inner certificates remain necessary
before interpreting such measurements as bounds for G2.

\subsection*{Complete verification of a new finite trial}
\label{app:v124-posterior}
The cumulative bundle adds \texttt{g2\_posterior\_trial}, containing
\nolinkurl{posterior_trial.py}, its imported helper scripts, the
saved remote-Gram reports, the frozen dyadic corrections and the
independent internal proof/code review. The main certificate uses
\nolinkurl{posterior_witness_M4096_s16.json.gz}; its $4080$ dyadic
pairs define the correction on indices $17$ through $4096$.
The original head vector and support-$256$ solve are unchanged.

The proposal uses $16$ steps of a finite preconditioned conjugate
gradient calculation at $256$ bits. This is only a way of choosing
coefficients. The separate verification routine reconstructs the
exact frozen vector, checks its dimensions and witness hashes,
and evaluates the unshifted complete residual at $768$ or $896$
bits. Its final gate includes the full moment remainder beyond
$J=65536$ and the old complete inner metric. It also checks nine
rows by independent dense assembly and checks the expanded energy
identity. No convergence theorem for the proposing iteration is
needed. A separate four-step trial also passes with lower bound
$2.9028\cdot10^{-20}$; it is not asserted to be the old degree-$4$
inverse polynomial. All trial histories are retained in the log.

\subsection*{Complete mixed-correlation certificate and head refinement}
\label{app:v123-mixed}
The cumulative bundle retains the exact outer and inner witnesses
from the preceding revision. The new files in
\texttt{g2\_mixed\_complete} are:
\begin{itemize}[nosep]
\item \nolinkurl{mixed_refinement.py}, for retained mixed vectors,
complete first-polynomial upper enclosures, and explicitly separate
higher-degree finite-prefix diagnostics;
\item \nolinkurl{mixed_directional_check.py}, for the complete
directional lower certificate of
Proposition~\ref{prop:v123-first-majorant-fails};
\item \nolinkurl{mixed_directional_witness_J65536.json}, specifying
the exact $496$-component probe by integer numerators over $2^{160}$;
\item the $768$- and $896$-bit interval reports and the independent
internal adverse review of the mixed, monotone-head and probe bounds.
\end{itemize}
The inner witness SHA-256 is checked against the saved positive-tail
report. The prefix inverse and its head pairings use the actual
unshifted residual and consistently shifted inner factors. The
certified probe norm is retained rather than rounded to one.
Both remote input and remote first-polynomial output enter the
error estimate. Replaying the analytic and interval inequality is
the sign certificate; a small floating residual is not.

The higher-polynomial experiments are retained in the research log.
Their positive finite-prefix margins do not certify the corresponding
complete corrections and are not used as propositions in this paper.
The strengthened head bound is proved analytically. The complete
$\lambda=4$ sign is established by the later finite-trial and
complement certificates; this does not certify those separate
higher-polynomial experiments.

\subsection*{Parity inverse refinement and its directional gate}
\label{app:v122-inverse}
The cumulative bundle's \texttt{g2\_inverse\_refinement} directory
contains \texttt{refine\_inverse.py}, both parity-constant interval
reports, both final directional-gate reports and the internal adverse
review. The constants replay at $320$ and $384$ bits; the directional
gate replays at $768$ and $896$ bits. The mathematical arguments of
Propositions~\ref{prop:v122-prefix-refinement} and
\ref{prop:v122-refined-structured} are independent of a numerical
positivity gate.

The exact direction uses the $17$ numerators and denominator $2^{160}$
from the archived \texttt{g2\_low\_schur} directory:
\begin{itemize}[nosep]
\item \nolinkurl{structured_ceiling_counterwitness.json} defines $v$;
\item \nolinkurl{low_witness_even_M256_b768.json.gz} defines the dyadic $X$.
\end{itemize}
These archived files define the inputs, not a new rounded eigenvector.
The actual matrix assembly retains its analytic archimedean-series
remainder and correct logarithmic diagonal.

For consecutive positive indices, four Arb polynomial products
evaluate the Toeplitz divided difference and the Hankel matrix action.
The first uses the convolution kernel $1/j$ for $j\ne0$, with zero
central coefficient, applied to $z_n$ and $b_nz_n$. The second
uses reversed input coefficients and the kernel $1/(n+m)$.
Thus the output is exactly
\[
 ((U_Jz)_n)=(d_n-ca_n)z_n+
 \sum_{m\ne n}\frac{b_n-b_m}{n-m}z_m+
 \sum_m\frac{b_n+b_m}{n+m}z_m
\]
in the even sector, including the opposite-sign contribution on its
parity diagonal. The polynomial products retain outward ball errors;
a dense $25$-row evaluation independently checks their indexing.
Only the computational method changes. The tail beyond $J=65536$
is still included by the complete moment bound and
\eqref{eq:v122-remote-H}. Earlier inconclusive cutoffs are retained
in the research log. No full mixed-term, matrix-head or G2 sign
certificate is asserted by this directional test.

\subsection*{Finite trial positivity and exact-source cancellation audit}
\label{app:v121-audit}
The cumulative bundle's \texttt{g2\_schur\_cancellation} directory
contains the exact witnesses, interval outputs and internal adverse
review for Propositions~\ref{prop:v121-alpha-zero} and
\ref{prop:v121-complement-cancellation}. These use the same actual
matrix assembly as the earlier certificates, with all physical scale
and parity factors retained. No externally reported F-68 or F-69
certificate is imported as a premise.

\texttt{certify\_finite\_K.py} reconstructs $K_X$ from its definition
for $\lambda=4$, split $16$ and solve support through $256$.
The files \texttt{finite\_K\_X\_even.json.gz} and
\texttt{finite\_K\_X\_odd.json.gz} contain the frozen exact dyadics.
Every interval $LDL^*$ pivot is positive in both the $768$- and
$896$-bit replays. Conservative lower bounds for the smallest pivot
are listed below; they are \emph{not} eigenvalue lower bounds.
\begin{center}
\begin{tabular}{@{}lll@{}}
\toprule Sector & Trial & Smallest pivot exceeds\\\midrule
Even & $X=0$ & $1.4513\cdot10^{-13}$\\
Odd & $X=0$ & $8.4162\cdot10^{-13}$\\
Even & Frozen solve & $1.2207\cdot10^{-58}$\\
Odd & Frozen solve & $1.3872\cdot10^{-56}$\\\bottomrule
\end{tabular}
\end{center}
This certifies $K_X$, not $S_\infty$ or the omitted inverse-residual
subtraction. The local $\alpha=0$ conclusion therefore does not
contradict previously failed lower Schur enclosures.

\nolinkurl{check_cancellation_subspace.py} uses the exact columns in
\nolinkurl{cancellation_two_mode_witness.json}, each with denominator
$2^{300}$. The $2\times2$ Gram positivity, nonzero constructed vector,
actual quadratic values and all matrix row-sum bounds are verified by
outward ball arithmetic. The source candidate's hash is checked
against the archived \texttt{pswf\_source\_certificate.json}; the
existing exact-source certificate is an explicit dependency.
Projection to the exact source complement is justified analytically,
as in the proof, rather than by a rounded orthogonality residual.

For clarity, the approximate intermediate energies are
\begin{center}
\begin{tabular}{@{}lrr@{}}
\toprule Contribution & Source candidate & Orthogonal trial $w$\\\midrule
Pure archimedean & $-1.4324413624$ & $-0.4195793594$\\
Pole & $+1.4877438175$ & $+0.8369031184$\\
Signed prime & $-0.0553024551$ & $-0.4173237590$\\
Archimedean plus pole & $+0.0553024551$ & $+0.4173237590$\\
Complete Weil & $5.3123817\cdot10^{-38}$ & $2.5844134\cdot10^{-31}$\\
\bottomrule
\end{tabular}
\end{center}
The displayed rounded components cannot resolve the tiny totals by
subtraction; the separate full-precision interval evaluations establish
the strict inequalities. Both precision replays use identical witness
hashes. The prior source identification and the finite-support equality
then supply the exact-source statement in the main text.

These are internal computer-assisted proofs with reproducible inputs,
not external peer verification. The witness proposals are numerical;
the stated sign gates and perturbation margins are interval-certified.
No growing-window convergence or new G2 sign conclusion follows.

\subsection*{Concentration commutator reproduction and scope}
The original \texttt{g2\_nested\_commutator} directory, adverse review,
\texttt{certify\_scaled\_commutator.py} and machine-readable certificate
are \emph{not archived}. The historical account reports exact rational
integration of a degree-four Taylor kernel and a stronger lower check
$\norm{[C_{1/10},C_{1/5}]}>1.198\cdot10^{-5}$. That check is not
presently replayable and is not used here as an available bound.
The complete main proof instead uses its printed degree-two estimate
exceeding $9\cdot10^{-6}$. Its passage to the actual source is analytic,
with Lemma~\ref{lem:v115-source-mass} explicit. This scoped obstruction
to automatic approximate commutation does not require the missing
stronger check and does not close the arithmetic hierarchy.

\subsection*{Disjoint-channel and primitive-transport reproduction}
The original \path{g2_primitive_transport} directory, its two adverse
reviews, \path{test_primitive_transport.py} and saved output are
\emph{not archived}. The disjoint-channel lower bound in
Proposition~\ref{prop:v134-disjoint}, exact drawdown optimization and
Dirichlet-model value $\pi^3/6$ are proved analytically in the manuscript.
They do not use a numerical singular-value estimate or the missing review.
The historical account reports finite-grid drawdown diagnostics at
$\lambda=3,4,5,6$, not interval certificates; the outputs are unavailable
and are not proof inputs. No $\lambda=8$ run was reported. The exact
joint-channel regrouping identity is likewise printed in the manuscript.

\section{Integration ledger for the complete previous manuscript}\label{app:ledger}
The table accounts for all 72 numbered theorems, lemmas, propositions and
corollaries in the supplied 89-page r11 source. Numbers in the first column
are historical identifiers; theorem numbering in this version is new.
A retained local implication does not certify an unproved upstream input.
The supplied audits were treated as arguments to check, not as authority.

All original subject areas are represented: Sections 1--19 are revised in
place, the original Section 20 is rebuilt with a proved compact-source passage
and restored older matrix/weighted estimates, and the original Section 21
is replaced by the consistent sampler and the rebuilt downstream analysis.
The original Sections 22--23 are replaced by the criterion and goals in
Section~\ref{sec:goals}. Invalid proofs are withdrawn instead of being
preserved as established results in an appendix.

\small
\begin{longtable}{@{}p{0.08\textwidth}p{0.85\textwidth}@{}}
\toprule Old & Disposition in v1.28 \\
\midrule\endfirsthead
\toprule Old & Disposition in v1.28 (continued) \\
\midrule\endhead
\bottomrule\endfoot
4.1 & \textbf{Pass.} Three-piece arithmetic source decomposition. Support covariance and a partition of unity identify the closure of the three-source sum; no closed-range theorem is asserted. \\[4pt]
4.2 & \textbf{Pass.} Arithmetic projection disappears in two adjoint channels. The three orthogonality channels remove the arithmetic projections exactly. \\[4pt]
6.1 & \textbf{Pass.} Finite-cutoff span positivity would imply RH. Full-span positivity is sufficient but stronger than the diagonal criterion; the research target is corrected to each evaluator. \\[4pt]
12.1 & \textbf{Repaired.} Exact multiplicity-$m$ leakage factorization. Off-line hypothesis restored; compact division and finite-jet density prove the missing kernel inclusion. Quantitative persistence still uses the evaluator assumption. \\[4pt]
13.1 & \textbf{Pass.} No generic domination between cross-scale and fixed-space defects. The finite-dimensional examples establish that generic projection geometry gives no defect domination. \\[4pt]
14.1 & \textbf{Repaired.} Exact generator range defect. Compact division and finite-jet density supply a proof. The range is the algebraic source-core power before closure. \\[4pt]
15.1 & \textbf{Pass.} Exact zero-incidence criterion. Pole removability proves zero incidence under the stated ambient Sonine division framework. The skew-extension sign is corrected nearby. \\[4pt]
17.1 & \textbf{Pass.} Exact diagonal contractions force orthogonality. Dense unimodular eigenvectors of a contraction force a unitary and phase orthogonality. \\[4pt]
17.2 & \textbf{Conditional.} Finite-mode asymptotic calibration. Retained under the explicit quantitative evaluator assumption. \\[4pt]
18.1 & \textbf{Conditional.} Exact off-line spectrum of the fixed arithmetic operator. Retained under an explicit closed Hilbert realization, domain, graph-core and spectral-completeness assumption; the construction is not supplied. \\[4pt]
19.1 & \textbf{Retained; finite-core alternative added.} The Riesz-window implication retains its realization hypothesis. Theorem~\ref{thm:v15-finite-metric} bypasses it on algebraic root spaces. The full-space projection now preserves coercivity and Weil transfer while annihilating source and shell; a small signed defect is still missing, and a fixed scalar Weil shift fails. \\[4pt]
20.1 & \textbf{Pass.} Endpoint-normalized radial bound. The fixed-order radial estimate is retained with its cited asymptotic input. \\[4pt]
20.2 & \textbf{Pass.} Away-boundary and exterior-energy bounds. Away-boundary and exterior-energy estimates are retained for each fixed radial order. \\[4pt]
20.3 & \textbf{Pass.} Uniform far field with controlled $c$-dependence. The radial ODE and integrable far-field perturbation argument are retained. \\[4pt]
20.4 & \textbf{Pass.} Endpoint-normalized radial lattice bound. The lattice summation argument and resonance control are retained. \\[4pt]
20.5 & \textbf{Pass for the explicit repaired source.} Fixed-order co-Poisson cross-tail. The normalized source, repaired tail, inversion-endpoint comparison and correlation estimates are proved. Theorem~\ref{thm:v14-radical} now supplies the global radical/semilocal-domain passage; no source assumption remains. \\[4pt]
20.6 & \textbf{Pass.} Quantitative endpoint-jump bound. The normalized source-tail bounds and the PNT remainder prove the rate; both translated-tail terms are retained. \\[4pt]
20.7 & \textbf{Pass.} Inversion symmetrization preserves the G1 envelope. Inversion preserves the chosen test class and transfers the proved envelope. \\[4pt]
20.8 & \textbf{Pass for Sobolev pairings.} Theorem~\ref{thm:v14-radical} closes the remaining domain input. This does not supply a growing-block operator-norm estimate. \\[4pt]
20.9 & \textbf{Pass.} Arithmetic BV structure and absence of a periodic boundary jump. The finite arithmetic sum has BV regularity and matching periodic endpoint representatives. \\[4pt]
20.10 & \textbf{Pass.} Fixed-cutoff boundary-preserving Galerkin approximation. Fixed-cutoff BV Fourier/form approximation is retained, with the logarithmic form weight and endpoint convergence distinguished. \\[4pt]
20.11 & \textbf{Pass.} Quantitative and cofinal finite-section transfer of G1. The finite-band approximation transfers the proved continuum estimate; its eventual threshold need not be economical. \\[4pt]
20.12 & \textbf{Pass.} Zero-independent high--low decay of the Weil matrix. The off-diagonal matrix estimate is retained and supplemented by the separately restored diagonal bound. \\[4pt]
20.13 & \textbf{Pass.} Explicit BV rate for the fixed-band Galerkin defect. BV coefficients and the off-diagonal estimate give the summable high-low tail. \\[4pt]
20.14 & \textbf{Pass.} Quantitative endpoint recovery by de la Vall\'ee--Poussin filtering. The near/far endpoint convolution estimate is retained. \\[4pt]
20.15 & \textbf{Pass.} Filtered BV transfer with explicit boundary control. The filtered transfer uses the now-proved continuum G1; zero local Lipschitz constant is handled without division. \\[4pt]
20.16 & \textbf{Pass.} High-frequency boundary repair has unavoidable square-root derivative cost. Weighted Cauchy--Schwarz gives the derivative cost; a general trace obstruction is added for the new sampler. \\[4pt]
20.17 & \textbf{Pass in the fixed-band sense.} Constructive zero-independent finite G1 closure. The exact-boundary construction transfers the now-proved continuum estimate with the stated deterministic cutoffs. No economical cutoff or growing-block uniformity is inferred. \\[4pt]
20.18 & \textbf{Repaired.} Form-Schur lower bound. The Schur error parameters are nonnegative. Rayleigh calibration and operator-domain overlap are separated from weak G1. \\[4pt]
21.1 & \textbf{Retained.} One-sided off-line Weil isotropy. Global one-sided isotropy is proved in the corrected convention within the centered transfer theorem. \\[4pt]
21.2 & \textbf{Replaced.} Off-line semilocal form transfer. The full-scale raw transfer is withdrawn. Centered strip transfer retains jumps; the explicit endpoint repair preserves its form limit. \\[4pt]
21.3 & \textbf{Rebuilt.} G2b under the off-line contradiction hypothesis. Restricted coercivity follows for the corrected sampler. A negative compact-test calibration is optional and needs its own form approximation. \\[4pt]
21.4 & \textbf{Rebuilt.} Restricted scalar-shift coercivity is automatic after form transfer. Positive scalar-shift coercivity follows from corrected form and ordinary Gram transfer; it is not independent positivity. \\[4pt]
21.5 & \textbf{Repaired.} Smooth-test continuity of the semilocal Weil pairing. The logarithmic diagonal is restored. Absolute matrix-series estimates and the weighted convolution proof replace the false uniform-diagonal bound. \\[4pt]
21.6 & \textbf{Rebuilt.} The proved G1 estimate applies to the full Hardy sampler. The restored weighted theorem gives weak G1 for centered Hardy samples using the now-proved source passage; the even endpoint correction is invisible to the zeroth pairing by parity. \\[4pt]
21.7 & \textbf{Scoped.} Finite G1 jet hierarchy on the full Hardy sampler. Sufficient smoothing controls each fixed finite derivative family. No growing-order uniformity or automatic parity preservation for all corrected jets is asserted. \\[4pt]
21.8 & \textbf{Retained.} Fixed residual-jet Gramians cannot supply G2 coercivity. A fixed finite Gramian of vanishing normalized residual pairings also vanishes; the fixed-family G1 input is now proved. \\[4pt]
21.9 & \textbf{Rebuilt.} Off-line coercivity forces a boundary-weighted sampler residual. The exact Lyapunov residual now includes the corrected graph term. Its coercive lower bound is retained. \\[4pt]
21.10 & \textbf{Rebuilt.} Exact sampler--prolate compatibility decomposition. The finite candidate decomposition uses the total source e z plus the graph defect. \\[4pt]
21.11 & \textbf{Rebuilt.} Commuting-correction cancellation and the scalar-shift no-go. Commuting and scalar cancellation remain exact with the total source. The old raw-transfer asymptotic is removed. \\[4pt]
21.12 & \textbf{Rebuilt.} Exact noncommuting G2 excess. The noncommuting excess includes the graph defect and still requires an independent upper estimate. \\[4pt]
21.13 & \textbf{Retained.} $W^2$ benchmark: positivity moves the gap to norm transfer. The positive square and rank-four commutator are exact; form smallness does not imply norm smallness. \\[4pt]
21.14 & \textbf{Repaired.} Reflected-partner $W^2$ coercivity at arbitrary multiplicity. The reflected top-root lower bound is proved for the corrected sampler; conjugation is on the second factor. This is not full-block coercivity. \\[4pt]
21.15 & \textbf{Rebuilt.} A boundary-annihilating G2 test isolates the exact sampler--prolate mismatch. Boundary annihilation is an exact finite construction when its denominator is nonzero. The proposed bounded endpoint ratio is now ruled out by Corollary~\ref{cor:v11-ratio}; bounded tests from a fixed input space lose the reflected signal by Theorem~\ref{thm:v11-fixed-tests}. \\[4pt]
21.16 & \textbf{Withdrawn.} Pole-aware radical transfer for the fixed Hardy sampler. The pole-aware transfer proof is insufficient. A replacement must bound paired errors over all zeros, including slow tails and actual cut boundaries. \\[4pt]
21.17 & \textbf{Conditional only.} The surviving simple-zero obstruction is a boundary commutator channel. The finite commutator decomposition survives; no nonzero beta limit is claimed. \\[4pt]
21.18 & \textbf{Rebuilt.} Multiplicity-robust reflected-partner block obstruction. The corrected top-Jordan relation contains an extra graph term. The cross-form lower bound survives. The original endpoint-annihilating test cannot retain it with bounded norm; a different test construction and graph estimates are required. \\[4pt]
21.20 & \textbf{Scoped.} Rank--two boundary traces exhaust the reflected G2 signal. Rank-two trace algebra survives, but it does not exhaust the signal without controlling the new graph pairing. \\[4pt]
21.21 & \textbf{Withdrawn.} Critical $L^{-1/2}$ scale of the surviving boundary trace. The critical nonzero beta limit is removed from the established chain. One former premise, a bounded ratio for the corrected resolvent, is disproved; the other graph and form obligations remain explicit. \\[4pt]
21.22 & \textbf{Retained.} Exact arithmetic realization of the surviving CCM boundary channel. The exact finite arithmetic identity is retained with physical scaling, integer index domain, and the first-slot convention. \\[4pt]
21.23 & \textbf{Conditional only.} G2 is an explicit finite prime/pole/archimedean forcing statement. The finite observable is explicit. Its asserted nonzero asymptotic is not established for the corrected test family. \\[4pt]
21.24 & \textbf{Scoped.} The present absolute PNT remainder is quantitatively insufficient for the G2 boundary channel. The absolute PNT estimate is retained as an upper-bound limitation, not a lower bound or impossibility theorem. \\[4pt]
21.26 & \textbf{Retained.} Hardy phase shifts preserve the global G2 signal while damping the Hardy endpoint. Global translation, Cauchy damping, and reflected cross-form invariance are retained in the fixed Fourier convention. \\[4pt]
21.27 & \textbf{Withdrawn.} A translated Hardy kernel recovers an absolute PNT gain at the continuum-kernel level. The periodic upper bound omitted a reflected image. The missing image and an explicit exponent comparison are now displayed. \\[4pt]
21.28 & \textbf{Scoped.} Reflected-pole alias obstruction to the moving Hardy phase. Threshold algebra is retained. A no-go conclusion requires relative alias remainder estimates that are not supplied. \\[4pt]
21.30 & \textbf{Retained.} Exact continuum normal form of the boundary-annihilated G2 test. The continuum meromorphic normal form and its residue are retained; it is not asserted to be the corrected finite endpoint-ratio limit. \\[4pt]
21.31 & \textbf{Repaired.} The reflected pole is forced by any cross-retaining holomorphic correction. Residue persistence is retained and second-slot scaling uses the conjugated scalar. \\[4pt]
21.32 & \textbf{Repaired.} Bounded Hardy filtering cannot suppress the reflected pole while retaining G2. The modulus inequality is proved; global admissibility and a nonzero endpoint factor are required for form interpretations. \\[4pt]
21.33 & \textbf{Scoped.} Absolute-PNT rescue by pole damping lies outside the audited G1 regime. Exponential norm cost describes this envelope-based estimate, not a necessary lower bound on actual transfer error. \\[4pt]
21.35 & \textbf{Retained.} Projective boundary-quotient invariance of the G1/G2 mismatch. Projective invariance is exact for a fixed endpoint-zero test. Its asymptotic consequence assumes both normalized G1 pairings vanish. \\[4pt]
21.36 & \textbf{Scoped.} The prolate core may be replaced by a pure boundary shell in the lower obstruction. A shell can replace the candidate in the exact paired identity; no shared nonzero G2 limit is imported. \\[4pt]
21.37 & \textbf{Scoped.} Moving Hardy evaluation aligns the graph functional. The moving Cauchy scalar converges to the graph functional on a fixed root space; moving-parameter uniform transfer remains open. \\[4pt]
21.38 & \textbf{Retained.} Candidate-adapted graph sampler: exact vector matching. Coefficientwise exact graph matching and scalar endpoint error are proved. Injectivity, norm, form and G1 compatibility are additional obligations. \\[4pt]
21.39 & \textbf{Retained.} Exact-radical factorization removes the zeta quotient but exposes an endpoint phase. Global Mellin divisibility and the endpoint/bulk integration-by-parts identity remain, with the Fourier-sign conversion explicit. \\[4pt]
21.40 & \textbf{Scoped.} Finite boundary-jet expansion of the exact matcher. The fixed-cutoff high-frequency jet expansion remains. Subtraction through zero frequency requires a prescription. \\[4pt]
21.42 & \textbf{Retained.} Exact periodic-resolvent formula for the adapted endpoint. The periodic ODE and parameter derivatives give the exact normalized Laplace-jet formula. \\[4pt]
21.43 & \textbf{Retained.} Boundary-perfect high-frequency matching has vanishing bulk mass. The shell exact matcher has asymptotically correct endpoints but vanishing ordinary mass, so it is not coercive. \\[4pt]
21.45 & \textbf{Rebuilt.} Parity projection gives a bounded noncommuting G2 benchmark. Parity coercivity uses only ordinary norm transfer and survives the corrected sampler. Its Lyapunov identity includes the graph defect. \\[4pt]
21.48 & \textbf{Repaired.} The exact-boundary shell is compatibility-efficient in the coercive pairing. Correct diagonal and row-sum bounds give shell weak estimates. They do not remove the commutator term in the graph pairing. \\[4pt]
21.49 & \textbf{Scoped.} One shell cutoff neutralizes both the G1 repair and its G2 pairing cost. Common-cutoff shell smallness is retained as an eventual-threshold argument; it controls auxiliary shell terms only. \\[4pt]
21.50 & \textbf{Replaced.} A joint exact-boundary G1/pairing diagonal. An explicit normalized common-cutoff selection lemma replaces ambiguous cofinality claims. It proves no missing core graph estimate. \\[4pt]
21.51 & \textbf{Scoped.} Scalar-shift positivity cannot remove the final defect. The reverse-triangle obstruction is retained only with its explicit full-residual and derivative-norm premises; weak G1 does not supply them. \\[4pt]
\end{longtable}
\normalsize
\paragraph{Earlier-version material.}
The separate correct diagonal bound, weighted-sampler convolution estimate,
arithmetic jump train, jump/continuous split, and the limited near-Slepian
jump cutoff were recovered from the older 65-page source. The 32-page v10
was checked for useful earlier structure. Its unsupported exact jet
Gram asymptotic, incomplete far-field comparison, and unseparated endpoint
jumps were not restored. The more careful one-sided Gram bound, radial ODE
argument, and endpoint-jump separation of the later source are retained.

\paragraph{Audit reconciliation.}
The supplied G1/G2 audit overstated the finite lower mechanism. The supplied
Claude audit checked useful local algebra but did not settle the analytic
transfers affected by the cut and diagonal errors. This revision adopts
locally sound calculations, corrects their conventions and scope, and
replaces the central transfer rather than treating either audit as a
certification of the whole paper.

\paragraph{Authorship and assistance.}
AI tools assisted with drafting, exploratory calculations, and adversarial
review. They are not mathematical authorities or sources of certification.
Responsibility for the manuscript's claims and any eventual publication
remains with its author. The reproducible finite checks are supplied separately. Exploratory
computations are distinguished from the rigorous finite interval
matrix and exact-source certificates in Appendix~\ref{app:numerics}.
The new endpoint and operator arguments received internal adversarial
review, including growing derivative orders, arithmetic threshold
resonance, and identification of the physical form closure. This is
not external referee verification.

\begin{thebibliography}{99}
\small
\setlength{\itemsep}{2pt}
\bibitem{BurnolExplicit2000}
J.-F. Burnol,
\emph{Sur les Formules Explicites I: analyse invariante},
C. R. Acad. Sci. Paris, Serie I, 331 (2000), 423--428;
arXiv:math/0101068.

\bibitem{Zhu2026}
X. Zhu,
\emph{Weil positivity in compact windows: a finite reduction,
certified two-sided bounds, and a Landau--Widom decay law},
arXiv:2608.24827v2 (2026).

\bibitem{Wiener1932}
N. Wiener,
\emph{Tauberian Theorems},
Annals of Mathematics, second series, 33 (1932), 1--100.

\bibitem{Burnol2001}
J.-F. Burnol,
\emph{Sur certains espaces de Hilbert de fonctions entieres, lies a la transformation de Fourier et aux fonctions L de Dirichlet et de Riemann},
arXiv:math/0105120v1 (2001).

\bibitem{Burnol2002Sonine}
J.-F. Burnol,
\emph{Sur les espaces de Sonine associes par de Branges a la transformation de Fourier},
arXiv:math/0208121v1 (2002).

\bibitem{BurnolCoPoisson2004}
J.-F. Burnol,
\emph{Entrelacement de co-Poisson},
arXiv:math/0407443 (2004).

\bibitem{Burnol2004}
J.-F. Burnol,
\emph{Two complete and minimal systems associated with the zeros of the Riemann zeta function},
arXiv:math/0203120v7; J. Theorie des Nombres de Bordeaux 16 (2004), 65--94.

\bibitem{Burnol2011}
J.-F. Burnol,
\emph{On the system of the functions $\zeta(s)/(s-\rho)^k$},
arXiv:1106.4751 (2011).

\bibitem{Livsic2012}
A. Aleman, R. T. W. Martin, and W. T. Ross,
\emph{On a theorem of Livsic},
arXiv:1209.4497 (2012).

\bibitem{ConnesConsaniMoscovici2025}
A. Connes, C. Consani, and H. Moscovici,
\emph{Zeta Spectral Triples},
arXiv:2511.22755 (2025).

\bibitem{ConnesVanSuijlekom2025}
A. Connes and W. D. van Suijlekom,
\emph{Quadratic Forms, Real Zeros and Echoes of the Spectral Action},
Commun. Math. Phys. 406 (2025), article 312.
\url{https://doi.org/10.1007/s00220-025-05493-1};
\url{https://arxiv.org/abs/2511.23257v1}.
Theorem~6.1 and Remark~4.3.

\bibitem{Dunster2017}
T. M. Dunster,
\emph{Asymptotics of Prolate Spheroidal Wave Functions},
arXiv:1601.00699v3 (2017).

\bibitem{Fuchs1964}
W. H. J. Fuchs,
\emph{On the eigenvalues of an integral equation arising in the theory of band-limited signals},
J. Math. Anal. Appl. 9 (1964), 317--330.

\bibitem{BonamiKaroui2014}
A. Bonami and A. Karoui,
\emph{Uniform bounds of prolate spheroidal wave functions and eigenvalues decay},
C. R. Math. Acad. Sci. Paris 352 (2014), 229--234.

\bibitem{Suzuki2026}
M. Suzuki,
\emph{Weil's quadratic form via the screw function},
arXiv:2606.09096v2 (2026).

% NS-17: bib_inserts.tex
\bibitem{Groskin2026}
A. Groskin,
\emph{A finite Guinand--Weil dictionary and archimedean tail order for the
truncated Weil quadratic form},
arXiv:2607.02828v3 (2026).
\url{https://arxiv.org/abs/2607.02828v3}.

\bibitem{KreinLanger2014}
M. G. Kre\u{\i}n and H. Langer,
\emph{Continuation of Hermitian positive definite functions and related
questions},
Integral Equations Operator Theory 78 (2014), 1--69.
\url{https://doi.org/10.1007/s00020-013-2091-z}.

\bibitem{BasorEhrhardt2002}
E. L. Basor and T. Ehrhardt,
\emph{Asymptotic formulas for the determinants of symmetric Toeplitz plus
Hankel matrices},
in: A. B\"ottcher, I. Gohberg, and P. Junghanns (eds.),
\emph{Toeplitz Matrices and Singular Integral Equations},
Operator Theory: Advances and Applications, Birkh\"auser, Basel, 2002, 61--90;
arXiv:math/0202078.
\url{https://doi.org/10.1007/978-3-0348-8199-9_5}.

\bibitem{ConnesConsani2023}
A. Connes and C. Consani,
\emph{Spectral triples and $\zeta$-cycles},
L'Enseignement Math. 69 (2023), 93--148.

\bibitem{SlepianTransition}
S. Karnik, J. Romberg, and M. A. Davenport,
\emph{Improved bounds for the eigenvalues of prolate spheroidal wave functions and discrete prolate spheroidal sequences},
Appl. Comput. Harmon. Anal. 55 (2021), 97--128.

\bibitem{Douglas}
R. G. Douglas,
\emph{On majorization, factorization, and range inclusion of operators on Hilbert space},
Proc. Amer. Math. Soc. 17 (1966), 413--415.

\bibitem{Hardy}
G. H. Hardy, \emph{Sur les z\'eros de la fonction $\zeta(s)$ de Riemann},
C. R. Acad. Sci. Paris 158 (1914), 1012--1014.

\bibitem{Trudgian2012}
T. S. Trudgian,
\emph{An improved upper bound for the argument of the Riemann
zeta-function on the critical line II},
\href{https://arxiv.org/abs/1208.5846v2}{arXiv:1208.5846v2} (2012),
Corollary 1 and equation (2.5).

\bibitem{PythonFlint2026}
The Python-FLINT and FLINT developers,
\emph{Python-FLINT 0.9.0 documentation}, real and complex ball arithmetic,
rigorous integration and matrix solves; used with FLINT 3.6.0.
\url{https://python-flint.readthedocs.io/en/stable/arb_mat.html};
\url{https://python-flint.readthedocs.io/en/stable/acb.html}.
Accessed September 20, 2026.

\bibitem{DLMFPSWF}
NIST Digital Library of Mathematical Functions, Chapter 30,
Spheroidal Wave Functions, Sections 30.2--30.3.
\url{https://dlmf.nist.gov/30.2};
\url{https://dlmf.nist.gov/30.3}.
Regular angular equation, spectral ordering and normalization convention.

\bibitem{NISTDigamma}
NIST Digital Library of Mathematical Functions, Chapter 5,
equations 5.9.15 and 5.15.1 (Binet's digamma formula and trigamma series),
\url{https://dlmf.nist.gov/5.9.E15},
\url{https://dlmf.nist.gov/5.15.E1}.

\bibitem{Teschl2009}
G. Teschl, \emph{Mathematical Methods in Quantum Mechanics},
Graduate Studies in Mathematics 99, American Mathematical Society, 2009,
Section~4.3.
\url{https://www.mat.univie.ac.at/~gerald/ftp/book-schroe/schroe.pdf}.
\bibitem{ChenWethLog}
H.~Chen and T.~Weth,
\emph{The Dirichlet Problem for the Logarithmic Laplacian},
arXiv:1710.03416v6 (2019),
\url{https://arxiv.org/abs/1710.03416v6}.

\bibitem{HLSLogBoundary}
V.~Hern\'andez-Santamar\'ia, L.~F.~L\'opez R\'ios, and A.~Salda\~na,
\emph{Optimal boundary regularity and a Hopf-type lemma for Dirichlet
problems involving the logarithmic Laplacian},
arXiv:2401.18033v2 (2024),
\url{https://arxiv.org/abs/2401.18033v2}.

\bibitem{NISTPrimeAsymptotics}
NIST Digital Library of Mathematical Functions, Section~27.12,
Asymptotic Formulas: Primes, prime number theorem and classical error
estimates. \url{https://dlmf.nist.gov/27.12}.

\bibitem{FrankSeiringer2008}
R. L. Frank and R. Seiringer,
\emph{Non-linear ground state representations and sharp Hardy inequalities},
J. Funct. Anal. 255 (2008), 3407--3430;
arXiv:0803.0503v2.
\url{https://arxiv.org/abs/0803.0503v2}.

\bibitem{ChenWang2013}
X. Chen and J. Wang,
\emph{Functional inequalities for nonlocal Dirichlet forms with finite
range jumps or large jumps},
arXiv:1212.6100v2 (2013).
\url{https://arxiv.org/abs/1212.6100v2}.

\bibitem{LandauPollak1961}
H. J. Landau and H. O. Pollak,
\emph{Prolate spheroidal wave functions, Fourier analysis and
uncertainty---II},
Bell System Tech. J. 40 (1961), 65--84.

\bibitem{Nazarov1994}
F. L. Nazarov,
\emph{Local estimates for exponential polynomials and their
applications to inequalities of the uncertainty principle type},
St. Petersburg Math. J. 5 (1994), 663--717.
\url{https://www.mathnet.ru/eng/aa397}.

\bibitem{Kovrijkine2000}
O. Kovrijkine,
\emph{Some results related to the Logvinenko--Sereda theorem},
arXiv:math/0012186v1 (2000).
\url{https://arxiv.org/abs/math/0012186v1}.

\bibitem{BrownHalmos1964}
A. Brown and P. R. Halmos,
\emph{Algebraic properties of Toeplitz operators},
J. Reine Angew. Math. 213 (1963/64), 89--102.
\url{https://doi.org/10.1515/crll.1964.213.89}.

\bibitem{Virtanen2022}
J. A. Virtanen,
\emph{On the product formula for Toeplitz and related operators},
arXiv:2205.12345v1 (2022).
\url{https://arxiv.org/abs/2205.12345v1}.

\bibitem{Bang1995}
H. H. Bang,
\emph{On the Bernstein--Nikolsky inequality II},
Tokyo J. Math. 18 (1995), 123--131.
\url{https://doi.org/10.3836/tjm/1270043613}.


\bibitem{FanPall1957}
K. Fan and G. Pall, \emph{Imbedding conditions for Hermitian and normal
matrices}, Canadian J. Math. 9 (1957), 298--304.
\url{https://doi.org/10.4153/CJM-1957-036-1}.

\bibitem{LiebLoss2001}
E. H. Lieb and M. Loss, \emph{Analysis}, second edition,
Graduate Studies in Mathematics 14, American Mathematical Society, 2001,
Section~1.14 (bathtub principle).
\url{https://www.ams.org/publications/authors/books/postpub/gsm-14-R}.

\bibitem{HarperZetaNotes}
A.~J. Harper, \emph{Riemann zeta function lecture notes 1},
Lemma~2.4 and Theorem~6.5,
\url{https://warwick.ac.uk/fac/sci/maths/people/staff/harper/zetacoursenotes1.pdf}.

\bibitem{DLMFZetaReflection}
NIST Digital Library of Mathematical Functions,
\S25.4, equations (25.4.3)--(25.4.4),
\url{https://dlmf.nist.gov/25.4}.

\bibitem{HasanalizadeShenWong2022}
E. Hasanalizade, Q. Shen and P.-J. Wong,
\emph{Counting zeros of the Riemann zeta function},
J. Number Theory 235 (2022), 219--241;
\url{https://arxiv.org/abs/2107.06506}.

\bibitem{Ingham1940}
A. E. Ingham, \emph{On the estimation of $N(\sigma,T)$},
Quart. J. Math. Oxford Ser. 11 (1940), 201--202.
\url{https://doi.org/10.1093/qmath/os-11.1.201}.

\bibitem{ChourasiyaSimonic2025}
S. Chourasiya and A. Simoni\v{c},
\emph{An explicit form of Ingham's zero density estimate},
arXiv:2507.15184v2 (2025), Introduction and Corollary 1.
\url{https://arxiv.org/abs/2507.15184v2}.

\bibitem{Aryan2019}
F. Aryan, \emph{On an extension of the Landau--Gonek formula},
arXiv:1902.05473 (2019), Introduction, equation (1.1).
\url{https://arxiv.org/abs/1902.05473}.

\bibitem{Montgomery1973}
H. L. Montgomery, \emph{The pair correlation of zeros of the zeta
function}, Proc. Sympos. Pure Math. 24 (1973), 181--193,
Theorem and conjecture in Section 1.
\url{https://www-personal.umich.edu/~hlm/paircor1.pdf}.

\bibitem{BaluyotEtAl2023}
S. A. C. Baluyot, D. A. Goldston, A. I. Suriajaya and
C. L. Turnage-Butterbaugh,
\emph{An unconditional Montgomery theorem for pair correlation of
zeros of the Riemann zeta function}, Acta Arith. 214 (2024), 357--376,
Theorem 1; \url{https://arxiv.org/abs/2306.04799}.

\bibitem{ChenHauerWeth2023}
H. Chen, D. Hauer, and T. Weth,
\emph{Extension problems for the logarithmic Laplacian},
arXiv:2312.15689v1 (2023), Theorems 1.2, 1.6, and 1.7,
\url{https://arxiv.org/abs/2312.15689v1}.

\bibitem{HarrachLinWeth2025}
B. Harrach, Y.-H. Lin, and T. Weth,
\emph{The Calder\'on problem for the logarithmic Schr\"odinger equation},
Journal of Differential Equations \textbf{444} (2025), 113665,
\url{https://doi.org/10.1016/j.jde.2025.113665}.


\bibitem{BaezDuarteAutocorrelation2003}
L. B\'aez-Duarte, M. Balazard, B. Landreau, and E. Saias,
\emph{Sur l'autocorr\'elation multiplicative de la fonction
``partie fractionnaire''}, arXiv:math/0306251v1 (2003), Proposition 89.
\url{https://arxiv.org/abs/math/0306251}.

\bibitem{HuxleyIvic2007}
M. N. Huxley and A. Ivi\'c,
\emph{Subconvexity for the Riemann zeta-function and the divisor problem},
Bulletin CXXXIV de l'Acad\'emie Serbe des Sciences et des Arts,
Sciences math\'ematiques 32 (2007), 13--32; Section 2, especially Eq. (2.1).
\url{https://arxiv.org/abs/math/0611809v3}.

\bibitem{Soundararajan2008}
K. Soundararajan, \emph{Extreme values of zeta and $L$-functions},
Math. Ann. 342 (2008), 467--486, Theorem 1.
\url{https://arxiv.org/abs/0708.3990v2}.


\bibitem{RosserSchoenfeld1962}
J. B. Rosser and L. Schoenfeld,
\emph{Approximate formulas for some functions of prime numbers},
Illinois J. Math. 6 (1962), 64--94, Corollary 1, Eq. (3.6).
\url{https://doi.org/10.1215/ijm/1255631807}.

\bibitem{BettinConreyFarmer2013}
S. Bettin, J. B. Conrey and D. W. Farmer,
\emph{An optimal choice of Dirichlet polynomials for the Nyman--Beurling criterion},
\url{https://arxiv.org/abs/1211.5191}, Theorem 1 (conditional on RH and an
additional inverse-derivative moment bound).

\bibitem{Ehm2024}
W. Ehm, \emph{On certain Gram matrices and their associated series},
arXiv:2405.06349v1 (2024), Theorem 2.1 and Proposition 5.2.
\url{https://arxiv.org/abs/2405.06349v1}.

\bibitem{BaezDuarte2000Arithmetic}
L. B\'aez-Duarte,
\emph{Arithmetical aspects of Beurling's real variable reformulation
of the Riemann hypothesis}, arXiv:math/0011254v1 (2000), Section 4.
\url{https://arxiv.org/abs/math/0011254v1}.

\bibitem{DLMFBernoulliFourier}
NIST Digital Library of Mathematical Functions,
\S24.8, Eq. (24.8.3), Bernoulli Fourier series.
\url{https://dlmf.nist.gov/24.8.E3}.

\bibitem{DLMFLaguerre}
NIST Digital Library of Mathematical Functions,
\S18.3, Table 18.3.1, Laguerre orthogonality.
\url{https://dlmf.nist.gov/18.3}.

\bibitem{DLMFGammaBounds}
NIST Digital Library of Mathematical Functions,
\S5.4, Eq. (5.4.13), and \S5.11, positive-real Stirling remainders.
\url{https://dlmf.nist.gov/5.4.E13}; \url{https://dlmf.nist.gov/5.11}.

\bibitem{BaezDuarteStrong}
L. B\'aez-Duarte,
\emph{A strengthening of the Nyman--Beurling criterion for the Riemann Hypothesis},
arXiv:math/0202141v2 (2002), especially Lemma 2.1 and Section 3.
\url{https://arxiv.org/abs/math/0202141v2}.


\begingroup\raggedright
\bibitem{BurnolNyman2001}
J.-F. Burnol, \emph{A note on Nyman's equivalent formulation of the
Riemann Hypothesis}, arXiv:math/9910055v3 (2001),
Theorem 1.2, Section 2 and Lemma 2.4.
\url{https://arxiv.org/abs/math/9910055v3}.

\bibitem{HasanalizadeShenWong2021}
E. Hasanalizade, Q. Shen and P.-J. Wong,
\emph{Counting zeros of the Riemann zeta function},
arXiv:2107.06506v1 (2021), Corollary 1.2.
\url{https://arxiv.org/abs/2107.06506v1}.

\bibitem{PlattTrudgian2021}
D. Platt and T. Trudgian,
\emph{The Riemann hypothesis is true up to $3\cdot10^{12}$},
Bull. London Math. Soc. \textbf{53} (2021), 792--797, Theorem 1;
arXiv:2004.09765.
\url{https://doi.org/10.1112/blms.12460}.
\par\endgroup

\end{thebibliography}

\end{document}
